\documentclass[oneside,10pt]{amsart}
\PassOptionsToPackage{hyperfootnotes=false}{hyperref}
\ifdefined\pdfvariable\pdfvariable suppressoptionalinfo 767\fi
\usepackage[no-math]{fontspec}
\setmainfont{Noto Serif}
\usepackage{amsmath,amssymb,amsthm}
\usepackage[all]{xy}
\usepackage{xstring,xparse,xcolor}
\definecolor{egalinkred}{rgb}{0.76,0.13,0.28}
\usepackage[
  linktocpage=true,
  colorlinks=true,
  hyperindex,
  hypertexnames=false,
  citecolor=blue,
  linkcolor=egalinkred
]{hyperref}
\usepackage[left=30mm,right=30mm,top=22mm,bottom=25mm]{geometry}
\usepackage{marginnote}
\usepackage[shortlabels]{enumitem}
\usepackage{longtable,multicol,microtype}
\usepackage[spanish,es-nodecimaldot]{babel}
\setlength{\emergencystretch}{5em}
\raggedbottom

\hypersetup{
  pdftitle={Elementos de geometría algebraica, EGA II — edición española},
  pdfauthor={Alexandre Grothendieck y Jean Dieudonné},
  pdfsubject={Estudio global elemental de algunas clases de morfismos}
}

% Statements. EGA numbers its statements in the text, so the environments
% are unnumbered and take the number as optional argument. For instance
% \begin{proposition}[2.1.3] starts with the bold heading "Proposition
% (2.1.3). ---" (in the language of the edition) and sets the statement in
% italics; definitions, examples, remarks and notation are set in roman;
% \begin{env}[2.1.4] starts a numbered paragraph with a bold "(2.1.4).".
\newtheoremstyle{egaitem}{}{}{\rmfamily}{}{\bfseries}{.}{ }{\thmnote{(#3)}}
\newtheoremstyle{egastatement}{}{}{\itshape}{}{\bfseries}{. --- }{ }{\thmname{#1}\thmnote{ (#3)}}
\newtheoremstyle{egadescription}{}{}{\rmfamily}{}{\bfseries}{. --- }{ }{\thmname{#1}\thmnote{ (#3)}}
\makeatletter
\newcommand*{\ega@statement}[3]{\theoremstyle{#1}\newtheorem*{#2}{#3}}
\ega@statement{egaitem}{env}{---}
\ega@statement{egastatement}{theorem}{Teorema}
\ega@statement{egastatement}{proposition}{Proposición}
\ega@statement{egastatement}{lemma}{Lema}
\ega@statement{egastatement}{corollary}{Corolario}
\ega@statement{egadescription}{definition}{Definición}
\ega@statement{egadescription}{example}{Ejemplo}
\ega@statement{egadescription}{remark}{Observación}
\ega@statement{egadescription}{remarks}{Observaciones}
\ega@statement{egadescription}{notation}{Notación}
\makeatother
\renewcommand{\proofname}{Demostración}

% Notation of the EGA editions: the abbreviations the transcription uses for alphabets, symbols, operators and
% arrows. Written for these editions and dedicated to the public domain under CC0 1.0.
\makeatletter
% a command that may or may not exist already is (re)defined in one step
\newcommand{\ega@define}[2]{\providecommand{#1}{}\renewcommand{#1}{#2}}
\newcommand{\ega@operator}[2]{\providecommand{#1}{}\renewcommand{#1}{\operatorname{#2}}}

% alphabets: sheaves, bold letters, categories
\ega@define{\sh}{\mathcal}
\ega@define{\bb}{\mathbf}
\ega@define{\cat}{\mathtt}
\ega@define{\C}{\cat{C}}
\ega@define{\Set}{\cat{Set}}
\ega@define{\CHom}{\cat{Hom}}
\ega@define{\op}{^{\cat{op}}}
\ega@define{\GL}{\bb{GL}}
\ega@define{\bbGamma}{\boldsymbol{\Gamma}}

% letters and symbols
\renewcommand{\leq}{\leqslant}
\renewcommand{\geq}{\geqslant}
\ega@define{\setmin}{-}
\ega@define{\emp}{\varnothing}
\ega@define{\vphi}{\phi}
\ega@define{\kres}{k}
\ega@define{\rad}{\mathfrak{r}}
\ega@define{\nilrad}{\mathfrak{N}}
\ega@define{\HH}{\mathrm{H}}
\ega@define{\CHH}{\check{\mathrm{H}}}
\ega@define{\RR}{\mathrm{R}}
\ega@define{\LL}{\mathrm{L}}
\ega@define{\red}{\mathrm{red}}
\providecommand{\dual}{}\renewcommand{\dual}[1]{{#1}^{\vee}}
% the tilde of the sheaf associated with a module, written after it: M\supertilde
\ega@define{\supertilde}{{\,\widetilde{\,}\,}}

% operators
\ega@operator{\rg}{rg}
\ega@operator{\gr}{gr}
\ega@operator{\Gr}{Gr}
\ega@operator{\Hom}{Hom}
\ega@operator{\Proj}{Proj}
\ega@operator{\Tor}{Tor}
\ega@operator{\Ext}{Ext}
\ega@operator{\Supp}{Supp}
\ega@operator{\Ker}{Ker}
\ega@operator{\Im}{Im}
\ega@operator{\Coker}{Coker}
\ega@operator{\Spec}{Spec}
\ega@operator{\Spm}{Spm}
\ega@operator{\Spf}{Spf}
\ega@operator{\grad}{grad}
\ega@operator{\dimc}{dimc}
\ega@operator{\codim}{codim}
\ega@operator{\id}{id}
\ega@operator{\Der}{Der}
\ega@operator{\Diff}{Diff}
\ega@operator{\Hyp}{Hyp}
\ega@operator{\Tr}{Tr}

% arrows: long in displays, short in text, scripts and subscripts
\newcommand{\ega@arrow}[2]{\mathchoice{#1}{#2}{#2}{#2}}
\let\ega@shortmapsto\mapsto
\renewcommand{\to}{\ega@arrow{\longrightarrow}{\rightarrow}}
\providecommand{\from}{}\renewcommand{\from}{\ega@arrow{\longleftarrow}{\leftarrow}}
\renewcommand{\mapsto}{\ega@arrow{\longmapsto}{\ega@shortmapsto}}
\ega@define{\isoto}{\simeq}
\ega@define{\simto}{\xrightarrow{\sim}}

% wide accents
\newcommand{\ega@wide}[2]{\renewcommand{#1}[1]{#2{##1}}}
\ega@wide{\hat}{\widehat}
\ega@wide{\tilde}{\widetilde}
\makeatother
\def\Frac{\operatorname{Frac}}
\def\shHom{\mathcal{H}\!\mathit{om}}
\def\shProj{\mathcal{P}\!\mathit{roj}}
\def\shAut{\mathcal{A}\!\mathit{ut}}
% A centred row of asterisks.
\newcommand*{\sectionbreak}{\begin{center}***\end{center}}
\newcommand{\oldpage}[2][]{\marginnote{\normalfont\footnotesize\textbf{#1}~|~#2}\ignorespaces}

\let\EGAOriginalLabel\label
\renewcommand{\label}[1]{\phantomsection\EGAOriginalLabel{#1}}
\makeatletter
\def\EGARefNumberAux#1.#2\@nil{#2}
\newcommand{\EGARefNumber}[1]{\begingroup\StrSubstitute{#1}{-fr}{}[\egarefkey]\StrSubstitute{\egarefkey}{-es}{}[\egarefkey]\StrBehind{\egarefkey}{.}\endgroup}
\newcommand{\EGAMaybeRefLink}[2]{%
  \ifcsname r@#1\endcsname\hyperref[#1]{#2}\else#2\fi}
\NewDocumentCommand{\sref}{O{} m O{}}{%
  {\normalfont(\IfBlankF{#1}{\textbf{#1},~}%
  \EGAMaybeRefLink{#2}{\EGARefNumber{#2}}%
  \IfBlankF{#3}{,~#3})}}
\makeatother
% \eref{eq:chapter.number} refers to a numbered display, printing "(number)".
\makeatletter
\newcommand*{\ega@labelnumber}[1]{\StrBehind{#1}{.}}
\newcommand{\eref}[1]{\hyperref[#1]{\normalfont{(\ega@labelnumber{#1})}}}
\makeatother


\newcommand{\EGAResetVolumeCounters}{%
  \setcounter{section}{0}\setcounter{subsection}{0}\setcounter{equation}{0}%
  \setcounter{figure}{0}\setcounter{table}{0}\setcounter{footnote}{0}}
\newcommand{\EGAVolumeTitlePage}[3]{%
  \clearpage\EGAResetVolumeCounters\phantomsection\label{section:#1}%
  \pdfbookmark[0]{#2: #3}{volume.#1}\thispagestyle{empty}\vspace*{\fill}%
  \begin{center}
    {\Huge\bfseries #2\par}\vspace{1.5em}{\LARGE #3\par}\vspace{2em}
    {\large Alexandre Grothendieck y Jean Dieudonné\par}\vspace{1.3em}
    {\large\bfseries Edición española actual\par}\vspace{1.3em}
  \end{center}
  \begin{minipage}{0.88\textwidth}\small
    Traducción, edición y composición producidas por OpenAI Codex — GPT-5.6 Sol, esfuerzo Ultra. No se atribuye revisión humana.\par\medskip
    Fuente: Alexandre Grothendieck y Jean Dieudonné, Publications Mathématiques de l'IHÉS 8 (1961), pp. 5–222; testigo digital NUMDAM PMIHES\_1961\_\_8\_\_5\_0.\par\medskip
    Edición corregida. La traducción según el impreso se ofrece por separado.\par\medskip
    Esta edición independiente conserva la autoría y la relación con la fuente histórica; no implica aval de los autores, de sus derechohabientes ni de NUMDAM.
  \end{minipage}
  \vspace*{\fill}\clearpage}

\begin{document}
% Fuente francesa: PMIHES_1961__8__5_0, printed p. 5.

\EGAVolumeTitlePage{ega2}{EGA II}{Estudio global elemental de algunas clases de morfismos}

% BEGIN RIGHTS AND AI AUTHORSHIP NOTICE
\begingroup
\small
\noindent\textbf{Fuentes y contribuciones}\par
La obra original es de Alexandre Grothendieck, con la colaboración de Jean Dieudonné.
La traducción española emplea como referencia y estructura TeX una edición inglesa
derivada de la traducción humana de \href{https://github.com/ryankeleti/ega}{ryankeleti/ega},
incluida la versión \href{https://github.com/ryankeleti/ega/tree/3d7afe28c99d1887ccdb85acf2ca1afd93fb1971}{3d7afe28}.
Sus colaboradores nombrados son Mateo Carmona, Tim Holzschuh, Tim Hosgood,
Ryan Keleti, Thiago Solovera \& Nery y Wayne Yang. La traducción inglesa y la
estructura heredada conservan su procedencia; no se atribuye a esos colaboradores
la traducción española ni las adiciones posteriores. FUENTES.md precisa las versiones
y los créditos por sección. La atribución no implica autorización ni aval.\par\medskip
La traducción y composición previas están documentadas como trabajo de OpenAI Codex,
GPT-5.6 Sol, esfuerzo Ultra. La composición y preparación de esta edición incluyen
OpenAI Codex, GPT-6 Astra, esfuerzo Ultra.\par\medskip
\noindent\textbf{Derechos, disponibilidad y autoría}

Los derechos sobre las obras originales y sobre las contribuciones protegidas a estas traducciones y transcripciones derivadas corresponden a sus respectivos titulares. Todo aquello sobre lo que el editor tenga derechos en estas ediciones queda dedicado al dominio público bajo CC0 1.0 (https://creativecommons.org/publicdomain/zero/1.0/deed.es). Las menciones de derechos de terceros y las autorizaciones ya concedidas siguen siendo aplicables.

Estas ediciones se ofrecen gratuitamente, con fines no comerciales y exclusivamente educativos, bajo la presunción de que los respectivos titulares de derechos no se opondrían a que estas traducciones y transcripciones estuvieran disponibles gratuitamente junto a los originales franceses en \href{https://www.numdam.org/}{NUMDAM}. Esta presunción no significa que se haya obtenido autorización; el acceso gratuito a los originales no autoriza por sí solo su traducción, adaptación o republicación.

Si usted es titular de derechos sobre algún material presentado aquí y se opone a su disponibilidad, puede \href{https://github.com/KokunoYumeto/ega-fga-es/issues}{abrir una incidencia en GitHub}. El responsable de la publicación retirará el material afectado a petición del titular de los derechos.

La traducción al español está producida por IA. Las correcciones, la edición y la composición también están producidas por IA. El responsable de la publicación no tradujo el texto. Se conservan la autoría y las atribuciones de las obras originales y de las contribuciones heredadas; no se atribuyen a la IA. Los modelos documentados se identifican en cada edición.
\par
\endgroup
\clearpage
% END RIGHTS AND AI AUTHORSHIP NOTICE


% EGA II SOURCE FRONT MATTER
\oldpage[II]{5}
\phantomsection
\label{chapter:II-fr}
\begin{center}
{\large CAPÍTULO II\par}
\vspace{1.5em}
{\Large\bfseries ESTUDIO GLOBAL ELEMENTAL\par}
{\Large\bfseries DE ALGUNAS CLASES DE MORFISMOS\par}
\vspace{1.8em}
\end{center}
\section*{Sumario}

\begin{longtable}{@{}p{0.08\textwidth}p{0.86\textwidth}@{}}
  \hyperref[section:II.1-fr]{\textsection1}. & Morfismos afines.\\
  \hyperref[section:II.2-fr]{\textsection2}. & Espectros primos homogéneos.\\
  \hyperref[section:II.3-fr]{\textsection3}. & Espectro homogéneo de un haz de álgebras graduadas.\\
  \hyperref[section:II.4-fr]{\textsection4}. & Fibrados proyectivos; haces amplios.\\
  \hyperref[section:II.5-fr]{\textsection5}. & Morfismos cuasiafines; morfismos cuasiproyectivos; morfismos propios; morfismos proyectivos.\\
  \hyperref[section:II.6-fr]{\textsection6}. & Morfismos enteros y morfismos finitos.\\
  \hyperref[section:II.7-fr]{\textsection7}. & Criterios valuativos.\\
  \hyperref[section:II.8-fr]{\textsection8}. & Esquemas obtenidos por explosión; conos proyectantes; clausura proyectiva.
\end{longtable}
\bigskip

Las diversas clases de morfismos estudiadas en este capítulo se tratan sin hacer gran uso de los métodos cohomológicos; en el capítulo~III se efectuará un estudio más profundo de ellas, utilizando esos métodos, y allí se emplearán principalmente las secciones~\textsection\textsection\hyperref[section:II.2-fr]{2}, \hyperref[section:II.4-fr]{4} y \hyperref[section:II.5-fr]{5} del capítulo~II.
La sección~\textsection\hyperref[section:II.8-fr]{8} puede omitirse en una primera lectura: ofrece algunos complementos al formalismo desarrollado en las secciones~\textsection\textsection\hyperref[section:II.1-fr]{1} a \hyperref[section:II.3-fr]{3}, que se reducen a aplicaciones sencillas de dicho formalismo, y haremos de ella un uso menos constante que de los demás resultados de este capítulo.
\bigskip
% Fuente francesa: PMIHES_1961__8__5_0, printed pp. 5-6.

\section{Morfismos afines}
\label{section:II.1-fr}

La mayoría de los resultados de esta sección son las contrapartidas ``globales'' de los del capítulo~I, \textsection1; por tanto, no son esencialmente nuevos y proporcionan simplemente un lenguaje cómodo para lo que sigue.
\bigskip

\subsection{$S$-preesquemas y $\mathcal{O}_S$-álgebras}
\label{subsection:II.1.1-fr}

\begin{env}[1.1.1]
\label{II.1.1.1-fr}
Sean $S$ un preesquema, $X$ un $S$-preesquema y $f:X\to S$ su morfismo estructural.
Se sabe por \sref[0]{0.4.2.4-fr} que la imagen directa $f_*(\mathcal{O}_X)$ es una $\mathcal{O}_S$-álgebra, que
\oldpage[II]{6}
denotaremos por $\mathcal{A}(X)$ cuando no haya riesgo de confusión; si $U$ es un abierto de $S$, se tiene
\[
  \mathcal{A}(f^{-1}(U))=\mathcal{A}(X)|U.
\]
De manera análoga, para todo $\mathcal{O}_X$-módulo $\mathcal{F}$ (resp. toda $\mathcal{O}_X$-álgebra $\mathcal{B}$), escribimos $\mathcal{A}(\mathcal{F})$ (resp. $\mathcal{A}(\mathcal{B})$) para la imagen directa $f_*(\mathcal{F})$ (resp. $f_*(\mathcal{B})$), que es un $\mathcal{A}(X)$-módulo (resp. una $\mathcal{A}(X)$-álgebra) y no solamente un $\mathcal{O}_S$-módulo (resp. una $\mathcal{O}_S$-álgebra).
\end{env}

\begin{env}[1.1.2]
\label{II.1.1.2-fr}
Sea $Y$ un segundo $S$-preesquema, sea $g:Y\to S$ su morfismo estructural y sea $h:X\to Y$ un $S$-morfismo; se tiene entonces el diagrama conmutativo
\[
  \xymatrix{
    X\ar[rr]^h\ar[rd]_f && Y\ar[ld]^g\\
    & S
  }
  \tag{1.1.2.1}\label{II.1.1.2.1-fr}
\]

Por definición, $h=(\psi,\theta)$, donde $\theta:\mathcal{O}_Y\to h_*(\mathcal{O}_X)=\psi_*(\mathcal{O}_X)$ es un homomorfismo de haces de anillos; se deduce así de \sref[0]{0.4.2.2-fr} un homomorfismo de $\mathcal{O}_S$-álgebras $g_*(\theta):g_*(\mathcal{O}_Y)\to g_*(h_*(\mathcal{O}_X))=f_*(\mathcal{O}_X)$, es decir, un homomorfismo de $\mathcal{O}_S$-álgebras $\mathcal{A}(Y)\to\mathcal{A}(X)$, que denotamos por $\mathcal{A}(h)$.
Si $h':Y\to Z$ es un segundo $S$-morfismo, es inmediato que $\mathcal{A}(h'\circ h)=\mathcal{A}(h)\circ\mathcal{A}(h')$.
Hemos definido así un \emph{funtor contravariante $\mathcal{A}(X)$} sobre la categoría de los $S$-preesquemas, con valores en la categoría de las $\mathcal{O}_S$-álgebras.

Sean ahora $\mathcal{F}$ un $\mathcal{O}_X$-módulo, $\mathcal{G}$ un $\mathcal{O}_Y$-módulo y $u:\mathcal{G}\to\mathcal{F}$ un $h$-morfismo, es decir, según \sref[0]{0.4.4.1-fr}, un homomorfismo de $\mathcal{O}_Y$-módulos $\mathcal{G}\to h_*(\mathcal{F})$.
Entonces $g_*(u):g_*(\mathcal{G})\to g_*(h_*(\mathcal{F}))=f_*(\mathcal{F})$ es un homomorfismo $\mathcal{A}(\mathcal{G})\to
\mathcal{A}(\mathcal{F})$ de $\mathcal{O}_S$-módulos, que denotamos por $\mathcal{A}(u)$; además, el par $(\mathcal{A}(h),\mathcal{A}(u))$ constituye un \emph{dihomomorfismo} del $\mathcal{A}(Y)$-módulo $\mathcal{A}(\mathcal{G})$ al $\mathcal{A}(X)$-módulo $\mathcal{A}(\mathcal{F})$.
\end{env}

\begin{env}[1.1.3]
\label{II.1.1.3-fr}
Fijado el preesquema $S$, se puede considerar los pares $(X,\mathcal{F})$, donde $X$ es un $S$-preesquema y $\mathcal{F}$ un $\mathcal{O}_X$-módulo, como objetos de una \emph{categoría}, definiendo un \emph{morfismo} $(X,\mathcal{F})\to(Y,\mathcal{G})$ como un par $(h,u)$, donde $h:X\to Y$ es un $S$-morfismo y $u:\mathcal{G}\to\mathcal{F}$ un $h$-morfismo.
Se puede decir entonces que $(\mathcal{A}(X),\mathcal{A}(\mathcal{F}))$ es un \emph{funtor contravariante} con valores en la categoría cuyos objetos son los pares formados por una $\mathcal{O}_S$-álgebra y un módulo sobre dicha álgebra, y cuyos morfismos son los dihomomorfismos.
\end{env}
% Fuente francesa: PMIHES_1961__8__5_0, printed pp. 6-7.

\subsection{Preesquemas afines sobre un preesquema}
\label{subsection:II.1.2-fr}

\begin{definition}[1.2.1]
\label{II.1.2.1-fr}
Sean $X$ un $S$-preesquema y $f:X\to S$ su morfismo estructural.
Se dice que $X$ es afín sobre $S$ si existe un recubrimiento $(S_\alpha)$ de $S$ por abiertos afines tal que, para todo $\alpha$, el preesquema inducido por $X$ sobre el abierto $f^{-1}(S_\alpha)$ es afín.
\end{definition}

\begin{example}[1.2.2]
\label{II.1.2.2-fr}
Todo subpreesquema cerrado de $S$ es un $S$-preesquema afín sobre $S$ (\sref[I]{I.4.2.3-fr} y \sref[I]{I.4.2.4-fr}).
\end{example}

\begin{remark}[1.2.3]
\label{II.1.2.3-fr}
Un preesquema $X$ afín sobre $S$ no es necesariamente un esquema afín, como muestra el ejemplo $X=S$ \sref{II.1.2.2-fr}.
Por otra parte, si un esquema afín $X$ es un $S$-preesquema, $X$ no es necesariamente afín sobre
\oldpage[II]{7}
$S$ (véase el Ejemplo~\sref{II.1.3.3-fr}).
Sin embargo, recordemos que, si $S$ es un \emph{esquema}, todo $S$-preesquema que sea un esquema afín es afín sobre $S$ \sref[I]{I.5.5.10-fr}.
\end{remark}

\begin{proposition}[1.2.4]
\label{II.1.2.4-fr}
Todo $S$-preesquema afín sobre $S$ es separado sobre $S$ \textup{(dicho de otro modo, es un $S$-\emph{esquema})}.
\end{proposition}

\begin{proof}
Esto se sigue inmediatamente de \sref[I]{I.5.5.5-fr} y \sref[I]{I.5.5.8-fr}.
\end{proof}
% Fuente francesa: PMIHES_1961__8__5_0, printed p. 7.

\begin{proposition}[1.2.5]
\label{II.1.2.5-fr}
Sean $X$ un $S$-esquema afín sobre $S$ y $f:X\to S$ su morfismo estructural.
Para todo abierto $U\subset S$, $f^{-1}(U)$ es afín sobre $U$.
\end{proposition}

\begin{proof}
En virtud de la Definición~\sref{II.1.2.1-fr}, se puede reducir al caso en que $S=\operatorname{Spec}(A)$ y $X=\operatorname{Spec}(B)$ son afines; entonces $f=({}^a\varphi,\widetilde{\varphi})$, donde $\varphi:A\to B$ es un homomorfismo.
Como los $D(g)$, para $g\in A$, forman una base de $S$, se puede reducir al caso en que $U=D(g)$; pero entonces se sabe que $f^{-1}(U)=D(\varphi(g))$ (\hyperref[I.1.2.2.2-fr]{\textbf{I}, 1.2.2.2}), de donde se sigue la proposición.
\end{proof}

\begin{proposition}[1.2.6]
\label{II.1.2.6-fr}
Sean $X$ un $S$-esquema afín sobre $S$ y $f:X\to S$ su morfismo estructural.
Para todo $\mathcal{O}_X$-módulo cuasicoherente $\mathcal{F}$, $f_*(\mathcal{F})$ es un $\mathcal{O}_S$-módulo cuasicoherente.
\end{proposition}

\begin{proof}
Teniendo en cuenta la Proposición~\sref{II.1.2.4-fr}, esto se sigue de \sref[I]{I.9.2.2-fr}[\emph{a}].
\end{proof}

En particular, la $\mathcal{O}_S$-álgebra $\mathcal{A}(X)=f_*(\mathcal{O}_X)$ es \emph{cuasicoherente}.
% Fuente francesa: PMIHES_1961__8__5_0, printed pp. 7-8.

\begin{proposition}[1.2.7]
\label{II.1.2.7-fr}
Sea $X$ un $S$-esquema afín sobre $S$.
Para todo $S$-preesquema $Y$, la aplicación $h\to\mathcal{A}(h)$ del conjunto $\operatorname{Hom}_S(Y,X)$ en el conjunto $\operatorname{Hom}(\mathcal{A}(X),\mathcal{A}(Y))$ \sref{II.1.1.2-fr} es biyectiva.
\end{proposition}

\begin{proof}
Sean $f:X\to S$ y $g:Y\to S$ los morfismos estructurales.
Supongamos primero que $S=\operatorname{Spec}(A)$ y $X=\operatorname{Spec}(B)$ son afines; hay que demostrar que, para todo homomorfismo $\omega:f_*(\mathcal{O}_X)\to g_*(\mathcal{O}_Y)$ de $\mathcal{O}_S$-álgebras, existe un único $S$-morfismo $h:Y\to X$ tal que $\mathcal{A}(h)=\omega$.
Por definición, para todo abierto $U\subset
S$, $\omega$ define un homomorfismo
$\omega_U=\Gamma(U,\omega):\Gamma(f^{-1}(U),\mathcal{O}_X)\to\Gamma(g^{-1}(U),\mathcal{O}_Y)$
de $\Gamma(U,\mathcal{O}_S)$-álgebras.
En particular, para $U=S$, se obtiene un homomorfismo $\varphi:\Gamma(X,\mathcal{O}_X)\to\Gamma(Y,\mathcal{O}_Y)$ de $\Gamma(S,\mathcal{O}_S)$-álgebras, al que corresponde un $S$-morfismo bien determinado $h:Y\to X$, puesto que $X$ es afín \sref[I]{I.2.2.4-fr}.
Queda por demostrar que $\mathcal{A}(h)=\omega$; dicho de otro modo, que, para todo abierto $U$ de una base de $S$, $\omega_U$ coincide con el homomorfismo de álgebras $\varphi_U$ correspondiente al $S$-morfismo $g^{-1}(U)\to f^{-1}(U)$ que es restricción de $h$.
Se puede limitar al caso en que $U=D(\lambda)$, con $\lambda\in A$; entonces, si $f=({}^a\rho,\widetilde{\rho})$, donde $\rho:A\to B$ es un homomorfismo de anillos, se tiene $f^{-1}(U)=D(\mu)$, donde $\mu=\rho(\lambda)$, y $\Gamma(f^{-1}(U),\mathcal{O}_X)$ es el anillo de fracciones $B_\mu$; el diagrama
\[
  \xymatrix{
    B\ar[r]^\varphi\ar[d] & \Gamma(Y,\mathcal{O}_Y)\ar[d]\\
    B_\mu\ar[r]_{\varphi_U} & \Gamma(g^{-1}(U),\mathcal{O}_Y)
  }
\]
es conmutativo, y lo mismo ocurre con el diagrama análogo en el que $\varphi_U$ se reemplaza por $\omega_U$; la igualdad $\varphi_U=\omega_U$ se sigue entonces de la propiedad universal de los anillos de fracciones \sref[0]{0.1.2.4-fr}.

Pasemos al caso general y sea $(S_\alpha)$ un recubrimiento de $S$ por abiertos afines
\oldpage[II]{8}
tales que los $f^{-1}(S_\alpha)$ sean afines.
Todo homomorfismo $\omega:\mathcal{A}(X)\to\mathcal{A}(Y)$ de $\mathcal{O}_S$-álgebras da entonces, por restricción, una familia de homomorfismos
\[
  \omega_\alpha:
  \mathcal{A}(f^{-1}(S_\alpha))\to
  \mathcal{A}(g^{-1}(S_\alpha))
\]
de $\mathcal{O}_{S_\alpha}$-álgebras y, por lo anterior, una familia de $S_\alpha$-morfismos $h_\alpha:g^{-1}(S_\alpha)\to f^{-1}(S_\alpha)$.
Basta ver que, para todo abierto afín $U$ de una base de $S_\alpha\cap S_\beta$, las restricciones de $h_\alpha$ y $h_\beta$ a $g^{-1}(U)$ coinciden, lo cual es evidente, puesto que, por lo anterior, ambas restricciones corresponden al homomorfismo $\mathcal{A}(X)|U\to\mathcal{A}(Y)|U$ que es restricción de $\omega$.
\end{proof}

\begin{corollary}[1.2.8]
\label{II.1.2.8-fr}
Sean $X$ e $Y$ dos $S$-esquemas afines sobre $S$.
Para que un $S$-morfismo $h:Y\to X$ sea un isomorfismo, es necesario y suficiente que $\mathcal{A}(h):\mathcal{A}(X)\to\mathcal{A}(Y)$ sea un isomorfismo.
\end{corollary}

\begin{proof}
Esto se sigue inmediatamente de la Proposición~\sref{II.1.2.7-fr} y del carácter funtorial de $\mathcal{A}(X)$.
\end{proof}
% Fuente francesa: PMIHES_1961__8__5_0, printed p. 8.

\subsection{Preesquema afín sobre $S$ asociado a una $\mathcal{O}_S$-álgebra}
\label{subsection:II.1.3-fr}


\begin{proposition}[1.3.1]
\label{II.1.3.1-fr}
Sea $S$ un preesquema.
Para toda $\mathcal{O}_S$-álgebra cuasicoherente $\mathcal{B}$, existe un preesquema $X$ afín sobre $S$, determinado salvo $S$-isomorfismo único, tal que $\mathcal{A}(X)=\mathcal{B}$.
\end{proposition}

\begin{proof}
La unicidad resulta del Corolario~\sref{II.1.2.8-fr}; demostremos la existencia de $X$.
Para todo abierto afín $U\subset S$, sea $X_U$ el preesquema $\operatorname{Spec}(\Gamma(U,\mathcal{B}))$; como $\Gamma(U,\mathcal{B})$ es una $\Gamma(U,\mathcal{O}_S)$-álgebra, $X_U$ es un $S$-preesquema \sref[I]{I.1.6.1-fr}.
Además, como $\mathcal{B}$ es cuasicoherente, la $\mathcal{O}_S$-álgebra $\mathcal{A}(X_U)$ se identifica canónicamente con $\mathcal{B}|U$ (\sref[I]{I.1.3.7-fr}, \sref[I]{I.1.3.13-fr} y \sref[I]{I.1.6.3-fr}).

Sea $V$ un segundo abierto afín de $S$, y sea $X_{U,V}$ el preesquema inducido por $X_U$ sobre $f_U^{-1}(U\cap V)$, donde $f_U$ denota el morfismo estructural $X_U\to S$; $X_{U,V}$ y $X_{V,U}$ son afines sobre $U\cap V$ \sref{II.1.2.5-fr}, y, por definición, $\mathcal{A}(X_{U,V})$ y $\mathcal{A}(X_{V,U})$ se identifican canónicamente con $\mathcal{B}|(U\cap V)$.
Existe, por tanto, según \sref{II.1.2.8-fr}, un $S$-isomorfismo canónico $\theta_{U,V}:X_{V,U}\to X_{U,V}$; además, si $W$ es un tercer abierto afín de $S$, y si $\theta'_{U,V}$, $\theta'_{V,W}$ y $\theta'_{U,W}$ son las restricciones de $\theta_{U,V}$, $\theta_{V,W}$ y $\theta_{U,W}$ a las imágenes inversas de $U\cap V\cap W$ en $X_V$, $X_W$ y $X_W$, respectivamente, por los morfismos estructurales, se tiene
$\theta'_{U,V}\circ\theta'_{V,W}=\theta'_{U,W}$.

Existe, en consecuencia, un preesquema $X$, un recubrimiento $(T_U)$ de $X$ por abiertos afines y, para cada $U$, un isomorfismo $\varphi_U:X_U\to T_U$, tales que $\varphi_U$ lleva $f_U^{-1}(U\cap V)$ sobre $T_U\cap T_V$ y
$\theta_{U,V}=\varphi_U^{-1}\circ\varphi_V$
\sref[I]{I.2.3.1-fr}.
El morfismo $g_U=f_U\circ\varphi_U^{-1}$ dota a $T_U$ de estructura de $S$-preesquema, y los morfismos $g_U$ y $g_V$ coinciden sobre $T_U\cap T_V$; por tanto, $X$ es un $S$-preesquema.
Además, por definición es claro que $X$ es afín sobre $S$ y que $\mathcal{A}(T_U)=\mathcal{B}|U$; por consiguiente, $\mathcal{A}(X)=\mathcal{B}$.
\end{proof}

Se dice que el $S$-esquema $X$ así definido está \emph{asociado a la $\mathcal{O}_S$-álgebra $\mathcal{B}$}, o que es el \emph{espectro} de $\mathcal{B}$, y se denota por $\operatorname{Spec}(\mathcal{B})$.
% Fuente francesa: PMIHES_1961__8__5_0, printed pp. 8-9.

\begin{corollary}[1.3.2]
\label{II.1.3.2-fr}
Sean $X$ un preesquema afín sobre $S$ y $f:X\to S$ su morfismo estructural.
Para todo abierto afín $U\subset S$, el preesquema inducido sobre $f^{-1}(U)$ es el esquema afín de anillo $\Gamma(U,\mathcal{A}(X))$.
\end{corollary}

\begin{proof}
\oldpage[II]{9}
Como, en virtud de \sref{II.1.2.6-fr} y \sref{II.1.3.1-fr}, se puede suponer que $X$ está asociado a una $\mathcal{O}_S$-álgebra, el corolario resulta de la construcción de $X$ descrita en \sref{II.1.3.1-fr}.
\end{proof}

\begin{example}[1.3.3]
\label{II.1.3.3-fr}
Sea $S$ el plano afín sobre un cuerpo $K$, donde se ha duplicado el punto $0$ \sref[I]{I.5.5.11-fr}; con las notaciones de \sref[I]{I.5.5.11-fr}, $S$ es la unión de dos abiertos afines $Y_1$ e $Y_2$; si $f$ es la inmersión abierta $Y_1\to S$, entonces $f^{-1}(Y_2)=Y_1\cap Y_2$ no es un abierto afín de $Y_1$ \emph{(loc. cit.)}, lo que proporciona un ejemplo de esquema afín que no es afín sobre $S$.
\end{example}

\begin{corollary}[1.3.4]
\label{II.1.3.4-fr}
Sea $S$ un esquema afín; para que un $S$-preesquema $X$ sea afín sobre $S$, es necesario y suficiente que $X$ sea un esquema afín.
\end{corollary}

\begin{corollary}[1.3.5]
\label{II.1.3.5-fr}
Sean $X$ un preesquema afín sobre un preesquema $S$ e $Y$ un $X$-preesquema.
Para que $Y$ sea afín sobre $S$, es necesario y suficiente que $Y$ sea afín sobre $X$.
\end{corollary}

\begin{proof}
Se reduce inmediatamente al caso en que $S$ es un esquema afín, y entonces también lo es $X$ \sref{II.1.3.4-fr}; las dos condiciones del enunciado expresan entonces que $Y$ es un esquema afín \sref{II.1.3.4-fr}.
\end{proof}

\begin{env}[1.3.6]
\label{II.1.3.6-fr}
Sea $X$ un preesquema afín sobre $S$.
Definir un preesquema $Y$ afín \emph{sobre $X$} equivale, en virtud de \sref{II.1.3.5-fr}, a darse un preesquema $Y$ afín \emph{sobre $S$} y un $S$-morfismo $g:Y\to X$; dicho de otro modo, según \sref{II.1.3.1-fr} y \sref{II.1.2.7-fr}, equivale a darse una $\mathcal{O}_S$-álgebra cuasicoherente $\mathcal{B}$ y un homomorfismo $\mathcal{A}(X)\to\mathcal{B}$ de $\mathcal{O}_S$-álgebras (que puede considerarse como la definición sobre $\mathcal{B}$ de una estructura de $\mathcal{A}(X)$-álgebra).
Si $f:X\to S$ es el morfismo estructural, se tiene entonces $\mathcal{B}=f_*(g_*(\mathcal{O}_Y))$.
\end{env}
% Fuente francesa: PMIHES_1961__8__5_0, printed p. 9.

\begin{corollary}[1.3.7]
\label{II.1.3.7-fr}
Sea $X$ un preesquema afín sobre $S$; para que $X$ sea de tipo finito sobre $S$, es necesario y suficiente que la $\mathcal{O}_S$-álgebra cuasicoherente $\mathcal{A}(X)$ sea de tipo finito \sref[I]{I.9.6.2-fr}.
\end{corollary}

\begin{proof}
Por definición \sref[I]{I.9.6.2-fr}, se puede reducir al caso en que $S$ es afín; entonces $X$ es un esquema afín \sref{II.1.3.4-fr}, y, si $S=\operatorname{Spec}(A)$ y $X=\operatorname{Spec}(B)$, $\mathcal{A}(X)$ es la $\mathcal{O}_S$-álgebra $\widetilde{B}$; como $\Gamma(U,\widetilde{B})=B$, el corolario resulta de \sref[I]{I.9.6.2-fr} y \sref[I]{I.6.3.3-fr}.
\end{proof}

\begin{corollary}[1.3.8]
\label{II.1.3.8-fr}
Sea $X$ un preesquema afín sobre $S$; para que $X$ sea reducido, es necesario y suficiente que la $\mathcal{O}_S$-álgebra cuasicoherente $\mathcal{A}(X)$ sea reducida \sref[0]{0.4.1.4-fr}.
\end{corollary}

\begin{proof}
La cuestión es local sobre $S$; en virtud de \sref{II.1.3.2-fr}, el corolario resulta de \sref[I]{I.5.1.1-fr} y \sref[I]{I.5.1.4-fr}.
\end{proof}
% Fuente francesa: PMIHES_1961__8__5_0, printed pp. 9-10.

\subsection{Haces cuasicoherentes sobre un preesquema afín sobre $S$}
\label{subsection:II.1.4-fr}

\begin{proposition}[1.4.1]
\label{II.1.4.1-fr}
Sean $X$ un preesquema afín sobre $S$, $Y$ un $S$-preesquema y $\mathcal{F}$ \textup{(resp. $\mathcal{G}$)} un $\mathcal{O}_X$-módulo cuasicoherente \textup{(resp.} un $\mathcal{O}_Y$-módulo\textup{)}.
La aplicación $(h,u)\to(\mathcal{A}(h),\mathcal{A}(u))$ del conjunto de los morfismos $(Y,\mathcal{G})\to(X,\mathcal{F})$ en el conjunto de los dihomomorfismos $(\mathcal{A}(X),\mathcal{A}(\mathcal{F}))\to
(\mathcal{A}(Y),\mathcal{A}(\mathcal{G}))$ \textup{(\sref{II.1.1.2-fr} y \sref{II.1.1.3-fr})} es biyectiva.
\end{proposition}

\begin{proof}
La demostración sigue exactamente el mismo procedimiento que la de \sref{II.1.2.7-fr}, utilizando \sref[I]{I.2.2.5-fr} y \sref[I]{I.2.2.4-fr}; se dejan los detalles al lector.
\end{proof}

\begin{corollary}[1.4.2]
\label{II.1.4.2-fr}
Si, además de las hipótesis de \sref{II.1.4.1-fr}, se supone que $Y$ es afín sobre $S$, para que $(h,u)$ sea un isomorfismo es necesario y suficiente que $(\mathcal{A}(h),\mathcal{A}(u))$ sea un diisomorfismo.
\end{corollary}

\oldpage[II]{10}
\begin{proposition}[1.4.3]
\label{II.1.4.3-fr}

Para todo par $(\mathcal{B},\mathcal{M})$ formado por una $\mathcal{O}_S$-álgebra cuasicoherente $\mathcal{B}$ y un $\mathcal{B}$-módulo cuasicoherente $\mathcal{M}$ \emph{(considerado como $\mathcal{O}_S$-módulo o como $\mathcal{B}$-módulo, lo cual es equivalente~\sref[I]{I.9.6.1-fr})}, existe un par $(X,\mathcal{F})$ formado por un preesquema $X$ afín sobre $S$ y un $\mathcal{O}_X$-módulo cuasicoherente $\mathcal{F}$, tal que $\mathcal{A}(X)=\mathcal{B}$ y $\mathcal{A}(\mathcal{F})=\mathcal{M}$; además, este par está determinado salvo isomorfismo único.
\end{proposition}

\begin{proof}
La unicidad resulta de \sref{II.1.4.1-fr} y \sref{II.1.4.2-fr}; la existencia se demuestra como en \sref{II.1.3.1-fr}, y dejamos los detalles al lector.
\end{proof}

Se denota por $\widetilde{\mathcal{M}}$ el $\mathcal{O}_X$-módulo $\mathcal{F}$ y se dice que está \emph{asociado} al $\mathcal{B}$-módulo cuasicoherente $\mathcal{M}$; para todo abierto afín $U\subset S$, $\widetilde{\mathcal{M}}|p^{-1}(U)$ (donde $p$ es el morfismo estructural $X\to S$) es canónicamente isomorfo a $(\Gamma(U,\mathcal{M}))^{\sim}$.

\begin{corollary}[1.4.4]
\label{II.1.4.4-fr}
En la categoría de los $\mathcal{B}$-módulos cuasicoherentes, $\widetilde{\mathcal{M}}$ es un funtor covariante aditivo exacto en $\mathcal{M}$, que conmuta con los límites inductivos y las sumas directas.
\end{corollary}

\begin{proof}
Se reduce inmediatamente al caso en que $S$ es afín; el corolario resulta entonces de \sref[I]{I.1.3.5-fr}, \sref[I]{I.1.3.9-fr} y \sref[I]{I.1.3.11-fr}.
\end{proof}
% Fuente francesa: PMIHES_1961__8__5_0, printed p. 10.

\begin{corollary}[1.4.5]
\label{II.1.4.5-fr}
Bajo las hipótesis de \sref{II.1.4.3-fr}, para que $\widetilde{\mathcal{M}}$ sea un $\mathcal{O}_X$-módulo de tipo finito, es necesario y suficiente que $\mathcal{M}$ sea un $\mathcal{B}$-módulo de tipo finito.
\end{corollary}

\begin{proof}
Se reduce inmediatamente al caso en que $S=\operatorname{Spec}(A)$ es un esquema afín.
Entonces $\mathcal{B}=\widetilde{B}$, donde $B$ es una $A$-álgebra de tipo finito (\hyperref[I.9.6.2-fr]{\textbf{I}, 9.6.2}), y $\mathcal{M}=\widetilde{M}$, donde $M$ es un $B$-módulo \sref[I]{I.1.3.13-fr}; \emph{sobre el preesquema $X$}, $\mathcal{O}_X$ está asociado al anillo $B$ y $\widetilde{\mathcal{M}}$ al $B$-módulo $M$; para que $\widetilde{\mathcal{M}}$ sea de tipo finito es, por tanto, necesario y suficiente que $M$ sea de tipo finito \sref[I]{I.1.3.13-fr}, de donde se sigue la afirmación.
\end{proof}

\begin{proposition}[1.4.6]
\label{II.1.4.6-fr}
Sean $Y$ un preesquema afín sobre $S$, y $X$ y $X'$ dos preesquemas afines sobre $Y$ \emph{(y, por tanto, también sobre $S$~\sref{II.1.3.5-fr})}.
Sean $\mathcal{B}=\mathcal{A}(Y)$, $\mathcal{A}=\mathcal{A}(X)$ y $\mathcal{A}'=\mathcal{A}(X')$.
Entonces $X\times_Y X'$ es afín sobre $Y$ \emph{(y, por tanto, también sobre $S$)}, y $\mathcal{A}(X\times_Y X')$ se identifica canónicamente con $\mathcal{A}\otimes_{\mathcal{B}}\mathcal{A}'$.
\end{proposition}

\begin{proof}
En virtud de \sref[I]{I.9.6.1-fr}, $\mathcal{A}\otimes_{\mathcal{B}}\mathcal{A}'$ es una $\mathcal{B}$-álgebra cuasicoherente y, por tanto, también una $\mathcal{O}_S$-álgebra cuasicoherente \sref[I]{I.9.6.1-fr}; sea $Z$ el espectro de $\mathcal{A}\otimes_{\mathcal{B}}\mathcal{A}'$ \sref{II.1.3.1-fr}.
Los $\mathcal{B}$-homomorfismos canónicos $\mathcal{A}\to\mathcal{A}\otimes_{\mathcal{B}}\mathcal{A}'$ y $\mathcal{A}'\to\mathcal{A}\otimes_{\mathcal{B}}\mathcal{A}'$ corresponden, según \sref{II.1.2.7-fr}, a $Y$-morfismos $p:Z\to X$ y $p':Z\to X'$.
Para ver que la terna $(Z,p,p')$ es un producto $X\times_Y X'$, se puede reducir al caso en que $S$ es un esquema afín de anillo $C$ \sref[I]{I.3.2.6.4-fr}.
Entonces $Y$, $X$ y $X'$ son esquemas afines \sref{II.1.3.4-fr}, cuyos anillos $B$, $A$ y $A'$ son $C$-álgebras tales que $\mathcal{B}=\widetilde{B}$, $\mathcal{A}=\widetilde{A}$ y $\mathcal{A}'=\widetilde{A}'$.
Se sabe entonces por \sref[I]{I.1.3.13-fr} que $\mathcal{A}\otimes_{\mathcal{B}}\mathcal{A}'$ se identifica canónicamente con la $\mathcal{O}_S$-álgebra $(A\otimes_B A')\supertilde$; por tanto, el anillo $A(Z)$ se identifica con $A\otimes_B A'$, y los morfismos $p$ y $p'$ corresponden a los homomorfismos canónicos $A\to A\otimes_B A'$ y $A'\to A\otimes_B A'$.
La proposición resulta entonces de \sref[I]{I.3.2.2-fr}.
\end{proof}
% Fuente francesa: PMIHES_1961__8__5_0, printed pp. 10-11.

\begin{corollary}[1.4.7]
\label{II.1.4.7-fr}
Sean $\mathcal{F}$ \textup{(resp. $\mathcal{F}'$)} un $\mathcal{O}_X$-módulo cuasicoherente \textup{(resp.} un $\mathcal{O}_{X'}$-módulo cuasicoherente); entonces $\mathcal{A}(\mathcal{F}\otimes_Y\mathcal{F}')$ se identifica canónicamente con $\mathcal{A}(\mathcal{F})\otimes_{\mathcal{A}(Y)}
\mathcal{A}(\mathcal{F}')$.
\end{corollary}

\begin{proof}
Se sabe que $\mathcal{F}\otimes_Y\mathcal{F}'$ es cuasicoherente sobre $X\times_Y X'$ \sref[I]{I.9.1.2-fr}.
Sean $g:Y\to S$, $f:X\to Y$ y $f':X'\to Y$ los morfismos estructurales, de modo que el morfismo estructural
\oldpage[II]{11}
$h:Z\to S$ es igual a $g\circ f\circ p$ y a $g\circ f'\circ p'$.
Se define un homomorfismo canónico
\[
  \mathcal{A}(\mathcal{F})\otimes_{\mathcal{A}(Y)}\mathcal{A}(\mathcal{F}')
  \longrightarrow \mathcal{A}(\mathcal{F}\otimes_Y\mathcal{F}')
\]
del modo siguiente: para todo abierto $U\subset S$, se tienen homomorfismos canónicos
$\Gamma(f^{-1}(g^{-1}(U)),\mathcal{F})\to\Gamma(h^{-1}(U),p^*(\mathcal{F}))$ y $\Gamma(f^{\prime-1}(g^{-1}(U)),\mathcal{F}')\to\Gamma(h^{-1}(U),p^{\prime*}(\mathcal{F}'))$
\sref[0]{0.4.4.3-fr}, de donde se deduce un homomorfismo canónico
\[
  \begin{aligned}
  &\Gamma(f^{-1}(g^{-1}(U)),\mathcal{F})
    \otimes_{\Gamma(g^{-1}(U),\mathcal{O}_Y)}
    \Gamma(f^{\prime-1}(g^{-1}(U)),\mathcal{F}')\longrightarrow\\
  &\hspace{17em}
    \Gamma(h^{-1}(U),p^*(\mathcal{F}))
    \otimes_{\Gamma(h^{-1}(U),\mathcal{O}_Z)}
    \Gamma(h^{-1}(U),p^{\prime*}(\mathcal{F}')).
  \end{aligned}
\]

Para ver que se ha definido un isomorfismo de $\mathcal{A}(Z)$-módulos, se puede reducir al caso en que $S$ (y, por tanto, $X$, $X'$, $Y$ y $X\times_Y X'$) son esquemas afines y, con las notaciones de \sref{II.1.4.6-fr}, $\mathcal{F}=\widetilde{M}$ y $\mathcal{F}'=\widetilde{M}'$, donde $M$ (resp. $M'$) es un $A$-módulo (resp. un $A'$-módulo).
Entonces $\mathcal{F}\otimes_Y\mathcal{F}'$ se identifica con el haz sobre $X\times_Y X'$ asociado al $(A\otimes_B A')$-módulo $M\otimes_B M'$ \sref[I]{I.9.1.3-fr}, y el corolario resulta de la identificación canónica de los $\mathcal{O}_S$-módulos $(M\otimes_B M')\supertilde$ y $\widetilde{M}\otimes_{\widetilde{B}}\widetilde{M}'$ (donde $M$, $M'$ y $B$ se consideran como $C$-módulos) (\sref[I]{I.1.3.13-fr} y \sref[I]{I.1.6.3-fr}).
\end{proof}

Aplicando en particular \sref{II.1.4.7-fr} al caso en que $X=Y$ y $\mathcal{F}'=\mathcal{O}_{X'}$, se ve que el $\mathcal{A}'$-módulo $\mathcal{A}(f^{\prime*}(\mathcal{F}))$ se identifica con $\mathcal{A}(\mathcal{F})\otimes_{\mathcal{B}}\mathcal{A}'$.
% Fuente francesa: PMIHES_1961__8__5_0, printed pp. 11-12.

\begin{env}[1.4.8]
\label{II.1.4.8-fr}
En particular, cuando $X=X'=Y$ (siendo $X$ afín sobre $S$), se ve que, si $\mathcal{F}$ y $\mathcal{G}$ son dos $\mathcal{O}_X$-módulos cuasicoherentes, se tiene
\[
  \mathcal{A}(\mathcal{F}\otimes_{\mathcal{O}_X}\mathcal{G})
  =\mathcal{A}(\mathcal{F})\otimes_{\mathcal{A}(X)}\mathcal{A}(\mathcal{G})
  \tag{1.4.8.1}\label{II.1.4.8.1-fr}
\]
salvo isomorfismo funtorial canónico.
Si, además, $\mathcal{F}$ admite una presentación finita, resulta de \sref[I]{I.1.6.3-fr} y \sref[I]{I.1.3.12-fr} que
\[
  \mathcal{A}(\mathcal{H}om_X(\mathcal{F},\mathcal{G}))
  =\mathcal{H}om_{\mathcal{A}(X)}
    (\mathcal{A}(\mathcal{F}),\mathcal{A}(\mathcal{G}))
  \tag{1.4.8.2}\label{II.1.4.8.2-fr}
\]
salvo isomorfismo canónico.
\end{env}

\begin{remark}[1.4.9]
\label{II.1.4.9-fr}
Si $X$ y $X'$ son dos preesquemas afines sobre $S$, la suma disjunta $X\amalg X'$ también es afín sobre $S$, pues la suma disjunta de dos esquemas afines es un esquema afín.
\end{remark}

\begin{proposition}[1.4.10]
\label{II.1.4.10-fr}
Sean $S$ un preesquema, $\mathcal{B}$ una $\mathcal{O}_S$-álgebra cuasicoherente y $X=\operatorname{Spec}(\mathcal{B})$.
Para todo haz cuasicoherente de ideales $\mathcal{J}$ de $\mathcal{B}$, $\widetilde{\mathcal{J}}$ es un haz cuasicoherente de ideales de $\mathcal{O}_X$, y el subpreesquema cerrado $Y$ de $X$ definido por $\widetilde{\mathcal{J}}$ es canónicamente isomorfo a $\operatorname{Spec}(\mathcal{B}/\mathcal{J})$.
\end{proposition}

\begin{proof}
De \sref[I]{I.4.2.3-fr} resulta inmediatamente que $Y$ es afín sobre $S$; en virtud de \sref{II.1.3.1-fr}, se puede reducir al caso en que $S$ es afín, y entonces la proposición resulta inmediatamente de \sref[I]{I.4.1.2-fr}.
\end{proof}

También puede expresarse el resultado de \sref{II.1.4.10-fr} diciendo que, si $h:\mathcal{B}\to\mathcal{B}'$ es un homomorfismo \emph{sobreyectivo} de $\mathcal{O}_S$-álgebras cuasicoherentes, ${}^a h:\operatorname{Spec}(\mathcal{B}')\to
\operatorname{Spec}(\mathcal{B})$ es una \emph{inmersión cerrada}.

\oldpage[II]{12}
\begin{proposition}[1.4.11]
\label{II.1.4.11-fr}

Sean $S$ un preesquema, $\mathcal{B}$ una $\mathcal{O}_S$-álgebra cuasicoherente y $X=\operatorname{Spec}(\mathcal{B})$.
Para todo haz cuasicoherente de ideales $\mathcal{K}$ de $\mathcal{O}_S$, se tiene (denotando por $f$ el morfismo estructural $X\to S$)
$f^*(\mathcal{K})\mathcal{O}_X=(\mathcal{K}\mathcal{B})\supertilde$
salvo isomorfismo canónico.
\end{proposition}

\begin{proof}
Como la cuestión es local sobre $S$, se puede reducir al caso en que $S=\operatorname{Spec}(A)$ es afín, y en este caso la proposición no es sino \sref[I]{I.1.6.9-fr}.
\end{proof}
% Fuente francesa: PMIHES_1961__8__5_0, printed p. 12.

\subsection{Cambio del preesquema de base}
\label{II.1.5-fr}

\begin{proposition}[1.5.1]
\label{II.1.5.1-fr}
Sea $X$ un preesquema afín sobre $S$.
Para toda extensión $g:S'\to S$ del preesquema de base, $X'=X_{(S')}=X\times_S S'$ es afín sobre $S'$.
\end{proposition}

\begin{proof}
Si $f'$ es la proyección $X'\to S'$, basta demostrar que $f^{\prime-1}(U')$ es un abierto afín para todo abierto afín $U'$ de $S'$ tal que $g(U')$ esté contenido en un abierto afín $U$ de $S$ \sref{II.1.2.1-fr}; se puede, por tanto, reducir al caso en que $S$ y $S'$ son afines, y basta demostrar que entonces $X'$ es un esquema afín \sref{II.1.3.4-fr}.
Pero entonces $X$ es un esquema afín \sref{II.1.3.4-fr}, y, si $A$, $A'$ y $B$ son respectivamente los anillos de $S$, $S'$ y $X$, se sabe que $X'$ es el esquema afín de anillo $A'\otimes_A B$ \sref[I]{I.3.2.2-fr}.
\end{proof}

\begin{corollary}[1.5.2]
\label{II.1.5.2-fr}
Bajo las hipótesis de \sref{II.1.5.1-fr}, sean $f:X\to S$ el morfismo estructural, y $f':X'\to S'$ y $g':X'\to X$ las proyecciones, de modo que el diagrama
\[
  \xymatrix{
    X\ar[d]_{f} & X'\ar[l]_{g'}\ar[d]^{f'}\\
    S & S'\ar[l]^g
  }
\]
es conmutativo.
Para todo $\mathcal{O}_X$-módulo cuasicoherente $\mathcal{F}$, existe un isomorfismo canónico de $\mathcal{O}_{S'}$-módulos
\[
  u:g^*(f_*(\mathcal{F}))\xrightarrow{\sim}
  f'_*(g^{\prime*}(\mathcal{F})).
  \tag{1.5.2.1}\label{II.1.5.2.1-fr}
\]
En particular, existe un isomorfismo canónico de $\mathcal{A}(X')$ sobre $g^*(\mathcal{A}(X))$.
\end{corollary}

\begin{proof}
Para definir $u$, basta definir un homomorfismo
\[
  v:f_*(\mathcal{F})\longrightarrow
  g_*(f'_*(g^{\prime*}(\mathcal{F})))
  =f_*(g'_*(g^{\prime*}(\mathcal{F})))
\]
y tomar $u=v^\sharp$ \sref[0]{0.4.4.3-fr}.
Tomamos $v=f_*(\rho)$, donde $\rho$ es el homomorfismo canónico $\mathcal{F}\to g'_*(g^{\prime*}(\mathcal{F}))$ \sref[0]{0.4.4.3-fr}.
Para demostrar que $u$ es un isomorfismo, se puede reducir al caso en que $S$ y $S'$ son afines y, por tanto, también $X$ y $X'$; con las notaciones de \sref{II.1.5.1-fr}, se tiene entonces $\mathcal{F}=\widetilde{M}$, donde $M$ es un $B$-módulo.
Se observa inmediatamente que $g^*(f_*(\mathcal{F}))$ y $f'_*(g^{\prime*}(\mathcal{F}))$ son ambos iguales al $\mathcal{O}_{S'}$-módulo asociado al $A'$-módulo $A'\otimes_A M$ (donde $M$ se considera como $A$-módulo), y que $u$ es el homomorfismo asociado a la identidad (\sref[I]{I.1.6.3-fr}, \sref[I]{I.1.6.5-fr} y \sref[I]{I.1.6.7-fr}).
\end{proof}
% Fuente francesa: PMIHES_1961__8__5_0, printed pp. 12-13.

\begin{remark}[1.5.3]
\label{II.1.5.3-fr}
No debe creerse que \sref{II.1.5.2-fr} siga siendo válido cuando ya no se supone que $X$ es afín sobre $S$, ni siquiera cuando $S'=\operatorname{Spec}(\boldsymbol{k}(s))$ ($s\in S$) y $S'\to S$ es el morfismo canónico \sref[I]{I.2.4.5-fr}, caso en el cual $X'$ no es sino la \emph{fibra $f^{-1}(s)$} \sref[I]{I.3.6.2-fr}.
Dicho de otro modo, cuando $X$ no es afín sobre $S$, la operación
\oldpage[II]{13}
``imagen directa de haces cuasicoherentes'' no conmuta con la operación de ``paso a las fibras''.
Sin embargo, veremos en el capítulo~III \sref[III]{III.4.2.4-fr} un resultado en este sentido, de carácter ``asintótico'', válido para haces \emph{coherentes} sobre $X$ cuando $f$ es propio (\hyperref[II.5.4-fr]{5.4}) y $S$ es noetheriano.
\end{remark}

\begin{corollary}[1.5.4]
\label{II.1.5.4-fr}
Para todo preesquema $X$ afín sobre $S$ y todo $s\in S$, la fibra $f^{-1}(s)$ (donde $f$ denota el morfismo estructural $X\to S$) es un esquema afín.
\end{corollary}

\begin{proof}
Basta aplicar \sref{II.1.5.1-fr} con $S'=\operatorname{Spec}(\boldsymbol{k}(s))$ y utilizar \sref{II.1.3.4-fr}.
\end{proof}

\begin{corollary}[1.5.5]
\label{II.1.5.5-fr}
Sean $X$ un $S$-preesquema y $S'$ un preesquema afín sobre $S$; entonces $X'=X_{(S')}$ es un preesquema afín sobre $X$.
Además, si $f:X\to S$ es el morfismo estructural, existe un isomorfismo canónico de $\mathcal{O}_X$-álgebras
$\mathcal{A}(X')\xrightarrow{\sim}f^*(\mathcal{A}(S'))$,
y, para todo $\mathcal{A}(S')$-módulo cuasicoherente $\mathcal{M}$, un diisomorfismo canónico
$f^*(\mathcal{M})\xrightarrow{\sim}
\mathcal{A}(f^{\prime*}(\widetilde{\mathcal{M}}))$,
donde $f'=f_{(S')}$ denota el morfismo estructural $X'\to S'$.
\end{corollary}

\begin{proof}
Basta intercambiar los papeles de $X$ y $S'$ en \sref{II.1.5.1-fr} y \sref{II.1.5.2-fr}.
\end{proof}
% Fuente francesa: PMIHES_1961__8__5_0, printed pp. 13-14.

\begin{env}[1.5.6]
\label{II.1.5.6-fr}
Sean ahora $S$ y $S'$ dos preesquemas, $q:S'\to S$ un morfismo, $\mathcal{B}$ (resp. $\mathcal{B}'$) una $\mathcal{O}_S$-álgebra cuasicoherente (resp. una $\mathcal{O}_{S'}$-álgebra cuasicoherente), y $u:\mathcal{B}\to\mathcal{B}'$ un $q$-morfismo (es decir, un homomorfismo $\mathcal{B}\to q_*(\mathcal{B}')$ de $\mathcal{O}_S$-álgebras).
Si $X=\operatorname{Spec}(\mathcal{B})$ y $X'=\operatorname{Spec}(\mathcal{B}')$, se obtiene canónicamente un morfismo
\[
  v=\operatorname{Spec}(u):X'\to X
\]
tal que el diagrama
\[
  \xymatrix{
    X'\ar[r]^v\ar[d] & X\ar[d]\\
    S'\ar[r]_q & S
  }
  \tag{1.5.6.1}\label{II.1.5.6.1-fr}
\]
es conmutativo (siendo las flechas verticales los morfismos estructurales).

En efecto, darse $u$ equivale a darse un homomorfismo de $\mathcal{O}_{S'}$-álgebras cuasicoherentes $u^\sharp:q^*(\mathcal{B})\to\mathcal{B}'$ \sref[0]{0.4.4.3-fr}; este define, por tanto, canónicamente un $S'$-morfismo
\[
  w:\operatorname{Spec}(\mathcal{B}')\to
  \operatorname{Spec}(q^*(\mathcal{B}))
\]
tal que $\mathcal{A}(w)=u^\sharp$ \sref{II.1.2.7-fr}.
Por otra parte, de \sref{II.1.5.2-fr} resulta que $\operatorname{Spec}(q^*(\mathcal{B}))$ se identifica canónicamente con $X\times_S S'$; el morfismo $v$ es la composición
$X'\xrightarrow{w}X\times_S S'\xrightarrow{p_1}X$
de $w$ con la primera proyección, y la conmutatividad de \sref{II.1.5.6.1-fr} resulta de las definiciones.

Sean $U$ (resp. $U'$) un abierto afín de $S$ (resp. de $S'$) tal que $q(U')\subset U$; sean $A=\Gamma(U,\mathcal{O}_S)$ y $A'=\Gamma(U',\mathcal{O}_{S'})$ sus anillos, y $B=\Gamma(U,\mathcal{B})$ y $B'=\Gamma(U',\mathcal{B}')$.
La restricción de $u$, como $(q|U')$-morfismo $\mathcal{B}|U\to\mathcal{B}'|U'$, corresponde a un dihomomorfismo de álgebras $B\to B'$; si $V$ y $V'$ son las imágenes inversas de $U$ y $U'$ en $X$ y $X'$, respectivamente, por los morfismos estructurales, el morfismo $V'\to V$ que es restricción de $v$ corresponde, según \sref[I]{I.1.7.3-fr}, al dihomomorfismo anterior.
\end{env}

\begin{env}[1.5.7]
\label{II.1.5.7-fr}
Bajo las mismas hipótesis que en \sref{II.1.5.6-fr}, sea $\mathcal{M}$ un $\mathcal{B}$-módulo cuasicoherente; existe entonces un isomorfismo canónico de $\mathcal{O}_{X'}$-módulos
\[
  v^*(\widetilde{\mathcal{M}})\xrightarrow{\sim}
  \bigl(q^*(\mathcal{M})\otimes_{q^*(\mathcal{B})}\mathcal{B}'\bigr)^{\sim}
  \tag{1.5.7.1}\label{II.1.5.7.1-fr}
\]

\oldpage[II]{14}
En efecto, el isomorfismo canónico \sref{II.1.5.2.1-fr} da un isomorfismo canónico de $p_1^*(\widetilde{\mathcal{M}})$ sobre el haz de $\operatorname{Spec}(q^*(\mathcal{B}))$ asociado al $q^*(\mathcal{B})$-módulo $q^*(\mathcal{M})$, y basta entonces aplicar \sref{II.1.4.7-fr}.
\end{env}
% Fuente francesa: PMIHES_1961__8__5_0, printed p. 14.

\subsection{Morfismos afines}
\label{II.1.6-fr}

\begin{env}[1.6.1]
\label{II.1.6.1-fr}
Diremos que un morfismo $f:X\to Y$ de preesquemas es \emph{afín} si define $X$ como preesquema afín sobre $Y$.
Las propiedades de los preesquemas afines sobre otro se traducen en este lenguaje del modo siguiente:
\end{env}

\begin{proposition}[1.6.2]
\label{II.1.6.2-fr}
\medskip\noindent
\begin{enumerate}
  \item[{\rm(i)}] Toda inmersión cerrada es afín.
  \item[{\rm(ii)}] La composición de dos morfismos afines es afín.
  \item[{\rm(iii)}] Si $f:X\to Y$ es un $S$-morfismo afín, entonces $f_{(S')}:X_{(S')}\to Y_{(S')}$ es afín para todo cambio de base $S'\to S$.
  \item[{\rm(iv)}] Si $f:X\to Y$ y $f':X'\to Y'$ son dos $S$-morfismos afines, entonces
    \[
      f\times_S f':X\times_S X'\to Y\times_S Y'
    \]
    es afín.
  \item[{\rm(v)}] Si $f:X\to Y$ y $g:Y\to Z$ son dos morfismos tales que $g\circ f$ es afín y $g$ es separado, entonces $f$ es afín.
  \item[{\rm(vi)}] Si $f$ es afín, también lo es $f_{\mathrm{red}}$.
\end{enumerate}
\end{proposition}

\begin{proof}
Por \sref[I]{I.5.5.12-fr}, basta demostrar (i), (ii) y (iii).
Ahora bien, (i) no es sino \sref{II.1.2.2-fr}, y (ii) no es sino \sref{II.1.3.5-fr}; finalmente, (iii) resulta de \sref{II.1.5.1-fr}, pues $X_{(S')}$ se identifica con el producto $X\times_Y Y_{(S')}$ \sref[I]{I.3.3.11-fr}.
\end{proof}

\begin{corollary}[1.6.3]
\label{II.1.6.3-fr}
Si $X$ es un esquema afín e $Y$ es un esquema, todo morfismo $X\to Y$ es afín.
\end{corollary}

\begin{proposition}[1.6.4]
\label{II.1.6.4-fr}
Sean $Y$ un preesquema localmente noetheriano y $f:X\to Y$ un morfismo de tipo finito.
Para que $f$ sea afín, es necesario y suficiente que $f_{\mathrm{red}}$ lo sea.
\end{proposition}

\begin{proof}
Por \sref{II.1.6.2-fr}[vi], solo tenemos que demostrar la suficiencia de la condición.
Basta demostrar que, si $Y$ es afín y noetheriano, entonces $X$ es afín; pero $Y_{\mathrm{red}}$ es entonces afín y, por hipótesis, lo mismo ocurre con $X_{\mathrm{red}}$.
Ahora bien, $X$ es noetheriano, por lo que la conclusión resulta de \sref[I]{I.6.1.7-fr}.
\end{proof}
% Fuente francesa: PMIHES_1961__8__5_0, printed pp. 14-15.

\subsection{Fibrado vectorial asociado a un haz de módulos}
\label{II.1.7-fr}

\begin{env}[1.7.1]
\label{II.1.7.1-fr}
Sean $A$ un anillo y $E$ un $A$-módulo.
Recordemos que el \emph{álgebra simétrica} de $E$, denotada por $\mathbf{S}(E)$ (o $\mathbf{S}_A(E)$), es el cociente del álgebra tensorial $\mathbf{T}(E)$ por el ideal bilateral generado por los elementos $x\otimes y-y\otimes x$, donde $x,y$ recorren $E$.
El álgebra $\mathbf{S}(E)$ se caracteriza por la siguiente propiedad universal: si $\sigma$ es la aplicación canónica $E\to\mathbf{S}(E)$ (obtenida componiendo $E\to\mathbf{T}(E)$ con la aplicación canónica $\mathbf{T}(E)\to\mathbf{S}(E)$), toda aplicación $A$-lineal $E\to B$, donde $B$ es una $A$-álgebra \emph{conmutativa}, se factoriza de manera única como
$E\xrightarrow{\sigma}\mathbf{S}(E)\xrightarrow{g}B$,
donde $g$ es un $A$-homomorfismo \emph{de álgebras}.
De esta caracterización se deduce inmediatamente que, para dos $A$-módulos $E$ y $F$, se tiene
\[
  \mathbf{S}(E\oplus F)=\mathbf{S}(E)\otimes\mathbf{S}(F)
\]
\oldpage[II]{15}
salvo isomorfismo canónico; además, $\mathbf{S}(E)$ es un funtor covariante en $E$, de la categoría de los $A$-módulos a la de las $A$-álgebras conmutativas; finalmente, la caracterización anterior muestra también que, si $E=\varinjlim E_\lambda$, se tiene $\mathbf{S}(E)=\varinjlim\mathbf{S}(E_\lambda)$ salvo isomorfismo canónico.
Por abuso de lenguaje, un producto $\sigma(x_1)\sigma(x_2)\ldots\sigma(x_n)$, donde $x_i\in E$, se denota a menudo por $x_1x_2\ldots x_n$ cuando no puede resultar confusión.
El álgebra $\mathbf{S}(E)$ es \emph{graduada}, siendo $\mathbf{S}_n(E)$ el conjunto de las combinaciones lineales de productos de $n$ elementos de $E$ ($n\geq 0$); el álgebra $\mathbf{S}(A)$ es canónicamente isomorfa al álgebra de polinomios $A[T]$ en una indeterminada, y $\mathbf{S}(A^n)$ es canónicamente isomorfa al álgebra de polinomios en $n$ indeterminadas $A[T_1,\ldots,T_n]$.
\end{env}

\begin{env}[1.7.2]
\label{II.1.7.2-fr}
Sea $\varphi$ un homomorfismo de anillos $A\to B$.
Si $F$ es un $B$-módulo, la aplicación canónica $F\to\mathbf{S}(F)$ da una aplicación canónica $F_{[\varphi]}\to\mathbf{S}(F)_{[\varphi]}$, que se factoriza, por tanto, como
$F_{[\varphi]}\to\mathbf{S}(F_{[\varphi]})\to
\mathbf{S}(F)_{[\varphi]}$
el homomorfismo canónico $\mathbf{S}(F_{[\varphi]})\to\mathbf{S}(F)_{[\varphi]}$ es sobreyectivo, pero no necesariamente biyectivo.
Si $E$ es un $A$-módulo, todo dihomomorfismo $E\to F$ (es decir, todo $A$-homomorfismo $E\to F_{[\varphi]}$) da así canónicamente un $A$-homomorfismo de álgebras
$\mathbf{S}(E)\to\mathbf{S}(F_{[\varphi]})\to
\mathbf{S}(F)_{[\varphi]}$,
es decir, un dihomomorfismo de álgebras $\mathbf{S}(E)\to\mathbf{S}(F)$.

Con las mismas notaciones, para todo $A$-módulo $E$, $\mathbf{S}(E\otimes_A B)$ se identifica canónicamente con el álgebra $\mathbf{S}(E)\otimes_A B$; esto resulta inmediatamente de la propiedad universal de $\mathbf{S}(E)$ \sref{II.1.7.1-fr}.
\end{env}

\begin{env}[1.7.3]
\label{II.1.7.3-fr}
Sea $R$ un subconjunto multiplicativo del anillo $A$.
Aplicando \sref{II.1.7.2-fr} al anillo $B=R^{-1}A$ y recordando que $R^{-1}E=E\otimes_A R^{-1}A$, se obtiene $\mathbf{S}(R^{-1}E)=R^{-1}\mathbf{S}(E)$ salvo isomorfismo canónico.
Además, si $R'\supset R$ es un segundo subconjunto multiplicativo de $A$, el diagrama
\[
  \xymatrix{
    R^{-1}E\ar[r]\ar[d] & {R'}^{-1}E\ar[d]\\
    \mathbf{S}(R^{-1}E)\ar[r] & \mathbf{S}({R'}^{-1}E)
  }
\]
es conmutativo.
\end{env}
% Fuente francesa: PMIHES_1961__8__5_0, printed pp. 15-16.

\begin{env}[1.7.4]
\label{II.1.7.4-fr}
Sea ahora $(S,\mathcal{A})$ un espacio anillado, y sea $\mathcal{E}$ un $\mathcal{A}$-módulo sobre $S$.
Asociando a cada abierto $U\subset S$ la $\Gamma(U,\mathcal{A})$-álgebra $\mathbf{S}(\Gamma(U,\mathcal{E}))$, se obtiene (por la funtorialidad de $\mathbf{S}(E)$, \sref{II.1.7.2-fr}) un prehaz de álgebras.
Su haz asociado, denotado por $\mathbf{S}(\mathcal{E})$ o $\mathbf{S}_{\mathcal{A}}(\mathcal{E})$, se llama \emph{álgebra simétrica sobre $\mathcal{A}$} del $\mathcal{A}$-módulo $\mathcal{E}$.
De \sref{II.1.7.1-fr} resulta inmediatamente que $\mathbf{S}(\mathcal{E})$ es solución de un problema universal: todo homomorfismo de $\mathcal{A}$-módulos $\mathcal{E}\to\mathcal{B}$, donde $\mathcal{B}$ es una $\mathcal{A}$-álgebra conmutativa, se factoriza de manera única como
$\mathcal{E}\to\mathbf{S}(\mathcal{E})\to\mathcal{B}$,
siendo la segunda flecha un homomorfismo de $\mathcal{A}$-álgebras.
Existe, por tanto, una correspondencia biyectiva entre los homomorfismos $\mathcal{E}\to\mathcal{B}$ de $\mathcal{A}$-módulos y los homomorfismos $\mathbf{S}(\mathcal{E})\to\mathcal{B}$ de $\mathcal{A}$-álgebras.
En particular, todo homomorfismo $u:\mathcal{E}\to\mathcal{F}$ de $\mathcal{A}$-módulos define un homomorfismo $\mathbf{S}(u):\mathbf{S}(\mathcal{E})\to\mathbf{S}(\mathcal{F})$ de $\mathcal{A}$-álgebras, y $\mathbf{S}(\mathcal{E})$ es así un funtor covariante en $\mathcal{E}$.

\oldpage[II]{16}
Por \sref{II.1.7.2-fr} y por el hecho de que $\mathbf{S}$ conmuta con los límites inductivos, se tiene $(\mathbf{S}(\mathcal{E}))_s=\mathbf{S}(\mathcal{E}_s)$ para todo punto $s\in S$.
Si $\mathcal{E}$ y $\mathcal{F}$ son dos $\mathcal{A}$-módulos, $\mathbf{S}(\mathcal{E}\oplus\mathcal{F})$ se identifica canónicamente con $\mathbf{S}(\mathcal{E})\otimes_{\mathcal{A}}\mathbf{S}(\mathcal{F})$, como se ve sobre los prehaces correspondientes.

Observemos también que $\mathbf{S}(\mathcal{E})$ es una $\mathcal{A}$-álgebra graduada, suma directa infinita de los $\mathbf{S}_n(\mathcal{E})$, donde el $\mathcal{A}$-módulo $\mathbf{S}_n(\mathcal{E})$ es el haz asociado al prehaz $U\to\mathbf{S}_n(\Gamma(U,\mathcal{E}))$.
Tomando en particular $\mathcal{E}=\mathcal{A}$, se ve que $\mathbf{S}_{\mathcal{A}}(\mathcal{A})$ se identifica con $\mathcal{A}[T]=\mathcal{A}\otimes_{\mathbf{Z}}\mathbf{Z}[T]$ ($T$ una indeterminada, y $\mathbf{Z}$ considerado como haz constante).
\end{env}

\begin{env}[1.7.5]
\label{II.1.7.5-fr}
Sea $(T,\mathcal{B})$ un segundo espacio anillado, y sea $f$ un morfismo $(S,\mathcal{A})\to(T,\mathcal{B})$.
Si $\mathcal{F}$ es un $\mathcal{B}$-módulo, $\mathbf{S}(f^*(\mathcal{F}))$ se identifica canónicamente con $f^*(\mathbf{S}(\mathcal{F}))$; en efecto, si $f=(\psi,\theta)$, se tiene por definición \sref[0]{0.4.3.1-fr}
\[
  \mathbf{S}(f^*(\mathcal{F}))
  =\mathbf{S}(\psi^*(\mathcal{F})
    \otimes_{\psi^*(\mathcal{B})}\mathcal{A})
  =\mathbf{S}(\psi^*(\mathcal{F}))
    \otimes_{\psi^*(\mathcal{B})}\mathcal{A}
\]
\sref{II.1.7.2-fr}.
Para todo abierto $U$ de $S$ y toda sección $h$ de $\mathbf{S}(\psi^*(\mathcal{F}))$ sobre $U$, $h$ coincide, en una vecindad $V$ de cada punto $s\in U$, con un elemento de $\mathbf{S}(\Gamma(V,\psi^*(\mathcal{F})))$; remitiéndose a la definición de $\psi^*(\mathcal{F})$ \sref[0]{0.3.7.1-fr} y teniendo en cuenta que todo elemento de $\mathbf{S}(E)$, para un módulo $E$, es combinación lineal de un número finito de productos de elementos de $E$, se ve que existe una vecindad $W$ de $\psi(s)$ en $T$, una sección $h'$ de $\mathbf{S}(\mathcal{F})$ sobre $W$ y una vecindad $V'\subset V\cap\psi^{-1}(W)$ de $s$ tales que $h$ coincide con $t\to h'(\psi(t))$ sobre $V'$; de aquí resulta la afirmación.
\end{env}

\begin{proposition}[1.7.6]
\label{II.1.7.6-fr}
Sean $A$ un anillo, $S=\operatorname{Spec}(A)$ su espectro primo y $\mathcal{E}=\widetilde{M}$ el $\mathcal{O}_S$-módulo asociado a un $A$-módulo $M$; entonces la $\mathcal{O}_S$-álgebra $\mathbf{S}(\mathcal{E})$ está asociada a la $A$-álgebra $\mathbf{S}(M)$.
\end{proposition}

\begin{proof}
Para todo $f\in A$, $\mathbf{S}(M_f)=(\mathbf{S}(M))_f$ \sref{II.1.7.3-fr}, y la proposición resulta así de la Definición~\sref[I]{I.1.3.4-fr}.
\end{proof}

\begin{corollary}[1.7.7]
\label{II.1.7.7-fr}
Si $S$ es un preesquema y $\mathcal{E}$ un $\mathcal{O}_S$-módulo cuasicoherente, la $\mathcal{O}_S$-álgebra $\mathbf{S}(\mathcal{E})$ es cuasicoherente.
Si, además, $\mathcal{E}$ es de tipo finito, cada uno de los $\mathcal{O}_S$-módulos $\mathbf{S}_n(\mathcal{E})$ es de tipo finito.
\end{corollary}

\begin{proof}
La primera afirmación es consecuencia inmediata de \sref{II.1.7.6-fr} y \sref[I]{I.1.4.1-fr}; la segunda resulta de que, si $E$ es un $A$-módulo de tipo finito, $\mathbf{S}_n(E)$ es un $A$-módulo de tipo finito; se aplica entonces \sref[I]{I.1.3.13-fr}.
\end{proof}
% Fuente francesa: PMIHES_1961__8__5_0, printed pp. 16-17.

\begin{definition}[1.7.8]
\label{II.1.7.8-fr}
Sea $\mathcal{E}$ un $\mathcal{O}_S$-módulo cuasicoherente.
Se llama fibrado vectorial sobre $S$ definido por $\mathcal{E}$, y se denota por $\mathbf{V}(\mathcal{E})$, al espectro \sref{II.1.3.1-fr} de la $\mathcal{O}_S$-álgebra cuasicoherente $\mathbf{S}(\mathcal{E})$.
\end{definition}

Por \sref{II.1.2.7-fr}, para todo $S$-preesquema $X$, existe una correspondencia biyectiva canónica entre los $S$-morfismos $X\to\mathbf{V}(\mathcal{E})$ y los homomorfismos de \emph{$\mathcal{O}_S$-álgebras $\mathbf{S}(\mathcal{E})\to\mathcal{A}(X)$}, y, por tanto, también entre estos $S$-morfismos y los homomorfismos de \emph{$\mathcal{O}_S$-módulos $\mathcal{E}\to\mathcal{A}(X)=f_*(\mathcal{O}_X)$} (donde $f$ es el morfismo estructural $X\to S$).
En particular:

\begin{env}[1.7.9]
\label{II.1.7.9-fr}
Tomemos como $X$ el subpreesquema de $S$ inducido sobre un \emph{abierto $U\subset S$}.
Entonces los $S$-morfismos $U\to\mathbf{V}(\mathcal{E})$ no son sino las $U$-secciones \sref[I]{I.2.5.5-fr} del $U$-preesquema inducido por $\mathbf{V}(\mathcal{E})$ sobre el abierto $p^{-1}(U)$ (donde $p$ es el morfismo estructural $\mathbf{V}(\mathcal{E})\to S$).
Por lo que se acaba de ver, estas $U$-secciones corresponden biyectivamente a los homomorfismos de $\mathcal{O}_S$-módulos $\mathcal{E}\to j_*(\mathcal{O}_S|U)$ (donde $j$ es la inyección canónica $U\to S$), o,
\oldpage[II]{17}
equivalentemente, por \sref[0]{0.4.4.3-fr}, a los homomorfismos de $(\mathcal{O}_S|U)$-módulos $j^*(\mathcal{E})=\mathcal{E}|U\to\mathcal{O}_S|U$.
Además, es inmediato que la restricción a un abierto $U'\subset U$ de un $S$-morfismo $U\to\mathbf{V}(\mathcal{E})$ corresponde a la restricción a $U'$ del homomorfismo correspondiente $\mathcal{E}|U\to\mathcal{O}_S|U$.
Se concluye que el \emph{haz de gérmenes de $S$-secciones} de $\mathbf{V}(\mathcal{E})$ se identifica canónicamente con el \emph{dual $\check{\mathcal{E}}$} de $\mathcal{E}$.

En particular, tomando $X=U=S$, el homomorfismo \emph{cero} $\mathcal{E}\to\mathcal{O}_S$ corresponde a una $S$-sección canónica de $\mathbf{V}(\mathcal{E})$, llamada \emph{$S$-sección cero} (cf. \sref{II.8.3.3-fr}).
\end{env}

\begin{env}[1.7.10]
\label{II.1.7.10-fr}
Tomemos ahora como $X$ el espectro $\{\xi\}$ de un cuerpo $K$; el morfismo estructural $f:X\to S$ corresponde entonces a un monomorfismo $\boldsymbol{k}(s)\to K$, donde $s=f(\xi)$ \sref[I]{I.2.4.6-fr}; los $S$-morfismos $\{\xi\}\to\mathbf{V}(\mathcal{E})$ no son sino los \emph{puntos geométricos de $\mathbf{V}(\mathcal{E})$ con valores en la extensión $K$ de $\boldsymbol{k}(s)$} \sref[I]{I.3.4.5-fr}, situados sobre los puntos de $p^{-1}(s)$.
El conjunto de estos puntos, que puede llamarse \emph{fibra geométrica racional sobre $K$} de $\mathbf{V}(\mathcal{E})$ \emph{sobre el punto $s$}, se identifica, por \sref{II.1.7.8-fr}, con el conjunto de homomorfismos de $\mathcal{O}_S$-módulos $\mathcal{E}\to f_*(\mathcal{O}_X)$, o equivalentemente, por \sref[0]{0.4.4.3-fr}, con el conjunto de homomorfismos de $\mathcal{O}_X$-módulos $f^*(\mathcal{E})\to\mathcal{O}_X=K$.
Pero, por definición \sref[0]{0.4.3.1-fr}, se tiene
$f^*(\mathcal{E})=\mathcal{E}_s\otimes_{\mathcal{O}_s}K
=\mathcal{E}^s\otimes_{\boldsymbol{k}(s)}K$, poniendo $\mathcal{E}^s=\mathcal{E}_s/\mathfrak{m}_s\mathcal{E}_s$;
la fibra geométrica de $\mathbf{V}(\mathcal{E})$ racional sobre $K$ y situada sobre $s$ se identifica, por tanto, con el \emph{dual} del \emph{$K$-espacio vectorial $\mathcal{E}^s\otimes_{\boldsymbol{k}(s)}K$}; si $\mathcal{E}^s$ o $K$ es de dimensión finita sobre $\boldsymbol{k}(s)$, este dual se identifica también con $(\mathcal{E}^s)\check{}\otimes_{\boldsymbol{k}(s)}K$, denotando por $(\mathcal{E}^s)\check{}$ el dual del $\boldsymbol{k}(s)$-espacio vectorial $\mathcal{E}^s$.
\end{env}

\begin{proposition}[1.7.11]
\label{II.1.7.11-fr}
\medskip\noindent
\begin{enumerate}
  \item[{\rm(i)}] $\mathbf{V}(\mathcal{E})$ es un funtor contravariante en $\mathcal{E}$, de la categoría de los $\mathcal{O}_S$-módulos cuasicoherentes a la categoría de los $S$-esquemas afines.
  \item[{\rm(ii)}] Si $\mathcal{E}$ es un $\mathcal{O}_S$-módulo de tipo finito, $\mathbf{V}(\mathcal{E})$ es de tipo finito sobre $S$.
  \item[{\rm(iii)}] Si $\mathcal{E}$ y $\mathcal{F}$ son dos $\mathcal{O}_S$-módulos cuasicoherentes, $\mathbf{V}(\mathcal{E}\oplus\mathcal{F})$ se identifica canónicamente con $\mathbf{V}(\mathcal{E})\times_S\mathbf{V}(\mathcal{F})$.
  \item[{\rm(iv)}] Sea $g:S'\to S$ un morfismo; para todo $\mathcal{O}_S$-módulo cuasicoherente $\mathcal{E}$, $\mathbf{V}(g^*(\mathcal{E}))$ se identifica canónicamente con $\mathbf{V}(\mathcal{E})_{(S')}
    =\mathbf{V}(\mathcal{E})\times_S S'$.
  \item[{\rm(v)}] Todo homomorfismo sobreyectivo $\mathcal{E}\to\mathcal{F}$ de $\mathcal{O}_S$-módulos cuasicoherentes corresponde a una inmersión cerrada $\mathbf{V}(\mathcal{F})\to\mathbf{V}(\mathcal{E})$.
\end{enumerate}
\end{proposition}

\begin{proof}
(i) es consecuencia inmediata de \sref{II.1.2.7-fr}, teniendo en cuenta que todo homomorfismo de $\mathcal{O}_S$-módulos $\mathcal{E}\to\mathcal{F}$ define canónicamente un homomorfismo de $\mathcal{O}_S$-álgebras $\mathbf{S}(\mathcal{E})\to\mathbf{S}(\mathcal{F})$.
(ii) resulta inmediatamente de la definición \sref[I]{I.6.3.1-fr} y de que, si $E$ es un $A$-módulo de tipo finito, $\mathbf{S}(E)$ es una $A$-álgebra de tipo finito.
Para demostrar (iii), basta partir del isomorfismo canónico $\mathbf{S}(\mathcal{E}\oplus\mathcal{F})
\xrightarrow{\sim}\mathbf{S}(\mathcal{E})
\otimes_{\mathcal{O}_S}\mathbf{S}(\mathcal{F})$ \sref{II.1.7.4-fr} y aplicar \sref{II.1.4.6-fr}.
De manera análoga, para demostrar (iv), basta partir del isomorfismo canónico $\mathbf{S}(g^*(\mathcal{E}))
\xrightarrow{\sim} g^*(\mathbf{S}(\mathcal{E}))$ \sref{II.1.7.5-fr} y aplicar \sref{II.1.5.2-fr}.
Finalmente, para establecer (v), basta observar que, si el homomorfismo $\mathcal{E}\to\mathcal{F}$ es sobreyectivo, también lo es el homomorfismo correspondiente $\mathbf{S}(\mathcal{E})\to\mathbf{S}(\mathcal{F})$ de $\mathcal{O}_S$-álgebras, y la conclusión resulta de \sref{II.1.4.10-fr}.
\end{proof}
% Fuente francesa: PMIHES_1961__8__5_0, printed pp. 17-19.

\begin{env}[1.7.12]
\label{II.1.7.12-fr}
Tomemos en particular $\mathcal{E}=\mathcal{O}_S$; el preesquema $\mathbf{V}(\mathcal{O}_S)$ es el $S$-esquema afín que es espectro de la $\mathcal{O}_S$-álgebra $\mathbf{S}(\mathcal{O}_S)$, la cual se identifica con la $\mathcal{O}_S$-álgebra $\mathcal{O}_S[T]=\mathcal{O}_S\otimes_{\mathbf{Z}}\mathbf{Z}[T]$
\oldpage[II]{18}
($T$ una indeterminada); esto es evidente cuando $S=\operatorname{Spec}(\mathbf{Z})$, en virtud de \sref{II.1.7.6-fr}, y de ahí se pasa al caso general considerando el morfismo estructural $S\to\operatorname{Spec}(\mathbf{Z})$ y utilizando \sref{II.1.7.11-fr}[iv].
En razón de este resultado, se escribe $\mathbf{V}(\mathcal{O}_S)=S[T]$, y se tiene así la fórmula
\[
  S[T]=S\otimes_{\mathbf{Z}}\mathbf{Z}[T].
  \tag{1.7.12.1}\label{II.1.7.12.1-fr}
\]

La identificación del haz de gérmenes de $S$-secciones de $S[T]$ con $\mathcal{O}_S$, ya vista en \sref[I]{I.3.3.15-fr}, reaparece aquí en un contexto más general como caso particular de \sref{II.1.7.9-fr}.
\end{env}

\begin{env}[1.7.13]
\label{II.1.7.13-fr}
Para todo $S$-preesquema $X$, se ha visto en \sref{II.1.7.8-fr} que $\operatorname{Hom}_S(X,S[T])$ se identifica canónicamente con $\operatorname{Hom}_{\mathcal{O}_S}(\mathcal{O}_S,\mathcal{A}(X))$, el cual es canónicamente isomorfo a $\Gamma(S,\mathcal{A}(X))$ y está, por tanto, provisto de una estructura de anillo; además, todo $S$-morfismo $h:X\to Y$ induce un homomorfismo de anillos
$\Gamma(\mathcal{A}(h)):\Gamma(S,\mathcal{A}(Y))
\to\Gamma(S,\mathcal{A}(X))$
\sref{II.1.1.2-fr}.
Con la estructura de anillo que acaba de definirse sobre $\operatorname{Hom}_S(X,S[T])$, puede considerarse, por tanto, $\operatorname{Hom}_S(X,S[T])$ como un \emph{funtor contravariante} en $X$, de la categoría de los $S$-preesquemas a la de los anillos.
Por otra parte, $\operatorname{Hom}_S(X,\mathbf{V}(\mathcal{E}))$ se identifica asimismo con $\operatorname{Hom}_{\mathcal{O}_S}(\mathcal{E},\mathcal{A}(X))$ (considerando $\mathcal{A}(X)$ como \emph{$\mathcal{O}_S$-módulo}); en consecuencia, se le puede dotar canónicamente de una estructura de \emph{módulo} sobre el anillo $\operatorname{Hom}_S(X,S[T])$, y se ve como antes que el par
\[
  \bigl(\operatorname{Hom}_S(X,S[T]),
  \operatorname{Hom}_S(X,\mathbf{V}(\mathcal{E}))\bigr)
\]
es un funtor contravariante en $X$, con valores en la categoría cuyos objetos son los pares $(A,M)$ formados por un anillo $A$ y un $A$-módulo $M$, siendo los morfismos los dihomomorfismos.

Interpretaremos estos hechos diciendo que $S[T]$ es un \emph{$S$-esquema de anillos} y que $\mathbf{V}(\mathcal{E})$ es un \emph{$S$-esquema de módulos} sobre el $S$-esquema de anillos $S[T]$ (cf. capítulo~\textbf 0, \textsection8).
\end{env}

\begin{env}[1.7.14]
\label{II.1.7.14-fr}
Vamos a ver que la estructura de $S$-esquema de módulos definida sobre el $S$-esquema $\mathbf{V}(\mathcal{E})$ permite reconstruir el $\mathcal{O}_S$-módulo $\mathcal{E}$ salvo isomorfismo único: para ello, mostraremos que $\mathcal{E}$ es canónicamente isomorfo a un $\mathcal{O}_S$-submódulo de $\mathbf{S}(\mathcal{E})=\mathcal{A}(\mathbf{V}(\mathcal{E}))$, definido mediante esta estructura.
En efecto, por \sref{II.1.7.4-fr}, el conjunto $\operatorname{Hom}_{\mathcal{O}_S}(\mathbf{S}(\mathcal{E}),
\mathcal{A}(X))$ de los homomorfismos de \emph{$\mathcal{O}_S$-álgebras} se identifica canónicamente con $\operatorname{Hom}_{\mathcal{O}_S}(\mathcal{E},\mathcal{A}(X))$, conjunto de los homomorfismos de \emph{$\mathcal{O}_S$-módulos}: si $h$ y $h'$ son dos elementos de este último conjunto, $s_i$ ($1\leq i\leq n$) son secciones de $\mathcal{E}$ sobre un abierto $U\subset S$, y $t$ es una sección de $\mathcal{A}(X)$ sobre $U$, se tiene por definición
\[
  (h+h')(s_1s_2\ldots s_n)
  =\prod_{i=1}^n\bigl(h(s_i)+h'(s_i)\bigr)
\]
y
\[
  (t.h)(s_1s_2\ldots s_n)=t^n\prod_{i=1}^n h(s_i).
\]

Dicho esto, si $z$ es una sección de $\mathbf{S}(\mathcal{E})$ sobre $U$, entonces $h\to h(z)$ es una aplicación de
$\operatorname{Hom}_S(X,\mathbf{V}(\mathcal{E}))
=\operatorname{Hom}_{\mathcal{O}_S}(\mathbf{S}(\mathcal{E}),
\mathcal{A}(X))$
en $\Gamma(U,\mathcal{A}(X))$.
Vamos a
\oldpage[II]{19}
mostrar que $\mathcal{E}$ se identifica con un submódulo de $\mathbf{S}(\mathcal{E})$ tal que \emph{para todo abierto $U\subset S$, toda sección $z$ de este $\mathcal{O}_S$-submódulo sobre $U$ y todo $S$-preesquema $X$, la aplicación $h\to h(z)$ de $\operatorname{Hom}_{\mathcal{O}_S|U}(\mathbf{S}(\mathcal{E})|U,
\mathcal{A}(X)|U)$ en $\Gamma(U,\mathcal{A}(X))$ es un homomorfismo de $\Gamma(U,\mathcal{A}(X))$-módulos}.

Es inmediato que $\mathcal{E}$ posee esta propiedad; para demostrar el recíproco, se puede reducir a probar que, cuando $S=\operatorname{Spec}(A)$ y $\mathcal{E}=\widetilde{M}$, toda sección $z$ de $\mathbf{S}(\mathcal{E})$ sobre $S$ que (para $U=S$) posea la propiedad anterior es necesariamente una sección de $\mathcal{E}$.
Se tiene entonces $z=\sum_{n=0}^{\infty}z_n$, donde $z_n\in\mathbf{S}_n(M)$, y se trata de demostrar que $z_n=0$ para $n\ne1$.
Pongamos $B=\mathbf{S}(M)$ y tomemos como $X$ el preesquema $\operatorname{Spec}(B[T])$, donde $T$ es una indeterminada.
El conjunto $\operatorname{Hom}_{\mathcal{O}_S}(\mathbf{S}(\mathcal{E}),
\mathcal{A}(X))$ se identifica con el conjunto de homomorfismos de anillos $h:B\to B[T]$ \sref[I]{I.1.3.13-fr}, y, por lo visto anteriormente, se tiene
$(T.h)(z)=\sum_{n=0}^{\infty}T^nh(z_n)$
La hipótesis sobre $z$ implica que
$\sum_{n=0}^{\infty}T^nh(z_n)=T.\sum_{n=0}^{\infty}h(z_n)$
para todo homomorfismo $h$.
Tomando en particular como $h$ la inyección canónica, se obtiene
$\sum_{n=0}^{\infty}T^nz_n=T.\sum_{n=0}^{\infty}z_n$,
lo que da la conclusión buscada $z_n=0$ para $n\ne1$.
\end{env}

\begin{proposition}[1.7.15]
\label{II.1.7.15-fr}
Sea $Y$ un preesquema cuyo espacio subyacente es noetheriano, o un esquema cuasicompacto.
Todo $Y$-esquema afín $X$ de tipo finito sobre $Y$ es $Y$-isomorfo a un $Y$-subesquema cerrado de un $Y$-esquema de la forma $\mathbf{V}(\mathcal{E})$, donde $\mathcal{E}$ es un $\mathcal{O}_Y$-módulo cuasicoherente de tipo finito.
\end{proposition}

\begin{proof}
La $\mathcal{O}_Y$-álgebra cuasicoherente $\mathcal{A}(X)$ es de tipo finito \sref{II.1.3.7-fr}.
Las hipótesis implican que $\mathcal{A}(X)$ está generada por un $\mathcal{O}_Y$-submódulo cuasicoherente de tipo finito $\mathcal{E}$ \sref[I]{I.9.6.5-fr}; por definición, esto implica que el homomorfismo canónico $\mathbf{S}(\mathcal{E})\to\mathcal{A}(X)$ que extiende canónicamente la inyección $\mathcal{E}\to\mathcal{A}(X)$ es \emph{sobreyectivo}; la conclusión resulta entonces de \sref{II.1.4.10-fr}.
\end{proof}
% Fuente francesa: PMIHES_1961__8__5_0, printed pp. 19-21.

\section{Espectros primos homogéneos}
\label{section:II.2-fr}

\subsection{Generalidades sobre los anillos y módulos graduados}
\label{subsection:II.2.1-fr}

\begin{notation}[2.1.1]
\label{II.2.1.1-fr}
Dado un anillo $S$ \emph{graduado en grados positivos}, denotaremos por $S_n$ la parte de $S$ formada por los elementos homogéneos de grado $n$ ($n\geq0$), y por $S_+$ la suma (directa) de los $S_n$ para $n>0$; se tiene $1\in S_0$, $S_0$ es un subanillo de $S$, $S_+$ es un ideal graduado de $S$, y $S$ es suma directa de $S_0$ y $S_+$.
Si $M$ es un módulo \emph{graduado} sobre $S$ (con grados positivos o negativos), se denota análogamente por $M_n$ el $S_0$-módulo formado por los elementos homogéneos de grado $n$ de $M$ (con $n\in\mathbf{Z}$).

Para todo entero $d>0$, se denota por $S^{(d)}$ la suma directa de los $S_{nd}$; considerando los elementos de $S_{nd}$ como homogéneos de grado $n$, los $S_{nd}$ definen sobre $S^{(d)}$ una estructura de anillo graduado.

Para todo entero $k$ tal que $0\leq k\leq d-1$, se denota por $M^{(d,k)}$ la suma directa
\oldpage[II]{20}
de los $M_{nd+k}$ ($n\in\mathbf{Z}$); es un $S^{(d)}$-módulo graduado cuando se consideran los elementos de $M_{nd+k}$ como homogéneos de grado $n$.
Se escribe $M^{(d)}$ en lugar de $M^{(d,0)}$.

Con las notaciones anteriores, para todo entero $n$ (positivo o negativo), denotamos por $M(n)$ el $S$-módulo graduado definido por $(M(n))_k=M_{n+k}$ para todo $k\in\mathbf{Z}$.
En particular, $S(n)$ será un $S$-módulo graduado tal que $(S(n))_k=S_{n+k}$, conviniendo en poner $S_n=0$ para $n<0$.
Se dice que un $S$-módulo graduado $M$ es \emph{libre} si es isomorfo, como módulo \emph{graduado}, a una suma directa de módulos de la forma $S(n)$; como $S(n)$ es un $S$-módulo cíclico generado por el elemento $1$ de $S$, considerado como elemento de grado $-n$, esto equivale a decir que $M$ admite una \emph{base} sobre $S$ formada por elementos \emph{homogéneos}.

Se dice que un $S$-módulo graduado $M$ \emph{admite una presentación finita} si existe una sucesión exacta $P\to Q\to M\to0$, donde $P$ y $Q$ son sumas directas finitas de módulos de la forma $S(n)$ y los homomorfismos son de grado $0$ (cf. \sref{II.2.1.2-fr}).
\end{notation}

\begin{env}[2.1.2]
\label{II.2.1.2-fr}
Sean $M$ y $N$ dos $S$-módulos graduados; se define sobre $M\otimes_S N$ una estructura de $S$-módulo \emph{graduado} del modo siguiente.
Sobre el producto tensorial $M\otimes_{\mathbf{Z}}N$ puede definirse una estructura de $\mathbf{Z}$-módulo graduado (donde $\mathbf{Z}$ está graduado por $\mathbf{Z}_0=\mathbf{Z}$, $\mathbf{Z}_n=0$ para $n\ne0$) poniendo
$(M\otimes_{\mathbf{Z}}N)_q=\bigoplus_{m+n=q}M_m\otimes_{\mathbf{Z}}N_n$
Como $M$ y $N$ son respectivamente sumas directas de los $M_m$ y los $N_n$, se sabe que $M\otimes_{\mathbf{Z}}N$ se identifica canónicamente con la suma directa de todos los $M_m\otimes_{\mathbf{Z}}N_n$.
Dicho esto, se tiene $M\otimes_S N=(M\otimes_{\mathbf{Z}}N)/P$, donde $P$ es el $\mathbf{Z}$-submódulo de $M\otimes_{\mathbf{Z}}N$ generado por los elementos $(xs)\otimes y-x\otimes(sy)$ para $x\in M$, $y\in N$ y $s\in S$; es claro que $P$ es un $\mathbf{Z}$-submódulo \emph{graduado} de $M\otimes_{\mathbf{Z}}N$, y se ve inmediatamente que, al pasar al cociente, se obtiene una estructura de $S$-módulo graduado sobre $M\otimes_S N$.

Para dos $S$-módulos graduados $M$ y $N$, recordemos que un homomorfismo $u:M\to N$ de $S$-módulos se dice \emph{de grado $k$} si $u(M_j)\subset N_{j+k}$ para todo $j\in\mathbf{Z}$.
Si $H_n$ denota el conjunto de todos los homomorfismos de grado $n$ de $M$ en $N$, se denota por $\operatorname{Hom}_S(M,N)$ la \emph{suma} (directa) de los $H_n$ ($n\in\mathbf{Z}$) en el $S$-módulo $H$ de todos los homomorfismos (de $S$-módulos) de $M$ en $N$; en general, $\operatorname{Hom}_S(M,N)$ no es igual a este último.
Sin embargo, se tiene $H=\operatorname{Hom}_S(M,N)$ cuando $M$ es \emph{de tipo finito}; en efecto, puede suponerse entonces que $M$ está generado por un número finito de elementos homogéneos $x_i$ ($1\leq i\leq n$), y todo homomorfismo $u\in H$ se escribe de manera única como $\sum_{k\in\mathbf{Z}}u_k$, donde, para cada $k$, $u_k(x_i)$ es la componente homogénea de grado $k+\deg(x_i)$ de $u(x_i)$ ($1\leq i\leq n$), lo que implica que $u_k=0$ salvo para un número finito de índices; por definición, se tiene $u_k\in H_k$, de donde resulta la conclusión.

Los elementos de grado $0$ de $\operatorname{Hom}_S(M,N)$ se llaman \emph{homomorfismos de $S$-módulos graduados}.
Es claro que $S_mH_n\subset H_{m+n}$, de modo que los $H_n$ definen sobre $\operatorname{Hom}_S(M,N)$ una estructura de $S$-módulo graduado.

De estas definiciones resulta inmediatamente que se tiene
\[
  \label{II.2.1.2.1-fr}
  M(m)\otimes_S N(n)=(M\otimes_S N)(m+n),
  \tag{2.1.2.1}
\]
\[
  \label{II.2.1.2.2-fr}
  \operatorname{Hom}_S(M(m),N(n))
  =(\operatorname{Hom}_S(M,N))(n-m)
  \tag{2.1.2.2}
\]
para dos $S$-módulos graduados $M$ y $N$.

\oldpage[II]{21}
Sean $S$ y $S'$ dos anillos graduados; un homomorfismo de \emph{anillos graduados $\varphi:S\to S'$} es un homomorfismo de anillos tal que $\varphi(S_n)\subset S'_n$ para todo $n\in\mathbf{Z}$ (dicho de otro modo, $\varphi$ debe ser un homomorfismo \emph{de grado $0$} de $\mathbf{Z}$-módulos graduados).
Tal homomorfismo dota a $S'$ de una estructura de $S$-módulo \emph{graduado}; provisto de esta estructura y de su estructura de anillo graduado, $S'$ se llama una \emph{$S$-álgebra graduada}.

Si $M$ es también un $S$-módulo graduado, el producto tensorial $M\otimes_S S'$ de $S$-módulos \emph{graduados} está provisto de manera natural de una estructura de $S'$-módulo \emph{graduado}, siendo la graduación la definida anteriormente.
\end{env}
% Fuente francesa: PMIHES_1961__8__5_0, printed p. 21.

\begin{lemma}[2.1.3]
\label{II.2.1.3-fr}
Sea $S$ un anillo graduado en grados positivos.
Para que un subconjunto $E$ de $S_+$ formado por elementos homogéneos genere $S_+$ como $S$-módulo, es necesario y suficiente que $E$ genere $S$ como $S_0$-álgebra.
\end{lemma}

\begin{proof}
La condición es evidentemente suficiente; demostremos que es necesaria.
Sea $E_n$ (resp. $E^n$) el conjunto de los elementos de $E$ de grado igual a $n$ (resp. de grado $\leq n$); basta demostrar, por inducción sobre $n>0$, que $S_n$ es el $S_0$-módulo generado por los elementos de grado $n$ que son productos de elementos de $E^n$.
Esto es evidente para $n=1$ en virtud de la hipótesis; esta muestra además que
$S_n=\sum_{p=0}^{n-1}S_pE_{n-p}$,
y el razonamiento por inducción es entonces inmediato.
\end{proof}

\begin{corollary}[2.1.4]
\label{II.2.1.4-fr}
Para que $S_+$ sea un ideal de tipo finito, es necesario y suficiente que $S$ sea una $S_0$-álgebra de tipo finito.
\end{corollary}

\begin{proof}
Siempre puede suponerse que un sistema finito de generadores de la $S_0$-álgebra $S$ (resp. del $S$-ideal $S_+$) está formado por elementos homogéneos, reemplazando cada uno de los generadores considerados por sus componentes homogéneas.
\end{proof}

\begin{corollary}[2.1.5]
\label{II.2.1.5-fr}
Para que $S$ sea noetheriano, es necesario y suficiente que $S_0$ sea noetheriano y que $S$ sea una $S_0$-álgebra de tipo finito.
\end{corollary}

\begin{proof}
La condición es evidentemente suficiente; es necesaria, pues $S_0$ es isomorfo a $S/S_+$ y $S_+$ debe ser un ideal de tipo finito \sref{II.2.1.4-fr}.
\end{proof}
% Fuente francesa: PMIHES_1961__8__5_0, printed pp. 21-22.

\begin{lemma}[2.1.6]
\label{II.2.1.6-fr}
Sea $S$ un anillo graduado en grados positivos que es una $S_0$-álgebra de tipo finito.
Sea $M$ un $S$-módulo graduado de tipo finito.
Entonces:
\begin{enumerate}
  \item[{\rm(i)}] Los $M_n$ son $S_0$-módulos de tipo finito, y existe un entero $n_0$ tal que $M_n=0$ para $n\leq n_0$.
  \item[{\rm(ii)}] Existen un entero $n_1$ y un entero $h>0$ tales que, para todo entero $n\geq n_1$, se tiene $M_{n+h}=S_hM_n$.
  \item[{\rm(iii)}] Para todo par de enteros $(d,k)$ tal que $d>0$ y $0\leq k\leq d-1$, $M^{(d,k)}$ es un $S^{(d)}$-módulo de tipo finito.
  \item[{\rm(iv)}] Para todo entero $d>0$, $S^{(d)}$ es una $S_0$-álgebra de tipo finito.
  \item[{\rm(v)}] Existe un entero $h>0$ tal que $S_{mh}=(S_h)^m$ para todo $m>0$.
  \item[{\rm(vi)}] Para todo entero $n>0$, existe un entero $m_0$ tal que $S_m\subset S_+^n$ para todo $m\geq m_0$.
\end{enumerate}
\end{lemma}

\begin{proof}
Se puede suponer que $S$ está generado (como $S_0$-álgebra) por elementos homogéneos $f_i$, de grados $h_i$ ($1\leq i\leq r$), y que $M$ está generado (como $S$-módulo) por elementos homogéneos $x_j$ de grados $k_j$ ($1\leq j\leq s$).
Es claro que $M_n$ está formado por combinaciones lineales,
\oldpage[II]{22}
con coeficientes en $S_0$, de elementos $f_1^{\alpha_1}\ldots f_r^{\alpha_r}x_j$ tales que los $\alpha_i$ son enteros $\geq 0$ que satisfacen
$k_j+\sum_i\alpha_i h_i=n$
Para cada $j$ solo hay un número finito de sistemas $(\alpha_i)$ que satisfacen esta ecuación, pues los $h_i$ son $>0$; de aquí resulta la primera afirmación de (i), y la segunda es evidente.

Por otra parte, sea $h$ el mínimo común múltiplo de los $h_i$, y pongamos $g_i=f_i^{h/h_i}$ ($1\leq i\leq r$), de modo que todos los $g_i$ son de grado $h$.
Sean $z_\mu$ los elementos de $M$ de la forma $f_1^{\alpha_1}\ldots f_r^{\alpha_r}x_j$ con $0\leq\alpha_i<h/h_i$ para $1\leq i\leq r$; estos elementos son finitos en número, y sea $n_1$ el mayor de sus grados.
Es claro que, para $n\geq n_1$, todo elemento de $M_{n+h}$ es combinación lineal de los $z_\mu$ cuyos coeficientes son monomios de grado $>0$ en los $g_i$; por tanto, se tiene $M_{n+h}=S_hM_n$, lo que establece (ii).

De igual modo, se ve (para todo $d>0$) que un elemento de $M^{(d,k)}$ es combinación lineal, con coeficientes en $S_0$, de elementos de la forma $g^df_1^{\alpha_1}\ldots f_r^{\alpha_r}x_j$, con $0\leq\alpha_i<d$, siendo $g$ un elemento homogéneo de $S$; de aquí resulta (iii).
(iv) resulta entonces de (iii) y de \sref{II.2.1.3-fr}, tomando $M=S_+$, pues $(S_+)^{(d)}=(S^{(d)})_+$.
La afirmación (v) se deduce de (ii) tomando $M=S$.
Finalmente, para un $n$ dado, solo hay un número finito de sistemas $(\alpha_i)$ tales que $\alpha_i\geq 0$ y $\sum_i\alpha_i<n$; si se elige $m_0$ estrictamente mayor que el mayor valor de $\sum_i\alpha_i h_i$ entre estos sistemas, se tiene $S_m\subset S_+^n$ para $m>m_0$, lo que demuestra (vi).
\end{proof}

\begin{corollary}[2.1.7]
\label{II.2.1.7-fr}
Si $S$ es noetheriano, también lo es $S^{(d)}$ para todo entero $d>0$.
\end{corollary}

\begin{proof}
Esto resulta de \sref{II.2.1.5-fr} y \sref{II.2.1.6-fr}[iv].
\end{proof}
% Fuente francesa: PMIHES_1961__8__5_0, printed pp. 22-23.

\begin{env}[2.1.8]
\label{II.2.1.8-fr}
Sea $\mathfrak{p}$ un ideal primo \emph{graduado} del anillo graduado $S$; $\mathfrak{p}$ es, por tanto, suma directa de los subgrupos $\mathfrak{p}_n=\mathfrak{p}\cap S_n$.
Supongamos que \emph{$\mathfrak{p}$ no contiene $S_+$}.
Entonces, si $f\in S_+$ no pertenece a $\mathfrak{p}$, la relación $f^nx\in\mathfrak{p}$ equivale a $x\in\mathfrak{p}$; en particular, si $f\in S_d$ ($d>0$), para todo $x\in S_{m-nd}$, la relación $f^nx\in\mathfrak{p}_m$ equivale a $x\in\mathfrak{p}_{m-nd}$.
\end{env}

\begin{proposition}[2.1.9]
\label{II.2.1.9-fr}
Sea $n_0$ un entero $>0$; para todo $n\geq n_0$, sea $\mathfrak{p}_n$ un subgrupo de $S_n$.
Para que exista un ideal primo graduado $\mathfrak{p}$ de $S$, que no contenga $S_+$ y tal que $\mathfrak{p}\cap S_n=\mathfrak{p}_n$ para todo $n\geq n_0$, es necesario y suficiente que se verifiquen las condiciones siguientes:
\begin{enumerate}
  \item[{\rm($1^\circ$)}] $S_m\mathfrak{p}_n\subset\mathfrak{p}_{m+n}$ para todo $m\geq 0$ y todo $n\geq n_0$.
  \item[{\rm($2^\circ$)}] Para $m\geq n_0$, $n\geq n_0$, $f\in S_m$ y $g\in S_n$, la relación $fg\in\mathfrak{p}_{m+n}$ implica $f\in\mathfrak{p}_m$ o $g\in\mathfrak{p}_n$.
  \item[{\rm($3^\circ$)}] $\mathfrak{p}_n\neq S_n$ para al menos un $n\geq n_0$.
\end{enumerate}
Además, el ideal primo graduado $\mathfrak{p}$ es entonces único.
\end{proposition}

\begin{proof}
Es evidente que las condiciones $1^\circ$ y $2^\circ$ son necesarias.
Además, si $\mathfrak{p}\not\supset S_+$, existe al menos un $k>0$ tal que $\mathfrak{p}\cap S_k\neq S_k$; si $f\in S_k$ no pertenece a $\mathfrak{p}$, la relación $\mathfrak{p}\cap S_n=S_n$ implica $\mathfrak{p}\cap S_{n-mk}=S_{n-mk}$ según \sref{II.2.1.8-fr}.
Por tanto, si $\mathfrak{p}\cap S_n=S_n$ para todo $n$ suficientemente grande, se tendría $\mathfrak{p}\supset S_+$, contra la hipótesis, lo que demuestra que $3^\circ$ es necesaria.

Recíprocamente, supongamos satisfechas las condiciones $1^\circ$, $2^\circ$ y $3^\circ$.
Observemos que, si para un entero $d\geq n_0$, $f\in S_d$ no pertenece a $\mathfrak{p}_d$, entonces, si $\mathfrak{p}$ existe, $\mathfrak{p}_m$, para $m<n_0$, es necesariamente igual al conjunto de los $x\in S_m$ tales que $f^rx\in\mathfrak{p}_{m+rd}$ para todos salvo un número finito de valores de $r$.
Esto demuestra ya que, si $\mathfrak{p}$ existe, es único.
Queda por demostrar que, si se definen los $\mathfrak{p}_m$ para $m<n_0$ mediante la condición anterior, entonces $\mathfrak{p}=\sum_{n=0}^{\infty}\mathfrak{p}_n$ es un ideal primo.
Observemos primero que, en virtud de $2^\circ$, para $m\geq n_0$, $\mathfrak{p}_m$ se define también como el conjunto de los $x\in S_m$ tales que $f^rx\in\mathfrak{p}_{m+rd}$ para todos salvo un número finito de valores de $r$.
Dicho
\oldpage[II]{23}
esto, si $g\in S_m$ y $x\in\mathfrak{p}_n$, se tiene $f^rgx\in\mathfrak{p}_{m+n+rd}$ para todos salvo un número finito de valores de $r$, y por tanto $gx\in\mathfrak{p}_{m+n}$, lo que demuestra que $\mathfrak{p}$ es un ideal de $S$.
Para establecer que este ideal es primo, es decir, que el anillo $S/\mathfrak{p}$, graduado por los subgrupos $S_n/\mathfrak{p}_n$, es un dominio de integridad, basta (considerando las componentes de mayor grado de dos elementos de $S/\mathfrak{p}$) demostrar que, si $x\in S_m$ e $y\in S_n$ son tales que $x\notin\mathfrak{p}_m$ e $y\notin\mathfrak{p}_n$, entonces $xy\notin\mathfrak{p}_{m+n}$.
En caso contrario, para $r$ suficientemente grande se tendría $f^{2r}xy\in\mathfrak{p}_{m+n+2rd}$; pero $f^ry\notin\mathfrak{p}_{n+rd}$ para todo $r>0$.
Resulta entonces de $2^\circ$ que, para todos salvo un número finito de valores de $r$, se tiene $f^rx\in\mathfrak{p}_{m+rd}$, y se concluye que $x\in\mathfrak{p}_m$, contra la hipótesis.
\end{proof}
% Fuente francesa: PMIHES_1961__8__5_0, printed p. 23.

\begin{env}[2.1.10]
\label{II.2.1.10-fr}
Diremos que un subconjunto $\mathfrak{J}$ de $S_+$ es un \emph{ideal de $S_+$} si es un ideal de $S$, y que $\mathfrak{J}$ es un \emph{ideal primo graduado de $S_+$} si es la intersección de $S_+$ con un ideal primo graduado de $S$ \emph{que no contiene $S_+$} (este ideal primo es también único según \sref{II.2.1.9-fr}).
Si $\mathfrak{J}$ es un ideal de $S_+$, el \emph{radical de $\mathfrak{J}$ en $S_+$} es el conjunto de los elementos de $S_+$ que tienen alguna potencia en $\mathfrak{J}$, es decir, el conjunto
$\mathfrak{r}_+(\mathfrak{J})=\mathfrak{r}(\mathfrak{J})\cap S_+$
En particular, el radical de $0$ en $S_+$ se llama \emph{nilradical} de $S_+$ y se denota por $\mathfrak{N}_+$: es el conjunto de los elementos nilpotentes de $S_+$.
Si $\mathfrak{J}$ es un ideal \emph{graduado} de $S_+$, su radical $\mathfrak{r}_+(\mathfrak{J})$ es un ideal \emph{graduado}: pasando al anillo cociente $S/\mathfrak{J}$, se puede reducir al caso $\mathfrak{J}=0$, y queda por ver que, si $x=x_h+x_{h+1}+\ldots+x_k$ es nilpotente, también lo son los $x_i\in S_i$ ($1\leq h\leq i\leq k$).
Puede suponerse $x_k\neq 0$, y la componente de mayor grado de $x^n$ es entonces $x_k^n$; por tanto, $x_k$ es nilpotente, y se razona después por inducción sobre $k$.
Se dice que el anillo graduado $S$ es \emph{esencialmente reducido} si $\mathfrak{N}_+=0$, es decir, si $S_+$ no contiene elementos nilpotentes $\neq 0$.
\end{env}

\begin{env}[2.1.11]
\label{II.2.1.11-fr}
Se observa que, si en el anillo graduado $S$ un elemento $x$ es divisor de $0$, también lo es su componente de mayor grado.
Se dice que un anillo $S$ es \emph{esencialmente íntegro} si el anillo $S_+$ (\emph{sin elemento unidad}) no contiene divisores de $0$ y es $\neq 0$; para ello basta que un elemento homogéneo $\neq 0$ de $S_+$ no sea divisor de $0$ en este anillo.
Es claro que, si $\mathfrak{p}$ es un ideal primo graduado de $S_+$, entonces $S/\mathfrak{p}$ es esencialmente íntegro.

Sea $S$ un anillo graduado esencialmente íntegro, y sea $x_0\in S_0$.
Si existe entonces \emph{un} elemento homogéneo $f\neq 0$ de $S_+$ tal que $x_0f=0$, se tiene $x_0S_+=0$, pues $(x_0g)f=(x_0f)g=0$ para todo $g\in S_+$, y la hipótesis implica así $x_0g=0$.
Para que $S$ sea íntegro, es necesario y suficiente que $S_0$ sea íntegro y que el anulador de $S_+$ en $S_0$ sea $0$.
\end{env}
% Fuente francesa: PMIHES_1961__8__5_0, printed pp. 23-24.

\subsection{Anillos de fracciones de un anillo graduado}
\label{subsection:II.2.2-fr}

\begin{env}[2.2.1]
\label{II.2.2.1-fr}
Sea $S$ un anillo graduado en grados positivos, y sea $f$ un elemento \emph{homogéneo} de $S$, de grado $d>0$; entonces el anillo de fracciones $S'=S_f$ está graduado, tomando como $S'_n$ el conjunto de los $x/f^k$, donde $x\in S_{n+kd}$ y $k\geq 0$ (se observa aquí que $n$ puede tomar valores negativos arbitrarios); denotamos el subanillo $S'_0=(S_f)_0$ de $S'$ formado por los elementos \emph{de grado $0$} con la notación $S_{(f)}$.

Si $f\in S_d$, los monomios $(f/1)^h$ en $S_f$ ($h$ entero positivo o negativo) forman una \emph{familia linealmente independiente} sobre el anillo $S_{(f)}$, y el conjunto de sus combinaciones lineales no es otro que
\oldpage[II]{24}
el anillo $(S^{(d)})_f$, que es por tanto \emph{isomorfo a $S_{(f)}[T,T^{-1}]=S_{(f)}\otimes_{\mathbf{Z}}\mathbf{Z}[T,T^{-1}]$} (donde $T$ es una indeterminada).
En efecto, si se tiene una relación
$\sum_{h=-a}^{b}z_h(f/1)^h=0$
con $z_h=x_h/f^m$, donde todos los $x_h$ pertenecen a $S_{md}$, esta relación equivale por definición a la existencia de un $k\geq a$ tal que
$\sum_{h=-a}^{b}f^{h+k}x_h=0$,
como los grados de los términos de esta suma son distintos, se tiene $f^{h+k}x_h=0$ para todo $h$, de donde $z_h=0$ para todo $h$.

Si $M$ es un $S$-módulo graduado, $M'=M_f$ es un $S_f$-módulo graduado, siendo $M'_n$ el conjunto de los $z/f^k$ con $z\in M_{n+kd}$ ($k\geq 0$); se denota por $M_{(f)}$ el conjunto de los elementos homogéneos de grado $0$ de $M'$.
Es inmediato que $M_{(f)}$ es un $S_{(f)}$-módulo y que se tiene
$(M^{(d)})_f=M_{(f)}\otimes_{S_{(f)}}(S^{(d)})_f$.
\end{env}
% Fuente francesa: PMIHES_1961__8__5_0, printed p. 24.

\begin{lemma}[2.2.2]
\label{II.2.2.2-fr}
Sean $d$ y $e$ dos enteros $>0$, $f\in S_d$ y $g\in S_e$.
Existe un isomorfismo canónico de anillos
\[
  S_{(fg)}\xrightarrow{\sim}(S_{(f)})_{g^d/f^e}~;
\]
si se identifican canónicamente estos dos anillos, existe un isomorfismo canónico de módulos
\[
  M_{(fg)}\xrightarrow{\sim}(M_{(f)})_{g^d/f^e}.
\]
\end{lemma}

\begin{proof}
En efecto, $fg$ divide a $f^e g^d$, y este último elemento divide a $(fg)^{de}$, por lo que los anillos graduados $S_{fg}$ y $S_{f^e g^d}$ se identifican canónicamente.
Por otra parte, $S_{f^e g^d}$ se identifica también con $(S_{f^e})_{g^d/1}$ \sref[0]{0.1.4.6-fr}; como $f^e/1$ es invertible en $S_{f^e}$, $S_{f^e g^d}$ se identifica asimismo con $(S_{f^e})_{g^d/f^e}$.
Ahora bien, el elemento $g^d/f^e$ es de grado $0$ en $S_{f^e}$; se concluye inmediatamente que el subanillo de $(S_{f^e})_{g^d/f^e}$ formado por los elementos de grado $0$ es $(S_{(f^e)})_{g^d/f^e}$, y como evidentemente se tiene $S_{(f^e)}=S_{(f)}$, esto demuestra la primera parte de la proposición; la segunda se establece del mismo modo.
\end{proof}

\begin{env}[2.2.3]
\label{II.2.2.3-fr}
Bajo las hipótesis de \sref{II.2.2.2-fr}, es claro que el homomorfismo canónico $S_f\to S_{fg}$ \sref[0]{0.1.4.1-fr}, que transforma $x/f^k$ en $g^k x/(fg)^k$, es de grado $0$ y, por restricción, da por tanto un \emph{homomorfismo canónico $S_{(f)}\to S_{(fg)}$}, tal que el diagrama
\[
  \xymatrix{
    & S_{(f)}\ar[dl]\ar[dr]\\
    S_{(fg)}\ar[rr]^-{\sim} && (S_{(f)})_{g^d/f^e}
  }
\]
es conmutativo.
Se define del mismo modo un homomorfismo canónico $M_{(f)}\to M_{(fg)}$.
\end{env}

\begin{lemma}[2.2.4]
\label{II.2.2.4-fr}
Si $f$ y $g$ son dos elementos homogéneos de $S_+$, el anillo $S_{(fg)}$ está generado por la unión de las imágenes canónicas de $S_{(f)}$ y $S_{(g)}$.
\end{lemma}

\begin{proof}
En virtud de \sref{II.2.2.2-fr}, basta ver que
$1/(g^d/f^e)=f^{d+e}/(fg)^d$
pertenece a la imagen canónica de $S_{(g)}$ en $S_{(fg)}$, lo que es evidente por definición.
\end{proof}
% Fuente francesa: PMIHES_1961__8__5_0, printed pp. 24-25.

\begin{proposition}[2.2.5]
\label{II.2.2.5-fr}
Sea $d$ un entero $>0$ y sea $f\in S_d$.
Entonces existe un isomorfismo canónico de anillos
$S_{(f)}\xrightarrow{\sim}S^{(d)}/(f-1)S^{(d)}$
si se identifican estos dos anillos mediante este isomorfismo, existe un isomorfismo canónico de módulos
$M_{(f)}\xrightarrow{\sim}M^{(d)}/(f-1)M^{(d)}$.
\end{proposition}

\begin{proof}
El primero de estos isomorfismos se define haciendo corresponder a $x/f^n$, donde $x\in S_{nd}$, el elemento $\overline{x}$, clase de $x$ módulo $(f-1)S^{(d)}$.
Esta aplicación está bien definida, pues se tiene la congruencia
$f^h x\equiv x$ (mod.~$(f-1)S^{(d)}$)
para todo $x\in S^{(d)}$; por tanto, si $f^h x=0$ para algún $h>0$,
\oldpage[II]{25}
se tiene $\overline{x}=0$.
Por otra parte, si $x\in S_{nd}$ es tal que $x=(f-1)y$, con
$y=y_{hd}+y_{(h+1)d}+\ldots+y_{kd}$
donde $y_{jd}\in S_{jd}$ y $y_{hd}\neq 0$, se tiene necesariamente $h=n$ y $x=-y_{hd}$, así como las relaciones $y_{(j+1)d}=fy_{jd}$ para $h\leq j\leq k-1$ y $fy_{kd}=0$, lo que da finalmente $f^{k-n+1}x=0$.
A toda clase $\overline{x}$ módulo $(f-1)S^{(d)}$ de un elemento $x\in S_{nd}$ puede, por tanto, hacerse corresponder el elemento $x/f^n$ de $S_{(f)}$, pues la observación precedente muestra que esta aplicación está bien definida.
Es inmediato que las dos aplicaciones así definidas son homomorfismos de anillos recíprocos entre sí.
Se procede exactamente del mismo modo para $M$.
\end{proof}

\begin{corollary}[2.2.6]
\label{II.2.2.6-fr}
Si $S$ es noetheriano, también lo es $S_{(f)}$ para todo $f$ homogéneo de grado $>0$.
\end{corollary}

\begin{proof}
Esto resulta inmediatamente de \sref{II.2.1.7-fr} y \sref{II.2.2.5-fr}.
\end{proof}

\begin{env}[2.2.7]
\label{II.2.2.7-fr}
Sea $T$ un subconjunto multiplicativo de $S_+$ formado por elementos \emph{homogéneos}; $T_0=T\cup\{1\}$ es entonces un subconjunto multiplicativo de $S$.
Como los elementos de $T_0$ son homogéneos, el anillo $T_0^{-1}S$ sigue estando graduado de manera evidente; se denota por $S_{(T)}$ el subanillo de $T_0^{-1}S$ formado por los elementos de grado $0$, es decir, los elementos de la forma $x/h$, donde $h\in T$ y $x$ es homogéneo de grado igual al de $h$.
Se sabe por \sref[0]{0.1.4.5-fr} que $T_0^{-1}S$ se identifica canónicamente con el límite inductivo de los anillos $S_f$, donde $f$ recorre $T$ (respecto de los homomorfismos canónicos $S_f\to S_{fg}$); como esta identificación respeta los grados, identifica $S_{(T)}$ con el \emph{límite inductivo} de los $S_{(f)}$ para $f\in T$.
Para todo $S$-módulo graduado $M$, se define análogamente el módulo $M_{(T)}$ (sobre el anillo $S_{(T)}$) formado por los elementos de grado $0$ de $T_0^{-1}M$, y se ve que este módulo es el límite inductivo de los $M_{(f)}$ para $f\in T$.

Si $\mathfrak{p}$ es un ideal primo graduado de $S_+$, se denotan respectivamente por $S_{(\mathfrak{p})}$ y $M_{(\mathfrak{p})}$ el anillo $S_{(T)}$ y el módulo $M_{(T)}$, donde $T$ es el conjunto de los elementos \emph{homogéneos} de $S_+$ que no pertenecen a $\mathfrak{p}$.
\end{env}
% Fuente francesa: PMIHES_1961__8__5_0, printed pp. 25-26.

\subsection{Espectro primo homogéneo de un anillo graduado}
\label{subsection:II.2.3-fr}

\begin{env}[2.3.1]
\label{II.2.3.1-fr}
Dado un anillo $S$ graduado en grados positivos, se llama \emph{espectro primo homogéneo} de $S$, y se denota por $\operatorname{Proj}(S)$, al conjunto de los ideales primos graduados de $S_+$ \sref{II.2.1.10-fr}, o, lo que es equivalente, al conjunto de los ideales primos graduados de $S$ \emph{que no contienen $S_+$}; vamos a definir una estructura de \emph{esquema} cuyo conjunto subyacente es $\operatorname{Proj}(S)$.
\end{env}

\begin{env}[2.3.2]
\label{II.2.3.2-fr}
Para todo subconjunto $E$ de $S$, sea $V_+(E)$ el conjunto de los ideales primos graduados de $S$ que contienen $E$ y no contienen $S_+$; este es, por tanto, el subconjunto $V(E)\cap\operatorname{Proj}(S)$ de $\operatorname{Spec}(S)$.
De \sref[I]{I.1.1.2-fr} se deduce:
\[
  \label{II.2.3.2.1-fr}
  V_+(0)=\operatorname{Proj}(S),\quad
  V_+(S)=V_+(S_+)=\varnothing
  \tag{2.3.2.1}
\]
\[
  \label{II.2.3.2.2-fr}
  V_+\left(\bigcup_\lambda E_\lambda\right)
  =\bigcap_\lambda V_+(E_\lambda)
  \tag{2.3.2.2}
\]
\[
  \label{II.2.3.2.3-fr}
  V_+(EE')=V_+(E)\cup V_+(E').
  \tag{2.3.2.3}
\]

No se modifica $V_+(E)$ al reemplazar $E$ por el ideal graduado generado por $E$; además, si $\mathfrak{J}$ es un ideal graduado de $S$, se tiene
\[
  \label{II.2.3.2.4-fr}
  V_+(\mathfrak{J})
  =V_+\left(\bigcup_{q\geq n}(\mathfrak{J}\cap S_q)\right)
  \tag{2.3.2.4}
\]
\oldpage[II]{26}
para todo $n>0$.
En efecto, si $\mathfrak{p}\in\operatorname{Proj}(S)$ contiene los elementos homogéneos de $\mathfrak{J}$ de grado $\geq n$, como por hipótesis existe un elemento homogéneo $f\in S_d$ que no está contenido en $\mathfrak{p}$, para todo $m\geq 0$ y todo $x\in S_m\cap\mathfrak{J}$ se tiene $f^r x\in\mathfrak{J}\cap S_{m+rd}$ para todo $r\geq0$; para todos salvo un número finito de valores de $r$, $m+rd\geq n$, y por tanto $f^r x\in\mathfrak{p}\cap S_{m+rd}$, lo que implica $x\in\mathfrak{p}\cap S_m$ \sref{II.2.1.8-fr}.

Finalmente, para todo ideal graduado $\mathfrak{J}$ de $S$, se tiene
\[
  \label{II.2.3.2.5-fr}
  V_+(\mathfrak{J})=V_+(\mathfrak{r}_+(\mathfrak{J})).
  \tag{2.3.2.5}
\]
\end{env}

\begin{env}[2.3.3]
\label{II.2.3.3-fr}
Por definición, los $V_+(E)$ son los subconjuntos cerrados de $X=\operatorname{Proj}(S)$ para la topología inducida por la topología espectral de $\operatorname{Spec}(S)$, que se llamará también \emph{topología espectral} sobre $X$.
Se pone, para todo $f\in S$,
\[
  \label{II.2.3.3.1-fr}
  D_+(f)=D(f)\cap\operatorname{Proj}(S)
  =\operatorname{Proj}(S)\setmin V_+(f)
  \tag{2.3.3.1}
\]
y, por consiguiente, para cualesquiera dos elementos $f$, $g$ de $S$ \sref[I]{I.1.1.9.1-fr},
\[
  \label{II.2.3.3.2-fr}
  D_+(fg)=D_+(f)\cap D_+(g).
  \tag{2.3.3.2}
\]
\end{env}

\begin{proposition}[2.3.4]
\label{II.2.3.4-fr}
Cuando $f$ recorre el conjunto de los elementos homogéneos de $S_+$, los $D_+(f)$ forman una base de la topología de $X=\operatorname{Proj}(S)$.
\end{proposition}

\begin{proof}
Resulta, en efecto, de \sref{II.2.3.2.2-fr} y \sref{II.2.3.2.4-fr} que todo subconjunto cerrado de $X$ es intersección de conjuntos de la forma $V_+(f)$, donde $f$ es homogéneo de grado $>0$.
\end{proof}
% Fuente francesa: PMIHES_1961__8__5_0, printed pp. 26-27.

\begin{env}[2.3.5]
\label{II.2.3.5-fr}
Sea $f$ un elemento \emph{homogéneo} de $S_+$, de grado $d>0$.
Para todo ideal primo graduado $\mathfrak{p}$ de $S$ que no contenga $f$, se sabe que el conjunto de los $x/f^n$, donde $x\in\mathfrak{p}$ y $n\geq 0$, es un ideal primo del anillo de fracciones $S_f$ \sref[0]{0.1.2.6-fr}; su intersección con $S_{(f)}$ es, por tanto, un ideal primo de este anillo, que denotamos por $\psi_f(\mathfrak{p})$: es el conjunto de los $x/f^n$ para $n\geq 0$ y $x\in\mathfrak{p}\cap S_{nd}$.
Se ha definido así una aplicación
\[
  \psi_f:D_+(f)\longrightarrow\operatorname{Spec}(S_{(f)})~;
\]
Además, si $g\in S_e$ es otro elemento homogéneo de $S_+$, se tiene un diagrama conmutativo
\[
  \label{II.2.3.5.1-fr}
  \xymatrix{
    D_+(f)\ar[r]^{\psi_f}
    & \operatorname{Spec}(S_{(f)})\\
    D_+(fg)\ar[u]\ar[r]_{\psi_{fg}}
    & \operatorname{Spec}(S_{(fg)})\ar[u]
  }
  \tag{2.3.5.1}
\]
donde la flecha vertical de la izquierda es la inclusión y la de la derecha es la aplicación ${}^a\omega_{fg,f}$ inducida por el homomorfismo canónico $\omega=\omega_{fg,f}:S_{(f)}\to S_{(fg)}$ \sref[I]{I.1.2.1-fr}.
En efecto, si $x/f^n\in\omega^{-1}(\psi_{fg}(\mathfrak{p}))$, con $fg\notin\mathfrak{p}$, entonces, por definición, $g^n x/(fg)^n\in\psi_{fg}(\mathfrak{p})$, y por tanto $g^n x\in\mathfrak{p}$, de donde $x\in\mathfrak{p}$; el recíproco es evidente.
\end{env}

\begin{proposition}[2.3.6]
\label{II.2.3.6-fr}
La aplicación $\psi_f$ es un homeomorfismo de $D_+(f)$ sobre $\operatorname{Spec}(S_{(f)})$.
\end{proposition}

\begin{proof}
En primer lugar, $\psi_f$ es continua; en efecto, si $h\in S_{nd}$, entonces $h/f^n\in\psi_f(\mathfrak{p})$ si y solo si $h\in\mathfrak{p}$, y por tanto $\psi_f^{-1}(D(h/f^n))=D_+(hf)$; nuestra afirmación resulta entonces de \sref{II.2.3.3.2-fr}.
Además, los $D_+(hf)$, donde $h$ recorre los conjuntos $S_{nd}$, forman una base de la topología de $D_+(f)$, en virtud de \sref{II.2.3.4-fr} y \sref{II.2.3.3.2-fr}.
Lo
\oldpage[II]{27}
anterior demuestra, teniendo en cuenta el axioma $(T_0)$, válido en $D_+(f)$ y en $\operatorname{Spec}(S_{(f)})$, que $\psi_f$ es inyectiva y que la aplicación inversa $\psi_f(D_+(f))\longrightarrow D_+(f)$ es continua.
Finalmente, para ver que $\psi_f$ es sobreyectiva, se observa que, si $\mathfrak{q}_0$ es un ideal primo de $S_{(f)}$ y, para todo $n>0$, se denota por $\mathfrak{p}_n$ el conjunto de los $x\in S_n$ tales que $x^d/f^n\in\mathfrak{q}_0$, entonces los $\mathfrak{p}_n$ verifican las condiciones de \sref{II.2.1.9-fr}.
En efecto, si $x\in S_n$, $y\in S_n$ son tales que $x^d/f^n\in\mathfrak{q}_0$ y $y^d/f^n\in\mathfrak{q}_0$, se tiene $(x+y)^{2d}/f^{2n}\in\mathfrak{q}_0$, de donde $(x+y)^d/f^n\in\mathfrak{q}_0$, pues $\mathfrak{q}_0$ es primo.
Esto demuestra que los $\mathfrak{p}_n$ son subgrupos de los $S_n$, y la verificación de las demás condiciones de \sref{II.2.1.9-fr} es inmediata, teniendo en cuenta que $\mathfrak{q}_0$ es primo.
Si $\mathfrak{p}$ es el ideal primo graduado de $S$ definido de este modo, se tiene efectivamente $\psi_f(\mathfrak{p})=\mathfrak{q}_0$, pues, para $x\in S_{nd}$, las relaciones $x/f^n\in\mathfrak{q}_0$ y $x^d/f^{nd}\in\mathfrak{q}_0$ son equivalentes, al ser $\mathfrak{q}_0$ primo.
\end{proof}
% Fuente francesa: PMIHES_1961__8__5_0, printed p. 27.

\begin{corollary}[2.3.7]
\label{II.2.3.7-fr}
Para que $D_+(f)=\varnothing$, es necesario y suficiente que $f$ sea nilpotente.
\end{corollary}

\begin{proof}
Para que $\operatorname{Spec}(S_{(f)})=\varnothing$, es necesario y suficiente que $S_{(f)}=0$, o, lo que es equivalente, que $1=0$ en $S_f$, lo que significa por definición que $f$ es nilpotente.
\end{proof}

\begin{corollary}[2.3.8]
\label{II.2.3.8-fr}
Sea $E$ un subconjunto de $S_+$.
Las condiciones siguientes son equivalentes:
\begin{enumerate}
  \item[\textit{(a)}] $V_+(E)=X=\operatorname{Proj}(S)$.
  \item[\textit{(b)}] Todo elemento de $E$ es nilpotente.
  \item[\textit{(c)}] Las componentes homogéneas de todo elemento de $E$ son nilpotentes.
\end{enumerate}
\end{corollary}

\begin{proof}
Es claro que \emph{(c)} implica \emph{(b)}, y que \emph{(b)} implica \emph{(a)}.
Si $\mathfrak{J}$ es el ideal graduado de $S$ generado por $E$, la condición \emph{(a)} equivale a exigir que $V_+(\mathfrak{J})=X$; \emph{a fortiori}, \emph{(a)} implica que todo elemento homogéneo $f\in\mathfrak{J}$ verifica $V_+(f)=X$, y por tanto $f$ es nilpotente en virtud de \sref{II.2.3.7-fr}.
\end{proof}

\begin{corollary}[2.3.9]
\label{II.2.3.9-fr}
Si $\mathfrak{J}$ es un ideal graduado de $S_+$, entonces $\mathfrak{r}_+(\mathfrak{J})$ es la intersección de los ideales primos graduados de $S_+$ que contienen $\mathfrak{J}$.
\end{corollary}

\begin{proof}
Considerando el anillo graduado $S/\mathfrak{J}$, se puede reducir al caso $\mathfrak{J}=0$.
Hay que demostrar que, si $f\in S_+$ no es nilpotente, existe un ideal primo graduado de $S$ que no contiene $f$; pero al menos una de las componentes homogéneas de $f$ no es nilpotente, y puede suponerse, por tanto, que $f$ es homogéneo.
La afirmación resulta entonces de \sref{II.2.3.7-fr}.
\end{proof}

\begin{env}[2.3.10]
\label{II.2.3.10-fr}
Para todo subconjunto $Y$ de $X=\operatorname{Proj}(S)$, sea $\mathfrak{j}_+(Y)$ el conjunto de los $f\in S_+$ tales que $Y\subset V_+(f)$; esto equivale a decir que $\mathfrak{j}_+(Y)=\mathfrak{j}(Y)\cap S_+$.
Entonces $\mathfrak{j}_+(Y)$ es un ideal de $S_+$ igual a su radical en $S_+$.
\end{env}

\begin{proposition}[2.3.11]
\label{II.2.3.11-fr}
\begin{enumerate}
  \item[\rm(i)] Para todo subconjunto $E$ de $S_+$, $\mathfrak{j}_+(V_+(E))$ es el radical en $S_+$ del ideal graduado de $S_+$ generado por $E$.
  \item[\rm(ii)] Para todo subconjunto $Y$ de $X$, $V_+(\mathfrak{j}_+(Y))=\overline{Y}$, clausura de $Y$ en $X$.
\end{enumerate}
\end{proposition}

\begin{proof}
\medskip\noindent
\begin{enumerate}
  \item[\rm(i)] Si $\mathfrak{J}$ es el ideal graduado de $S_+$ generado por $E$, se tiene $V_+(E)=V_+(\mathfrak{J})$, y la afirmación resulta entonces de \sref{II.2.3.9-fr}.
  \item[\rm(ii)] Como $V_+(\mathfrak{J})=\bigcap_{f\in\mathfrak{J}}V_+(f)$, la relación $Y\subset V_+(\mathfrak{J})$ implica $Y\subset V_+(f)$ para todo $f\in\mathfrak{J}$ y, por consiguiente, $\mathfrak{j}_+(Y)\supset\mathfrak{J}$, de donde $V_+(\mathfrak{j}_+(Y))\subset V_+(\mathfrak{J})$; esto demuestra (ii) por la definición de los subconjuntos cerrados.
\end{enumerate}
\end{proof}
% Fuente francesa: PMIHES_1961__8__5_0, printed pp. 27-28.

\begin{corollary}[2.3.12]
\label{II.2.3.12-fr}
Los subconjuntos cerrados $Y$ de $X=\operatorname{Proj}(S)$ se corresponden biyectivamente con los ideales graduados de $S_+$ iguales a su radical en $S_+$ mediante las aplicaciones que invierten las inclusiones $Y\to\mathfrak{j}_+(Y)$ y $\mathfrak{J}\to
V_+(\mathfrak{J})$.
A la unión $Y_1\cup Y_2$ de dos subconjuntos cerrados de $X$ corresponde
\oldpage[II]{28}
$\mathfrak{j}_+(Y_1)\cap\mathfrak{j}_+(Y_2)$, y a la intersección de una familia arbitraria $(Y_\lambda)$ de subconjuntos cerrados corresponde el radical en $S_+$ de la suma de los $\mathfrak{j}_+(Y_\lambda)$.
\end{corollary}

\begin{corollary}[2.3.13]
\label{II.2.3.13-fr}
Sea $\mathfrak{J}$ un ideal graduado de $S_+$.
Para que $V_+(\mathfrak{J})=\varnothing$, es necesario y suficiente que todo elemento de $S_+$ tenga una potencia en $\mathfrak{J}$.
\end{corollary}

El corolario anterior puede expresarse también en cualquiera de las dos formas equivalentes siguientes:

\begin{corollary}[2.3.14]
Sea $(f_\alpha)$ una familia de elementos homogéneos de $S_+$.
Para que los $D_+(f_\alpha)$ formen una cubierta de $X=\operatorname{Proj}(S)$, es necesario y suficiente que todo elemento de $S_+$ tenga alguna potencia en el ideal generado por los $f_\alpha$.
\label{II.2.3.14-fr}
\end{corollary}

\begin{corollary}[2.3.15]
\label{II.2.3.15-fr}
Sea $(f_\alpha)$ una familia de elementos homogéneos de $S_+$, y sea $f$ un elemento de $S_+$.
Entonces son equivalentes:
\begin{enumerate}
  \item[\textit{(a)}] $D_+(f)\subset\bigcup_\alpha D_+(f_\alpha)$;
  \item[\textit{(b)}] $V_+(f)\supset\bigcap_\alpha V_+(f_\alpha)$;
  \item[\textit{(c)}] $f$ tiene alguna potencia en el ideal generado por los $f_\alpha$.
\end{enumerate}
\end{corollary}

\begin{corollary}[2.3.16]
\label{II.2.3.16-fr}
Para que $X=\operatorname{Proj}(S)$ sea vacío, es necesario y suficiente que todo elemento de $S_+$ sea nilpotente.
\end{corollary}

\begin{corollary}[2.3.17]
\label{II.2.3.17-fr}
En la correspondencia biyectiva descrita en \sref{II.2.3.12-fr}, a los subconjuntos cerrados \emph{irreducibles} de $X$ corresponden los ideales graduados \emph{primos} de $S_+$.
\end{corollary}

\begin{proof}
Si $Y=Y_1\cup Y_2$, donde $Y_1$ e $Y_2$ son subconjuntos cerrados propios de $Y$, se tiene
\[
  \mathfrak{j}_+(Y)
  =\mathfrak{j}_+(Y_1)\cap\mathfrak{j}_+(Y_2)
\]
siendo los ideales $\mathfrak{j}_+(Y_1)$ y $\mathfrak{j}_+(Y_2)$ distintos de $\mathfrak{j}_+(Y)$; por tanto, $\mathfrak{j}_+(Y)$ no es primo.
Recíprocamente, si $\mathfrak{J}$ es un ideal graduado no primo de $S_+$, existen dos elementos $f$, $g$ de $S_+$ tales que $f\notin\mathfrak{J}$ y $g\notin\mathfrak{J}$, pero $fg\in\mathfrak{J}$.
Entonces $V_+(f)\not\supset V_+(\mathfrak{J})$ y $V_+(g)\not\supset V_+(\mathfrak{J})$, pero $V_+(\mathfrak{J})\subset V_+(f)\cup V_+(g)$ por \sref{II.2.3.2.3-fr}; se concluye así que $V_+(\mathfrak{J})$ es la unión de los subconjuntos cerrados $V_+(f)\cap V_+(\mathfrak{J})$ y $V_+(g)\cap V_+(\mathfrak{J})$, que son distintos de $V_+(\mathfrak{J})$.
\end{proof}
% Fuente francesa: PMIHES_1961__8__5_0, printed pp. 28-29.

\subsection{La estructura de esquema sobre $\operatorname{Proj}(S)$}
\label{subsection:II.2.4-fr}

\begin{env}[2.4.1]
\label{II.2.4.1-fr}
Sean $f$ y $g$ dos elementos homogéneos de $S_+$; consideremos los esquemas afines
$Y_f=\operatorname{Spec}(S_{(f)})$, $Y_g=\operatorname{Spec}(S_{(g)})$ y $Y_{fg}=\operatorname{Spec}(S_{(fg)})$.
En virtud de \sref{II.2.2.2-fr}, el morfismo
$w_{fg,f}=({}^a\omega_{fg,f},\widetilde{\omega}_{fg,f})$
de $Y_{fg}$ en $Y_f$, correspondiente al homomorfismo canónico $\omega_{fg,f}:S_{(f)}\to S_{(fg)}$, es una \emph{inmersión abierta} \sref[I]{I.1.3.6-fr}.
Mediante el homeomorfismo inverso de $\psi_f:D_+(f)\to Y_f$ \sref{II.2.3.6-fr}, puede transportarse a $D_+(f)$ la estructura de esquema afín de $Y_f$; en virtud de la conmutatividad del diagrama \sref{II.2.3.5.1-fr}, el esquema afín $D_+(fg)$ se identifica así con el esquema inducido sobre el subconjunto abierto $D_+(fg)$ del espacio subyacente al esquema afín $D_+(f)$.
Es claro entonces, teniendo en cuenta \sref{II.2.3.4-fr}, que $X=\operatorname{Proj}(S)$ queda provisto de una estructura de \emph{preesquema} única cuya restricción a cada $D_+(f)$ es el esquema afín que acaba de definirse.
Además:
\end{env}

\begin{proposition}[2.4.2]
\label{II.2.4.2-fr}
El preesquema $\operatorname{Proj}(S)$ es un esquema.
\end{proposition}

\begin{proof}
Basta \sref[I]{I.5.5.6-fr} demostrar que, cualesquiera que sean $f$, $g$ homogéneos en $S_+$, $D_+(f)\cap D_+(g)=D_+(fg)$ es afín y su anillo está generado por las imágenes canónicas de los anillos de $D_+(f)$ y $D_+(g)$.
El primer punto es evidente por definición, y el segundo resulta de \sref{II.2.2.4-fr}.
\end{proof}

\oldpage[II]{29}
Siempre que se hable del espectro primo homogéneo $\operatorname{Proj}(S)$ como de un \emph{esquema}, se entenderá que está provisto de la estructura que acaba de definirse.

\begin{example}[2.4.3]
\label{II.2.4.3-fr}
Tomemos $S=K[T_1,T_2]$, donde $K$ es un cuerpo, $T_1$ y $T_2$ son indeterminadas, y $S$ está graduado por el grado total.
Resulta de \sref{II.2.3.14-fr} que $\operatorname{Proj}(S)$ es la unión de $D_+(T_1)$ y $D_+(T_2)$; se ve inmediatamente que estos esquemas afines son canónicamente isomorfos a $K[T]$, y que $\operatorname{Proj}(S)$ se obtiene pegando estos dos esquemas afines como se describe en \sref[I]{I.2.3.2-fr} (cf. \sref{II.7.4.14-fr}).
\end{example}
% Fuente francesa: PMIHES_1961__8__5_0, printed p. 29.

\begin{proposition}[2.4.4]
\label{II.2.4.4-fr}
Sea $S$ un anillo graduado en grados positivos, y sea $X$ el esquema $\operatorname{Proj}(S)$.
\begin{enumerate}
  \item[\rm(i)] Si $\mathfrak{N}_+$ es el nilradical de $S_+$ \sref{II.2.1.10-fr}, el esquema $X_{\mathrm{red}}$ es canónicamente isomorfo a $\operatorname{Proj}(S/\mathfrak{N}_+)$; en particular, si $S$ es esencialmente reducido, $\operatorname{Proj}(S)$ es reducido.
  \item[\rm(ii)] Supongamos que $S$ es esencialmente reducido. Entonces, para que $X$ sea íntegro, es necesario y suficiente que $S$ sea esencialmente íntegro.
\end{enumerate}
\end{proposition}

\begin{proof}
\medskip\noindent
\begin{enumerate}
  \item[\rm(i)] Sea $\overline{S}$ el anillo graduado $S/\mathfrak{N}_+$, y denotemos por $x\to\overline{x}$ el homomorfismo canónico $S\to\overline{S}$ de grado $0$.
  Para todo $f\in S_d$ ($d>0$), el homomorfismo canónico $S_f\to\overline{S}_{\overline{f}}$ \sref[0]{0.1.5.1-fr} es sobreyectivo y de grado $0$ y, por restricción, da por tanto un homomorfismo sobreyectivo $S_{(f)}\to\overline{S}_{(\overline{f})}$.
  Si se supone $f\notin\mathfrak{N}_+$, se comprueba inmediatamente que $\overline{S}_{(\overline{f})}$ es reducido y que el núcleo del homomorfismo anterior es el nilradical de $S_{(f)}$, es decir,
  $\overline{S}_{(\overline{f})}=(S_{(f)})_{\mathrm{red}}$.
  A este homomorfismo corresponde, por tanto, una inmersión cerrada $D_+(\overline{f})\to D_+(f)$ que identifica $D_+(\overline{f})$ con $(D_+(f))_{\mathrm{red}}$ \sref[I]{I.5.1.2-fr} y que es, en particular, un homeomorfismo de los espacios subyacentes a estos dos esquemas afines.
  Además, si $g\notin\mathfrak{N}_+$ es otro elemento homogéneo de $S_+$, el diagrama
  \[
      \xymatrix{
        S_{(f)}\ar[r]\ar[d]
        & \overline{S}_{(\overline{f})}\ar[d]\\
        S_{(fg)}\ar[r]
        & \overline{S}_{(\overline{f}\overline{g})}
      }
    \]
  es conmutativo.
  Como, además, los $D_+(f)$, para $f$ homogéneo de grado $>0$ con $f\notin\mathfrak{N}_+$, forman una cubierta de $X=\operatorname{Proj}(S)$ \sref{II.2.3.7-fr}, se ve que los morfismos $D_+(\overline{f})\to D_+(f)$ son las restricciones de una inmersión cerrada $\operatorname{Proj}(\overline{S})\to\operatorname{Proj}(S)$ que es un homeomorfismo de los espacios subyacentes; de aquí resulta la conclusión \sref[I]{I.5.1.2-fr}.
  \item[\rm(ii)] Supongamos que $S$ es esencialmente íntegro, es decir, que $(0)$ es un ideal primo graduado de $S_+$ distinto de $S_+$; entonces $X$ es reducido por (i) e irreducible por \sref{II.2.3.17-fr}.
  Recíprocamente, supongamos que $S$ es esencialmente reducido y que $X$ es íntegro.
  Para todo $f\neq0$ homogéneo en $S_+$ se tiene entonces $D_+(f)\neq\varnothing$ \sref{II.2.3.7-fr}; la hipótesis de que $X$ es irreducible implica que $D_+(f)\cap D_+(g)\neq\varnothing$ para $f$, $g$ homogéneos y $\neq0$ en $S_+$.
  Por tanto, $fg\neq0$ en virtud de \sref{II.2.3.3.2-fr}, y se concluye que $S_+$ no tiene divisores de cero, de donde resulta la primera afirmación.
\end{enumerate}
\end{proof}
% Fuente francesa: PMIHES_1961__8__5_0, printed pp. 29-30.

\begin{env}[2.4.5]
\label{II.2.4.5-fr}
Dado un anillo conmutativo $A$, recordemos que se dice que un anillo graduado $S$ es una \emph{$A$-álgebra graduada} si está provisto de una estructura de $A$-álgebra tal que cada uno de sus subgrupos $S_n$ sea un $A$-módulo.
Para ello basta, por lo demás, que $S_0$ sea
\oldpage[II]{30}
una $A$-álgebra; dicho de otro modo, se define sobre $S$ una estructura de $A$-álgebra graduada definiendo una estructura de $A$-álgebra sobre $S_0$ y poniendo, para $\alpha\in A$, $x\in S_n$, $\alpha.x=(\alpha.1)x$.
\end{env}

\begin{proposition}[2.4.6]
\label{II.2.4.6-fr}
Supongamos que $S$ es una $A$-álgebra graduada.
Entonces, sobre $X=\operatorname{Proj}(S)$, el haz estructural $\mathcal{O}_X$ es un haz de $A$-álgebras (considerándose $A$ como haz constante sobre $X$); dicho de otro modo, $X$ es un esquema sobre $\operatorname{Spec}(A)$.
\end{proposition}

\begin{proof}
Basta observar que, para todo $f$ homogéneo en $S_+$, $S_{(f)}$ es un álgebra sobre $A$, y que el diagrama
\[
  \xymatrix{
    S_{(f)}\ar[rr] && S_{(fg)}\\
    & A\ar[ul]\ar[ur] &
  }
\]
es conmutativo para $f$, $g$ homogéneos en $S_+$.
\end{proof}

\begin{proposition}[2.4.7]
\label{II.2.4.7-fr}
Sea $S$ un anillo graduado en grados positivos.
\begin{enumerate}
  \item[\rm(i)] Para todo entero $d>0$, existe un isomorfismo canónico del esquema $\operatorname{Proj}(S)$ al esquema $\operatorname{Proj}(S^{(d)})$.
  \item[\rm(ii)] Sea $S'$ el anillo graduado tal que $S'_0=\mathbf{Z}$ y $S'_n=S_n$ (considerado como $\mathbf{Z}$-módulo) para $n>0$. Existe entonces un isomorfismo canónico del esquema $\operatorname{Proj}(S)$ al esquema $\operatorname{Proj}(S')$.
\end{enumerate}
\end{proposition}

\begin{proof}
\medskip\noindent
\begin{enumerate}
  \item[\rm(i)] Demostremos primero que la aplicación $\mathfrak{p}\to\mathfrak{p}\cap S^{(d)}$ es una biyección de $\operatorname{Proj}(S)$ sobre $\operatorname{Proj}(S^{(d)})$.
  En efecto, supongamos dado un ideal primo graduado $\mathfrak{p}'\in\operatorname{Proj}(S^{(d)})$, y pongamos $\mathfrak{p}_{nd}=\mathfrak{p}'\cap S_{nd}$ ($n\geq0$).
  Para todo $n>0$ que no sea múltiplo de $d$, definamos $\mathfrak{p}_n$ como el conjunto de los $x\in S_n$ tales que $x^d\in\mathfrak{p}_{nd}$.
  Si $x\in\mathfrak{p}_n$, $y\in\mathfrak{p}_n$, se tiene $(x+y)^{2d}\in\mathfrak{p}_{2nd}$ y, por tanto, $(x+y)^d\in\mathfrak{p}_{nd}$, pues $\mathfrak{p}'$ es primo.
  Es inmediato que los $\mathfrak{p}_n$ así definidos para todo $n\geq0$ verifican las condiciones de \sref{II.2.1.9-fr}; existe, por tanto, un único ideal primo $\mathfrak{p}\in\operatorname{Proj}(S)$ tal que $\mathfrak{p}\cap S^{(d)}=\mathfrak{p}'$.
  Como, para todo $f$ homogéneo en $S_+$, se tiene $V_+(f)=V_+(f^d)$ \sref{II.2.3.2.3-fr}, se ve que la biyección anterior es un \emph{homeomorfismo} de espacios topológicos.
  Finalmente, con las mismas notaciones, $S_{(f)}$ y $S_{(f^d)}$ se identifican canónicamente \sref{II.2.2.2-fr}; por consiguiente, $\operatorname{Proj}(S)$ y $\operatorname{Proj}(S^{(d)})$ se identifican canónicamente como \emph{esquemas}.
  \item[\rm(ii)] Si a cada $\mathfrak{p}\in\operatorname{Proj}(S)$ se asocia el único ideal primo $\mathfrak{p}'\in\operatorname{Proj}(S')$ tal que $\mathfrak{p}'\cap S_n=\mathfrak{p}\cap S_n$ para todo $n>0$ \sref{II.2.1.9-fr}, es claro que se ha definido un homeomorfismo canónico $\operatorname{Proj}(S)\xrightarrow{\sim}\operatorname{Proj}(S')$ de los espacios subyacentes, pues $V_+(f)$ es el mismo conjunto para $S$ y $S'$ cuando $f$ es un elemento homogéneo de $S_+$.
  Como, además, $S_{(f)}=S'_{(f)}$, $\operatorname{Proj}(S)$ y $\operatorname{Proj}(S')$ se identifican como \emph{esquemas}.
\end{enumerate}
\end{proof}

\begin{corollary}[2.4.8]
\label{II.2.4.8-fr}
Si $S$ es una $A$-álgebra graduada y $S'_A$ es la $A$-álgebra graduada tal que $(S'_A)_0=A$ y $(S'_A)_n=S_n$ para $n>0$, existe un isomorfismo canónico de $\operatorname{Proj}(S)$ a $\operatorname{Proj}(S'_A)$.
\end{corollary}

\begin{proof}
En efecto, estos dos esquemas son ambos canónicamente isomorfos a $\operatorname{Proj}(S')$, con las notaciones de \sref{II.2.4.7-fr}[(ii)].
\end{proof}
% Fuente francesa: PMIHES_1961__8__5_0, printed pp. 30-31.

\subsection{El haz asociado a un módulo graduado}
\label{subsection:II.2.5-fr}

\begin{env}[2.5.1]
\label{II.2.5.1-fr}
Sea $M$ un $S$-módulo \emph{graduado}.
Para todo $f$ homogéneo en $S_+$, $M_{(f)}$ es entonces un $S_{(f)}$-módulo y le corresponde, por tanto, un haz cuasicoherente asociado $(M_{(f)})\supertilde$ sobre el esquema afín $\operatorname{Spec}(S_{(f)})$, identificado con $D_+(f)$ \sref[I]{I.1.3.4-fr}.
\end{env}

\oldpage[II]{31}
\begin{proposition}[2.5.2]
\label{II.2.5.2-fr}
Existe sobre $X=\operatorname{Proj}(S)$ un único $\mathcal{O}_X$-módulo cuasicoherente $\widetilde{M}$ tal que, para $f$ homogéneo en $S_+$, se tenga $\Gamma(D_+(f),\widetilde{M})=M_{(f)}$, siendo el homomorfismo de restricción $\Gamma(D_+(f),\widetilde{M})\to
\Gamma(D_+(fg),\widetilde{M})$, para $f$, $g$ homogéneos en $S_+$, el homomorfismo canónico $M_{(f)}\to M_{(fg)}$ \sref{II.2.2.3-fr}.
\end{proposition}

\begin{proof}
Supongamos $f\in S_d$ y $g\in S_e$.
Como $D_+(fg)$ se identifica con el espectro primo de $(S_{(f)})_{g^d/f^e}$ en virtud de \sref{II.2.2.2-fr}, la restricción a $D_+(fg)$ del haz $(M_{(f)})\supertilde$ sobre $D_+(f)$ se identifica canónicamente con el haz asociado al módulo $(M_{(f)})_{g^d/f^e}$ \sref[I]{I.1.3.6-fr} y, por tanto, también con $(M_{(fg)})\supertilde$ \sref{II.2.2.2-fr}.
Se concluye así que existe un isomorfismo canónico
\[
  \theta_{g,f}:(M_{(f)})\supertilde|D_+(fg)
  \xrightarrow{\sim}(M_{(g)})\supertilde|D_+(fg)
\]
tal que, si $h$ es un tercer elemento homogéneo de $S_+$, se tiene $\theta_{f,h}=\theta_{f,g}\circ\theta_{g,h}$ en $D_+(fgh)$.
Por consiguiente \sref[0]{0.3.3.1-fr}, existe sobre $X$ un $\mathcal{O}_X$-módulo cuasicoherente $\mathcal{F}$ y, para cada $f$ homogéneo en $S_+$, un isomorfismo $\eta_f$ de $\mathcal{F}|D_+(f)$ sobre $(M_{(f)})\supertilde$ tal que $\theta_{g,f}=\eta_g\circ\eta_f^{-1}$.
Si se considera entonces el haz $\mathcal{G}$ asociado al prehaz (sobre la base de la topología de $X$ formada por los $D_+(f)$) definido por $D_+(f)\to M_{(f)}$, tomando como homomorfismos de restricción los homomorfismos canónicos $M_{(f)}\to M_{(fg)}$, lo anterior demuestra que $\mathcal{F}$ y $\mathcal{G}$ son isomorfos, teniendo en cuenta \sref[I]{I.1.3.7-fr}.
El haz $\mathcal{G}$ se denota por $\widetilde{M}$ y verifica efectivamente las condiciones del enunciado.
En particular, se tiene $\widetilde{S}=\mathcal{O}_X$.
\end{proof}

\begin{definition}[2.5.3]
\label{II.2.5.3-fr}
Se dice que el $\mathcal{O}_X$-módulo cuasicoherente $\widetilde{M}$ definido en \sref{II.2.5.2-fr} está asociado al $S$-módulo graduado $M$.
\end{definition}

Recordemos que los $S$-módulos graduados forman una categoría cuando se restringen los homomorfismos arbitrarios de módulos graduados a los homomorfismos \emph{de grado $0$}.
Con esta convención:

\begin{proposition}[2.5.4]
\label{II.2.5.4-fr}
El funtor $\widetilde{M}$ en $M$ es un funtor covariante aditivo exacto de la categoría de los $S$-módulos graduados en la de los $\mathcal{O}_X$-módulos cuasicoherentes, y conmuta con los límites inductivos y las sumas directas.
\end{proposition}

\begin{proof}
Como estas propiedades son locales, basta demostrarlas para los haces de la forma $\widetilde{M}|D_+(f)=(M_{(f)})\supertilde$.
Ahora bien, el funtor de localización $M_f$ en $M$, el funtor de grado cero $N_0$ en $N$ sobre los $S_f$-módulos graduados, y el funtor de haz asociado $\widetilde{P}$ en $P$ sobre los $S_{(f)}$-módulos poseen estas tres propiedades: exactitud y conmutación con los límites inductivos y las sumas directas \sref[I]{I.1.3.5-fr} \sref[I]{I.1.3.9-fr}; de aquí resulta la proposición.
\end{proof}

Se denota por $\widetilde{u}$ el homomorfismo $\widetilde{M}\to\widetilde{N}$ correspondiente a un homomorfismo de grado $0$, $u:M\to N$.
Se deduce inmediatamente de \sref{II.2.5.4-fr} que los resultados de \sref[I]{I.1.3.9-fr} y \sref[I]{I.1.3.10-fr} siguen siendo válidos para los $S$-módulos graduados y los homomorfismos de grado $0$ (con el sentido dado aquí a $\widetilde{M}$), pues las demostraciones son puramente formales.

\begin{proposition}[2.5.5]
\label{II.2.5.5-fr}
Para todo $\mathfrak{p}\in X=\operatorname{Proj}(S)$, se tiene $\widetilde{M}_{\mathfrak{p}}=M_{(\mathfrak{p})}$.
\end{proposition}

\begin{proof}
Por definición, $\widetilde{M}_{\mathfrak{p}}=
\varinjlim\Gamma(D_+(f),\widetilde{M})$, donde $f$ recorre el conjunto de los elementos homogéneos de $S_+$ tales que $f\notin\mathfrak{p}$.
Como $\Gamma(D_+(f),\widetilde{M})=M_{(f)}$, la proposición resulta de la definición de $M_{(\mathfrak{p})}$ \sref{II.2.2.7-fr}.
\end{proof}
% Fuente francesa: PMIHES_1961__8__5_0, printed p. 32.

\oldpage[II]{32}
En particular, el \emph{anillo local} $(\mathcal{O}_X)_\mathfrak{p}$ no es otro que el anillo $S_{(\mathfrak{p})}$, conjunto de las fracciones $x/f$ donde $f$ es homogéneo en $S_+$ y no pertenece a $\mathfrak{p}$, y $x$ es homogéneo del mismo grado que $f$.

Más particularmente, si $S$ es \emph{esencialmente íntegro}, de modo que $\operatorname{Proj}(S)=X$ es \emph{íntegro} \sref{II.2.4.4-fr}, y si $\xi=(0)$ es el punto genérico de $X$, el \emph{cuerpo de funciones racionales $R(X)=\mathcal{O}_\xi$} es el cuerpo formado por los elementos $fg^{-1}$, donde $f$ y $g$ son homogéneos del mismo grado en $S_+$ y $g\neq0$.

\begin{proposition}[2.5.6]
\label{II.2.5.6-fr}
Si, para todo $z\in M$ y todo $f$ homogéneo en $S_+$, existe alguna potencia de $f$ que anula $z$, entonces $\widetilde{M}=0$.
Esta condición suficiente es también necesaria cuando la $S_0$-álgebra $S$ está generada por el conjunto $S_1$ de los elementos homogéneos de grado $1$.
\end{proposition}

\begin{proof}
La condición $\widetilde{M}=0$ equivale a $M_{(f)}=0$ para todo $f$ homogéneo en $S_+$.
Por otra parte, si $f\in S_d$, decir que $M_{(f)}=0$ significa que, para todo $z\in M$ homogéneo y de grado múltiplo de $d$, existe una potencia $f^n$ tal que $f^nz=0$.
Si $M_{(f)}=0$ para todo $f\in S_1$, la condición del enunciado queda, por tanto, verificada para todo $f\in S_1$; queda \emph{a fortiori} verificada para todo $f$ homogéneo en $S_+$ cuando $S_1$ genera $S$, pues todo elemento homogéneo de $S_+$ es entonces combinación lineal de productos de elementos de $S_1$.
\end{proof}

\begin{proposition}[2.5.7]
\label{II.2.5.7-fr}
Sea $d$ un entero $>0$, y sea $f\in S_d$.
Entonces, para todo $n\in\mathbf{Z}$, el $(\mathcal{O}_X|D_+(f))$-módulo $(S(nd))\supertilde|D_+(f)$ es canónicamente isomorfo a $\mathcal{O}_X|D_+(f)$.
\end{proposition}

\begin{proof}
En efecto, la multiplicación por el elemento invertible $(f/1)^n$ de $S_f$ es una biyección de $S_{(f)}=(S_f)_0$ sobre
$(S_f)_{nd}=(S_f(nd))_0=(S(nd)_f)_0=S(nd)_{(f)}$
Dicho de otro modo, los $S_{(f)}$-módulos $S_{(f)}$ y $S(nd)_{(f)}$ son canónicamente isomorfos, de donde resulta la proposición.
\end{proof}

\begin{corollary}[2.5.8]
\label{II.2.5.8-fr}
Sobre el subconjunto abierto $U=\bigcup_{f\in S_d}D_+(f)$, la restricción del $\mathcal{O}_X$-módulo $(S(nd))\supertilde$ es un $(\mathcal{O}_X|U)$-módulo invertible \sref[0]{0.5.4.1-fr}.
\end{corollary}

\begin{corollary}[2.5.9]
\label{II.2.5.9-fr}
Si el ideal $S_+$ de $S$ está generado por el conjunto $S_1$ de los elementos homogéneos de grado $1$, el $\mathcal{O}_X$-módulo $(S(n))\supertilde$ es invertible para todo $n\in\mathbf{Z}$.
\end{corollary}

\begin{proof}
Basta observar que $X=\bigcup_{f\in S_1}D_+(f)$ en virtud de la hipótesis \sref{II.2.3.14-fr}, y aplicar \sref{II.2.5.8-fr} con $U=X$.
\end{proof}

\begin{env}[2.5.10]
\label{II.2.5.10-fr}
Pondremos, en todo lo que sigue de esta sección,
\[
\label{II.2.5.10.1-fr}
  \mathcal{O}_X(n)=(S(n))\supertilde
\tag{2.5.10.1}
\]
para todo $n\in\mathbf{Z}$ y, para todo subconjunto abierto $U$ de $X$ y todo $(\mathcal{O}_X|U)$-módulo $\mathcal{F}$,
\[
\label{II.2.5.10.2-fr}
  \mathcal{F}(n)=
  \mathcal{F}\otimes_{\mathcal{O}_X|U}(\mathcal{O}_X(n)|U)
\tag{2.5.10.2}
\]
para todo $n\in\mathbf{Z}$.
Si el ideal $S_+$ está generado por $S_1$, el funtor $\mathcal{F}(n)$ es \emph{exacto} en $\mathcal{F}$ para todo $n\in\mathbf{Z}$, pues $\mathcal{O}_X(n)$ es entonces un $\mathcal{O}_X$-módulo \emph{invertible}.
\end{env}
% Fuente francesa: PMIHES_1961__8__5_0, printed pp. 32-34.

\begin{env}[2.5.11]
\label{II.2.5.11-fr}
Sean $M$ y $N$ dos $S$-módulos graduados.
Para todo $f\in S_d$ ($d>0$), se define un homomorfismo canónico funtorial de $S_{(f)}$-módulos
\[
\label{II.2.5.11.1-fr}
  \lambda_f:M_{(f)}\otimes_{S_{(f)}}N_{(f)}
  \longrightarrow(M\otimes_S N)_{(f)}
\tag{2.5.11.1}
\]
\oldpage[II]{33}
componiendo el homomorfismo $M_{(f)}\otimes_{S_{(f)}}N_{(f)}\to M_f\otimes_{S_f}N_f$ (procedente de las inyecciones $M_{(f)}\to M_f$, $N_{(f)}\to N_f$ y $S_{(f)}\to S_f$) con el isomorfismo canónico $M_f\otimes_{S_f}N_f\xrightarrow{\sim}(M\otimes_S N)_f$ \sref[0]{0.1.3.4-fr}, y observando que, por la definición del producto tensorial de dos módulos graduados, este último isomorfismo conserva los grados.
Para $x\in M_{md}$ e $y\in N_{nd}$ ($m\geq0$, $n\geq0$), se tiene, por tanto,
\[
  \lambda_f((x/f^m)\otimes(y/f^n))=(x\otimes y)/f^{m+n}.
\]

Resulta inmediatamente de esta definición que, si $g\in S_e$ ($e>0$), el diagrama
\[
  \xymatrix{
    M_{(f)}\otimes_{S_{(f)}}N_{(f)}\ar[r]^{\lambda_f}\ar[d]
    & (M\otimes_S N)_{(f)}\ar[d]\\
    M_{(fg)}\otimes_{S_{(fg)}}N_{(fg)}\ar[r]_{\lambda_{fg}}
    & (M\otimes_S N)_{(fg)}
  }
\]
(donde la flecha vertical de la derecha es el homomorfismo canónico y la de la izquierda procede de los homomorfismos canónicos) es conmutativo.
Los homomorfismos $\lambda$ inducen, por tanto, un homomorfismo canónico funtorial de $\mathcal{O}_X$-módulos
\[
\label{II.2.5.11.2-fr}
  \lambda:\widetilde{M}\otimes_{\mathcal{O}_X}\widetilde{N}
  \longrightarrow(M\otimes_S N)\supertilde.
\tag{2.5.11.2}
\]

Consideremos, en particular, dos ideales graduados $\mathfrak{J}$, $\mathfrak{K}$ de $S$.
Como $\widetilde{\mathfrak{J}}$ y $\widetilde{\mathfrak{K}}$ son haces de ideales de $\mathcal{O}_X$, se tiene un homomorfismo canónico $\widetilde{\mathfrak{J}}\otimes_{\mathcal{O}_X}
\widetilde{\mathfrak{K}}\to\mathcal{O}_X$; el diagrama
\[
\label{II.2.5.11.3-fr}
  \xymatrix{
    \widetilde{\mathfrak{J}}\otimes_{\mathcal{O}_X}
    \widetilde{\mathfrak{K}}\ar[rr]^\lambda\ar[dr]
    && (\mathfrak{J}\otimes_S\mathfrak{K})\supertilde\ar[dl]\\
    & \mathcal{O}_X &
  }
\tag{2.5.11.3}
\]
es entonces conmutativo.
En efecto, se puede reducir a verificarlo sobre cada subconjunto abierto $D_+(f)$ (para $f$ homogéneo en $S_+$), y esto resulta inmediatamente de la definición \sref{II.2.5.11.1-fr} de $\lambda_f$ y de \sref[I]{I.1.3.13-fr}.

Observemos finalmente que, si $M$, $N$, $P$ son tres $S$-módulos graduados, el diagrama
\[
\label{II.2.5.11.4-fr}
  \xymatrix{
    \widetilde{M}\otimes_{\mathcal{O}_X}\widetilde{N}
    \otimes_{\mathcal{O}_X}\widetilde{P}
    \ar[r]^{\lambda\otimes1}\ar[d]_{1\otimes\lambda}
    & (M\otimes_S N)\supertilde\otimes_{\mathcal{O}_X}\widetilde{P}
    \ar[d]^\lambda\\
    \widetilde{M}\otimes_{\mathcal{O}_X}(N\otimes_S P)\supertilde
    \ar[r]_\lambda
    & (M\otimes_S N\otimes_S P)\supertilde
  }
\tag{2.5.11.4}
\]
\oldpage[II]{34}
es conmutativo.
Basta de nuevo verificarlo sobre cada subconjunto abierto $D_+(f)$, y esto resulta inmediatamente de las definiciones y de \sref[I]{I.1.3.13-fr}.
\end{env}
% Fuente francesa: PMIHES_1961__8__5_0, printed pp. 34-35.

\begin{env}[2.5.12]
\label{II.2.5.12-fr}
Bajo las hipótesis de \sref{II.2.5.11-fr}, se define un homomorfismo canónico funtorial de $S_{(f)}$-módulos
\[
\label{II.2.5.12.1-fr}
  \mu_f:\bigl(\operatorname{Hom}_S(M,N)\bigr)_{(f)}
  \longrightarrow
  \operatorname{Hom}_{S_{(f)}}(M_{(f)},N_{(f)})
\tag{2.5.12.1}
\]
haciendo corresponder a $u/f^n$, donde $u$ es un homomorfismo de grado $nd$, el homomorfismo $M_{(f)}\to N_{(f)}$ que transforma $x/f^m$ ($x\in M_{md}$) en $u(x)/f^{m+n}$.
Para $g\in S_e$ ($e>0$), se tiene de nuevo un diagrama conmutativo:
\[
  \xymatrix{
    \bigl(\operatorname{Hom}_S(M,N)\bigr)_{(f)}
    \ar[r]^{\mu_f}\ar[d]
    & \operatorname{Hom}_{S_{(f)}}(M_{(f)},N_{(f)})\ar[d]\\
    \bigl(\operatorname{Hom}_S(M,N)\bigr)_{(fg)}
    \ar[r]_{\mu_{fg}}
    & \operatorname{Hom}_{S_{(fg)}}(M_{(fg)},N_{(fg)})
  }
\]
(donde la flecha vertical de la izquierda es el homomorfismo canónico y la de la derecha procede de los homomorfismos canónicos).
Se concluye así de nuevo, teniendo en cuenta \sref[I]{I.1.3.8-fr}, que los $\mu_f$ definen un homomorfismo canónico funtorial de $\mathcal{O}_X$-módulos
\[
\label{II.2.5.12.2-fr}
  \mu:(\operatorname{Hom}_S(M,N))\supertilde
  \longrightarrow
  \mathcal{H}om_{\mathcal{O}_X}(\widetilde{M},\widetilde{N}).
\tag{2.5.12.2}
\]
\end{env}

\begin{proposition}[2.5.13]
\label{II.2.5.13-fr}
Supongamos que el ideal $S_+$ está generado por $S_1$.
Entonces, el homomorfismo $\lambda$ \sref{II.2.5.11.2-fr} es un isomorfismo; lo mismo ocurre con el homomorfismo $\mu$ \sref{II.2.5.12.2-fr} cuando el $S$-módulo graduado $M$ admite una presentación finita \sref{II.2.1.1-fr}.
\end{proposition}

\begin{proof}
Como $X$ es la unión de los $D_+(f)$ para $f\in S_1$ \sref{II.2.3.14-fr}, se reduce a demostrar que $\lambda_f$ y $\mu_f$ son isomorfismos, bajo las hipótesis consideradas, cuando $f$ es homogéneo y \emph{de grado $1$}.
Ahora bien, puede definirse entonces una aplicación $\mathbf{Z}$-bilineal $M_m\times N_n\longrightarrow M_{(f)}\otimes_{S_{(f)}}N_{(f)}$ haciendo corresponder a $(x,y)$ el elemento $(x/f^m)\otimes(y/f^n)$ (si $m<0$, escribimos $x/f^m$ para significar $f^{-m}x/1$).
Estas aplicaciones definen una aplicación $\mathbf{Z}$-lineal $M\otimes_{\mathbf{Z}}N\longrightarrow M_{(f)}\otimes_{S_{(f)}}N_{(f)}$ y, si $s\in S_q$, esta aplicación transforma $(sx)\otimes y$ en
$(s/f^q)\bigl((x/f^m)\otimes(y/f^n)\bigr)$
para $x\in M_m$ e $y\in N_n$.
Se obtiene así un dihomomorfismo de módulos
$\gamma_f:M\otimes_S N\longrightarrow M_{(f)}\otimes_{S_{(f)}}N_{(f)}$
relativo al homomorfismo canónico $S\to S_{(f)}$ que transforma $s\in S_q$ en $s/f^q$.
Supongamos además que, para un elemento $\sum_i(x_i\otimes y_i)$ de $M\otimes_S N$, donde $x_i$, $y_i$ son homogéneos de grados respectivos $m_i$, $n_i$, se tiene
$f^r\sum_i(x_i\otimes y_i)=0$,
es decir, $\sum_i(f^rx_i\otimes y_i)=0$.
Se deduce entonces de \sref[0]{0.1.3.4-fr} que
$\sum_i(f^rx_i/f^{m_i+r})\otimes(y_i/f^{n_i})=0$,
es decir, $\gamma_f(\sum_i(x_i\otimes y_i))=0$.
Por consiguiente, $\gamma_f$ se factoriza en
$M\otimes_S N\longrightarrow(M\otimes_S N)_f\xrightarrow{\gamma'_f}M_{(f)}\otimes_{S_{(f)}}N_{(f)}$
Si $\lambda'_f$ es la restricción de $\gamma'_f$
\oldpage[II]{35}
a $(M\otimes_S N)_{(f)}$, se verifica inmediatamente que $\lambda_f$ y $\lambda'_f$ son $S_{(f)}$-homomorfismos recíprocos, de donde resulta la primera parte de la proposición.

Para demostrar la segunda parte, supongamos que $M$ es el conúcleo de un homomorfismo $P\to Q$ de $S$-módulos graduados, siendo $P$ y $Q$ sumas directas de un número finito de módulos de la forma $S(n)$.
Utilizando la exactitud por la izquierda de $\operatorname{Hom}_S(L,N)$ en $L$ y la exactitud de $M_{(f)}$ en $M$, se reduce inmediatamente a demostrar que $\mu_f$ es un isomorfismo cuando $M=S(n)$.
Ahora bien, para todo $z$ homogéneo en $N$, sea $u_z$ el homomorfismo de $S(n)$ en $N$ tal que $u_z(1)=z$; se ve inmediatamente que $\eta:z\to u_z$ es un isomorfismo de grado $0$ de $N(-n)$ sobre $\operatorname{Hom}_S(S(n),N)$.
Le corresponde un isomorfismo
\[
  \eta_f:\bigl(N(-n)\bigr)_{(f)}
  \longrightarrow
  \bigl(\operatorname{Hom}_S(S(n),N)\bigr)_{(f)}.
\]

Por otra parte, sea $\eta'_f$ el isomorfismo $N_{(f)}\to\operatorname{Hom}_{S_{(f)}}(S(n)_{(f)},N_{(f)})$ que hace corresponder a cada $z'\in N_{(f)}$ el homomorfismo $v_{z'}$ tal que
$v_{z'}(s/f^k)=sz'/f^{n+k}$
para $s\in S_{n+k}=(S(n))_k$.
Se comprueba fácilmente que la aplicación compuesta
\[
  \bigl(N(-n)\bigr)_{(f)}
  \xrightarrow{\eta_f}
  \bigl(\operatorname{Hom}_S(S(n),N)\bigr)_{(f)}
  \xrightarrow{\mu_f}
  \operatorname{Hom}_{S_{(f)}}\bigl(S(n)_{(f)},N_{(f)}\bigr)
  \xrightarrow{{\eta'_f}^{-1}}N_{(f)}
\]
es el isomorfismo $z/f^h\to z/f^{h-n}$ de $\bigl(N(-n)\bigr)_{(f)}$ sobre $N_{(f)}$; por tanto, $\mu_f$ es un isomorfismo.
\end{proof}
% Fuente francesa: PMIHES_1961__8__5_0, printed pp. 35-36.

Si el ideal $S_+$ está generado por $S_1$, se deduce de \sref{II.2.5.13-fr} que, para todo ideal graduado $\mathfrak{J}$ de $S$ y todo $S$-módulo graduado $M$, se tiene
\[
\label{II.2.5.13.1-fr}
  \widetilde{\mathfrak{J}}.\widetilde{M}
  =(\mathfrak{J}.M)\supertilde
\tag{2.5.13.1}
\]
salvo isomorfismo canónico.
Esto resulta, en efecto, de la conmutatividad del diagrama
\[
  \xymatrix{
    \widetilde{\mathfrak{J}}\otimes_{\mathcal{O}_X}\widetilde{M}
    \ar[rr]^\lambda\ar[dr]
    && (\mathfrak{J}\otimes_S M)\supertilde\ar[dl]\\
    & \widetilde{M} &
  }
\]
que se verifica como el de \sref{II.2.5.11.3-fr}.

\begin{corollary}[2.5.14]
\label{II.2.5.14-fr}
Supongamos que $S$ está generado por $S_1$.
Para cualesquiera $m$, $n$ en $\mathbf{Z}$, se tiene entonces:
\[
\label{II.2.5.14.1-fr}
  \mathcal{O}_X(m)\otimes_{\mathcal{O}_X}\mathcal{O}_X(n)
  =\mathcal{O}_X(m+n)
\tag{2.5.14.1}
\]
\[
\label{II.2.5.14.2-fr}
  \mathcal{O}_X(n)=\bigl(\mathcal{O}_X(1)\bigr)^{\otimes n}
\tag{2.5.14.2}
\]
salvo isomorfismos canónicos.
\end{corollary}

\begin{proof}
La primera fórmula resulta de \sref{II.2.5.13-fr} y de la existencia del isomorfismo canónico de grado $0$
$S(m)\otimes_S S(n)\xrightarrow{\sim}S(m+n)$,
que transforma el elemento $1\otimes 1$ (donde el primer $1$ está en $(S(m))_{-m}$ y el segundo en $(S(n))_{-n}$) en el elemento $1\in(S(m+n))_{-(m+n)}$.
Basta entonces demostrar la segunda fórmula para $n=-1$ y, en virtud de \sref{II.2.5.13-fr}, esto se reduce a ver que $\operatorname{Hom}_S(S(1),S)$ es canónicamente isomorfo a $S(-1)$, lo que resulta inmediatamente de las definiciones \sref{II.2.1.2-fr} y del hecho de que $S(1)$ es un $S$-módulo cíclico.
\end{proof}

\oldpage[II]{36}
\begin{corollary}[2.5.15]
\label{II.2.5.15-fr}
Supongamos que $S$ está generado por $S_1$.
Entonces, para todo $S$-módulo graduado $M$ y todo $n\in\mathbf{Z}$, se tiene
\[
\label{II.2.5.15.1-fr}
  (M(n))\supertilde=\widetilde{M}(n)
\tag{2.5.15.1}
\]
salvo isomorfismo canónico.
\end{corollary}

\begin{proof}
Esto resulta de las definiciones \sref{II.2.5.10.2-fr} y \sref{II.2.5.10.1-fr}, de la Proposición \sref{II.2.5.13-fr}, y de la existencia de un isomorfismo canónico $M(n)\xrightarrow{\sim}M\otimes_S S(n)$ de grado $0$ que hace corresponder a todo $z\in(M(n))_h=M_{n+h}$ el elemento
$z\otimes1\in M_{n+h}\otimes(S(n))_{-n}\subset(M\otimes_S S(n))_h$.
\end{proof}
% Fuente francesa: PMIHES_1961__8__5_0, printed p. 36.

\begin{env}[2.5.16]
\label{II.2.5.16-fr}
Denotemos por $S'$ el anillo graduado tal que $S'_0=\mathbf{Z}$ y $S'_n=S_n$ para $n>0$.
Entonces, si $f\in S_d$ ($d>0$), se tiene $(S(n))_{(f)}=(S'(n))_{(f)}$ para todo $n\in\mathbf{Z}$, pues un elemento de $(S'(n))_{(f)}$ es de la forma $x/f^k$, con $x\in S'_{n+kd}$ ($k>0$), y siempre puede tomarse $k$ tal que $n+kd\neq0$.
Como $X=\operatorname{Proj}(S)$ y $X'=\operatorname{Proj}(S')$ se identifican canónicamente \sref{II.2.4.7-fr}[(ii)], se ve que, para todo $n\in\mathbf{Z}$, $\mathcal{O}_X(n)$ y $\mathcal{O}_{X'}(n)$ son las imágenes uno de otro por la identificación anterior.

Observemos, además, que, para todo $d>0$ y todo $n\in\mathbf{Z}$, se tiene
\[
  (S^{(d)}(n))_h=S_{(n+h)d}=(S(nd))_{hd}
\]
para $f\in S_d$, y por tanto $(S^{(d)}(n))_{(f)}=(S(nd))_{(f)}$.
Se sabe que los esquemas $X=\operatorname{Proj}(S)$ y $X^{(d)}=\operatorname{Proj}(S^{(d)})$ se identifican canónicamente \sref{II.2.4.7-fr}[(i)].
Lo anterior muestra que, si la $S_0$-álgebra $S^{(d)}$ está generada por $S_d$, entonces $\mathcal{O}_X(nd)$ y $\mathcal{O}_{X^{(d)}}(n)$ son las imágenes uno de otro por esta identificación, para todo $n\in\mathbf{Z}$.
\end{env}

\begin{proposition}[2.5.17]
\label{II.2.5.17-fr}
Sea $d$ un entero $>0$, y sea $U=\bigcup_{f\in S_d}D_+(f)$.
La restricción a $U$ del homomorfismo canónico
$\mathcal{O}_X(nd)\otimes_{\mathcal{O}_X}\mathcal{O}_X(-nd)
\to\mathcal{O}_X$
es un isomorfismo para todo entero $n$.
\end{proposition}

\begin{proof}
En virtud de \sref{II.2.5.16-fr}, se puede limitar al caso $d=1$, y la conclusión resulta entonces de la demostración de \sref{II.2.5.13-fr}.
\end{proof}
% Fuente francesa: PMIHES_1961__8__5_0, printed pp. 36-37.

\subsection{El $S$-módulo graduado asociado a un haz sobre $\operatorname{Proj}(S)$}
\label{subsection:II.2.6-fr}

\emph{En toda esta sección se supone que el ideal $S_+$ de $S$ está generado por el conjunto $S_1$ de los elementos homogéneos de grado $1$.}

\begin{env}[2.6.1]
\label{II.2.6.1-fr}
El $\mathcal{O}_X$-módulo $\mathcal{O}_X(1)$ es \emph{invertible} \sref{II.2.5.9-fr}; se pone entonces, para todo $\mathcal{O}_X$-módulo $\mathcal{F}$ \sref[0]{0.5.4.6-fr},
\[
\label{II.2.6.1.1-fr}
  \Gamma_*(\mathcal{F})
  =\Gamma_*(\mathcal{O}_X(1),\mathcal{F})
  =\bigoplus_{n\in\mathbf{Z}}\Gamma(X,\mathcal{F}(n))
\tag{2.6.1.1}
\]
teniendo en cuenta \sref{II.2.5.14.2-fr}.
Recordemos \sref[0]{0.5.4.6-fr} que $\Gamma_*(\mathcal{O}_X)$ está provisto de una estructura de \emph{anillo graduado}, y $\Gamma_*(\mathcal{F})$ de una estructura de \emph{$\Gamma_*(\mathcal{O}_X)$-módulo graduado}.

Como $\mathcal{O}_X(n)$ es localmente libre, $\Gamma_*(\mathcal{F})$ es un funtor covariante aditivo \emph{exacto por la izquierda} en $\mathcal{F}$; en particular, si $\mathcal{J}$ es un haz de ideales de $\mathcal{O}_X$, $\Gamma_*(\mathcal{J})$ se identifica canónicamente con un \emph{ideal graduado} de $\Gamma_*(\mathcal{O}_X)$.
\end{env}

\begin{env}[2.6.2]
\label{II.2.6.2-fr}
Sea $M$ un $S$-módulo graduado.
Para todo $f\in S_d$ ($d>0$), $x\to x/1$ es un homomorfismo de grupos abelianos $M_0\to M_{(f)}$ y, como $M_{(f)}$ se identifica canónicamente
\oldpage[II]{37}
con $\Gamma(D_+(f),\widetilde{M})$, se obtiene así un homomorfismo de grupos abelianos
$\alpha_0^f:M_0\to\Gamma(D_+(f),\widetilde{M})$.
Es claro que, para todo $g\in S_e$ ($e>0$), el diagrama
\[
  \xymatrix{
    & \Gamma(D_+(f),\widetilde{M})\ar[dd]\\
    M_0\ar[ur]^{\alpha_0^f}\ar[dr]_{\alpha_0^{fg}}\\
    & \Gamma(D_+(fg),\widetilde{M})
  }
\]
es conmutativo.
Esto significa que, para todo $x\in M_0$, las secciones $\alpha_0^f(x)$ y $\alpha_0^g(x)$ de $\widetilde{M}$ coinciden sobre $D_+(f)\cap D_+(g)$ y, por consiguiente, existe una única sección $\alpha_0(x)\in\Gamma(X,\widetilde{M})$ cuya restricción a cada $D_+(f)$ es $\alpha_0^f(x)$.
Se ha definido así, sin utilizar la hipótesis de que $S$ esté generado por $S_1$, un homomorfismo de grupos abelianos
\[
\label{II.2.6.2.1-fr}
  \alpha_0:M_0\longrightarrow\Gamma(X,\widetilde{M}).
\tag{2.6.2.1}
\]

Aplicando este resultado al $S$-módulo graduado $M(n)$, para cada $n\in\mathbf{Z}$, se obtiene, para cada $n\in\mathbf{Z}$, un homomorfismo de grupos abelianos
\[
\label{II.2.6.2.2-fr}
  \alpha_n:M_n=(M(n))_0\longrightarrow\Gamma(X,\widetilde{M}(n))
\tag{2.6.2.2}
\]
teniendo en cuenta \sref{II.2.5.15-fr}; de aquí se obtiene un homomorfismo funtorial (de grado $0$) de grupos abelianos graduados
\[
\label{II.2.6.2.3-fr}
  \alpha:M\longrightarrow\Gamma_*(\widetilde{M})
\tag{2.6.2.3}
\]
(denotado también por $\alpha_M$) que, en cada $M_n$, coincide con $\alpha_n$.

Si se toma, en particular, $M=S$, se comprueba inmediatamente, teniendo en cuenta la definición \sref[0]{0.5.4.6-fr} de la multiplicación en $\Gamma_*(\mathcal{O}_X)$, que $\alpha:S\to\Gamma_*(\mathcal{O}_X)$ es un homomorfismo de anillos graduados y que, para todo $S$-módulo graduado $M$, \sref{II.2.6.2.3-fr} es un dihomomorfismo de módulos graduados.
\end{env}
% Fuente francesa: PMIHES_1961__8__5_0, printed pp. 37-38.

\begin{proposition}[2.6.3]
\label{II.2.6.3-fr}
Para todo $f\in S_d$ ($d>0$), $D_+(f)$ es idéntico al conjunto de los $\mathfrak{p}\in X$ donde la sección $\alpha_d(f)$ de $\mathcal{O}_X(d)$ no se anula \sref[0]{0.5.5.2-fr}.
\end{proposition}

\begin{proof}
Como $X=\bigcup_{g\in S_1}D_+(g)$ por hipótesis, basta demostrar que, para todo $g\in S_1$, el conjunto de los $\mathfrak{p}\in D_+(g)$ donde $\alpha_d(f)$ no se anula es idéntico a $D_+(fg)$.
Ahora bien, la restricción de $\alpha_d(f)$ a $D_+(g)$ es, por definición, la sección que corresponde al elemento $f/1$ de $(S(d))_{(g)}$.
Mediante el isomorfismo canónico $(S(d))_{(g)}\xrightarrow{\sim}S_{(g)}$ \sref{II.2.5.7-fr}, esta sección de $\mathcal{O}_X(d)$ sobre $D_+(g)$ se identifica con la sección de $\mathcal{O}_X$ sobre $D_+(g)$ que corresponde al elemento $f/g^d$ de $S_{(g)}$.
Decir que esta sección se anula en $\mathfrak{p}\in D_+(g)$ significa que $f/g^d\in\mathfrak{q}$, donde $\mathfrak{q}$ es el ideal primo de $S_{(g)}$ correspondiente a $\mathfrak{p}$ \sref{II.2.3.6-fr}; por definición, esto equivale a $f\in\mathfrak{p}$, de donde resulta la proposición.
\end{proof}

\begin{env}[2.6.4]
\label{II.2.6.4-fr}
Sea ahora $\mathcal{F}$ un $\mathcal{O}_X$-módulo, y pongamos $M=\Gamma_*(\mathcal{F})$.
En virtud de la existencia del homomorfismo de anillos graduados $\alpha:S\to\Gamma_*(\mathcal{O}_X)$, puede considerarse $M$ como un $S$-módulo graduado.
Para todo $f\in S_d$ ($d>0$), resulta de \sref{II.2.6.3-fr} que la restricción a $D_+(f)$ de la sección $\alpha_d(f)$ de $\mathcal{O}_X(d)$ es invertible.
Lo mismo ocurre, por tanto, con la restricción a $D_+(f)$ de la sección $\alpha_d(f^n)$ de $\mathcal{O}_X(nd)$, para todo $n>0$.
Sea entonces $z\in M_{nd}=\Gamma(X,\mathcal{F}(nd))$ ($n>0$).
Si existe un entero $k\geq0$ tal que la restricción a $D_+(f)$ de $f^kz$, es decir, la
\oldpage[II]{38}
sección $(z|D_+(f))(\alpha_d(f^k)|D_+(f))$ de $\mathcal{F}((n+k)d)$, sea nula, entonces, por la observación anterior, se tiene también $z|D_+(f)=0$.
Esto muestra que se ha definido un $S_{(f)}$-homomorfismo $\beta_f:M_{(f)}\to\Gamma(D_+(f),\mathcal{F})$ haciendo corresponder al elemento $z/f^n$ la sección
$(z|D_+(f))(\alpha_d(f^n)|D_+(f))^{-1}$
de $\mathcal{F}$ sobre $D_+(f)$.
Se comprueba inmediatamente que el diagrama
\[
\label{II.2.6.4.1-fr}
  \xymatrix{
    M_{(f)}\ar[r]^{\beta_f}\ar[d]
    & \Gamma(D_+(f),\mathcal{F})\ar[d]\\
    M_{(fg)}\ar[r]_{\beta_{fg}}
    & \Gamma(D_+(fg),\mathcal{F})
  }
\tag{2.6.4.1}
\]
es conmutativo para $g\in S_e$ ($e>0$).
Si se recuerda que $M_{(f)}$ se identifica canónicamente con $\Gamma(D_+(f),\widetilde{M})$ y que los $D_+(f)$ forman una base de la topología de $X$ \sref{II.2.3.4-fr}, se ve que los $\beta_f$ proceden de un único homomorfismo canónico de $\mathcal{O}_X$-módulos
\[
\label{II.2.6.4.2-fr}
  \beta:(\Gamma_*(\mathcal{F}))\supertilde
  \longrightarrow\mathcal{F}
\tag{2.6.4.2}
\]
(denotado también por $\beta_{\mathcal{F}}$), que es evidentemente funtorial.
\end{env}
% Fuente francesa: PMIHES_1961__8__5_0, printed p. 38.

\begin{proposition}[2.6.5]
\label{II.2.6.5-fr}
Sea $M$ un $S$-módulo graduado, y sea $\mathcal{F}$ un $\mathcal{O}_X$-módulo.
Entonces, los homomorfismos compuestos
\[
\label{II.2.6.5.1-fr}
  \widetilde{M}\xrightarrow{\widetilde{\alpha}}
  (\Gamma_*(\widetilde{M}))\supertilde
  \xrightarrow{\beta}\widetilde{M}
\tag{2.6.5.1}
\]
\[
\label{II.2.6.5.2-fr}
  \Gamma_*(\mathcal{F})\xrightarrow{\alpha}
  \Gamma_*((\Gamma_*(\mathcal{F}))\supertilde)
  \xrightarrow{\Gamma_*(\beta)}\Gamma_*(\mathcal{F})
\tag{2.6.5.2}
\]
son los isomorfismos identidad.
\end{proposition}

\begin{proof}
La verificación de \sref{II.2.6.5.1-fr} es local: sobre un abierto $D_+(f)$, resulta inmediatamente de las definiciones y del hecho de que $\beta$, aplicado a haces cuasicoherentes, está determinado por su acción sobre las secciones por encima de $D_+(f)$ \sref[I]{I.1.3.8-fr}.
La afirmación \sref{II.2.6.5.2-fr} se verifica separadamente en cada grado.
Si se pone $M=\Gamma_*(\mathcal{F})$, se tiene $M_n=\Gamma(X,\mathcal{F}(n))$ y
$(\Gamma_*(\widetilde{M}))_n=\Gamma(X,\widetilde{M}(n))
=\Gamma(X,(M(n))\supertilde)$.
Ahora bien, si $f\in S_1$ y $z\in M_n$, $\alpha_n^f(z)$ es el elemento $z/1$ de $(M(n))_{(f)}$, igual a $(f/1)^n(z/f^n)$.
Le corresponde mediante $\beta_f$ la sección
\[
  ((\alpha_1(f))^n|D_+(f))
  ((z|D_+(f))((\alpha_1(f))^n|D_+(f))^{-1})
\]
sobre $D_+(f)$, es decir, la restricción de $z$ a $D_+(f)$, lo que termina la demostración de \sref{II.2.6.5.2-fr}.
\end{proof}
% Fuente francesa: PMIHES_1961__8__5_0, printed pp. 38-40.


\subsection{Condiciones de finitud}
\label{subsection:II.2.7-fr}

\begin{proposition}[2.7.1]
\label{II.2.7.1-fr}
\begin{enumerate}
  \item[{\rm(i)}] Si $S$ es un anillo graduado noetheriano, entonces $X=\operatorname{Proj}(S)$ es un esquema noetheriano.
  \item[{\rm(ii)}] Si $S$ es una $A$-álgebra graduada de tipo finito, entonces $X=\operatorname{Proj}(S)$ es un esquema de tipo finito sobre $Y=\operatorname{Spec}(A)$.
\end{enumerate}
\end{proposition}

\oldpage[II]{39}
\begin{proof}
\medskip\noindent
\begin{enumerate}
  \item[{\rm(i)}] Si $S$ es noetheriano, el ideal $S_+$ admite un sistema finito de generadores homogéneos $f_i$ ($1\leq i\leq p$); por tanto, \sref{II.2.3.14-fr}, el espacio subyacente $X$ es la unión de los $D_+(f_i)=\operatorname{Spec}(S_{(f_i)})$, y todo se reduce entonces a ver que cada uno de los $S_{(f_i)}$ es noetheriano, lo que resulta de \sref{II.2.2.6-fr}.
  \item[{\rm(ii)}] La hipótesis implica que $S_0$ es una $A$-álgebra de tipo finito y que $S$ es una $S_0$-álgebra de tipo finito, de modo que $S_+$ es un ideal de tipo finito \sref{II.2.1.4-fr}.
    Se reduce entonces, como en (i), a demostrar que, para todo $f\in S_d$, $S_{(f)}$ es una $A$-álgebra de tipo finito.
    En virtud de \sref{II.2.2.5-fr}, basta demostrar que $S^{(d)}$ es una $A$-álgebra de tipo finito, lo que resulta de \sref{II.2.1.6-fr}.
\end{enumerate}
\end{proof}

\begin{env}[2.7.2]
\label{II.2.7.2-fr}
En lo que sigue, consideraremos las siguientes condiciones de finitud para un $S$-módulo graduado $M$:
\begin{enumerate}
  \item[{\rm(TF)}] \emph{Existe un entero $n$ tal que el submódulo $\bigoplus_{k\geq n}M_k$ es un $S$-módulo de tipo finito.}
  \item[{\rm(TN)}] \emph{Existe un entero $n$ tal que $M_k=0$ para $k\geq n$.}
\end{enumerate}

Si $M$ satisface (TN), se tiene $M_{(f)}=0$ para todo $f$ homogéneo en $S_+$ y, por tanto, $\widetilde{M}=0$.

Sean $M$ y $N$ dos $S$-módulos graduados; se dice que un homomorfismo $u:M\to N$ de grado $0$ es (TN)-\emph{inyectivo} (resp. (TN)-\emph{sobreyectivo}, (TN)-\emph{biyectivo}) si existe un entero $n$ tal que $u_k:M_k\to N_k$ es inyectivo (resp. sobreyectivo, biyectivo) para $k\geq n$.
Decir que $u$ es (TN)-inyectivo (resp. (TN)-sobreyectivo) se reduce, pues, a decir que $\operatorname{Ker}u$ (resp. $\operatorname{Coker}u$) satisface (TN).
En virtud de \sref{II.2.5.4-fr}, si $u$ es (TN)-inyectivo (resp. (TN)-sobreyectivo, (TN)-biyectivo), entonces $\widetilde{u}$ es inyectivo (resp. sobreyectivo, biyectivo); si $u$ es (TN)-biyectivo, se dice también que $u$ es un (TN)-\emph{isomorfismo}.
\end{env}

\begin{proposition}[2.7.3]
\label{II.2.7.3-fr}
Sea $S$ un anillo graduado tal que el ideal $S_+$ sea de tipo finito, y sea $M$ un $S$-módulo graduado.
\begin{enumerate}
  \item[{\rm(i)}] Si $M$ satisface la condición \textup{(TF)}, el $\mathcal{O}_X$-módulo $\widetilde{M}$ es de tipo finito.
  \item[{\rm(ii)}] Supongamos que $M$ satisface \textup{(TF)}; para que $\widetilde{M}=0$, es necesario y suficiente que $M$ satisfaga \textup{(TN)}.
\end{enumerate}
\end{proposition}

\begin{proof}
Se acaba de ver que la condición (TN) implica que $\widetilde{M}=0$.
Si $M$ satisface (TF), el submódulo graduado $M'=\bigoplus_{k\geq n}M_k$, que es por hipótesis de tipo finito, es tal que $M/M'$ satisface (TN); por consiguiente, $(M/M')\supertilde=0$, y la exactitud del funtor $\widetilde{M}$ \sref{II.2.5.4-fr} implica que $\widetilde{M}=\widetilde{M}'$; para demostrar que $\widetilde{M}$ es de tipo finito, se puede reducir, pues, al caso en que $M$ es \emph{de tipo finito}.
Como la cuestión es local, basta demostrar que $M_{(f)}$ es un $S_{(f)}$-módulo de tipo finito \sref[I]{I.1.3.9-fr}; pero $M^{(d)}$ es un $S^{(d)}$-módulo de tipo finito \sref{II.2.1.6-fr}[iii], y nuestra afirmación resulta entonces de \sref{II.2.2.5-fr}.

Supongamos ahora que $M$ satisface (TF) y que $\widetilde{M}=0$; entonces, con las mismas notaciones que antes, se tiene $\widetilde{M}'=0$, y la condición (TN) para $M'$ equivale a la condición (TN) para $M$, de modo que, para demostrar que $\widetilde{M}=0$ implica que $M$ satisface (TN), se puede limitar de nuevo al caso en que $M$ está generado por un número finito de elementos homogéneos $x_i$ ($1\leq i\leq p$); sea, por otra parte, $(f_j)_{1\leq j\leq q}$ un sistema de generadores homogéneos del ideal $S_+$.
Por hipótesis, $M_{(f_j)}=0$ para todo $j$ y, por tanto, existe un entero $n$ tal que $f_j^nx_i=0$ cualesquiera que sean $i$ y $j$.
Sea $n_j$ el grado de $f_j$, y sea $m$ el mayor valor de $\sum_jr_jn_j$ sobre el conjunto finito de sistemas de enteros no negativos $(r_j)$ tales que $\sum_jr_j\leq nq$; es claro entonces
\oldpage[II]{40}
que, si $k>m$, se tiene $S_kx_i=0$ para todo $i$; si $h$ es el mayor de los grados de los $x_i$, se concluye que $M_k=0$ para $k>h+m$, lo que termina la demostración.
\end{proof}
% Fuente francesa: PMIHES_1961__8__5_0, printed p. 40.

\begin{corollary}[2.7.4]
\label{II.2.7.4-fr}
Sea $S$ un anillo graduado tal que el ideal $S_+$ sea de tipo finito; para que $X=\operatorname{Proj}(S)=\varnothing$, es necesario y suficiente que exista $n$ tal que $S_k=0$ para $k\geq n$.
\end{corollary}

\begin{proof}
La condición $X=\varnothing$ equivale, en efecto, a $\mathcal{O}_X=\widetilde{S}=0$, y $S$ es un $S$-módulo monógeno.
\end{proof}

\begin{theorem}[2.7.5]
\label{II.2.7.5-fr}
Se supone que el ideal $S_+$ está generado por un número finito de elementos homogéneos de grado $1$; sea $X=\operatorname{Proj}(S)$.
Entonces, para todo $\mathcal{O}_X$-módulo cuasicoherente $\mathcal{F}$, el homomorfismo canónico $\beta:(\Gamma_*(\mathcal{F}))\supertilde
\to\mathcal{F}$ \sref{II.2.6.4-fr} es un isomorfismo.
\end{theorem}

\begin{proof}
Si $S_+$ está generado por un número finito de elementos $f_i\in S_1$, entonces $X$ es la unión de los subespacios $\operatorname{Spec}(S_{(f_i)})$ \sref{II.2.3.6-fr}, que son cuasicompactos, y por tanto $X$ es cuasicompacto; además, es un esquema \sref{II.2.4.2-fr}; según \sref[I]{I.9.3.2-fr}, \sref{II.2.5.14.2-fr} y \sref{II.2.6.3-fr}, se tiene, para todo $f\in S_d$ ($d>0$), un isomorfismo canónico $(\Gamma_*(\mathcal{F}))_{(\alpha_d(f))}\xrightarrow{\sim}
\Gamma(D_+(f),\mathcal{F})$; por otra parte, por definición, $(\Gamma_*(\mathcal{F}))_{(\alpha_d(f))}$ (donde $\Gamma_*(\mathcal{F})$ se considera como $\Gamma_*(\mathcal{O}_X)$-módulo) no es otro que $(\Gamma_*(\mathcal{F}))_{(f)}$ (donde $\Gamma_*(\mathcal{F})$ se considera como $S$-módulo); si se remite a la definición \sref[I]{I.9.3.1-fr} del isomorfismo canónico anterior, se ve que coincide con el homomorfismo $\beta_f$, de donde resulta el teorema.
\end{proof}

\begin{remark}[2.7.6]
\label{II.2.7.6-fr}
Si se supone que el anillo graduado $S$ es \emph{noetheriano}, la condición de \sref{II.2.7.5-fr} se cumple \emph{ipso facto} en cuanto se supone que el ideal $S_+$ está generado por el conjunto $S_1$ de los elementos homogéneos de grado $1$.
\end{remark}
% Fuente francesa: PMIHES_1961__8__5_0, printed pp. 40-41.

\begin{corollary}[2.7.7]
\label{II.2.7.7-fr}
Bajo las hipótesis de \sref{II.2.7.5-fr}, todo $\mathcal{O}_X$-módulo cuasicoherente $\mathcal{F}$ es isomorfo a un $\mathcal{O}_X$-módulo de la forma $\widetilde{M}$, donde $M$ es un $S$-módulo graduado.
\end{corollary}

\begin{corollary}[2.7.8]
\label{II.2.7.8-fr}
Bajo las hipótesis de \sref{II.2.7.5-fr}, todo $\mathcal{O}_X$-módulo cuasicoherente de tipo finito $\mathcal{F}$ es isomorfo a un $\mathcal{O}_X$-módulo de la forma $\widetilde{N}$, donde $N$ es un $S$-módulo graduado de tipo finito.
\end{corollary}

\begin{proof}
Se puede suponer que $\mathcal{F}=\widetilde{M}$, donde $M$ es un $S$-módulo graduado \sref{II.2.7.7-fr}.
Sea $(f_\lambda)_{\lambda\in L}$ un sistema de generadores homogéneos de $M$; para todo subconjunto finito $H$ de $L$, sea $M_H$ el submódulo graduado de $M$ generado por los $f_\lambda$ tales que $\lambda\in H$; es claro que $M$ es el límite inductivo de sus submódulos $M_H$ y, por tanto, $\mathcal{F}$ es el límite inductivo de sus sub-$\mathcal{O}_X$-módulos $\widetilde{M}_H$ \sref{II.2.5.4-fr}.
Pero, como $\mathcal{F}$ es de tipo finito y el espacio subyacente de $X$ es cuasicompacto, resulta de \sref[0]{0.5.2.3-fr} que $\mathcal{F}=\widetilde{M}_H$ para algún subconjunto finito $H$ de $L$.
\end{proof}

\begin{corollary}[2.7.9]
\label{II.2.7.9-fr}
Bajo las hipótesis de \sref{II.2.7.5-fr}, sea $\mathcal{F}$ un $\mathcal{O}_X$-módulo cuasicoherente de tipo finito.
Existe entonces un entero $n_0$ tal que, para todo $n\geq n_0$, $\mathcal{F}(n)$ es isomorfo a un cociente de un $\mathcal{O}_X$-módulo de la forma $\mathcal{O}_X^k$ (para cierto $k>0$ que depende de $n$) y, por consiguiente, está generado por un número finito de sus secciones sobre $X$ \sref[0]{0.5.1.1-fr}.
\end{corollary}

\begin{proof}
En virtud de \sref{II.2.7.8-fr}, se puede suponer que $\mathcal{F}=\widetilde{M}$, donde $M$ es un cociente de una suma directa finita de $S$-módulos de la forma $S(m_i)$; en virtud de \sref{II.2.5.4-fr}, se puede, pues, limitar al caso en que $M=S(m)$, de modo que $\mathcal{F}(n)=(S(m+n))\supertilde
=\mathcal{O}_X(m+n)$ \sref{II.2.5.15-fr}.
Basta entonces demostrar
% Fuente francesa: PMIHES_1961__8__5_0, printed p. 41.

\oldpage[II]{41}
  \begin{lemma}[2.7.9.1]
  \label{II.2.7.9.1-fr}
  Bajo las hipótesis de \sref{II.2.7.5-fr}, para todo $n\geq0$, existe un entero $k$ (que depende de $n$) y un homomorfismo sobreyectivo $\mathcal{O}_X^k\to\mathcal{O}_X(n)$.
  \end{lemma}

Basta \sref{II.2.7.2-fr} demostrar que, para un $k$ conveniente, existe un homomorfismo (TN)-sobreyectivo $u$ \emph{de grado $0$} del $S$-módulo graduado producto $S^k$ en $S(n)$.
Ahora bien, $(S(n))_0=S_n$ y, por hipótesis, $S_h=S_1^h$ para todo $h>0$, de modo que $SS_n=S_n+S_{n+1}+\ldots$.
Como $S_n$ es un $S_0$-módulo de tipo finito (\sref{II.2.1.5-fr} y \sref{II.2.1.6-fr}[i]), consideremos un sistema de generadores $(a_i)_{1\leq i\leq k}$ de este módulo; el homomorfismo $u$ hará corresponder $a_i$ al $i$-ésimo elemento de la base canónica de $S^k$ ($1\leq i\leq k$); como $\operatorname{Coker}u$ se identifica entonces con $(S(n))_{-n}+\ldots+(S(n))_{-1}$, $u$ responde efectivamente a la cuestión.
\end{proof}

\begin{corollary}[2.7.10]
\label{II.2.7.10-fr}
Bajo las hipótesis de \sref{II.2.7.5-fr}, sea $\mathcal{F}$ un $\mathcal{O}_X$-módulo cuasicoherente de tipo finito.
Existe entonces un entero $n_0$ tal que, para todo $n\geq n_0$, $\mathcal{F}$ es isomorfo a un cociente de un $\mathcal{O}_X$-módulo de la forma $(\mathcal{O}_X(-n))^k$ (para cierto $k$ que depende de $n$).
\end{corollary}

\begin{proposition}[2.7.11]
\label{II.2.7.11-fr}
Supongamos que se cumplen las hipótesis de \sref{II.2.7.5-fr}, y sea $M$ un $S$-módulo graduado.
Entonces:
\begin{enumerate}
  \item[{\rm(i)}] El homomorfismo canónico $\widetilde{\alpha}:\widetilde{M}\to
    (\Gamma_*(\widetilde{M}))\supertilde$ es un isomorfismo.
  \item[{\rm(ii)}] Sea $\mathcal{G}$ un sub-$\mathcal{O}_X$-módulo cuasicoherente de $\widetilde{M}$, y sea $N$ el sub-$S$-módulo graduado de $M$ dado por la imagen inversa de $\Gamma_*(\mathcal{G})$ mediante $\alpha$.
    Entonces $\widetilde{N}=\mathcal{G}$ (identificándose $\widetilde{N}$, en virtud de \sref{II.2.5.4-fr}, con un sub-$\mathcal{O}_X$-módulo de $\widetilde{M}$).
\end{enumerate}
\end{proposition}

\begin{proof}
Como $\beta:(\Gamma_*(\widetilde{M}))\supertilde
\to\widetilde{M}$ es un isomorfismo (en virtud de \sref{II.2.7.5-fr}), $\widetilde{\alpha}$ es el isomorfismo inverso (en virtud de \sref{II.2.6.5.1-fr}), de donde resulta (i).
Sea $P$ el submódulo graduado $\alpha(M)$ de $\Gamma_*(\widetilde{M})$;
como $\widetilde{M}$ es un funtor exacto \sref{II.2.5.4-fr}, la imagen de $\widetilde{M}$ mediante $\widetilde{\alpha}$ es igual a $\widetilde{P}$ y, por tanto, en virtud de (i), $\widetilde{P}=(\Gamma_*(\widetilde{M}))\supertilde$.
Pongamos $Q=\Gamma_*(\mathcal{G})\cap P$; en virtud de lo que precede y de \sref{II.2.5.4-fr}, se tiene $\widetilde{Q}=(\Gamma_*(\mathcal{G}))\supertilde$, y la restricción de $\beta$ a $\widetilde{Q}$ es, pues, un \emph{isomorfismo} de este $\mathcal{O}_X$-módulo sobre $\mathcal{G}$ por \sref{II.2.7.5-fr}.
Pero, por definición de $N$ y por \sref{II.2.5.4-fr}, la restricción del isomorfismo $\widetilde{\alpha}$ a $\widetilde{N}$ es un isomorfismo de $\widetilde{N}$ sobre $\widetilde{Q}$, de donde resulta la conclusión por \sref{II.2.6.5.1-fr}.
\end{proof}
% Fuente francesa: PMIHES_1961__8__5_0, printed pp. 41-42.

\subsection{Comportamientos funtoriales}
\label{subsection:II.2.8-fr}

\begin{env}[2.8.1]
\label{II.2.8.1-fr}
Sean $S$ y $S'$ dos anillos graduados en grados positivos, y sea $\varphi:S'\to S$ un homomorfismo de anillos graduados.
Designaremos por $G(\varphi)$ la parte abierta de $X=\operatorname{Proj}(S)$ complementaria de $V_+(\varphi(S'_+))$ o, lo que equivale a lo mismo, la unión de los $D_+(\varphi(f'))$, donde $f'$ recorre el conjunto de los elementos homogéneos de $S'_+$.
La restricción a $G(\varphi)$ de la aplicación continua ${}^a\varphi$ de $\operatorname{Spec}(S)$ en $\operatorname{Spec}(S')$ \sref[I]{I.1.2.1-fr} es, pues, una aplicación continua de $G(\varphi)$ en $\operatorname{Proj}(S')$, que seguiremos denotando, por abuso de lenguaje, por ${}^a\varphi$.
Si $f'\in S'_+$ es homogéneo, se tiene
\[
\label{II.2.8.1.1-fr}
  {}^a\varphi^{-1}(D_+(f'))=D_+(\varphi(f'))
\tag{2.8.1.1}
\]
teniendo en cuenta que ${}^a\varphi$ aplica $G(\varphi)$ en $\operatorname{Proj}(S')$, así como \sref[I]{I.1.2.2.2-fr}.
El homomorfismo $\varphi$ define también canónicamente (con las mismas notaciones) un homomorfismo de anillos graduados $S'_{f'}\to S_f$, de donde se obtiene, por restricción a los elementos de grado $0$,
\oldpage[II]{42}
un homomorfismo $S'_{(f')}\to S_{(f)}$, que denotaremos por $\varphi_{(f)}$; le corresponde \sref[I]{I.1.6.1-fr} un morfismo de esquemas afines $({}^a\varphi_{(f)},\widetilde{\varphi}_{(f)}):
\operatorname{Spec}(S_{(f)})\to\operatorname{Spec}(S'_{(f')})$.
Si se identifica canónicamente $\operatorname{Spec}(S_{(f)})$ con el esquema inducido por $\operatorname{Proj}(S)$ sobre $D_+(f)$ \sref{II.2.3.6-fr}, se ha definido así un morfismo $\Phi_f:D_+(f)\to D_+(f')$, y ${}^a\varphi_{(f)}$ se identifica con la restricción de ${}^a\varphi$ a $D_+(f)$.
Por otra parte, es inmediato que, si $g'$ es otro elemento homogéneo de $S'_+$ y $g=\varphi(g')$, el diagrama
\[
  \xymatrix{
    D_+(f)\ar[r]^{\Phi_f} & D_+(f')\\
    D_+(fg)\ar[u]\ar[r]_{\Phi_{fg}} & D_+(f'g')\ar[u]
  }
\]
es conmutativo, en virtud de la conmutatividad del diagrama
\[
  \xymatrix{
    S'_{(f')}\ar[r]^{\varphi_{(f)}}\ar[d]_{\omega_{f'g',f'}}
    & S_{(f)}\ar[d]^{\omega_{fg,f}}\\
    S'_{(f'g')}\ar[r]_{\varphi_{(fg)}} & S_{(fg)}
  }
\]
Teniendo en cuenta la definición de $G(\varphi)$ y \sref{II.2.3.3.2-fr}, se ve, pues, que:
\end{env}
% Fuente francesa: PMIHES_1961__8__5_0, printed pp. 42-43.

\begin{proposition}[2.8.2]
\label{II.2.8.2-fr}
Dado un homomorfismo de anillos graduados $\varphi:S'\to S$, existe un único morfismo $({}^a\varphi,\widetilde{\varphi})$ del preesquema inducido $G(\varphi)$ en $\operatorname{Proj}(S')$ \textup{(llamado} asociado a $\varphi$ y denotado por $\operatorname{Proj}(\varphi)$) tal que, para todo elemento homogéneo $f'\in S'_+$, la restricción de este morfismo a $D_+(\varphi(f'))$ coincide con el morfismo asociado al homomorfismo $S'_{(f')}\to S_{(\varphi(f'))}$ correspondiente a $\varphi$.
\end{proposition}

\begin{proof}
Con las notaciones anteriores, si $f'\in S'_d$, el diagrama
\[
\label{II.2.8.2.1-fr}
  \xymatrix{
    S'_{(f')}\ar[r]^{\varphi_{(f)}}\ar[d]_{\wr}
    & S_{(f)}\ar[d]_{\wr}\\
    {S'}^{(d)}/(f'-1){S'}^{(d)}\ar[r]
    & S^{(d)}/(f-1)S^{(d)}
  }
\tag{2.8.2.1}
\]
es conmutativo (siendo las flechas verticales los isomorfismos de \sref{II.2.2.5-fr}).
\end{proof}

\begin{corollary}[2.8.3]
\label{II.2.8.3-fr}
\begin{enumerate}
  \item[{\rm(i)}] El morfismo $\operatorname{Proj}(\varphi)$ es afín.
  \item[{\rm(ii)}] Si $\operatorname{Ker}(\varphi)$ es nilpotente (y, en particular, si $\varphi$ es inyectivo), el morfismo $\operatorname{Proj}(\varphi)$ es dominante.
\end{enumerate}
\end{corollary}

\begin{proof}
La afirmación (i) es consecuencia inmediata de \sref{II.2.8.2-fr} y \sref{II.2.8.1.1-fr}.
La afirmación (ii) resulta de que, si $\operatorname{Ker}(\varphi)$ es nilpotente, para todo $f'$ homogéneo en $S'_+$ se comprueba inmediatamente que $\operatorname{Ker}(\varphi_f)$ (con $f=\varphi(f')$) es nilpotente y, por tanto, también lo es $\operatorname{Ker}(\varphi_{(f)})$; se aplican entonces \sref{II.2.8.2-fr} y \sref[I]{I.1.2.7-fr}.
\end{proof}

\oldpage[II]{43}
Se observa que, en general, hay morfismos de $\operatorname{Proj}(S)$ en $\operatorname{Proj}(S')$ que no son afines y que, por consiguiente, no provienen de homomorfismos de anillos graduados $S'\to S$; un ejemplo es el morfismo estructural $\operatorname{Proj}(S)\to\operatorname{Spec}(A)$ cuando $A$ es un cuerpo (identificándose entonces $\operatorname{Spec}(A)$ con $\operatorname{Proj}(A[T])$, cf.~\sref{II.3.1.7-fr}); esto resulta, en efecto, de \sref[I]{I.2.3.2-fr}.
% Fuente francesa: PMIHES_1961__8__5_0, printed p. 43.

\begin{env}[2.8.4]
\label{II.2.8.4-fr}
Sea $S''$ un tercer anillo graduado en grados positivos, sea $\varphi':S''\to S'$ un homomorfismo de anillos graduados, y pongamos $\varphi''=\varphi\circ\varphi'$.
En virtud de \sref{II.2.8.1.1-fr} y de la fórmula ${}^a\varphi''=({}^a\varphi')\circ({}^a\varphi)$, se comprueba inmediatamente que $G(\varphi'')\subset G(\varphi)$ y que, si $\Phi$, $\Phi'$ y $\Phi''$ son los morfismos asociados respectivamente a $\varphi$, $\varphi'$ y $\varphi''$, se tiene $\Phi''=\Phi'\circ(\Phi|G(\varphi''))$.
\end{env}

\begin{env}[2.8.5]
\label{II.2.8.5-fr}
Supongamos que $S$ (resp. $S'$) es una $A$-álgebra graduada (resp. una $A'$-álgebra graduada), y sea $\psi:A'\to A$ un homomorfismo de anillos tal que el diagrama
\[
\label{II.2.8.5.1-fr}
  \xymatrix{
    A'\ar[r]^{\psi}\ar[d] & A\ar[d]\\
    S'\ar[r]_{\varphi} & S
  }
\tag{2.8.5.1}
\]
sea conmutativo.
Se pueden considerar entonces $G(\varphi)$ (resp. $\operatorname{Proj}(S')$) como un esquema sobre $\operatorname{Spec}(A)$ (resp. $\operatorname{Spec}(A')$); si $\Phi$ (resp. $\Psi$) es el morfismo asociado a $\varphi$ (resp. a $\psi$), el diagrama
\[
  \xymatrix{
    G(\varphi)\ar[r]^{\Phi}\ar[d] & \operatorname{Proj}(S')\ar[d]\\
    \operatorname{Spec}(A)\ar[r]_{\Psi} & \operatorname{Spec}(A')
  }
\]
es conmutativo: basta, en efecto, demostrarlo para la restricción de $\Phi$ a $D_+(f)$, donde $f=\varphi(f')$, con $f'$ homogéneo en $S'_+$; esto resulta entonces de la conmutatividad del diagrama
\[
  \xymatrix{
    A'\ar[r]^{\psi}\ar[d] & A\ar[d]\\
    S'_{(f')}\ar[r]_{\varphi_{(f)}} & S_{(f)}
  }
\]
\end{env}
% Fuente francesa: PMIHES_1961__8__5_0, printed pp. 43-44.

\begin{env}[2.8.6]
\label{II.2.8.6-fr}
Sea ahora $M$ un $S$-módulo graduado, y consideremos el $S'$-módulo $M_{[\varphi]}$, que está evidentemente graduado.
Sea $f'$ homogéneo en $S'_+$, y sea $f=\varphi(f')$; se sabe \sref[0]{0.1.5.2-fr} que existe un isomorfismo canónico $(M_{[\varphi]})_{f'}\xrightarrow{\sim}(M_f)_{[\varphi_f]}$, y es inmediato que este isomorfismo conserva los grados, de donde se obtiene un isomorfismo canónico $(M_{[\varphi]})_{(f')}\xrightarrow{\sim}
(M_{(f)})_{[\varphi_{(f)}]}$.
A este isomorfismo corresponde canónicamente un isomorfismo de haces $(M_{[\varphi]})\supertilde|D_+(f')
\xrightarrow{\sim}(\Phi_f)_*(\widetilde{M}|D_+(f))$ (\sref{II.2.5.2-fr} y \sref[I]{I.1.6.3-fr}).
Además,
\oldpage[II]{44}
si $g'$ es otro elemento homogéneo de $S'_+$ y $g=\varphi(g')$, el diagrama
\[
  \xymatrix{
    (M_{[\varphi]})_{(f')}\ar[r]^{\sim}\ar[d]
    & (M_{(f)})_{[\varphi_{(f)}]}\ar[d]\\
    (M_{[\varphi]})_{(f'g')}\ar[r]^{\sim}
    & (M_{(fg)})_{[\varphi_{(fg)}]}
  }
\]
es conmutativo, de donde se concluye inmediatamente que el isomorfismo
\[
  (M_{[\varphi]})\supertilde|D_+(f'g')
  \xrightarrow{\sim}(\Phi_{fg})_*(\widetilde{M}|D_+(fg))
\]
es la restricción a $D_+(f'g')$ del isomorfismo $(M_{[\varphi]})\supertilde|D_+(f')
\xrightarrow{\sim}(\Phi_f)_*(\widetilde{M}|D_+(f))$.
Como $\Phi_f$ es la restricción a $D_+(f)$ del morfismo $\Phi$, se ve, teniendo en cuenta \sref{II.2.8.1.1-fr} y poniendo $X'=\operatorname{Proj}(S')$, que:
\end{env}

\begin{proposition}[2.8.7]
\label{II.2.8.7-fr}
Existe un isomorfismo canónico funtorial del $\mathcal{O}_{X'}$-módulo $(M_{[\varphi]})\supertilde$ sobre el $\mathcal{O}_{X'}$-módulo $\Phi_*(\widetilde{M}|G(\varphi))$.
\end{proposition}

Se deduce inmediatamente una aplicación canónica funtorial del conjunto de los $\varphi$-morfismos $M'\to M$ de un $S'$-módulo graduado en el $S$-módulo graduado $M$, en el conjunto de los $\Phi$-morfismos $\widetilde{M}'\to\widetilde{M}|G(\varphi)$.
Con las notaciones de \sref{II.2.8.4-fr}, si $M''$ es un $S''$-módulo graduado, al compuesto de un $\varphi$-morfismo $M'\to M$ y de un $\varphi'$-morfismo $M''\to M'$ corresponde canónicamente el compuesto de $\widetilde{M}'|G(\varphi')\to\widetilde{M}|G(\varphi'')$ y $\widetilde{M}''\to\widetilde{M}'|G(\varphi')$.
% Fuente francesa: PMIHES_1961__8__5_0, printed pp. 44-45.

\begin{proposition}[2.8.8]
\label{II.2.8.8-fr}
Bajo las hipótesis de \sref{II.2.8.1-fr}, sea $M'$ un $S'$-módulo graduado.
Existe entonces un homomorfismo canónico funtorial $\nu$ del $(\mathcal{O}_X|G(\varphi))$-módulo $\Phi^*(\widetilde{M}')$ en el $(\mathcal{O}_X|G(\varphi))$-módulo $(M'\otimes_{S'}S)\supertilde|G(\varphi)$.
Si el ideal $S'_+$ está generado por $S'_1$, entonces $\nu$ es un isomorfismo.
\end{proposition}

\begin{proof}
En efecto, para $f'\in S'_d$ ($d>0$), se define un homomorfismo canónico funtorial de $S_{(f)}$-módulos (donde $f=\varphi(f')$)
\[
\label{II.2.8.8.1-fr}
  \nu_f:M'_{(f')}\otimes_{S'_{(f')}}S_{(f)}
  \longrightarrow(M'\otimes_{S'}S)_{(f)}
\tag{2.8.8.1}
\]
componiendo el homomorfismo $M'_{(f')}\otimes_{S'_{(f')}}S_{(f)}\to
M'_{f'}\otimes_{S'_{f'}}S_f$ con el isomorfismo canónico $M'_{f'}\otimes_{S'_{f'}}S_f\xrightarrow{\sim}(M'\otimes_{S'}S)_f$ \sref[0]{0.1.5.4-fr}, y observando que este último conserva los grados.
Se comprueba inmediatamente la compatibilidad de los $\nu_f$ con los operadores de restricción de $D_+(f)$ a $D_+(fg)$ (para todo $g'\in S'_+$ y $g=\varphi(g')$), de donde se obtiene la definición del homomorfismo
\[
  \nu:\Phi^*(\widetilde{M}')
  \longrightarrow(M'\otimes_{S'}S)\supertilde|G(\varphi)
\]
teniendo en cuenta \sref[I]{I.1.6.5-fr}.
Para demostrar la segunda afirmación, basta probar que $\nu_f$ es un isomorfismo para todo $f'\in S'_1$, pues $G(\varphi)$ es entonces la unión de los $D_+(\varphi(f'))$.
Se define primero una aplicación $\mathbf{Z}$-bilineal $M'_m\times S_n\to M'_{(f')}\otimes_{S'_{(f')}}S_{(f)}$ haciendo corresponder a $(x',s)$ el elemento $(x'/{f'}^m)\otimes(s/f^n)$ (con el convenio de que $x'/{f'}^m$ es ${f'}^{-m}x'/1$ 
\oldpage[II]{45}
cuando $m<0$).

Se comprueba, como en la demostración de \sref{II.2.5.13-fr}, que esta aplicación da lugar a un dihomomorfismo de módulos
\[
  \eta_f:M'\otimes_{S'}S
  \longrightarrow M'_{(f')}\otimes_{S'_{(f')}}S_{(f)}.
\]
Además, si, para $r>0$, se tiene $f^r\sum_i(x'_i\otimes s_i)=0$, esto se escribe también $\sum_i({f'}^rx'_i\otimes s_i)=0$, de donde, por \sref[0]{0.1.5.4-fr}, $\sum_i({f'}^rx'_i/{f'}^{m_i+r})\otimes(s_i/f^{n_i})=0$, es decir, $\eta_f(\sum_i x'_i\otimes s_i)=0$, lo que demuestra que $\eta_f$ se factoriza en \[
  M'\otimes_{S'}S\longrightarrow(M'\otimes_{S'}S)_f
  \xrightarrow{\eta'_f}M'_{(f')}\otimes_{S'_{(f')}}S_{(f)}~;
\]; se comprueba finalmente que $\eta'_f$ y $\nu_f$ son isomorfismos inversos entre sí.

En particular, resulta de \sref{II.2.1.2.1-fr} que se tiene un homomorfismo canónico
\[
\label{II.2.8.8.2-fr}
  \Phi^*(\mathcal{O}_{X'}(n))\xrightarrow{\sim}
  \mathcal{O}_X(n)|G(\varphi)
\tag{2.8.8.2}
\]
para todo $n\in\mathbf{Z}$.
\end{proof}
% Fuente francesa: PMIHES_1961__8__5_0, printed p. 45.

\begin{env}[2.8.9]
\label{II.2.8.9-fr}
Sean $A$ y $A'$ dos anillos, y sea $\psi:A'\to A$ un homomorfismo de anillos que define un morfismo $\Psi:\operatorname{Spec}(A)\to\operatorname{Spec}(A')$.
Sea $S'$ una $A'$-álgebra graduada en grados positivos, y pongamos $S=S'\otimes_{A'}A$, que es evidentemente una $A$-álgebra graduada por los $S'_n\otimes_{A'}A$; la aplicación $\varphi:s'\to s'\otimes1$ es entonces un homomorfismo de anillos graduados que hace conmutativo el diagrama \sref{II.2.8.5.1-fr}.
Como $S_+$ es aquí el $A$-módulo generado por $\varphi(S'_+)$, se tiene $G(\varphi)=\operatorname{Proj}(S)=X$; por tanto, poniendo $X'=\operatorname{Proj}(S')$, se obtiene el diagrama conmutativo
\[
\label{II.2.8.9.1-fr}
  \xymatrix{
    X\ar[r]^{\Phi}\ar[d]_p & X'\ar[d]\\
    Y\ar[r]_{\Psi} & Y'
  }
\tag{2.8.9.1}
\]

Sea ahora $M'$ un $S'$-módulo graduado, y pongamos $M=M'\otimes_{A'}A=M'\otimes_{S'}S$.
En estas condiciones:
\end{env}

\begin{proposition}[2.8.10]
\label{II.2.8.10-fr}
El diagrama \sref{II.2.8.9.1-fr} identifica el esquema $X$ con el producto $X'\times_{Y'}Y$; además, el homomorfismo canónico $\nu:\Phi^*(\widetilde{M}')\to\widetilde{M}$ \sref{II.2.8.8-fr} es un isomorfismo.
\end{proposition}

\begin{proof}
La primera afirmación quedará demostrada si se prueba que, para todo $f'$ homogéneo en $S'_+$, poniendo $f=\varphi(f')$, las restricciones de $\Phi$ y $p$ a $D_+(f)$ identifican este esquema con el producto $D_+(f')\times_{Y'}Y$ \sref[I]{I.3.2.6.2-fr}; dicho de otro modo, basta \sref[I]{I.3.2.2-fr} demostrar que $S_{(f)}$ se identifica canónicamente con $S'_{(f')}\otimes_{A'}A$, lo que resulta inmediatamente del isomorfismo canónico, que conserva los grados, $S_f\xrightarrow{\sim}S'_{f'}\otimes_{A'}A$ \sref[0]{0.1.5.4-fr}.
La segunda afirmación resulta entonces de que $M'_{(f')}\otimes_{S'_{(f')}}S_{(f)}$ se identifica, por lo anterior, con $M'_{(f')}\otimes_{A'}A$, y este último con $M_{(f)}$, pues $M_f$ se identifica canónicamente con $M'_{f'}\otimes_{A'}A$ mediante un isomorfismo que conserva los grados.
\end{proof}
% Fuente francesa: PMIHES_1961__8__5_0, printed pp. 45-46.

\begin{corollary}[2.8.11]
\label{II.2.8.11-fr}
Para todo entero $n\in\mathbf{Z}$, $\widetilde{M}(n)$ se identifica con $\Phi^*(\widetilde{M}'(n))=\widetilde{M}'(n)\otimes_{Y'}\mathcal{O}_Y$; en particular, $\mathcal{O}_X(n)=\Phi^*(\mathcal{O}_{X'}(n))
=\mathcal{O}_{X'}(n)\otimes_{Y'}\mathcal{O}_Y$.
\end{corollary}

\begin{proof}
Esto resulta de \sref{II.2.8.10-fr} y \sref{II.2.5.15-fr}.
\end{proof}

\oldpage[II]{46}

\begin{env}[2.8.12]
\label{II.2.8.12-fr}
Bajo las hipótesis de \sref{II.2.8.9-fr}, para $f'\in S'_d$ ($d>0$) y $f=\varphi(f')$, el diagrama
\[
\label{II.2.8.12.1-fr}
  \xymatrix{
    M'_{(f')}\ar[r]^{\sim}\ar[d]
    & M'^{(d)}/(f'-1)M'^{(d)}\ar[d]\\
    M_{(f)}\ar[r]^{\sim}
    & M^{(d)}/(f-1)M^{(d)}
  }
\tag{2.8.12.1}
\]
(cf. \sref{II.2.2.5-fr}) es conmutativo.
\end{env}

\begin{env}[2.8.13]
\label{II.2.8.13-fr}
Conservemos las notaciones y las hipótesis de \sref{II.2.8.9-fr}, y sea $\mathcal{F}'$ un $\mathcal{O}_{X'}$-módulo; si se pone $\mathcal{F}=\Phi^*(\mathcal{F}')$, se tiene, para todo $n\in\mathbf{Z}$, $\mathcal{F}(n)=\Phi^*(\mathcal{F}'(n))$, en virtud de \sref{II.2.8.11-fr} y \sref[0]{0.4.3.3-fr}.
Por consiguiente, según \sref[0]{0.3.7.1-fr}, se tiene un homomorfismo canónico
\[
  \Gamma(\rho):\Gamma(X',\mathcal{F}'(n))
  \longrightarrow\Gamma(X,\mathcal{F}(n))
\]
que da un dihomomorfismo canónico de módulos graduados
\[
  \Gamma_*(\mathcal{F}')\longrightarrow
  \Gamma_*(\mathcal{F}).
\]

Supongamos que el ideal $S_+$ está generado por $S_1$ y que $\mathcal{F}'=\widetilde{M}'$, de modo que $\mathcal{F}=\widetilde{M}$ con $M=M'\otimes_{A'}A$.
Si $f'$ es homogéneo en $S'_+$ y $f=\varphi(f')$, se ha visto que $M_{(f)}=M'_{(f')}\otimes_{A'}A$, y el diagrama
\[
  \xymatrix{
    M'_0\ar[r]\ar[d]
    & M'_{(f')}\ar[d]
    & =\Gamma(D_+(f'),\widetilde{M}')\\
    M_0\ar[r]
    & M_{(f)}
    & =\Gamma(D_+(f),\widetilde{M})
  }
\]
es, pues, conmutativo; se concluye inmediatamente de esta observación y de la definición del homomorfismo $\alpha$ \sref{II.2.6.2-fr} que el diagrama
\[
\label{II.2.8.13.1-fr}
  \xymatrix{
    M'\ar[r]^{\alpha_{M'}}\ar[d]
    & \Gamma_*(\widetilde{M}')\ar[d]\\
    M\ar[r]_{\alpha_M}
    & \Gamma_*(\widetilde{M})
  }
\tag{2.8.13.1}
\]
es conmutativo.
Del mismo modo, el diagrama
\[
\label{II.2.8.13.2-fr}
  \xymatrix{
    (\Gamma_*(\mathcal{F}'))\supertilde\ar[r]^{\beta_{\mathcal{F}'}}
      \ar[d]
    & \mathcal{F}'\ar[d]\\
    (\Gamma_*(\mathcal{F}))\supertilde\ar[r]_{\beta_{\mathcal{F}}}
    & \mathcal{F}
  }
\tag{2.8.13.2}
\]
es conmutativo (siendo la flecha vertical de la derecha el $\Phi$-morfismo canónico $\mathcal{F}'\to\Phi^*(\mathcal{F}')=\mathcal{F}$).
\end{env}
% Fuente francesa: PMIHES_1961__8__5_0, printed p. 47.

\oldpage[II]{47}

\begin{env}[2.8.14]
\label{II.2.8.14-fr}
Conservando todavía las notaciones y las hipótesis de \sref{II.2.8.9-fr}, sea $N'$ otro $S'$-módulo graduado, y sea $N=N'\otimes_{A'}A$.
Es inmediato que los dihomomorfismos canónicos $M'\to M$ y $N'\to N$ dan un dihomomorfismo $M'\otimes_{S'}N'\to M\otimes_S N$ (relativo al homomorfismo canónico de anillos $S'\to S$) y, por consiguiente, también un $S$-homomorfismo de grado $0$, \[
  (M'\otimes_{S'}N')\otimes_{A'}A\longrightarrow M\otimes_S N,
\], al que corresponde (teniendo en cuenta \sref{II.2.8.10-fr}) un homomorfismo de $\mathcal{O}_X$-módulos
\[
\label{II.2.8.14.1-fr}
  \Phi^*((M'\otimes_{S'}N')\supertilde)
  \longrightarrow(M\otimes_S N)\supertilde.
\tag{2.8.14.1}
\]

Además, se comprueba inmediatamente que el diagrama
\[
\label{II.2.8.14.2-fr}
  \xymatrix{
    \Phi^*(\widetilde{M}'\otimes_{\mathcal{O}_{X'}}\widetilde{N}')
      \ar[r]^{\sim}\ar[d]_{\Phi^*(\lambda)}
    & \widetilde{M}\otimes_{\mathcal{O}_X}\widetilde{N}
      \ar[d]^{\lambda}
    & =\Phi^*(\widetilde{M}')\otimes_{\mathcal{O}_X}
      \Phi^*(\widetilde{N}')\\
    \Phi^*((M'\otimes_{S'}N')\supertilde)\ar[r]
    & (M\otimes_S N)\supertilde
  }
\tag{2.8.14.2}
\]
es conmutativo, siendo la primera fila el isomorfismo canónico \sref[0]{0.4.3.3-fr}.
Si el ideal $S'_+$ está generado por $S'_1$, es claro que $S_+$ está generado por $S_1$, y las dos flechas verticales de \sref{II.2.8.14.2-fr} son entonces isomorfismos \sref{II.2.5.13-fr}; lo mismo ocurre, por tanto, con \sref{II.2.8.14.1-fr}.

Se tiene igualmente un dihomomorfismo canónico $\operatorname{Hom}_{S'}(M',N')\to\operatorname{Hom}_S(M,N)$ haciendo corresponder a un homomorfismo $u'$ de grado $k$ el homomorfismo $u'\otimes1$, que también es de grado $k$; de aquí se deduce de nuevo un $S$-homomorfismo de grado $0$
\[
  (\operatorname{Hom}_{S'}(M',N'))\otimes_{A'}A
  \longrightarrow\operatorname{Hom}_S(M,N)
\]
de donde se obtiene un homomorfismo de $\mathcal{O}_X$-módulos
\[
\label{II.2.8.14.3-fr}
  \Phi^*((\operatorname{Hom}_{S'}(M',N'))\supertilde)
  \longrightarrow(\operatorname{Hom}_S(M,N))\supertilde.
\tag{2.8.14.3}
\]

Además, el diagrama
\[
  \xymatrix{
    \Phi^*((\operatorname{Hom}_{S'}(M',N'))\supertilde)\ar[r]
      \ar[d]_{\Phi^*(\mu)}
    & (\operatorname{Hom}_S(M,N))\supertilde\ar[d]^{\mu}\\
    \Phi^*(\mathcal{H}om_{\mathcal{O}_{X'}}
      (\widetilde{M}',\widetilde{N}'))\ar[r]
    & \mathcal{H}om_{\mathcal{O}_X}(\widetilde{M},\widetilde{N})
  }
\]
es conmutativo (siendo la segunda fila horizontal el homomorfismo canónico \sref[0]{0.4.4.6-fr}).
\end{env}
% Fuente francesa: PMIHES_1961__8__5_0, printed p. 48.

\oldpage[II]{48}

\begin{env}[2.8.15]
\label{II.2.8.15-fr}
Con las notaciones y las hipótesis de \sref{II.2.8.1-fr}, resulta de \sref{II.2.4.7-fr} que no se cambia el morfismo $\Phi$, salvo isomorfismo, cuando se sustituyen $S_0$ y $S'_0$ por $\mathbf{Z}$ y $\varphi_0$ por la aplicación identidad y, por tanto, cuando se sustituyen $S$ y $S'$ por $S^{(d)}$ y ${S'}^{(d)}$, respectivamente ($d>0$), y $\varphi$ por su restricción $\varphi^{(d)}$ a $S^{(d)}$.
\end{env}
% Fuente francesa: PMIHES_1961__8__5_0, printed p. 48.

\subsection{Subesquemas cerrados de un esquema $\operatorname{Proj}(S)$}
\label{subsection:II.2.9-fr}

\begin{env}[2.9.1]
\label{II.2.9.1-fr}
Si $\varphi:S\to S'$ es un homomorfismo de anillos graduados, se dice que $\varphi$ es (TN)-\emph{sobreyectivo} (resp. (TN)-\emph{inyectivo}, (TN)-\emph{biyectivo}) si existe un entero $n$ tal que, para $k\geq n$, $\varphi_k:S_k\to S'_k$ es \emph{sobreyectivo} (resp. \emph{inyectivo}, \emph{biyectivo}).
Resulta entonces de \sref{II.2.8.15-fr} que el estudio de $\Phi$ puede reducirse al caso en que $\varphi$ es sobreyectivo (resp. inyectivo, biyectivo).
En lugar de decir que $\varphi$ es (TN)-biyectivo, se dice también que es un (TN)-\emph{isomorfismo}.
\end{env}

\begin{proposition}[2.9.2]
\label{II.2.9.2-fr}
Sea $S$ un anillo graduado en grados positivos, y sea $X=\operatorname{Proj}(S)$.
\begin{enumerate}
  \item[{\rm(i)}] Si $\varphi:S\to S'$ es un homomorfismo \textup{(TN)}-sobreyectivo de anillos graduados, el morfismo correspondiente $\Phi$ \sref{II.2.8.1-fr} está definido sobre todo $\operatorname{Proj}(S')$ y es una inmersión cerrada de $\operatorname{Proj}(S')$ en $X$.
    Si $\mathfrak{J}$ es el núcleo de $\varphi$, el subesquema cerrado de $X$ asociado a $\Phi$ está definido por el haz cuasicoherente de ideales $\widetilde{\mathfrak{J}}$ de $\mathcal{O}_X$.
  \item[{\rm(ii)}] Supongamos además que el ideal $S_+$ está generado por un número finito de elementos homogéneos de grado $1$.
    Sea $X'$ un subesquema cerrado de $X$, definido por un haz cuasicoherente de ideales $\mathcal{J}$ de $\mathcal{O}_X$.
    Sea $\mathfrak{J}$ el ideal graduado de $S$, imagen inversa de $\Gamma_*(\mathcal{J})$ mediante el homomorfismo canónico $\alpha:S\to\Gamma_*(\mathcal{O}_X)$ \sref{II.2.6.2-fr}, y pongamos $S'=S/\mathfrak{J}$.
    Entonces $X'$ es el subesquema asociado a la inmersión cerrada $\operatorname{Proj}(S')\to X$ correspondiente al homomorfismo canónico de anillos graduados $S\to S'$.
\end{enumerate}
\end{proposition}

\begin{proof}
\medskip\noindent
\begin{enumerate}
  \item[{\rm(i)}] Puede suponerse que $\varphi$ es sobreyectivo \sref{II.2.9.1-fr}.
    Como, por hipótesis, $\varphi(S_+)$ genera $S'_+$, se tiene $G(\varphi)=\operatorname{Proj}(S')$.
    Por otra parte, la segunda afirmación se verifica localmente sobre $X$; sea, pues, $f$ un elemento homogéneo de $S_+$, y pongamos $f'=\varphi(f)$.
    Como $\varphi$ es un homomorfismo sobreyectivo de anillos graduados, se comprueba inmediatamente que $\varphi_{(f')}:S_{(f)}\to S'_{(f')}$ es sobreyectivo y que su núcleo es $\mathfrak{J}_{(f)}$, lo que termina de demostrar (i) \sref[I]{I.4.2.3-fr}.
  \item[{\rm(ii)}] En virtud de (i), se reduce a comprobar que el homomorfismo $\widetilde{j}:\widetilde{\mathfrak{J}}\to\mathcal{O}_X$ deducido de la inyección canónica $j:\mathfrak{J}\to S$ es un isomorfismo de $\widetilde{\mathfrak{J}}$ sobre $\mathcal{J}$, lo que resulta de \sref{II.2.7.11-fr}.
\end{enumerate}
\end{proof}

Se observa que $\mathfrak{J}$ es el \emph{mayor} de los ideales graduados $\mathfrak{J}'$ de $S$ tales que $\widetilde{j}(\widetilde{\mathfrak{J}}')=\mathcal{J}$, pues se comprueba inmediatamente, remontándose a las definiciones \sref{II.2.6.2-fr}, que esta relación implica $\alpha(\mathfrak{J}')\subset\Gamma_*(\mathcal{J})$.
% Fuente francesa: PMIHES_1961__8__5_0, printed pp. 48-49.

\begin{corollary}[2.9.3]
\label{II.2.9.3-fr}
Se supone que se cumplen las hipótesis de \sref{II.2.9.2-fr}[(i)] y, además, que el ideal $S_+$ está generado por $S_1$; entonces $\Phi^*((S(n))\supertilde)$ es canónicamente isomorfo a $(S'(n))\supertilde$ para todo $n\in\mathbf{Z}$ y, por consiguiente, $\Phi^*(\mathcal{F}(n))$ es canónicamente isomorfo a $\Phi^*(\mathcal{F})(n)$ para todo $\mathcal{O}_X$-módulo $\mathcal{F}$.
\end{corollary}

\begin{proof}
Es un caso particular de \sref{II.2.8.8-fr}, teniendo en cuenta \sref{II.2.5.10.2-fr}.
\end{proof}

\begin{corollary}[2.9.4]
\label{II.2.9.4-fr}
Se supone que se cumplen las hipótesis de \sref{II.2.9.2-fr}[(ii)].
Para que el subpreesquema cerrado $X'$ de $X$ sea íntegro, es necesario y suficiente que el ideal graduado $\mathfrak{J}$ sea primo en $S$.
\end{corollary}

\oldpage[II]{49}

\begin{proof}
Como $X'$ es isomorfo a $\operatorname{Proj}(S/\mathfrak{J})$, la condición es suficiente en virtud de \sref{II.2.4.4-fr}.
Para ver que es necesaria, consideremos la sucesión exacta $0\to\mathcal{J}\to\mathcal{O}_X\to\mathcal{O}_X/\mathcal{J}$, que da la sucesión exacta
\[
  0\longrightarrow\Gamma_*(\mathcal{J})
  \longrightarrow\Gamma_*(\mathcal{O}_X)
  \longrightarrow\Gamma_*(\mathcal{O}_X/\mathcal{J}).
\]

Basta demostrar que, si $f\in S_m$ y $g\in S_n$ son tales que la imagen en $\Gamma_*(\mathcal{O}_X/\mathcal{J})$ de $\alpha_{n+m}(fg)$ es nula, entonces una de las imágenes de $\alpha_m(f)$ o de $\alpha_n(g)$ es nula.
Ahora bien, por definición, estas imágenes son secciones de $(\mathcal{O}_X/\mathcal{J})$-módulos invertibles $\mathcal{L}=(\mathcal{O}_X/\mathcal{J})(m)$ y $\mathcal{L}'=(\mathcal{O}_X/\mathcal{J})(n)$ sobre el esquema íntegro $X'$; la hipótesis implica que el producto de estas dos secciones es nulo en $\mathcal{L}\otimes\mathcal{L}'$ (\sref{II.2.9.3-fr} y \sref{II.2.5.14.1-fr}) y, por tanto, una de ellas es nula en virtud de \sref[I]{I.7.4.4-fr}.
\end{proof}

\begin{corollary}[2.9.5]
\label{II.2.9.5-fr}
Sean $A$ un anillo, $M$ un $A$-módulo, $S$ una $A$-álgebra graduada generada por el conjunto $S_1$ de los elementos homogéneos de grado $1$, $u:M\to S_1$ un homomorfismo sobreyectivo de $A$-módulos, y $\overline{u}:\mathbf{S}(M)\to S$ el homomorfismo (de $A$-álgebras) del álgebra simétrica $\mathbf{S}(M)$ de $M$ en $S$ que prolonga $u$.
Entonces el morfismo correspondiente a $\overline{u}$ es una inmersión cerrada de $\operatorname{Proj}(S)$ en $\operatorname{Proj}(\mathbf{S}(M))$.
\end{corollary}

\begin{proof}
En efecto, $\overline{u}$ es sobreyectivo por hipótesis, y basta aplicar \sref{II.2.9.2-fr}.
\end{proof}
% Fuente francesa: PMIHES_1961__8__5_0, printed pp. 49-50.

\section{Espectro homogéneo de un haz de álgebras graduadas}
\label{section:II.3-fr}

\subsection{Espectro homogéneo de una $\mathcal{O}_Y$-álgebra graduada cuasicoherente}
\label{subsection:II.3.1-fr}
\label{II.3.1-fr}

\begin{env}[3.1.1]
\label{II.3.1.1-fr}
Sean $Y$ un preesquema, $\mathcal{S}$ una $\mathcal{O}_Y$-álgebra graduada y $\mathcal{M}$ un $\mathcal{S}$-módulo graduado.
Si $\mathcal{S}$ es \emph{cuasicoherente}, cada una de sus componentes homogéneas $\mathcal{S}_n$ es un $\mathcal{O}_Y$-módulo \emph{cuasicoherente}, pues es imagen de $\mathcal{S}$ mediante un homomorfismo de $\mathcal{S}$ en sí misma (\sref[I]{I.1.3.8-fr} y \sref[I]{I.1.3.9-fr}); del mismo modo, si $\mathcal{M}$ es cuasicoherente como $\mathcal{O}_Y$-módulo, sus componentes homogéneas $\mathcal{M}_n$ también lo son, y recíprocamente.
Si $d$ es un entero $>0$, se denota por $\mathcal{S}^{(d)}$ la suma directa de los $\mathcal{O}_Y$-módulos $\mathcal{S}_{nd}$ ($n\in\mathbf{Z}$), que es cuasicoherente si $\mathcal{S}$ lo es \sref[I]{I.1.3.9-fr}; para todo entero $k$ tal que $0\leq k\leq d-1$, se denota por $\mathcal{M}^{(d,k)}$ (o $\mathcal{M}^{(d)}$, para $k=0$) la suma directa de los $\mathcal{M}_{nd+k}$ ($n\in\mathbf{Z}$), que es un $\mathcal{S}^{(d)}$-módulo graduado y es cuasicoherente si $\mathcal{S}$ y $\mathcal{M}$ son cuasicoherentes \sref[I]{I.9.6.1-fr}.
Se denota por $\mathcal{M}(n)$ el $\mathcal{S}$-módulo graduado tal que $(\mathcal{M}(n))_k=\mathcal{M}_{n+k}$ para todo $k\in\mathbf{Z}$; si $\mathcal{S}$ y $\mathcal{M}$ son cuasicoherentes, $\mathcal{M}(n)$ es un $\mathcal{S}$-módulo graduado cuasicoherente \sref[I]{I.9.6.1-fr}.

Diremos que $\mathcal{M}$ es un $\mathcal{S}$-módulo graduado \emph{de tipo finito} (resp. que admite una \emph{presentación finita}) si, para todo $y\in Y$, existe una vecindad abierta $U$ de $y$ y enteros $n_i$ (resp. enteros $m_i$ y $n_j$) tales que existe un homomorfismo sobreyectivo de grado $0$, \[
  \bigoplus_{i=1}^r(\mathcal{S}(n_i)|U)\longrightarrow\mathcal{M}|U
\] (resp. tales que $\mathcal{M}|U$ es isomorfo al conúcleo de un homomorfismo de grado $0$, \[
  \bigoplus_{i=1}^r(\mathcal{S}(m_i)|U)
  \longrightarrow\bigoplus_{j=1}^s(\mathcal{S}(n_j)|U)).
\]).

Sea $U$ un abierto afín de $Y$, y sea $A=\Gamma(U,\mathcal{O}_Y)$ su anillo; por hipótesis, la $(\mathcal{O}_Y|U)$-álgebra graduada $\mathcal{S}|U$ es isomorfa a $\widetilde{S}$, donde $S=\Gamma(U,\mathcal{S})$ es una $A$-álgebra graduada \sref[I]{I.1.4.3-fr};
\oldpage[II]{50}
pongamos $X_U=\operatorname{Proj}(\Gamma(U,\mathcal{S}))$.
Sea $U'\subset U$ otro abierto afín de $Y$, de anillo $A'$, y sea $j:U'\to U$ la inyección canónica, que corresponde al homomorfismo de restricción $A\to A'$; se tiene $\mathcal{S}|U'=j^*(\mathcal{S}|U)$ y, por consiguiente, $S'=\Gamma(U',\mathcal{S})$ se identifica canónicamente con $S\otimes_A A'$ \sref[I]{I.1.6.4-fr}.
Se concluye de \sref{II.2.8.10-fr} que $X_{U'}$ se identifica canónicamente con $X_U\times_U U'$ y, por tanto, también con $f_U^{-1}(U')$, donde $f_U$ denota el morfismo estructural $X_U\to U$ \sref[I]{I.4.4.1-fr}.
Denotaremos por $\sigma_{U',U}$ el isomorfismo canónico $f_U^{-1}(U')\xrightarrow{\sim}X_{U'}$ así definido, y por $\rho_{U',U}$ la inmersión abierta $X_{U'}\to X_U$ obtenida componiendo $\sigma_{U',U}^{-1}$ con la inyección canónica $f_U^{-1}(U')\to X_U$.
Es inmediato que, si $U''\subset U'$ es otro abierto afín de $Y$, se tiene $\rho_{U'',U}=\rho_{U',U}\circ\rho_{U'',U'}$.
\end{env}
% Fuente francesa: PMIHES_1961__8__5_0, printed p. 50.

\begin{proposition}[3.1.2]
\label{II.3.1.2-fr}
Sea $Y$ un preesquema.
Para toda $\mathcal{O}_Y$-álgebra graduada cuasicoherente en grados positivos $\mathcal{S}$, existe un único preesquema $X$ sobre $Y$ (salvo $Y$-isomorfismo) que posee la propiedad siguiente: si $f$ es el morfismo estructural $X\to Y$, para todo abierto afín $U$ de $Y$ existe un isomorfismo $\eta_U$ del preesquema inducido $f^{-1}(U)$ sobre $X_U=\operatorname{Proj}(\Gamma(U,\mathcal{S}))$ tal que, si $V$ es otro abierto afín de $Y$ contenido en $U$, el diagrama
\[
\label{II.3.1.2.1-fr}
  \xymatrix{
    f^{-1}(V)\ar[r]^{\overset{\eta_V}{\sim}}\ar[d]
    & X_V\ar[d]^{\rho_{V,U}}\\
    f^{-1}(U)\ar[r]^{\sim}_{\eta_U}
    & X_U
  }
\tag{3.1.2.1}
\]
es conmutativo.
\end{proposition}

\begin{proof}
Para dos abiertos afines $U$ y $V$ de $Y$, sea $X_{U,V}$ el preesquema inducido sobre $f_U^{-1}(U\cap V)$ por $X_U$; vamos a definir un $Y$-isomorfismo $\theta_{U,V}:X_{V,U}\xrightarrow{\sim}X_{U,V}$.
Para ello, consideremos un abierto afín $W\subset U\cap V$: componiendo los isomorfismos
\[
  f_U^{-1}(W)\xrightarrow{\sigma_{W,U}}X_W
  \xrightarrow{\sigma_{W,V}^{-1}}f_V^{-1}(W),
\]
se obtiene un isomorfismo $\tau_W$, y se comprueba inmediatamente que, si $W'\subset W$ es un abierto afín, $\tau_{W'}$ es la restricción a $f_U^{-1}(W')$ de $\tau_W$; los $\tau_W$ son, pues, las restricciones de un $Y$-isomorfismo $\theta_{V,U}$.
Además, si $U$, $V$ y $W$ son tres abiertos afines de $Y$, y si $\theta'_{U,V}$, $\theta'_{V,W}$ y $\theta'_{U,W}$ son las restricciones respectivas de $\theta_{U,V}$, $\theta_{V,W}$ y $\theta_{U,W}$ a las imágenes inversas de $U\cap V\cap W$ en $X_V$, $X_W$ y $X_W$, resulta de las definiciones anteriores que $\theta'_{U,V}\circ\theta'_{V,W}=\theta'_{U,W}$.
La existencia de un $X$ que satisfaga las propiedades enunciadas resulta entonces de \sref[I]{I.2.3.1-fr}; su unicidad, salvo $Y$-isomorfismo, es trivial, teniendo en cuenta \sref{II.3.1.2.1-fr}.
\end{proof}

\begin{env}[3.1.3]
\label{II.3.1.3-fr}
Diremos que el preesquema $X$ definido en \sref{II.3.1.2-fr} es el \emph{espectro homogéneo} de la $\mathcal{O}_Y$-álgebra graduada cuasicoherente $\mathcal{S}$, y lo denotaremos por $\operatorname{Proj}(\mathcal{S})$.
Es inmediato que $\operatorname{Proj}(\mathcal{S})$ está \emph{separado sobre $Y$} (\sref{II.2.4.2-fr} y \sref[I]{I.5.5.5-fr}); si $\mathcal{S}$ es una $\mathcal{O}_Y$-álgebra \emph{de tipo finito} \sref[I]{I.9.6.2-fr}, entonces $\operatorname{Proj}(\mathcal{S})$ es \emph{de tipo finito} sobre $Y$ (\sref{II.2.7.1-fr}[(ii)] y \sref[I]{I.6.3.1-fr}).

Si $f$ es el morfismo estructural $X\to Y$, es inmediato que, para todo preesquema inducido por $Y$ sobre una parte abierta $U$ de $Y$, $f^{-1}(U)$ se identifica con el espectro homogéneo $\operatorname{Proj}(\mathcal{S}|U)$.
\end{env}
% Fuente francesa: PMIHES_1961__8__5_0, printed pp. 50-51.

\begin{proposition}[3.1.4]
\label{II.3.1.4-fr}
Sea $f\in\Gamma(Y,\mathcal{S}_d)$, con $d>0$.
Existe una parte abierta $X_f$ del espacio subyacente a $X=\operatorname{Proj}(\mathcal{S})$ que posee la propiedad siguiente: para todo abierto afín $U$ de $Y$, se tiene
\oldpage[II]{51}
 $X_f\cap\varphi^{-1}(U)=D_+(f|U)$ en $\varphi^{-1}(U)$ identificado con $X_U=\operatorname{Proj}(\Gamma(U,\mathcal{S}))$, donde $\varphi$ denota el morfismo estructural $X\to Y$.

Además, el preesquema inducido sobre $X_f$ es afín sobre $Y$ y es canónicamente isomorfo a $\operatorname{Spec}(\mathcal{S}^{(d)}/(f-1)\mathcal{S}^{(d)})$ \sref{II.1.3.1-fr}.
\end{proposition}

\begin{proof}
Se tiene $f|U\in\Gamma(U,\mathcal{S}_d)=(\Gamma(U,\mathcal{S}))_d$.
Si $U$ y $U'$ son abiertos afines de $Y$ tales que $U'\subset U$, $f|U'$ es la imagen de $f|U$ mediante el homomorfismo de restricción
\[
  \Gamma(U,\mathcal{S})\longrightarrow\Gamma(U',\mathcal{S}),
\]
por tanto, $D_+(f|U')$ es igual (con las notaciones de \sref{II.3.1.1-fr}) al preesquema inducido sobre la imagen inversa $\rho_{U',U}^{-1}(D_+(f|U))$ en $X_{U'}$ \sref{II.2.8.1-fr}; de aquí resulta la primera afirmación.
Además, el preesquema inducido sobre $D_+(f|U)$ por $X_U$ se identifica canónicamente con $\operatorname{Spec}((\Gamma(U,\mathcal{S}))_{(f|U)})$, siendo estas identificaciones compatibles con los homomorfismos de restricción \sref{II.2.8.1-fr}; la segunda afirmación resulta entonces de \sref{II.2.2.5-fr} y de la conmutatividad del diagrama \sref{II.2.8.2.1-fr}.
\end{proof}

Se dice también que $X_f$ (como abierto del espacio subyacente $X$) es el conjunto de los $x\in X$ donde $f$ \emph{no se anula}.

\begin{corollary}[3.1.5]
\label{II.3.1.5-fr}
Si $f\in\Gamma(Y,\mathcal{S}_d)$ y $g\in\Gamma(Y,\mathcal{S}_e)$, se tiene
\[
\label{II.3.1.5.1-fr}
  X_{fg}=X_f\cap X_g.
\tag{3.1.5.1}
\]
\end{corollary}

\begin{proof}
Basta considerar la intersección de los dos miembros con un conjunto $\varphi^{-1}(U)$, donde $U$ es un abierto afín de $Y$, y aplicar la fórmula \sref{II.2.3.3.2-fr}.
\end{proof}

\begin{corollary}[3.1.6]
\label{II.3.1.6-fr}
Sea $(f_\alpha)$ una familia de secciones de $\mathcal{S}$ sobre $Y$ tal que $f_\alpha\in\Gamma(Y,\mathcal{S}_{d_\alpha})$; si el haz de ideales de $\mathcal{S}$ generado por esta familia \sref[0]{0.5.1.1-fr} contiene todos los $\mathcal{S}_n$ a partir de cierto rango, el espacio subyacente $X$ es la unión de los $X_{f_\alpha}$.
\end{corollary}

\begin{proof}
Para todo abierto afín $U$ de $Y$, $\varphi^{-1}(U)$ es la unión de los $X_{f_\alpha}\cap\varphi^{-1}(U)$ \sref{II.2.3.14-fr}.
\end{proof}
% Fuente francesa: PMIHES_1961__8__5_0, printed p. 51.

\begin{corollary}[3.1.7]
\label{II.3.1.7-fr}
Sea $\mathcal{A}$ una $\mathcal{O}_Y$-álgebra cuasicoherente; pongamos
\[
  \mathcal{S}=\mathcal{A}[T]
  =\mathcal{A}\otimes_{\mathbf{Z}}\mathbf{Z}[T]
\]
donde $T$ es una indeterminada (considerándose $\mathbf{Z}$ y $\mathbf{Z}[T]$ como haces simples sobre $Y$).
Entonces $X=\operatorname{Proj}(\mathcal{S})$ se identifica canónicamente con $\operatorname{Spec}(\mathcal{A})$.
En particular, $\operatorname{Proj}(\mathcal{O}_Y[T])$ se identifica con $Y$.
\end{corollary}

\begin{proof}
Aplicando \sref{II.3.1.6-fr} a la única sección $f\in\Gamma(Y,\mathcal{S})$ que es igual a $T$ en cada punto de $Y$, se ve que $X_f=X$.
Además, aquí se tiene $d=1$, y $\mathcal{S}^{(1)}/(f-1)\mathcal{S}^{(1)}
=\mathcal{S}/(f-1)\mathcal{S}$ es canónicamente isomorfa a $\mathcal{A}$, de donde resulta el corolario \sref{II.3.1.4-fr}.
\end{proof}

Sea $g\in\Gamma(Y,\mathcal{O}_Y)$; si se toma $\mathcal{S}=\mathcal{O}_Y[T]$, se tiene entonces $g\in\Gamma(Y,\mathcal{S}_0)$; sea
\[
  h=gT\in\Gamma(Y,\mathcal{S}_1).
\]
Si $X=\operatorname{Proj}(\mathcal{S})$, la identificación canónica definida en \sref{II.3.1.7-fr} identifica $X_h$ con el abierto $Y_g$ de $Y$ (en el sentido de \sref[0]{0.5.5.2-fr}): en efecto, se puede limitar al caso en que $Y=\operatorname{Spec}(A)$ es afín, y todo se reduce entonces (teniendo en cuenta \sref{II.2.2.5-fr}) al hecho de que el anillo de fracciones $A_g$ se identifica canónicamente con $A[T]/(gT-1)A[T]$ \sref[0]{0.1.2.3-fr}.
% Fuente francesa: PMIHES_1961__8__5_0, printed pp. 51-52.

\begin{proposition}[3.1.8]
\label{II.3.1.8-fr}
Sea $\mathcal{S}$ una $\mathcal{O}_Y$-álgebra cuasicoherente graduada en grados positivos.
Entonces:
\begin{enumerate}
  \item[(i)] Para todo $d>0$, existe un $Y$-isomorfismo canónico de $\operatorname{Proj}(\mathcal{S})$ sobre $\operatorname{Proj}(\mathcal{S}^{(d)})$.
\oldpage[II]{52}
  \item[(ii)] Sea $\mathcal{S}'$ la $\mathcal{O}_Y$-álgebra graduada suma directa de $\mathcal{O}_Y$ y de los $\mathcal{S}_n$ ($n>0$); entonces $\operatorname{Proj}(\mathcal{S}')$ y $\operatorname{Proj}(\mathcal{S})$ son canónicamente $Y$-isomorfos.
  \item[(iii)] Sea $\mathcal{L}$ un $\mathcal{O}_Y$-módulo invertible \sref[0]{0.5.4.1-fr}, y sea $\mathcal{S}_{(\mathcal{L})}$ la $\mathcal{O}_Y$-álgebra graduada suma directa de los $\mathcal{S}_d\otimes\mathcal{L}^{\otimes d}$ ($d\geq0$); entonces $\operatorname{Proj}(\mathcal{S})$ y $\operatorname{Proj}(\mathcal{S}_{(\mathcal{L})})$ son canónicamente $Y$-isomorfos.
\end{enumerate}
\end{proposition}

\begin{proof}
En cada uno de los tres casos, basta definir el isomorfismo localmente sobre $Y$, pues la verificación de las compatibilidades con las operaciones de restricción de un abierto a otro menor es trivial.
Puede suponerse, pues, que $Y$ es afín, y entonces (i) resulta de \sref{II.2.4.7-fr}[(i)] y (ii) de \sref{II.2.4.8-fr}.
En cuanto a (iii), si se supone además que $\mathcal{L}$ es isomorfo a $\mathcal{O}_Y$ (lo que está permitido, pues la cuestión es local sobre $Y$), el isomorfismo entre $\operatorname{Proj}(\mathcal{S})$ y $\operatorname{Proj}(\mathcal{S}_{(\mathcal{L})})$ es evidente; para definir un isomorfismo \emph{canónico}, sean $Y=\operatorname{Spec}(A)$ y $\mathcal{S}=\widetilde{S}$, donde $S$ es una $A$-álgebra graduada, y sea $c$ un generador del $A$-módulo libre $L$ tal que $\mathcal{L}=\widetilde{L}$; entonces, para todo $n>0$, $x_n\to x_n\otimes c^{\otimes n}$ es un $A$-isomorfismo de $S_n$ sobre $S_n\otimes L^{\otimes n}$, y estos $A$-isomorfismos definen un $A$-isomorfismo de álgebras graduadas $\varphi_c:S\to S_{(L)}
=\bigoplus_{n\geq0}S_n\otimes L^{\otimes n}$.
Sea entonces $f\in S_+$ homogéneo de grado $d$; para todo $x\in S_{nd}$ se tiene
$(x\otimes c^{nd})/(f\otimes c^d)^n
=(x\otimes(\varepsilon c)^{nd})/(f\otimes(\varepsilon c)^d)^n$
para toda unidad $\varepsilon\in A$, lo que muestra que el isomorfismo $S_{(f)}\to(S_{(L)})_{(f\otimes c^d)}$ deducido de $\varphi_c$ es \emph{independiente} del generador $c$ considerado de $L$, y termina así la demostración.
\end{proof}
% Fuente francesa: PMIHES_1961__8__5_0, printed pp. 52-53.

\begin{env}[3.1.9]
\label{II.3.1.9-fr}
Recordemos (\sref[0]{0.4.1.3-fr} y \sref[I]{I.1.3.14-fr}) que, para que la $\mathcal{O}_Y$-álgebra graduada cuasicoherente $\mathcal{S}$ esté \emph{generada por el $\mathcal{O}_Y$-módulo $\mathcal{S}_1$}, es necesario y suficiente que exista una cubierta $(U_\alpha)$ de $Y$ por abiertos afines tales que el álgebra graduada $\Gamma(U_\alpha,\mathcal{S})$ sobre $\Gamma(U_\alpha,\mathcal{S}_0)$ esté generada por el conjunto $\Gamma(U_\alpha,\mathcal{S}_1)$ de sus elementos homogéneos de grado $1$.
Para todo abierto $V$ de $Y$, $\mathcal{S}|V$ está entonces generada por el $(\mathcal{O}_Y|V)$-módulo $\mathcal{S}_1|V$.
\end{env}

\begin{proposition}[3.1.10]
\label{II.3.1.10-fr}
Supongamos que existe una cubierta abierta afín finita $(U_i)$ de $Y$ tal que cada álgebra graduada $\Gamma(U_i,\mathcal{S})$ es de tipo finito sobre $\Gamma(U_i,\mathcal{O}_Y)$.
Entonces existe $d>0$ tal que $\mathcal{S}^{(d)}$ está generada por $\mathcal{S}_d$, siendo $\mathcal{S}_d$ un $\mathcal{O}_Y$-módulo de tipo finito.
\end{proposition}

\begin{proof}
En efecto, resulta de \sref{II.2.1.6-fr}[(v)] que, para cada $i$, existe un entero $m_i$ tal que $\Gamma(U_i,\mathcal{S}_{nm_i})=(\Gamma(U_i,\mathcal{S}_{m_i}))^n$ para todo $n>0$; basta tomar como $d$ un múltiplo común de todos los $m_i$, teniendo en cuenta \sref{II.2.1.6-fr}[(i)].
\end{proof}

\begin{corollary}[3.1.11]
\label{II.3.1.11-fr}
Bajo las hipótesis de \sref{II.3.1.10-fr}, $\operatorname{Proj}(\mathcal{S})$ es $Y$-isomorfo a un espectro homogéneo $\operatorname{Proj}(\mathcal{S}')$, donde $\mathcal{S}'$ es una $\mathcal{O}_Y$-álgebra graduada generada por $\mathcal{S}'_1$, suponiéndose que $\mathcal{S}'_1$ es un $\mathcal{O}_Y$-módulo de tipo finito.
\end{corollary}

\begin{proof}
Basta tomar $\mathcal{S}'=\mathcal{S}^{(d)}$, donde $d$ satisface la propiedad de \sref{II.3.1.10-fr}, y aplicar \sref{II.3.1.8-fr}[(i)].
\end{proof}

\begin{env}[3.1.12]
\label{II.3.1.12-fr}
Si $\mathcal{S}$ es una $\mathcal{O}_Y$-álgebra cuasicoherente graduada en grados positivos, se sabe \sref[I]{I.5.1.1-fr} que su \emph{nilradical} $\mathcal{N}$ es un $\mathcal{O}_Y$-módulo cuasicoherente; diremos que $\mathcal{N}_+=\mathcal{N}\cap\mathcal{S}_+$ es el \emph{nilradical de $\mathcal{S}_+$}; es un $\mathcal{S}_0$-módulo graduado cuasicoherente, pues se puede reducir inmediatamente al caso en que $Y$ es afín, y la afirmación resulta entonces de \sref{II.2.1.10-fr}.
Para todo $y\in Y$, $(\mathcal{N}_+)_y$ es entonces el nilradical de $(\mathcal{S}_+)_y=(\mathcal{S}_y)_+$ \sref[I]{I.5.1.1-fr}.
Se dice que la $\mathcal{O}_Y$-álgebra graduada $\mathcal{S}$ es \emph{esencialmente reducida} si $\mathcal{N}_+=0$, lo que equivale
\oldpage[II]{53}
a decir que $\mathcal{S}_y$ es una $\mathcal{O}_y$-álgebra graduada esencialmente reducida para todo $y\in Y$.
Para toda $\mathcal{O}_Y$-álgebra graduada $\mathcal{S}$, $\mathcal{S}/\mathcal{N}_+$ es esencialmente reducida.

Diremos que $\mathcal{S}$ es \emph{íntegra} si $\mathcal{S}_y$ es un anillo íntegro para todo $y\in Y$ y si, además, $(\mathcal{S}_y)_+=(\mathcal{S}_+)_y\neq0$ para todo $y\in Y$.
\end{env}
% Fuente francesa: PMIHES_1961__8__5_0, printed p. 53.

\begin{proposition}[3.1.13]
\label{II.3.1.13-fr}
Sea $\mathcal{S}$ una $\mathcal{O}_Y$-álgebra graduada en grados positivos.
Si $X=\operatorname{Proj}(\mathcal{S})$, el $Y$-esquema $X_{\mathrm{red}}$ es canónicamente isomorfo a $\operatorname{Proj}(\mathcal{S}/\mathcal{N}_+)$; en particular, si $\mathcal{S}$ es esencialmente reducida, $X$ es reducido.
\end{proposition}

\begin{proof}
El hecho de que $X'=\operatorname{Proj}(\mathcal{S}/\mathcal{N}_+)$ sea reducido resulta inmediatamente de \sref{II.2.4.4-fr}[(i)], pues la propiedad es local; además, para todo abierto afín $U\subset Y$, ${\varphi'}^{-1}(U)$ es igual a $(\varphi^{-1}(U))_{\mathrm{red}}$ (denotando por $\varphi$ y $\varphi'$ los morfismos estructurales $X\to Y$ y $X'\to Y$, respectivamente); se comprueba inmediatamente que los $U$-morfismos canónicos ${\varphi'}^{-1}(U)\to\varphi^{-1}(U)$ son compatibles con las operaciones de restricción y definen, por tanto, una inmersión cerrada $X'\to X$ que es un homeomorfismo de los espacios subyacentes; de aquí resulta la conclusión \sref[I]{I.5.1.2-fr}.
\end{proof}
% Fuente francesa: PMIHES_1961__8__5_0, printed pp. 53-54.

\begin{proposition}[3.1.14]
\label{II.3.1.14-fr}
Sean $Y$ un preesquema íntegro y $\mathcal{S}$ una $\mathcal{O}_Y$-álgebra graduada cuasicoherente tal que $\mathcal{S}_0=\mathcal{O}_Y$.
\begin{enumerate}
  \item[(i)] Si $\mathcal{S}$ es íntegra \sref{II.3.1.12-fr}, entonces $X=\operatorname{Proj}(\mathcal{S})$ es íntegro y el morfismo estructural $\varphi:X\to Y$ es dominante.
  \item[(ii)] Supongamos además que $\mathcal{S}$ es esencialmente reducida.
    Entonces, recíprocamente, si $X$ es íntegro y $\varphi$ es dominante, $\mathcal{S}$ es íntegra.
\end{enumerate}
\end{proposition}

\begin{proof}
\medskip\noindent
\begin{enumerate}
  \item[(i)] Si $(U_\alpha)$ es una base de $Y$ formada por abiertos afines no vacíos, basta demostrar la proposición cuando se sustituye $Y$ por uno de los $U_\alpha$ y $\mathcal{S}$ por $\mathcal{S}|U_\alpha$: en efecto, por una parte, resultará que los espacios subyacentes $\varphi^{-1}(U_\alpha)$ son abiertos irreducibles (y, por tanto, no vacíos) de $X$ tales que $\varphi^{-1}(U_\alpha)\cap\varphi^{-1}(U_\beta)\neq\varnothing$ para todo par de índices (pues $U_\alpha\cap U_\beta$ contiene algún $U_\gamma$), de modo que $X$ es irreducible \sref[0]{0.2.1.4-fr}; por otra parte, $X$ será reducido, pues esta propiedad es local, y por tanto $X$ será íntegro y $\varphi(X)$ será denso en $Y$.

    Supongamos, pues, que $Y=\operatorname{Spec}(A)$, donde $A$ es íntegro \sref[I]{I.5.1.4-fr}, y que $\mathcal{S}=\widetilde{S}$, donde $S$ es una $A$-álgebra graduada; la hipótesis es que, para todo $y\in Y$, $\widetilde{S}_y=S_y$ es un anillo graduado íntegro tal que $(S_y)_+\neq0$.
    Basta demostrar que $S$ es un anillo \emph{íntegro}, pues entonces se tendrá $S_+\neq0$ y podrá aplicarse \sref{II.2.4.4-fr}[(ii)].
    Sean $f$, $g$ dos elementos $\neq0$ de $S$, y supongamos que $fg=0$; para todo $y\in Y$ se tiene entonces $(f/1)(g/1)=0$ en $S_y$ y, por hipótesis, $f/1=0$ o $g/1=0$.
    Supongamos, por ejemplo, que $f/1=0$ en $S_y$; esto significa que existe $a\in A$ tal que $a\notin\mathfrak{j}_y$ y $af=0$; se tiene entonces, para \emph{todo} $z\in Y$, $(a/1)(f/1)=0$ en el anillo \emph{íntegro} $S_z$ y, como $a/1\neq0$ (pues $A$ es íntegro), $f/1=0$, lo que implica $f=0$.
  \item[(ii)] Como la cuestión es local sobre $Y$, puede suponerse de nuevo que $Y=\operatorname{Spec}(A)$, con $A$ íntegro, y que $\mathcal{S}=\widetilde{S}$.
    Por hipótesis, para todo $y\in Y$, $(S_y)_+$ no contiene ningún elemento nilpotente $\neq0$, y lo mismo ocurre con $(S_0)_y=A_y$ por hipótesis; por tanto, $S_y$ es un anillo reducido para todo $y\in Y$, y se concluye primero que el propio $S$ es reducido \sref[I]{I.5.1.1-fr}.
    La hipótesis de que $X$ es íntegro implica, por otra parte, que $S$ es esencialmente íntegro \sref{II.2.4.4-fr}[(ii)], y todo se reduce a demostrar que el anulador $\mathfrak{J}$ de $S_+$ en $A=S_0$ es
\oldpage[II]{54}
 $0$ \sref{II.2.1.11-fr}.
    Si no fuera así, se tendría $(S_h)_+=0$ para algún $h\neq0$ en $\mathfrak{J}$ y, por tanto, \sref{II.3.1.1-fr}, $\varphi^{-1}(D(h))=\varnothing$, y $\varphi(X)$ no sería denso en $Y$, contrariamente a la hipótesis (pues $D(h)\neq\varnothing$, ya que $h$ no es nilpotente).

\end{enumerate}
\end{proof}
% Fuente francesa: PMIHES_1961__8__5_0, printed pp. 54-55.

\subsection{Haz sobre $\operatorname{Proj}(\mathcal{S})$ asociado a un $\mathcal{S}$-módulo graduado}
\label{subsection:II.3.2-fr}

\begin{env}[3.2.1]
\label{II.3.2.1-fr}
Sean $Y$ un preesquema, $\mathcal{S}$ una $\mathcal{O}_Y$-álgebra cuasicoherente graduada en grados positivos y $\mathcal{M}$ un $\mathcal{S}$-módulo graduado cuasicoherente (sobre $(Y,\mathcal{O}_Y)$ o, lo que equivale a lo mismo \sref[I]{I.9.6.1-fr}, sobre el espacio anillado $(Y,\mathcal{S})$).
Con las notaciones de \sref{II.3.1.1-fr}, denotaremos por $\widetilde{\mathcal{M}}_U$ el $\mathcal{O}_{X_U}$-módulo cuasicoherente $(\Gamma(U,\mathcal{M}))\supertilde$; para $U'\subset U$, $\Gamma(U',\mathcal{M})$ se identifica canónicamente con $\Gamma(U,\mathcal{M})\otimes_A A'$ \sref[I]{I.1.6.4-fr}; por tanto, se tiene $\widetilde{\mathcal{M}}_{U'}
=\rho_{U',U}^*(\widetilde{\mathcal{M}}_U)$ \sref{II.2.8.11-fr}.
\end{env}

\begin{proposition}[3.2.2]
\label{II.3.2.2-fr}
Existe sobre $\operatorname{Proj}(\mathcal{S})=X$ un único $\mathcal{O}_X$-módulo cuasicoherente $\widetilde{\mathcal{M}}$ tal que, para todo abierto afín $U$ de $Y$, se tenga $\eta_U^*((\Gamma(U,\mathcal{M}))\supertilde)
=\widetilde{\mathcal{M}}|f^{-1}(U)$ (donde $\eta_U$ denota el isomorfismo $f^{-1}(U)\xrightarrow{\sim}\operatorname{Proj}(\Gamma(U,\mathcal{S}))$), siendo $f$ el morfismo estructural $X\to Y$.
\end{proposition}

\begin{proof}
Como $\rho_{U',U}$ se identifica con el morfismo de inyección $f^{-1}(U')\to f^{-1}(U)$ \sref{II.3.1.2.1-fr}, la proposición resulta inmediatamente de la relación $\widetilde{\mathcal{M}}_{U'}=\rho_{U',U}^*(\widetilde{\mathcal{M}}_U)$ y del principio de pegado de haces \sref[0]{0.3.3.1-fr}.
\end{proof}

Se dice que $\widetilde{\mathcal{M}}$ es el $\mathcal{O}_X$-módulo \emph{asociado al} $\mathcal{S}$-módulo graduado cuasicoherente $\mathcal{M}$.

\begin{proposition}[3.2.3]
\label{II.3.2.3-fr}
Sea $\mathcal{M}$ un $\mathcal{S}$-módulo graduado cuasicoherente, y sea $f\in\Gamma(Y,\mathcal{S}_d)$ ($d>0$).
Si $\xi_f$ es el isomorfismo canónico de $X_f$ sobre el $Y$-preesquema $Z_f=\operatorname{Spec}(\mathcal{S}^{(d)}/(f-1)\mathcal{S}^{(d)})$ \sref{II.3.1.4-fr}, entonces $(\xi_f)_*(\widetilde{\mathcal{M}}|X_f)$ es el $\mathcal{O}_{Z_f}$-módulo $(\mathcal{M}^{(d)}/(f-1)\mathcal{M}^{(d)})\supertilde$ \sref{II.1.4.3-fr}.
\end{proposition}

\begin{proof}
Como la cuestión es local sobre $Y$, se reduce inmediatamente a \sref{II.2.2.5-fr}, teniendo en cuenta la conmutatividad del diagrama \sref{II.2.8.12.1-fr}.
\end{proof}

\begin{proposition}[3.2.4]
\label{II.3.2.4-fr}
El $\mathcal{O}_X$-módulo $\widetilde{\mathcal{M}}$ es un funtor covariante aditivo exacto en $\mathcal{M}$, de la categoría de los $\mathcal{S}$-módulos graduados cuasicoherentes en la de los $\mathcal{O}_X$-módulos cuasicoherentes, y conmuta con los límites inductivos y las sumas directas.
\end{proposition}

\begin{proof}
Como la cuestión es local sobre $Y$, se reduce a \sref[I]{I.1.3.11-fr}, \sref[I]{I.1.3.9-fr} y \sref{II.2.5.4-fr}.
\end{proof}

En particular, si $\mathcal{N}$ es un sub-$\mathcal{S}$-módulo graduado cuasicoherente de $\mathcal{M}$, $\widetilde{\mathcal{N}}$ se identifica canónicamente con un sub-$\mathcal{O}_X$-módulo cuasicoherente de $\widetilde{\mathcal{M}}$; más particularmente, para todo haz graduado cuasicoherente $\mathcal{J}$ de ideales de $\mathcal{S}$, $\widetilde{\mathcal{J}}$ es un haz cuasicoherente de ideales de $\mathcal{O}_X$.

Si $\mathcal{M}$ es un $\mathcal{S}$-módulo graduado cuasicoherente y $\mathcal{I}$ es un haz cuasicoherente de ideales de $\mathcal{O}_Y$, entonces $\mathcal{I}\mathcal{M}$ es un sub-$\mathcal{S}$-módulo graduado cuasicoherente de $\mathcal{M}$, y se tiene
\[
\label{II.3.2.4.1-fr}
  (\mathcal{I}\mathcal{M})\supertilde
  =\mathcal{I}.\widetilde{\mathcal{M}}
\tag{3.2.4.1}
\]
(donde el segundo miembro se define como en \sref[0]{0.4.3.5-fr}).
Basta verificar esta fórmula en el caso en que $Y=\operatorname{Spec}(A)$ es afín, $\mathcal{S}=\widetilde{S}$, donde $S$ es una $A$-álgebra graduada, $\mathcal{M}=\widetilde{M}$, 
\oldpage[II]{55}
donde $M$ es un $S$-módulo graduado, y $\mathcal{I}=\widetilde{\mathfrak{J}}$, donde $\mathfrak{J}$ es un ideal de $A$.

Para todo elemento homogéneo $f$ de $S_+$, la restricción a $D_+(f)=\operatorname{Spec}(S_{(f)})$ del primer miembro de \sref{II.3.2.4.1-fr} está asociada a $(\mathfrak{J}M)_{(f)}=\mathfrak{J}.M_{(f)}$, y lo mismo ocurre con la restricción del segundo miembro, en virtud de \sref[I]{I.1.3.13-fr} y \sref[I]{I.1.6.9-fr}.
% Fuente francesa: PMIHES_1961__8__5_0, printed p. 55.

\begin{proposition}[3.2.5]
\label{II.3.2.5-fr}
Sea $f\in\Gamma(Y,\mathcal{S}_d)$ ($d>0$).
Sobre el abierto $X_f$, el $(\mathcal{O}_X|X_f)$-módulo $(\mathcal{S}(nd))\supertilde|X_f$ es canónicamente isomorfo a $\mathcal{O}_X|X_f$ para todo $n\in\mathbf{Z}$.
En particular, si la $\mathcal{O}_Y$-álgebra $\mathcal{S}$ está generada por $\mathcal{S}_1$ \sref{II.3.1.9-fr}, los $\mathcal{O}_X$-módulos $(\mathcal{S}(n))\supertilde$ son invertibles para todo $n\in\mathbf{Z}$.
\end{proposition}

\begin{proof}
En efecto, para todo abierto afín $U$ de $Y$, se ha definido en \sref{II.2.5.7-fr} un isomorfismo canónico de $(\mathcal{S}(nd))\supertilde|(X_f\cap\varphi^{-1}(U))$ sobre $\mathcal{O}_X|(X_f\cap\varphi^{-1}(U))$, teniendo en cuenta \sref{II.3.1.4-fr} (donde $\varphi$ es el morfismo estructural $X\to Y$); es inmediato comprobar que estos isomorfismos son compatibles con la restricción de $U$ a un abierto afín $U'\subset U$, de donde resulta la primera afirmación.
Para demostrar la segunda, basta observar que, si $\mathcal{S}$ está generada por $\mathcal{S}_1$, existe una cubierta $(U_\alpha)$ de $Y$ por abiertos afines tales que $\Gamma(U_\alpha,\mathcal{S})$ está generada por $\Gamma(U_\alpha,\mathcal{S})_1=\Gamma(U_\alpha,\mathcal{S}_1)$; puede aplicarse entonces el resultado de \sref{II.2.5.9-fr}, pues la propiedad de ser invertible es local.
\end{proof}

Se pondrá de nuevo, para todo $n\in\mathbf{Z}$,
\[
\label{II.3.2.5.1-fr}
  \mathcal{O}_X(n)=(\mathcal{S}(n))\supertilde
\tag{3.2.5.1}
\]
y, para todo $\mathcal{O}_X$-módulo $\mathcal{F}$,
\[
\label{II.3.2.5.2-fr}
  \mathcal{F}(n)=\mathcal{F}\otimes_{\mathcal{O}_X}\mathcal{O}_X(n).
\tag{3.2.5.2}
\]

Resulta inmediatamente de estas definiciones que, para todo abierto $U$ de $Y$, se tiene
\[
  ((\mathcal{S}|U)(n))\supertilde
  =\mathcal{O}_X(n)|f^{-1}(U)
\]
donde $f$ es el morfismo estructural $X\to Y$.
% Fuente francesa: PMIHES_1961__8__5_0, printed pp. 55-56.

\begin{proposition}[3.2.6]
\label{II.3.2.6-fr}
Sean $\mathcal{M}$ y $\mathcal{N}$ dos $\mathcal{S}$-módulos graduados cuasicoherentes.
Existe un homomorfismo canónico funtorial (en $\mathcal{M}$ y $\mathcal{N}$)
\[
\label{II.3.2.6.1-fr}
  \lambda:\widetilde{\mathcal{M}}
  \otimes_{\mathcal{O}_X}\widetilde{\mathcal{N}}
  \longrightarrow(\mathcal{M}\otimes_{\mathcal{S}}\mathcal{N})\supertilde
\tag{3.2.6.1}
\]
y un homomorfismo canónico funtorial (en $\mathcal{M}$ y $\mathcal{N}$)
\[
\label{II.3.2.6.2-fr}
  \mu:(\mathcal{H}om_{\mathcal{S}}(\mathcal{M},\mathcal{N}))\supertilde
  \longrightarrow
  \mathcal{H}om_{\mathcal{O}_X}
  (\widetilde{\mathcal{M}},\widetilde{\mathcal{N}}).
\tag{3.2.6.2}
\]

Además, si $\mathcal{S}$ está generada por $\mathcal{S}_1$ \sref{II.3.1.9-fr}, $\lambda$ es un isomorfismo; si, además, $\mathcal{M}$ admite una presentación finita \sref{II.3.1.1-fr}, $\mu$ es un isomorfismo.
\end{proposition}

\begin{proof}
Los homomorfismos $\lambda$ y $\mu$ se han definido en \sref{II.2.5.11.2-fr} y \sref{II.2.5.12.2-fr} cuando $Y$ es afín; como estas definiciones son locales, se transportan inmediatamente al caso general considerado aquí, teniendo en cuenta \sref{II.2.8.14-fr}.
\end{proof}

\begin{corollary}[3.2.7]
\label{II.3.2.7-fr}
Si $\mathcal{S}$ está generada por $\mathcal{S}_1$, se tiene, cualesquiera que sean $m$, $n$ en $\mathbf{Z}$,
\[
\label{II.3.2.7.1-fr}
  \mathcal{O}_X(m)\otimes_{\mathcal{O}_X}\mathcal{O}_X(n)
  =\mathcal{O}_X(m+n)
\tag{3.2.7.1}
\]
\[
\label{II.3.2.7.2-fr}
  \mathcal{O}_X(n)=(\mathcal{O}_X(1))^{\otimes n}
\tag{3.2.7.2}
\]
salvo isomorfismos canónicos.
\end{corollary}

\oldpage[II]{56}

\begin{corollary}[3.2.8]
\label{II.3.2.8-fr}
Supongamos que $\mathcal{S}$ está generada por $\mathcal{S}_1$. Para todo $\mathcal{S}$-módulo graduado $\mathcal{M}$ y todo $n\in\mathbf{Z}$ se tiene
\[
\label{II.3.2.8.1-fr}
  (\mathcal{M}(n))^{\sim}=\widetilde{\mathcal{M}}(n)
\tag{3.2.8.1}
\]
salvo isomorfismo canónico.
\end{corollary}

\begin{proof}
Esto resulta de las propiedades correspondientes en el caso en que $Y$ es afín (\sref{II.2.5.14-fr} y \sref{II.2.5.15-fr}), así como de \sref{II.2.8.11-fr}.
\end{proof}
% Fuente francesa: PMIHES_1961__8__5_0, printed p. 56.

\begin{remarks}[3.2.9]
\label{II.3.2.9-fr}
\medskip\noindent
\begin{enumerate}
  \item Si $\mathcal{S}=\mathcal{A}[T]$, siendo $\mathcal{A}$ una $\mathcal{O}_Y$-álgebra cuasicoherente \sref{II.3.1.7-fr}, se comprueba inmediatamente que todos los $\mathcal{O}_X$-módulos invertibles $\mathcal{O}_X(n)$ son canónicamente isomorfos a $\mathcal{O}_X$.

    Además, sea $\mathcal{N}$ un $\mathcal{A}$-módulo cuasicoherente, y pongamos $\mathcal{M}=\mathcal{N}\otimes_{\mathcal{A}}\mathcal{A}[T]$.
    Resulta entonces de \sref{II.3.2.3-fr} y \sref{II.3.1.7-fr} que, en la identificación canónica de $X=\operatorname{Proj}(\mathcal{A}[T])$ con $X'=\operatorname{Spec}(\mathcal{A})$, el $\mathcal{O}_X$-módulo $\widetilde{\mathcal{M}}$ se identifica con el $\mathcal{O}_{X'}$-módulo $\widetilde{\mathcal{N}}$ asociado a $\mathcal{N}$ (en el sentido de \sref{II.1.4.3-fr}).
  \item Sea $\mathcal{S}$ una $\mathcal{O}_Y$-álgebra graduada cualquiera, y sea $\mathcal{S}'$ la $\mathcal{O}_Y$-álgebra graduada tal que $\mathcal{S}'_0=\mathcal{O}_Y$ y $\mathcal{S}'_n=\mathcal{S}_n$ para todo $n>0$; entonces el isomorfismo canónico de $X=\operatorname{Proj}(\mathcal{S})$ sobre $X'=\operatorname{Proj}(\mathcal{S}')$ \sref{II.3.1.8-fr}[(ii)] identifica $\mathcal{O}_X(n)$ con $\mathcal{O}_{X'}(n)$ para todo $n\in\mathbf{Z}$.
    Esto resulta de la misma proposición para el caso afín \sref{II.2.5.16-fr} y del hecho de que estas identificaciones, para los abiertos afines de $Y$, conmutan con las operaciones de restricción.
    Del mismo modo, sea $X^{(d)}=\operatorname{Proj}(\mathcal{S}^{(d)})$; entonces el isomorfismo canónico de $X$ sobre $X^{(d)}$ \sref{II.3.1.8-fr}[(i)] identifica $\mathcal{O}_X(nd)$ con $\mathcal{O}_{X^{(d)}}(n)$ para todo $n\in\mathbf{Z}$.
\end{enumerate}
\end{remarks}

\begin{proposition}[3.2.10]
\label{II.3.2.10-fr}
Sea $\mathcal{L}$ un $\mathcal{O}_Y$-módulo invertible, y sea $g$ el isomorfismo canónico $X_{(\mathcal{L})}=\operatorname{Proj}(\mathcal{S}_{(\mathcal{L})})
\to X=\operatorname{Proj}(\mathcal{S})$ \sref{II.3.1.8-fr}[(iii)].
Entonces, para todo $n\in\mathbf{Z}$, $g_*(\mathcal{O}_{X_{(\mathcal{L})}}(n))$ es canónicamente isomorfo a $\mathcal{O}_X(n)\otimes_Y\mathcal{L}^{\otimes n}$.
\end{proposition}

\begin{proof}
Supongamos primero que $Y$ es afín, de anillo $A$, y que $\mathcal{L}=\widetilde{L}$, donde $L$ es un $A$-módulo libre de rango uno.
Con las notaciones de la demostración de \sref{II.3.1.8-fr}[(iii)], se define, para $f\in S_d$, un isomorfismo de $(S(n))_{(f)}\otimes_A L^{\otimes n}$ sobre $(S_{(L)}(n))_{(f\otimes c^d)}$ haciendo corresponder a $(x/f^k)\otimes c^n$, donde $x\in S_{kd+n}$, el elemento $(x\otimes c^{n+kd})/(f\otimes c^d)^k$; es inmediato que este isomorfismo es independiente del generador $c$ elegido de $L$; además, los isomorfismos así definidos para cada $f\in S_+$ son compatibles con los operadores de restricción $D_+(f)\to D_+(fg)$.
Finalmente, en el caso general, se ve sin dificultad, remontándose a las definiciones \sref{II.3.1.1-fr}, que los isomorfismos así definidos para cada abierto afín $U$ de $Y$ son compatibles con el paso de $U$ a un abierto afín $U'\subset U$.
\end{proof}
% Fuente francesa: PMIHES_1961__8__5_0, printed pp. 56-57.

\subsection{Módulo graduado $\mathcal{S}$ asociado a un haz sobre $\operatorname{Proj}(\mathcal{S})$}
\label{subsection:II.3.3-fr}

\emph{En todo este apartado se supone que la $\mathcal{O}_Y$-álgebra graduada $\mathcal{S}$ está engendrada por $\mathcal{S}_1$} \sref{II.3.1.9-fr}.
Recordemos que, en virtud de \sref{II.3.1.8-fr}[(i)], esta restricción no es esencial cuando se imponen las condiciones de finitud de \sref{II.3.1.10-fr}.

\begin{env}[3.3.1]
\label{II.3.3.1-fr}
Sea $p$ el morfismo estructural $X=\operatorname{Proj}(\mathcal{S})\to Y$.
Para todo $\mathcal{O}_X$-módulo $\mathcal{F}$, pongamos
\[
\label{II.3.3.1.1-fr}
  \bbGamma_*(\mathcal{F})
  =\bigoplus_{n\in\mathbf{Z}}p_*(\mathcal{F}(n))
\tag{3.3.1.1}
\]
\oldpage[II]{57}
y, en particular,
\[
\label{II.3.3.1.2-fr}
  \bbGamma_*(\mathcal{O}_X)
  =\bigoplus_{n\in\mathbf{Z}}p_*(\mathcal{O}_X(n)).
\tag{3.3.1.2}
\]

Se sabe \sref[0]{0.4.2.2-fr} que existe un homomorfismo canónico
\[
  p_*(\mathcal{F})\otimes_{\mathcal{O}_Y}p_*(\mathcal{G})
  \longrightarrow
  p_*(\mathcal{F}\otimes_{\mathcal{O}_X}\mathcal{G})
\]
para dos $\mathcal{O}_X$-módulos $\mathcal{F}$ y $\mathcal{G}$; se deduce, pues, de \sref{II.3.2.7.1-fr} que $\bbGamma_*(\mathcal{O}_X)$ está provisto de una estructura de $\mathcal{O}_Y$-álgebra graduada, y \sref{II.3.2.5.2-fr} define asimismo sobre $\bbGamma_*(\mathcal{F})$ una estructura de módulo graduado sobre $\bbGamma_*(\mathcal{O}_X)$.

En virtud de \sref{II.3.2.5-fr} y de la exactitud a izquierda del funtor $p_*$ \sref[0]{0.4.2.1-fr}, $\bbGamma_*(\mathcal{F})$ es un funtor covariante aditivo y \emph{exacto por la izquierda} en $\mathcal{F}$ en la categoría de los $\mathcal{O}_X$-módulos en la categoría de los $\mathcal{O}_Y$-módulos graduados (donde los morfismos son los homomorfismos de grado~$0$).
En particular, si $\mathcal{J}$ es un haz de ideales de $\mathcal{O}_X$, entonces $\bbGamma_*(\mathcal{J})$ se identifica con un haz graduado de ideales de $\bbGamma_*(\mathcal{O}_X)$.
\end{env}
% Fuente francesa: PMIHES_1961__8__5_0, printed pp. 57-58.

\begin{env}[3.3.2]
\label{II.3.3.2-fr}
Sea $\mathcal{M}$ un $\mathcal{S}$-módulo graduado cuasicoherente.
Para todo abierto afín $U$ de $Y$, se ha definido en \sref{II.2.6.2-fr} un homomorfismo de grupos abelianos
\[
  \alpha_{0,U}:\Gamma(U,\mathcal{M}_0)
  \longrightarrow\Gamma(p^{-1}(U),\widetilde{\mathcal{M}}).
\]
Es inmediato que estos homomorfismos conmutan con las operaciones de restricción \sref{II.2.8.13.1-fr} y definen, por tanto (sin recurrir a la hipótesis de que $\mathcal{S}$ está engendrada por $\mathcal{S}_1$), un homomorfismo de haces de grupos abelianos
\[
\label{II.3.3.2.1-fr}
  \alpha_0:\mathcal{M}_0\longrightarrow p_*(\widetilde{\mathcal{M}}).
\tag{3.3.2.1}
\]

Aplicando este resultado a cada uno de los $\mathcal{M}_n=(\mathcal{M}(n))_0$, y teniendo en cuenta \sref{II.3.2.8.1-fr}, se define un homomorfismo de haces de grupos abelianos
\[
\label{II.3.3.2.2-fr}
  \alpha_n:\mathcal{M}_n\longrightarrow
  p_*(\widetilde{\mathcal{M}}(n))
  \qquad\text{para todo }n\in\mathbf{Z},
\tag{3.3.2.2}
\]
de donde se obtiene un homomorfismo funtorial (de grado~$0$) de haces graduados de grupos abelianos
\[
\label{II.3.3.2.3-fr}
  \alpha:\mathcal{M}\longrightarrow\bbGamma_*(\widetilde{\mathcal{M}})
\tag{3.3.2.3}
\]
(denotado también por $\alpha_{\mathcal{M}}$).

Tomando en particular $\mathcal{M}=\mathcal{S}$, se comprueba que $\alpha:\mathcal{S}\to\bbGamma_*(\mathcal{O}_X)$ es un homomorfismo de $\mathcal{O}_Y$-álgebras graduadas, y que \sref{II.3.3.2.3-fr} es un dihomomorfismo de módulos graduados relativo a este homomorfismo de álgebras graduadas.

Observemos además que a cada $\alpha_n$ le corresponde \sref[0]{0.4.4.3-fr} un homomorfismo canónico de $\mathcal{O}_X$-módulos
\[
\label{II.3.3.2.4-fr}
  \alpha_n^\sharp:p^*(\mathcal{M}_n)
  \longrightarrow\widetilde{\mathcal{M}}(n).
\tag{3.3.2.4}
\]

Se comprueba sin dificultad que este homomorfismo no es otro que el que corresponde funtorialmente \sref{II.3.2.4-fr} al homomorfismo canónico (de grado~$0$) de $\mathcal{O}_Y$-módulos \emph{graduados}
\[
\label{II.3.3.2.5-fr}
  \mathcal{M}_n\otimes_{\mathcal{O}_Y}\mathcal{S}
  \longrightarrow\mathcal{M}(n)
\tag{3.3.2.5}
\]
\oldpage[II]{58}
donde la graduación del segundo miembro procede naturalmente de la de $\mathcal{M}$.
En efecto, se puede limitar al caso en que $Y=\operatorname{Spec}(A)$ es afín, $\mathcal{M}=\widetilde{M}$ y $\mathcal{S}=\widetilde{S}$, estando la $A$-álgebra graduada $S$ engendrada por $S_1$, de modo que, cuando $f$ recorre $S_1$, los $D_+(f)$ forman un recubrimiento de $X$.
Remontándose a las definiciones \sref{II.2.6.2-fr}, se ve entonces, teniendo en cuenta \sref[I]{I.1.6.7-fr}, que la restricción a $D_+(f)$ del homomorfismo \sref{II.3.3.2.4-fr} corresponde \sref[I]{I.1.3.8-fr} al homomorfismo de $S_{(f)}$-módulos
$M_n\otimes_A S_{(f)}\to(S(n))_{(f)}$
que hace corresponder a $x\otimes1$, donde $x\in M_n$, el elemento $x/1$; esto demuestra nuestra aserción.
\end{env}
% Fuente francesa: PMIHES_1961__8__5_0, printed p. 58.

\begin{proposition}[3.3.3]
\label{II.3.3.3-fr}
Para toda sección $f\in\Gamma(Y,\mathcal{S}_d)$, donde $d>0$, $X_f$ es idéntico al conjunto de los puntos de $X$ donde $\alpha_d(f)$, considerada como sección de $\mathcal{O}_X(d)$, no se anula \sref[0]{0.5.5.2-fr}.
\end{proposition}

\begin{proof}
( $\alpha_d(f)$ es una sección de $p_*(\mathcal{O}_X(d))$ sobre $Y$, pero, por definición, tal sección es también una sección de $\mathcal{O}_X(d)$ sobre $X$ \sref[0]{0.4.2.1-fr}.)
La definición de $X_f$ \sref{II.3.1.4-fr} reduce al caso en que $Y$ es afín, que ya fue tratado en \sref{II.2.6.3-fr}.
\end{proof}

\begin{env}[3.3.4]
\label{II.3.3.4-fr}
Supondremos en lo sucesivo, además de la hipótesis del comienzo de este apartado, que, para todo $\mathcal{O}_X$-módulo cuasicoherente $\mathcal{F}$, los $p_*(\mathcal{F}(n))$ son cuasicoherentes sobre $Y$, y que, por consiguiente, $\bbGamma_*(\mathcal{F})=\bigoplus_{n\in\mathbf{Z}}p_*(\mathcal{F}(n))$ es también un $\mathcal{O}_Y$-módulo cuasicoherente (\sref[I]{I.1.4.1-fr} y \sref[I]{I.1.3.9-fr}); esta circunstancia se producirá siempre, en particular, si $X$ es de tipo finito sobre $Y$ \sref[I]{I.9.2.2-fr}.
Se concluye que $(\bbGamma_*(\mathcal{F}))^{\sim}$ está definido y es un $\mathcal{O}_X$-módulo cuasicoherente.
Para todo abierto afín $U$ de $Y$, se tiene (\sref[I]{I.1.3.9-fr} y \sref{II.2.5.4-fr})
\begin{align*}
  \left(\Gamma\left(U,\bigoplus_{n\in\mathbf{Z}}
  p_*(\mathcal{F}(n))\right)\right)^{\sim}
  &=\bigoplus_{n\in\mathbf{Z}}
  \bigl(\Gamma(U,p_*(\mathcal{F}(n)))\bigr)^{\sim}\\
  &=\bigoplus_{n\in\mathbf{Z}}
  \bigl(\Gamma(p^{-1}(U),\mathcal{F}(n))\bigr)^{\sim}\\
  &=\left(\bigoplus_{n\in\mathbf{Z}}
  \Gamma(p^{-1}(U),\mathcal{F}(n))\right)^{\sim}\\
  &=\bigl(\Gamma_*(\mathcal{F}|p^{-1}(U))\bigr)^{\sim}
\end{align*}
y, por tanto, en virtud de \sref{II.2.6.4-fr}, se tiene un homomorfismo canónico
\[
  \beta_U:\left(\Gamma\left(U,\bigoplus_{n\in\mathbf{Z}}
  p_*(\mathcal{F}(n))\right)\right)^{\sim}
  \longrightarrow\mathcal{F}|p^{-1}(U).
\]

Además, la conmutatividad de \sref{II.2.8.13.2-fr} muestra que estos homomorfismos conmutan con las operaciones de restricción sobre $Y$; se obtiene, pues, un homomorfismo canónico funtorial
\[
\label{II.3.3.4.1-fr}
  \beta:(\bbGamma_*(\mathcal{F}))^{\sim}\longrightarrow\mathcal{F}
\tag{3.3.4.1}
\]
(denotado también por $\beta_{\mathcal{F}}$) para los $\mathcal{O}_X$-módulos cuasicoherentes.
\end{env}

\begin{proposition}[3.3.5]
\label{II.3.3.5-fr}
Sean $\mathcal{M}$ un $\mathcal{S}$-módulo graduado cuasicoherente y $\mathcal{F}$ un $\mathcal{O}_X$-módulo cuasicoherente; los homomorfismos compuestos
\[
\label{II.3.3.5.1-fr}
  \widetilde{\mathcal{M}}
  \xrightarrow{\widetilde{\alpha}}
  (\bbGamma_*(\widetilde{\mathcal{M}}))^{\sim}
  \xrightarrow{\beta}\widetilde{\mathcal{M}}
\tag{3.3.5.1}
\]
\[
\label{II.3.3.5.2-fr}
  \bbGamma_*(\mathcal{F})\xrightarrow{\alpha}
  \bbGamma_*((\bbGamma_*(\mathcal{F}))^{\sim})
  \xrightarrow{\bbGamma_*(\beta)}\bbGamma_*(\mathcal{F})
\tag{3.3.5.2}
\]
son los isomorfismos identidad.
\end{proposition}

\begin{proof}
La cuestión es local sobre $Y$, y se reduce a \sref{II.2.6.5-fr}.
\end{proof}
% Fuente francesa: PMIHES_1961__8__5_0, printed p. 59.

\oldpage[II]{59}
\subsection{Condiciones de finitud}
\label{subsection:II.3.4-fr}


\begin{proposition}[3.4.1]
\label{II.3.4.1-fr}
Sean $Y$ un preesquema y $\mathcal{S}$ una $\mathcal{O}_Y$-álgebra cuasicoherente engendrada por $\mathcal{S}_1$ \sref{II.3.1.9-fr}; se supone además que $\mathcal{S}_1$ es de tipo finito.
Entonces $X=\operatorname{Proj}(\mathcal{S})$ es de tipo finito sobre $Y$.
\end{proposition}

\begin{proof}
Como la cuestión es local sobre $Y$, se puede suponer que $Y$ es afín de anillo $A$; se tiene entonces $\mathcal{S}=\widetilde{S}$, donde $S=\Gamma(Y,\mathcal{S})$, y, por hipótesis, $S$ es una $A$-álgebra engendrada por $S_1=\Gamma(Y,\mathcal{S}_1)$, pudiéndose suponer además que $S_1$ es un $A$-módulo de tipo finito (\sref[I]{I.1.3.9-fr} y \sref[I]{I.1.3.12-fr}).
Por consiguiente, $S$ es una $A$-álgebra graduada de tipo finito, y se reduce a \sref{II.2.7.1-fr}[(ii)].
\end{proof}

\begin{env}[3.4.2]
\label{II.3.4.2-fr}
Sea $\mathcal{S}$ una $\mathcal{O}_Y$-álgebra graduada cuasicoherente; para un $\mathcal{S}$-módulo graduado cuasicoherente $\mathcal{M}$, consideraremos las condiciones de finitud siguientes:
\begin{enumerate}
  \item[(\textbf{TF})] Existe un entero $n$ tal que el $\mathcal{S}$-módulo $\bigoplus_{k\geq n}\mathcal{M}_k$ es de tipo finito.
  \item[(\textbf{TN})] Existe un entero $n$ tal que $\mathcal{M}_k=0$ para $k\geq n$.
\end{enumerate}

Si $\mathcal{M}$ satisface (\textbf{TN}), se tiene $\widetilde{\mathcal{M}}=0$, pues la cuestión es local sobre $Y$ \sref{II.2.7.2-fr}.

Sean $\mathcal{M}$ y $\mathcal{N}$ dos $\mathcal{S}$-módulos graduados cuasicoherentes; se dice que un homomorfismo $u:\mathcal{M}\to\mathcal{N}$ de grado~$0$ es \emph{(\textbf{TN})-inyectivo} (resp. \emph{(\textbf{TN})-sobreyectivo}, \emph{(\textbf{TN})-biyectivo}) si existe un entero $n$ tal que $u_k:\mathcal{M}_k\to\mathcal{N}_k$ es inyectivo (resp. sobreyectivo, biyectivo) para $k\geq n$; entonces $\widetilde{u}:\widetilde{\mathcal{M}}\to\widetilde{\mathcal{N}}$ es inyectivo (resp. sobreyectivo, biyectivo) en virtud de \sref{II.2.7.2-fr}, por ser la cuestión local sobre $Y$ y teniendo en cuenta \sref[I]{I.1.3.9-fr}; cuando $u$ es (\textbf{TN})-biyectivo, se dice también que $u$ es un (\textbf{TN})-isomorfismo.
\end{env}

\begin{proposition}[3.4.3]
\label{II.3.4.3-fr}
Sean $Y$ un preesquema y $\mathcal{S}$ una $\mathcal{O}_Y$-álgebra graduada cuasicoherente engendrada por $\mathcal{S}_1$, suponiéndose que $\mathcal{S}_1$ es de tipo finito.
Sea $\mathcal{M}$ un $\mathcal{S}$-módulo graduado cuasicoherente.
\begin{enumerate}
  \item[(i)] Si $\mathcal{M}$ satisface la condición \textup{(\textbf{TF})}, $\widetilde{\mathcal{M}}$ es de tipo finito.
  \item[(ii)] Supongamos que $\mathcal{M}$ satisface \textup{(\textbf{TF})}; para que $\widetilde{\mathcal{M}}=0$, es necesario y suficiente que $\mathcal{M}$ satisfaga \textup{(\textbf{TN})}.
\end{enumerate}
\end{proposition}

\begin{proof}
Como las cuestiones son locales sobre $Y$, se reduce al caso en que $Y$ es afín de anillo $A$, $\mathcal{S}=\widetilde{S}$, donde $S$ es una $A$-álgebra graduada tal que el ideal $S_+$ es de tipo finito, y $\mathcal{M}=\widetilde{M}$, donde $M$ es un $S$-módulo graduado; la proposición resulta entonces de \sref{II.2.7.3-fr}.
\end{proof}
% Fuente francesa: PMIHES_1961__8__5_0, printed pp. 59-60.

\begin{theorem}[3.4.4]
\label{II.3.4.4-fr}
Sean $Y$ un preesquema y $\mathcal{S}$ una $\mathcal{O}_Y$-álgebra graduada cuasicoherente engendrada por $\mathcal{S}_1$, suponiéndose que $\mathcal{S}_1$ es de tipo finito; sea $X=\operatorname{Proj}(\mathcal{S})$.
Para todo $\mathcal{O}_X$-módulo cuasicoherente $\mathcal{F}$, el homomorfismo canónico $\beta$ \sref{II.3.3.4-fr} es un isomorfismo.
\end{theorem}

\begin{proof}
Observemos ante todo que $\beta$ está definido en virtud de \sref{II.3.4.1-fr}.
Para ver que $\beta$ es un isomorfismo, se reduce al caso en que $Y$ es afín de anillo $A$, $\mathcal{S}=\widetilde{S}$, donde $S$ es una $A$-álgebra graduada engendrada por $S_1$, y $S_1$ es un $A$-módulo de tipo finito.
Basta entonces aplicar \sref{II.2.7.5-fr}.
\end{proof}

\begin{corollary}[3.4.5]
\label{II.3.4.5-fr}
Bajo las hipótesis de \sref{II.3.4.4-fr}, todo $\mathcal{O}_X$-módulo cuasicoherente $\mathcal{F}$ es isomorfo a un $\mathcal{O}_X$-módulo de la forma $\widetilde{\mathcal{M}}$, donde $\mathcal{M}$ es un $\mathcal{S}$-módulo graduado cuasicoherente.
Si, además, $\mathcal{F}$ es de tipo finito, y se supone que $Y$ es un esquema cuasicompacto o que el espacio subyacente a $Y$ es noetheriano, puede suponerse que $\mathcal{M}$ es de tipo finito.
\end{corollary}

\oldpage[II]{60}
\begin{proof}
La primera afirmación resulta inmediatamente de \sref{II.3.4.4-fr}, tomando $\mathcal{M}=\mathbf{\Gamma}_*(\mathcal{F})$.
Para establecer la segunda, basta demostrar que $\mathcal{M}$ es el límite inductivo de sus sub-$\mathcal{S}$-módulos \emph{graduados} de tipo finito $\mathcal{N}_\lambda$: en efecto, resultará de ello que $\widetilde{\mathcal{M}}$ es el límite inductivo de los $\widetilde{\mathcal{N}}_\lambda$ \sref{II.3.2.4-fr}, y, por tanto, $\mathcal{F}$ es el límite inductivo de los $\beta(\widetilde{\mathcal{N}}_\lambda)$; como $X$ es cuasicompacto (\sref{II.3.4.1-fr} y \sref[I]{I.6.3.1-fr}) y $\mathcal{F}$ es de tipo finito, $\mathcal{F}$ será necesariamente igual a uno de los $\beta(\widetilde{\mathcal{N}}_\lambda)$ \sref[0]{0.5.2.3-fr}.

Para definir los $\mathcal{N}_\lambda$ que tienen a $\mathcal{M}$ por límite inductivo, basta considerar, para cada $n\in\mathbf{Z}$, el $\mathcal{O}_Y$-módulo cuasicoherente $\mathcal{M}_n$, que es el límite inductivo de sus sub-$\mathcal{O}_Y$-módulos $\mathcal{M}_n^{(\mu_n)}$ de tipo finito, en virtud de las hipótesis sobre $Y$ \sref[I]{I.9.4.9-fr}; es inmediato que $\mathcal{P}_{\mu_n}=\mathcal{S}\cdot\mathcal{M}_n^{(\mu_n)}$ es un $\mathcal{S}$-módulo graduado de tipo finito, y se comprueba inmediatamente que, tomando por $\mathcal{N}_\lambda$ las sumas finitas de $\mathcal{S}$-módulos de la forma $\mathcal{P}_{\mu_n}$, se obtienen los objetos buscados.
\end{proof}
% Fuente francesa: PMIHES_1961__8__5_0, printed p. 60.

\begin{corollary}[3.4.6]
\label{II.3.4.6-fr}
Supongamos satisfechas las hipótesis de \sref{II.3.4.4-fr}, y supongamos además que el espacio subyacente a $Y$ es cuasicompacto; sea $\mathcal{F}$ un $\mathcal{O}_X$-módulo cuasicoherente de tipo finito.
Existe un entero $n_0$ tal que, para todo $n\geq n_0$, el homomorfismo canónico
$\sigma:p^*(p_*(\mathcal{F}(n)))\to\mathcal{F}(n)$
\sref[0]{0.4.4.3-fr} es sobreyectivo.
\end{corollary}

\begin{proof}
Para todo $y\in Y$, sea $U$ una vecindad abierta afín de $y$ en $Y$.
Existe un entero $n_0(U)$ tal que, para todo $n\geq n_0(U)$, $\mathcal{F}(n)|p^{-1}(U)$ está engendrado por un número finito de sus secciones sobre $p^{-1}(U)$ \sref{II.2.7.9-fr}; pero estas secciones son las imágenes canónicas de secciones de $p^*(p_*(\mathcal{F}(n)))$ sobre $p^{-1}(U)$ (\sref[0]{0.3.7.1-fr} y \sref[0]{0.4.4.3-fr}), por lo que $\mathcal{F}(n)|p^{-1}(U)$ es igual a la imagen canónica de $p^*(p_*(\mathcal{F}(n)))|p^{-1}(U)$.
Finalmente, como $Y$ es cuasicompacto, existe un recubrimiento finito de $Y$ por abiertos afines $U_i$, y tomando por $n_0$ el mayor de los $n_0(U_i)$ se termina la demostración.
\end{proof}

\par\medskip\noindent\textit{Observaciones (3.4.7). ---\ }
\label{II.3.4.7-fr}
Si $p=(\psi,\theta):X\to Y$ es un morfismo de espacios anillados y $\mathcal{F}$ un $\mathcal{O}_X$-módulo, el hecho de que el homomorfismo canónico $\sigma:p^*(p_*(\mathcal{F}))\to\mathcal{F}$ sea sobreyectivo se explicita de la manera siguiente \sref[0]{0.4.4.1-fr}: para todo $x\in X$ y toda sección $s$ de $\mathcal{F}$ sobre una vecindad abierta $V$ de $x$, existen una vecindad abierta $U$ de $p(x)$ en $Y$, un número finito de secciones $t_i$ ($1\leq i\leq m$) de $\mathcal{F}$ sobre $p^{-1}(U)$, una vecindad $W\subset V\cap p^{-1}(U)$ de $x$ y secciones $a_i$ ($1\leq i\leq m$) de $\mathcal{O}_X$ sobre $W$, tales que
\[
  s|W=\sum_i a_i\cdot(t_i|W).
\]
Si $Y$ es un \emph{esquema afín} y $p_*(\mathcal{F})$ es \emph{cuasicoherente}, esta condición equivale a que $\mathcal{F}$ esté \emph{engendrado por sus secciones sobre $X$} \sref[0]{0.5.5.1-fr}: en efecto, si $Y=\operatorname{Spec}(A)$, puede suponerse que $U=D(f)$, con $f\in A$; existen entonces un entero $n>0$ y secciones $s_i$ de $\mathcal{F}$ sobre $X$ tales que $t_i$ es la restricción a $p^{-1}(U)$ de $s_ig^n$, donde $g=\theta(f)$ (aplicando \sref[I]{I.1.4.1-fr} a $p_*(\mathcal{F})$); como $g$ es invertible sobre $p^{-1}(U)$, se tiene, pues,
\[
  s|W=\sum_i b_i\cdot(s_i|W)
\]
donde $b_i=a_i(g|W)^{-n}$, de donde resulta nuestra afirmación.
Cuando $Y$ es afín, el corolario \sref{II.3.4.6-fr} vuelve a dar \sref{II.2.7.9-fr}.
% Fuente francesa: PMIHES_1961__8__5_0, printed p. 61.


Se concluye que, cuando $Y$ es un preesquema cualquiera, las tres condiciones 
\oldpage[II]{61}
siguientes, para un $\mathcal{O}_X$-módulo cuasicoherente $\mathcal{F}$ tal que $p_*(\mathcal{F})$ sea \emph{cuasicoherente}, son equivalentes:
\begin{enumerate}
  \item[a)] \emph{El homomorfismo canónico $\sigma:p^*(p_*(\mathcal{F}))\to\mathcal{F}$ es sobreyectivo.}
  \item[b)] \emph{Existe un $\mathcal{O}_Y$-módulo cuasicoherente $\mathcal{G}$ y un homomorfismo sobreyectivo $p^*(\mathcal{G})\to\mathcal{F}$.}
  \item[c)] \emph{Para todo abierto afín $U$ de $Y$, $\mathcal{F}|p^{-1}(U)$ está engendrado por sus secciones sobre $p^{-1}(U)$.}
\end{enumerate}

En efecto, acaba de demostrarse la equivalencia entre \emph{a)} y \emph{c)}.
Por otra parte, es claro que \emph{a)} implica \emph{b)}, pues $p_*(\mathcal{F})$ es cuasicoherente por hipótesis.
Recíprocamente, todo homomorfismo $u:p^*(\mathcal{G})\to\mathcal{F}$ se factoriza en
$p^*(\mathcal{G})\to p^*(p_*(\mathcal{F}))\xrightarrow{\sigma}\mathcal{F}$
\sref[0]{0.3.5.4.4-fr}; por tanto, si $u$ es sobreyectivo, también lo es $\sigma$, lo que demuestra que \emph{b)} implica \emph{a)}.


\begin{corollary}[3.4.8]
\label{II.3.4.8-fr}
Supongamos satisfechas las hipótesis de \sref{II.3.4.4-fr}, y supongamos además que $Y$ es un esquema cuasicompacto o que el espacio subyacente a $Y$ es noetheriano.
Sea $\mathcal{F}$ un $\mathcal{O}_X$-módulo cuasicoherente de tipo finito; existe entonces un entero $n_0$ tal que, para todo $n\geq n_0$ $\mathcal{F}$ es isomorfo a un cociente de un $\mathcal{O}_X$-módulo de la forma $(p^*(\mathcal{G}))(-n)$, donde $\mathcal{G}$ es un $\mathcal{O}_Y$-módulo cuasicoherente de tipo finito (dependiente de $n$).
\end{corollary}

\begin{proof}
Como el morfismo estructural $X\to Y$ es separado y de tipo finito, $p_*(\mathcal{F}(n))$ es cuasicoherente \sref[I]{I.9.2.2-fr}[\emph{b)}] y, por tanto, límite inductivo de sus sub-$\mathcal{O}_Y$-módulos cuasicoherentes de tipo finito, en virtud de las hipótesis sobre $Y$ \sref[I]{I.9.4.9-fr}.
Se deduce entonces de \sref{II.3.4.6-fr}, \sref[0]{0.4.3.2-fr} y \sref[0]{0.5.2.3-fr} que $\mathcal{F}(n)$ es la imagen canónica de un $\mathcal{O}_X$-módulo de la forma $p^*(\mathcal{G})$, donde $\mathcal{G}$ es un sub-$\mathcal{O}_Y$-módulo cuasicoherente de tipo finito de $p_*(\mathcal{F}(n))$; el corolario resulta entonces de \sref{II.3.2.5.2-fr} y \sref{II.3.2.7.1-fr}.
\end{proof}
% Fuente francesa: PMIHES_1961__8__5_0, printed pp. 61-62.

\subsection{Comportamiento funtorial}
\label{subsection:II.3.5-fr}

\begin{env}[3.5.1]
\label{II.3.5.1-fr}
Sean $Y$ un preesquema, y $\mathcal{S}$ y $\mathcal{S}'$ dos $\mathcal{O}_Y$-álgebras graduadas cuasicoherentes en grados positivos; pongamos $X=\operatorname{Proj}(\mathcal{S})$ y $X'=\operatorname{Proj}(\mathcal{S}')$, y sean $p$ y $p'$ los morfismos estructurales de $X$ y $X'$ en $Y$.
Sea $\varphi:\mathcal{S}'\to\mathcal{S}$ un $\mathcal{O}_Y$-homomorfismo de álgebras graduadas.
Para todo abierto afín $U$ de $Y$, pongamos $S_U=\Gamma(U,\mathcal{S})$ y $S'_U=\Gamma(U,\mathcal{S}')$; el homomorfismo $\varphi$ define un homomorfismo $\varphi_U:S'_U\to S_U$ de $A_U$-álgebras graduadas, donde $A_U=\Gamma(U,\mathcal{O}_Y)$.
Le corresponden en $p^{-1}(U)$ un conjunto abierto $G(\varphi_U)$ y un morfismo $\Phi_U:G(\varphi_U)\to p'^{-1}(U)$ \sref{II.2.8.1-fr}.
Además, si $V\subset U$ es un abierto afín, el diagrama
\[
\label{II.3.5.1.1-fr}
  \xymatrix{
    S'_U\ar[r]^{\varphi_U}\ar[d]
    & S_U\ar[d]\\
    S'_V\ar[r]_{\varphi_V}
    & S_V
  }
\tag{3.5.1.1}
\]
es conmutativo, y se ve inmediatamente, por definición \sref{II.2.8.1-fr}, que
\[
  G(\varphi_V)=G(\varphi_U)\cap p^{-1}(V)
\]
y que $\Phi_V$ es la restricción de $\Phi_U$ a $G(\varphi_V)$.
Así, se ha definido un subconjunto abierto $G(\varphi)$ de $X$ tal que $G(\varphi)\cap p^{-1}(U)=G(\varphi_U)$ para todo abierto afín $U\subset Y$, y un morfismo $Y$-
\oldpage[II]{62}
\emph{afín} $\Phi:G(\varphi)\to X'$, que se dice \emph{asociado} a $\varphi$ y que denotaremos por $\operatorname{Proj}(\varphi)$.

Si, para todo $y\in Y$, existe una vecindad afín $U$ de $y$ tal que el $\Gamma(U,\mathcal{O}_Y)$-módulo $\Gamma(U,\mathcal{S}_+)$ está engendrado por $\varphi(\Gamma(U,\mathcal{S}'_+))$, entonces $G(\varphi_U)=p^{-1}(U)$ y, por consiguiente, $G(\varphi)=X$.
\end{env}

\begin{proposition}[3.5.2]
\label{II.3.5.2-fr}
\medskip\noindent
\begin{enumerate}
  \item[(i)] Si $\mathcal{M}$ es un $\mathcal{S}$-módulo graduado cuasicoherente, existe un isomorfismo canónico funtorial del $\mathcal{O}_{X'}$-módulo $(\mathcal{M}_{[\varphi]})^{\sim}$ sobre el $\mathcal{O}_{X'}$-módulo $\Phi_*(\widetilde{\mathcal{M}}|G(\varphi))$.
  \item[(ii)] Si $\mathcal{M}'$ es un $\mathcal{S}'$-módulo graduado cuasicoherente, existe un homomorfismo canónico funtorial $\nu$ del $(\mathcal{O}_X|G(\varphi))$-módulo $\Phi^*(\widetilde{\mathcal{M}'})$ en el $(\mathcal{O}_X|G(\varphi))$-módulo $(\mathcal{M}'\otimes_{\mathcal{S}'}\mathcal{S})^{\sim}|G(\varphi)$.
    Si $\mathcal{S}'$ está engendrada por $\mathcal{S}'_1$, $\nu$ es un isomorfismo.
\end{enumerate}
\end{proposition}

\begin{proof}
Los homomorfismos en cuestión están efectivamente definidos ya cuando $Y$ es afín (\sref{II.2.8.7-fr} y \sref{II.2.8.8-fr}), y, en el caso general, basta comprobar que son compatibles con la restricción de un abierto afín de $Y$ a un abierto más pequeño, lo que resulta inmediatamente de la conmutatividad de \sref{II.3.5.1.1-fr}.
\end{proof}

En particular, para todo $n\in\mathbf{Z}$, se tiene un homomorfismo canónico
\[
\label{II.3.5.2.1-fr}
  \Phi^*(\mathcal{O}_{X'}(n))\longrightarrow
  \mathcal{O}_X(n)|G(\varphi).
\tag{3.5.2.1}
\]
% Fuente francesa: PMIHES_1961__8__5_0, printed p. 62.

\begin{proposition}[3.5.3]
\label{II.3.5.3-fr}
Sean $Y$ y $Y'$ preesquemas, $\psi:Y'\to Y$ un morfismo y $\mathcal{S}$ una $\mathcal{O}_Y$-álgebra graduada cuasicoherente; pongamos $\mathcal{S}'=\psi^*(\mathcal{S})$.
Entonces el $Y'$-esquema $X'=\operatorname{Proj}(\mathcal{S}')$ se identifica canónicamente con $\operatorname{Proj}(\mathcal{S})\times_Y Y'$.
Además, si $\mathcal{M}$ es un $\mathcal{S}$-módulo graduado cuasicoherente, el $\mathcal{O}_{X'}$-módulo $(\psi^*(\mathcal{M}))^{\sim}$ se puede identificar con $\widetilde{\mathcal{M}}\otimes_Y\mathcal{O}_{Y'}$.
\end{proposition}

\begin{proof}
Observemos ante todo que $\psi^*(\mathcal{S})$ y $\psi^*(\mathcal{M})$ son $\mathcal{O}_{Y'}$-módulos cuasicoherentes, al igual que sus componentes homogéneos \sref[0]{0.5.1.4-fr}.
Sea $U$ un abierto afín de $Y$, $U'\subset\psi^{-1}(U)$ un abierto afín de $Y'$, y sean $A$ y $A'$ los anillos de $U$ y $U'$, respectivamente; entonces $\mathcal{S}|U=\widetilde{S}$, donde $S$ es una $A$-álgebra graduada, y $\mathcal{S}'|U'$ se puede identificar con $(S\otimes_A A')^{\sim}$ \sref[I]{I.1.6.5-fr}; la primera afirmación resulta entonces de \sref{II.2.8.10-fr} y \sref[I]{I.3.2.6.2-fr}, pues se comprueba inmediatamente que la proyección $\operatorname{Proj}(\mathcal{S}'|U')\to
\operatorname{Proj}(\mathcal{S}|U)$ definida por la identificación precedente es compatible con las operaciones de restricción sobre $U$ y $U'$, y define por tanto efectivamente un morfismo $\operatorname{Proj}(\mathcal{S}')\to\operatorname{Proj}(\mathcal{S})$.
Sean ahora
\[
  q:\operatorname{Proj}(\mathcal{S})\to Y,
  \qquad q':\operatorname{Proj}(\mathcal{S}')\to Y'
\]
los morfismos estructurales; $q'^{-1}(U')$ se puede identificar entonces con $q^{-1}(U)\times_U U'$, y los dos haces $(\psi^*(\mathcal{M}))^{\sim}|q'^{-1}(U')$ y $(\widetilde{\mathcal{M}}\otimes_Y\mathcal{O}_{Y'})|q'^{-1}(U')$ se identifican ambos canónicamente con $(M\otimes_A A')^{\sim}$, donde se ha puesto $M=\Gamma(U,\mathcal{M})$, en virtud de \sref{II.2.8.10-fr} y \sref[I]{I.1.6.5-fr}; de aquí resulta la segunda afirmación, pues se comprueba de nuevo inmediatamente la compatibilidad de las identificaciones precedentes con las operaciones de restricción.
\end{proof}

\begin{corollary}[3.5.4]
\label{II.3.5.4-fr}
Con las notaciones de \sref{II.3.5.3-fr}, $\mathcal{O}_{X'}(n)$ se identifica canónicamente con $\mathcal{O}_X(n)\otimes_Y\mathcal{O}_{Y'}$ para todo $n\in\mathbf{Z}$ (donde $X=\operatorname{Proj}(\mathcal{S})$).
\end{corollary}

\begin{proof}
En efecto, con las notaciones de \sref{II.3.5.3-fr}, es claro que $\psi^*(\mathcal{S}(n))=\mathcal{S}'(n)$ para todo $n\in\mathbf{Z}$.
\end{proof}
% Fuente francesa: PMIHES_1961__8__5_0, printed pp. 62-63.

\begin{env}[3.5.5]
\label{II.3.5.5-fr}
Manteniendo las notaciones anteriores, designemos por $\Psi$ la proyección canónica $X'\to X$, y pongamos $\mathcal{M}'=\psi^*(\mathcal{M})$; suponemos además que $\mathcal{S}$ está engendrada por $\mathcal{S}_1$ 
\oldpage[II]{63}
y que $X$ es de tipo finito sobre $Y$.

Se deduce entonces que $\mathcal{S}'$ está engendrada por $\mathcal{S}'_1$ (como se ve reduciéndose al caso en que $Y$ y $Y'$ son afines), y que $X'$ es de tipo finito sobre $Y'$ \sref[I]{I.6.3.4-fr}.
Sea $\mathcal{F}$ un $\mathcal{O}_X$-módulo y pongamos $\mathcal{F}'=\Psi^*(\mathcal{F})$; se deduce entonces de \sref{II.3.5.4-fr} y \sref[0]{0.4.3.3-fr} que $\mathcal{F}'(n)=\Psi^*(\mathcal{F}(n))$ para todo $n\in\mathbf{Z}$.
Se define además un $\psi$-homomorfismo canónico $\theta_n:q_*(\mathcal{F}(n))\to q'_*(\mathcal{F}'(n))$ de la manera siguiente: dada la conmutatividad del diagrama
\[
  \xymatrix{
    X\ar[d]_{q}
    & X'\ar[l]_{\Psi}\ar[d]^{q'}\\
    Y
    & Y'\ar[l]^{\psi}
  }
\]
basta definir un homomorfismo $q_*(\mathcal{F}(n))\to
\psi_*(q'_*(\Psi^*(\mathcal{F}(n))))
=q_*(\Psi_*(\Psi^*(\mathcal{F}(n))))$, y basta tomar el homomorfismo $\theta_n=q_*(\rho_n)$, donde $\rho_n$ es el homomorfismo canónico $\mathcal{F}(n)\to\Psi_*(\Psi^*(\mathcal{F}(n)))$ \sref[0]{0.4.4.3-fr}.
Es inmediato que, para todo abierto afín $U$ de $Y$ y todo abierto afín $U'$ de $Y'$ tal que $U'\subset\psi^{-1}(U)$, el homomorfismo $\theta_n$ da, sobre las secciones, el homomorfismo canónico \sref[0]{0.3.7.2-fr} $\Gamma(q^{-1}(U),\mathcal{F}(n))\to
\Gamma(q'^{-1}(U'),\mathcal{F}'(n))$.
La conmutatividad de \sref{II.2.8.13.2-fr} muestra entonces que, si $\mathcal{F}$ es cuasicoherente, el diagrama
\[
  \xymatrix{
    \mathcal{F}\ar[r]^{\rho}
    & \mathcal{F}'\\
    (\mathbf{\Gamma}_*(\mathcal{F}))^{\sim}
      \ar[u]^{\beta_{\mathcal{F}}}\ar[r]_{\widetilde{\theta}}
    & (\mathbf{\Gamma}_*(\mathcal{F}'))^{\sim}
      \ar[u]_{\beta_{\mathcal{F}'}}
  }
\]
es conmutativo (siendo la flecha horizontal superior el $\Psi$-morfismo canónico $\mathcal{F}\to\Psi^*(\mathcal{F})$).

Del mismo modo, la conmutatividad de \sref{II.2.8.13.1-fr} muestra que el diagrama
\[
  \xymatrix{
    \mathbf{\Gamma}_*(\widetilde{\mathcal{M}})\ar[r]^{\theta}
    & \mathbf{\Gamma}_*(\widetilde{\mathcal{M}'})\\
    \mathcal{M}\ar[r]_{\rho}\ar[u]^{\alpha_{\mathcal{M}}}
    & \mathcal{M}'\ar[u]_{\alpha_{\mathcal{M}'}}
  }
\]
es conmutativo (siendo la flecha horizontal inferior el $\psi$-morfismo canónico $\mathcal{M}\to\psi^*(\mathcal{M})$).
\end{env}
% Fuente francesa: PMIHES_1961__8__5_0, printed pp. 63-64.

\begin{env}[3.5.6]
\label{II.3.5.6-fr}
Consideremos ahora dos preesquemas $Y$, $Y'$, un morfismo $g:Y'\to Y$, una $\mathcal{O}_Y$-álgebra (resp. $\mathcal{O}_{Y'}$-álgebra) graduada cuasicoherente $\mathcal{S}$ (resp. $\mathcal{S}'$), y un $g$-morfismo de álgebras graduadas $u:\mathcal{S}\to\mathcal{S}'$, es decir, un $\mathcal{O}_Y$-homomorfismo de álgebras graduadas $\mathcal{S}\to g_*(\mathcal{S}')$; ya se sabe que esto equivale a dar un $\mathcal{O}_{Y'}$-homomorfismo de álgebras graduadas $u^\sharp:g^*(\mathcal{S})\to\mathcal{S}'$.
Así, se obtiene canónicamente de $u^\sharp$ un $Y'$-morfismo $W=\operatorname{Proj}(u^\sharp):G(u^\sharp)\to
\operatorname{Proj}(g^*(\mathcal{S}))$, donde $G(u^\sharp)$ es un subconjunto abierto de $X'=\operatorname{Proj}(\mathcal{S}')$ \sref{II.3.5.1-fr}.
Por otra parte, $X''=\operatorname{Proj}(g^*(\mathcal{S}))$ se identifica 
\oldpage[II]{64}
canónicamente con $X\times_Y Y'$, donde $X=\operatorname{Proj}(\mathcal{S})$ \sref{II.3.5.3-fr};

componiendo con $\operatorname{Proj}(u^\sharp)$ la primera proyección $p:X\times_Y Y'\to X$, se obtiene, pues, un morfismo $v:G(u^\sharp)\to X$, que denotaremos por $\operatorname{Proj}(u)$, y tal que el diagrama
\[
  \xymatrix{
    G(u^\sharp)\ar[r]^{v}\ar[d]
    & X\ar[d]\\
    Y'\ar[r]_{g}
    & Y
  }
\]
es conmutativo.

Además, para todo $\mathcal{S}$-módulo graduado cuasicoherente $\mathcal{M}$, se tiene un $v$-morfismo canónico
\[
\label{II.3.5.6.1-fr}
  \upsilon:\widetilde{\mathcal{M}}\longrightarrow
  (g^*(\mathcal{M})\otimes_{g^*(\mathcal{S})}\mathcal{S}')^{\sim}
  |G(u^\sharp).
\tag{3.5.6.1}
\]

En efecto, $\upsilon^\sharp$ se obtiene componiendo los homomorfismos
\[
  v^*(\widetilde{\mathcal{M}})
  =W^*(p^*(\widetilde{\mathcal{M}}))
  \longrightarrow W^*((g^*(\mathcal{M}))^{\sim})
  \longrightarrow
  (g^*(\mathcal{M})\otimes_{g^*(\mathcal{S})}\mathcal{S}')^{\sim}
  |G(u^\sharp)
\]
donde la primera flecha procede del isomorfismo \sref{II.3.5.3-fr} y la segunda es el homomorfismo \sref{II.3.5.2-fr}[(ii)]; si $\mathcal{S}$ está engendrada por $\mathcal{S}_1$, se deduce de \sref{II.3.5.2-fr} que $\upsilon^\sharp$ es un \emph{isomorfismo}.

Como caso particular de \sref{II.3.5.6.1-fr}, se tiene, para todo $n\in\mathbf{Z}$, un $v$-morfismo canónico
\[
\label{II.3.5.6.2-fr}
  \upsilon:\mathcal{O}_X(n)\longrightarrow
  \mathcal{O}_{X'}(n)|G(u^\sharp).
\tag{3.5.6.2}
\]
\end{env}
% Fuente francesa: PMIHES_1961__8__5_0, printed pp. 64-65.

\subsection{Subpreesquemas cerrados de un preesquema $\operatorname{Proj}(\mathcal{S})$}
\label{subsection:II.3.6-fr}

\begin{env}[3.6.1]
\label{II.3.6.1-fr}
Sea $Y$ un preesquema, y $\varphi:\mathcal{S}\to\mathcal{S}'$ un homomorfismo de $\mathcal{O}_Y$-álgebras graduadas cuasicoherentes, de grado $0$.
Se dice que $\varphi$ es \emph{\textbf{(TN)}-sobreyectivo} (resp. \emph{\textbf{(TN)}-inyectivo}, \emph{\textbf{(TN)}-biyectivo}) si existe un entero $n$ tal que, para todo $k\geq n$, $\varphi_k:\mathcal{S}_k\to\mathcal{S}'_k$ es sobreyectivo (resp. inyectivo, biyectivo).
Cuando así sucede, se reduce el estudio del morfismo correspondiente $\Phi:\operatorname{Proj}(\mathcal{S}')\to\operatorname{Proj}(\mathcal{S})$ al caso en que $\varphi$ es \emph{sobreyectivo} (resp. \emph{inyectivo}, \emph{biyectivo}).
Esto se demuestra como en \sref{II.2.9.1-fr} (que es el caso particular en que $Y$ es afín), utilizando \sref{II.3.1.8-fr}.
En lugar de decir que $\varphi$ es \textbf{(TN)}-biyectivo, se dirá también que es un \emph{\textbf{(TN)}-isomorfismo}.
\end{env}

\begin{proposition}[3.6.2]
\label{II.3.6.2-fr}
Sea $Y$ un preesquema y $\mathcal{S}$ una $\mathcal{O}_Y$-álgebra graduada cuasicoherente, y sea $X=\operatorname{Proj}(\mathcal{S})$.
\begin{enumerate}
  \item[(i)] Si $\varphi:\mathcal{S}\to\mathcal{S}'$ es un homomorfismo \textup{\textbf{(TN)}}-sobreyectivo de $\mathcal{O}_Y$-álgebras graduadas, entonces el morfismo correspondiente $\Phi=\operatorname{Proj}(\varphi)$ \sref{II.3.5.1-fr} está definido en todo $\operatorname{Proj}(\mathcal{S}')$ y es una inmersión cerrada de $\operatorname{Proj}(\mathcal{S}')$ en $X$.
    Si $\mathcal{J}$ es el núcleo de $\varphi$, el subpreesquema cerrado de $X$ asociado a $\Phi$ está definido por el haz cuasicoherente de ideales $\widetilde{\mathcal{J}}$ en $\mathcal{O}_X$.
  \item[(ii)] Supongamos además que $\mathcal{S}_0=\mathcal{O}_Y$, que $\mathcal{S}$ está engendrada por $\mathcal{S}_1$, y que $\mathcal{S}_1$ es de tipo finito.
    Sea $X'$ un subpreesquema cerrado de $X=\operatorname{Proj}(\mathcal{S})$, definido por un haz cuasicoherente de ideales $\mathcal{I}$ en $\mathcal{O}_X$.
    
    Sea $\mathcal{J}$ 
\oldpage[II]{65}
el haz graduado cuasicoherente de ideales de $\mathcal{S}$ dado por la imagen inversa de $\mathbf{\Gamma}_*(\mathcal{I})$ mediante el homomorfismo canónico $\alpha:\mathcal{S}\to\mathbf{\Gamma}_*(\mathcal{O}_X)$ \sref{II.3.3.2-fr}, y pongamos $\mathcal{S}'=\mathcal{S}/\mathcal{J}$.
    Entonces $X'$ es el subpreesquema asociado \sref[I]{I.4.2.1-fr} a la inmersión cerrada $\operatorname{Proj}(\mathcal{S}')\to X$ correspondiente al homomorfismo canónico $\mathcal{S}\to\mathcal{S}'$ de $\mathcal{O}_Y$-álgebras graduadas.
\end{enumerate}
\end{proposition}
% Fuente francesa: PMIHES_1961__8__5_0, printed pp. 65.

\begin{proof}
\medskip\noindent
\begin{enumerate}
  \item[(i)] Se puede suponer que $\varphi$ es sobreyectivo \sref{II.3.6.1-fr}.
    Entonces, para todo abierto afín $U$ de $Y$, $\Gamma(U,\mathcal{S})\to\Gamma(U,\mathcal{S}')$ es sobreyectivo \sref[I]{I.1.3.9-fr}, de modo que $G(\varphi)=X$ por \sref{II.3.5.1-fr}.
    Se reduce inmediatamente a demostrar la proposición en el caso en que $Y$ es afín, y entonces resulta de \sref{II.2.9.2-fr}[(i)].
  \item[(ii)] Se puede reducir a demostrar que el homomorfismo $\widetilde{\mathcal{J}}\to\mathcal{O}_X$ inducido por la inyección canónica $\mathcal{J}\to\mathcal{S}$ es un isomorfismo de $\widetilde{\mathcal{J}}$ sobre $\mathcal{I}$; como la cuestión es local sobre $Y$, se puede tomar $Y$ afín de anillo $A$, lo que implica que $\mathcal{S}=\widetilde{S}$, donde $S$ es una $A$-álgebra graduada engendrada por $S_1$, con $S_1$ de tipo finito sobre $A$.
    Basta entonces aplicar \sref{II.2.9.2-fr}[(ii)].
\end{enumerate}
\end{proof}

\begin{corollary}[3.6.3]
\label{II.3.6.3-fr}
Bajo las condiciones de \sref{II.3.6.2-fr}[(i)], y suponiendo además que $\mathcal{S}$ está engendrada por $\mathcal{S}_1$.
Entonces $\Phi^*(\mathcal{O}_X(n))$ se identifica canónicamente con $\mathcal{O}_{X'}(n)$ para todo $n\in\mathbf{Z}$.
\end{corollary}

\begin{proof}
Se ha definido un isomorfismo canónico de este tipo cuando $Y$ es afín \sref{II.2.9.3-fr}; en el caso general, basta demostrar que los isomorfismos así definidos para cada abierto afín $U$ de $Y$ son compatibles con el paso de $U$ a un abierto afín $U'\subset U$, lo que es inmediato.
\end{proof}

\begin{corollary}[3.6.4]
\label{II.3.6.4-fr}
Sean $Y$ un preesquema, $\mathcal{S}$ una $\mathcal{O}_Y$-álgebra graduada cuasicoherente engendrada por $\mathcal{S}_1$, $\mathcal{M}$ un $\mathcal{O}_Y$-módulo cuasicoherente, $u$ un $\mathcal{O}_Y$-homomorfismo sobreyectivo $\mathcal{M}\to\mathcal{S}_1$, y $\overline{u}:\mathbf{S}_{\mathcal{O}_Y}(\mathcal{M})\to\mathcal{S}$ el homomorfismo de $\mathcal{O}_Y$-álgebras graduadas que prolonga $u$ \sref{II.1.7.4-fr}.
Entonces el morfismo correspondiente a $\overline{u}$ es una inmersión cerrada de $\operatorname{Proj}(\mathcal{S})$ en $\operatorname{Proj}(\mathbf{S}_{\mathcal{O}_Y}(\mathcal{M}))$.
\end{corollary}

\begin{proof}
En efecto, $\overline{u}$ es sobreyectivo por hipótesis, y se aplica \sref{II.3.6.1-fr}[(i)].
\end{proof}
% Fuente francesa: PMIHES_1961__8__5_0, printed pp. 65-66.

\subsection{Morfismos de un preesquema a un espectro homogéneo}
\label{subsection:II.3.7-fr}

\begin{env}[3.7.1]
\label{II.3.7.1-fr}
Sean $q:X\to Y$ un morfismo de preesquemas, $\mathcal{L}$ un $\mathcal{O}_X$-módulo invertible y $\mathcal{S}$ una $\mathcal{O}_Y$-álgebra graduada cuasicoherente positivamente graduada;
entonces $q^*(\mathcal{S})$ es una $\mathcal{O}_X$-álgebra graduada cuasicoherente positivamente graduada.
Consideremos la $\mathcal{O}_X$-álgebra graduada cuasicoherente $\mathcal{S}'=\bigoplus_{n\geq0}\mathcal{L}^{\otimes n}$;
supongamos que tenemos un homomorfismo de $\mathcal{O}_X$-álgebras graduadas
\[
  \psi:q^*(\mathcal{S})\longrightarrow
  \mathcal{S}'=\bigoplus_{n\geq0}\mathcal{L}^{\otimes n}
\]
que equivale a dar un $q$-morfismo de álgebras graduadas
\[
  \psi^\flat:\mathcal{S}\longrightarrow q_*(\mathcal{S}').
\]

Se sabe que $\operatorname{Proj}(\mathcal{S}')$ se identifica canónicamente con $X$ (\sref{II.3.1.7-fr} y \sref{II.3.1.8-fr}[(iii)]);
se obtiene así canónicamente a partir de $\psi$ un subconjunto abierto $G(\psi)$ de $X$ y un $Y$-morfismo
\[
\label{II.3.7.1.1-fr}
  r_{\mathcal{L},\psi}:G(\psi)\longrightarrow
  \operatorname{Proj}(\mathcal{S})=P
\tag{3.7.1.1}
\]
\oldpage[II]{66}
que llamaremos el morfismo asociado a $\mathcal{L}$ y $\psi$;
recordemos (\sref{II.3.5.6-fr}) que este morfismo se obtiene componiendo con la primera proyección $\pi:\operatorname{Proj}(q^*(\mathcal{S}))=P\times_YX\to P$ el $Y$-morfismo
\[
  \tau=\operatorname{Proj}(\psi):G(\psi)\longrightarrow
  \operatorname{Proj}(q^*(\mathcal{S})).
\]

\end{env}

\begin{env}[3.7.2]
\label{II.3.7.2-fr}
Describamos explícitamente $r=r_{\mathcal{L},\psi}$ en el caso en que $Y=\operatorname{Spec}(A)$ sea afín, de modo que $\mathcal{S}=\widetilde{S}$, donde $S$ es una $A$-álgebra positivamente graduada.
Supongamos en primer lugar que $X=\operatorname{Spec}(B)$ también sea afín y que $\mathcal{L}=\widetilde{L}$, donde $L$ es un $B$-módulo libre de rango~$1$.
Entonces $q^*(\mathcal{S})=(S\otimes_AB)^{\sim}$ \sref[I]{I.1.6.5-fr};
si $c$ es un generador de $L$, entonces $\psi_n:q^*(\mathcal{S}_n)\to\mathcal{L}^{\otimes n}$ corresponde a un homomorfismo $w_n:s\otimes b\to bv_n(s)c^{\otimes n}$ de $S_n\otimes_AB$ a $L^{\otimes n}$, donde $v_n:S_n\to B$ es un homomorfismo de $A$-módulos, siendo los $v_n$ conjuntamente un homomorfismo de álgebras $S\to B$.
Sea $f\in S_d$ (donde $d>0$), y pongamos $g=v_d(f)$;
se tiene $\pi^{-1}(D_+(f))=D_+(f\otimes1)$ por \sref{II.2.8.10-fr} y la identificación de $D_+(f)$ con $\operatorname{Spec}(S_{(f)})$ \sref{II.2.3.6-fr};
por otra parte, la fórmula \sref{II.2.8.1.1-fr} muestra (teniendo en cuenta la identificación canónica de $X$ con $\operatorname{Proj}(\mathcal{S}')$), que
\[
  \tau^{-1}(D_+(f\otimes1))=D(g)
\]
de donde
\[
\label{II.3.7.2.1-fr}
  r^{-1}(D_+(f))=D(g).
\tag{3.7.2.1}
\]

Además, el morfismo $\tau=\operatorname{Proj}(\psi)$, restringido a $D(g)$, corresponde al homomorfismo que envía $(s\otimes1)/(f\otimes1)^n$ (donde $s\in S_{nd}$) a $v_{nd}(s)/g^n$ \sref{II.2.8.1-fr}, y la proyección $\pi$, restringida a $D_+(f\otimes1)$, corresponde al homomorfismo $s/f^n\to(s\otimes1)/(f\otimes1)^n$;
se concluye así que $r$, restringido a $D(g)$, corresponde al homomorfismo de $A$-álgebras $\omega:S_{(f)}\to B_g$ tal que $\omega(s/f^n)=v_{nd}(s)/g^n$ (donde $s\in S_{nd}$ y $n>0$).
Pasando al caso en que $X$ es arbitrario (pero $Y$ sigue siendo afín), se obtiene así, teniendo en cuenta \sref{II.2.8.1-fr}, lo siguiente:
\end{env}

\begin{proposition}[3.7.3]
\label{II.3.7.3-fr}
Si $Y=\operatorname{Spec}(A)$ es afín y $\mathcal{S}=\widetilde{S}$, donde $S$ es una $A$-álgebra graduada, entonces, para todo $f\in S_d=\Gamma(Y,\mathcal{S}_d)$, se tiene
\[
\label{II.3.7.3.1-fr}
  r_{\mathcal{L},\psi}^{-1}(D_+(f))=X_{\psi^\flat(f)}
  \qquad\bigl(\text{donde }\psi^\flat(f)\in
  \Gamma(X,\mathcal{L}^{\otimes d})\bigr)
\tag{3.7.3.1}
\] y la restricción $X_{\psi^\flat(f)}\to D_+(f)=\operatorname{Spec}(S_{(f)})$ de $r_{\mathcal{L},\psi}$ corresponde \sref[I]{I.2.2.4-fr} al homomorfismo de álgebras
\[
\label{II.3.7.3.2-fr}
  \psi_{(f)}^\flat:S_{(f)}\longrightarrow
  \Gamma(X_{\psi^\flat(f)},\mathcal{O}_X)
\tag{3.7.3.2}
\]
tal que, para $s\in S_{nd}=\Gamma(Y,\mathcal{S}_{nd})$
\[
\label{II.3.7.3.3-fr}
  \psi_{(f)}^\flat(s/f^n)
  =(\psi^\flat(s)|X_{\psi^\flat(f)})
   (\psi^\flat(f)|X_{\psi^\flat(f)})^{-n}.
\tag{3.7.3.3}
\]
\end{proposition}
% Fuente francesa: PMIHES_1961__8__5_0, printed pp. 66-67.

Diremos que $r_{\mathcal{L},\psi}$ está definido en todas partes si $G(\psi)=X$.
Para que así sea, es evidentemente necesario y suficiente que $G(\psi)\cap q^{-1}(U)=q^{-1}(U)$ para todo abierto afín $U\subset Y$;
en otras palabras, la cuestión es \emph{local} sobre $Y$.
Si $Y$ es afín, $G(\psi)$ es la unión de los $r^{-1}(D_+(f))$ para los $f$ homogéneos de $S_+$ \sref{II.2.8.1-fr};
por \sref{II.3.7.3.1-fr}, los $X_{\psi^\flat(f)}$ deben formar un \emph{recubrimiento} de $X$;
en otras palabras:

\begin{corollary}[3.7.4]
\label{II.3.7.4-fr}
Bajo las hipótesis de \sref{II.3.7.3-fr}, para que $r_{\mathcal{L},\psi}$ esté definido en todas partes, es necesario y suficiente que, para todo $x\in X$, exista un entero $n>0$ y una sección $s$ de $\mathcal{S}_n$ sobre $Y$ tal que, si se pone $t=\psi^\flat(s)\in\Gamma(X,\mathcal{L}^{\otimes n})$, entonces $t(x)\neq0$.
\end{corollary}

\oldpage[II]{67}
Se observa que esta condición siempre se satisface si $\psi$ es \emph{\textbf{(TN)}-sobreyectivo}.

Del mismo modo, la cuestión de saber si $r_{\mathcal{L},\psi}$ es \emph{dominante} es local sobre $Y$, y se tiene:

\begin{corollary}[3.7.5]
\label{II.3.7.5-fr}
Bajo las hipótesis de \sref{II.3.7.3-fr}, para que $r_{\mathcal{L},\psi}$ sea dominante, es necesario y suficiente que, para todo entero $n>0$, toda sección $s\in S_n$ tal que $\psi^\flat(s)\in\Gamma(X,\mathcal{L}^{\otimes n})$ sea localmente nilpotente sea ella misma nilpotente.
\end{corollary}

\begin{proof}
Hay que expresar, en efecto, que $r_{\mathcal{L},\psi}^{-1}(D_+(s))$ no es vacío siempre que $D_+(s)$ no sea vacío, y el corolario se sigue así de \sref{II.3.7.3.1-fr} y \sref{II.2.3.7-fr}.
\end{proof}

\begin{proposition}[3.7.6]
\label{II.3.7.6-fr}
Sean $q:X\to Y$ un morfismo, $\mathcal{L}$ un $\mathcal{O}_X$-módulo invertible, $\mathcal{S}$ y $\mathcal{S}'$ $\mathcal{O}_Y$-álgebras graduadas cuasicoherentes, $u:\mathcal{S}'\to\mathcal{S}$ un homomorfismo de álgebras graduadas y $\psi:q^*(\mathcal{S})\to\bigoplus_{n\geq0}\mathcal{L}^{\otimes n}$ un homomorfismo de álgebras graduadas;
sea $\psi'=\psi\circ q^*(u)$ el homomorfismo compuesto.
Si $r_{\mathcal{L},\psi'}$ está definido en todas partes, también lo está $r_{\mathcal{L},\psi}$;
si $u$ es \textup{\textbf{(TN)}}-sobreyectivo, entonces, si $r_{\mathcal{L},\psi'}$ es dominante, también lo es $r_{\mathcal{L},\psi}$;
recíprocamente, si $u$ es \textup{\textbf{(TN)}}-inyectivo, entonces, si $r_{\mathcal{L},\psi}$ es dominante, también lo es $r_{\mathcal{L},\psi'}$.
\end{proposition}

\begin{proof}
Se sabe que $G(\psi')\subset G(\psi)$ \sref{II.2.8.4-fr}, de donde se obtiene la primera afirmación;
si $u$ es \textbf{(TN)}-sobreyectivo, $\operatorname{Proj}(u):\operatorname{Proj}(\mathcal{S})\to
\operatorname{Proj}(\mathcal{S}')$ está definido en todas partes y es una inmersión cerrada;
como $r_{\mathcal{L},\psi'}$ es la composición de $\operatorname{Proj}(u)$ y de la restricción de $r_{\mathcal{L},\psi}$ a $G(\psi')$, se concluye que, si $r_{\mathcal{L},\psi'}$ es dominante, también lo es $r_{\mathcal{L},\psi}$.
Finalmente, si $u$ es \textbf{(TN)}-inyectivo, se sabe que $\operatorname{Proj}(u)$ es un morfismo dominante de $G(u)$ a $\operatorname{Proj}(\mathcal{S}')$ \sref{II.2.8.3-fr};
como $G(\psi')$ es la imagen inversa de $G(u)$ mediante $r_{\mathcal{L},\psi}$, se ve que, si $r_{\mathcal{L},\psi}$ es dominante, también lo es $r_{\mathcal{L},\psi'}$.
\end{proof}
% Fuente francesa: PMIHES_1961__8__5_0, printed pp. 67-68.

\begin{proposition}[3.7.7]
\label{II.3.7.7-fr}
Sea $Y$ un preesquema cuasicompacto, $q:X\to Y$ un morfismo cuasicompacto, $\mathcal{L}$ un $\mathcal{O}_X$-módulo invertible y $\mathcal{S}$ una $\mathcal{O}_Y$-álgebra graduada cuasicoherente dada por el límite inductivo filtrante de un sistema inductivo $(\mathcal{S}^\lambda)$ de $\mathcal{O}_Y$-álgebras graduadas cuasicoherentes.
Sean $\varphi_\lambda:\mathcal{S}^\lambda\to\mathcal{S}$ el homomorfismo canónico, $\psi:q^*(\mathcal{S})\to\bigoplus_{n\geq0}\mathcal{L}^{\otimes n}$ un homomorfismo de álgebras graduadas, y pongamos $\psi_\lambda=\psi\circ q^*(\varphi_\lambda)$.
Para que $r_{\mathcal{L},\psi}$ esté definido en todas partes, es necesario y suficiente que exista algún $\lambda$ tal que $r_{\mathcal{L},\psi_\lambda}$ esté definido en todas partes;
entonces $r_{\mathcal{L},\psi_\mu}$ está definido en todas partes para $\mu\geq\lambda$.
\end{proposition}

\begin{proof}
La condición es suficiente por \sref{II.3.7.6-fr}.
Recíprocamente, supongamos que $r_{\mathcal{L},\psi}$ está definido en todas partes;
se puede reducir al caso en que $Y$ es afín, pues, si para todo abierto afín $U\subset Y$ existe $\lambda(U)$ tal que la restricción de $r_{\mathcal{L},\psi_{\lambda(U)}}$ a $q^{-1}(U)$ esté definida en todas partes, bastará ($Y$ siendo cuasicompacto) recubrir $Y$ por un número finito de abiertos afines $U_i$ y tomar $\lambda\geq\lambda(U_i)$ para todos los índices $i$, por \sref{II.3.7.6-fr}.
Si $Y$ es afín, la hipótesis implica que, para todo $x\in X$, existe una sección $s^{(x)}$ de $S_n$ tal que, si se pone $t^{(x)}=\psi^\flat(s^{(x)})$, entonces $t^{(x)}(x)\neq0$ (donde $t^{(x)}$ se considera como una sección de $\mathcal{L}^{\otimes n}$ sobre $X$), lo que implica que $t^{(x)}(z)\neq0$ para todo $z$ en una vecindad $V(x)$ de $x$.
Recubramos $X$ por un número finito de $V(x_i)$ y sean $s^{(i)}$ las secciones correspondientes de $S$;
existe entonces algún $\lambda$ tal que las $s^{(i)}$ sean todas de la forma $\varphi_\lambda(s_\lambda^{\prime(i)})$, con $s_\lambda^{\prime(i)}\in S^\lambda$ para todo $i$;
se sigue de \sref{II.3.7.4-fr} que $r_{\mathcal{L},\psi_\lambda}$ está definido en todas partes.
La última afirmación es una consecuencia trivial de \sref{II.3.7.6-fr}.
\end{proof}

\begin{corollary}[3.7.8]
\label{II.3.7.8-fr}
Bajo las hipótesis de \sref{II.3.7.7-fr}, si los $r_{\mathcal{L},\psi_\lambda}$ son dominantes, también lo es $r_{\mathcal{L},\psi}$;
la recíproca es cierta si los $\varphi_\lambda$ son todos inyectivos.
\end{corollary}

\oldpage[II]{68}
\begin{proof}
La segunda afirmación es un caso particular de \sref{II.3.7.6-fr};
para demostrar que $r_{\mathcal{L},\psi}$ es dominante, se puede restringir al caso en que $Y$ es afín;
si $s\in S$ es tal que $\psi^\flat(s)$ es localmente nilpotente, como se puede escribir $s=\varphi_\lambda(s_\lambda)$ para al menos un $\lambda$, se concluye de las hipótesis y de \sref{II.3.7.5-fr} que $s_\lambda$ es nilpotente, y por tanto también lo es $s$, por lo que se aplica el criterio de \sref{II.3.7.5-fr}.
\end{proof}
% Fuente francesa: PMIHES_1961__8__5_0, printed pp. 68-69.

\begin{namedtheorem}[3.7.9]{Observaciones}
\label{II.3.7.9-fr}
\medskip\noindent
\begin{enumerate}
  \item[(i)] Con la notación de \sref{II.3.7.1-fr}, y teniendo en cuenta \sref{II.3.2.10-fr}, se tiene, para todo $n\in\mathbf{Z}$, un homomorfismo canónico
    \[
    \label{II.3.7.9.1-fr}
      \theta:r_{\mathcal{L},\psi}^*(\mathcal{O}_P(n))\longrightarrow
      \mathcal{L}^{\otimes n}
    \tag{3.7.9.1}
    \]
    definido de manera general en \sref{II.3.5.6.2-fr}.
    Se ve inmediatamente que, bajo las condiciones de \sref{II.3.7.3-fr}, la restricción de este homomorfismo a $X_{\psi^\flat(f)}$ envía $s/f^k$ (donde $s\in S_{n+kd}$) al elemento $(\psi^\flat(s)|X_{\psi^\flat(f)})
    (\psi^\flat(f)|X_{\psi^\flat(f)})^{-k}$.
  \item[(ii)] Sea $\mathcal{F}$ un $\mathcal{O}_X$-módulo cuasicoherente, y supongamos que $q$ es cuasicompacto y separado, de modo que, para todo $n\geq0$, $q_*(\mathcal{F}\otimes\mathcal{L}^{\otimes n})$ es un $\mathcal{O}_Y$-módulo cuasicoherente \sref[I]{I.9.2.2-fr}.
    Sea $\mathcal{M}'=\bigoplus_{n\geq0}\mathcal{F}\otimes
    \mathcal{L}^{\otimes n}$, que es un $\mathcal{S}'$-módulo graduado cuasicoherente, y consideremos su imagen $\mathcal{M}=q_*(\mathcal{M}')=
    \bigoplus_{n\geq0}q_*(\mathcal{F}\otimes\mathcal{L}^{\otimes n})$ (que es un $\mathcal{S}$-módulo cuasicoherente, mediante el homomorfismo $\psi^\flat$).
    Vamos a demostrar la existencia de un homomorfismo canónico de $\mathcal{O}_X$-módulos
    \[
    \label{II.3.7.9.2-fr}
      \xi:r_{\mathcal{L},\psi}^*(\widetilde{\mathcal{M}})
      \longrightarrow\mathcal{F}|G(\psi).
    \tag{3.7.9.2}
    \]

    En efecto, ya se ha definido \sref{II.3.5.6.1-fr} un homomorfismo canónico
    \[
    \label{II.3.7.9.3-fr}
      r_{\mathcal{L},\psi}^*(\widetilde{\mathcal{M}})
      \longrightarrow
      (q^*(\mathcal{M})\otimes_{q^*(\mathcal{S})}\mathcal{S}')^{\sim}
      |G(\psi),
    \tag{3.7.9.3}
    \]
    donde el miembro de la derecha se considera como un haz cuasicoherente sobre $\operatorname{Proj}(\mathcal{S}')$.
    También se tiene un homomorfismo canónico
    \[
    \label{II.3.7.9.4-fr}
      q^*(q_*(\mathcal{M}'))\otimes_{q^*(\mathcal{S})}\mathcal{S}'
      \longrightarrow\mathcal{M}'
    \tag{3.7.9.4}
    \]
    para todo $\mathcal{S}'$-módulo graduado cuasicoherente $\mathcal{M}'$: para todo abierto $U$ de $X$, toda sección $t'$ de $q^*(q_*(\mathcal{M}'_h))$ sobre $U$ y toda sección $b'$ de $\mathcal{S}'_k$ sobre $U$, se hace corresponder a $t'\otimes b'$ la sección $b'\sigma(t')$ de $\mathcal{M}'_{h+k}$, donde $\sigma(t')$ es la sección de $\mathcal{M}'_h$ sobre $U$ que corresponde canónicamente \sref[0]{0.4.4.3-fr} a $t'$.
    Se obtiene así un homomorfismo canónico
    \[
    \label{II.3.7.9.5-fr}
      (q^*(q_*(\mathcal{M}'))\otimes_{q^*(\mathcal{S})}\mathcal{S}')^{\sim}
      |G(\psi)\longrightarrow\widetilde{\mathcal{M}}'|G(\psi)
    \tag{3.7.9.5}
    \]
    y, finalmente, puesto que $\widetilde{\mathcal{M}}'$ se identifica canónicamente con $\mathcal{F}$ \sref{II.3.2.9-fr}[(i)], se obtiene efectivamente el homomorfismo canónico buscado.

    Bajo las condiciones de \sref{II.3.7.3-fr}, la restricción de este homomorfismo a $X_{\psi^\flat(f)}$ actúa de la manera siguiente:
    dada una sección $t_{nd}$ de $\mathcal{F}\otimes\mathcal{L}^{\otimes nd}$ sobre $X$, si $t'_{nd}$ es la sección $t_{nd}$ considerada como sección de $q_*(\mathcal{F}\otimes\mathcal{L}^{\otimes nd})$ sobre $Y$, envía el elemento $t'_{nd}/f^n$ a la sección $(t_{nd}|X_{\psi^\flat(f)})
    (\psi^\flat(f)|X_{\psi^\flat(f)})^{-n}$ de $\mathcal{F}$ sobre $X_{\psi^\flat(f)}$.
\end{enumerate}
\end{namedtheorem}
% Fuente francesa: PMIHES_1961__8__5_0, printed pp. 69-70.

\oldpage[II]{69}
\subsection{Criterios para la inmersión en un espectro homogéneo}
\label{subsection:II.3.8-fr}

\begin{env}[3.8.1]
\label{II.3.8.1-fr}
Con la notación de \sref{II.3.7.1-fr}, la cuestión de saber si $r_{\mathcal{L},\psi}$ es una inmersión (resp. una inmersión abierta, una inmersión cerrada) es claramente \emph{local sobre $Y$}.
\end{env}

\begin{proposition}[3.8.2]
\label{II.3.8.2-fr}
Bajo las hipótesis de \sref{II.3.7.3-fr}, para que $r_{\mathcal{L},\psi}$ esté definido en todas partes y sea una inmersión, es necesario y suficiente que exista una familia de secciones $s_\alpha\in S_{n_\alpha}$ (donde $n_\alpha>0$) tal que, si se pone $f_\alpha=\psi^\flat(s_\alpha)$, se satisfagan las condiciones siguientes:
\begin{enumerate}
  \item[\textup{(i)}] Los $X_{f_\alpha}$ forman un recubrimiento de $X$.
  \item[\textup{(ii)}] Los $X_{f_\alpha}$ son abiertos afines.
  \item[\textup{(iii)}] Para todo $\alpha$ y todo $t\in\Gamma(X_{f_\alpha},\mathcal{O}_X)$, existe un entero $m>0$ y algún $s\in S_{mn_\alpha}$ tales que $t=(\psi^\flat(s)|X_{f_\alpha})(f_\alpha|X_{f_\alpha})^{-m}$.
\end{enumerate}

Para que $r_{\mathcal{L},\psi}$ esté definido en todas partes y sea una inmersión abierta, es necesario y suficiente que exista una familia $(s_\alpha)$ que satisfaga las condiciones \textup{(i), (ii) y (iii)} anteriores, así como:
\begin{enumerate}
  \item[\textup{(iv)}] Para todo $n>0$ y toda $s\in S_{nn_\alpha}$ tal que $\psi^\flat(s)|X_{f_\alpha}=0$, existe un entero $k>0$ tal que $s_\alpha^ks=0$.
\end{enumerate}

Para que $r_{\mathcal{L},\psi}$ esté definido en todas partes y sea una inmersión cerrada, es necesario y suficiente que exista una familia $(s_\alpha)$ que satisfaga las condiciones \textup{(i), (ii) y (iii)} anteriores, así como:
\begin{enumerate}
  \item[\textup{(v)}] Los $D_+(s_\alpha)$ forman un recubrimiento de $P=\operatorname{Proj}(S)$.
\end{enumerate}
\end{proposition}
% Fuente francesa: PMIHES_1961__8__5_0, printed pp. 69-70.

\begin{proof}
Para que $r$ sea una inmersión (resp. una inmersión cerrada), es necesario y suficiente que exista un recubrimiento de $r(G(\psi))$ (resp. de $P$) por los $D_+(s_\alpha)$ tal que, si se pone $V_\alpha=r^{-1}(D_+(s_\alpha))$, la restricción de $r$ a $V_\alpha$ sea una \emph{inmersión cerrada} de $V_\alpha$ en $D_+(s_\alpha)$ \sref[I]{I.4.2.4-fr}.
La condición (i) expresa que $r$ está definido en todas partes y también que los $D_+(s_\alpha)$ recubren $r(X)$, por \sref{II.3.7.3.1-fr};
como $D_+(s_\alpha)$ es afín, las condiciones (ii) y (iii) expresan que la restricción de $r$ a $X_{f_\alpha}$ es una inmersión cerrada en $D_+(s_\alpha)$, por \sref[I]{I.4.2.3-fr};
finalmente, puesto que las condiciones (iii) y (iv) expresan que $\psi_{(s_\alpha)}^\flat$ es biyectivo (usando la notación de \sref{II.3.7.3.2-fr}), las condiciones (ii), (iii) y (iv) expresan que la restricción de $r$ a $X_{f_\alpha}$ es un isomorfismo sobre $D_+(s_\alpha)$ para todo $\alpha$, y por tanto las condiciones (i), (ii), (iii) y (iv) juntas expresan que $r$ es una inmersión abierta.
\end{proof}

\begin{corollary}[3.8.3]
\label{II.3.8.3-fr}
Bajo las hipótesis de \sref{II.3.7.6-fr}, si $r_{\mathcal{L},\psi'}$ está definido en todas partes y es una inmersión, también lo está $r_{\mathcal{L},\psi}$.
Si además se supone que $u$ es \textup{\textbf{(TN)}}-sobreyectivo, entonces, si $r_{\mathcal{L},\psi'}$ es una inmersión abierta \textup{(resp.} cerrada\textup{)}, también lo es $r_{\mathcal{L},\psi}$.
\end{corollary}

\begin{proof}
Por \sref{II.3.8.2-fr}, existe una familia $s'_\alpha\in S'_{n_\alpha}$ tal que, si se pone $f_\alpha=\psi'^\flat(s'_\alpha)$, se satisfacen las condiciones (i), (ii) y (iii).
Pero si se pone $s_\alpha=u(s'_\alpha)$, también se tiene $f_\alpha=\psi^\flat(s_\alpha)$, y si $t=(\psi'^\flat(s')|X_{f_\alpha})(f_\alpha|X_{f_\alpha})^{-m}$, también $t=(\psi^\flat(s)|X_{f_\alpha})(f_\alpha|X_{f_\alpha})^{-m}$ tomando $s=u(s')$, de donde se obtiene la primera afirmación.
La segunda afirmación se sigue inmediatamente del hecho de que $\operatorname{Proj}(u)$ es una inmersión cerrada.
\end{proof}
% Fuente francesa: PMIHES_1961__8__5_0, printed pp. 70-71.

\begin{proposition}[3.8.4]
\label{II.3.8.4-fr}
Supongamos satisfechas las hipótesis de \sref{II.3.7.7-fr} y, además, que $q:X\to Y$ es un morfismo de tipo finito.
Entonces, para que $r_{\mathcal{L},\psi}$ esté definido en todas partes y sea una inmersión, es necesario y suficiente que exista algún $\lambda$ tal que $r_{\mathcal{L},\psi_\lambda}$ esté definido en todas partes y sea una inmersión;
en tal caso, $r_{\mathcal{L},\psi_\mu}$ también está definido en todas partes y es una inmersión para todo $\mu\geq\lambda$.
\end{proposition}

\oldpage[II]{70}
\begin{proof}
Teniendo en cuenta \sref{II.3.8.3-fr}, basta demostrar que, si $r_{\mathcal{L},\psi}$ está definido en todas partes y es una inmersión, entonces $r_{\mathcal{L},\psi_\lambda}$ también lo está para al menos un $\lambda$.
Por el mismo argumento que en \sref{II.3.7.7-fr}, usando la cuasicompacidad de $Y$, se reduce primero al caso en que $Y$ es afín.
Como $X$ es cuasicompacto, \sref{II.3.8.2-fr} muestra entonces la existencia de una familia \emph{finita} $(s_i)$ de elementos de $S$ ($s_i\in S_{n_i}$) que satisfacen las condiciones (i), (ii) y (iii).
El morfismo $X_{f_i}\to Y$ (donde $f_i=\psi^\flat(s_i)$) es de tipo finito: es un morfismo de esquemas afines y, por tanto, es cuasicompacto \sref[I]{I.6.6.1-fr}, y también localmente de tipo finito, puesto que $q$ es de tipo finito \sref[I]{I.6.3.2-fr}; la conclusión se sigue entonces de \sref[I]{I.6.6.3-fr}.
El anillo $B_i$ de $X_{f_i}$ es, por tanto, una $A$-álgebra de tipo finito \sref[I]{I.6.3.3-fr};
sea $(t_{ij})$ una familia de generadores de esta álgebra.
Por hipótesis, existen elementos $s'_{ij}\in S_{m_{ij}n_i}$ tales que
\[
  t_{ij}=(\psi^\flat(s'_{ij})|X_{f_i})
  (\psi^\flat(s_i)|X_{f_i})^{-m_{ij}}.
\]
También por hipótesis, existe algún $\lambda$ y elementos $s_{i\lambda}\in S_{n_i}^\lambda$ y $s'_{ij\lambda}\in S_{m_{ij}n_i}^\lambda$ cuyas imágenes mediante $\varphi_\lambda$ son, respectivamente, $s_i$ y $s'_{ij}$;
es claro que $r_{\mathcal{L},\psi_\lambda}$ satisface las condiciones (i), (ii) y (iii) de \sref{II.3.8.2-fr}.
\end{proof}
% Fuente francesa: PMIHES_1961__8__5_0, printed pp. 70-71.

\begin{proposition}[3.8.5]
\label{II.3.8.5-fr}
Sea $Y$ un esquema cuasicompacto, o un preesquema cuyo espacio subyacente sea noetheriano, $q:X\to Y$ un morfismo \emph{de tipo finito}, $\mathcal{L}$ un $\mathcal{O}_X$-módulo invertible, $\mathcal{S}$ una $\mathcal{O}_Y$-álgebra graduada cuasicoherente y $\psi:q^*(\mathcal{S})\to\bigoplus_{n\geq0}\mathcal{L}^{\otimes n}$ un homomorfismo de álgebras graduadas.
Para que $r_{\mathcal{L},\psi}$ esté definido en todas partes y sea una inmersión, es necesario y suficiente que existan un entero $n>0$ y un sub-$\mathcal{O}_Y$-módulo cuasicoherente \emph{de tipo finito} $\mathcal{E}$ de $\mathcal{S}_n$ tales que
\begin{enumerate}
  \item[\textup{a)}] el homomorfismo $\psi_n\circ q^*(j_n):q^*(\mathcal{E})\to\mathcal{L}^{\otimes n}$ (donde $j_n:\mathcal{E}\to\mathcal{S}$ es la inyección canónica) sea sobreyectivo; y
  \item[\textup{b)}] si se denota por $\mathcal{S}'$ la sub-$\mathcal{O}_Y$-álgebra (graduada) de $\mathcal{S}$ engendrada por $\mathcal{E}$, y por $\psi'$ el homomorfismo $\psi\circ q^*(j')$ (donde $j'$ es la inyección de $\mathcal{S}'$ en $\mathcal{S}$), $r_{\mathcal{L},\psi'}$ esté definido en todas partes y sea una inmersión.
\end{enumerate}

Si esto se cumple, todo sub-$\mathcal{O}_Y$-módulo cuasicoherente $\mathcal{E}'$ de $\mathcal{S}_n$ que contenga a $\mathcal{E}$ posee también la misma propiedad, al igual que la imagen de $\bigotimes^k\mathcal{E}$ en $\mathcal{S}_{kn}$ para todo $k>0$.
\end{proposition}
% Fuente francesa: PMIHES_1961__8__5_0, printed pp. 70-71.

\begin{proof}
El hecho de que la condición sea suficiente, y las dos últimas afirmaciones, son casos particulares de \sref{II.3.8.3-fr}, teniendo en cuenta el isomorfismo canónico entre $\operatorname{Proj}(\mathcal{S})$ y $\operatorname{Proj}(\mathcal{S}^{(k)})$ \sref{II.3.1.8-fr}.

Sea $(U_i)$ un recubrimiento abierto finito de $Y$ por abiertos afines y pongamos $A_i=A(U_i)$.
Como $q^{-1}(U_i)$ es cuasicompacto, la hipótesis de que $r_{\mathcal{L},\psi}$ esté definido en todas partes y sea una inmersión implica, junto con \sref{II.3.8.2-fr}, la existencia de una familia finita $(s_{ij})$ de elementos de $S^{(i)}=\Gamma(U_i,\mathcal{S})$ ($s_{ij}\in S_{n_{ij}}^{(i)}$) que satisfacen las condiciones (i), (ii) y (iii).
Se ve, como en la demostración de \sref{II.3.8.4-fr}, que el morfismo $X_{f_{ij}}\to U_i$ (donde $f_{ij}=\psi^\flat(s_{ij})$) es de tipo finito, y que el anillo $B_{ij}$ de $X_{f_{ij}}$ es, por tanto, una $A_i$-álgebra de tipo finito, cuyos generadores tienen la forma $(\psi^\flat(t_{ijk})|X_{f_{ij}})(f_{ij}|X_{f_{ij}})^{-m_{ijk}}$, donde $t_{ijk}\in S_{m_{ijk}n_{ij}}^{(i)}$.
Sea $n$ un múltiplo común de los $m_{ijk}n_{ij}$;
reemplazando (para cada $(i,j,k)$) $s_{ij}$ por una potencia $s_{ij}^\rho$ tal que $\rho m_{ijk}n_{ij}=n$, y multiplicando $t_{ijk}$ por $s_{ij}^{\rho-m_{ijk}}$, se puede suponer que, para cada $i$, los $s_{ij}$ y $t_{ijk}$ pertenecen a $S_n^{(i)}$ y que $m_{ijk}=1$.
Sea $E_i$ el sub-$A_i$-módulo de $S^{(i)}$ generado por estos elementos (para $i$ fijo).
Entonces existe un sub-$\mathcal{O}_Y$-módulo coherente de tipo 
\oldpage[II]{71}
finito $\mathcal{E}_i$ de $\mathcal{S}_n$ de tipo finito tal que $\mathcal{E}_i|U_i=(E_i)^{\sim}$ \sref[I]{I.9.4.7-fr}.

Es claro que el sub-$\mathcal{O}_Y$-módulo $\mathcal{E}$ de $\mathcal{S}_n$ dado por la suma de los $\mathcal{E}_i$ es el objeto buscado (pues cada sección $f_{ij}$ es tal que, para todo $x\in X_{f_{ij}}$, existe una vecindad afín $V\subset X_{f_{ij}}$ de $x$ tal que $f|V$ es una base de $\Gamma(V,\mathcal{L}^{\otimes n})$).
\end{proof}
% Fuente francesa: PMIHES_1961__8__5_0, printed pp. 71-72.

\section{Fibrados proyectivos; haces amplios}
\label{section:II.4-fr}

\subsection{Definición de los fibrados proyectivos}
\label{subsection:II.4.1-fr}

\begin{definition}[4.1.1]
\label{II.4.1.1-fr}
Sea $Y$ un preesquema, $\mathcal{E}$ un $\mathcal{O}_Y$-módulo cuasicoherente y $\mathbf{S}_{\mathcal{O}_Y}(\mathcal{E})$ la $\mathcal{O}_Y$-álgebra simétrica de $\mathcal{E}$ \sref{II.1.7.4-fr}, que es cuasicoherente \sref{II.1.7.7-fr}.
Se define el \emph{fibrado proyectivo sobre $Y$ definido por $\mathcal{E}$}, denotado por $\mathbf{P}(\mathcal{E})$, como el $Y$-esquema $P=\operatorname{Proj}(\mathbf{S}_{\mathcal{O}_Y}(\mathcal{E}))$.
El $\mathcal{O}_P$-módulo $\mathcal{O}_P(1)$ se denomina el haz fundamental sobre $P$.
\end{definition}

Cuando $Y$ es afín de anillo $A$, se tiene $\mathcal{E}=\widetilde{E}$ donde $E$ es un $A$-módulo, y se escribe entonces $\mathbf{P}(E)$ en lugar de $\mathbf{P}(\widetilde{E})$.

Cuando se toma $\mathcal{E}=\mathcal{O}_Y^n$, se escribe $\mathbf{P}_Y^{n-1}$ en lugar de $\mathbf{P}(\mathcal{E})$;
si además $Y$ es afín de anillo $A$, se escribe también $\mathbf{P}_A^{n-1}$ en lugar de $\mathbf{P}_Y^{n-1}$.
Como $\mathbf{S}_{\mathcal{O}_Y}(\mathcal{O}_Y)$ se identifica canónicamente con $\mathcal{O}_Y[T]$ \sref{II.1.7.4-fr}, $\mathbf{P}_Y^0$ se identifica canónicamente con $Y$ \sref{II.3.1.7-fr};
el ejemplo~\sref{II.2.4.3-fr} es entonces exactamente $\mathbf{P}_K^1$.

\begin{env}[4.1.2]
\label{II.4.1.2-fr}
Sean $\mathcal{E}$ y $\mathcal{F}$ $\mathcal{O}_Y$-módulos cuasicoherentes;
sea $u:\mathcal{E}\to\mathcal{F}$ un $\mathcal{O}_Y$-homomorfismo;
le corresponde canónicamente un homomorfismo $\mathbf{S}(u):\mathbf{S}_{\mathcal{O}_Y}(\mathcal{E})\to
\mathbf{S}_{\mathcal{O}_Y}(\mathcal{F})$ de $\mathcal{O}_Y$-álgebras graduadas \sref{II.1.7.4-fr}.
Si $u$ es \emph{sobreyectivo}, también lo es $\mathbf{S}(u)$, y por tanto, según \sref{II.3.6.2-fr}[(i)], $\operatorname{Proj}(\mathbf{S}(u))$ es una \emph{inmersión cerrada} $\mathbf{P}(\mathcal{F})\to\mathbf{P}(\mathcal{E})$, que denotamos por $\mathbf{P}(u)$.
Se puede decir así que $\mathbf{P}(\mathcal{E})$ es un \emph{funtor contravariante} en $\mathcal{E}$, a condición de tomar como morfismos únicamente los homomorfismos \emph{sobreyectivos} de $\mathcal{O}_Y$-módulos cuasicoherentes.

Suponiendo todavía que $u$ es sobreyectivo, y poniendo $P=\mathbf{P}(\mathcal{E})$, $Q=\mathbf{P}(\mathcal{F})$ y $j=\mathbf{P}(u)$, se tiene, salvo isomorfismo,
\[
\label{II.4.1.2.1-fr}
  j^*(\mathcal{O}_P(n))=\mathcal{O}_Q(n)
  \qquad\text{para todo }n\in\mathbf{Z}
\tag{4.1.2.1}
\]
por \sref{II.3.6.3-fr}.
\end{env}

\begin{env}[4.1.3]
\label{II.4.1.3-fr}
Sea ahora $\psi:Y'\to Y$ un morfismo, y sea $\mathcal{E}'=\psi^*(\mathcal{E})$;
entonces $\mathbf{S}_{\mathcal{O}_{Y'}}(\mathcal{E}')=
\psi^*(\mathbf{S}_{\mathcal{O}_Y}(\mathcal{E}))$ \sref{II.1.7.5-fr};
por tanto, según \sref{II.3.5.3-fr},
\[
\label{II.4.1.3.1-fr}
  \mathbf{P}(\psi^*(\mathcal{E}))=
  \mathbf{P}(\mathcal{E})\times_YY'
\tag{4.1.3.1}
\]
salvo un isomorfismo canónico;
además, claramente se tiene
\[
  \psi^*((\mathbf{S}_{\mathcal{O}_Y}(\mathcal{E}))(n))
  =(\mathbf{S}_{\mathcal{O}_{Y'}}(\mathcal{E}'))(n)
\]
para todo $n\in\mathbf{Z}$, de donde, poniendo $P=\mathbf{P}(\mathcal{E})$ y $P'=\mathbf{P}(\mathcal{E}')$, se tiene, por \sref{II.3.5.4-fr}, salvo isomorfismo,
\[
\label{II.4.1.3.2-fr}
  \mathcal{O}_{P'}(n)=\mathcal{O}_P(n)\otimes_Y\mathcal{O}_{Y'}
  \qquad\text{para todo }n\in\mathbf{Z}.
\tag{4.1.3.2}
\]
\end{env}
% Fuente francesa: PMIHES_1961__8__5_0, printed pp. 72-73.

\oldpage[II]{72}
\begin{proposition}[4.1.4]
\label{II.4.1.4-fr}
Sea $\mathcal{L}$ un $\mathcal{O}_Y$-módulo invertible.
Para todo $\mathcal{O}_Y$-módulo cuasicoherente $\mathcal{E}$, existe un $Y$-isomorfismo canónico $i_{\mathcal{L}}:\mathbf{P}(\mathcal{E})\xrightarrow{\sim}
\mathbf{P}(\mathcal{E}\otimes\mathcal{L})$;
además, si se pone $P=\mathbf{P}(\mathcal{E})$ y $Q=\mathbf{P}(\mathcal{E}\otimes\mathcal{L})$, entonces $i_{\mathcal{L}}^*(\mathcal{O}_Q(n))$ es canónicamente isomorfo a $\mathcal{O}_P(n)\otimes_Y\mathcal{L}^{\otimes n}$ para todo $n\in\mathbf{Z}$.
\end{proposition}

\begin{proof}
Observemos primero que, si $A$ es un anillo, $E$ un $A$-módulo y $L$ un $A$-módulo \emph{libre de rango uno}, se define canónicamente un homomorfismo de $A$-módulos
\[
  \mathbf{S}_n(E\otimes L)\longrightarrow
  \mathbf{S}_n(E)\otimes L^{\otimes n}
\]
enviando $(x_1\otimes y_1)\ldots(x_n\otimes y_n)$ al elemento
\[
  (x_1x_2\ldots x_n)\otimes(y_1\otimes y_2\otimes\ldots\otimes y_n)
  \qquad(x_i\in E,\ y_i\in L\text{ para }1\leq i\leq n)~;
\]
se ve inmediatamente (restringiéndose al caso $L=A$) que este homomorfismo es de hecho un isomorfismo.
Se obtiene así un isomorfismo canónico de $A$-álgebras graduadas $\mathbf{S}_A(E\otimes L)\xrightarrow{\sim}
\bigoplus_{n\geq0}\mathbf{S}_n(E)\otimes L^{\otimes n}$.
Volviendo a las condiciones de (4.1.4), las observaciones precedentes permiten definir un isomorfismo canónico de $\mathcal{O}_Y$-álgebras graduadas
\[
\label{II.4.1.4.1-fr}
  \mathbf{S}_{\mathcal{O}_Y}(\mathcal{E}\otimes_{\mathcal{O}_Y}\mathcal{L})
  \xrightarrow{\sim}
  \bigoplus_{n\geq0}\mathbf{S}_n(\mathcal{E})
  \otimes_{\mathcal{O}_Y}\mathcal{L}^{\otimes n}
\tag{4.1.4.1}
\]
definiendo este isomorfismo como un isomorfismo de prehaces y teniendo en cuenta \sref{II.1.7.4-fr}, \sref[I]{I.1.3.9-fr} y \sref[I]{I.1.3.12-fr}.
La proposición se sigue entonces de \sref{II.3.1.8-fr}[(iii)] y \sref{II.3.2.10-fr}.
\end{proof}

\begin{env}[4.1.5]
\label{II.4.1.5-fr}
Bajo las hipótesis de \sref{II.4.1.1-fr}, pongamos $P=\mathbf{P}(\mathcal{E})$, y denotemos por $p$ el morfismo estructural $P\to Y$.
Puesto que, por definición, $\mathcal{E}=(\mathbf{S}_{\mathcal{O}_Y}(\mathcal{E}))_1$, se tiene un homomorfismo canónico $\alpha_1:\mathcal{E}\to p_*(\mathcal{O}_P(1))$ \sref{II.3.3.2.2-fr}, y por tanto \sref[0]{0.4.4.3-fr} también un homomorfismo canónico
\[
\label{II.4.1.5.1-fr}
  \alpha_1^\sharp:p^*(\mathcal{E})\longrightarrow\mathcal{O}_P(1).
\tag{4.1.5.1}
\]
\end{env}

\begin{proposition}[4.1.6]
\label{II.4.1.6-fr}
El homomorfismo canónico \sref{II.4.1.5.1-fr} es sobreyectivo.
\end{proposition}

\begin{proof}
Se ha visto, en \sref{II.3.3.2-fr}, que $\alpha_1^\sharp$ corresponde funtorialmente al homomorfismo canónico $\mathcal{E}\otimes_{\mathcal{O}_Y}\mathbf{S}_{\mathcal{O}_Y}(\mathcal{E})\longrightarrow(\mathbf{S}_{\mathcal{O}_Y}(\mathcal{E}))(1)$;
como, por definición, $\mathcal{E}$ genera $\mathbf{S}_{\mathcal{O}_Y}(\mathcal{E})$, este homomorfismo es sobreyectivo, de donde se obtiene la conclusión por \sref{II.3.2.4-fr}.
\end{proof}
% Fuente francesa: PMIHES_1961__8__5_0, printed pp. 72-73.

\subsection{Morfismos de un preesquema a un fibrado proyectivo}
\label{subsection:II.4.2-fr}

\begin{env}[4.2.1]
\label{II.4.2.1-fr}
Conservando la notación de \sref{II.4.1.5-fr}, sea $X$ un $Y$-preesquema, $q:X\to Y$ el morfismo estructural, y sea $r:X\to P$ un $Y$-\emph{morfismo} tal que el diagrama siguiente sea conmutativo:
\[
  \xymatrix{
    P \ar[d]_p & X \ar[l]_r \ar[dl]^q
    \\Y
  }
\]
\oldpage[II]{73}

Puesto que el funtor $r^*$ es exacto por la derecha \sref[0]{0.4.3.1-fr}, se deduce del homomorfismo \sref{II.4.1.5.1-fr}, que es sobreyectivo \sref{II.4.1.6-fr}, un homomorfismo sobreyectivo
\[
  r^*(\alpha_1^\sharp):r^*(p^*(\mathcal{E}))\longrightarrow
  r^*(\mathcal{O}_P(1)).
\]

Pero $r^*(p^*(\mathcal{E}))=q^*(\mathcal{E})$, y $r^*(\mathcal{O}_P(1))$ es localmente isomorfo a $r^*(\mathcal{O}_P)=\mathcal{O}_X$; en otras palabras, es un haz \emph{invertible} $\mathcal{L}_r$ sobre $\mathcal{O}_X$.
Así, a partir de $r$ se ha definido un homomorfismo \emph{sobreyectivo} canónico de $\mathcal{O}_X$-módulos
\[
\label{II.4.2.1.1-fr}
  \varphi_r:q^*(\mathcal{E})\longrightarrow\mathcal{L}_r.
\tag{4.2.1.1}
\]

Cuando $Y=\operatorname{Spec}(A)$ es afín y $\mathcal{E}=\widetilde{E}$, este homomorfismo se explicita además de la forma siguiente: dado $f\in E$, se deduce de \sref{II.2.6.3-fr} que
\[
\label{II.4.2.1.2-fr}
  r^{-1}(D_+(f))=X_{\varphi_r^\flat(f)}.
\tag{4.2.1.2}
\]

Sea ahora $V$ un abierto afín de $X$ contenido en $r^{-1}(D_+(f))$, y sea $B$ su anillo, que es una $A$-álgebra;
pongamos $S=\mathbf{S}_A(E)$;
la restricción de $r$ a $V$ corresponde a un $A$-homomorfismo $\omega:S_{(f)}\to B$, y se tiene $q^*(\mathcal{E})|V=(E\otimes_A B)\supertilde$ y $\mathcal{L}_r|V=\widetilde{L}_r$, de donde $L_r=(S(1))_{(f)}\otimes_{S_{(f)}}B_{[\omega]}$ \sref[I]{I.1.6.5-fr}.
La restricción de $\varphi_r$ a $q^*(\mathcal{E})|V$ corresponde al $B$-homomorfismo $u:E\otimes_A B\to L_r$, que envía $x\otimes1$ a $(x/1)\otimes1=(f/1)\otimes\omega(x/f)$.
La extensión canónica de $\varphi_r$ a un homomorfismo de $\mathcal{O}_X$-álgebras
\[
  \psi_r:q^*(\mathbf{S}(\mathcal{E}))=
  \mathbf{S}(q^*(\mathcal{E}))\longrightarrow
  \mathbf{S}(\mathcal{L}_r)=\bigoplus_{n\geq0}\mathcal{L}_r^{\otimes n}
\]
es, por tanto, tal que la restricción de $\psi_r$ a $q^*(\mathbf{S}_n(\mathcal{E}))|V$ corresponde al homomorfismo
$\mathbf{S}_n(E\otimes_A B)=\mathbf{S}_n(E)\otimes_A B\to L_r^{\otimes n}$ que envía $s\otimes1$ a $(f/1)^{\otimes n}\otimes\omega(s/f^n)$.
\end{env}

\begin{env}[4.2.2]
\label{II.4.2.2-fr}
Recíprocamente, supongamos que se nos dan un morfismo $q:X\to Y$, un $\mathcal{O}_X$-módulo invertible $\mathcal{L}$ y un $\mathcal{O}_Y$-módulo cuasicoherente $\mathcal{E}$;
cada homomorfismo $\varphi:q^*(\mathcal{E})\to\mathcal{L}$ determina canónicamente un homomorfismo de $\mathcal{O}_X$-álgebras cuasicoherentes
\[
  \psi:\mathbf{S}(q^*(\mathcal{E}))=q^*(\mathbf{S}(\mathcal{E}))
  \longrightarrow\bigoplus_{n\geq0}\mathcal{L}^{\otimes n}
\]
y por tanto, según \sref{II.3.7.1-fr}, un $Y$-morfismo $r_{\mathcal{L},\psi}:G(\psi)\to
\operatorname{Proj}(\mathbf{S}(\mathcal{E}))=\mathbf{P}(\mathcal{E})$, que denotamos por $r_{\mathcal{L},\varphi}$.
Si $\varphi$ es \emph{sobreyectivo}, también lo es $\psi$, y por tanto, según \sref{II.3.7.4-fr}, $r_{\mathcal{L},\varphi}$ está \emph{definido en todas partes}.
Además, con la notación de \sref{II.4.2.1-fr} y (4.2.2):
\end{env}
% Fuente francesa: PMIHES_1961__8__5_0, printed pp. 73-74.

\begin{proposition}[4.2.3]
\label{II.4.2.3-fr}
Dados un morfismo $q:X\to Y$ y un $\mathcal{O}_Y$-módulo cuasicoherente $\mathcal{E}$, las asignaciones $r\to(\mathcal{L}_r,\varphi_r)$ y $(\mathcal{L},\varphi)\to r_{\mathcal{L},\varphi}$ establecen una correspondencia biyectiva entre el conjunto de $Y$-morfismos $r:X\to\mathbf{P}(\mathcal{E})$ y el conjunto de clases de equivalencia de pares $(\mathcal{L},\varphi)$ formados por un $\mathcal{O}_X$-módulo invertible $\mathcal{L}$ y un homomorfismo sobreyectivo $\varphi:q^*(\mathcal{E})\to\mathcal{L}$, donde dos pares $(\mathcal{L},\varphi)$ y $(\mathcal{L}',\varphi')$ son equivalentes si existe un $\mathcal{O}_X$-isomorfismo $\tau:\mathcal{L}\xrightarrow{\sim}\mathcal{L}'$ tal que $\varphi'=\tau\circ\varphi$.
\end{proposition}

\begin{proof}
Comencemos con un $Y$-morfismo $r:X\to\mathbf{P}(\mathcal{E})$, construyamos $\mathcal{L}_r$ y $\varphi_r$ según \sref{II.4.2.1-fr}, y sea $r'=r_{\mathcal{L}_r,\varphi_r}$;
se sigue inmediatamente de \sref{II.4.2.1-fr} y \sref{II.3.7.2-fr} que los morfismos $r$ y $r'$ son idénticos (tomando como generador de $\mathcal{L}_r$ el elemento $(f/1)\otimes1$ para definir los homomorfismos $v_n$ de \sref{II.3.7.2-fr}).
Recíprocamente, tomemos un par $(\mathcal{L},\varphi)$ y construyamos
\oldpage[II]{74}
$r''=r_{\mathcal{L},\varphi}$, y después $\mathcal{L}_{r''}$ y $\varphi_{r''}$;
mostraremos que existe un isomorfismo canónico $\tau:\mathcal{L}_{r''}\xrightarrow{\sim}\mathcal{L}$ tal que $\varphi=\tau\circ\varphi_{r''}$;
para definirlo, se puede restringir al caso en que $Y=\operatorname{Spec}(A)$ y $X=\operatorname{Spec}(B)$ sean afines, y (con la notación de \sref{II.4.2.1-fr} y \sref{II.3.7.2-fr}) asociar a cada elemento $(x/1)\otimes1$ de $L_{r''}$ (donde $x\in E$) el elemento $v_1(x)c$ de $L$.
Se ve inmediatamente que $\tau$ no depende del generador $c$ elegido de $L$;
como $v_1$ es sobreyectivo por hipótesis, para demostrar que $\tau$ es un isomorfismo basta probar que, si $x/1=0$ en $(S(1))_{(f)}$, entonces $v_1(x)/1=0$ en $B_g$;
pero la primera igualdad implica que $f^nx=0$ en $\mathbf{S}_{n+1}(E)$ para algún $n$, y esto implica que $v_{n+1}(f^nx)=g^nv_1(x)=0$ en $B$, de donde se obtiene la conclusión.
Finalmente, es inmediato que, para dos pares equivalentes $(\mathcal{L},\varphi)$ y $(\mathcal{L}',\varphi')$, se tiene $r_{\mathcal{L},\varphi}=r_{\mathcal{L}',\varphi'}$.
\end{proof}

En particular, para $X=Y$:
\begin{theorem}[4.2.4]
\label{II.4.2.4-fr}
El conjunto de secciones sobre $Y$ de $\mathbf{P}(\mathcal{E})$ está en correspondencia biyectiva canónica con el conjunto de sub-$\mathcal{O}_Y$-módulos cuasicoherentes $\mathcal{F}$ de $\mathcal{E}$ tales que $\mathcal{E}/\mathcal{F}$ es invertible.
\end{theorem}

Se observa que esta propiedad corresponde a la definición clásica de «espacio proyectivo» como conjunto de hiperplanos de un espacio vectorial (el caso clásico corresponde a $Y=\operatorname{Spec}(K)$, donde $K$ es un cuerpo, y $\mathcal{E}=\widetilde{E}$, donde $E$ es un espacio vectorial de dimensión finita sobre $K$; los $\mathcal{F}$ que tienen la propiedad descrita en \sref{II.4.2.4-fr} corresponden entonces a los hiperplanos de $E$, y ya se sabe que las secciones sobre $Y$ de $\mathbf{P}(\mathcal{E})$ son entonces los \emph{puntos de $\mathbf{P}(\mathcal{E})$ racionales sobre} $K$ \sref[I]{I.3.4.5-fr}).
% Fuente francesa: PMIHES_1961__8__5_0, printed pp. 74-76.

\begin{remark}[4.2.5]
\label{II.4.2.5-fr}
Como existe una correspondencia biyectiva canónica entre los $Y$-morfismos de $X$ a $P$ y sus morfismos gráficos, es decir, las secciones sobre $X$ de $P\times_YX$ \sref[I]{I.3.3.14-fr}, se ve que, recíprocamente, \sref{II.4.2.3-fr} puede deducirse de \sref{II.4.2.4-fr}.
Denotemos por $\operatorname{Hyp}_Y(X,\mathcal{E})$ el conjunto de sub-$\mathcal{O}_X$-módulos cuasicoherentes $\mathcal{F}$ de $\mathcal{E}\otimes_Y\mathcal{O}_X=q^*(\mathcal{E})$ tales que $q^*(\mathcal{E})/\mathcal{F}$ es un $\mathcal{O}_X$-módulo invertible.
Si $g:X'\to X$ es un $Y$-morfismo, la exactitud por la derecha de $g^*$ da
$g^*(q^*(\mathcal{E})/\mathcal{F})=g^*(q^*(\mathcal{E}))/g^*(\mathcal{F})$,
por lo que este último haz es invertible. Así, $\operatorname{Hyp}_Y(X,\mathcal{E})$ es un \emph{funtor contravariante} en la categoría de los $Y$-preesquemas.
Se puede, por tanto, interpretar el teorema \sref{II.4.2.4-fr} como la definición de un \emph{isomorfismo canónico} entre los funtores $\operatorname{Hom}_Y(X,\mathbf{P}(\mathcal{E}))$ y $\operatorname{Hyp}_Y(X,\mathcal{E})$, ambos contravariantes en $X$ sobre la categoría de los $Y$-preesquemas.
Esto proporciona también una caracterización del fibrado proyectivo $P=\mathbf{P}(\mathcal{E})$ mediante la siguiente \emph{propiedad universal}, mucho más próxima a la intuición geométrica que las construcciones de §§2 y 3:
para todo morfismo $q:X\to Y$ y todo $\mathcal{O}_X$-módulo invertible $\mathcal{L}$ que sea un cociente de $\mathcal{E}\otimes_{\mathcal{O}_Y}\mathcal{O}_X$, existe un único $Y$-morfismo $r:X\to P$ tal que $\mathcal{L}=r^*(\mathcal{O}_P(1))$.

Se verá más adelante que también se puede definir, entre otras cosas, esquemas «grassmannianos».
\end{remark}

\begin{corollary}[4.2.6]
\label{II.4.2.6-fr}
Supongamos que todo $\mathcal{O}_Y$-módulo invertible es trivial \sref[I]{I.2.4.8-fr}.
Sea $V$ el grupo $\operatorname{Hom}_{\mathcal{O}_Y}(\mathcal{E},\mathcal{O}_Y)$, considerado como módulo sobre el anillo $A=\Gamma(Y,\mathcal{O}_Y)$, y sea $V^*$ el subconjunto de $V$ formado por los homomorfismos sobreyectivos.
Entonces el conjunto de secciones sobre $Y$ de $\mathbf{P}(\mathcal{E})$ se identifica canónicamente con $V^*/A^*$, donde $A^*$ es el grupo de las unidades de $A$.
\end{corollary}

\oldpage[II]{75}
En particular:
\begin{enumerate}
  \item El corolario \sref{II.4.2.6-fr} se aplica siempre que $Y$ sea un \emph{esquema local} \sref[I]{I.2.4.8-fr}.
    Sea $Y$ un preesquema arbitrario, $y$ un punto de $Y$ y $Y'=\operatorname{Spec}(\boldsymbol{k}(y))$;
    entonces la fibra $p^{-1}(y)$ de $\mathbf{P}(\mathcal{E})$ puede identificarse, por \sref{II.4.1.3.1-fr}, con $\mathbf{P}(\mathcal{E}^y)$, donde $\mathcal{E}^y=\mathcal{E}_y\otimes_{\mathcal{O}_y}\boldsymbol{k}(y)
  =\mathcal{E}_y/\mathfrak{m}_y\mathcal{E}_y$ se considera como espacio vectorial sobre $\boldsymbol{k}(y)$.
    Más generalmente, si $K$ es una extensión de $\boldsymbol{k}(y)$, entonces $p^{-1}(y)\otimes_{\boldsymbol{k}(y)}K$ puede identificarse con $\mathbf{P}(\mathcal{E}^y\otimes_{\boldsymbol{k}(y)}K)$.
    El corolario \sref{II.4.2.6-fr} muestra entonces que el conjunto de los \emph{puntos geométricos de $\mathbf{P}(\mathcal{E})$ con valores en la extensión $K$ de $\boldsymbol{k}(y)$} \sref[I]{I.3.4.5-fr}, que también se puede llamar la \emph{fibra geométrica racional sobre} $K$ de $\mathbf{P}(\mathcal{E})$ \emph{sobre el punto $y$}, puede identificarse con el \emph{espacio proyectivo} asociado al \emph{dual} del espacio vectorial sobre $K$ $\mathcal{E}^y\otimes_{\boldsymbol{k}(y)}K$.
  \item Supongamos que $Y$ es afín de anillo $A$ y, además, que todo $\mathcal{O}_Y$-módulo invertible es trivial;
    tomemos además $\mathcal{E}=\mathcal{O}_Y^n$;
    entonces, en \sref{II.4.2.6-fr}, $V$ puede identificarse con $A^n$ \sref[I]{I.1.3.8-fr}, y $V^*$ con el conjunto de sistemas $(f_i)_{1\leq i\leq n}$ de elementos de $A$ que generan el ideal $A$;
    dos sistemas cualesquiera de este tipo definen la misma sección sobre $Y$ de $\mathbf{P}_Y^{n-1}=\mathbf{P}_A^{n-1}$, o, en otras palabras, \emph{el mismo punto de $\mathbf{P}_A^{n-1}$ con valores en $A$}, si y solo si uno de ellos se obtiene del otro multiplicándolo por un elemento invertible de $A$.
\end{enumerate}

Estas propiedades justifican la terminología «fibrado proyectivo» para $\mathbf{P}(\mathcal{E})$.
Se observa que la definición de «espacio proyectivo» que obtenemos así es de hecho \emph{dual} a la definición clásica;
esto nos lo impone la necesidad de poder definir $\mathbf{P}(\mathcal{E})$ para un $\mathcal{O}_Y$-módulo cuasicoherente \emph{cualquiera} $\mathcal{E}$, no necesariamente localmente libre.
% Fuente francesa: PMIHES_1961__8__5_0, printed pp. 76.

\begin{remark}[4.2.7]
\label{II.4.2.7-fr}
Veremos, en el capítulo~V, que, si $Y$ es no vacío, conexo y localmente noetheriano, y si $\mathcal{E}$ es localmente libre de rango al menos $2$, entonces, poniendo $P=\mathbf{P}(\mathcal{E})$, todo $\mathcal{O}_P$-módulo invertible $\mathcal{L}$ es isomorfo a un $\mathcal{O}_P$-módulo de la forma $\mathcal{L}'\otimes_{\mathcal{O}_Y}\mathcal{O}_P(m)$, donde $\mathcal{L}'$ es un $\mathcal{O}_Y$-módulo invertible, bien definido salvo isomorfismo, y $m$ es un entero bien definido.
En otras palabras, $\mathrm{H}^1(P,\mathcal{O}_P^*)$ es isomorfo a $\mathbf{Z}\times\mathrm{H}^1(Y,\mathcal{O}_Y^*)$ \sref[0]{0.5.4.7-fr}.
También veremos (\sref[III]{III.2.1.15-fr}, teniendo en cuenta \sref[0]{0.5.4.10-fr}) que $p_*(\mathcal{L})=0$ si $m<0$, y que $p_*(\mathcal{L})$ es isomorfo a $\mathcal{L}'\otimes_{\mathcal{O}_Y}
(\mathbf{S}_{\mathcal{O}_Y}(\mathcal{E}))_m$ si $m\geq0$.
Si $\mathcal{F}$ es un $\mathcal{O}_Y$-módulo cuasicoherente, todo $Y$-morfismo $\mathbf{P}(\mathcal{E})\to\mathbf{P}(\mathcal{F})$ está determinado por los datos de un $\mathcal{O}_Y$-módulo invertible $\mathcal{L}'$, un entero $m\geq0$ y un $\mathcal{O}_Y$-homomorfismo $\psi:\mathcal{F}\to\mathcal{L}'\otimes_{\mathcal{O}_Y}
(\mathbf{S}_{\mathcal{O}_Y}(\mathcal{E}))_m$ tal que el homomorfismo correspondiente $\psi^\sharp$ de $\mathcal{O}_{\mathbf{P}(\mathcal{E})}$-módulos sea sobreyectivo.
También veremos que, si el $Y$-morfismo en cuestión es un isomorfismo, entonces $m=1$ y $\mathcal{F}$ es isomorfo a $\mathcal{E}\otimes_{\mathcal{O}_Y}\mathcal{L}'$ (la recíproca de \sref{II.4.1.4-fr}).
Esto permitirá determinar el haz de gérmenes de automorfismos de $\mathbf{P}(\mathcal{E})$ como el cociente del haz de grupos $\mathcal{A}ut(\mathcal{E})$ (que es localmente isomorfo a $\mathbf{GL}(n,\mathcal{O}_Y)$ si $\mathcal{E}$ es de rango $n$) por $\mathcal{O}_Y^*$.
\end{remark}

\begin{env}[4.2.8]
\label{II.4.2.8-fr}
Conservando la notación de \sref{II.4.2.1-fr}, sea $u:X'\to X$ un morfismo;
si el $Y$-morfismo $r:X\to P$ corresponde al homomorfismo $\varphi:q^*(\mathcal{E})\to\mathcal{L}$, entonces el $Y$-morfismo $r\circ u$ corresponde a $u^*(\varphi):u^*(q^*(\mathcal{E}))\to u^*(\mathcal{L})$, como se sigue inmediatamente de las definiciones.
\end{env}

\begin{env}[4.2.9]
\label{II.4.2.9-fr}
Sean $\mathcal{E}$ y $\mathcal{F}$ $\mathcal{O}_Y$-módulos cuasicoherentes, $v:\mathcal{E}\to\mathcal{F}$ un homomorfismo sobreyectivo y $j=\mathbf{P}(v)$ la inmersión cerrada correspondiente $\mathbf{P}(\mathcal{F})\to\mathbf{P}(\mathcal{E})$ \sref{II.4.1.2-fr}.
Si el $Y$-morfismo $r:X\to\mathbf{P}(\mathcal{F})$ corresponde al homomorfismo $\varphi:q^*(\mathcal{F})\to\mathcal{L}$, entonces el
\oldpage[II]{76}
$Y$-morfismo $j\circ r$ corresponde a $q^*(\mathcal{E})\xrightarrow{q^*(v)}q^*(\mathcal{F})
\xrightarrow{\varphi}\mathcal{L}$;
esto se sigue también de la definición dada en \sref{II.4.2.1-fr}.
\end{env}
% Fuente francesa: PMIHES_1961__8__5_0, printed p. 76.

\begin{env}[4.2.10]
\label{II.4.2.10-fr}
Sea $\psi:Y'\to Y$ un morfismo, y pongamos $\mathcal{E}'=\psi^*(\mathcal{E})$.
Si el $Y$-morfismo $r:X\to P$ corresponde al homomorfismo $\varphi:q^*(\mathcal{E})\to\mathcal{L}$, entonces el $Y'$-morfismo
\[
  r_{(Y')}:X_{(Y')}\longrightarrow P'=\mathbf{P}(\mathcal{E}')
\]
corresponde a $\varphi_{(Y')}:q_{(Y')}^*(\mathcal{E}')=
  q^*(\mathcal{E})\otimes_Y\mathcal{O}_{Y'}\to
\mathcal{L}\otimes_Y\mathcal{O}_{Y'}$.
En efecto, por \sref{II.4.1.3.1-fr}, se tiene el diagrama conmutativo
\[
  \xymatrix{
    Y' \ar[d]
    & P'=P_{(Y')} \ar[l]_{p_{(Y')}} \ar[d]^{u}
    & X_{(Y')} \ar[l]_{r_{(Y')}} \ar[d]^{v}
    \\Y
    & P \ar[l]^{p}
    & X \ar[l]^{r}
  }
\]

Por \sref{II.4.1.3.2-fr}, se tiene
\[
  (r_{(Y')})^*(\mathcal{O}_{P'}(1))
  =(r_{(Y')})^*(u^*(\mathcal{O}_P(1)))
  =v^*(r^*(\mathcal{O}_P(1)))
  =v^*(\mathcal{L})=\mathcal{L}\otimes_Y\mathcal{O}_{Y'}~;
\]
por otra parte, $u^*(\alpha_1^\sharp)$ es exactamente el homomorfismo canónico $\alpha_1^\sharp:(p_{(Y')})^*(\mathcal{E}')\to\mathcal{O}_{P'}(1)$;
esto puede verse calculando explícitamente los homomorfismos canónicos $\alpha_1^\sharp$ para $P$ y $P'$ como en \sref{II.4.1.6-fr}.
De aquí se obtiene nuestra afirmación.
\end{env}
% Fuente francesa: PMIHES_1961__8__5_0, printed pp. 76-77.

\subsection{El morfismo de Segre}
\label{subsection:II.4.3-fr}

\begin{env}[4.3.1]
\label{II.4.3.1-fr}
Sea $Y$ un preesquema, y sean $\mathcal{E}$ y $\mathcal{F}$ $\mathcal{O}_Y$-módulos cuasicoherentes;
pongamos $P_1=\mathbf{P}(\mathcal{E})$ y $P_2=\mathbf{P}(\mathcal{F})$, y denotemos por $p_1:P_1\to Y$ y $p_2:P_2\to Y$ los morfismos estructurales.
Sea $Q=P_1\times_YP_2$, y sean $q_1:Q\to P_1$ y $q_2:Q\to P_2$ las proyecciones canónicas;
entonces el $\mathcal{O}_Q$-módulo
$\mathcal{L}=\mathcal{O}_{P_1}(1)\otimes_Y\mathcal{O}_{P_2}(1)
=q_1^*(\mathcal{O}_{P_1}(1))\otimes_{\mathcal{O}_Q}
q_2^*(\mathcal{O}_{P_2}(1))$ es invertible, por ser el producto tensorial de dos $\mathcal{O}_Q$-módulos invertibles \sref[0]{0.5.4.4-fr}.
Además, si $r=p_1\circ q_1=p_2\circ q_2$ es el morfismo estructural $Q\to Y$, entonces
$r^*(\mathcal{E}\otimes_{\mathcal{O}_Y}\mathcal{F})
=q_1^*(p_1^*(\mathcal{E}))\otimes_{\mathcal{O}_Q}
q_2^*(p_2^*(\mathcal{F}))$ \sref[0]{0.4.3.3-fr};
los homomorfismos canónicos sobreyectivos \sref{II.4.1.5.1-fr}, $p_1^*(\mathcal{E})\to\mathcal{O}_{P_1}(1)$ y $p_2^*(\mathcal{F})\to\mathcal{O}_{P_2}(1)$, dan así, tomando su producto tensorial, un homomorfismo canónico
\[
\label{II.4.3.1.1-fr}
  s:r^*(\mathcal{E}\otimes_{\mathcal{O}_Y}\mathcal{F})\longrightarrow
  \mathcal{L}
\tag{4.3.1.1}
\]
evidentemente sobreyectivo;
se deduce de él, por \sref{II.4.2.2-fr}, un morfismo canónico, llamado \emph{morfismo de Segre}:
\[
\label{II.4.3.1.2-fr}
  \varsigma:\mathbf{P}(\mathcal{E})\times_Y\mathbf{P}(\mathcal{F})
  \longrightarrow
  \mathbf{P}(\mathcal{E}\otimes_{\mathcal{O}_Y}\mathcal{F}).
\tag{4.3.1.2}
\]

Explicitemos el morfismo $\varsigma$ más explícitamente cuando $Y=\operatorname{Spec}(A)$ es afín, $\mathcal{E}=\widetilde{E}$ y $\mathcal{F}=\widetilde{F}$, donde $E$ y $F$ son $A$-módulos, de modo que $\mathcal{E}\otimes_{\mathcal{O}_Y}\mathcal{F}=(E\otimes_A F)\supertilde$ \sref[I]{I.1.3.12-fr};
pongamos $R=\mathbf{S}_A(E)$, $S=\mathbf{S}_A(F)$ y $T=\mathbf{S}_A(E\otimes_A F)$;
sean $f\in E$ y $g\in F$, y consideremos el abierto afín
\[
  D_+(f)\times_YD_+(g)=\operatorname{Spec}(B)
\]
\oldpage[II]{77}
de $Q$, donde $B=R_{(f)}\otimes_AS_{(g)}$;
la restricción de $\mathcal{L}$ a este abierto afín es $\widetilde{L}$, donde
\[
  L=(R(1))_{(f)}\otimes_A(S(1))_{(g)}
\]
y el elemento $c=(f/1)\otimes(g/1)$ es un generador de $L$ considerado como $B$-módulo libre \sref{II.2.5.7-fr}.
El homomorfismo \sref{II.4.3.1.1-fr} corresponde al homomorfismo
\[
  (x\otimes y)\otimes b\to b((x/1)\otimes(y/1))
\]
de $(E\otimes_A F)\otimes_AB$ a $L$.
Con la notación de \sref{II.3.7.2-fr}, se tiene así $v_1(x\otimes y)=(x/f)\otimes(y/g)$;
la restricción de $\varsigma$ a $D_+(f)\times_YD_+(g)$ es un morfismo de este esquema afín a $D_+(f\otimes g)$, correspondiente al homomorfismo de anillos $\omega:T_{(f\otimes g)}\to R_{(f)}\otimes_AS_{(g)}$ definido por
\[
\label{II.4.3.1.3-fr}
  \omega((x\otimes y)/(f\otimes g))=(x/f)\otimes(y/g)
\tag{4.3.1.3}
\]
para $x\in E$ y $y\in F$.
\end{env}
% Fuente francesa: PMIHES_1961__8__5_0, printed p. 77.

\begin{env}[4.3.2]
\label{II.4.3.2-fr}
Se sigue de \sref{II.4.2.3-fr} que se tiene un isomorfismo canónico
\[
\label{II.4.3.2.1-fr}
  \tau:\varsigma^*(\mathcal{O}_P(1))\xrightarrow{\sim}
  \mathcal{O}_{P_1}(1)\otimes_Y\mathcal{O}_{P_2}(1)
\tag{4.3.2.1}
\]
donde se pone $P=\mathbf{P}(\mathcal{E}\otimes_{\mathcal{O}_Y}\mathcal{F})$.
Mostraremos que, para $x\in\Gamma(Y,\mathcal{E})$ e $y\in\Gamma(Y,\mathcal{F})$, se tiene
\[
\label{II.4.3.2.2-fr}
  \tau(\alpha_1(x\otimes y))=\alpha_1(x)\otimes\alpha_1(y).
\tag{4.3.2.2}
\]

En efecto, se puede restringir al caso en que $Y$ es afín, y entonces, con la notación de \sref{II.4.3.1-fr} y \sref{II.2.6.2-fr}, se tiene $\alpha_1^{f\otimes g}(x\otimes y)=(x\otimes y)/1$ en $(T(1))_{(f\otimes g)}$, $\alpha_1^f(x)=x/1$ en $(R(1))_{(f)}$ y $\alpha_1^g(y)=y/1$ en $(S(1))_{(g)}$.
La definición de $\tau$ dada en \sref{II.4.2.3-fr} y el cálculo de $v_1$ hecho en \sref{II.4.3.1-fr} demuestran inmediatamente la afirmación \sref{II.4.3.2.2-fr}.
De aquí se obtiene la igualdad
\[
\label{II.4.3.2.3-fr}
  \varsigma^{-1}(P_{x\otimes y})=(P_1)_x\times_Y(P_2)_y
\tag{4.3.2.3}
\]
con la notación de \sref{II.3.1.4-fr}.
En efecto, teniendo en cuenta \sref{II.3.3.3-fr}, la igualdad \sref{II.4.3.2.2-fr} (reduciéndose al caso afín, con ayuda de \sref[I]{I.3.2.7-fr} y \sref[I]{I.3.2.3-fr}) reduce a demostrar el lema siguiente:

\begin{lemma}[4.3.2.4]
\label{II.4.3.2.4-fr}
Sean $B$ y $B'$ $A$-álgebras, y sean $Y=\operatorname{Spec}(A)$, $Z=\operatorname{Spec}(B)$ y $Z'=\operatorname{Spec}(B')$;
cualesquiera que sean $t\in B$, $t'\in B'$, se tiene $D(t\otimes t')=D(t)\times_YD(t')$.
\end{lemma}

\begin{proof}
En efecto, si $p:Z\times_YZ'\to Z$ y $p':Z\times_YZ'\to Z'$ son las proyecciones canónicas, se sigue de \sref[I]{I.1.2.2.2-fr} que $p^{-1}(D(t))=D(t\otimes1)$ y $p'^{-1}(D(t'))=D(1\otimes t')$;
la conclusión resulta de \sref[I]{I.3.2.7-fr} y \sref[I]{I.1.1.9.1-fr}, pues $(t\otimes1)(1\otimes t')=t\otimes t'$.
\end{proof}
\end{env}
% Fuente francesa: PMIHES_1961__8__5_0, printed pp. 77-78.

\begin{proposition}[4.3.3]
\label{II.4.3.3-fr}
El morfismo de Segre es una inmersión cerrada.
\end{proposition}

\begin{proof}
Como la cuestión es local sobre $Y$, se puede restringir al caso en que $Y$ es afín.
Con la notación de \sref{II.4.3.1-fr} y \sref{II.4.3.2-fr}, los $D_+(f\otimes g)$ forman entonces una base de la topología de $P$, pues los $f\otimes g$ generan $T$ cuando $f$ recorre $E$ y $g$ recorre $F$.
Por \sref{II.4.3.2.3-fr}, se sabe también que $\varsigma^{-1}(D_+(f\otimes g))=D_+(f)\times_YD_+(g)$.
Basta, pues, por \sref[I]{I.4.2.4-fr}, demostrar que la restricción de $\varsigma$ a $D_+(f)\times_YD_+(g)$ es una inmersión cerrada en $D_+(f\otimes g)$.
Pero, con la misma notación, la igualdad \sref{II.4.3.1.3-fr} muestra que $\omega$ es \emph{sobreyectivo}, lo que termina la demostración.
\end{proof}

\begin{env}[4.3.4]
\label{II.4.3.4-fr}
El morfismo de Segre es \emph{funtorial} en $\mathcal{E}$ y $\mathcal{F}$, si solo se considera
\oldpage[II]{78}
homomorfismos \emph{sobreyectivos} de $\mathcal{O}_Y$-módulos cuasicoherentes.
En efecto, hay que demostrar que, si $\mathcal{E}\to\mathcal{E}'$ es un $\mathcal{O}_Y$-homomorfismo sobreyectivo, el diagrama
\[
  \xymatrix{
    \mathbf{P}(\mathcal{E}')\times\mathbf{P}(\mathcal{F})
      \ar[r]^{j\times1} \ar[d]_{\varsigma}
    & \mathbf{P}(\mathcal{E})\times\mathbf{P}(\mathcal{F})
      \ar[d]^{\varsigma}
    \\\mathbf{P}(\mathcal{E}'\otimes\mathcal{F}) \ar[r]
    & \mathbf{P}(\mathcal{E}\otimes\mathcal{F})
  }
\]
es conmutativo, donde $j$ denota la inmersión cerrada canónica $\mathbf{P}(\mathcal{E}')\to\mathbf{P}(\mathcal{E})$.
Pongamos $P_1'=\mathbf{P}(\mathcal{E}')$ y conservemos la notación de \sref{II.4.3.1-fr};
entonces $j\times1$ es una inmersión cerrada \sref[I]{I.4.3.1-fr} y, salvo isomorfismo,
\[
  (j\times1)^*(\mathcal{O}_{P_1}(1)\otimes\mathcal{O}_{P_2}(1))
  =j^*(\mathcal{O}_{P_1}(1))\otimes\mathcal{O}_{P_2}(1)
  =\mathcal{O}_{P_1'}(1)\otimes\mathcal{O}_{P_2}(1)
\]
por \sref{II.4.1.2.1-fr} y \sref[I]{I.9.1.5-fr};
nuestra afirmación se sigue entonces de \sref{II.4.2.8-fr} y \sref{II.4.2.9-fr}.
\end{env}
% Fuente francesa: PMIHES_1961__8__5_0, printed p. 78.

\begin{env}[4.3.5]
\label{II.4.3.5-fr}
Con la notación de \sref{II.4.3.1-fr}, sea $\psi:Y'\to Y$ un morfismo, y pongamos $\mathcal{E}'=\psi^*(\mathcal{E})$ y $\mathcal{F}'=\psi^*(\mathcal{F})$;
entonces el morfismo de Segre $\mathbf{P}(\mathcal{E}')\times\mathbf{P}(\mathcal{F}')
\to\mathbf{P}(\mathcal{E}'\otimes\mathcal{F}')$ puede identificarse con $\varsigma_{(Y')}$.
En efecto, conservando la notación de \sref{II.4.3.1-fr}, pongamos $P_1'=\mathbf{P}(\mathcal{E}')$ y $P_2'=\mathbf{P}(\mathcal{F}')$;
se sabe por \sref{II.4.1.3.1-fr} que $P_i'$ puede identificarse con $(P_i)_{(Y')}$ ($i=1,2$), y por tanto el morfismo estructural $P_1'\times P_2'\to Y'$ puede identificarse con $r_{(Y')}$.
Además, $\mathcal{E}'\otimes\mathcal{F}'$ puede identificarse con $\psi^*(\mathcal{E}\otimes\mathcal{F})$, y por tanto $\mathbf{P}(\mathcal{E}'\otimes\mathcal{F}')$ puede identificarse con $(\mathbf{P}(\mathcal{E}\otimes\mathcal{F}))_{(Y')}$ \sref{II.4.1.3.1-fr}.
Finalmente, $\mathcal{O}_{P_1'}(1)\otimes_{Y'}\mathcal{O}_{P_2'}(1)=\mathcal{L}'$ puede identificarse con $\mathcal{L}\otimes_Y\mathcal{O}_{Y'}$ por \sref{II.4.1.3.2-fr} y \sref[I]{I.9.1.11-fr}.
El homomorfismo canónico $(r_{(Y')})^*(\mathcal{E}'\otimes\mathcal{F}')\to\mathcal{L}'$ puede identificarse entonces con $s_{(Y')}$, y nuestra afirmación se sigue de \sref{II.4.2.10-fr}.
\end{env}

\begin{remark}[4.3.6]
\label{II.4.3.6-fr}
El preesquema \emph{suma} de $\mathbf{P}(\mathcal{E})$ y $\mathbf{P}(\mathcal{F})$ es asimismo canónicamente isomorfo a un \emph{subpreesquema cerrado de $\mathbf{P}(\mathcal{E}\oplus\mathcal{F})$}.
En efecto, los homomorfismos sobreyectivos $\mathcal{E}\oplus\mathcal{F}\to\mathcal{E}$ y $\mathcal{E}\oplus\mathcal{F}\to\mathcal{F}$ corresponden a inmersiones cerradas $\mathbf{P}(\mathcal{E})\to\mathbf{P}(\mathcal{E}\oplus\mathcal{F})$ y $\mathbf{P}(\mathcal{F})\to\mathbf{P}(\mathcal{E}\oplus\mathcal{F})$;
todo se reduce entonces a demostrar que los espacios subyacentes de los subpreesquemas cerrados de $\mathbf{P}(\mathcal{E}\oplus\mathcal{F})$ así obtenidos tienen intersección vacía.
Como la cuestión es local sobre $Y$, se puede adoptar la notación de \sref{II.4.3.1-fr};
pero $\mathbf{S}_n(E)$ y $\mathbf{S}_n(F)$ pueden identificarse con submódulos de $\mathbf{S}_n(E\oplus F)$ cuya intersección se reduce a $0$;
si $\mathfrak{p}$ es un ideal primo graduado de $\mathbf{S}(E)$ tal que $\mathfrak{p}\cap\mathbf{S}_n(E)\neq\mathbf{S}_n(E)$ para todo $n\geq0$, existe un ideal primo graduado correspondiente de $\mathbf{S}(E\oplus F)$ cuya intersección con $\mathbf{S}_n(E)$ es $\mathfrak{p}\cap\mathbf{S}_n(E)$, pero que también \emph{contiene} $\mathbf{S}_+(F)$, como se ve inmediatamente;
por tanto, dos puntos de $\operatorname{Proj}(\mathbf{S}(E))$ y $\operatorname{Proj}(\mathbf{S}(F))$, respectivamente, no pueden tener la misma imagen en $\operatorname{Proj}(\mathbf{S}(E\oplus F))$.
\end{remark}
% Fuente francesa: PMIHES_1961__8__5_0, printed pp. 78-79.

\subsection{Inmersiones en fibrados proyectivos; haces muy amplios}
\label{subsection:II.4.4-fr}

\begin{proposition}[4.4.1]
\label{II.4.4.1-fr}
Sea $Y$ un esquema cuasicompacto, o un preesquema cuyo espacio subyacente sea noetheriano, sea $q:X\to Y$ un morfismo \emph{de tipo finito}, y sea $\mathcal{L}$ un $\mathcal{O}_X$-módulo invertible.
\begin{enumerate}
  \item[\rm{(i)}] Sea $\mathcal{S}$ una $\mathcal{O}_Y$-álgebra graduada cuasicoherente en grados positivos, y sea $\psi:q^*(\mathcal{S})\to\bigoplus_{n\geq0}\mathcal{L}^{\otimes n}$ un homomorfismo de álgebras graduadas.
    Para que $r_{\mathcal{L},\psi}$ esté definido en todas partes y sea una inmersión, es necesario y
\oldpage[II]{79}
    suficiente que exista un entero $n>0$ y un sub-$\mathcal{O}_Y$-módulo cuasicoherente \emph{de tipo finito} $\mathcal{E}$ de $\mathcal{S}_n$ tales que el homomorfismo $\psi'=\psi_n\circ q^*(j):q^*(\mathcal{E})
  \to\mathcal{L}^{\otimes n}=\mathcal{L}'$ (donde $j$ es la inyección $\mathcal{E}\to\mathcal{S}_n$) sea sobreyectivo y el morfismo $r_{\mathcal{L}',\psi'}:X\to\mathbf{P}(\mathcal{E})$ sea una inmersión.
  \item[\rm{(ii)}] Sea $\mathcal{F}$ un $\mathcal{O}_Y$-módulo cuasicoherente, y sea $\varphi:q^*(\mathcal{F})\to\mathcal{L}$ un homomorfismo sobreyectivo.
    Para que el morfismo $r_{\mathcal{L},\varphi}$ sea una inmersión $X\to\mathbf{P}(\mathcal{F})$, es necesario y suficiente que exista un sub-$\mathcal{O}_Y$-módulo cuasicoherente \emph{de tipo finito} $\mathcal{E}$ de $\mathcal{F}$ tal que el homomorfismo $\varphi'=\varphi\circ q^*(j):q^*(\mathcal{E})\to\mathcal{L}$ (donde $j:\mathcal{E}\to\mathcal{F}$ es la inyección canónica) sea sobreyectivo y el morfismo $r_{\mathcal{L},\varphi'}:X\to\mathbf{P}(\mathcal{E})$ sea una inmersión.
\end{enumerate}
\end{proposition}

\begin{proof}
\medskip\noindent
\begin{enumerate}
  \item[\rm{(i)}] El hecho de que $r_{\mathcal{L},\psi}$ esté definido en todas partes y sea una inmersión equivale, por \sref{II.3.8.5-fr}, a la existencia de un $n>0$ y de un $\mathcal{E}$ tales que, si $\mathcal{S}'$ es la subálgebra de $\mathcal{S}$ generada por $\mathcal{E}$, el homomorfismo $q^*(\mathcal{E})\to\mathcal{L}^{\otimes n}$ sea sobreyectivo y el morfismo $r_{\mathcal{L},\psi'}:X\to\operatorname{Proj}(\mathcal{S}')$ esté definido en todas partes y sea una inmersión.
    Por otra parte, se tiene un homomorfismo canónico sobreyectivo $\mathbf{S}(\mathcal{E})\to\mathcal{S}'$, al cual corresponde una inmersión cerrada $\operatorname{Proj}(\mathcal{S}')\to\mathbf{P}(\mathcal{E})$ \sref{II.3.6.2-fr};
    de donde resulta la conclusión.
  \item[\rm{(ii)}] Como $\mathcal{F}$ es el límite inductivo de sus submódulos cuasicoherentes de tipo finito $\mathcal{E}_\lambda$ \sref[I]{I.9.4.9-fr}, $\mathbf{S}(\mathcal{F})$ es el límite inductivo de los $\mathbf{S}(\mathcal{E}_\lambda)$;
    la conclusión resulta entonces de \sref{II.3.8.4-fr}, observando que se puede tomar todos los $n_i$ de la demostración de \sref{II.3.8.4-fr} iguales a $1$:
    en efecto, suponiendo $Y$ afín, si $r=r_{\mathcal{L},\varphi}$ es una inmersión, entonces $r(X)$ es un subespacio cuasicompacto de $\mathbf{P}(\mathcal{F})$ que puede cubrirse por un número finito de abiertos de $\mathbf{P}(\mathcal{F})$ de la forma $D_+(f)$, con $f\in F$, tales que $D_+(f)\cap r(X)$ sea cerrado.
\end{enumerate}
\end{proof}
% Fuente francesa: PMIHES_1961__8__5_0, printed p. 79.

\begin{definition}[4.4.2]
\label{II.4.4.2-fr}
Sean $Y$ un preesquema y $q:X\to Y$ un morfismo.
Se dice que un $\mathcal{O}_X$-módulo invertible $\mathcal{L}$ es \emph{muy amplio} relativamente a $q$ \textup{(o} muy amplio relativamente a $Y$\textup{, o simplemente} muy amplio\textup{, si no resulta confusión alguna)} si existe un $\mathcal{O}_Y$-módulo cuasicoherente $\mathcal{E}$ y una $Y$-inmersión $i$ de $X$ en $P=\mathbf{P}(\mathcal{E})$ tales que $\mathcal{L}$ sea isomorfo a $i^*(\mathcal{O}_P(1))$.
\end{definition}

Es equivalente \sref{II.4.2.3-fr} decir que existe un $\mathcal{O}_Y$-módulo cuasicoherente $\mathcal{E}$ y un homomorfismo \emph{sobreyectivo} $\varphi:q^*(\mathcal{E})\to\mathcal{L}$ tal que $r_{\mathcal{L},\varphi}:X\to\mathbf{P}(\mathcal{E})$ sea una \emph{inmersión}.

Se observa que la existencia de un $\mathcal{O}_X$-módulo muy amplio (para $Y$) implica que $q$ es \emph{separado} (\sref{II.3.1.3-fr} y \sref[I]{I.5.5.1-fr}[(i) y (ii)]).

\begin{corollary}[4.4.3]
\label{II.4.4.3-fr}
Supongamos que existe una $\mathcal{O}_Y$-álgebra graduada cuasicoherente $\mathcal{S}$, generada por $\mathcal{S}_1$, y una $Y$-inmersión $i:X\to P=\operatorname{Proj}(\mathcal{S})$ tal que $\mathcal{L}$ sea isomorfo a $i^*(\mathcal{O}_P(1))$;
entonces $\mathcal{L}$ es muy amplio relativamente a $q$.
\end{corollary}

\begin{proof}
Si $\mathcal{F}=\mathcal{S}_1$, el homomorfismo canónico $\mathbf{S}(\mathcal{F})\to\mathcal{S}$ es sobreyectivo, y por tanto, componiendo la inmersión cerrada correspondiente $\operatorname{Proj}(\mathcal{S})\to\mathbf{P}(\mathcal{F})$ \sref{II.3.6.2-fr} con la inmersión $i$, se obtiene una inmersión $j:X\to\mathbf{P}(\mathcal{F})=P'$ tal que $\mathcal{L}$ es isomorfo a $j^*(\mathcal{O}_{P'}(1))$ \sref{II.3.6.3-fr}.
\end{proof}
% Fuente francesa: PMIHES_1961__8__5_0, printed pp. 79-80.

\begin{proposition}[4.4.4]
\label{II.4.4.4-fr}
Sean $q:X\to Y$ un morfismo cuasicompacto y $\mathcal{L}$ un $\mathcal{O}_X$-módulo invertible.
Las propiedades siguientes son equivalentes:
\begin{enumerate}
  \item[\rm{(a)}] $\mathcal{L}$ es muy amplio relativamente a $q$.
  \item[\rm{(b)}] $q_*(\mathcal{L})$ es cuasicoherente, el homomorfismo canónico $\sigma:q^*(q_*(\mathcal{L}))\to\mathcal{L}$ es sobreyectivo, y el morfismo $r_{\mathcal{L},\sigma}:X\to\mathbf{P}(q_*(\mathcal{L}))$ es una inmersión.
\end{enumerate}
\end{proposition}

\begin{proof}
Como $q$ es cuasicompacto, se sabe que $q_*(\mathcal{L})$ es cuasicoherente si $q$ es separado \sref[I]{I.9.2.2-fr}.

\oldpage[II]{80}
Se sabe por \sref{II.3.4.7-fr} que la existencia de un homomorfismo sobreyectivo $\varphi:q^*(\mathcal{E})\to\mathcal{L}$ (siendo $\mathcal{E}$ un $\mathcal{O}_Y$-módulo cuasicoherente) implica que $\sigma$ es sobreyectivo;
además, a la factorización $q^*(\mathcal{E})\to q^*(q_*(\mathcal{L}))\xrightarrow{\sigma}\mathcal{L}$ de $\varphi$ corresponde canónicamente una factorización
\[
  q^*(\mathbf{S}(\mathcal{E}))
    \longrightarrow q^*(\mathbf{S}(q_*(\mathcal{L})))
    \longrightarrow\bigoplus_{n\geq0}\mathcal{L}^{\otimes n}
\]
y por tanto, por \sref{II.3.8.3-fr}, la hipótesis de que $r_{\mathcal{L},\varphi}$ es una inmersión implica que también lo es $j=r_{\mathcal{L},\sigma}$;
además, por \sref{II.4.2.4-fr}, $\mathcal{L}$ es isomorfo a $j^*(\mathcal{O}_{P'}(1))$, donde $P'=\mathbf{P}(q_*(\mathcal{L}))$.
Se ve, pues, que \emph{(a)} y \emph{(b)} son equivalentes.
\end{proof}

\begin{corollary}[4.4.5]
\label{II.4.4.5-fr}
Supongamos que $q$ es cuasicompacto.
Para que $\mathcal{L}$ sea muy amplio relativamente a $Y$, es necesario y suficiente que exista una cubierta abierta $(U_\alpha)$ de $Y$ tal que $\mathcal{L}|q^{-1}(U_\alpha)$ sea muy amplio relativamente a $U_\alpha$ para todo $\alpha$.
\end{corollary}

\begin{proof}
En efecto, la condición \emph{(b)} de \sref{II.4.4.4-fr} es local sobre $Y$.
\end{proof}
% Fuente francesa: PMIHES_1961__8__5_0, printed p. 80.

\begin{proposition}[4.4.6]
\label{II.4.4.6-fr}
Sea $Y$ un esquema cuasicompacto, o un preesquema cuyo espacio subyacente sea noetheriano, sea $q:X\to Y$ un morfismo \emph{de tipo finito}, y sea $\mathcal{L}$ un $\mathcal{O}_X$-módulo invertible.
Entonces las condiciones \textup{(a)} y \textup{(b)} de \sref{II.4.4.4-fr} son equivalentes a las siguientes:
\begin{enumerate}
  \item[\rm{(a')}] Existe un $\mathcal{O}_Y$-módulo cuasicoherente $\mathcal{E}$ \emph{de tipo finito} y un homomorfismo sobreyectivo $\varphi:q^*(\mathcal{E})\to\mathcal{L}$ tal que $r_{\mathcal{L},\varphi}$ sea una inmersión.
  \item[\rm{(b')}] Existe un sub-$\mathcal{O}_Y$-módulo cuasicoherente $\mathcal{E}$ de $q_*(\mathcal{L})$, \emph{de tipo finito}, que posee las propiedades enunciadas en la condición~\textup{(a')}.
\end{enumerate}
\end{proposition}

\begin{proof}
Es claro que \emph{(a')} o \emph{(b')} implica \emph{(a)};
por otra parte, \emph{(a)} implica \emph{(a')} por \sref{II.4.4.1-fr}, y del mismo modo \emph{(b)} implica \emph{(b')}.
\end{proof}

\begin{corollary}[4.4.7]
\label{II.4.4.7-fr}
Supongamos que $Y$ es un esquema cuasicompacto o un preesquema noetheriano.
Si $\mathcal{L}$ es muy amplio para $q$, existe una $\mathcal{O}_Y$-álgebra graduada cuasicoherente $\mathcal{S}$ tal que $\mathcal{S}_1$ es de tipo finito y genera $\mathcal{S}$, y una $Y$-inmersión \emph{abierta dominante} $i:X\to P=\operatorname{Proj}(\mathcal{S})$ tal que $\mathcal{L}$ es isomorfo a $i^*(\mathcal{O}_P(1))$.
\end{corollary}

\begin{proof}
En efecto, se satisface la condición~(\emph{b')} de \sref{II.4.4.6-fr};
el morfismo estructural $p:\mathbf{P}(\mathcal{E})=P'\to Y$ es entonces separado y de tipo finito \sref{II.3.1.3-fr}, y por tanto $P'$ es un esquema cuasicompacto (resp. un preesquema noetheriano) si $Y$ es un esquema cuasicompacto (resp. un preesquema noetheriano).
Sea $Z$ la clausura \sref[I]{I.9.5.10-fr} del subpreesquema $X'$ de $P'$ asociado a la inmersión $j=r_{\mathcal{L},\varphi}$ de $X$ en $P'$;
es claro que $j$ se factoriza en la inyección canónica $Z\to P'$ y en una inmersión abierta dominante $i:X\to Z$.
Pero $Z$ puede identificarse con un preesquema $\operatorname{Proj}(\mathcal{S})$, donde $\mathcal{S}$ es una $\mathcal{O}_Y$-álgebra graduada cociente de $\mathbf{S}(\mathcal{E})$ por un haz graduado cuasicoherente de ideales \sref{II.3.6.2-fr}, y es claro que $\mathcal{S}_1$ es de tipo finito y genera $\mathcal{S}$;
además, $\mathcal{O}_Z(1)$ es la imagen inversa de $\mathcal{O}_{P'}(1)$ por la inyección canónica \sref{II.3.6.3-fr}, y por tanto $\mathcal{L}=i^*(\mathcal{O}_Z(1))$.
\end{proof}
% Fuente francesa: PMIHES_1961__8__5_0, printed pp. 80-81.

\begin{proposition}[4.4.8]
\label{II.4.4.8-fr}
Sean $q:X\to Y$ un morfismo, $\mathcal{L}$ un $\mathcal{O}_X$-módulo muy amplio relativamente a $q$, y $\mathcal{L}'$ un $\mathcal{O}_X$-módulo invertible tales que existe un $\mathcal{O}_Y$-módulo cuasicoherente $\mathcal{E}'$ y un homomorfismo sobreyectivo $q^*(\mathcal{E}')\to\mathcal{L}'$.
Entonces $\mathcal{L}\otimes_{\mathcal{O}_X}\mathcal{L}'$ es muy amplio relativamente a $q$.
\end{proposition}

\begin{proof}
La hipótesis implica la existencia de un $Y$-morfismo $r':X\to P'=\mathbf{P}(\mathcal{E}')$ tal que $\mathcal{L}'=r'^*(\mathcal{O}_{P'}(1))$ \sref{II.4.2.1-fr}.
Existe, por hipótesis, un $\mathcal{O}_Y$-módulo cuasicoherente $\mathcal{E}$ y una
\oldpage[II]{81}
$Y$-inmersión $r:X\to P=\mathbf{P}(\mathcal{E})$ tal que $\mathcal{L}=r^*(\mathcal{O}_P(1))$.
Pongamos $Q=\mathbf{P}(\mathcal{E}\otimes\mathcal{E}')$ y consideremos el morfismo de Segre $\varsigma:P\times_YP'\to Q$ \sref{II.4.3.1-fr}.
Como $r$ es una inmersión, también lo es $(r,r')_Y:X\to P\times_YP'$ \sref[I]{I.5.3.14-fr};
pero, como $\varsigma$ es una inmersión \sref{II.4.3.3-fr}, también lo es $r'':X\xrightarrow{(r,r')}P\times_YP'\xrightarrow{\varsigma}Q$.
Por otra parte, por \sref{II.4.3.2.1-fr}, $\varsigma^*(\mathcal{O}_Q(1))$ es isomorfo a $\mathcal{O}_P(1)\otimes_Y\mathcal{O}_{P'}(1)$, y por tanto, por \sref[I]{I.9.1.4-fr}, $r''^*(\mathcal{O}_Q(1))$ es isomorfo a $\mathcal{L}\otimes\mathcal{L}'$, lo que demuestra la proposición.
\end{proof}

\begin{corollary}[4.4.9]
\label{II.4.4.9-fr}
Sea $q:X\to Y$ un morfismo.
\begin{enumerate}
  \item Sea $\mathcal{L}$ un $\mathcal{O}_X$-módulo invertible y $\mathcal{K}$ un $\mathcal{O}_Y$-módulo invertible.
    Para que $\mathcal{L}$ sea muy amplio relativamente a $q$, es necesario y suficiente que $\mathcal{L}\otimes q^*(\mathcal{K})$ lo sea.
  \item Si $\mathcal{L}$ y $\mathcal{L}'$ son dos $\mathcal{O}_X$-módulos muy amplios relativamente a $q$, también lo es $\mathcal{L}\otimes\mathcal{L}'$;
    en particular, $\mathcal{L}^{\otimes n}$ es muy amplio relativamente a $q$ para todo $n>0$.
\end{enumerate}
\end{corollary}

\begin{proof}
La afirmación~(ii) es consecuencia inmediata de \sref{II.4.4.8-fr}, al igual que la necesidad de la condición~(i);
recíprocamente, si $\mathcal{L}\otimes q^*(\mathcal{K})$ es muy amplio, también lo es $(\mathcal{L}\otimes q^*(\mathcal{K}))\otimes q^*(\mathcal{K}^{-1})$ por lo anterior, y este último $\mathcal{O}_X$-módulo es isomorfo a $\mathcal{L}$ (\sref[0]{0.5.4.3-fr} y \sref[0]{0.5.4.5-fr}).
\end{proof}
% Fuente francesa: PMIHES_1961__8__5_0, printed p. 81.

\begin{proposition}[4.4.10]
\label{II.4.4.10-fr}
\medskip\noindent
\begin{enumerate}
  \item[\rm{(i)}] Para todo preesquema $Y$, todo $\mathcal{O}_Y$-módulo invertible $\mathcal{L}$ es muy amplio relativamente al morfismo identidad $1_Y$.
  \item[\rm{(i \emph{bis})}] Sean $f:X\to Y$ un morfismo y $j:X'\to X$ una inmersión.
    Si $\mathcal{L}$ es un $\mathcal{O}_X$-módulo muy amplio relativamente a $f$, entonces $j^*(\mathcal{L})$ es muy amplio relativamente a $f\circ j$.
  \item[\rm{(ii)}] Sean $Z$ un preesquema cuasicompacto, $f:X\to Y$ un morfismo de tipo finito, $g:Y\to Z$ un morfismo cuasicompacto, $\mathcal{L}$ un $\mathcal{O}_X$-módulo muy amplio relativamente a $f$, y $\mathcal{K}$ un $\mathcal{O}_Y$-módulo muy amplio relativamente a $g$.
    Entonces existe un entero $n_0>0$ tal que $\mathcal{L}\otimes f^*(\mathcal{K}^{\otimes n})$ es muy amplio relativamente a $g\circ f$ para todo $n\geq n_0$.
  \item[\rm{(iii)}] Sean $f:X\to Y$ y $g:Y'\to Y$ dos morfismos, y pongamos $X'=X_{(Y')}$.
    Si $\mathcal{L}$ es un $\mathcal{O}_X$-módulo muy amplio relativamente a $f$, entonces $\mathcal{L}'=\mathcal{L}\otimes_Y\mathcal{O}_{Y'}$ es un $\mathcal{O}_{X'}$-módulo muy amplio relativamente a $f_{(Y')}$.
  \item[\rm{(iv)}] Sean $f_i:X_i\to Y_i$ ($i=1,2$) dos $S$-morfismos.
    Si $\mathcal{L}_i$ es un $\mathcal{O}_{X_i}$-módulo muy amplio relativamente a $f_i$ ($i=1,2$), entonces $\mathcal{L}_1\otimes_S\mathcal{L}_2$ es muy amplio relativamente a $f_1\times_Sf_2$.
  \item[\rm{(v)}] Sean $f:X\to Y$ y $g:Y\to Z$ dos morfismos.
    Si un $\mathcal{O}_X$-módulo $\mathcal{L}$ es muy amplio relativamente a $g\circ f$, entonces $\mathcal{L}$ es muy amplio relativamente a $f$.
  \item[\rm{(vi)}] Sean $f:X\to Y$ un morfismo y $j$ la inyección canónica $X_{\mathrm{red}}\to X$.
    Si $\mathcal{L}$ es un $\mathcal{O}_X$-módulo muy amplio relativamente a $f$, entonces $j^*(\mathcal{L})$ es muy amplio relativamente a $f_{\mathrm{red}}$.
\end{enumerate}
\end{proposition}
% Fuente francesa: PMIHES_1961__8__5_0, printed pp. 81-82.

\begin{proof}
La propiedad~(i~\emph{bis}) se sigue inmediatamente de la definición \sref{II.4.4.2-fr}, y es inmediato que (vi) se deduce formalmente de (i~\emph{bis}) y (v), mediante un razonamiento calcado del de la demostración de \sref[I]{I.5.5.12-fr}, que dejamos al lector.
Para demostrar (v), consideremos, como en \sref[I]{I.5.5.12-fr}, la factorización $X\xrightarrow{\Gamma_f}X\times_ZY\xrightarrow{p_2}Y$, donde $p_2=(g\circ f)\times1_Y$.
Se sigue de la hipótesis y de (i) y (iv) que $\mathcal{L}\otimes_{\mathcal{O}_Z}\mathcal{O}_Y$ es muy amplio para $p_2$;
por otra parte, $\mathcal{L}=\Gamma_f^*(\mathcal{L}\otimes_{\mathcal{O}_Z}\mathcal{O}_Y)$ \sref[I]{I.9.1.4-fr}, y $\Gamma_f$ es una inmersión \sref[I]{I.5.3.11-fr};
se puede, pues, aplicar (i~\emph{bis}).

\oldpage[II]{82}
Para demostrar (i), se aplica la definición \sref{II.4.4.2-fr} con $\mathcal{E}=\mathcal{L}$ observando que entonces $\mathbf{P}(\mathcal{E})$ puede identificarse con $Y$ \sref{II.4.1.4-fr}.

Demostremos ahora (iii).
Existe un $\mathcal{O}_Y$-módulo cuasicoherente $\mathcal{E}$ y una $Y$-inmersión $i:X\to\mathbf{P}(\mathcal{E})=P$ tales que $\mathcal{L}=i^*(\mathcal{O}_P(1))$;
si se pone $\mathcal{E}'=g^*(\mathcal{E})$, entonces $\mathcal{E}'$ es un $\mathcal{O}_{Y'}$-módulo cuasicoherente, y se tiene $P'=\mathbf{P}(\mathcal{E}')=P_{(Y')}$ \sref{II.4.1.3.1-fr}, que $i_{(Y')}$ es una inmersión de $X_{(Y')}$ en $P'$ \sref[I]{I.4.3.2-fr}, y que $\mathcal{L}'$ es isomorfo a $(i_{(Y')})^*(\mathcal{O}_{P'}(1))$ \sref{II.4.2.10-fr}.

Para demostrar (iv), observemos que existe, por hipótesis, una $Y_i$-inmersión $r_i:X_i\to P_i=\mathbf{P}(\mathcal{E}_i)$, donde $\mathcal{E}_i$ es un $\mathcal{O}_{Y_i}$-módulo cuasicoherente, y $\mathcal{L}_i=r_i^*(\mathcal{O}_{P_i}(1))$ ($i=1,2$);
$r_1\times_Sr_2$ es una $S$-inmersión de $X_1\times_SX_2$ en $P_1\times_SP_2$ \sref[I]{I.4.3.1-fr}, y la imagen inversa de $\mathcal{O}_{P_1}(1)\otimes_S\mathcal{O}_{P_2}(1)$ por esta inmersión es $\mathcal{L}_1\otimes_S\mathcal{L}_2$.
Pongamos ahora $T=Y_1\times_SY_2$, y sean $p_1,p_2$ las proyecciones respectivas de $T$ sobre $Y_1,Y_2$.
Si se pone $P_i'=\mathbf{P}(p_i^*(\mathcal{E}_i))$ ($i=1,2$), entonces $P_i'=P_i\times_{Y_i}T$, por \sref{II.4.1.3.1-fr}, y por tanto
\[
  P_1'\times_TP_2'
  =(P_1\times_{Y_1}T)\times_T(P_2\times_{Y_2}T)
  =P_1\times_{Y_1}(T\times_{Y_2}P_2)
  =P_1\times_{Y_1}(Y_1\times_SP_2)
  =P_1\times_SP_2
\]
salvo isomorfismo canónico.
Asimismo, $\mathcal{O}_{P_i'}(1)=\mathcal{O}_{P_i}(1)\otimes_{Y_i}\mathcal{O}_T$ \sref{II.4.1.3.2-fr}, y un cálculo análogo (basado en particular en \sref[I]{I.9.1.9.1-fr} y \sref[I]{I.9.1.2-fr}) muestra que, en la identificación precedente, $\mathcal{O}_{P_1'}(1)\otimes_T\mathcal{O}_{P_2'}(1)$ puede identificarse con $\mathcal{O}_{P_1}(1)\otimes_S\mathcal{O}_{P_2}(1)$.
Se puede, pues, considerar $r_1\times_Sr_2$ como una $T$-inmersión de $X_1\times_SX_2$ en $P_1'\times_TP_2'$, cuya imagen inversa de $\mathcal{O}_{P_1'}(1)\otimes_T\mathcal{O}_{P_2'}(1)$ es $\mathcal{L}_1\otimes_S\mathcal{L}_2$.
Se termina entonces el razonamiento como en \sref{II.4.4.8-fr}, utilizando el morfismo de Segre.
% Fuente francesa: PMIHES_1961__8__5_0, printed p. 82.

Solo queda demostrar (ii).
Se puede restringir primero al caso en que $Z$ es un esquema afín, pues, en general, existe una cubierta finita $(U_i)$ de $Z$ por abiertos afines;
si la proposición estuviera demostrada para $\mathcal{K}|g^{-1}(U_i)$, $\mathcal{L}|f^{-1}(g^{-1}(U_i))$ y un entero $n_i$, bastaría tomar por $n_0$ el mayor de los $n_i$ para demostrar la proposición para $\mathcal{K}$ y $\mathcal{L}$ \sref{II.4.4.5-fr}.
La hipótesis implica que $f$ y $g$ son morfismos separados, y por tanto $X$ e $Y$ son \emph{esquemas} cuasicompactos.

Por \sref{II.4.4.6-fr}, existe una inmersión $r:X\to P=\mathbf{P}(\mathcal{E})$, donde $\mathcal{E}$ es un $\mathcal{O}_Y$-módulo cuasicoherente \emph{de tipo finito}, y $\mathcal{L}=r^*(\mathcal{O}_P(1))$.
Veremos que existe un $\mathcal{O}_P$-módulo $\mathcal{M}$ \emph{muy amplio} relativamente al morfismo compuesto $P\xrightarrow{h}Y\xrightarrow{g}Z$, tal que $\mathcal{O}_P(1)$ es isomorfo a $\mathcal{M}\otimes_Y\mathcal{K}^{\otimes(-m)}$ para algún entero $m$.
Para $n\geq m+1$, $\mathcal{O}_P(1)\otimes_Y\mathcal{K}^{\otimes n}$ será entonces muy amplio para $Z$, por hipótesis y por (iv) aplicado a los morfismos $h:P\to Y$ y $1_Y$;
como $r$ es una inmersión y $\mathcal{L}\otimes f^*(\mathcal{K}^{\otimes n})
=r^*(\mathcal{O}_P(1)\otimes_Y\mathcal{K}^{\otimes n})$, nuestra conclusión se seguirá entonces de (i~\emph{bis}).
Para demostrar nuestra afirmación sobre $\mathcal{O}_P(1)$, utilizaremos el lema siguiente:

\begin{lemma}[4.4.10.1]
\label{II.4.4.10.1-fr}
Sea $Z$ un esquema cuasicompacto, o un preesquema cuyo espacio subyacente sea noetheriano, y sean $g:Y\to Z$ un morfismo cuasicompacto, $\mathcal{K}$ un $\mathcal{O}_Y$-módulo invertible muy amplio relativamente a $g$, y $\mathcal{E}$ un $\mathcal{O}_Y$-módulo cuasicoherente de tipo finito.
Entonces existe un entero $m_0$ tal que, para todo $m\geq m_0$, $\mathcal{E}$ es isomorfo a un cociente de un $\mathcal{O}_Y$-módulo de la forma $g^*(\mathcal{F})\otimes\mathcal{K}^{\otimes(-m)}$, donde $\mathcal{F}$ es un $\mathcal{O}_Z$-módulo cuasicoherente de tipo finito (dependiente de $m$).
\end{lemma}

Este lema se demostrará en \sref{II.4.5.10.1-fr};
el lector puede comprobar que \sref{II.4.4.10-fr} no se utiliza en el n\textsuperscript{o}~4.5.
% Fuente francesa: PMIHES_1961__8__5_0, printed p. 83.

\oldpage[II]{83}
Admitiendo este lema, existe una inmersión cerrada $j_1$ de $P$ en
\[
  P_1=\mathbf{P}(g^*(\mathcal{F})\otimes\mathcal{K}^{\otimes(-m)})
\]
tal que $\mathcal{O}_P(1)$ es isomorfo a $j_1^*(\mathcal{O}_{P_1}(1))$ \sref{II.4.1.2-fr}.
Ahora bien, existe un isomorfismo de $P_1$ sobre $P_2=\mathbf{P}(g^*(\mathcal{F}))$ que identifica $\mathcal{O}_{P_1}(1)$ con $\mathcal{O}_{P_2}(1)\otimes_Y\mathcal{K}^{\otimes(-m)}$ \sref{II.4.1.4-fr};
se tiene, pues, una inmersión cerrada $j_2:P\to P_2$ tal que $\mathcal{O}_P(1)$ es isomorfo a $j_2^*(\mathcal{O}_{P_2}(1))\otimes_Y\mathcal{K}^{\otimes(-m)}$.
Finalmente, $P_2$ puede identificarse con $P_3\times_ZY$, donde $P_3=\mathbf{P}(\mathcal{F})$, y $\mathcal{O}_{P_2}(1)$ con $\mathcal{O}_{P_3}(1)\otimes_Z\mathcal{O}_Y$ \sref{II.4.1.3-fr}.
Por definición, $\mathcal{O}_{P_3}(1)$ es muy amplio para $Z$;
como también lo es $\mathcal{K}$, se deduce de (iv) que $\mathcal{O}_{P_2}(1)\otimes_Y\mathcal{K}$ es muy amplio para $Z$;
también lo es $\mathcal{M}=j_2^*(\mathcal{O}_{P_2}(1)\otimes_Y\mathcal{K})$ por (i~\emph{bis}), y $\mathcal{O}_P(1)$ es isomorfo a $\mathcal{M}\otimes_Y\mathcal{K}^{\otimes(-m-1)}$, lo que termina la demostración.
\end{proof}
% Fuente francesa: PMIHES_1961__8__5_0, printed p. 83.

\begin{proposition}[4.4.11]
\label{II.4.4.11-fr}
Sean $f:X\to Y$ y $f':X'\to Y$ dos morfismos, sea $X''$ el preesquema suma $X\amalg X'$, y sea $f''$ el morfismo $X''\to Y$ que coincide con $f$ \textup{(resp. $f'$)} sobre $X$ \textup{(resp. $X'$)}.
Sea $\mathcal{L}$ \textup{(resp. $\mathcal{L}'$)} un $\mathcal{O}_X$-módulo invertible \textup{(resp.} un $\mathcal{O}_{X'}$-módulo invertible\textup{)}, y sea $\mathcal{L}''$ el $\mathcal{O}_{X''}$-módulo invertible que coincide con $\mathcal{L}$ en $X$ y con $\mathcal{L}'$ en $X'$.
Para que $\mathcal{L}''$ sea muy amplio relativamente a $f''$, es necesario y suficiente que $\mathcal{L}$ sea muy amplio relativamente a $f$ y que $\mathcal{L}'$ sea muy amplio relativamente a $f'$.
\end{proposition}

\begin{proof}
Se puede restringir inmediatamente al caso en que $Y$ es afín.
Si $\mathcal{L}''$ es muy amplio, también lo son $\mathcal{L}$ y $\mathcal{L}'$, por \sref{II.4.4.10-fr}[(i~\emph{bis})].
Recíprocamente, si $\mathcal{L}$ y $\mathcal{L}'$ son muy amplios, se sigue inmediatamente de la definición \sref{II.4.4.2-fr} y de \sref{II.4.3.6-fr} que $\mathcal{L}''$ es muy amplio.
\end{proof}
% Fuente francesa: PMIHES_1961__8__5_0, printed p. 83.

\subsection{Haces amplios}
\label{subsection:II.4.5-fr}

\begin{env}[4.5.1]
\label{II.4.5.1-fr}
Dados un preesquema $X$ y un $\mathcal{O}_X$-módulo invertible $\mathcal{L}$, pondremos, para todo $\mathcal{O}_X$-módulo $\mathcal{F}$ (cuando no pueda haber confusión sobre $\mathcal{L}$), $\mathcal{F}(n)=\mathcal{F}\otimes_{\mathcal{O}_X}\mathcal{L}^{\otimes n}$ ($n\in\mathbf{Z}$);
pondremos además $S=\bigoplus_{n\geq0}\Gamma(X,\mathcal{L}^{\otimes n})$ (subanillo graduado del anillo $\Gamma_*(\mathcal{L})$ definido en \sref[0]{0.5.4.6-fr}).
Si se considera $X$ como un $\mathbf{Z}$-preesquema y se denota por $p$ el morfismo estructural $X\to\operatorname{Spec}(\mathbf{Z})$, existe una correspondencia biyectiva entre los homomorfismos $p^*(\widetilde{S})\longrightarrow\bigoplus_{n\geq0}\mathcal{L}^{\otimes n}$ de $\mathcal{O}_X$-álgebras graduadas y los endomorfismos del anillo graduado $S$ \sref[I]{I.2.2.5-fr};
el homomorfismo $\varepsilon:p^*(\widetilde{S})\longrightarrow\bigoplus_{n\geq0}\mathcal{L}^{\otimes n}$ que corresponde así al automorfismo \emph{identidad} de $S$ se llama \emph{canónico}.
Le corresponde \sref{II.3.7.1-fr} un morfismo $G(\varepsilon)\to\operatorname{Proj}(S)$ que también se llama \emph{canónico}.
\end{env}
% Fuente francesa: PMIHES_1961__8__5_0, printed pp. 83-84.

\begin{theorem}[4.5.2]
\label{II.4.5.2-fr}
Sean $X$ un esquema cuasicompacto o un preesquema cuyo espacio subyacente sea noetheriano, $\mathcal{L}$ un $\mathcal{O}_X$-módulo invertible, y $S$ el anillo graduado $\bigoplus_{n\geq0}\Gamma(X,\mathcal{L}^{\otimes n})$.
Las condiciones siguientes son equivalentes:
\begin{enumerate}
  \item[\rm{(a)}] Cuando $f$ recorre el conjunto de los elementos homogéneos de $S_+$, los $X_f$ forman una base de la topología de $X$.
  \item[\rm{(a')}] Cuando $f$ recorre el conjunto de los elementos homogéneos de $S_+$, los $X_f$ que son afines forman una cubierta de $X$.
  \item[\rm{(b)}] El morfismo canónico $G(\varepsilon)\to\operatorname{Proj}(S)$ \sref{II.4.5.1-fr} está definido en todas partes y es una inmersión abierta dominante.
\oldpage[II]{84}
  \item[\rm{(b')}] El morfismo canónico $G(\varepsilon)\to\operatorname{Proj}(S)$ está definido en todas partes y es un homeomorfismo del espacio subyacente a $X$ sobre un subespacio de $\operatorname{Proj}(S)$.
  \item[\rm{(c)}] Para todo $\mathcal{O}_X$-módulo cuasicoherente $\mathcal{F}$, si se denota por $\mathcal{F}_n$ el sub-$\mathcal{O}_X$-módulo de $\mathcal{F}(n)$ generado por las secciones de $\mathcal{F}(n)$ sobre $X$, entonces $\mathcal{F}$ es la suma de los sub-$\mathcal{O}_X$-módulos $\mathcal{F}_n(-n)$ para los enteros $n>0$.
  \item[\rm{(c')}] La propiedad~\rm{(c)} se cumple para todo haz cuasicoherente de ideales de $\mathcal{O}_X$.
\end{enumerate}

Además, si $(f_\alpha)$ es una familia de elementos homogéneos de $S_+$ tal que los $X_{f_\alpha}$ son afines, la restricción a $\bigcup_\alpha X_{f_\alpha}$ del morfismo canónico $X\to\operatorname{Proj}(S)$ es un isomorfismo de $\bigcup_\alpha X_{f_\alpha}$ sobre $\bigcup_\alpha(\operatorname{Proj}(S))_{f_\alpha}$.
\end{theorem}
% Fuente francesa: PMIHES_1961__8__5_0, printed p. 84.

\begin{proof}
Es claro que \emph{(b)} implica \emph{(b')}, y que \emph{(b')} implica \emph{(a)} por \sref{II.3.7.3.1-fr} (teniendo en cuenta que $\varepsilon^\flat$ es la identidad).
La condición~\emph{(a)} implica \emph{(a')}, pues todo $x\in X$ posee una vecindad afín $U$ tal que $\mathcal{L}|U$ es isomorfo a $\mathcal{O}_X|U$;
si $f\in\Gamma(X,\mathcal{L}^{\otimes n})$ es tal que $x\in X_f\subset U$, entonces $X_f$ es también el conjunto de los $x'\in U$ tales que $(f|U)(x')\neq0$, y es por tanto un abierto afín \sref[I]{I.1.3.6-fr}.
Para demostrar que \emph{(a')} implica \emph{(b)}, basta demostrar la última afirmación del teorema y verificar además que, si $X=\bigcup_\alpha X_{f_\alpha}$, se satisface la condición~(iv) de \sref{II.3.8.2-fr}.
Este último punto se sigue inmediatamente de \sref[I]{I.9.3.1-fr}[(i)].
En cuanto a la última afirmación de \sref{II.4.5.2-fr}, como $X_{f_\alpha}$ es la imagen inversa de $(\operatorname{Proj}(S))_{f_\alpha}$ por $G(\varepsilon)\to\operatorname{Proj}(S)$, basta aplicar \sref[I]{I.9.3.2-fr}.
Por tanto, \emph{(a)}, \emph{(a')}, \emph{(b)} y \emph{(b')} son equivalentes.

Para demostrar que \emph{(a')} implica \emph{(c)}, observemos que, si $X_f$ es afín (con $f\in S_k$), entonces $\mathcal{F}|X_f$ está generado por sus secciones sobre $X_f$ \sref[I]{I.1.3.9-fr};
por otra parte, por \sref[I]{I.9.3.1-fr}[(ii)], una sección tal $s$ es de la forma $(t|X_f)\otimes(f|X_f)^{-m}$, donde $t\in\Gamma(X,\mathcal{F}(km))$;
por definición, $t$ es también una sección de $\mathcal{F}_{km}$, de modo que $s$ es una sección de $\mathcal{F}_{km}(-km)$ sobre $X_f$, lo que demuestra \emph{(c)}.
Es claro que \emph{(c)} implica \emph{(c')}, y solo queda demostrar que \emph{(c')} implica \emph{(a)}.
Sea, pues, $U$ una vecindad abierta de $x\in X$, y sea $\mathcal{J}$ un haz cuasicoherente de ideales de $\mathcal{O}_X$ que define un subpreesquema cerrado de $X$ cuyo espacio subyacente es $X\setmin U$ \sref[I]{I.5.2.1-fr}.
La hipótesis~\emph{(c')} implica que existe un entero $n>0$ y una sección $f$ de $\mathcal{J}(n)$ sobre $X$ tal que $f(x)\neq0$.
Pero evidentemente $f\in S_n$ y $x\in X_f\subset U$, lo que demuestra \emph{(a)}.
\end{proof}
% Fuente francesa: PMIHES_1961__8__5_0, printed p. 84.

Cuando $X$ es un preesquema cuyo espacio subyacente es noetheriano, las condiciones equivalentes de \sref{II.4.5.2-fr} implican que $X$ es un \emph{esquema}, pues es isomorfo a un subpreesquema del esquema $P=\operatorname{Proj}(S)$, por \sref{II.4.5.2-fr}[\emph{(b)}].

\begin{definition}[4.5.3]
\label{II.4.5.3-fr}
Se dice que un $\mathcal{O}_X$-módulo invertible $\mathcal{L}$ es \emph{amplio} si $X$ es un esquema cuasicompacto y se satisfacen las condiciones equivalentes de \sref{II.4.5.2-fr}.
\end{definition}

Se sigue evidentemente del criterio~\emph{(a)} de \sref{II.4.5.2-fr} que, si $\mathcal{L}$ es un $\mathcal{O}_X$-módulo amplio, entonces, para todo abierto cuasicompacto $U$ de $X$, $\mathcal{L}|U$ es un $(\mathcal{O}_X|U)$-módulo amplio.

Se sigue de la demostración de \sref{II.4.5.2-fr} que los $X_f$ \emph{afines} forman una \emph{base} de la topología de $X$.
Además:
% Fuente francesa: PMIHES_1961__8__5_0, printed pp. 84-85.

\begin{corollary}[4.5.4]
\label{II.4.5.4-fr}
Sea $\mathcal{L}$ un $\mathcal{O}_X$-módulo amplio.
Para todo subespacio finito $Z$ de $X$ y toda vecindad $U$ de $Z$, existe un entero $n$ y un $f\in\Gamma(X,\mathcal{L}^{\otimes n})$ tales que $X_f$ es una vecindad afín de $Z$ contenida en $U$.
\end{corollary}

\begin{proof}
\oldpage[II]{85}
Por \sref{II.4.5.2-fr}[\emph{(b)}], basta demostrar que, para todo subconjunto finito $Z'$ de $\operatorname{Proj}(S)$ y toda vecindad abierta $U$ de $Z'$, existe un elemento homogéneo $f\in S_+$ tal que $Z'\subset(\operatorname{Proj}(S))_f\subset U$ \sref{II.2.4.1-fr}.
Ahora bien, por definición, el conjunto cerrado $Y$, complementario de $U$ en $\operatorname{Proj}(S)$, es de la forma $V_+(\mathfrak{J})$, donde $\mathfrak{J}$ es un ideal graduado de $S$ que no contiene $S_+$ \sref{II.2.3.2-fr};
por otra parte, los puntos de $Z'$ son, por definición, ideales primos graduados $\mathfrak{p}_i$ de $S_+$ que no contienen $\mathfrak{J}$ \sref{II.2.3.1-fr}.
Existe, pues, un elemento $f\in\mathfrak{J}$ que no pertenece a ninguno de los $\mathfrak{p}_i$ (Bourbaki, \emph{Alg. comm.}, cap.~II, \S1, n.º~1, prop.~2), y, como los $\mathfrak{p}_i$ son graduados, el razonamiento de \emph{loc. cit.} muestra que incluso se puede tomar $f$ homogéneo;
este elemento satisface entonces lo requerido.
\end{proof}
% Fuente francesa: PMIHES_1961__8__5_0, printed p. 85.

\begin{proposition}[4.5.5]
\label{II.4.5.5-fr}
Supongamos que $X$ es un esquema cuasicompacto o un preesquema cuyo espacio subyacente es noetheriano.
Entonces las condiciones~\textup{(a)} a \textup{(c')} de \sref{II.4.5.2-fr} son equivalentes también a las siguientes:
\begin{enumerate}
  \item[\rm{(d)}] Para todo $\mathcal{O}_X$-módulo cuasicoherente $\mathcal{F}$ de tipo finito, existe un entero $n_0$ tal que, para todo $n\geq n_0$, $\mathcal{F}(n)$ está generado por sus secciones sobre $X$.
  \item[\rm{(d')}] Para todo $\mathcal{O}_X$-módulo cuasicoherente $\mathcal{F}$ de tipo finito, existen enteros $n>0$ y $k>0$ tales que $\mathcal{F}$ es isomorfo a un cociente del $\mathcal{O}_X$-módulo $\mathcal{L}^{\otimes(-n)}\otimes\mathcal{O}_X^k$.
  \item[\rm{(d'')}] La propiedad~\textup{(d')} se cumple para todo haz cuasicoherente de ideales de $\mathcal{O}_X$ de tipo finito.
\end{enumerate}
\end{proposition}

\begin{proof}
Como $X$ es cuasicompacto, si un $\mathcal{O}_X$-módulo cuasicoherente de tipo finito $\mathcal{F}$ es tal que $\mathcal{F}(n)$ (que es de tipo finito) está generado por sus secciones sobre $X$, entonces $\mathcal{F}(n)$ está generado por un número \emph{finito} de estas secciones \sref[0]{0.5.2.3-fr};
por tanto, \emph{(d)} implica \emph{(d')}, y es claro que \emph{(d')} implica \emph{(d'')}.
Como todo $\mathcal{O}_X$-módulo cuasicoherente $\mathcal{G}$ es el límite inductivo de sus sub-$\mathcal{O}_X$-módulos de tipo finito \sref[I]{I.9.4.9-fr}, para verificar la condición~\emph{(c')} de \sref{II.4.5.2-fr} basta hacerlo para un haz cuasicoherente de ideales de $\mathcal{O}_X$ que sea de tipo finito, y por tanto \emph{(d'')} implica \emph{(c')}.
Solo queda demostrar que, si $\mathcal{L}$ es amplio, se satisface la propiedad~\emph{(d)}.
Consideremos una cubierta finita de $X$ por $X_{f_i}$ ($f_i\in S_{n_i}$), que se puede suponer afines;
reemplazando los $f_i$ por potencias convenientes (lo que no altera los $X_{f_i}$), se puede suponer que todos los $n_i$ son iguales a un mismo entero $m$.
El haz $\mathcal{F}|X_{f_i}$, que es de tipo finito por hipótesis, está generado por un número finito de sus secciones $h_{ij}$ sobre $X_{f_i}$ \sref[I]{I.1.3.13-fr};
existe, pues, un entero $k_0$ tal que la sección $h_{ij}\otimes f_i^{\otimes k_0}$ se prolonga a una sección de $\mathcal{F}(k_0m)$ sobre $X$ para todo par $(i,j)$ \sref[I]{I.9.3.1-fr}.
\emph{A fortiori}, los $h_{ij}\otimes f_i^{\otimes k}$ se prolongan a secciones de $\mathcal{F}(km)$ sobre $X$ para todo $k\geq k_0$, y, para estos valores de $k$, $\mathcal{F}(km)$ está generado por sus secciones sobre $X$.
Para todo $p$ tal que $0<p<m$, $\mathcal{F}(p)$ es también de tipo finito, y por tanto existe un entero $k_p$ tal que $\mathcal{F}(p)(km)=\mathcal{F}(p+km)$ está generado por sus secciones sobre $X$ para todo $k\geq k_p$.
Tomando por $n_0$ el mayor de los $k_pm$, se concluye que $\mathcal{F}(n)$ está generado por sus secciones sobre $X$ para todo $n\geq n_0$, pues tal $n$ es de la forma $n=km+p$, con $k\geq k_p$ y $0\leq p<m$.
\end{proof}
% Fuente francesa: PMIHES_1961__8__5_0, printed pp. 85-86.

\begin{proposition}[4.5.6]
\label{II.4.5.6-fr}
Sean $X$ un esquema cuasicompacto y $\mathcal{L}$ un $\mathcal{O}_X$-módulo invertible.
\begin{enumerate}
  \item[\rm{(i)}] Sea $n$ un entero $>0$. Para que $\mathcal{L}$ sea amplio, es necesario y suficiente que $\mathcal{L}^{\otimes n}$ sea amplio.
  \item[\rm{(ii)}] Sea $\mathcal{L}'$ un $\mathcal{O}_X$-módulo invertible tal que, para todo $x\in X$, existe un entero $n>0$
\oldpage[II]{86}
    y una sección $s'$ de $\mathcal{L}'^{\otimes n}$ sobre $X$ tal que $s'(x)\neq0$.
    Entonces, si $\mathcal{L}$ es amplio, también lo es $\mathcal{L}\otimes\mathcal{L}'$.
\end{enumerate}
\end{proposition}

\begin{proof}
La propiedad~(i) es una consecuencia evidente del criterio~\emph{(a)} de \sref{II.4.5.2-fr}, pues $X_{f^{\otimes n}}=X_f$.
Por otra parte, si $\mathcal{L}$ es amplio, para todo $x\in X$ y toda vecindad $U$ de $x$, existe un $m>0$ y un $f\in\Gamma(X,\mathcal{L}^{\otimes m})$ tales que $x\in X_f\subset U$ \sref{II.4.5.2-fr}[\emph{(a)}];
si $f'\in\Gamma(X,\mathcal{L}'^{\otimes n})$ es tal que $f'(x)\neq0$, entonces $s(x)\neq0$ para $s=f^{\otimes n}\otimes f'^{\otimes m}\in
\Gamma(X,(\mathcal{L}\otimes\mathcal{L}')^{\otimes mn})$, y por tanto $x\in X_s\subset X_f\subset U$, lo que demuestra que $\mathcal{L}\otimes\mathcal{L}'$ es amplio \sref{II.4.5.2-fr}[\emph{(a)}].
\end{proof}

\begin{corollary}[4.5.7]
\label{II.4.5.7-fr}
El producto tensorial de dos $\mathcal{O}_X$-módulos amplios es amplio.
\end{corollary}
% Fuente francesa: PMIHES_1961__8__5_0, printed p. 86.

\begin{corollary}[4.5.8]
\label{II.4.5.8-fr}
Sean $\mathcal{L}$ un $\mathcal{O}_X$-módulo amplio y $\mathcal{L}'$ un $\mathcal{O}_X$-módulo invertible;
entonces existe un entero $n_0>0$ tal que $\mathcal{L}^{\otimes n}\otimes\mathcal{L}'$ es amplio y está generado por sus secciones sobre $X$ para $n\geq n_0$.
\end{corollary}

\begin{proof}
Se sigue de \sref{II.4.5.5-fr} que existe un entero $m_0$ tal que $\mathcal{L}^{\otimes m}\otimes\mathcal{L}'$ está generado por sus secciones sobre $X$ para todo $m\geq m_0$;
por \sref{II.4.5.6-fr}, se puede tomar entonces $n_0=m_0+1$.
\end{proof}

\begin{remark}[4.5.9]
\label{II.4.5.9-fr}
Sea $P=\mathrm{H}^1(X,\mathcal{O}_X^*)$ el grupo de las clases de $\mathcal{O}_X$-módulos invertibles \sref[0]{0.5.4.7-fr}, y sea $P^+$ el subconjunto de $P$ formado por las clases de haces amplios.
Supongamos que $P^+$ es \emph{no vacío}.
Se sigue entonces de \sref{II.4.5.7-fr} y \sref{II.4.5.8-fr} que
\[
  P^++P^+\subset P^+
  \qquad\text{y}\qquad
  P^+-P^+=P
\]
o, dicho de otro modo, $P^+\cup\{0\}$ es el conjunto de los elementos \emph{positivos} de $P$ para una estructura de \emph{preorden} sobre $P$ compatible con su estructura de grupo, y es incluso \emph{arquimediana}, por \sref{II.4.5.8-fr}.
Por esta razón se dice a veces ``haz positivo'' en lugar de haz amplio, y ``haz negativo'' para el inverso de un haz amplio (pero no utilizaremos esta terminología).
\end{remark}
% Fuente francesa: PMIHES_1961__8__5_0, printed pp. 86-87.

\begin{proposition}[4.5.10]
\label{II.4.5.10-fr}
Sean $Y$ un esquema afín, $q:X\to Y$ un morfismo separado cuasicompacto y $\mathcal{L}$ un $\mathcal{O}_X$-módulo invertible.
\begin{enumerate}
  \item[\rm{(i)}] Si $\mathcal{L}$ es muy amplio relativamente a $q$, entonces $\mathcal{L}$ es amplio.
  \item[\rm{(ii)}] Supongamos además que el morfismo $q$ es \emph{de tipo finito}.
    Para que $\mathcal{L}$ sea amplio, es necesario y suficiente que posea una de las propiedades equivalentes siguientes:
    \begin{enumerate}
      \item[\rm{(e)}] Existe $n_0>0$ tal que, para todo entero $n\geq n_0$, $\mathcal{L}^{\otimes n}$ es muy amplio relativamente a $q$.
      \item[\rm{(e')}] Existe $n>0$ tal que $\mathcal{L}^{\otimes n}$ es muy amplio relativamente a $q$.
    \end{enumerate}
\end{enumerate}
\end{proposition}

\begin{proof}
La primera afirmación se sigue de la definición \sref{II.4.4.2-fr} de un $\mathcal{O}_X$-módulo muy amplio:
si $A$ es el anillo de $Y$, existen un $A$-módulo $E$ y un homomorfismo sobreyectivo
\[
  \psi:q^*((\mathbf{S}(E))^{\sim})\longrightarrow
  \bigoplus_{n\geq0}\mathcal{L}^{\otimes n}
\]
tales que $i=r_{\mathcal{L},\psi}$ es una inmersión definida en todo punto $X\to P=\mathbf{P}(\widetilde{E})$ y $\mathcal{L}=i^*(\mathcal{O}_P(1))$;
como los $D_+(f)$, para $f$ homogéneo en $(\mathbf{S}(E))_+$, forman una base de la topología de $P$, y como $i^{-1}(D_+(f))=X_{\psi^\flat(f)}$ por \sref{II.3.7.3.1-fr}, se ve que se satisface la condición~\emph{(a)} de \sref{II.4.5.2-fr}, y por tanto $\mathcal{L}$ es amplio.

Demostremos ahora que, si $q$ es de tipo finito y $\mathcal{L}$ es amplio, se satisface la condición~\emph{(e)}.
En primer lugar, se sigue del criterio~\emph{(b)} de \sref{II.4.5.2-fr} y de \sref{II.4.4.1-fr}[(i)] que existe
\oldpage[II]{87}
un entero $k_0>0$ tal que $\mathcal{L}^{\otimes k_0}$ es muy amplio relativamente a $q$.
Por otra parte, en virtud de \sref{II.4.5.5-fr}, existe un entero $m_0$ tal que, para todo $m\geq m_0$, $\mathcal{L}^{\otimes m}$ está generado por sus secciones sobre $X$.
Pongamos $n_0=k_0+m_0$;
si $n\geq n_0$, se puede escribir $n=k_0+m$, con $m\geq m_0$, de donde $\mathcal{L}^{\otimes n}
=\mathcal{L}^{\otimes k_0}\otimes\mathcal{L}^{\otimes m}$.
Como $\mathcal{L}^{\otimes m}$ está generado por sus secciones sobre $X$, se sigue de \sref{II.4.4.8-fr} y \sref{II.3.4.7-fr} que $\mathcal{L}^{\otimes n}$ es muy amplio relativamente a $q$.
Finalmente, está claro que \emph{(e)} implica \emph{(e')}, y que \emph{(e')} implica que $\mathcal{L}$ es amplio en virtud de (i) y de \sref{II.4.5.6-fr}[(i)].

  \begin{env}[4.5.10.1]
  \label{II.4.5.10.1-fr}
  [\emph{Demostración del lema}~\sref{II.4.4.10.1-fr}].
  Pongamos $\mathcal{E}(n)=\mathcal{E}\otimes\mathcal{K}^{\otimes n}$;
  como $g$ es separado \sref{II.4.4.2-fr}, el razonamiento de \sref{II.3.4.8-fr} se aplica y reduce la afirmación a demostrar que el homomorfismo canónico $g^*(g_*(\mathcal{E}(n)))\to\mathcal{E}(n)$ es sobreyectivo para todo $n$ suficientemente grande.
  Además, como $Z$ es cuasicompacto, el razonamiento de \sref{II.3.4.6-fr} reduce a demostrar la afirmación cuando $Z$ es afín.
  Pero entonces $\mathcal{K}$ es amplio por \sref{II.4.5.10-fr}[(i)], y la conclusión se sigue de \sref{II.4.5.5-fr}[\emph{(d')}].
  \end{env}
\end{proof}
% Fuente francesa: PMIHES_1961__8__5_0, printed p. 87.

\begin{corollary}[4.5.11]
\label{II.4.5.11-fr}
Sean $Y$ un esquema afín, $q:X\to Y$ un morfismo separado de tipo finito, $\mathcal{L}$ un $\mathcal{O}_X$-módulo amplio y $\mathcal{L}'$ un $\mathcal{O}_X$-módulo invertible.
Existe un entero $n_0$ tal que, para todo $n\geq n_0$, $\mathcal{L}^{\otimes n}\otimes\mathcal{L}'$ es muy amplio relativamente a $q$.
\end{corollary}

\begin{proof}
Existe $m_0$ tal que, para $m\geq m_0$, $\mathcal{L}^{\otimes m}\otimes\mathcal{L}'$ está generado por sus secciones sobre $X$ \sref{II.4.5.8-fr};
existe asimismo $k_0$ tal que $\mathcal{L}^{\otimes k}$ es muy amplio relativamente a $q$ para $k\geq k_0$.
Entonces $\mathcal{L}^{\otimes(k+m_0)}\otimes\mathcal{L}'$ es muy amplio si $k\geq k_0$ (\sref{II.4.4.8-fr} y \sref{II.3.4.7-fr}).
\end{proof}

\begin{remark}[4.5.12]
\label{II.4.5.12-fr}
Se ignora si la hipótesis de que un $\mathcal{O}_X$-módulo $\mathcal{L}$ sea tal que $\mathcal{L}^{\otimes n}$ es muy amplio (relativamente a $q$) implica la misma conclusión para $\mathcal{L}^{\otimes(n+1)}$.
\end{remark}
% Fuente francesa: PMIHES_1961__8__5_0, printed pp. 87-88.

\begin{proposition}[4.5.13]
\label{II.4.5.13-fr}
Sean $X$ un preesquema cuasicompacto, $Z$ un subpreesquema cerrado de $X$ definido por un haz cuasicoherente \emph{nilpotente} $\mathcal{J}$ de ideales de $\mathcal{O}_X$, y $j$ la inyección canónica $Z\to X$.
Para que un $\mathcal{O}_X$-módulo invertible $\mathcal{L}$ sea amplio, es necesario y suficiente que $\mathcal{L}'=j^*(\mathcal{L})$ sea un $\mathcal{O}_Z$-módulo amplio.
\end{proposition}

\begin{proof}
La condición es \emph{necesaria}.
En efecto, para toda sección $f$ de $\mathcal{L}^{\otimes n}$ sobre $X$, sea $f'$ su imagen canónica $f\otimes1$, que es una sección de $\mathcal{L}'^{\otimes n}=\mathcal{L}^{\otimes n}
\otimes_{\mathcal{O}_X}(\mathcal{O}_X/\mathcal{J})$ sobre el espacio $Z$ (que es idéntico a $X$);
está claro que $X_f=Z_{f'}$, y por tanto el criterio~\emph{(a)} de \sref{II.4.5.2-fr} muestra que $\mathcal{L}'$ es amplio.

Para ver que la condición es \emph{suficiente}, observemos ante todo que se puede reducir al caso en que $\mathcal{J}^2=0$, considerando la sucesión (finita) de preesquemas $X_k=(X,\mathcal{O}_X/\mathcal{J}^{k+1})$, cada uno de los cuales es un subpreesquema cerrado del siguiente, definido por un haz de ideales de cuadrado nulo.
Además, $X$ es un esquema, ya que $X_{\mathrm{red}}$ lo es por hipótesis (\sref{II.4.5.3-fr} y \sref[I]{I.5.5.1-fr}).
El criterio~\emph{(a)} de \sref{II.4.5.2-fr} muestra que basta demostrar el
  \begin{lemma}[4.5.13.1]
  \label{II.4.5.13.1-fr}
  Bajo las hipótesis de \sref{II.4.5.13-fr}, supongamos además que $\mathcal{J}$ es de cuadrado nulo;
  siendo $\mathcal{L}$ un $\mathcal{O}_X$-módulo invertible, sea $g$ una sección de $\mathcal{L}'^{\otimes n}$ sobre $Z$ tal que $Z_g$ es afín.
  Entonces existe un entero $m>0$ tal que $g^{\otimes m}$ es la imagen canónica de una sección $f$ de $\mathcal{L}^{\otimes nm}$ sobre $X$.
  \end{lemma}

  Se tiene la sucesión exacta de $\mathcal{O}_X$-módulos
  \[
  0\longrightarrow\mathcal{J}(n)\longrightarrow
  \mathcal{O}_X(n)=\mathcal{L}^{\otimes n}\longrightarrow
  \mathcal{O}_Z(n)=\mathcal{L}'^{\otimes n}\longrightarrow0
\]
\oldpage[II]{88}
  ya que $\mathcal{F}(n)$ es un funtor exacto en $\mathcal{F}$;
  de donde la sucesión exacta de cohomología
  \[
  0\longrightarrow\Gamma(X,\mathcal{J}(n))\longrightarrow
  \Gamma(X,\mathcal{L}^{\otimes n})\longrightarrow
  \Gamma(X,\mathcal{L}'^{\otimes n})\xrightarrow{\partial}
  \mathrm{H}^1(X,\mathcal{J}(n))
\]
  que envía, en particular, $g$ a un elemento $\partial g\in\mathrm{H}^1(X,\mathcal{J}(n))$.
  Observemos que, como $\mathcal{J}^2=0$, $\mathcal{J}$ puede considerarse como un $\mathcal{O}_Z$-módulo cuasicoherente y, para todo $k$, $\mathcal{L}'^{\otimes k}\otimes_{\mathcal{O}_Z}\mathcal{J}(n)
=\mathcal{J}(n+k)$;
  para toda sección $s\in\Gamma(X,\mathcal{L}'^{\otimes k})$, la multiplicación tensorial por $s$ es, por tanto, un homomorfismo $\mathcal{J}(n)\to\mathcal{J}(n+k)$ de $\mathcal{O}_Z$-módulos, que induce un homomorfismo $\mathrm{H}^i(X,\mathcal{J}(n))\xrightarrow{s}
\mathrm{H}^i(X,\mathcal{J}(n+k))$ de grupos de cohomología.
% Fuente francesa: PMIHES_1961__8__5_0, printed pp. 88-89.

  Dicho esto, veremos que
  \begin{equation}
\label{II.4.5.13.2-fr}
  g^{\otimes m}\otimes\partial g=0
\tag{4.5.13.2}
\end{equation}
  para $m>0$ suficientemente grande.
  En efecto, $Z_g$ es un \emph{abierto afín} de $Z$, y por tanto $\mathrm{H}^1(Z_g,\mathcal{J}(n))=0$ cuando $\mathcal{J}(n)$ se considera como un \emph{$\mathcal{O}_Z$-módulo} \sref[I]{I.5.1.9.2-fr}.
  En particular, si se pone $g'=g|Z_g$ y se considera su imagen por la aplicación $\partial:\Gamma(Z_g,\mathcal{L}'^{\otimes n})\to
\mathrm{H}^1(Z_g,\mathcal{J}(n))$, se tiene $\partial g'=0$.
  Para explicitar esta relación, observemos que, en grado~$1$, la cohomología de un haz de grupos abelianos coincide con su cohomología de \v{C}ech (G, II, 5.9);
  para calcular $\partial g$, hay que considerar, pues, una cubierta abierta $(U_\alpha)$ suficientemente fina de $X$, que se puede suponer \emph{finita} y formada por abiertos afines, tomar, para cada $\alpha$, una sección $g_\alpha\in\Gamma(U_\alpha,\mathcal{L}^{\otimes n})$ cuya imagen canónica en $\Gamma(U_\alpha,\mathcal{L}'^{\otimes n})$ sea $g|U_\alpha$, y considerar la clase del cociclo $(g_{\alpha|\beta}-g_{\beta|\alpha})$, donde $g_{\alpha|\beta}$ denota la restricción de $g_\alpha$ a $U_\alpha\cap U_\beta$ (este cociclo toma valores en $\mathcal{J}(n)$).
  Se puede suponer además que $\partial g'$ se calcula del mismo modo, por medio de la cubierta formada por los $U_\alpha\cap Z_g$ y de las restricciones $g_\alpha|(U_\alpha\cap Z_g)$ (reemplazando, si es necesario, $(U_\alpha)$ por una cubierta más fina);
  la relación $\partial g'=0$ implica entonces que existe, para cada $\alpha$, una sección $h_\alpha\in\Gamma(U_\alpha\cap Z_g,\mathcal{J}(n))$ tal que
  $(g_{\alpha|\beta}-g_{\beta|\alpha})
|(U_\alpha\cap U_\beta\cap Z_g)
=h_{\alpha|\beta}-h_{\beta|\alpha}$
  donde $h_{\alpha|\beta}$ denota la restricción de $h_\alpha$ a $U_\alpha\cap U_\beta\cap Z_g$ (G, II, 5.11).
  Existe entonces un entero $m>0$ tal que $g^{\otimes m}\otimes h_\alpha$ es la restricción a $U_\alpha\cap Z_g$ de una sección $t_\alpha\in\Gamma(U_\alpha,\mathcal{J}(n+nm))$ para todo $\alpha$ \sref[I]{I.9.3.1-fr};
  por tanto, $g^{\otimes m}\otimes(g_{\alpha|\beta}-g_{\beta|\alpha})
=t_{\alpha|\beta}-t_{\beta|\alpha}$ para todo par de índices, lo que demuestra \sref{II.4.5.13.2-fr}.

  Observemos, por otra parte, que, si $s\in\Gamma(X,\mathcal{O}_Z(p))$ y $t\in\Gamma(X,\mathcal{O}_Z(q))$, se tiene, en el grupo $\mathrm{H}^1(X,\mathcal{J}(p+q))$,
  \begin{equation}
\label{II.4.5.13.3-fr}
  \partial(s\otimes t)=(\partial s)\otimes t+s\otimes(\partial t).
\tag{4.5.13.3}
\end{equation}

  En efecto, se puede calcular de nuevo ambos miembros considerando una cubierta abierta $(U_\alpha)$ de $X$ y, para cada $\alpha$, una sección $s_\alpha\in\Gamma(U_\alpha,\mathcal{O}_X(p))$ (resp. $t_\alpha\in\Gamma(U_\alpha,\mathcal{O}_X(q))$) cuya imagen canónica en $\Gamma(U_\alpha,\mathcal{O}_Z(p))$ (resp. $\Gamma(U_\alpha,\mathcal{O}_Z(q))$) sea $s|U_\alpha$ (resp. $t|U_\alpha$);
  la relación~\sref{II.4.5.13.3-fr} se sigue entonces de las relaciones
  \[
  (s_{\alpha|\beta}\otimes t_{\alpha|\beta})
  -(s_{\beta|\alpha}\otimes t_{\beta|\alpha})
  =(s_{\alpha|\beta}-s_{\beta|\alpha})\otimes t_{\alpha|\beta}
  +s_{\beta|\alpha}\otimes(t_{\alpha|\beta}-t_{\beta|\alpha})
\]
  con las mismas notaciones que antes.
  Por inducción sobre $k$, se obtiene, por tanto,
  \begin{equation}
\label{II.4.5.13.4-fr}
  \partial(g^{\otimes k})=(kg^{\otimes(k-1)})\otimes(\partial g)
\tag{4.5.13.4}
\end{equation}
\oldpage[II]{89}
  y se concluye de \sref{II.4.5.13.2-fr} y \sref{II.4.5.13.4-fr} que $\partial(g^{\otimes(m+1)})=0$;
  por tanto, $g^{\otimes(m+1)}$ es la imagen canónica de una sección $f$ de $\mathcal{L}^{\otimes n(m+1)}$ sobre $X$, lo que demuestra \sref{II.4.5.13-fr}.
\end{proof}
% Fuente francesa: PMIHES_1961__8__5_0, printed p. 89.

\begin{corollary}[4.5.14]
\label{II.4.5.14-fr}
Sean $X$ un esquema noetheriano y $j$ la inyección canónica $X_{\mathrm{red}}\to X$.
Para que un $\mathcal{O}_X$-módulo invertible $\mathcal{L}$ sea amplio, es necesario y suficiente que $j^*(\mathcal{L})$ sea un $\mathcal{O}_{X_{\mathrm{red}}}$-módulo amplio.
\end{corollary}

\begin{proof}
Esto se sigue de \sref[I]{I.6.1.6-fr}.
\end{proof}
% Fuente francesa: PMIHES_1961__8__5_0, printed p. 89.

\subsection{Haces relativamente amplios}
\label{II.4.6-fr}

\begin{definition}[4.6.1]
\label{II.4.6.1-fr}
Sean $f:X\to Y$ un morfismo cuasicompacto y $\mathcal{L}$ un $\mathcal{O}_X$-módulo invertible.
Se dice que $\mathcal{L}$ es \emph{amplio relativamente a $f$}, o \emph{relativamente a $Y$}, o \emph{$f$-amplio}, o \emph{$Y$-amplio} \textup{(o incluso simplemente \emph{amplio}, si no puede resultar ninguna confusión con la noción definida en \sref{II.4.5.3-fr})} si existe una cubierta $(U_\alpha)$ de $Y$ por abiertos afines tal que, poniendo $X_\alpha=f^{-1}(U_\alpha)$, $\mathcal{L}|X_\alpha$ es un $\mathcal{O}_{X_\alpha}$-módulo amplio para todo $\alpha$.
\end{definition}

La existencia de un $\mathcal{O}_X$-módulo $f$-amplio implica necesariamente que $f$ es \emph{separado} (\sref{II.4.5.3-fr} y \sref[I]{I.5.5.5-fr}).

\begin{proposition}[4.6.2]
\label{II.4.6.2-fr}
Sean $f:X\to Y$ un morfismo cuasicompacto y $\mathcal{L}$ un $\mathcal{O}_X$-módulo invertible.
Si $\mathcal{L}$ es muy amplio relativamente a $f$, entonces es amplio relativamente a $f$.
\end{proposition}

\begin{proof}
Esto se sigue del carácter local (sobre $Y$) de la noción de haz muy amplio \sref{II.4.4.5-fr}, de la definición \sref{II.4.6.1-fr} y del criterio~\sref{II.4.5.10-fr}[(i)].
\end{proof}
% Fuente francesa: PMIHES_1961__8__5_0, printed pp. 89-90.

\begin{proposition}[4.6.3]
\label{II.4.6.3-fr}
Sean $f:X\to Y$ un morfismo cuasicompacto, $\mathcal{L}$ un $\mathcal{O}_X$-módulo invertible y $\mathcal{S}$ la $\mathcal{O}_Y$-álgebra graduada $\bigoplus_{n\geq0}f_*(\mathcal{L}^{\otimes n})$.
Las condiciones siguientes son equivalentes:
\begin{enumerate}
  \item[\rm{(a)}] $\mathcal{L}$ es $f$-amplio.
  \item[\rm{(b)}] $\mathcal{S}$ es cuasicoherente, y el homomorfismo canónico $\sigma:f^*(\mathcal{S})\to
  \bigoplus_{n\geq0}\mathcal{L}^{\otimes n}$ \sref[0]{0.4.4.3-fr} es tal que el $Y$-morfismo $r_{\mathcal{L},\sigma}:G(\sigma)\to\operatorname{Proj}(\mathcal{S})=P$ está definido en todo punto y es una inmersión abierta dominante.
  \item[\rm{(b')}] El morfismo $f$ es separado, y el $Y$-morfismo $r_{\mathcal{L},\sigma}$ está definido en todo punto y es un homeomorfismo del espacio subyacente a $X$ sobre un subespacio de $\operatorname{Proj}(\mathcal{S})$.
\end{enumerate}

Además, cuando se satisfacen estas condiciones, para todo $n\in\mathbf{Z}$, el homomorfismo canónico
\begin{equation}
\label{II.4.6.3.1-fr}
  r_{\mathcal{L},\sigma}^*(\mathcal{O}_P(n))
  \longrightarrow\mathcal{L}^{\otimes n}
\tag{4.6.3.1}
\end{equation}
definido en \sref{II.3.7.9.1-fr} es un isomorfismo.

Finalmente, para todo $\mathcal{O}_X$-módulo cuasicoherente $\mathcal{F}$, si se pone $\mathcal{M}=\bigoplus_{n\geq0}f_*(\mathcal{F}\otimes
\mathcal{L}^{\otimes n})$, el homomorfismo canónico
\begin{equation}
\label{II.4.6.3.2-fr}
  r_{\mathcal{L},\sigma}^*(\widetilde{\mathcal{M}})
  \longrightarrow\mathcal{F}
\tag{4.6.3.2}
\end{equation}
definido en \sref{II.3.7.9.2-fr} es un isomorfismo.
\end{proposition}

\begin{proof}
Se ha observado que \emph{(a)} implica que $f$ es separado, y por tanto que $\mathcal{S}$ es cuasicoherente \sref[I]{I.9.2.2-fr}[\emph{(a)}].
Como el hecho de que $r_{\mathcal{L},\sigma}$ sea una inmersión definida en todo punto tiene carácter local sobre $Y$, para demostrar que \emph{(a)} implica \emph{(b)} se puede suponer que $Y$ es afín y $\mathcal{L}$ amplio;
\oldpage[II]{90}
la afirmación se sigue entonces de \sref{II.4.5.2-fr}[\emph{(b)}].
Está claro que \emph{(b)} implica \emph{(b')};
finalmente, para demostrar que \emph{(b')} implica \emph{(a)}, basta considerar una cubierta de $Y$ por abiertos afines $U_\alpha$ y aplicar el criterio~\sref{II.4.5.2-fr}[\emph{(b')}] a cada haz $\mathcal{L}|f^{-1}(U_\alpha)$.

Para las dos últimas afirmaciones, se utiliza que $\sigma^\flat$ es aquí la identidad y la explicitación de los homomorfismos \sref{II.3.7.9.1-fr} y \sref{II.3.7.9.2-fr};
de ello se sigue inmediatamente que \sref{II.4.6.3.1-fr} es un isomorfismo.
En cuanto a \sref{II.4.6.3.2-fr}, se puede reducir al caso en que $Y$ es afín, y por tanto $\mathcal{L}$ amplio;
está claro que el homomorfismo \sref{II.4.6.3.2-fr} es inyectivo, y el criterio~\sref{II.4.5.2-fr}[\emph{(c)}] muestra que es sobreyectivo, de donde se sigue la conclusión.
\end{proof}
% Fuente francesa: PMIHES_1961__8__5_0, printed p. 90.

\begin{corollary}[4.6.4]
\label{II.4.6.4-fr}
Sea $(U_\alpha)$ una cubierta abierta de $Y$.
Para que $\mathcal{L}$ sea amplio relativamente a $Y$, es necesario y suficiente que $\mathcal{L}|f^{-1}(U_\alpha)$ sea amplio relativamente a $U_\alpha$ para todo $\alpha$.
\end{corollary}

\begin{proof}
En efecto, la condición~\emph{(b)} es local sobre $Y$.
\end{proof}

\begin{corollary}[4.6.5]
\label{II.4.6.5-fr}
Sea $\mathcal{K}$ un $\mathcal{O}_Y$-módulo invertible.
Para que $\mathcal{L}$ sea $Y$-amplio, es necesario y suficiente que $\mathcal{L}\otimes f^*(\mathcal{K})$ lo sea.
\end{corollary}

\begin{proof}
Es una consecuencia evidente de \sref{II.4.6.4-fr}, tomando los $U_\alpha$ de modo que $\mathcal{K}|U_\alpha$ sea isomorfo a $\mathcal{O}_Y|U_\alpha$ para todo $\alpha$.
\end{proof}

\begin{corollary}[4.6.6]
\label{II.4.6.6-fr}
Supongamos que $Y$ es afín;
para que $\mathcal{L}$ sea $Y$-amplio, es necesario y suficiente que $\mathcal{L}$ sea amplio.
\end{corollary}

\begin{proof}
Es una consecuencia inmediata de la definición \sref{II.4.6.1-fr} y de los criterios \sref{II.4.6.3-fr}[\emph{(b)}] y \sref{II.4.5.2-fr}[\emph{(b)}], ya que aquí $\operatorname{Proj}(\mathcal{S})
=\operatorname{Proj}(\Gamma(Y,\mathcal{S}))$ por definición.
\end{proof}

\begin{corollary}[4.6.7]
\label{II.4.6.7-fr}
Sea $f:X\to Y$ un morfismo cuasicompacto.
Supongamos que existen un $\mathcal{O}_Y$-módulo cuasicoherente $\mathcal{E}$ y un $Y$-morfismo $g:X\to P=\mathbf{P}(\mathcal{E})$ que es un homeomorfismo del espacio subyacente a $X$ sobre un subespacio de $P$;
entonces $\mathcal{L}=g^*(\mathcal{O}_P(1))$ es $Y$-amplio.
\end{corollary}

\begin{proof}
Se puede suponer que $Y$ es afín;
el corolario se sigue entonces del criterio~\sref{II.4.5.2-fr}[\emph{(a)}], de la fórmula~\sref{II.3.7.3.1-fr} y de \sref{II.4.2.3-fr}.
\end{proof}
% Fuente francesa: PMIHES_1961__8__5_0, printed pp. 90-91.

\begin{proposition}[4.6.8]
\label{II.4.6.8-fr}
Sean $X$ un esquema cuasicompacto o un preesquema cuyo espacio subyacente sea noetheriano, y $f:X\to Y$ un morfismo separado cuasicompacto.
Para que un $\mathcal{O}_X$-módulo invertible $\mathcal{L}$ sea $f$-amplio, es necesario y suficiente que se satisfaga una de las condiciones equivalentes siguientes:
\begin{enumerate}
  \item[\rm{(c)}] Para todo $\mathcal{O}_X$-módulo cuasicoherente $\mathcal{F}$ de tipo finito, existe un entero $n_0>0$ tal que, para todo $n\geq n_0$, el homomorfismo canónico $\sigma:f^*(f_*(\mathcal{F}\otimes\mathcal{L}^{\otimes n}))\to
  \mathcal{F}\otimes\mathcal{L}^{\otimes n}$ es sobreyectivo.
  \item[\rm{(c')}] Para todo haz cuasicoherente $\mathcal{J}$ de ideales de $\mathcal{O}_X$ de tipo finito, existe un entero $n>0$ tal que el homomorfismo canónico $\sigma:f^*(f_*(\mathcal{J}\otimes\mathcal{L}^{\otimes n}))\to
  \mathcal{J}\otimes\mathcal{L}^{\otimes n}$ es sobreyectivo.
\end{enumerate}
\end{proposition}

\begin{proof}
Como $X$ es cuasicompacto, también lo es $f(X)$, y existe, por tanto, una cubierta finita $(U_i)$ de $f(X)$ por abiertos afines de $Y$.
Para demostrar la condición~\emph{(c)} cuando $\mathcal{L}$ es $f$-amplio, se puede reemplazar $Y$ por los $U_i$ y $X$ por los $f^{-1}(U_i)$, pues, si se obtiene para cada $i$ un entero $n_i$ tal que \emph{(c)} se satisfaga (para $U_i$, $f^{-1}(U_i)$ y $\mathcal{L}|f^{-1}(U_i)$) para todo $n\geq n_i$, basta tomar como $n_0$ el mayor de los $n_i$ para obtener \emph{(c)} para $Y$, $X$ y $\mathcal{L}$.
Pero, si $Y$ es afín, la condición~\emph{(c)} se sigue de \sref{II.4.5.5-fr}[\emph{(d)}], teniendo en cuenta \sref{II.4.6.6-fr}.
Es trivial que \emph{(c)} implica \emph{(c')}.
Finalmente, para demostrar que \emph{(c')} implica que $\mathcal{L}$ es $f$-amplio, se puede reducir de nuevo al caso en que $Y$ es afín:
en efecto, todo haz cuasicoherente $\mathcal{J}_i$ de ideales de $\mathcal{O}_X|f^{-1}(U_i)$ de tipo finito es la restricción de un haz cuasicoherente de ideales de $\mathcal{O}_X$ de tipo finito \sref[I]{I.9.4.7-fr}, y la hipótesis~\emph{(c')} implica que
\oldpage[II]{91}
$\mathcal{J}_i\otimes(\mathcal{L}^{\otimes n}|f^{-1}(U_i))$ está generado por sus secciones (teniendo en cuenta \sref[I]{I.9.2.2-fr} y \sref{II.3.4.7-fr});
basta entonces aplicar el criterio~\sref{II.4.5.5-fr}[\emph{(d'')}].
\end{proof}
% Fuente francesa: PMIHES_1961__8__5_0, printed p. 91.

\begin{proposition}[4.6.9]
\label{II.4.6.9-fr}
Sean $f:X\to Y$ un morfismo cuasicompacto y $\mathcal{L}$ un $\mathcal{O}_X$-módulo invertible.
\begin{enumerate}
  \item[\rm{(i)}] Sea $n$ un entero $>0$.
    Para que $\mathcal{L}$ sea $f$-amplio, es necesario y suficiente que $\mathcal{L}^{\otimes n}$ sea $f$-amplio.
  \item[\rm{(ii)}] Sea $\mathcal{L}'$ un $\mathcal{O}_X$-módulo invertible, y supongamos que existe un entero $n>0$ tal que el homomorfismo canónico $\sigma:f^*(f_*(\mathcal{L}'^{\otimes n}))\to
  \mathcal{L}'^{\otimes n}$ es sobreyectivo.
    Entonces, si $\mathcal{L}$ es $f$-amplio, también lo es $\mathcal{L}\otimes\mathcal{L}'$.
\end{enumerate}
\end{proposition}

\begin{proof}
Se puede reducir inmediatamente al caso en que $Y$ es afín, y la proposición es entonces una consecuencia inmediata de \sref{II.4.5.6-fr}.
\end{proof}

\begin{corollary}[4.6.10]
\label{II.4.6.10-fr}
El producto tensorial de dos $\mathcal{O}_X$-módulos $f$-amplios es $f$-amplio.
\end{corollary}
% Fuente francesa: PMIHES_1961__8__5_0, printed p. 91.

\begin{proposition}[4.6.11]
\label{II.4.6.11-fr}
Sean $Y$ un preesquema cuasicompacto, $f:X\to Y$ un morfismo de tipo finito y $\mathcal{L}$ un $\mathcal{O}_X$-módulo invertible.
Para que $\mathcal{L}$ sea $f$-amplio, es necesario y suficiente que posea una de las propiedades equivalentes siguientes:
\begin{enumerate}
  \item[\rm{(d)}] Existe $n_0>0$ tal que, para todo entero $n\geq n_0$, $\mathcal{L}^{\otimes n}$ es muy amplio relativamente a $f$.
  \item[\rm{(d')}] Existe $n>0$ tal que $\mathcal{L}^{\otimes n}$ es muy amplio relativamente a $f$.
\end{enumerate}
\end{proposition}

\begin{proof}
Si $\mathcal{L}$ es amplio relativamente a $f$, existe una cubierta \emph{finita} $(U_i)$ de $Y$ por abiertos afines tal que los $\mathcal{L}|f^{-1}(U_i)$ son amplios.
Se concluye entonces de \sref{II.4.5.10-fr} que existe un entero $n_0$ tal que $\mathcal{L}^{\otimes n}|f^{-1}(U_i)$ es muy amplio relativamente a $f^{-1}(U_i)\to U_i$ para todo $n\geq n_0$ y todo $i$, y por tanto $\mathcal{L}^{\otimes n}$ es muy amplio relativamente a $f$ \sref{II.4.4.5-fr}.
Recíprocamente, \emph{(d')} implica ya que $\mathcal{L}^{\otimes n}$ es $f$-amplio \sref{II.4.6.2-fr}, y por tanto también lo es $\mathcal{L}$ \sref{II.4.6.9-fr}[(i)].
\end{proof}

\begin{corollary}[4.6.12]
\label{II.4.6.12-fr}
Sean $Y$ un preesquema cuasicompacto, $f:X\to Y$ un morfismo de tipo finito, y $\mathcal{L}$ y $\mathcal{L}'$ dos $\mathcal{O}_X$-módulos invertibles.
Si $\mathcal{L}$ es $f$-amplio, existe $n_0$ tal que $\mathcal{L}^{\otimes n}\otimes\mathcal{L}'$ es muy amplio relativamente a $f$ para todo $n\geq n_0$.
\end{corollary}

\begin{proof}
Se razona como en \sref{II.4.6.11-fr}, utilizando una cubierta abierta afín finita de $Y$ y \sref{II.4.5.11-fr}.
\end{proof}
% Fuente francesa: PMIHES_1961__8__5_0, printed pp. 91-92.

\begin{proposition}[4.6.13]
\label{II.4.6.13-fr}
\begin{enumerate}
  \item[\rm{(i)}] Para todo preesquema $Y$, todo $\mathcal{O}_Y$-módulo invertible $\mathcal{L}$ es amplio relativamente al morfismo identidad $1_Y$.
  \item[\rm{(i \emph{bis})}] Sean $f:X\to Y$ un morfismo cuasicompacto y $j:X'\to X$ un morfismo cuasicompacto que es un homeomorfismo del espacio subyacente a $X'$ sobre un subespacio de $X$.
    Si $\mathcal{L}$ es un $\mathcal{O}_X$-módulo amplio relativamente a $f$, entonces $j^*(\mathcal{L})$ es amplio relativamente a $f\circ j$.
  \item[\rm{(ii)}] Sean $Z$ un preesquema cuasicompacto, $f:X\to Y$ y $g:Y\to Z$ dos morfismos cuasicompactos, $\mathcal{L}$ un $\mathcal{O}_X$-módulo amplio relativamente a $f$, y $\mathcal{K}$ un $\mathcal{O}_Y$-módulo amplio relativamente a $g$.
    Entonces existe un entero $n_0>0$ tal que $\mathcal{L}\otimes f^*(\mathcal{K}^{\otimes n})$ es amplio relativamente a $g\circ f$ para todo $n\geq n_0$.
  \item[\rm{(iii)}] Sean $f:X\to Y$ un morfismo cuasicompacto y $g:Y'\to Y$ un morfismo, y pongamos $X'=X_{(Y')}$.
    Si $\mathcal{L}$ es un $\mathcal{O}_X$-módulo amplio relativamente a $f$, entonces $\mathcal{L}'=\mathcal{L}\otimes_Y\mathcal{O}_{Y'}$ es un $\mathcal{O}_{X'}$-módulo amplio relativamente a $f_{(Y')}$.
  \item[\rm{(iv)}] Sean $f_i:X_i\to Y_i$ ($i=1,2$) dos $S$-morfismos cuasicompactos.
    Si $\mathcal{L}_i$ es un $\mathcal{O}_{X_i}$-módulo amplio relativamente a $f_i$ ($i=1,2$), entonces $\mathcal{L}_1\otimes_S\mathcal{L}_2$ es amplio relativamente a $f_1\times_S f_2$.
  \item[\rm{(v)}] Sean $f:X\to Y$ y $g:Y\to Z$ dos morfismos tales que $g\circ f$ sea cuasicompacto.
    Si un
\oldpage[II]{92}
    $\mathcal{O}_X$-módulo $\mathcal{L}$ es amplio relativamente a $g\circ f$, y si $g$ es separado o el espacio subyacente a $X$ es localmente noetheriano, entonces $\mathcal{L}$ es amplio relativamente a $f$.
  \item[\rm{(vi)}] Sean $f:X\to Y$ un morfismo cuasicompacto y $j$ la inyección canónica $X_{\mathrm{red}}\to X$.
    Si $\mathcal{L}$ es un $\mathcal{O}_X$-módulo amplio relativamente a $f$, entonces $j^*(\mathcal{L})$ es amplio relativamente a $f_{\mathrm{red}}$.
\end{enumerate}
\end{proposition}
% Fuente francesa: PMIHES_1961__8__5_0, printed p. 92.

\begin{proof}
Observemos ante todo que (v) y (vi) se deducen de (i), (i~\emph{bis}) y (iv) mediante el mismo razonamiento que en \sref{II.4.4.10-fr}, utilizando \sref{II.4.6.4-fr} en lugar de \sref{II.4.4.5-fr};
dejamos al lector los detalles del razonamiento.
La afirmación~(i) es una consecuencia trivial de \sref{II.4.4.10-fr}[(i)] y \sref{II.4.6.2-fr}.
Para demostrar (i~\emph{bis}), (iii) y (iv), utilizaremos el lema siguiente:
  \begin{lemma}[4.6.13.1]
  \label{II.4.6.13.1-fr}
  \begin{enumerate}
    \item[\rm{(i)}] Sean $u:Z\to S$ un morfismo, $\mathcal{L}$ un $\mathcal{O}_S$-módulo invertible, $s$ una sección de $\mathcal{L}$ sobre $S$, y $s'$ la sección canónicamente correspondiente de $u^*(\mathcal{L})=\mathcal{L}'$ sobre $Z$.
      Entonces $Z_{s'}=u^{-1}(S_s)$.
    \item[\rm{(ii)}] Sean $Z$ y $Z'$ dos $S$-preesquemas, $p$ y $p'$ las proyecciones de $T=Z\times_S Z'$, $\mathcal{L}$ \textup{(resp. $\mathcal{L}'$)} un $\mathcal{O}_Z$-módulo invertible \textup{(resp.} un $\mathcal{O}_{Z'}$-módulo invertible\textup{)}, $t$ \textup{(resp. $t'$)} una sección de $\mathcal{L}$ \textup{(resp. $\mathcal{L}'$)} sobre $Z$ \textup{(resp. $Z'$)}, y $s$ \textup{(resp. $s'$)} la sección correspondiente de $p^*(\mathcal{L})$ \textup{(resp. $p'^*(\mathcal{L}')$)} sobre $Z\times_S Z'$.
      Entonces $T_{s\otimes s'}=Z_t\times_S Z'_{t'}$.
  \end{enumerate}
  \end{lemma}

  \begin{proof}
  Se sigue de las definiciones que se puede reducir al caso en que todos los preesquemas considerados son afines.
  Además, en (i), se puede suponer que $\mathcal{L}=\mathcal{O}_S$;
  la afirmación~(i) se sigue entonces inmediatamente de \sref[I]{I.1.2.2.2-fr}.
  Del mismo modo, en (ii), se puede reducir al caso en que $\mathcal{L}=\mathcal{O}_Z$ y $\mathcal{L}'=\mathcal{O}_{Z'}$, y la afirmación se reduce entonces al lema~\sref{II.4.3.2.4-fr}.
  \end{proof}

Demostremos ahora (i~\emph{bis}).
Se puede suponer que $Y$ es afín \sref{II.4.6.4-fr}, y por tanto que $\mathcal{L}$ es amplio \sref{II.4.6.6-fr};
cuando $s$ recorre la unión de los conjuntos $\Gamma(X,\mathcal{L}^{\otimes n})$ ($n>0$), los $X_s$ forman una base de la topología de $X$ \sref{II.4.5.2-fr}[\emph{(a)}], y por tanto, por hipótesis, los $j^{-1}(X_s)$ forman una base de la topología de $X'$;
se concluye así, por el lema~\sref{II.4.6.13.1-fr}[(i)] y \sref{II.4.5.2-fr}[\emph{(a)}], que $j^*(\mathcal{L})$ es amplio.

Demostremos a continuación (iii).
Se puede suponer de nuevo que $Y$ e $Y'$ son afines \sref{II.4.6.4-fr}, de donde se sigue que la proyección $h:X'\to X$ es afín \sref{II.1.5.5-fr}.
Como $\mathcal{L}$ es amplio \sref{II.4.6.6-fr}, cuando $s$ recorre el conjunto de las secciones de los $\mathcal{L}^{\otimes n}$ ($n>0$) sobre $X$ tales que $X_s$ es afín, los $X_s$ cubren $X$ \sref{II.4.5.2-fr}[\emph{(a')}], y por tanto los $h^{-1}(X_s)$ son afines \sref{II.1.2.5-fr} y cubren $X'$;
se sigue de nuevo del lema~\sref{II.4.6.13.1-fr}[(i)] y de \sref{II.4.5.2-fr}[\emph{(a')}] que $\mathcal{L}'$ es amplio, ya que el morfismo $f_{(Y')}$ es cuasicompacto \sref[I]{I.6.6.4-fr}[(iii)].
% Fuente francesa: PMIHES_1961__8__5_0, printed pp. 92-93.

Para demostrar (iv), observemos ante todo que $f_1\times_S f_2$ es cuasicompacto \sref[I]{I.6.6.4-fr}[(iv)].
Se puede suponer además que $S$, $Y_1$ e $Y_2$ son afines (\sref{II.4.6.4-fr} y \sref[I]{I.3.2.7-fr}), y por tanto que $\mathcal{L}_i$ es amplio ($i=1,2$) \sref{II.4.6.6-fr}.
Los abiertos $(X_1)_{s_1}\times_S(X_2)_{s_2}$ forman una cubierta de $X_1\times_S X_2$, donde $s_i$ recorre las secciones de $\mathcal{L}_i^{\otimes n_i}$ tales que $(X_i)_{s_i}$ es afín \sref{II.4.5.2-fr}[\emph{(a')}].
Reemplazando $s_1$ y $s_2$ por potencias adecuadas, lo cual no cambia los $(X_i)_{s_i}$, se puede suponer que $n_1=n_2$.
Se deduce entonces de \sref{II.4.6.13.1-fr}[(ii)] y \sref{II.4.5.2-fr}[\emph{(a')}] que $\mathcal{L}_1\otimes_S\mathcal{L}_2$ es amplio, de donde se sigue la afirmación, ya que $Y_1\times_S Y_2$ es afín \sref{II.4.6.6-fr}.

Solo queda demostrar (ii).
Mediante el mismo razonamiento que en \sref{II.4.4.10-fr}, pero utilizando aquí \sref{II.4.6.4-fr}, se puede reducir al caso en que $Z$ es afín.
Como $\mathcal{K}$ es entonces amplio y $Y$ cuasicompacto, existe un número finito de secciones $s_i\in\Gamma(Y,\mathcal{K}^{\otimes k_i})$ tales que los $Y_{s_i}$
\oldpage[II]{93}
son \emph{afines} y cubren $Y$ \sref{II.4.5.2-fr}[\emph{(a')}];
reemplazando los $s_i$ por potencias adecuadas, se puede suponer además que todos los $k_i$ son iguales a un mismo entero $k$.
Sean $s'_i$ las secciones de $f^*(\mathcal{K}^{\otimes k})$ sobre $X$ que corresponden canónicamente a los $s_i$, de modo que los $X_{s'_i}=f^{-1}(Y_{s_i})$ \sref{II.4.6.13.1-fr}[(i)] cubren $X$.
Como $\mathcal{L}|X_{s'_i}$ es amplio (\sref{II.4.6.4-fr} y \sref{II.4.6.6-fr}), existe, para cada $i$, un número finito de secciones $t_{ij}\in\Gamma(X,\mathcal{L}^{\otimes n_{ij}})$ tales que los $X_{t_{ij}}$ son afines, están contenidos en $X_{s'_i}$ y cubren $X_{s'_i}$ \sref{II.4.5.2-fr}[\emph{(a')}];
se puede suponer asimismo que todos los $n_{ij}$ son iguales a un mismo entero $n$.
Dicho esto, $X$ es separado y cuasicompacto, y existen, por tanto, un entero $m>0$ y, para cada $(i,j)$, una sección
\[
  u_{ij}\in\Gamma(X,\mathcal{L}^{\otimes n}\otimes_X
  f^*(\mathcal{K}^{\otimes mk}))
\]
tal que $t_{ij}\otimes s_i'^{\otimes m}$ es la restricción a $X_{s'_i}$ de $u_{ij}$ \sref[I]{I.9.3.1-fr};
además, $X_{u_{ij}}=X_{t_{ij}}$, y por tanto los $X_{u_{ij}}$ son afines y cubren $X$.
Se puede suponer también que $m$ es de la forma $nr$;
si se pone $n_0=rk$, se ve por \sref{II.4.5.2-fr}[\emph{(a')}] que $\mathcal{L}\otimes_{\mathcal{O}_X}f^*(\mathcal{K}^{\otimes n_0})$ es amplio.
Además, existe $h_0>0$ tal que $\mathcal{K}^{\otimes h}$ está generado por sus secciones sobre $Y$ para todo $h\geq h_0$ \sref{II.4.5.5-fr};
\emph{a fortiori}, $f^*(\mathcal{K}^{\otimes h})$ está generado por sus secciones sobre $X$ para todo $h\geq h_0$, por definición de las imágenes inversas (\sref[0]{0.3.7.1-fr} y \sref{0.4.4.1-fr}).
Se concluye así que $\mathcal{L}\otimes f^*(\mathcal{K}^{\otimes n_0+h})$ es amplio para todo $h\geq h_0$ \sref{II.4.5.6-fr}, lo que termina la demostración.
\end{proof}
% Fuente francesa: PMIHES_1961__8__5_0, printed p. 93.

\begin{remark}[4.6.14]
\label{II.4.6.14-fr}
Bajo las condiciones de (ii), no debe concluirse que $\mathcal{L}\otimes f^*(\mathcal{K})$ es amplio para $g\circ f$;
en efecto, como $\mathcal{L}\otimes f^*(\mathcal{K}^{-1})$ también es amplio para $f$ \sref{II.4.6.5-fr}, se concluiría que $\mathcal{L}$ es amplio para $g\circ f$;
tomando en particular como $g$ el morfismo identidad, \emph{todo} $\mathcal{O}_X$-módulo invertible sería amplio para $f$, lo cual no ocurre en general (véanse \sref{II.5.1.6-fr}, \sref{II.5.3.4-fr}[(i)] y \sref{II.5.3.1-fr}).
\end{remark}

\begin{proposition}[4.6.15]
\label{II.4.6.15-fr}
Sean $f:X\to Y$ un morfismo cuasicompacto, $\mathcal{J}$ un haz cuasicoherente localmente nilpotente de ideales de $\mathcal{O}_X$, $Z$ el subpreesquema cerrado de $X$ definido por $\mathcal{J}$, y $j:Z\to X$ la inyección canónica.
Para que un $\mathcal{O}_X$-módulo invertible $\mathcal{L}$ sea amplio para $f$, es necesario y suficiente que $j^*(\mathcal{L})$ sea amplio para $f\circ j$.
\end{proposition}

\begin{proof}
Como la cuestión es local sobre $Y$ \sref{II.4.6.4-fr}, se puede suponer que $Y$ es afín;
como $X$ es entonces cuasicompacto, se puede suponer que $\mathcal{J}$ es nilpotente.
Teniendo en cuenta \sref{II.4.6.6-fr}, la proposición coincide entonces exactamente con \sref{II.4.5.13-fr}.
\end{proof}

\begin{corollary}[4.6.16]
\label{II.4.6.16-fr}
Sean $X$ un preesquema localmente noetheriano y $f:X\to Y$ un morfismo cuasicompacto.
Para que un $\mathcal{O}_X$-módulo invertible $\mathcal{L}$ sea amplio para $f$, es necesario y suficiente que su imagen inversa $\mathcal{L}'$ por la inyección canónica $X_{\mathrm{red}}\to X$ sea amplia para $f_{\mathrm{red}}$.
\end{corollary}

\begin{proof}
Ya se ha visto que la condición es necesaria \sref{II.4.6.13-fr}[(vi)];
recíprocamente, si se satisface, se puede reducir, para demostrar que $\mathcal{L}$ es amplio para $f$, al caso en que $Y$ es afín \sref{II.4.6.4-fr};
entonces $Y_{\mathrm{red}}$ también es afín, y por tanto $\mathcal{L}'$ es amplio \sref{II.4.6.6-fr}, y también lo es $\mathcal{L}$ por \sref{II.4.5.13-fr}, ya que $X$ es entonces noetheriano y $X_{\mathrm{red}}$ un subpreesquema cerrado de $X$ definido por un haz cuasicoherente nilpotente de ideales \sref[I]{I.6.1.6-fr}.
\end{proof}
% Fuente francesa: PMIHES_1961__8__5_0, printed pp. 93-94.

\begin{proposition}[4.6.17]
\label{II.4.6.17-fr}
Con la notación y las hipótesis de \sref{II.4.4.11-fr}, para que $\mathcal{L}''$ sea amplio relativamente a $f''$, es necesario y suficiente que $\mathcal{L}$ sea amplio relativamente a $f$ y $\mathcal{L}'$ amplio relativamente a $f'$.
\end{proposition}

\oldpage[II]{94}
\begin{proof}
La necesidad de la condición se sigue de \sref{II.4.6.13-fr}[(i~\emph{bis})].
Para ver que la condición es suficiente, se puede reducir al caso en que $Y$ es afín, y entonces el hecho de que $\mathcal{L}''$ es amplio se sigue del criterio~\sref{II.4.5.2-fr}[\emph{(a)}] aplicado a $\mathcal{L}$, $\mathcal{L}'$ y $\mathcal{L}''$, observando que una sección de $\mathcal{L}$ sobre $X$ puede prolongarse (por $0$) a una sección de $\mathcal{L}''$ sobre $X''$.
\end{proof}

\begin{proposition}[4.6.18]
\label{II.4.6.18-fr}
Sean $Y$ un preesquema cuasicompacto, $\mathcal{S}$ una $\mathcal{O}_Y$-álgebra graduada cuasicoherente de tipo finito, $X=\operatorname{Proj}(\mathcal{S})$, y $f:X\to Y$ el morfismo estructural.
Entonces $f$ es de tipo finito, y existe un entero $d>0$ tal que $\mathcal{O}_X(d)$ es invertible y $f$-amplio.
\end{proposition}

\begin{proof}
Por \sref{II.3.1.10-fr}, existe un entero $d>0$ tal que $\mathcal{S}^{(d)}$ está generada por $\mathcal{S}_d$.
Se sabe que, bajo el isomorfismo canónico entre $X$ y $X^{(d)}=\operatorname{Proj}(\mathcal{S}^{(d)})$, $\mathcal{O}_X(d)$ se identifica con $\mathcal{O}_{X^{(d)}}(1)$ \sref{II.3.2.9-fr}[(ii)].
Se ve así que se puede reducir al caso en que $\mathcal{S}$ está generada por $\mathcal{S}_1$;
la proposición se sigue entonces de \sref{II.4.4.3-fr} y \sref{II.4.6.2-fr} (teniendo en cuenta que $f$ es un morfismo de tipo finito \sref{II.3.4.1-fr}).
\end{proof}
% Fuente francesa: PMIHES_1961__8__5_0, printed pp. 94-95.

\section{Morfismos cuasiafines; morfismos cuasiproyectivos; morfismos propios; morfismos proyectivos}
\label{section:II.5-fr}

\subsection{Morfismos cuasiafines}
\label{subsection:II.5.1-fr}

\begin{definition}[5.1.1]
\label{II.5.1.1-fr}
Se llama esquema cuasiafín a todo esquema isomorfo a un subesquema inducido sobre un abierto cuasicompacto de un esquema afín.
Se dice que un morfismo $f:X\to Y$ es cuasiafín, o también que $X$ (considerado como $Y$-preesquema mediante $f$) es un $Y$-esquema cuasiafín, si existe una cubierta $(U_\alpha)$ de $Y$ por abiertos afines tal que los $f^{-1}(U_\alpha)$ son esquemas cuasiafines.
\end{definition}

Está claro que todo morfismo cuasiafín es \emph{separado} (\sref[I]{I.5.5.5-fr} y \sref[I]{I.5.5.8-fr}) y \emph{cuasicompacto} \sref[I]{I.6.6.1-fr};
todo morfismo afín es evidentemente cuasiafín.

Recordemos que, para todo preesquema $X$, si se pone $A=\Gamma(X,\mathcal{O}_X)$, el homomorfismo identidad $A\to A=\Gamma(X,\mathcal{O}_X)$ define un morfismo $X\to\operatorname{Spec}(A)$, llamado \emph{canónico} \sref[I]{I.2.2.4-fr};
no es sino el morfismo canónico definido en \sref{II.4.5.1-fr} para el caso particular en que $\mathcal{L}=\mathcal{O}_X$, si se recuerda que $\operatorname{Proj}(A[T])$ se identifica canónicamente con $\operatorname{Spec}(A)$ \sref{II.3.1.7-fr}.

\begin{proposition}[5.1.2]
\label{II.5.1.2-fr}
Sean $X$ un esquema cuasicompacto o un preesquema cuyo espacio subyacente sea noetheriano, y $A$ el anillo $\Gamma(X,\mathcal{O}_X)$.
Las condiciones siguientes son equivalentes:
\begin{enumerate}
  \item[{\rm(a)}] $X$ es un esquema cuasiafín.
  \item[{\rm(b)}] El morfismo canónico $u:X\to\operatorname{Spec}(A)$ es una inmersión abierta.
  \item[{\rm(b')}] El morfismo canónico $u:X\to\operatorname{Spec}(A)$ es un homeomorfismo de $X$ sobre un subespacio del espacio subyacente a $\operatorname{Spec}(A)$.
  \item[{\rm(c)}] El $\mathcal{O}_X$-módulo $\mathcal{O}_X$ es muy amplio relativamente a $u$ \sref{II.4.4.2-fr}.
  \item[{\rm(c')}] El $\mathcal{O}_X$-módulo $\mathcal{O}_X$ es amplio \sref{II.4.5.1-fr}.
  \item[{\rm(d)}] Cuando $f$ recorre $A$, los $X_f$ forman una base de la topología de $X$.
  \item[{\rm(d')}] Cuando $f$ recorre $A$, aquellos $X_f$ que son afines forman una cubierta de $X$.
\oldpage[II]{95}
  \item[{\rm(e)}] Todo $\mathcal{O}_X$-módulo cuasicoherente está generado por sus secciones sobre $X$.
  \item[{\rm(e')}] Todo haz cuasicoherente de ideales de $\mathcal{O}_X$ de tipo finito está generado por sus secciones sobre $X$.
\end{enumerate}
\end{proposition}

\begin{proof}
Está claro que \emph{(b)} implica \emph{(a)}, y que \emph{(a)} implica \emph{(c)} por \sref{II.4.4.4-fr}[\emph{(b)}] aplicado al morfismo identidad (teniendo en cuenta la observación que precede a esta proposición);
además, \emph{(c)} implica \emph{(c')} \sref{II.4.5.10-fr}[(i)], y \emph{(c')}, \emph{(b)} y \emph{(b')} son equivalentes por (\hyperref[II.4.5.2-fr]{4.5.2, criterios \emph{b)} y \emph{b')}}).
Finalmente, \emph{(c')} equivale a cada una de las condiciones \emph{(d)}, \emph{(d')}, \emph{(e)} y \emph{(e')} por (\hyperref[II.4.5.2-fr]{4.5.2, criterios \emph{a)}, \emph{a')}, \emph{c)}}) y (\hyperref[II.4.5.5-fr]{4.5.5, criterio \emph{d'')}}).
\end{proof}

Se observa además que, con la notación anterior, aquellos $X_f$ que son afines forman una \emph{base} de la topología de $X$, y que el morfismo canónico $u$ es \emph{dominante} \sref{II.4.5.2-fr}.
% Fuente francesa: PMIHES_1961__8__5_0, printed p. 95.

\begin{corollary}[5.1.3]
\label{II.5.1.3-fr}
Sea $X$ un preesquema cuasicompacto.
Si existe un morfismo $v:X\to Y$ de $X$ a un esquema afín $Y$ que es un homeomorfismo de $X$ sobre un subespacio abierto de $Y$, entonces $X$ es cuasiafín.
\end{corollary}

\begin{proof}
Existe una familia $(g_\alpha)$ de secciones de $\mathcal{O}_Y$ sobre $Y$ tal que los $D(g_\alpha)$ forman una base de la topología de $v(X)$;
si $v=(\psi,\theta)$ y se pone $f_\alpha=\theta(g_\alpha)$, entonces $X_{f_\alpha}=\psi^{-1}(D(g_\alpha))$ \sref[I]{I.2.2.4.1-fr}, de modo que los $X_{f_\alpha}$ forman una base de la topología de $X$, y se satisface el criterio~\sref{II.5.1.2-fr}[\emph{(d)}].
\end{proof}

\begin{corollary}[5.1.4]
\label{II.5.1.4-fr}
Si $X$ es un esquema cuasiafín, entonces \emph{todo} $\mathcal{O}_X$-módulo invertible es muy amplio (relativamente al morfismo canónico) y, \emph{a fortiori}, amplio.
\end{corollary}

\begin{proof}
Un módulo tal $\mathcal{L}$ está generado por sus secciones sobre $X$ \sref{II.5.1.2-fr}[\emph{(e)}], de modo que $\mathcal{L}\otimes\mathcal{O}_X=\mathcal{L}$ es muy amplio \sref{II.4.4.8-fr}.
\end{proof}

\begin{corollary}[5.1.5]
\label{II.5.1.5-fr}
Sea $X$ un preesquema cuasicompacto.
Si existe un $\mathcal{O}_X$-módulo invertible $\mathcal{L}$ tal que $\mathcal{L}$ y $\mathcal{L}^{-1}$ son amplios, entonces $X$ es un esquema cuasiafín.
\end{corollary}

\begin{proof}
En efecto, $\mathcal{O}_X=\mathcal{L}\otimes\mathcal{L}^{-1}$ es entonces amplio \sref{II.4.5.7-fr}.
\end{proof}
% Fuente francesa: PMIHES_1961__8__5_0, printed pp. 95-96.

\begin{proposition}[5.1.6]
\label{II.5.1.6-fr}
Sea $f:X\to Y$ un morfismo cuasicompacto.
Las condiciones siguientes son equivalentes:
\begin{enumerate}
  \item[{\rm(a)}] El morfismo $f$ es cuasiafín.
  \item[{\rm(b)}] La $\mathcal{O}_Y$-álgebra $f_*(\mathcal{O}_X)=\mathcal{A}(X)$ es cuasicoherente, y el morfismo canónico $X\to\operatorname{Spec}(\mathcal{A}(X))$ correspondiente al homomorfismo identidad $\mathcal{A}(X)\to\mathcal{A}(X)$ \sref{II.1.2.7-fr} es una inmersión abierta.
  \item[{\rm(b')}] La $\mathcal{O}_Y$-álgebra $\mathcal{A}(X)$ es cuasicoherente, y el morfismo canónico $X\to\operatorname{Spec}(\mathcal{A}(X))$ es un homeomorfismo de $X$ sobre un subespacio de $\operatorname{Spec}(\mathcal{A}(X))$.
  \item[{\rm(c)}] El $\mathcal{O}_X$-módulo $\mathcal{O}_X$ es muy amplio para $f$.
  \item[{\rm(c')}] El $\mathcal{O}_X$-módulo $\mathcal{O}_X$ es amplio para $f$.
  \item[{\rm(d)}] El morfismo $f$ es separado y, para todo $\mathcal{O}_X$-módulo cuasicoherente $\mathcal{F}$, el homomorfismo canónico $\sigma:f^*(f_*(\mathcal{F}))\to\mathcal{F}$ \sref[0]{0.4.4.3-fr} es sobreyectivo.
\end{enumerate}

Además, cuando $f$ es cuasiafín, todo $\mathcal{O}_X$-módulo invertible $\mathcal{L}$ es muy amplio relativamente a $f$.
\end{proposition}

\begin{proof}
La equivalencia de \emph{(a)} y \emph{(c')} se sigue del carácter local sobre $Y$ de la $f$-amplitud \sref{II.4.6.4-fr}, de la definición~\sref{II.5.1.1-fr} y del criterio~\sref{II.5.1.2-fr}[\emph{(c')}].
Las demás propiedades son locales sobre $Y$
\oldpage[II]{96}
y se siguen, por tanto, inmediatamente de \sref{II.5.1.2-fr} y \sref{II.5.1.4-fr}, teniendo en cuenta que $f_*(\mathcal{F})$ es cuasicoherente cuando $f$ es separado \sref[I]{I.9.2.2-fr}[\emph{(a)}].
\end{proof}
% Fuente francesa: PMIHES_1961__8__5_0, printed p. 96.

\begin{corollary}[5.1.7]
\label{II.5.1.7-fr}
Sea $f:X\to Y$ un morfismo cuasiafín.
Para todo abierto $U$ de $Y$, la restricción $f^{-1}(U)\to U$ de $f$ es cuasiafín.
\end{corollary}

\begin{corollary}[5.1.8]
\label{II.5.1.8-fr}
Sean $Y$ un esquema afín y $f:X\to Y$ un morfismo cuasicompacto.
Para que $f$ sea cuasiafín, es necesario y suficiente que $X$ sea un esquema cuasiafín.
\end{corollary}

\begin{proof}
Es una consecuencia inmediata de \sref{II.5.1.6-fr} y \sref{II.4.6.6-fr}.
\end{proof}

\begin{corollary}[5.1.9]
\label{II.5.1.9-fr}
Sean $Y$ un esquema cuasicompacto o un preesquema cuyo espacio subyacente sea noetheriano, y $f:X\to Y$ un morfismo de \emph{tipo finito}.
Si $f$ es cuasiafín, existe una sub-$\mathcal{O}_Y$-álgebra cuasicoherente $\mathcal{B}$ de $\mathcal{A}(X)=f_*(\mathcal{O}_X)$ de \emph{tipo finito} \sref[I]{I.9.6.2-fr} tal que el morfismo $X\to\operatorname{Spec}(\mathcal{B})$ correspondiente a la inyección canónica $\mathcal{B}\to\mathcal{A}(X)$ \textup{(\hyperref[II.1.2.7-fr]{1.2.7})} es una inmersión.
Además, toda sub-$\mathcal{O}_Y$-álgebra cuasicoherente $\mathcal{B}'$ de tipo finito de $\mathcal{A}(X)$ que contenga a $\mathcal{B}$ posee la misma propiedad.
\end{corollary}

\begin{proof}
En efecto, $\mathcal{A}(X)$ es el límite inductivo de sus sub-$\mathcal{O}_Y$-álgebras cuasicoherentes de tipo finito \sref[I]{I.9.6.5-fr};
el resultado es entonces un caso particular de \sref{II.3.8.4-fr}, teniendo en cuenta la identificación de $\operatorname{Spec}(\mathcal{A}(X))$ con $\operatorname{Proj}(\mathcal{A}(X)[T])$ \sref{II.3.1.7-fr}.
\end{proof}
% Fuente francesa: PMIHES_1961__8__5_0, printed p. 96.

\begin{proposition}[5.1.10]
\label{II.5.1.10-fr}
\medskip\noindent
\begin{enumerate}
  \item[{\rm(i)}] Todo morfismo cuasicompacto $X\to Y$ que sea un homeomorfismo del espacio subyacente a $X$ sobre un subespacio del espacio subyacente a $Y$ (en particular, toda inmersión cerrada) es cuasiafín.
  \item[{\rm(ii)}] El compuesto de dos morfismos cuasiafines es cuasiafín.
  \item[{\rm(iii)}] Si $f:X\to Y$ es un $S$-morfismo cuasiafín, entonces $f_{(S')}:X_{(S')}\to Y_{(S')}$ es un morfismo cuasiafín para toda extensión $S'\to S$ del preesquema de base.
  \item[{\rm(iv)}] Si $f:X\to Y$ y $g:X'\to Y'$ son dos $S$-morfismos cuasiafines, entonces $f\times_S g$ es cuasiafín.
  \item[{\rm(v)}] Si $f:X\to Y$ y $g:Y\to Z$ son dos morfismos tales que $g\circ f$ es cuasiafín, y si $g$ es separado o el espacio subyacente a $X$ es localmente noetheriano, entonces $f$ es cuasiafín.
  \item[{\rm(vi)}] Si $f$ es un morfismo cuasiafín, también lo es $f_{\mathrm{red}}$.
\end{enumerate}
\end{proposition}

\begin{proof}
Teniendo en cuenta el criterio~\sref{II.5.1.6-fr}[\emph{(c')}], las afirmaciones (i), (iii), (iv), (v) y (vi) se siguen inmediatamente, respectivamente, de (\hyperref[II.4.6.13-fr]{4.6.13, (i \emph{bis}), (iii), (iv), (v) y (vi)}).
Para demostrar (ii), se puede reducir al caso en que $Z$ es afín, y entonces la afirmación se sigue directamente de aplicar \sref{II.4.6.13-fr}[(ii)] a $\mathcal{L}=\mathcal{O}_X$ y $\mathcal{K}=\mathcal{O}_Y$.
\end{proof}

\begin{remark}[5.1.11]
\label{II.5.1.11-fr}
Sean $f:X\to Y$ y $g:Y\to Z$ dos morfismos tales que $X\times_ZY$ sea localmente noetheriano.
Entonces la inmersión gráfica $\Gamma_f:X\to X\times_ZY$ es cuasiafín, pues es cuasicompacta \sref[I]{I.6.3.5-fr}, y \sref[I]{I.5.5.12-fr} muestra que, en (v), la conclusión sigue siendo válida si se suprime la hipótesis de que $g$ es separado.
\end{remark}

\begin{proposition}[5.1.12]
\label{II.5.1.12-fr}
Sean $f:X\to Y$ un morfismo cuasicompacto y $g:X'\to X$ un morfismo cuasiafín.
Si $\mathcal{L}$ es un $\mathcal{O}_X$-módulo amplio para $f$, entonces $g^*(\mathcal{L})$ es un $\mathcal{O}_{X'}$-módulo amplio para $f\circ g$.
\end{proposition}

\begin{proof}
Como $\mathcal{O}_{X'}$ es muy amplio para $g$ y la cuestión es local sobre $Y$ \sref{II.4.6.4-fr}, se sigue de \sref{II.4.6.13-fr}[(ii)] que existe (para $Y$ afín) un entero $n$ tal que
\[
  g^*(\mathcal{L}^{\otimes n})=(g^*(\mathcal{L}))^{\otimes n}
\]
es amplio para $f\circ g$, y por tanto $g^*(\mathcal{L})$ es amplio para $f\circ g$ \sref{II.4.6.9-fr}.
\end{proof}
% Fuente francesa: PMIHES_1961__8__5_0, printed p. 97.

\oldpage[II]{97}
\subsection{El criterio de Serre}
\label{subsection:II.5.2-fr}

\begin{theorem}[5.2.1]
\label{II.5.2.1-fr}
\emph{(Criterio de Serre).}
Sea $X$ un esquema cuasicompacto o un preesquema cuyo espacio subyacente sea noetheriano.
Las condiciones siguientes son equivalentes:
\begin{enumerate}
  \item[{\rm(a)}] $X$ es un esquema afín.
  \item[{\rm(b)}] Existe una familia de elementos $f_\alpha\in A=\Gamma(X,\mathcal{O}_X)$ tal que los $X_{f_\alpha}$ son afines y el ideal generado por los $f_\alpha$ en $A$ es igual a $A$.
  \item[{\rm(c)}] El funtor $\Gamma(X,\mathcal{F})$ es exacto en $\mathcal{F}$ en la categoría de los $\mathcal{O}_X$-módulos cuasicoherentes; dicho de otro modo, si
    \[
    0\to\mathcal{F}'\to\mathcal{F}\to\mathcal{F}''\to0 \tag{*}
  \]
    es una sucesión exacta de $\mathcal{O}_X$-módulos cuasicoherentes, entonces la sucesión
    \[
    0\to\Gamma(X,\mathcal{F}')\to\Gamma(X,\mathcal{F})
      \to\Gamma(X,\mathcal{F}'')\to0
  \]
    también es exacta.
  \item[{\rm(c')}] La propiedad~{\rm(c)} se satisface para toda sucesión exacta \textup{(*)} de $\mathcal{O}_X$-módulos cuasicoherentes tal que $\mathcal{F}$ es isomorfo a un $\mathcal{O}_X$-submódulo de un producto finito $\mathcal{O}_X^n$.
  \item[{\rm(d)}] $\mathrm{H}^1(X,\mathcal{F})=0$ para todo $\mathcal{O}_X$-módulo cuasicoherente $\mathcal{F}$.
  \item[\rm{d')}] $\mathrm{H}^1(X,\mathcal{J})=0$ para todo haz cuasicoherente de ideales $\mathcal{J}$ de $\mathcal{O}_X$.
\end{enumerate}
\end{theorem}
% Fuente francesa: PMIHES_1961__8__5_0, printed pp. 97-98.

\begin{proof}
Es evidente que \emph{(a)} implica \emph{(b)};
por otra parte, \emph{(b)} implica que los $X_{f_\alpha}$ cubren $X$, pues, por hipótesis, la sección $1$ es combinación lineal de los $f_\alpha$, y los $D(f_\alpha)$ cubren entonces $\operatorname{Spec}(A)$.
La última afirmación de \sref{II.4.5.2-fr} implica, por tanto, que $X\to\operatorname{Spec}(A)$ es un isomorfismo.

Se sabe que \emph{(a)} implica \emph{(c)} \sref[I]{I.1.3.11-fr}, y que \emph{(c)} implica trivialmente \emph{(c')}.
Demostremos ahora que \emph{(c')} implica \emph{(b)}.
Ante todo, \emph{(c')} implica que, para todo punto \emph{cerrado} $x\in X$ y toda vecindad abierta $U$ de $x$, existe $f\in A$ tal que $x\in X_f\subset U$.
Sea $\mathcal{J}$ (resp. $\mathcal{J}'$) el haz cuasicoherente de ideales de $\mathcal{O}_X$ que define el subpreesquema cerrado reducido de $X$ cuyo espacio subyacente es $X\setmin U$ (resp. $(X\setmin U)\cup\{x\}$) \sref[I]{I.5.2.1-fr};
está claro que $\mathcal{J}'\subset\mathcal{J}$, y que $\mathcal{J}''=\mathcal{J}/\mathcal{J}'$ es un $\mathcal{O}_X$-módulo cuasicoherente cuyo soporte es $\{x\}$ y tal que $\mathcal{J}''_x=\boldsymbol{k}(x)$.
La hipótesis~\emph{(c')} aplicada a la sucesión exacta $0\to\mathcal{J}'\to\mathcal{J}\to\mathcal{J}''\to0$ muestra que $\Gamma(X,\mathcal{J})\to\Gamma(X,\mathcal{J}'')$ es sobreyectivo.
La sección de $\mathcal{J}''$ cuyo germen en $x$ es $1_x$ es, por tanto, la imagen de una sección $f\in\Gamma(X,\mathcal{J})\subset\Gamma(X,\mathcal{O}_X)$, y se tiene, por definición, $f(x)=1_x$ y $f(y)=0$ en $X\setmin U$, lo que demuestra nuestra afirmación.
Además, si $U$ es afín, también lo es $X_f$ \sref[I]{I.1.3.6-fr}; por tanto, la unión de los $X_f$ que son afines ($f\in A$) es un abierto $Z$ que contiene \emph{todos los puntos cerrados} de $X$;
como $X$ es un espacio de Kolmogórov cuasicompacto, necesariamente $Z=X$ \sref[0]{0.2.1.3-fr}.
Como $X$ es cuasicompacto, existe un número \emph{finito} de elementos $f_i\in A$ ($1\leq i\leq n$) tales que los $X_{f_i}$ son afines y cubren $X$.
Consideremos entonces el homomorfismo $\mathcal{O}_X^n\to\mathcal{O}_X$ definido por las secciones $f_i$ \sref[0]{0.5.1.1-fr};
como, para todo $x\in X$, al menos uno de los $(f_i)_x$ es invertible, este homomorfismo es \emph{sobreyectivo}, y se tiene, por tanto, una sucesión exacta $0\to\mathcal{R}\to\mathcal{O}_X^n\to\mathcal{O}_X\to0$, donde $\mathcal{R}$ es un $\mathcal{O}_X$-submódulo cuasicoherente de $\mathcal{O}_X^n$.
Se sigue entonces
\oldpage[II]{98}
de (\emph{c')} que el homomorfismo correspondiente $\Gamma(X,\mathcal{O}_X^n)\to\Gamma(X,\mathcal{O}_X)$ es sobreyectivo, lo que demuestra \emph{(b)}.

Finalmente, \emph{(a)} implica \emph{(d)} \sref[I]{I.5.1.9.2-fr}, y \emph{(d)} implica trivialmente \emph{(d')}.
Queda por demostrar que \emph{(d')} implica \emph{(c')}.
Pero, si $\mathcal{F}'$ es un $\mathcal{O}_X$-submódulo cuasicoherente de $\mathcal{O}_X^n$, la filtración $0\subset\mathcal{O}_X\subset\mathcal{O}_X^2\subset\ldots\subset
\mathcal{O}_X^n$ define una filtración de $\mathcal{F}'$ dada por los $\mathcal{F}'_k=\mathcal{F}'\cap\mathcal{O}_X^k$ ($0\leq k\leq n$), que son $\mathcal{O}_X$-módulos cuasicoherentes \sref[I]{I.4.1.1-fr}, y $\mathcal{F}'_{k+1}/\mathcal{F}'_k$ es isomorfo a un $\mathcal{O}_X$-submódulo cuasicoherente de $\mathcal{O}_X^{k+1}/\mathcal{O}_X^k=\mathcal{O}_X$, es decir, a un haz cuasicoherente de ideales de $\mathcal{O}_X$.
La hipótesis~\emph{(d')} implica, por tanto, que $\mathrm{H}^1(X,\mathcal{F}'_{k+1}/\mathcal{F}'_k)=0$;
la sucesión exacta de cohomología $\mathrm{H}^1(X,\mathcal{F}'_k)\to\mathrm{H}^1(X,\mathcal{F}'_{k+1})
\to\mathrm{H}^1(X,\mathcal{F}'_{k+1}/\mathcal{F}'_k)=0$ permite entonces demostrar por inducción sobre $k$ que $\mathrm{H}^1(X,\mathcal{F}'_k)=0$ para todo $k$.
\end{proof}
% Fuente francesa: PMIHES_1961__8__5_0, printed p. 98.

\begin{remark}[5.2.1.1]
\label{II.5.2.1.1-fr}
Cuando $X$ es un preesquema \emph{noetheriano}, puede reemplazarse ``cuasicoherente'' por ``coherente'' en los enunciados de \emph{(c')} y \emph{(d')}.
En efecto, en la demostración de que \emph{(c')} implica \emph{(b)}, $\mathcal{J}$ y $\mathcal{J}'$ son entonces haces \emph{coherentes} de ideales y, además, todo submódulo cuasicoherente de un módulo coherente es coherente \sref[I]{I.6.1.1-fr};
de donde se sigue la conclusión.
\end{remark}
% Fuente francesa: PMIHES_1961__8__5_0, printed p. 98.

\begin{corollary}[5.2.2]
\label{II.5.2.2-fr}
Sea $f:X\to Y$ un morfismo separado cuasicompacto.
Las condiciones siguientes son equivalentes:
\begin{enumerate}
  \item[{\rm(a)}] El morfismo $f$ es afín.
  \item[{\rm(b)}] El funtor $f_*$ es exacto en la categoría de los $\mathcal{O}_X$-módulos cuasicoherentes.
  \item[{\rm(c)}] Para todo $\mathcal{O}_X$-módulo cuasicoherente $\mathcal{F}$, se tiene $\mathrm{R}^1f_*(\mathcal{F})=0$.
  \item[{\rm(c')}] Para todo haz cuasicoherente de ideales $\mathcal{J}$ de $\mathcal{O}_X$, se tiene $\mathrm{R}^1f_*(\mathcal{J})=0$.
\end{enumerate}
\end{corollary}

\begin{proof}
Todas estas condiciones son locales sobre $Y$, por definición del funtor $\mathrm{R}^1f_*$ (T,~3.7.3), de modo que se puede suponer que $Y$ es afín.
Si $f$ es afín, entonces $X$ es afín, y la propiedad~\emph{(b)} no es sino \sref[I]{I.1.6.4-fr}.
Recíprocamente, demostremos que \emph{(b)} implica \emph{(a)}:
para todo $\mathcal{O}_X$-módulo cuasicoherente $\mathcal{F}$, $f_*(\mathcal{F})$ es un $\mathcal{O}_Y$-módulo cuasicoherente \sref[I]{I.9.2.2-fr}[\emph{(a)}].
Por hipótesis, el funtor $f_*(\mathcal{F})$ es exacto en $\mathcal{F}$, y el funtor $\Gamma(Y,\mathcal{G})$ es exacto en $\mathcal{G}$ (en la categoría de los $\mathcal{O}_Y$-módulos cuasicoherentes), ya que $Y$ es afín \sref[I]{I.1.3.11-fr};
por tanto, $\Gamma(Y,f_*(\mathcal{F}))=\Gamma(X,\mathcal{F})$ es exacto en $\mathcal{F}$, lo que demuestra nuestra afirmación por \sref{II.5.2.1-fr}[\emph{(c)}].

Si $f$ es afín, $f^{-1}(U)$ es afín para todo abierto afín $U$ de $Y$ \sref{II.1.3.2-fr}, y por tanto $\mathrm{H}^1(f^{-1}(U),\mathcal{F})=0$ \sref{II.5.2.1-fr}[\emph{(d)}], lo cual implica, por definición, que $\mathrm{R}^1f_*(\mathcal{F})=0$.
Finalmente, supongamos que se satisface la condición~\emph{(c')};
la sucesión exacta de términos de grados bajos de la sucesión espectral de Leray (G,~II,~4.17.1 e I,~4.5.1) da, en particular, la sucesión exacta
\[
  0\to\mathrm{H}^1(Y,f_*(\mathcal{J}))\to
  \mathrm{H}^1(X,\mathcal{J})\to
  \mathrm{H}^0(Y,\mathrm{R}^1f_*(\mathcal{J})).
\]
Como $Y$ es afín y $f_*(\mathcal{J})$ es cuasicoherente \sref[I]{I.9.2.2-fr}[\emph{(a)}], se tiene $\mathrm{H}^1(Y,f_*(\mathcal{J}))=0$ \sref{II.5.2.1-fr};
la hipótesis~\emph{(c')} implica entonces que $\mathrm{H}^1(X,\mathcal{J})=0$, y se concluye por \sref{II.5.2.1-fr} que $X$ es un esquema afín.
\end{proof}
% Fuente francesa: PMIHES_1961__8__5_0, printed pp. 98-99.

\begin{corollary}[5.2.3]
\label{II.5.2.3-fr}
Si $f:X\to Y$ es un morfismo afín, entonces, para todo $\mathcal{O}_X$-módulo cuasicoherente $\mathcal{F}$, el homomorfismo canónico $\mathrm{H}^1(Y,f_*(\mathcal{F}))\to\mathrm{H}^1(X,\mathcal{F})$ es biyectivo.
\end{corollary}

\oldpage[II]{99}
\begin{proof}
Se tiene la sucesión exacta
\[
  0\to\mathrm{H}^1(Y,f_*(\mathcal{F}))\to
  \mathrm{H}^1(X,\mathcal{F})\to
  \mathrm{H}^0(Y,\mathrm{R}^1f_*(\mathcal{F}))
\]
de términos de grados bajos de la sucesión espectral de Leray, y la conclusión se sigue de \sref{II.5.2.2-fr}.
\end{proof}

\begin{remark}[5.2.4]
\label{II.5.2.4-fr}
En el capítulo~III,~\textsection1, demostraremos que, si $X$ es afín, entonces $\mathrm{H}^i(X,\mathcal{F})=0$ para todo $i>0$ y todo $\mathcal{O}_X$-módulo cuasicoherente $\mathcal{F}$.
\end{remark}
% Fuente francesa: PMIHES_1961__8__5_0, printed p. 99.

\subsection{Morfismos cuasiproyectivos}
\label{subsection:II.5.3-fr}

\begin{definition}[5.3.1]
\label{II.5.3.1-fr}
Se dice que un morfismo $f:X\to Y$ es cuasiproyectivo, o que $X$ (considerado como $Y$-preesquema mediante $f$) es cuasiproyectivo sobre $Y$, o que $X$ es un $Y$-esquema cuasiproyectivo, si $f$ es de tipo finito y existe un $\mathcal{O}_X$-módulo invertible $f$-amplio.
\end{definition}

Se observa que esta noción \emph{no es local sobre $Y$}:
los contraejemplos de Nagata~\cite{II-26} y de Hironaka muestran que, incluso si $X$ e $Y$ son esquemas algebraicos no singulares sobre un cuerpo algebraicamente cerrado, todo punto de $Y$ puede tener una vecindad afín $U$ tal que $f^{-1}(U)$ sea cuasiproyectivo sobre $U$, sin que $f$ sea cuasiproyectivo.

Se observa que todo morfismo cuasiproyectivo es necesariamente \emph{separado} \sref{II.4.6.1-fr}.
Cuando $Y$ es cuasicompacto, es equivalente decir que $f$ es cuasiproyectivo o que $f$ es de tipo finito y existe un $\mathcal{O}_X$-módulo \emph{muy amplio} relativamente a $f$ (\sref{II.4.6.2-fr} y \sref{II.4.6.11-fr}).
Además:

\begin{proposition}[5.3.2]
\label{II.5.3.2-fr}
Sean $Y$ un esquema cuasicompacto o un preesquema cuyo espacio subyacente sea noetheriano, y $X$ un $Y$-preesquema.
Las condiciones siguientes son equivalentes:
\begin{enumerate}
  \item[{\rm(a)}] $X$ es un $Y$-esquema cuasiproyectivo.
  \item[{\rm(b)}] $X$ es de tipo finito sobre $Y$, y existe un $\mathcal{O}_Y$-módulo cuasicoherente $\mathcal{E}$ de tipo finito tal que $X$ es $Y$-isomorfo a un subpreesquema de $\mathbf{P}(\mathcal{E})$.
  \item[{\rm(c)}] $X$ es de tipo finito sobre $Y$, y existe una $\mathcal{O}_Y$-álgebra graduada cuasicoherente $\mathcal{S}$ tal que $\mathcal{S}_1$ es de tipo finito y genera $\mathcal{S}$, y tal que $X$ es $Y$-isomorfo a un subpreesquema inducido sobre un abierto denso en todas partes de $\operatorname{Proj}(\mathcal{S})$.
\end{enumerate}
\end{proposition}

\begin{proof}
Esto se sigue inmediatamente de la observación anterior y de \sref{II.4.4.3-fr}, \sref{II.4.4.6-fr} y \sref{II.4.4.7-fr}.
\end{proof}

Se observa que, cuando $Y$ es un preesquema \emph{noetheriano}, pueden suprimirse, en las condiciones~\emph{(b)} y \emph{(c)} de \sref{II.5.3.2-fr}, la hipótesis de que $X$ es de tipo finito sobre $Y$, pues entonces se satisface automáticamente \sref[I]{I.6.3.5-fr}.
% Fuente francesa: PMIHES_1961__8__5_0, printed p. 99.

\begin{corollary}[5.3.3]
\label{II.5.3.3-fr}
Sea $Y$ un esquema cuasicompacto tal que existe un $\mathcal{O}_Y$-módulo amplio $\mathcal{L}$ \sref{II.4.5.3-fr}.
Para que un $Y$-esquema $X$ sea cuasiproyectivo, es necesario y suficiente que sea de tipo finito sobre $Y$ y, además, isomorfo a un $Y$-subesquema de un fibrado proyectivo de la forma $\mathbf{P}_Y^r$.
\end{corollary}

\begin{proof}
Si $\mathcal{E}$ es un $\mathcal{O}_Y$-módulo cuasicoherente de tipo finito, entonces $\mathcal{E}$ es isomorfo a un cociente de un $\mathcal{O}_Y$-módulo $\mathcal{L}^{\otimes(-n)}\otimes_{\mathcal{O}_Y}\mathcal{O}_Y^k$ \sref{II.4.5.5-fr}, y por tanto $\mathbf{P}(\mathcal{E})$ es isomorfo a un subesquema cerrado de $\mathbf{P}_Y^{k-1}$ (\sref{II.4.1.2-fr} y \sref{II.4.1.4-fr}).
\end{proof}
% Fuente francesa: PMIHES_1961__8__5_0, printed pp. 99-100.

\begin{proposition}[5.3.4]
\label{II.5.3.4-fr}
\medskip\noindent
\begin{enumerate}
  \item[{\rm(i)}] Todo morfismo cuasiafín de tipo finito (y, en particular, toda inmersión cuasicompacta o todo morfismo afín de tipo finito) es cuasiproyectivo.
  \item[{\rm(ii)}] Si $f:X\to Y$ y $g:Y\to Z$ son cuasiproyectivos, y si $Z$ es cuasicompacto, entonces $g\circ f$ es cuasiproyectivo.
\oldpage[II]{100}
  \item[{\rm(iii)}] Si $f:X\to Y$ es un $S$-morfismo cuasiproyectivo, entonces $f_{(S')}:X_{(S')}\to Y_{(S')}$ es cuasiproyectivo para toda extensión $S'\to S$ del preesquema de base.
  \item[{\rm(iv)}] Si $f:X\to Y$ y $g:X'\to Y'$ son dos $S$-morfismos cuasiproyectivos, entonces $f\times_Sg$ es cuasiproyectivo.
  \item[{\rm(v)}] Si $f:X\to Y$ y $g:Y\to Z$ son morfismos tales que $g\circ f$ es cuasiproyectivo, y si $g$ es separado o $X$ es localmente noetheriano, entonces $f$ es cuasiproyectivo.
  \item[{\rm(vi)}] Si $f$ es un morfismo cuasiproyectivo, también lo es $f_{\mathrm{red}}$.
\end{enumerate}
\end{proposition}

\begin{proof}
(i) se sigue de \sref{II.5.1.6-fr} y \sref{II.5.1.10-fr}[(i)].
Las demás afirmaciones son consecuencias inmediatas de la definición~\sref{II.5.3.1-fr}, de las propiedades de los morfismos de tipo finito \sref[I]{I.6.3.4-fr} y de \sref{II.4.6.13-fr}.
\end{proof}

\begin{remark}[5.3.5]
\label{II.5.3.5-fr}
Se observa que puede ocurrir que $f_{\mathrm{red}}$ sea cuasiproyectivo sin que $f$ lo sea, incluso suponiendo que $Y$ es el espectro de un álgebra de rango finito sobre $\mathbf{C}$ y que $f$ es propio.
\end{remark}

\begin{corollary}[5.3.6]
\label{II.5.3.6-fr}
Si $X$ y $X'$ son dos $Y$-esquemas cuasiproyectivos, entonces $X\amalg X'$ es un $Y$-esquema cuasiproyectivo.
\end{corollary}

\begin{proof}
Esto se sigue de \sref{II.4.6.18-fr}.
\end{proof}
% Fuente francesa: PMIHES_1961__8__5_0, printed pp. 100-101.

\subsection{Morfismos propios y morfismos universalmente cerrados}
\label{subsection:II.5.4-fr}
\label{II.5.4-fr}

\begin{definition}[5.4.1]
\label{II.5.4.1-fr}
Se dice que un morfismo de preesquemas $f:X\to Y$ es propio si satisface las dos condiciones siguientes:
\begin{enumerate}
  \item[(a)] $f$ es separado y de tipo finito;
  \item[(b)] para todo preesquema $Y'$ y todo morfismo $Y'\to Y$ la proyección $f_{(Y')}:X\times_YY'\to Y'$ es un morfismo cerrado \sref[I]{I.2.2.6-fr}.
\end{enumerate}

En tal caso, se dice también que $X$ (considerado como un $Y$-preesquema cuyo morfismo estructural es $f$) es propio sobre $Y$.
\end{definition}

Es inmediato que las condiciones~\emph{(a)} y \emph{(b)} son \emph{locales sobre} $Y$.
Para mostrar que la imagen de un subconjunto cerrado $Z$ de $X\times_YY'$ por la proyección $q:X\times_YY'\to Y'$ es cerrada en $Y'$, basta ver que $q(Z)\cap U'$ es cerrado en $U'$ para todo abierto afín $U'$ de $Y'$;
como $q(Z)\cap U'=q(Z\cap q^{-1}(U'))$, y como $q^{-1}(U')$ puede identificarse con $X\times_YU'$ \sref[I]{I.4.4.1-fr}, se ve que, para satisfacer la condición~\emph{(b)} de la Definición~\sref{II.5.4.1-fr}, se puede limitar al caso en que $Y'$ sea un esquema \emph{afín}.
Se verá más adelante \sref{II.5.6.3-fr} que, si $Y$ es localmente noetheriano, puede incluso bastar con demostrar \emph{(b)} cuando $Y'$ es de tipo finito sobre $Y$.

Es claro que todo morfismo propio es \emph{cerrado}.
% Fuente francesa: PMIHES_1961__8__5_0, printed p. 101.

\begin{proposition}[5.4.2]
\label{II.5.4.2-fr}
\medskip\noindent
\begin{enumerate}
  \item[{\rm(i)}] Toda inmersión cerrada es un morfismo propio.
  \item[{\rm(ii)}] La composición de dos morfismos propios es propia.
  \item[{\rm(iii)}] Si $X$ e $Y$ son $S$-preesquemas y $f:X\to Y$ es un $S$-morfismo propio, entonces $f_{(S')}:X_{(S')}\to Y_{(S')}$ es propio para toda extensión $S'\to S$ del preesquema de base.
  \item[{\rm(iv)}] Si $f:X\to Y$ y $g:X'\to Y'$ son $S$-morfismos propios, el $S$-morfismo $f\times_Sg:X\times_SX'\to Y\times_SY'$ es propio.
\end{enumerate}
\end{proposition}

\oldpage[II]{101}
\begin{proof}
Basta demostrar (i), (ii) y (iii) \sref[I]{I.3.5.1-fr}.
En cada uno de los tres casos, la verificación de la condición~\emph{(a)} de \sref{II.5.4.1-fr} se deduce de resultados anteriores (\sref[I]{I.5.5.1-fr} y \sref[I]{I.6.3.4-fr}); queda por verificar la condición~\emph{(b)}.
Esto es inmediato en el caso (i), pues, si $X\to Y$ es una inmersión cerrada, también lo es $X\times_YY'\to Y\times_YY'=Y'$ (\sref[I]{I.4.3.2-fr} y \sref{II.3.3.3-fr}).
Para demostrar (ii), consideremos dos morfismos propios $X\to Y$ e $Y\to Z$, y un morfismo $Z'\to Z$.
Se puede escribir $X\times_ZZ'=X\times_Y(Y\times_ZZ')$ \sref[I]{I.3.3.9.1-fr}; por tanto, la proyección $X\times_ZZ'\to Z'$ se factoriza como $X\times_Y(Y\times_ZZ')\to Y\times_ZZ'\to Z'$.
Teniendo en cuenta la observación inicial, (ii) se deduce entonces de que la composición de dos morfismos cerrados es cerrada.
Por último, para todo morfismo $S'\to S$, se puede identificar $X_{(S')}$ con $X\times_YY_{(S')}$ \sref[I]{I.3.3.11-fr}; para todo morfismo $Z\to Y_{(S')}$, se puede escribir
\[
  X_{(S')}\times_{Y_{(S')}}Z
  =(X\times_YY_{(S')})\times_{Y_{(S')}}Z=X\times_YZ~;
\]
como, por hipótesis, $X\times_YZ\to Z$ es cerrado, esto demuestra (iii).
\end{proof}
% Fuente francesa: PMIHES_1961__8__5_0, printed p. 101.

\begin{corollary}[5.4.3]
\label{II.5.4.3-fr}
Sean $f:X\to Y$ y $g:Y\to Z$ dos morfismos tales que $g\circ f$ sea propio.
\begin{enumerate}
  \item[{\rm(i)}] Si $g$ es separado, $f$ es propio.
  \item[{\rm(ii)}] Si $g$ es separado y de tipo finito, y si $f$ es sobreyectivo, entonces $g$ es propio.
\end{enumerate}
\end{corollary}

\begin{proof}
(i) se deduce de \sref{II.5.4.2-fr} por el procedimiento general \sref[I]{I.5.5.12-fr}.
Para demostrar (ii), solo hay que verificar que se satisface la condición~\emph{(b)} de la Definición~\sref{II.5.4.1-fr}.
Para todo morfismo $Z'\to Z$, el diagrama
\[
  \xymatrix{
    X\times_ZZ'\ar[r]^{f\times1_{Z'}}\ar[dr]_p &
    Y\times_ZZ'\ar[d]^{p'}\\
    & Z'
  }
\]
(donde $p$ y $p'$ son las proyecciones) es conmutativo \sref[I]{I.3.2.1-fr};
además, $f\times1_{Z'}$ es sobreyectivo porque $f$ lo es \sref[I]{I.3.5.2-fr}, y $p$ es un morfismo cerrado por hipótesis.
Todo subconjunto cerrado $F$ de $Y\times_ZZ'$ es, por tanto, la imagen por $f\times1_{Z'}$ de algún subconjunto cerrado $E$ de $X\times_ZZ'$; así, $p'(F)=p(E)$ es cerrado en $Z'$ por hipótesis, de donde el corolario.
\end{proof}
% Fuente francesa: PMIHES_1961__8__5_0, printed p. 101.

\begin{corollary}[5.4.4]
\label{II.5.4.4-fr}
Si $X$ es un preesquema propio sobre $Y$, y $\mathcal{S}$ es una $\mathcal{O}_Y$-álgebra cuasicoherente, entonces todo $Y$-morfismo $f:X\to\operatorname{Proj}(\mathcal{S})$ es propio (y, \emph{a fortiori}, cerrado).
\end{corollary}

\begin{proof}
El morfismo estructural $p:\operatorname{Proj}(\mathcal{S})\to Y$ es separado, y $p\circ f$ es propio por hipótesis.
\end{proof}

\begin{corollary}[5.4.5]
\label{II.5.4.5-fr}
Sea $f:X\to Y$ un morfismo separado de tipo finito.
Sean $(X_i)_{1\leq i\leq n}$ \textup{(resp. $(Y_i)_{1\leq i\leq n}$)} una familia finita de subpreesquemas cerrados de $X$ \textup{(resp. de $Y$)}, y $j_i$ \textup{(resp. $h_i$)} la inyección canónica $X_i\to X$ \textup{(resp. $Y_i\to Y$)}.
Supongamos que el espacio subyacente de $X$ es la unión de los $X_i$ y que, para todo $i$, existe un morfismo $f_i:X_i\to Y_i$ tal que el diagrama
\[
  \xymatrix{
    X_i\ar[r]^{f_i}\ar[d]_{j_i} & Y_i\ar[d]^{h_i}\\
    X\ar[r]_f & Y
  }
\]
es conmutativo.
Entonces, para que $f$ sea propio, es necesario y suficiente que cada uno de los $f_i$ lo sea.
\end{corollary}
% Fuente francesa: PMIHES_1961__8__5_0, printed p. 102.

\oldpage[II]{102}
\begin{proof}
Si $f$ es propio, también lo es $f\circ j_i$, pues $j_i$ es una inmersión cerrada \sref{II.5.4.2-fr};
como $h_i$ es una inmersión cerrada y, por tanto, un morfismo separado, $f_i$ es propio por \sref{II.5.4.3-fr}.
Recíprocamente, supongamos que cada uno de los $f_i$ es propio y consideremos el preesquema $Z$ dado por la \emph{suma} de los $X_i$; sea $u$ el morfismo $Z\to X$ que se reduce a $j_i$ sobre cada $X_i$.
La restricción de $f\circ u$ a cada $X_i$ es igual a $f\circ j_i=h_i\circ f_i$ y, por tanto, es propia, pues tanto $h_i$ como $f_i$ lo son \sref{II.5.4.2-fr};
se sigue entonces inmediatamente de la Definición~\sref{II.5.4.1-fr} que $f\circ u$ es propio.
Pero, como $u$ es sobreyectivo por hipótesis, se concluye que $f$ es propio por \sref{II.5.4.3-fr}.
\end{proof}

\begin{corollary}[5.4.6]
\label{II.5.4.6-fr}
Sea $f:X\to Y$ un morfismo separado de tipo finito; para que $f$ sea propio, es necesario y suficiente que $f_{\mathrm{red}}:X_{\mathrm{red}}\to Y_{\mathrm{red}}$ sea propio.
\end{corollary}

\begin{proof}
Es un caso particular de \sref{II.5.4.5-fr}, con $n=1$, $X_1=X_{\mathrm{red}}$ e $Y_1=Y_{\mathrm{red}}$ \sref[I]{I.5.1.5-fr}.
\end{proof}
% Fuente francesa: PMIHES_1961__8__5_0, printed p. 102.

\begin{env}[5.4.7]
\label{II.5.4.7-fr}
Si $X$ e $Y$ son preesquemas noetherianos y $f:X\to Y$ es un morfismo separado de tipo finito, para demostrar que $f$ es propio se puede reducir al caso de morfismos \emph{dominantes} de preesquemas \emph{íntegros}.
En efecto, sean $X_i$ ($1\leq i\leq n$) las componentes irreducibles (finitas en número) de $X$ y consideremos para cada $i$ el único subpreesquema cerrado reducido de $X$ que tiene $X_i$ por espacio subyacente, que seguiremos denotando por $X_i$ \sref[I]{I.5.2.1-fr}.
Sea $Y_i$ el único subpreesquema cerrado reducido de $Y$ que tiene $\overline{f(X_i)}$ por espacio subyacente.
Si $g_i$ (resp. $h_i$) es el morfismo de inyección $X_i\to X$ (resp. $Y_i\to Y$), se concluye que $f\circ g_i=h_i\circ f_i$, donde $f_i$ es un morfismo dominante $X_i\to Y_i$ \sref[I]{I.5.2.2-fr};
se está entonces en condiciones de aplicar \sref{II.5.4.5-fr}, y para que $f$ sea propio es necesario y suficiente que todos los $f_i$ lo sean.
\end{env}
% Fuente francesa: PMIHES_1961__8__5_0, printed p. 102.

\begin{corollary}[5.4.8]
\label{II.5.4.8-fr}
Sean $X$ e $Y$ dos $S$-preesquemas separados de tipo finito sobre $S$, y sea $f:X\to Y$ un $S$-morfismo.
Para que $f$ sea propio, es necesario y suficiente que, para todo $S$-preesquema $S'$, el morfismo $f\times_S1_{S'}:X\times_SS'\to Y\times_SS'$ sea cerrado.
\end{corollary}

\begin{proof}
Observemos primero que, si $g:X\to S$ y $h:Y\to S$ son los morfismos estructurales, se tiene por definición $g=h\circ f$; por tanto, $f$ es separado y de tipo finito (\sref[I]{I.5.5.1-fr} y \sref[I]{I.6.3.4-fr}).
Si $f$ es propio, también lo es $f\times_S1_{S'}$ \sref{II.5.4.2-fr}; \emph{a fortiori}, $f\times_S1_{S'}$ es cerrado.
Recíprocamente, supongamos satisfechas las condiciones del enunciado y sea $Y'$ un $Y$-preesquema;
$Y'$ puede considerarse también como un $S$-preesquema y, como $Y\to S$ es separado, $X\times_YY'$ se identifica con un subpreesquema cerrado de $X\times_SY'$ \sref[I]{I.5.4.2-fr}.
En el diagrama conmutativo
\[
  \xymatrix{
    X\times_YY'\ar[r]^{f\times1_{Y'}}\ar[d] &
    Y\times_YY'=Y'\ar[d]\\
    X\times_SY'\ar[r]_{f\times_S1_{Y'}} & Y\times_SY'
  }
\]
las flechas verticales son inmersiones cerradas; se sigue inmediatamente que, si $f\times_S1_{Y'}$ es un morfismo cerrado, también lo es el cambio de base $f\times1_{Y'}$.
\end{proof}
% Fuente francesa: PMIHES_1961__8__5_0, printed pp. 102-103.

\begin{remark}[5.4.9]
\label{II.5.4.9-fr}
Diremos que un morfismo $f:X\to Y$ es \emph{universalmente cerrado} si satisface la condición~\emph{(b)} de la Definición~\sref{II.5.4.1-fr}.
El lector observará que,
\oldpage[II]{103}
en \sref{II.5.4.2-fr} a \sref{II.5.4.8-fr}, se puede reemplazar en todas partes ``propio'' por ``universalmente cerrado'' sin alterar la validez de los resultados (y, en las hipótesis de \sref{II.5.4.3-fr}, \sref{II.5.4.5-fr}, \sref{II.5.4.6-fr} y \sref{II.5.4.8-fr}, se pueden omitir las condiciones de finitud).
\end{remark}
% Fuente francesa: PMIHES_1961__8__5_0, printed p. 103.

\begin{env}[5.4.10]
\label{II.5.4.10-fr}
Sea $f:X\to Y$ un morfismo de tipo finito.
Diremos que un subconjunto cerrado $Z$ de $X$ es \emph{propio sobre $Y$} (o \emph{$Y$-propio}, o \emph{propio para $f$}) si la restricción de $f$ a un subpreesquema cerrado de $X$ cuyo espacio subyacente es $Z$ \sref[I]{I.5.2.1-fr} es un morfismo \emph{propio}.
Como esta restricción es entonces separada, se sigue de \sref{II.5.4.6-fr} y \sref[I]{I.5.5.1-fr}[vi] que la propiedad precedente \emph{no depende} del subpreesquema cerrado de $X$ que tiene $Z$ por espacio subyacente.
Si $g:X'\to X$ es un morfismo \emph{propio}, entonces $g^{-1}(Z)$ es un subconjunto \emph{propio} de $X'$:
si $T$ es un subpreesquema de $X$ cuyo espacio subyacente es $Z$, basta observar que la restricción de $g$ al subpreesquema cerrado $g^{-1}(T)$ de $X'$ es un morfismo propio $g^{-1}(T)\to T$ por \sref{II.5.4.2-fr}[iii], y aplicar después \sref{II.5.4.2-fr}[ii].
Por otra parte, si $X''$ es un $Y$-\emph{esquema} de tipo finito y $u:X\to X''$ es un $Y$-morfismo, entonces $u(Z)$ es un subconjunto \emph{propio} de $X''$;
en efecto, tomemos como $T$ el subpreesquema cerrado reducido de $X$ cuyo espacio subyacente es $Z$;
como la restricción de $f$ a $T$ es propia, también lo es la restricción de $u$ a $T$ \sref{II.5.4.3-fr}[(i)], por lo que $u(Z)$ es cerrado en $X''$;
sea $T''$ un subpreesquema cerrado de $X''$ cuyo espacio subyacente es $u(Z)$ \sref[I]{I.5.2.1-fr}, de modo que $u|T$ se factorice como \[
  T\xrightarrow{v}T''\xrightarrow{j}X''
\], donde $j$ es la inyección canónica \sref[I]{I.5.2.2-fr}; entonces $v$ es propio y sobreyectivo (\hyperref[II.5.4.5-fr]{5.4.5});
si $g$ es la restricción a $T''$ del morfismo estructural $X''\to Y$, entonces $g$ es separado y de tipo finito, y se tiene $f|T=g\circ v$;
se sigue entonces de \sref{II.5.4.3-fr}[(ii)] que $g$ es propio, de donde nuestra afirmación.
\end{env}
% Fuente francesa: PMIHES_1961__8__5_0, printed p. 103.

De estas observaciones resulta, en particular, que, si $Z$ es un subconjunto $Y$-propio de $X$, entonces:
\begin{enumerate}
  \item para todo subpreesquema cerrado $X'$ de $X$, $Z\cap X'$ es un subconjunto $Y$-propio de $X'$;
  \item si $X$ es un subpreesquema de un $Y$-esquema de tipo finito $X''$, entonces $Z$ es también un subconjunto $Y$-propio de $X''$ (y, en particular, es \emph{cerrado en $X''$}).
\end{enumerate}
% Fuente francesa: PMIHES_1961__8__5_0, printed pp. 103-104.

\subsection{Morfismos proyectivos}
\label{subsection:II.5.5-fr}

\begin{proposition}[5.5.1]
\label{II.5.5.1-fr}
Sea $X$ un $Y$-preesquema.
Las condiciones siguientes son equivalentes:
\begin{enumerate}
  \item[{\rm(a)}] $X$ es $Y$-isomorfo a un subpreesquema \emph{cerrado} de un fibrado proyectivo $\mathbf{P}(\mathcal{E})$, donde $\mathcal{E}$ es un $\mathcal{O}_Y$-módulo cuasicoherente de tipo finito.
  \item[{\rm(b)}] Existe una $\mathcal{O}_Y$-álgebra graduada cuasicoherente $\mathcal{S}$ tal que $\mathcal{S}_1$ es de tipo finito y genera $\mathcal{S}$, y tal que $X$ es $Y$-isomorfo a $\operatorname{Proj}(\mathcal{S})$.
\end{enumerate}
\end{proposition}

\begin{proof}
La condición~\emph{(a)} implica \emph{(b)} en virtud de \sref{II.3.6.2-fr}[ii]: si $\mathcal{J}$ es un haz graduado cuasicoherente de ideales de $\mathbf{S}(\mathcal{E})$, la $\mathcal{O}_Y$-álgebra graduada cuasicoherente $\mathcal{S}=\mathbf{S}(\mathcal{E})/\mathcal{J}$ está generada por $\mathcal{S}_1$, y este último, imagen canónica de $\mathcal{E}$, es un $\mathcal{O}_Y$-módulo de tipo finito.
La condición~\emph{(b)} implica \emph{(a)} en virtud de \sref{II.3.6.2-fr}, aplicado al caso en que $\mathcal{M}\to\mathcal{S}_1$ es la aplicación identidad.
\end{proof}

\oldpage[II]{104}
\begin{definition}[5.5.2]
\label{II.5.5.2-fr}
Se dice que un $Y$-preesquema $X$ es proyectivo sobre $Y$, o que es un $Y$-esquema proyectivo, si satisface las condiciones equivalentes~(a) y (b) de \sref{II.5.5.1-fr}.
Se dice que un morfismo $f:X\to Y$ es proyectivo si hace de $X$ un $Y$-esquema proyectivo.
\end{definition}

Es claro que, si $f:X\to Y$ es proyectivo, existe un $\mathcal{O}_X$-módulo \emph{muy amplio} relativamente a $f$ \sref{II.4.4.2-fr}.
% Fuente francesa: PMIHES_1961__8__5_0, printed p. 104.

\begin{theorem}[5.5.3]
\label{II.5.5.3-fr}
\medskip\noindent
\begin{enumerate}
  \item[{\rm(i)}] Todo morfismo proyectivo es cuasiproyectivo y propio.
  \item[{\rm(ii)}] Recíprocamente, sea $Y$ un esquema cuasicompacto o un preesquema cuyo espacio subyacente sea noetheriano; entonces todo morfismo $f:X\to Y$ que sea cuasiproyectivo y propio es proyectivo.
\end{enumerate}
\end{theorem}

\begin{proof}
\medskip\noindent
\begin{enumerate}
  \item[(i)] Es claro que, si $f:X\to Y$ es proyectivo, entonces es de tipo finito y cuasiproyectivo (y, por tanto, en particular separado); además, se sigue inmediatamente de \sref{II.5.5.1-fr}[\emph{b}] y \sref{II.3.5.3-fr} que, si $f$ es proyectivo, también lo es $f\times_Y1_{Y'}:X\times_YY'\to Y'$ para todo morfismo $Y'\to Y$.
    Para demostrar que $f$ es universalmente cerrado, basta entonces probar que todo morfismo proyectivo $f$ es \emph{cerrado}.
    Como la cuestión es local sobre $Y$, se puede suponer que $Y=\operatorname{Spec}(A)$ y, por \sref{II.5.5.1-fr}, que $X=\operatorname{Proj}(S)$, donde $S$ es una $A$-álgebra graduada generada por un número finito de elementos de $S_1$.
    Para todo $y\in Y$, la fibra $f^{-1}(y)$ se identifica con $\operatorname{Proj}(S)\times_Y\operatorname{Spec}(\boldsymbol{k}(y))$ \sref[I]{I.3.6.1-fr} y, por tanto, también con $\operatorname{Proj}(S\otimes_A\boldsymbol{k}(y))$ \sref{II.2.8.10-fr};
    así, $f^{-1}(y)$ es vacía si y solo si $S\otimes_A\boldsymbol{k}(y)$ satisface la condición~(TN) \sref{II.2.7.4-fr}, es decir, si $S_n\otimes_A\boldsymbol{k}(y)=0$ para $n$ suficientemente grande.
    Pero, como $(S_n)_y$ es un $\mathcal{O}_y$-módulo de tipo finito, la condición precedente implica que $(S_n)_y=0$ para $n$ suficientemente grande, por el lema de Nakayama.
    Si $\mathfrak a_n$ es el anulador en $A$ del $A$-módulo $S_n$, la condición precedente implica asimismo que $\mathfrak a_n\not\subset\mathfrak j_y$ para $n$ suficientemente grande \sref[0]{0.1.7.4-fr}.
    Ahora bien, como $S_nS_1=S_{n+1}$ por hipótesis, se tiene $\mathfrak a_n\subset\mathfrak a_{n+1}$; si $\mathfrak a$ es la unión de los $\mathfrak a_n$, se ve entonces que $f(X)=V(\mathfrak a)$, lo cual demuestra que $f(X)$ es cerrado en $Y$.
    Si $X'$ es ahora un subconjunto cerrado cualquiera de $X$, existe un subpreesquema cerrado de $X$ cuyo espacio subyacente es $X'$ \sref[I]{I.5.2.1-fr}, y es claro por \sref{II.5.5.1-fr}[\emph{a}] que el morfismo $X'\to X\xrightarrow{f}Y$ es proyectivo; por tanto, $f(X')$ es cerrado en $Y$.
  \item[(ii)] La hipótesis sobre $Y$ y el hecho de que $f$ sea cuasiproyectivo implican la existencia de un $\mathcal{O}_Y$-módulo cuasicoherente de tipo finito $\mathcal{E}$ y de una $Y$-inmersión $j:X\to\mathbf{P}(\mathcal{E})$ \sref{II.5.3.2-fr}.
    Pero, como $f$ es propio, $j$ es \emph{cerrada} por \sref{II.5.4.4-fr}, y por tanto $f$ es proyectivo.
\end{enumerate}
\end{proof}
% Fuente francesa: PMIHES_1961__8__5_0, printed pp. 104-105.

\par\medskip\noindent\textit{Observaciones (5.5.4). ---\ }
\label{II.5.5.4-fr}
\medskip\noindent
\begin{enumerate}
  \item[(i)] Sea $f:X\to Y$ un morfismo tal que: 1.º $f$ sea \emph{propio}; 2.º exista un $\mathcal{O}_X$-módulo $\mathcal{L}$ \emph{muy amplio} relativamente a $f$; 3.º el $\mathcal{O}_Y$-módulo cuasicoherente $\mathcal{E}=f_*(\mathcal{L})$ sea \emph{de tipo finito}.
    Entonces $f$ es un morfismo \emph{proyectivo}: en efecto, por \sref{II.4.4.4-fr}, existe una $Y$-inmersión $r:X\to\mathbf{P}(\mathcal{E})$ y, como $f$ es propio, $r$ es una inmersión \emph{cerrada} \sref{II.5.4.4-fr}.
    Se verá en el Capítulo~III, \textsection3, que, cuando $Y$ es \emph{localmente noetheriano}, la tercera condición precedente es consecuencia de las otras dos; por tanto, en este caso, las dos primeras condiciones \emph{caracterizan} los morfismos proyectivos y, si $Y$ es cuasicompacto, se puede reemplazar la segunda condición por la hipótesis de que existe un $\mathcal{O}_X$-módulo \emph{amplio} relativamente a $f$ \sref{II.4.6.11-fr}.
  \item[(ii)] Sea $Y$ un esquema cuasicompacto tal que existe un $\mathcal{O}_Y$-módulo amplio.
    Para que un $Y$-esquema $X$ sea \emph{proyectivo}, es necesario y suficiente que sea $Y$-isomorfo a un $Y$-subesquema \emph{cerrado} de un fibrado proyectivo de la forma $\mathbf{P}_Y^r$.
    La condición es claramente suficiente.
\oldpage[II]{105}
    Recíprocamente, si $X$ es proyectivo sobre $Y$, es cuasiproyectivo; por tanto, existe una $Y$-inmersión $j$ de $X$ en algún $\mathbf{P}_Y^r$ \sref{II.5.3.3-fr}, que es \emph{cerrada} por \sref{II.5.4.4-fr} y \sref{II.5.5.3-fr}.
  \item[(iii)] El razonamiento de \sref{II.5.5.3-fr} muestra que, para todo preesquema $Y$ y todo entero $r\geq0$, el morfismo estructural $\mathbf{P}_Y^r\to Y$ es \emph{sobreyectivo}; en efecto, si se pone $\mathcal{S}=\mathbf{S}_{\mathcal{O}_Y}(\mathcal{O}_Y^{r+1})$, se tiene evidentemente $\mathcal{S}_y=\mathbf{S}_{\boldsymbol{k}(y)}((\boldsymbol{k}(y))^{r+1})$ \sref{II.1.7.3-fr}, y por tanto $(\mathcal{S}_n)_y\neq0$ para todo $y\in Y$ y todo $n\geq0$.
  \item[(iv)] De los ejemplos de Nagata~\cite{II-26} se deduce que existen morfismos propios que no son cuasiproyectivos.
\end{enumerate}

% Fuente francesa: PMIHES_1961__8__5_0, printed p. 105.

\begin{proposition}[5.5.5]
\label{II.5.5.5-fr}
\medskip\noindent
\begin{enumerate}
  \item[{\rm(i)}] Toda inmersión cerrada es un morfismo proyectivo.
  \item[{\rm(ii)}] Si $f:X\to Y$ y $g:Y\to Z$ son morfismos proyectivos, y si $Z$ es un esquema cuasicompacto o un preesquema cuyo espacio subyacente es noetheriano, entonces $g\circ f$ es proyectivo.
  \item[{\rm(iii)}] Si $f:X\to Y$ es un $S$-morfismo proyectivo, entonces $f_{(S')}:X_{(S')}\to Y_{(S')}$ es proyectivo para toda extensión $S'\to S$ del preesquema de base.
  \item[{\rm(iv)}] Si $f:X\to Y$ y $g:X'\to Y'$ son $S$-morfismos proyectivos, también lo es $f\times_Sg$.
  \item[{\rm(v)}] Si $g\circ f$ es un morfismo proyectivo y $g$ es separado, entonces $f$ es proyectivo.
  \item[{\rm(vi)}] Si $f$ es proyectivo, también lo es $f_{\mathrm{red}}$.
\end{enumerate}
\end{proposition}
% Fuente francesa: PMIHES_1961__8__5_0, printed p. 105.

\begin{proof}
(i) se deduce inmediatamente de \sref{II.3.1.7-fr}.
Aquí hay que demostrar (iii) y (iv) por separado, a causa de la restricción impuesta a $Z$ en (ii) (cf.~\sref[I]{I.3.5.1-fr}).
Para demostrar (iii), se reduce al caso en que $S=Y$ \sref[I]{I.3.3.11-fr}, y la afirmación se deduce entonces inmediatamente de \sref{II.5.5.1-fr}[\emph{b}] y \sref{II.3.5.3-fr}.
Para demostrar (iv), se reduce inmediatamente al caso en que $X=\mathbf{P}(\mathcal{E})$ y $X'=\mathbf{P}(\mathcal{E}')$, donde $\mathcal{E}$ (resp. $\mathcal{E}'$) es un $\mathcal{O}_Y$-módulo (resp. un $\mathcal{O}_{Y'}$-módulo) cuasicoherente de tipo finito.
Sean $p$ y $p'$ las proyecciones canónicas de $T=Y\times_SY'$ sobre $Y$ e $Y'$, respectivamente; por \sref{II.4.1.3.1-fr}, se tiene $\mathbf{P}(p^*(\mathcal{E}))=\mathbf{P}(\mathcal{E})\times_YT$ y $\mathbf{P}(p'^*(\mathcal{E}'))=\mathbf{P}(\mathcal{E}')\times_{Y'}T$; de donde
\begin{align*}
\mathbf{P}(p^*(\mathcal{E}))\times_T\mathbf{P}(p'^*(\mathcal{E}'))
&=(\mathbf{P}(\mathcal{E})\times_YT)\times_T
(T\times_{Y'}\mathbf{P}(\mathcal{E}'))\\
&=\mathbf{P}(\mathcal{E})\times_Y
(T\times_{Y'}\mathbf{P}(\mathcal{E}'))
=\mathbf{P}(\mathcal{E})\times_S\mathbf{P}(\mathcal{E}')
\end{align*}
al sustituir $T$ por $Y\times_SY'$ y utilizar \sref[I]{I.3.3.9.1-fr}.
Pero $p^*(\mathcal{E})$ y $p'^*(\mathcal{E}')$ son de tipo finito sobre $T$ \sref[0]{0.5.2.4-fr}, y por tanto también lo es $p^*(\mathcal{E})\otimes_{\mathcal{O}_T}p'^*(\mathcal{E}')$;
como \[
\mathbf{P}(p^*(\mathcal{E}))\times_T\mathbf{P}(p'^*(\mathcal{E}'))
\] se identifica con un subpreesquema cerrado de $\mathbf{P}(p^*(\mathcal{E})\otimes_{\mathcal{O}_T}p'^*(\mathcal{E}'))$ \sref{II.4.3.3-fr}, esto termina de demostrar (iv).
Para demostrar (v) y (vi), se puede aplicar \sref[I]{I.5.5.13-fr}, pues todo subpreesquema cerrado de un $Y$-esquema proyectivo es un $Y$-esquema proyectivo por \sref{II.5.5.1-fr}[\emph{a}].

Queda por demostrar (ii); en virtud de la hipótesis sobre $Z$, esto se deduce de \sref{II.5.5.3-fr}, \sref{II.5.3.4-fr}[ii] y \sref{II.5.4.2-fr}[ii].
\end{proof}
% Fuente francesa: PMIHES_1961__8__5_0, printed p. 105.

\begin{proposition}[5.5.6]
\label{II.5.5.6-fr}
Si $X$ y $X'$ son $Y$-esquemas proyectivos, entonces $X\amalg X'$ es un $Y$-esquema proyectivo.
\end{proposition}

\begin{proof}
Es una consecuencia evidente de \sref{II.5.5.2-fr} y \sref{II.4.3.6-fr}.
\end{proof}
% Fuente francesa: PMIHES_1961__8__5_0, printed pp. 105-106.

\begin{proposition}[5.5.7]
\label{II.5.5.7-fr}
Sea $X$ un $Y$-esquema proyectivo, y sea $\mathcal{L}$ un $\mathcal{O}_X$-módulo $Y$-amplio; entonces, para toda sección $f$ de $\mathcal{L}$ sobre $X$, $X_f$ es afín sobre $Y$.
\end{proposition}

\oldpage[II]{106}
\begin{proof}
Como la cuestión es local sobre $Y$, se puede suponer que $Y=\operatorname{Spec}(A)$; además, $X_{f^{\otimes n}}=X_f$, por lo que, reemplazando $\mathcal{L}$ por un $\mathcal{L}^{\otimes n}$ conveniente, se puede suponer que $\mathcal{L}$ es muy amplio relativamente al morfismo estructural $q:X\to Y$ \sref{II.4.6.11-fr}.
El homomorfismo canónico $\sigma:q^*(q_*(\mathcal{L}))\to\mathcal{L}$ es entonces sobreyectivo, y el morfismo correspondiente
\[
r=r_{\mathcal{L},\sigma}:X\to P=\mathbf{P}(q_*(\mathcal{L}))
\]
es una inmersión tal que $\mathcal{L}=r^*(\mathcal{O}_P(1))$ \sref{II.4.4.4-fr}; además, como $X$ es propio sobre $Y$, la inmersión $r$ es cerrada \sref{II.5.4.4-fr}.
Pero, por definición, $f\in\Gamma(Y,q_*(\mathcal{L}))$, y $\sigma^\flat$ es la identidad de $q_*(\mathcal{L})$; se sigue entonces de la Fórmula~\sref{II.3.7.3.1-fr} que $X_f=r^{-1}(D_+(f))$;
por tanto, $X_f$ es un subpreesquema cerrado del esquema afín $D_+(f)$, y es también un esquema afín.
\end{proof}
% Fuente francesa: PMIHES_1961__8__5_0, printed p. 106.

En el caso particular en que $Y=X$, se obtiene (teniendo en cuenta \sref{II.4.6.13-fr}[i]) el corolario siguiente, cuya demostración directa es, por lo demás, inmediata:
\begin{corollary}[5.5.8]
\label{II.5.5.8-fr}
Sea $X$ un preesquema y sea $\mathcal{L}$ un $\mathcal{O}_X$-módulo invertible.
Para toda sección $f$ de $\mathcal{L}$ sobre $X$, $X_f$ es afín sobre $X$ (y, por tanto, es también un esquema afín cuando $X$ es un esquema afín).
\end{corollary}
% Fuente francesa: PMIHES_1961__8__5_0, printed p. 106.

\subsection{Lema de Chow}
\label{subsection:II.5.6-fr}

\begin{theorem}[5.6.1]
\label{II.5.6.1-fr}
\emph{(Lema de Chow)}.
Sean $S$ un preesquema y $X$ un $S$-esquema de tipo finito.
Se supone satisfecha una de las condiciones siguientes:
\begin{enumerate}
  \item[{\rm(a)}] $S$ es noetheriano;
  \item[{\rm(b)}] $S$ es un esquema cuasicompacto y $X$ tiene un número finito de componentes irreducibles.
\end{enumerate}
En estas condiciones:
\begin{enumerate}
  \item[{\rm(i)}] existe un $S$-esquema $X'$ \emph{cuasiproyectivo} sobre $S$ y un $S$-morfismo $f:X'\to X$ \emph{proyectivo} y \emph{sobreyectivo};
  \item[{\rm(ii)}] se pueden elegir $X'$ y $f$ de modo que exista un abierto $U\subset X$ para el cual $U'=f^{-1}(U)$ sea denso en $X'$ y la restricción de $f$ a $U'$ sea un isomorfismo $U'\xrightarrow{\sim}U$;
  \item[{\rm(iii)}] si $X$ es reducido \textup{(resp.} irreducible, íntegro), se puede suponer que $X'$ es reducido \textup{(resp.} irreducible, íntegro).
\end{enumerate}
\end{theorem}
% Fuente francesa: PMIHES_1961__8__5_0, printed pp. 106-107.

\begin{proof}
La demostración se divide en varias etapas.
\begin{enumerate}
  \item[(A)] Se puede reducir primero al caso en que $X$ es \emph{irreducible}.
    En efecto, bajo la hipótesis~(\emph{a)}, $X$ es noetheriano; por tanto, bajo cualquiera de las dos hipótesis, las componentes irreducibles $X_i$ de $X$ son finitas en número.
    Si el teorema se ha demostrado para cada uno de los subpreesquemas cerrados reducidos de $X$ cuyos espacios subyacentes son los $X_i$, y si $X'_i$ y $f_i:X'_i\to X_i$ son, respectivamente, el preesquema y el morfismo correspondientes a $X_i$, entonces el preesquema $X'$ dado por la \emph{suma} de los $X'_i$ y el morfismo $f:X'\to X$ cuya restricción a cada $X'_i$ es $j_i\circ f_i$ (donde $j_i$ es la inyección canónica $X_i\to X$) satisfacen la conclusión del teorema.
    Es inmediato que $X'$ es reducido si todos los $X'_i$ lo son; además, se puede satisfacer (ii) tomando como $U$ la unión de los conjuntos $U_i\cap\complement(\bigcup_{j\neq i}X_j)$.
    Por último, como los $X'_i$ son cuasiproyectivos sobre $S$, también lo es $X'$
\oldpage[II]{107}
    \sref{II.5.3.6-fr}; asimismo, los morfismos $X'_i\to X$ son proyectivos por (\hyperref[II.5.5.5-fr]{5.5.5, (i) y (ii)}), de modo que $f$ es proyectivo \sref{II.5.5.6-fr}, y es evidentemente sobreyectivo por definición.
% Fuente francesa: PMIHES_1961__8__5_0, printed p. 107.

  \item[(B)] Supongamos ahora que $X$ es \emph{irreducible}.
    Como el morfismo estructural $r:X\to S$ es de tipo finito, existe una cubierta finita $(S_i)$ de $S$ por abiertos afines y, para cada $i$, una cubierta finita $(T_{ij})$ de $r^{-1}(S_i)$ por abiertos afines; los morfismos $T_{ij}\to S_i$ son afines y de tipo finito, y por tanto cuasiproyectivos \sref{II.5.3.4-fr}[i];
    como, bajo ambas hipótesis~\emph{(a)} y \emph{(b)}, la inmersión $S_i\to S$ es cuasicompacta, también es cuasiproyectiva \sref{II.5.3.4-fr}[i], de modo que la restricción de $r$ a $T_{ij}$ es un morfismo cuasiproyectivo \sref{II.5.3.4-fr}[ii].
    Denotemos los $T_{ij}$ por $U_k$ ($(1\leq k\leq n)$).
    Para cada índice $k$, existe una inmersión abierta $\varphi_k:U_k\to P_k$, donde $P_k$ es proyectivo sobre $S$ (\sref{II.5.3.2-fr} y \sref{II.5.5.2-fr}).
    Sea $U=\bigcap_kU_k$; como $X$ es irreducible y los $U_k$ son no vacíos, $U$ es no vacío y, por tanto, denso en $X$; las restricciones a $U$ de los $\varphi_k$ definen un morfismo
    \[
\varphi:U\to P=P_1\times_SP_2\times_S\ldots\times_SP_n
\]
    tal que los diagramas
    \[
\tag{5.6.1.1}
\label{II.5.6.1.1-fr}
\xymatrix{
U\ar[r]^\varphi\ar[d]_{j_k} & P\ar[d]^{p_k}\\
U_k\ar[r]_{\varphi_k} & P_k
}
\]
    son conmutativos, donde $j_k$ es la inyección canónica $U\to U_k$ y $p_k$ la proyección canónica $P\to P_k$.
% Fuente francesa: PMIHES_1961__8__5_0, printed p. 107.

    Si $j$ es la inyección canónica $U\to X$, el morfismo $\psi=(j,\varphi)_S:U\to X\times_SP$ es una \emph{inmersión} \sref[I]{I.5.3.14-fr}.
    Bajo la hipótesis~(\emph{a)}, $X\times_SP$ es localmente noetheriano (\sref{II.3.4.1-fr}, (\hyperref[I.6.3.7-fr]{\textbf{I}, 6.3.7 y 6.3.8}));
    bajo la hipótesis~\emph{(b)}, $X\times_SP$ es un esquema cuasicompacto ((\hyperref[I.5.5.1-fr]{\textbf{I}, 5.5.1 y 6.6.4}));
    en ambos casos existe la \emph{clausura} $X'$ en $X\times_SP$ del subpreesquema $Z$ asociado a $\psi$ (y cuyo espacio subyacente es $\psi(U)$), y $\psi$ se factoriza como
    \[
\tag{5.6.1.2}
\label{II.5.6.1.2-fr}
\psi:U\xrightarrow{\psi'}X'\xrightarrow{h}X\times_SP
\]
    donde $\psi'$ es una \emph{inmersión abierta} y $h$ una \emph{inmersión cerrada} \sref[I]{I.9.5.10-fr}.
    Sean $q_1:X\times_SP\to X$ y $q_2:X\times_SP\to P$ las proyecciones canónicas; pongamos
    \[
\tag{5.6.1.3}
\label{II.5.6.1.3-fr}
f:X'\xrightarrow{h}X\times_SP\xrightarrow{q_1}X
\]
    \[
\tag{5.6.1.4}
\label{II.5.6.1.4-fr}
g:X'\xrightarrow{h}X\times_SP\xrightarrow{q_2}P.
\]
    Veremos que $X'$ y $f$ satisfacen la conclusión del teorema.
% Fuente francesa: PMIHES_1961__8__5_0, printed pp. 107-108.

  \item[(C)] Demostremos primero que $f$ es \emph{proyectivo} y \emph{sobreyectivo}, y que la restricción de $f$ a $U'=f^{-1}(U)$ es un \emph{isomorfismo} de $U'$ sobre $U$.
    Como los $P_k$ son proyectivos sobre $S$, también lo es $P$ \sref{II.5.5.5-fr}[iv]; por tanto, $X\times_SP$ es proyectivo sobre $X$ \sref{II.5.5.5-fr}[iii], y también lo es $X'$, que es un subpreesquema cerrado de $X\times_SP$.
    Además, $f\circ\psi'=q_1\circ(h\circ\psi')=q_1\circ\psi=j$, de modo que $f(X')$ contiene el abierto denso $U$ de $X$; pero $f$ es un morfismo \emph{cerrado} \sref{II.5.5.3-fr}, por lo que $f(X')=X$.
    Observemos ahora que $q_1^{-1}(U)=U\times_SP$ está inducido sobre un abierto de $X\times_SP$ y que, por definición, el preesquema $U'=h^{-1}(U\times_SP)$ es el subpreesquema inducido por $X'$ sobre el abierto $U'$; es, por tanto, la clausura \emph{relativa a}
\oldpage[II]{108}
    $U\times_SP$ del preesquema $Z$ \sref[I]{I.9.5.8-fr}.
    Pero la inmersión $\psi$ se factoriza como $U\xrightarrow{\Gamma_\varphi}U\times_SP\xrightarrow{j\times1}X\times_SP$ y, como $P$ es separado sobre $S$, el morfismo grafo $\Gamma_\varphi$ es una inmersión cerrada \sref[I]{I.5.4.3-fr}; por tanto, $Z$ es un subpreesquema \emph{cerrado} de $U\times_SP$, de donde $U'=Z$.
    Como $\psi$ es una inmersión, la restricción de $f$ a $U'$ es un isomorfismo sobre $U$, inverso de $\psi'$; finalmente, por definición de $X'$, $U'$ es denso en $X'$.
% Fuente francesa: PMIHES_1961__8__5_0, printed p. 108.

  \item[(D)] Demostremos ahora que $g$ es una \emph{inmersión}; esto implicará que $X'$ es \emph{cuasiproyectivo} sobre $S$, pues $P$ es proyectivo sobre $S$.
    Pongamos
    \begin{align*}
V_k&=\varphi_k(U_k) &&\text{(abierto de $P_k$)}\\
W_k&=p_k^{-1}(V_k) &&\text{(abierto de $P$)}\\
&\left.\begin{aligned}U'_k&=f^{-1}(U_k)\\U''_k&=g^{-1}(W_k)\end{aligned}\right\}
&&\text{(abiertos de $X'$)}
\end{align*}
    Es claro que los $U'_k$ forman una cubierta abierta de $X'$; veremos primero que ocurre lo mismo con los $U''_k$, demostrando que $U'_k\subset U''_k$.
    Para ello, bastará demostrar que el diagrama
    \[
\tag{5.6.1.5}
\label{II.5.6.1.5-fr}
\xymatrix{
U'_k\ar[r]^{g|U'_k}\ar[d]_{f|U'_k} & P\ar[d]^{p_k}\\
U_k\ar[r]_{\varphi_k} & P_k
}
\]
    es conmutativo.
    Pero el preesquema $U'_k=h^{-1}(U_k\times_SP)$ es el subpreesquema inducido por $X'$ sobre el abierto $U'_k$ y es, por tanto, la clausura de $Z=U'\subset U'_k$ relativa a $U'_k$ \sref[I]{I.9.5.8-fr}.
    Para demostrar la conmutatividad de \sref{II.5.6.1.5-fr}, basta entonces (puesto que $P_k$ es un $S$-esquema) demostrar que, al componer el diagrama con la inyección canónica $U'\to U'_k$ (o, equivalentemente, gracias al isomorfismo de $U'$ con $U$, con $\psi$), se obtiene un diagrama conmutativo \sref[I]{I.9.5.6-fr}.
    Pero, por definición, el diagrama así obtenido no es otro que \sref{II.5.6.1.1-fr}, de donde nuestra afirmación.
% Fuente francesa: PMIHES_1961__8__5_0, printed pp. 108-109.

    Los $W_k$ forman, por tanto, una cubierta abierta de $g(X')$; para demostrar que $g$ es una inmersión, basta ver que cada una de las restricciones $g|U''_k$ es una inmersión en $W_k$ \sref[I]{I.4.2.4-fr}.
    Consideremos para ello el morfismo $u_k:W_k\xrightarrow{p_k}V_k\xrightarrow{\varphi_k^{-1}}U_k\to X$; como $X$ es separado sobre $S$, el morfismo grafo $\Gamma_{u_k}:W_k\to X\times_SW_k$ es una inmersión cerrada \sref[I]{I.5.4.3-fr}, de modo que el grafo $T_k=\Gamma_{u_k}(W_k)$ es un subpreesquema cerrado de $X\times_SW_k$;
    si mostramos que este subpreesquema contiene a $U'$, contendrá también al subpreesquema inducido por $X'$ sobre el abierto $U''_k$ de $X'$, por (\hyperref[I.9.5.8-fr]{\textbf{I}, 9.5.8}).
    Como la restricción de $q_2$ a $T_k$ es un isomorfismo sobre $W_k$, la restricción de $g$ a $U''_k$ será una inmersión en $W_k$, y nuestra afirmación quedará demostrada.
    Sea $v_k$ la inyección canónica $U'\to X\times_SW_k$; debemos demostrar que existe un morfismo $w_k:U'\to W_k$ tal que $v_k=\Gamma_{u_k}\circ w_k$.
    Por definición del producto, basta probar que $q_1\circ v_k=u_k\circ q_2\circ v_k$ \sref[I]{I.3.2.1-fr}, o bien, componiendo a la derecha
\oldpage[II]{109}
    con el isomorfismo $\psi':U\to U'$, que $q_1\circ\psi=u_k\circ q_2\circ\psi$.
    Pero, como $q_1\circ\psi=j$ y $q_2\circ\psi=\varphi$, nuestra afirmación se deduce de la conmutatividad de \sref{II.5.6.1.1-fr}, teniendo en cuenta la definición de $u_k$.
% Fuente francesa: PMIHES_1961__8__5_0, printed p. 109.

  \item[(E)] Es claro que, como $U$ y, por tanto, $U'$ son irreducibles, también lo es el $X'$ de la construcción precedente; el morfismo $f$ es entonces \emph{birracional} \sref[I]{I.2.2.9-fr}.
    Si, además, $X$ es reducido, también lo es $U'$, y por tanto $X'$ es reducido \sref[I]{I.9.5.9-fr}.
    Esto termina la demostración.
\end{enumerate}
\end{proof}
% Fuente francesa: PMIHES_1961__8__5_0, printed p. 109.

\begin{corollary}[5.6.2]
\label{II.5.6.2-fr}
Se supone satisfecha una de las hipótesis \emph{(a)} y \emph{(b)} de \sref{II.5.6.1-fr}.
Para que $X$ sea propio sobre $S$, es necesario y suficiente que exista un esquema $X'$ proyectivo sobre $S$ y un $S$-morfismo sobreyectivo $f:X'\to X$ (que es entonces proyectivo por \sref{II.5.5.5-fr}[v]).
Cuando esto ocurre, se puede elegir además $f$ de modo que exista un abierto denso $U$ de $X$ para el cual la restricción de $f$ a $f^{-1}(U)$ sea un isomorfismo $f^{-1}(U)\xrightarrow{\sim}U$ y $f^{-1}(U)$ sea denso en $X'$.
Si, además, $X$ es irreducible \textup{(resp.} reducido), se puede suponer que $X'$ también lo es; cuando $X$ y $X'$ son irreducibles, $f$ es un morfismo birracional.
\end{corollary}

\begin{proof}
La condición es suficiente por \sref{II.5.5.3-fr} y \sref{II.5.4.3-fr}[ii].
Es necesaria porque, con la notación de \sref{II.5.6.1-fr}, si $X$ es propio sobre $S$, entonces $X'$ es propio sobre $S$: en efecto, es proyectivo sobre $X$ y, por tanto, propio sobre $X$ \sref{II.5.5.3-fr}, y nuestra afirmación se deduce de \sref{II.5.4.2-fr}[ii]; además, como $X'$ es cuasiproyectivo sobre $S$, es proyectivo sobre $S$ por \sref{II.5.5.3-fr}.
\end{proof}
% Fuente francesa: PMIHES_1961__8__5_0, printed p. 109.

\begin{corollary}[5.6.3]
\label{II.5.6.3-fr}
Sean $S$ un preesquema localmente noetheriano y $X$ un $S$-esquema de tipo finito sobre $S$, con morfismo estructural $f_0:X\to S$.
Para que $X$ sea propio sobre $S$, es necesario y suficiente que, para todo morfismo \emph{de tipo finito} $S'\to S$, $(f_0)_{(S')}:X_{(S')}\to S'$ sea un morfismo cerrado.
Basta incluso que esta condición se verifique únicamente para todo $S$-preesquema de la forma $S'=S\otimes_{\mathbf Z}\mathbf Z[T_1,\ldots,T_n]$ (donde $T_i$ son indeterminadas).
\end{corollary}

\begin{proof}
Como la condición es claramente necesaria, demostremos que es suficiente.
Puesto que la cuestión es local sobre $S$ y $S'$ \sref{II.5.4.1-fr}, se puede suponer que $S$ y $S'$ son afines y noetherianos.
Por el lema de Chow, existe un $S$-esquema proyectivo $P$, una inmersión $j:X'\to P$ y un morfismo proyectivo y sobreyectivo $f:X'\to X$, tales que el diagrama
\[
\xymatrix{
X\ar[d]_{f_0} & X'\ar[l]_f\ar[d]^j\\
S & P\ar[l]^r
}
\]
es conmutativo.
Como $P$ es de tipo finito sobre $S$, la primera hipótesis implica que la proyección $q_2:X\times_SP\to P$ es un morfismo \emph{cerrado}.
Pero la inmersión $j$ se factoriza en $j=q_2\circ(f\times1)\circ\Gamma_j$, donde $\Gamma_j:X'\to X'\times_SP$ es el morfismo grafo, que es una inmersión cerrada puesto que $P$ es separado sobre $S$ (\hyperref[I.5.4.3-fr]{I, 5.4.3}). Ahora bien, $f$, al ser proyectivo, es propio (\hyperref[II.5.5.3-fr]{5.5.3}), por lo que $f\times1$ es cerrado. Se concluye que $j$ es una inmersión cerrada y, por consiguiente, es propio (\hyperref[II.5.4.2-fr]{5.4.2, (i)}).

Además, el morfismo estructural $r:P\to S$ es proyectivo y, por tanto, propio \sref{II.5.5.3-fr}; así, $f_0\circ f=r\circ j$ es propio \sref{II.5.4.2-fr}[ii];
finalmente, como $f$ es sobreyectivo, $f_0$ es propio por \sref{II.5.4.3-fr}.
% Fuente francesa: PMIHES_1961__8__5_0, printed pp. 109-110.

Para demostrar la proposición utilizando únicamente la segunda hipótesis, basta mostrar que esta implica la primera.
Pero, si $S'$ es afín y de tipo finito sobre $S=\operatorname{Spec}(A)$,
\oldpage[II]{110}
se tiene $S'=\operatorname{Spec}(A[c_1,\ldots,c_n])$ \sref[I]{I.6.3.3-fr}, de modo que $S'$ es isomorfo a un subpreesquema cerrado de $S''=\operatorname{Spec}(A[T_1,\ldots,T_n])$ (donde $T_i$ son indeterminadas).
En el diagrama conmutativo
\[
\xymatrix{
X\times_SS'\ar[r]^{1_X\times j}\ar[d]_{(f_0)_{(S')}} &
X\times_SS''\ar[d]^{(f_0)_{(S'')}}\\
S'\ar[r]_j & S''
}
\]
tanto $j$ como $1_X\times j$ son inmersiones cerradas \sref[I]{I.4.3.1-fr}, y $(f_0)_{(S'')}$ es cerrado por hipótesis; por tanto, $(f_0)_{(S')}$ también es cerrado.
\end{proof}
% Fuente francesa: PMIHES_1961__8__5_0, printed p. 110.

\section{Morfismos enteros y morfismos finitos}
\label{section:II.6-fr}


\subsection{Preesquemas enteros sobre otro preesquema}
\label{subsection:II.6.1-fr}

\begin{definition}[6.1.1]
\label{II.6.1.1-fr}
Sea $X$ un $S$-preesquema, con morfismo estructural $f:X\to S$.
Se dice que $X$ es entero sobre $S$, o que $f$ es un morfismo entero, si existe una cubierta $(S_\alpha)$ de $S$ por abiertos afines tal que, para todo $\alpha$, el preesquema inducido $f^{-1}(S_\alpha)$ sea un esquema afín cuyo anillo $B_\alpha$ sea un álgebra entera sobre el anillo $A_\alpha$ de $S_\alpha$.
Se dice que $X$ es finito sobre $S$, o que $f$ es un morfismo finito, si $X$ es entero sobre $S$ y de tipo finito sobre $S$.
\end{definition}

Si $S$ es afín, de anillo $A$, se dice también «entero (resp. finito) sobre $A$» en lugar de «entero (resp. finito) sobre $S$».

\begin{env}[6.1.2]
\label{II.6.1.2-fr}
Es claro que, si $X$ es entero sobre $S$, entonces es \emph{afín} sobre $S$.
Para que un preesquema $X$ afín sobre $S$ sea entero (resp. finito) sobre $S$, es necesario y suficiente que la $\mathcal{O}_S$-álgebra cuasicoherente asociada $\mathcal{A}(X)$ sea tal que exista una cubierta $(S_\alpha)$ de $S$ por abiertos afines con la propiedad de que, para todo $\alpha$, $\Gamma(S_\alpha,\mathcal{A}(X))$ sea un álgebra entera (resp. entera y de tipo finito) sobre $\Gamma(S_\alpha,\mathcal{O}_S)$.
Una $\mathcal{O}_S$-álgebra cuasicoherente con esta propiedad se dice \emph{entera} (resp. \emph{finita}) sobre $\mathcal{O}_S$.
Dar un preesquema entero (resp. finito) sobre $S$ equivale, pues, \sref{II.1.3.1-fr} a dar una $\mathcal{O}_S$-álgebra cuasicoherente entera (resp. finita) sobre $\mathcal{O}_S$.
Se señala que una $\mathcal{O}_S$-álgebra cuasicoherente $\mathcal{B}$ es finita si y solo si es un \emph{$\mathcal{O}_S$-módulo de tipo finito} \sref[I]{I.1.3.9-fr};
equivalentemente, $\mathcal{B}$ es una $\mathcal{O}_S$-álgebra \emph{entera} y \emph{de tipo finito}, pues un álgebra entera y de tipo finito sobre un anillo $A$ es un $A$-módulo de tipo finito.
\end{env}
% Fuente francesa: PMIHES_1961__8__5_0, printed p. 110.

\begin{proposition}[6.1.3]
\label{II.6.1.3-fr}
Sea $S$ un preesquema localmente noetheriano.
Para que un preesquema $X$ afín sobre $S$ sea finito sobre $S$, es necesario y suficiente que la $\mathcal{O}_S$-álgebra $\mathcal{A}(X)$ sea coherente.
\end{proposition}

\begin{proof}
Teniendo en cuenta la observación precedente, esto se reduce a señalar que, si $S$ es localmente noetheriano, los $\mathcal{O}_S$-módulos cuasicoherentes de tipo finito son precisamente los $\mathcal{O}_S$-módulos coherentes \sref[I]{I.1.5.1-fr}.
\end{proof}

\begin{proposition}[6.1.4]
\label{II.6.1.4-fr}
Sea $X$ un preesquema entero \textup{(resp. \emph{finito})} sobre $S$, con morfismo estructural $f:X\to S$.
Entonces, para todo abierto afín $U\subset S$, de anillo $A$, $f^{-1}(U)$ es un esquema afín cuyo anillo $B$ es un álgebra entera \textup{(resp. \emph{finita})} sobre $A$.
\end{proposition}
% Fuente francesa: PMIHES_1961__8__5_0, printed p. 111.

\oldpage[II]{111}
\begin{proof}
Demostremos primero el lema siguiente:

  \begin{lemma}[6.1.4.1]
  \label{II.6.1.4.1-fr}
  Sean $A$ un anillo, $M$ un $A$-módulo y $(g_i)_{1\leq i\leq m}$ una familia finita de elementos de $A$ tal que los $D(g_i)$ ($1\leq i\leq m$) cubren $\operatorname{Spec}(A)$.
  Si, para todo $i$, $M_{g_i}$ es un $A_{g_i}$-módulo de tipo finito, entonces $M$ es un $A$-módulo de tipo finito.
  \end{lemma}

  \begin{proof}
  Se puede suponer que $M_{g_i}$ admite una familia finita de generadores $(m_{ij}/g_i^n)$, con $m_{ij}\in M$ y siendo $n$ el mismo para todos los índices $i$.
  Mostraremos que los $m_{ij}$ forman una familia de generadores de $M$.
  Sea $M'$ el sub-$A$-módulo de $M$ generado por estos elementos y sea $m$ un elemento de $M$.
  Por hipótesis, para cada $i$ existen $a_{ij}\in A$ y un entero $p$ (independiente de $i$) tales que, en $M_{g_i}$, $m/1=(\sum_j a_{ij}m_{ij})/g_i^p$;
  esto implica que existe un entero $r\geq p$ tal que, para todo $i$, se tiene $g_i^r m\in M'$.
  Ahora bien, como los $D(g_i^r)=D(g_i)$ cubren $\operatorname{Spec}(A)$, el ideal de $A$ generado por los $g_i^r$ es igual a $A$; dicho de otro modo, existen elementos $a_i\in A$ tales que $\sum_i a_i g_i^r=1$.
  Por tanto, $m=(\sum_i a_i g_i^r)m\in M'$, de donde se sigue el lema.
  \end{proof}
% Fuente francesa: PMIHES_1961__8__5_0, printed p. 111.

Ahora bien, ya se sabe \sref{II.1.3.2-fr} que $f^{-1}(U)$ es afín.
Si $\varphi$ es el homomorfismo $A\to B$ correspondiente a $f$, existe una cubierta finita de $U$ por abiertos $D(g_i)$ (donde $g_i\in A$) tal que, si $h_i=\varphi(g_i)$, $B_{h_i}$ sea un álgebra entera (resp. entera y finita) sobre $A_{g_i}$.
En efecto, existe una cubierta de $U$ por abiertos afines $V_\alpha\subset U$ tal que, si $A_\alpha=A(V_\alpha)$, $B_\alpha=A(f^{-1}(V_\alpha))$ sea un álgebra entera (resp. finita) sobre $A_\alpha$.
Todo $x\in U$ pertenece a algún $V_\alpha$, por lo que existe $g\in A$ tal que $x\in D(g)\subset V_\alpha$;
si $g_\alpha$ es la imagen de $g$ en $A_\alpha$, entonces $A(D(g))=A_g=(A_\alpha)_{g_\alpha}$;
sea $h=\varphi(g)$, y sea $h_\alpha$ la imagen de $g_\alpha$ en $B_\alpha$;
se tiene
\[
A(D(h))=B_h=(B_\alpha)_{h_\alpha}
\]
y, como $B_\alpha$ es entera (resp. finita) sobre $A_\alpha$, $(B_\alpha)_{h_\alpha}$ es entera (resp. finita) sobre $(A_\alpha)_{g_\alpha}$.
Basta ahora utilizar que $U$ es cuasicompacto para obtener la cubierta buscada.

Se supone primero que los $B_{h_i}$ son enteros y finitos sobre los $A_{g_i}$.
Como, en cuanto $A_{g_i}$-módulo, $B_{h_i}$ también se escribe $B_{g_i}$, el lema~\sref{II.6.1.4.1-fr} muestra que, en este caso, $B$ es un $A$-módulo de tipo finito.

Supongamos ahora únicamente que cada $B_{h_i}$ es entero sobre $A_{g_i}$;
sea $b\in B$, y sea $C$ la sub-$A$-álgebra de $B$ generada por $b$.
Para todo $i$, $C_{h_i}$ es el álgebra sobre $A_{g_i}$ generada por $b/1$ en $B_{h_i}$;
de la hipótesis se sigue que cada $C_{h_i}$ es un $A_{g_i}$-módulo de tipo finito y, por tanto, \sref{II.6.1.4.1-fr} $C$ es un $A$-módulo de tipo finito, lo que prueba que $B$ es entero sobre $A$.
\end{proof}
% Fuente francesa: PMIHES_1961__8__5_0, printed p. 111.

\begin{proposition}[6.1.5]
\label{II.6.1.5-fr}
\medskip\noindent
\begin{enumerate}
  \item[{\rm(i)}] Una inmersión cerrada es finita (y, \emph{a fortiori}, entera).
  \item[{\rm(ii)}] La composición de dos morfismos finitos \textup{(resp. \emph{enteros})} es finita \textup{(resp. \emph{entera})}.
  \item[{\rm(iii)}] Si $f:X\to Y$ es un $S$-morfismo finito \textup{(resp. \emph{entero})}, entonces $f_{(S')}:X_{(S')}\to Y_{(S')}$ es finito \textup{(resp. \emph{entero})} para toda extensión $S'\to S$ del preesquema de base.
  \item[{\rm(iv)}] Si $f:X\to Y$ y $g:X'\to Y'$ son morfismos finitos \textup{(resp. \emph{enteros})} de $S$-preesquemas, entonces $f\times_S g:X\times_S X'\to Y\times_S Y'$ es finito (resp. entero).
  \item[{\rm(v)}] Si $f:X\to Y$ y $g:Y\to Z$ son morfismos tales que $g\circ
f$ es finito \textup{(resp. \emph{entero})}, y si $g$ es separado, entonces $f$ es finito \textup{(resp. \emph{entero})}.
  \item[{\rm(vi)}] Si $f:X\to Y$ es un morfismo finito \textup{(resp. \emph{entero})}, entonces $f_{\mathrm{red}}$ es finito (resp. entero).
\end{enumerate}
\end{proposition}
% Fuente francesa: PMIHES_1961__8__5_0, printed p. 112.

\oldpage[II]{112}
\begin{proof}
En virtud de \sref[I]{I.5.5.12-fr}, basta demostrar (i), (ii) y (iii).
Para probar que una inmersión cerrada $X\to S$ es finita, se puede limitar al caso en que $S=\operatorname{Spec}(A)$, y todo se reduce entonces a señalar que un anillo cociente $A/\mathfrak{J}$ es un $A$-módulo cíclico.
Para probar que la composición de dos morfismos finitos (resp. enteros) $X\to Y$, $Y\to Z$ es finita (resp. entera), se puede suponer de nuevo que $Z$ (y, por tanto, $X$ e $Y$ \sref{II.1.3.4-fr}) es afín; la afirmación equivale entonces a decir que, si $B$ es un $A$-álgebra finita (resp. entera) y $C$ un $B$-álgebra finita (resp. entera), entonces $C$ es un $A$-álgebra finita (resp. entera), lo que es inmediato.
Por último, para demostrar (iii), se puede limitar al caso en que $S=Y$, puesto que $X_{(S')}$ se identifica con $X\times_Y Y_{(S')}$ \sref[I]{I.3.3.11-fr};
se puede suponer además que $S=\operatorname{Spec}(A)$ y $S'=\operatorname{Spec}(A')$.
Entonces $X$ es afín, de anillo $B$ \sref{II.1.3.4-fr}, y $X_{(S')}$ es afín, de anillo $A'\otimes_A B$; basta señalar que, si $B$ es un $A$-álgebra finita (resp. entera), entonces $A'\otimes_A B$ es un $A'$-álgebra finita (resp. entera).
\end{proof}
% Fuente francesa: PMIHES_1961__8__5_0, printed p. 112.

Señalemos también que, si $X$ e $Y$ son $S$-preesquemas finitos (resp. enteros) sobre $S$, su \emph{suma} $X\amalg Y$ es un preesquema finito (resp. entero) sobre $S$, pues esto se reduce a mostrar que, si $B$ y $C$ son $A$-álgebras finitas (resp. enteras) sobre $A$, entonces también lo es $B\times C$.

\begin{corollary}[6.1.6]
\label{II.6.1.6-fr}
Si $X$ es un preesquema entero \textup{(resp. \emph{finito})} sobre $S$, entonces, para todo abierto $U\subset S$, $f^{-1}(U)$ es entero \textup{(resp. \emph{finito})} sobre $U$.
\end{corollary}

\begin{proof}
Es un caso particular de \sref{II.6.1.5-fr}[(iii)].
\end{proof}

\begin{corollary}[6.1.7]
\label{II.6.1.7-fr}
Sea $f:X\to Y$ un morfismo finito.
Entonces, para todo $y\in Y$, la fibra $f^{-1}(y)$ es un esquema algebraico finito sobre $\boldsymbol{k}(y)$ y, \emph{a fortiori}, su espacio subyacente es discreto y finito.
\end{corollary}

\begin{proof}
En efecto, en cuanto $\boldsymbol{k}(y)$-preesquema, $f^{-1}(y)$ se identifica con $X\times_Y\operatorname{Spec}(\boldsymbol{k}(y))$ \sref[I]{I.3.6.1-fr}, que es finito sobre $\operatorname{Spec}(\boldsymbol{k}(y))$ \sref{II.6.1.5-fr}[(iii)];
es, por tanto, un esquema afín cuyo anillo es un álgebra de rango finito sobre $\boldsymbol{k}(y)$ \sref{II.6.1.4-fr}.
El corolario se sigue entonces de \sref[I]{I.6.4.4-fr}.
\end{proof}

\begin{corollary}[6.1.8]
\label{II.6.1.8-fr}
Sean $X$ y $S$ preesquemas íntegros, y $f:X\to S$ un morfismo \emph{dominante}.
Si $f$ es entero \textup{(resp. \emph{finito})}, el cuerpo $R(X)$ de funciones racionales sobre $X$ es algebraico \textup{(resp. \emph{algebraico de grado finito})} sobre el cuerpo $R(S)$ de funciones racionales sobre $S$.
\end{corollary}

\begin{proof}
Sea $s$ el punto genérico de $S$;
el $\boldsymbol{k}(s)$-preesquema $f^{-1}(s)$ es entero (resp. finito) sobre $\operatorname{Spec}(\boldsymbol{k}(s))$ \sref{II.6.1.5-fr}[(iii)] y contiene, por hipótesis, el punto genérico $x$ de $X$;
el anillo local de $x$ en $f^{-1}(s)$ es igual a $\boldsymbol{k}(x)$ \sref[I]{I.3.6.5-fr} y es, por tanto, un anillo local de un álgebra entera (resp. finita) sobre $\boldsymbol{k}(s)$ \sref{II.6.1.4-fr}, de donde se sigue el corolario.
\end{proof}
% Fuente francesa: PMIHES_1961__8__5_0, printed p. 112.

\begin{remark}[6.1.9]
\label{II.6.1.9-fr}
La hipótesis de que $g$ sea \emph{separado} es esencial para la validez de \sref{II.6.1.5-fr}[(v)]: si $Y$ no es separado sobre $Z$, la identidad $1_Y$ es el morfismo compuesto $Y\xrightarrow{\Delta_Y}Y\times_ZY\xrightarrow{p_1}Y$, pero $\Delta_Y$ no es un morfismo entero, como se sigue de \sref{II.6.1.10-fr}:
\end{remark}

\begin{proposition}[6.1.10]
\label{II.6.1.10-fr}
Todo morfismo entero es universalmente cerrado.
\end{proposition}

\begin{proof}
Sea $f:X\to Y$ un morfismo entero;
en virtud de \sref{II.6.1.5-fr}[(iii)], basta demostrar que $f$ es \emph{cerrado}.
Sea $Z$ un subconjunto cerrado de $X$;
% Fuente francesa: PMIHES_1961__8__5_0, printed p. 113.


existe entonces un subpreesquema de $X$ 
\oldpage[II]{113}
cuyo espacio subyacente es $Z$ \sref[I]{I.5.2.1-fr}, y se sigue de \sref{II.6.1.5-fr}[(i) y (ii)] que basta probar que $f(X)$ es \emph{cerrado} en $Y$.
En virtud de \sref{II.6.1.5-fr}[(vi)], se puede suponer que $X$ e $Y$ son \emph{reducidos};
además, si $T$ es el subpreesquema cerrado reducido de $Y$ cuyo espacio subyacente es $\overline{f(X)}$ \sref[I]{I.5.2.1-fr}, se sabe que $f$ se factoriza como $X\xrightarrow{g}T\xrightarrow{j}Y$, donde $j$ es el morfismo de inyección \sref[I]{I.5.2.2-fr}; como $j$ es separado \sref[I]{I.5.5.1-fr}[(i)], se sigue de \sref{II.6.1.5-fr}[(v)] que $g$ es un morfismo entero.
Se puede suponer, pues, que $f(X)$ es \emph{denso} en $Y$.
Finalmente, como la cuestión es local sobre $Y$, se puede limitar al caso en que $Y=\operatorname{Spec}(A)$.
Entonces $X=\operatorname{Spec}(B)$, donde $B$ es un $A$-álgebra entera sobre $A$ \sref{II.6.1.4-fr};
además, $A$ es reducido \sref[I]{I.5.1.4-fr}, y la hipótesis de que $f(X)$ sea denso en $Y$ implica que el homomorfismo $\varphi:A\to B$ correspondiente a $f$ es \emph{inyectivo} \sref[I]{I.1.2.7-fr}.
En estas condiciones, decir que $f(X)=Y$ equivale a decir que todo ideal primo de $A$ es la intersección con $A$ de un ideal primo de $B$, lo cual es precisamente el primer teorema de Cohen--Seidenberg (\cite[t.~I, p.~257, th.~3]{I-13}).
\end{proof}

\begin{corollary}[6.1.11]
\label{II.6.1.11-fr}
Todo morfismo finito $f:X\to Y$ es proyectivo.
\end{corollary}

\begin{proof}
Como $f$ es afín, $\mathcal{O}_X$ es un $\mathcal{O}_X$-módulo \emph{muy amplio} relativamente a $f$ \sref{II.5.1.2-fr};
además, $f_*(\mathcal{O}_X)$ es un $\mathcal{O}_Y$-módulo cuasicoherente \emph{de tipo finito} \sref{II.6.1.2-fr};
por último, $f$ es separado, de tipo finito y universalmente cerrado \sref{II.6.1.10-fr}, y satisface, por tanto, las condiciones del criterio~\sref{II.5.5.4-fr}[(i)].
\end{proof}
% Fuente francesa: PMIHES_1961__8__5_0, printed p. 113.

\begin{proposition}[6.1.12]
\label{II.6.1.12-fr}
Sea $f:X'\to X$ un morfismo finito, y sea $\mathcal{B}=f_*(\mathcal{O}_{X'})$ (que es una $\mathcal{O}_X$-álgebra cuasicoherente y un $\mathcal{O}_X$-módulo de tipo finito).
Sea $\mathcal{F}'$ un $\mathcal{O}_{X'}$-módulo cuasicoherente;
para que $\mathcal{F}'$ sea localmente libre de rango~$r$, es necesario y suficiente que $f_*(\mathcal{F}')$ sea un $\mathcal{B}$-módulo localmente libre de rango~$r$.
\end{proposition}

\begin{proof}
Es claro que, si $f_*(\mathcal{F}')|U$ es isomorfo a $\mathcal{B}^r|U$ ($U$ abierto de $X$), entonces $\mathcal{F}'|f^{-1}(U)$ es isomorfo a $\mathcal{O}_{X'}^r|f^{-1}(U)$ \sref{II.1.4.2-fr}.
Recíprocamente, supongamos que $\mathcal{F}'$ es localmente libre de rango~$r$; mostraremos que $f_*(\mathcal{F}')$ es localmente isomorfo a $\mathcal{B}^r$ en cuanto $\mathcal{B}$-módulo.
Sea $x$ un punto de $X$;
cuando $U$ recorre un sistema fundamental de vecindades afines de $x$, $f^{-1}(U)$ recorre un sistema fundamental de vecindades afines \sref{II.1.2.5-fr} del conjunto finito $f^{-1}(x)$, pues $f$ es cerrado \sref{II.6.1.10-fr}.
La proposición se sigue entonces del lema siguiente:
\end{proof}

\begin{lemma}[6.1.12.1]
\label{II.6.1.12.1-fr}
Sean $Y$ un preesquema, $\mathcal{E}$ un $\mathcal{O}_Y$-módulo localmente libre de rango~$r$, y $Z$ un subconjunto finito de $Y$ contenido en un abierto afín $V$.
Entonces existe una vecindad $U\subset V$ de $Z$ tal que $\mathcal{E}|U$ sea isomorfo a $\mathcal{O}_Y^r|U$.
\end{lemma}
% Fuente francesa: PMIHES_1961__8__5_0, printed pp. 113-114.

\begin{proof}
Se puede suponer evidentemente que $Y$ es afín;
para todo $z_i\in Z$, existe en la clausura $\overline{\{z_i\}}$ al menos un punto cerrado $z'_i$ \sref[0]{0.2.1.3-fr};
si $Z'$ es el conjunto de los $z'_i$, toda vecindad de $Z'$ es una vecindad de $Z$, y se puede suponer, por tanto, que $Z$ es discreto y cerrado en $Y$.
Consideremos el subpreesquema cerrado reducido de $Y$ cuyo espacio subyacente es $Z$ \sref[I]{I.5.2.1-fr}, y sea $j:Z\to Y$ la inyección canónica;
$j^*(\mathcal{E})=\mathcal{E}\otimes_Y
\mathcal{O}_Z$ es localmente libre de rango~$r$ sobre el esquema discreto $Z$ y es, por tanto, isomorfo a $\mathcal{O}_Z^r$;
dicho de otro modo, existen $r$ secciones $s_i$ ($1\leq i\leq r$) de $\mathcal{E}\otimes_Y\mathcal{O}_Z$ sobre $Z$ tales que el homomorfismo $\mathcal{O}_Z^r\to\mathcal{E}\otimes_Y\mathcal{O}_Z$ definido por estas secciones sea biyectivo.
Pero $Y=\operatorname{Spec}(A)$ es afín, $Z$ está definido por un ideal $\mathfrak{J}$ de $A$, y se tiene $\mathcal{E}=\widetilde{M}$, donde $M$ es un $A$-módulo;
los $s_i$ son elementos 
\oldpage[II]{114}
de $M\otimes_A(A/\mathfrak{J})$ y, por tanto, son imágenes de $r$ elementos $t_i\in M=\Gamma(Y,\mathcal{E})$.

Para todo $z_j\in Z$, existe entonces una vecindad $V_j$ de $z_j$ tal que las restricciones de los $t_i$ a $V_j$ definan un isomorfismo $\mathcal{O}_Y^r|V_j\to\mathcal{E}|V_j$ \sref[0]{0.5.5.4-fr};
la vecindad $U$ dada por la unión de los $V_j$ es la vecindad buscada.
\end{proof}
% Fuente francesa: PMIHES_1961__8__5_0, printed p. 114.

\begin{proposition}[6.1.13]
\label{II.6.1.13-fr}
Sean $g:X'\to X$ un morfismo entero de preesquemas, $Y$ un preesquema normal y localmente íntegro, y $f$ una aplicación racional de $Y$ en $X'$ tal que $g\circ f$ sea una aplicación racional definida en todas partes \textup{\sref[I]{I.7.2.1-fr}}.
Entonces $f$ está definida en todas partes.
\end{proposition}

\begin{proof}
Si $f_1$ y $f_2$ son morfismos (de abiertos densos de $Y$ en $X'$) pertenecientes a la clase $f$, es claro que $g\circ f_1$ y $g\circ f_2$ son morfismos equivalentes, lo que justifica la notación $g\circ f$ para su clase de equivalencia.
Se recuerda también que, si se supone además que $Y$ es \emph{localmente noetheriano}, la hipótesis de que $Y$ sea normal implica ya que $Y$ es localmente íntegro \sref[I]{I.6.1.13-fr}.

Para demostrar \sref{II.6.1.13-fr}, señalemos primero que la cuestión es local sobre $Y$, de modo que se puede suponer que en la clase $g\circ f$ existe un \emph{morfismo} $h:Y\to X$.
Consideremos la imagen inversa $Y'=X'_{(h)}=X'_{(Y)}$, y señalemos que el morfismo $g'=g_{(Y)}:Y'\to Y$ es \emph{entero} \sref{II.6.1.5-fr}[(iii)].
En vista de la correspondencia entre las aplicaciones racionales de $Y$ en $X'$ y las $Y$-secciones racionales de $Y'$ \sref[I]{I.7.1.2-fr}, basta probar el caso particular de \sref{II.6.1.13-fr} en que $X=Y$; dicho de otro modo, el siguiente:
\end{proof}

\begin{corollary}[6.1.14]
\label{II.6.1.14-fr}
Sean $X$ un preesquema normal y localmente íntegro, $g:X'\to X$ un morfismo \emph{entero}, y $f$ una $X$-sección racional de $X'$.
Entonces $f$ está definida en todas partes.
\end{corollary}

\begin{proof}
Como la cuestión es local sobre $X$, se puede suponer que $X$ es íntegro, y entonces $f$ se identifica con un morfismo de un abierto $U$ de $X$ en $X'$ \sref[I]{I.7.2.2-fr} que es una $U$-sección de $g^{-1}(U)$.
Como $g$ es separado, $f$ es una inmersión cerrada de $U$ en $g^{-1}(U)$ \sref[I]{I.5.4.6-fr};
sea $Z$ el subpreesquema cerrado de $g^{-1}(U)$ asociado a $f$ \sref[I]{I.4.2.1-fr}, que es isomorfo a $U$ y, por tanto, íntegro;
sea $X_1$ el subpreesquema reducido de $X'$ cuyo espacio subyacente es la clausura $\overline{Z}$ de $Z$ en $X'$ \sref[I]{I.5.2.1-fr};
$Z$ es entonces un subpreesquema inducido sobre un abierto de $X_1$ \sref[I]{I.5.2.3-fr} y, como es irreducible, también lo es $X_1$, que es, por tanto, íntegro.
El morfismo $f$ puede considerarse entonces como una $X$-sección racional de $X_1$;
como la restricción de $g$ a $X_1$ es un morfismo entero \sref{II.6.1.5-fr}[(i) y (ii)], se reduce finalmente a demostrar \sref{II.6.1.14-fr} en el caso particular en que $X'=X_1$; dicho de otro modo, el siguiente:
\end{proof}
% Fuente francesa: PMIHES_1961__8__5_0, printed p. 114.

\begin{corollary}[6.1.15]
\label{II.6.1.15-fr}
Sean $X$ un preesquema íntegro normal, $X'$ un preesquema íntegro y $g:X'\to X$ un morfismo \emph{entero}.
Si existe una $X$-sección racional $f$ de $X'$, entonces $g$ es un isomorfismo.
\end{corollary}

\begin{proof}
Como la cuestión es local sobre $X$, se puede suponer que $X$ es afín, de anillo íntegro $A$, y entonces $X'$ es afín, de anillo íntegro $A'$ entero sobre $A$ \sref{II.6.1.4-fr};
además, el razonamiento de \sref{II.6.1.14-fr} muestra que existe un abierto denso de $X$ isomorfo a un abierto denso de $X'$, de modo que $A$ y $A'$ tienen el mismo cuerpo de fracciones.
Por otra parte, en virtud de \sref[I]{I.8.2.1.1-fr} y de la hipótesis de que los $\mathcal{O}_x$ son íntegramente cerrados, el anillo $A$ es íntegramente cerrado; por tanto, $A'=A$, lo que termina la demostración de \sref{II.6.1.13-fr}.
\end{proof}
% Fuente francesa: PMIHES_1961__8__5_0, printed pp. 114-115.

\subsection{Morfismos cuasifinitos}
\label{subsection:II.6.2-fr}

\begin{proposition}[6.2.1]
\label{II.6.2.1-fr}
Sean $f:X\to Y$ un morfismo localmente de tipo finito y $x$ un punto de $X$.
Las condiciones siguientes son equivalentes:
\oldpage[II]{115}
\begin{enumerate}
  \item[{\rm(a)}] El punto $x$ está aislado en su fibra $f^{-1}(f(x))$.
  \item[{\rm(b)}] El anillo $\mathcal{O}_x$ es un $\mathcal{O}_{f(x)}$-módulo cuasifinito \textup{\sref[0]{0.7.4.1-fr}}.
\end{enumerate}
\end{proposition}

\begin{proof}
Como la cuestión es manifiestamente local sobre $X$ y sobre $Y$, se puede suponer que $X=\operatorname{Spec}(A)$ e $Y=\operatorname{Spec}(B)$ son afines, siendo $A$ un $B$-álgebra de tipo finito \sref[I]{I.6.3.3-fr}.
Además, se puede reemplazar $X$ por $X\times_Y\operatorname{Spec}(\mathcal{O}_{f(x)})$ sin modificar ni la fibra $f^{-1}(f(x))$ ni el anillo local $\mathcal{O}_x$ \sref[I]{I.3.6.5-fr};
se puede suponer, por tanto, que $B$ es un anillo local (igual a $\mathcal{O}_{f(x)}$).
Si $\mathfrak{n}$ es el ideal maximal de $B$, entonces $f^{-1}(f(x))$ es un esquema afín de anillo $A/\mathfrak{n}A$, de tipo finito sobre $\boldsymbol{k}(f(x))=B/\mathfrak{n}$ \sref[I]{I.6.4.11-fr}.
Dicho esto, si se satisface \emph{(a)}, se puede suponer además que $f^{-1}(f(x))$ se reduce al único punto $x$;
por tanto, $A/\mathfrak{n}A$ es de rango finito sobre $B/\mathfrak{n}$ \sref[I]{I.6.4.4-fr}, o, dicho de otro modo, $A$ es un $B$-módulo cuasifinito.
Recíprocamente, si ocurre así, entonces $f^{-1}(f(x))$ es un esquema afín artiniano y, por tanto, discreto \sref[I]{I.6.4.4-fr};
así, $x$ está aislado en su fibra, lo que muestra que \emph{(b)} implica \emph{(a)}.
\end{proof}
% Fuente francesa: PMIHES_1961__8__5_0, printed p. 115.

\begin{corollary}[6.2.2]
\label{II.6.2.2-fr}
Sea $f:X\to Y$ un morfismo de tipo finito.
Las condiciones siguientes son equivalentes:
\begin{enumerate}
  \item[{\rm(a)}] Todo punto $x\in X$ está aislado en su fibra $f^{-1}(f(x))$ \textup{(o, dicho de otro modo, el subespacio $f^{-1}(f(x))$ es \emph{discreto})}.
  \item[{\rm(b)}] Para todo $x\in X$, el preesquema $f^{-1}(f(x))$ es un $\boldsymbol{k}(f(x))$-preesquema finito.
  \item[{\rm(c)}] Para todo $x\in X$, el anillo $\mathcal{O}_x$ es un $\mathcal{O}_{f(x)}$-módulo cuasifinito.
\end{enumerate}
\end{corollary}

\begin{proof}
La equivalencia de \emph{(a)} y \emph{(c)} se sigue de \sref{II.6.2.1-fr}.
Como $f^{-1}(f(x))$ es un $\boldsymbol{k}(f(x))$-preesquema algebraico \sref[I]{I.6.4.11-fr}, la equivalencia de \emph{(a)} y \emph{(b)} se sigue de \sref[I]{I.6.4.4-fr}.
\end{proof}

\begin{definition}[6.2.3]
\label{II.6.2.3-fr}
Si $f:X\to Y$ es un morfismo de tipo finito que satisface las condiciones equivalentes de \textup{\sref{II.6.2.2-fr}}, se dice que $f$ es cuasifinito, o que $X$ es \emph{cuasifinito sobre $Y$}.
\end{definition}

Es claro que todo morfismo \emph{finito} es cuasifinito \sref{II.6.1.8-fr}.
% Fuente francesa: PMIHES_1961__8__5_0, printed p. 115.

\begin{proposition}[6.2.4]
\label{II.6.2.4-fr}
\medskip\noindent
\begin{enumerate}
  \item[{\rm(i)}] Una inmersión $X\to Y$ que sea cerrada, o tal que $X$ sea noetheriano, es un morfismo cuasifinito.
  \item[{\rm(ii)}] Si $f:X\to Y$ y $g:Y\to Z$ son morfismos cuasifinitos, entonces $g\circ f$ es cuasifinito.
  \item[{\rm(iii)}] Si $X$ e $Y$ son $S$-preesquemas y $f:X\to Y$ es un $S$-morfismo cuasifinito, entonces $f_{(S')}:X_{(S')}\to Y_{(S')}$ es cuasifinito para toda extensión de la base $g:S'\to S$.
  \item[{\rm(iv)}] Si $f:X\to Y$ y $g:X'\to Y'$ son $S$-morfismos cuasifinitos, entonces
    \[
f\times_S g:X\times_S X'\to Y\times_S Y'
\]
    es cuasifinito.
  \item[{\rm(v)}] Sean $f:X\to Y$ y $g:Y\to Z$ morfismos tales que $g\circ
f$ sea cuasifinito; si, además, $g$ es separado, o $X$ es noetheriano, o $X\times_ZY$ es localmente noetheriano, entonces $f$ es cuasifinito.
  \item[{\rm(vi)}] Si $f$ es cuasifinito, entonces $f_{\mathrm{red}}$ es cuasifinito.
\end{enumerate}
\end{proposition}
% Fuente francesa: PMIHES_1961__8__5_0, printed pp. 115-116.

\begin{proof}
Si $f:X\to Y$ es una inmersión, toda fibra tiene a lo sumo un punto, y la afirmación~(i) se sigue entonces de (\hyperref[I.6.3.4-fr]{\textbf{I}, 6.3.4, (i) y 6.3.5}).
Para demostrar (ii), se señala primero que $h=g\circ f$ es de tipo finito \sref[I]{I.6.3.4-fr}[\textup{(ii)}];
además, si $z=h(x)$ e $y=f(x)$, entonces $y$ está aislado en $g^{-1}(z)$, por lo que existe una vecindad abierta $V$ de $y$ en $Y$ que no contiene ningún punto de $g^{-1}(z)$ distinto de $y$;
por tanto, $f^{-1}(V)$ es una vecindad abierta de $x$ que no contiene ningún punto 
\oldpage[II]{116}
de los $f^{-1}(y')$, donde $y'\neq y$ pertenece a $g^{-1}(z)$;

como $x$ está aislado en $f^{-1}(y)$, está aislado en $h^{-1}(z)=f^{-1}(g^{-1}(z))$.
Para demostrar (iii), se puede reducir al caso en que $Y=S$ \sref[I]{I.3.3.11-fr};
se señala de nuevo, en primer lugar, que $f'=f_{(S')}$ es de tipo finito \sref[I]{I.6.3.4-fr}[(iii)];
además, si $x'\in X'=X_{(S')}$, y si se pone $y'=f'(x')$ e $y=g(y')$, entonces $f'^{-1}(y')$ se identifica con $f^{-1}(y)\otimes_{\boldsymbol{k}(y)}\boldsymbol{k}(y')$ \sref[I]{I.3.6.5-fr};
como $f^{-1}(y)$ es de rango finito sobre $\boldsymbol{k}(y)$ por hipótesis, $f'^{-1}(y')$ es de rango finito sobre $\boldsymbol{k}(y')$ y, por tanto, discreto.
Las afirmaciones (iv), (v) y (vi) se deducen de (i), (ii) y (iii) por el método general \sref[I]{I.5.5.12-fr}, salvo cuando la hipótesis de (v) no es «$g$ es separado»;
en estos casos, se señala primero que, si $x$ está aislado en $f^{-1}(g^{-1}(g(f(x))))$, está \emph{a fortiori} aislado en $f^{-1}(f(x))$;
el hecho de que $f$ sea de tipo finito se sigue entonces de \sref[I]{I.6.3.6-fr}.
\end{proof}
% Fuente francesa: PMIHES_1961__8__5_0, printed p. 116.

\begin{proposition}[6.2.5]
\label{II.6.2.5-fr}
Sean $A$ un anillo local noetheriano \emph{completo}, $X$ un $A$-esquema localmente de tipo finito, y $x$ un punto de $X$ situado sobre el punto cerrado $y$ de $Y=\operatorname{Spec}(A)$.
Supongamos que $x$ está aislado en su fibra $f^{-1}(y)$ (donde $f$ es el morfismo estructural $X\to Y$).
Entonces $\mathcal{O}_x$ es un $A$-módulo de tipo finito, y $X$ es $Y$-isomorfo a la suma \textup{(\textbf{I}, 3.1)} de $X'=\operatorname{Spec}(\mathcal{O}_x)$ (que es un $Y$-esquema finito) y un $A$-esquema $X''$.
\end{proposition}

\begin{proof}
Se sigue de \sref{II.6.2.1-fr} que $\mathcal{O}_x$ es un $A$-módulo cuasifinito.
Como $\mathcal{O}_x$ es noetheriano \sref[I]{I.6.3.7-fr}, y el homomorfismo $A\to\mathcal{O}_x$ es local, la hipótesis de que $A$ es \emph{completo} implica que $\mathcal{O}_x$ es un $A$-módulo \emph{de tipo finito} \sref[0]{0.7.4.3-fr}.
Sea $X'=\operatorname{Spec}(\mathcal{O}_x)$ el esquema local de $X$ en el punto $x$ \sref[I]{I.2.4.1-fr}, y sea $g:X'\to X$ el morfismo canónico.
Como el compuesto $X'\xrightarrow{g}X\xrightarrow{f}Y$ es finito \sref{II.6.1.1-fr}, y como $f$ es separado, $g$ es finito \sref{II.6.1.5-fr}[(v)]; por tanto, $g(X')$ es \emph{cerrado} en $X$ \sref{II.6.1.10-fr}.
Además, como $g$ es de tipo finito, $g$ es un isomorfismo local en el punto cerrado $x'$ de $X'$ por la definición de $g$ y \sref[I]{I.6.5.4-fr};
pero, como $X'$ es la única vecindad abierta de $x'$, esto implica que $g$ es una inmersión abierta y, por tanto, $g(X')$ es también \emph{abierto} en $X$, lo que termina la demostración.
\end{proof}
% Fuente francesa: PMIHES_1961__8__5_0, printed p. 116.

\begin{corollary}[6.2.6]
\label{II.6.2.6-fr}
Sean $A$ un anillo local noetheriano \emph{completo}, $Y=\operatorname{Spec}(A)$ y $f:X\to Y$ un morfismo separado y cuasifinito.
Entonces $X$ es $Y$-isomorfo a una suma $X'\amalg X''$, donde $X'$ es un $Y$-esquema finito y $X''$ un $Y$-esquema cuasifinito tal que, si $y$ es el punto cerrado de $Y$, entonces $X''\cap f^{-1}(y)=\varnothing$.
\end{corollary}

\begin{proof}
En efecto, la fibra $f^{-1}(y)$ es finita y discreta por hipótesis, y el corolario se sigue entonces, por inducción sobre el número de puntos de esta fibra, de \sref{II.6.2.5-fr}.
\end{proof}

\begin{remark}[6.2.7]
\label{II.6.2.7-fr}
En el Capítulo~IV veremos que, si $Y$ es \emph{localmente noetheriano}, todo morfismo separado y cuasifinito $X\to Y$ es necesariamente \emph{cuasiafín}.
\end{remark}
% Fuente francesa: PMIHES_1961__8__5_0, printed pp. 116-117.

\subsection{Clausura integral de un preesquema}
\label{subsection:II.6.3-fr}

\begin{proposition}[6.3.1]
\label{II.6.3.1-fr}
Sean $(X,\mathcal{A})$ un espacio anillado, $\mathcal{B}$ una $\mathcal{A}$-álgebra (conmutativa), y $f$ una sección de $\mathcal{B}$ sobre $X$.
Las propiedades siguientes son equivalentes:
\begin{enumerate}
  \item[{\rm(a)}] El sub-$\mathcal{A}$-módulo de $\mathcal{B}$ generado por los $f^n$, para $n\geq 0$ \textup{\sref[0]{0.5.1.1-fr}}, es de tipo finito.
  \item[{\rm(b)}] Existe una sub-$\mathcal{A}$-álgebra $\mathcal{C}$ de $\mathcal{B}$ que es un $\mathcal{A}$-módulo de tipo finito y tal que $f\in\Gamma(X,\mathcal{C})$.
  \item[{\rm(c)}] Para todo $x\in X$, $f_x$ es entero sobre $\mathcal{A}_x$.
\end{enumerate}
\end{proposition}

\oldpage[II]{117}
\begin{proof}
Como el sub-$\mathcal{A}$-módulo de $\mathcal{B}$ generado por los $f^n$ es una $\mathcal{A}$-álgebra, es claro que \emph{(a)} implica \emph{(b)}.
Por otra parte, \emph{(b)} implica que, para todo $x\in X$, el $\mathcal{A}_x$-módulo $\mathcal{C}_x$ es de tipo finito, lo que implica que todo elemento del álgebra $\mathcal{C}_x$, y en particular $f_x$, es entero sobre $\mathcal{A}_x$.
Finalmente, si, para algún $x\in X$, se tiene una relación de la forma
\[
f_x^n+(a_1)_xf_x^{n-1}+\ldots+(a_n)_x=0
\]
donde los $a_i$ ($1\leq i\leq n$) son secciones de $\mathcal{A}$ sobre una vecindad abierta $U$ de $x$, entonces la sección $f^n|U+a_1.f^{n-1}|U+\ldots+a_n$ es nula sobre una vecindad $V\subset U$ de $x$; se sigue inmediatamente que los $f^k|V$ ($k\geq0$) son combinaciones lineales con coeficientes en $\Gamma(V,\mathcal{A})$ de los $f^j|V$ para $0\leq j\leq n-1$.
Se concluye así que \emph{(c)} implica \emph{(a)}.
\end{proof}

Cuando se satisfacen las condiciones equivalentes de \sref{II.6.3.1-fr}, se dice que la sección $f$ es \emph{entera sobre $\mathcal{A}$}.
% Fuente francesa: PMIHES_1961__8__5_0, printed p. 117.

\begin{corollary}[6.3.2]
\label{II.6.3.2-fr}
Bajo las hipótesis de \textup{\sref{II.6.3.1-fr}}, existe un único sub-$\mathcal{A}$-módulo $\mathcal{A}'$ de $\mathcal{B}$ tal que, para todo $x\in X$, $\mathcal{A}'_x$ sea el conjunto de gérmenes $f_x\in\mathcal{B}_x$ enteros sobre $\mathcal{A}_x$.
Para todo abierto $U\subset X$, las secciones de $\mathcal{A}'$ sobre $U$ son las secciones $f\in\Gamma(U,\mathcal{B})$ enteras sobre $\mathcal{A}|U$.
\end{corollary}

\begin{proof}
La existencia de $\mathcal{A}'$ es inmediata, tomando como $\Gamma(U,\mathcal{A}')$ el conjunto de los $f\in\Gamma(U,\mathcal{B})$ tales que $f_x$ sea entero sobre $\mathcal{A}_x$ para todo $x\in U$.
La segunda afirmación se sigue inmediatamente de \sref{II.6.3.1-fr}.
\end{proof}

Es claro que $\mathcal{A}'$ es una sub-$\mathcal{A}$-álgebra de $\mathcal{B}$;
se dice que es la \emph{clausura integral de $\mathcal{A}$ en $\mathcal{B}$}.
% Fuente francesa: PMIHES_1961__8__5_0, printed p. 117.

\begin{env}[6.3.3]
\label{II.6.3.3-fr}
Sean $(X,\mathcal{A})$ e $(Y,\mathcal{B})$ dos espacios anillados, y sea
\[
g=(\psi,\theta):X\to Y
\]
un morfismo.
Sea $\mathcal{C}$ (resp. $\mathcal{D}$) una $\mathcal{A}$-álgebra (resp. una $\mathcal{B}$-álgebra), y sea
\[
u:\mathcal{D}\to\mathcal{C}
\]
un $g$-morfismo \sref[0]{0.4.4.1-fr}.
Entonces, si $\mathcal{A}'$ (resp. $\mathcal{B}'$) es la clausura integral de $\mathcal{A}$ (resp. $\mathcal{B}$) en $\mathcal{C}$ (resp. $\mathcal{D}$), la \emph{restricción} de $u$ a $\mathcal{B}'$ es un $g$-morfismo
\[
u':\mathcal{B}'\to\mathcal{A}'.
\]

En efecto, si $j$ es la inyección canónica $\mathcal{B}'\to\mathcal{D}$, basta mostrar que
\[
v=u^\sharp\circ g^*(j):g^*(\mathcal{B}')\to\mathcal{C}
\]
envía $g^*(\mathcal{B}')$ en $\mathcal{A}'$.
Pero todo elemento de $(g^*(\mathcal{B}'))_x=\mathcal{B}'_{\psi(x)}\otimes_{\mathcal{B}_{\psi(x)}}
\mathcal{A}_x$ es entero sobre $\mathcal{A}_x$ por la definición de $\mathcal{B}'$, y lo mismo ocurre, por tanto, con su imagen por $v_x$, lo que demuestra nuestra afirmación.
\end{env}
% Fuente francesa: PMIHES_1961__8__5_0, printed pp. 117-118.

\begin{proposition}[6.3.4]
\label{II.6.3.4-fr}
Sean $X$ un preesquema y $\mathcal{A}$ una $\mathcal{O}_X$-álgebra cuasicoherente.
Entonces la clausura integral $\mathcal{O}'_X$ de $\mathcal{O}_X$ en $\mathcal{A}$ es una $\mathcal{O}_X$-álgebra cuasicoherente y, para todo abierto afín $U$ de $X$, $\Gamma(U,\mathcal{O}'_X)$ es la clausura integral de $\Gamma(U,\mathcal{O}_X)$ en $\Gamma(U,\mathcal{A})$.
\end{proposition}

\begin{proof}
Se puede limitar al caso en que $X=\operatorname{Spec}(B)$ es afín y $\mathcal{A}=\widetilde{A}$, donde $A$ es un
\oldpage[II]{118}
 $B$-álgebra;

sea $B'$ la clausura integral de $B$ en $A$.
Todo se reduce entonces a probar que, para todo $x\in X$, todo elemento de $A_x$ entero sobre $B_x$ pertenece necesariamente a $B'_x$, lo que se sigue de que, para todo anillo conmutativo $C$, las operaciones de tomar la clausura integral en un $C$-álgebra y de pasar a un anillo de fracciones (respecto de un subconjunto multiplicativo de $C$) conmutan \cite[t.~I, pp.~261 y 257]{I-13}.
\end{proof}

El $X$-esquema $X'=\operatorname{Spec}(\mathcal{O}'_X)$ se denomina entonces la \emph{clausura integral de $X$ respecto de $\mathcal{A}$} (o \emph{respecto de $\operatorname{Spec}(\mathcal{A})$});
es claro que $X'$ es \emph{entero} sobre $X$ \sref{II.6.1.2-fr}.

Se deduce inmediatamente de \sref{II.6.3.4-fr} que, si $f:X'\to X$ es el morfismo estructural, entonces, para todo abierto $U$ de $X$, $f^{-1}(U)$ es la \emph{clausura integral del preesquema inducido por $X$ sobre $U$, respecto de $\mathcal{A}|U$}.
% Fuente francesa: PMIHES_1961__8__5_0, printed p. 118.

\begin{env}[6.3.5]
\label{II.6.3.5-fr}
Sean $X$ e $Y$ preesquemas, $f:X\to Y$ un morfismo, $\mathcal{A}$ (resp. $\mathcal{B}$) una $\mathcal{O}_X$-álgebra cuasicoherente (resp. una $\mathcal{O}_Y$-álgebra cuasicoherente), y sea $u:\mathcal{B}\to\mathcal{A}$ un $f$-morfismo.
Se ha visto \sref{II.6.3.3-fr} que se obtiene un $f$-morfismo $u':\mathcal{O}'_Y\to\mathcal{O}'_X$, donde $\mathcal{O}'_X$ (resp. $\mathcal{O}'_Y$) es la clausura integral de $\mathcal{O}_X$ (resp. $\mathcal{O}_Y$) en $\mathcal{A}$ (resp. $\mathcal{B}$).
Entonces, si $X'$ (resp. $Y'$) es la clausura integral de $X$ (resp. $Y$) respecto de $\mathcal{A}$ (resp. $\mathcal{B}$), se deduce canónicamente de $u$ un morfismo $f'=\operatorname{Spec}(u'):X'\to Y'$ \sref{II.1.5.6-fr} que hace conmutativo el diagrama
\[
\label{II.6.3.5.1-fr}
\xymatrix{
X'\ar[r]^{f'}\ar[d]&Y'\ar[d]\\
X\ar[r]_{f}&Y
}
\tag{6.3.5.1}
\]
\end{env}
% Fuente francesa: PMIHES_1961__8__5_0, printed p. 118.

\begin{env}[6.3.6]
\label{II.6.3.6-fr}
Supongamos que $X$ solo tiene un número \emph{finito} de componentes irreducibles $X_i$ ($1\leq i\leq r$), de puntos genéricos $\xi_i$, y consideremos, en particular, la clausura integral de $X$ respecto de una $\mathcal{R}(X)$-\emph{álgebra} cuasicoherente $\mathcal{A}$ (cuasicoherente tanto en cuanto $\mathcal{O}_X$-álgebra como en cuanto $\mathcal{R}(X)$-álgebra, pues ambas condiciones son equivalentes).
Se sabe \sref[I]{I.7.3.5-fr} que $\mathcal{A}$ es la suma directa de $r$ $\mathcal{O}_X$-álgebras cuasicoherentes $\mathcal{A}_i$, donde el soporte de $\mathcal{A}_i$ está contenido en $X_i$, y el haz inducido por $\mathcal{A}_i$ sobre $X_i$ es un haz simple cuyo tallo $A_i$ es un álgebra sobre $\mathcal{O}_{\xi_i}$.
Es claro entonces \sref{II.6.3.4-fr} que la clausura integral $\mathcal{O}'_X$ de $\mathcal{O}_X$ en $\mathcal{A}$ es la \emph{suma directa} de las clausuras integrales $\mathcal{O}^{(i)}_X$ de $\mathcal{O}_X$ en cada una de las $\mathcal{A}_i$; por tanto, la clausura integral $X'=\operatorname{Spec}(\mathcal{O}'_X)$ de $X$ respecto de $\mathcal{A}$ es el $X$-esquema dado por la \emph{suma} de los $\operatorname{Spec}(\mathcal{O}^{(i)}_X)=X'_i$ ($1\leq i\leq r$).

Supongamos además que la $\mathcal{O}_X$-álgebra $\mathcal{A}$ es \emph{reducida}, o, equivalentemente, que cada una de las álgebras $A_i$ es reducida y puede considerarse, por tanto, como un álgebra sobre el \emph{cuerpo} $\boldsymbol{k}(\xi_i)$ (igual al cuerpo de funciones racionales del subpreesquema reducido de $X$ cuyo espacio subyacente es $X_i$);
entonces \sref{II.1.3.8-fr} cada uno de los $X'_i$ es un $X$-esquema \emph{reducido}, y $X'$ es también la clausura integral de $X_{\mathrm{red}}$.
Supongamos además que cada una de las álgebras $A_i$ es una \emph{suma directa de un número finito de cuerpos $K_{ij}$} ($1\leq j\leq s_i$);
si $\mathcal{K}_{ij}$ es la subálgebra de $\mathcal{A}_i$ correspondiente a $K_{ij}$, es claro que $\mathcal{O}^{(i)}_X$ es la \emph{suma directa} de las clausuras integrales $\mathcal{O}^{(ij)}_X$ de $\mathcal{O}_X$ en cada una de las $\mathcal{K}_{ij}$.
Entonces $X'_i$ es la \emph{suma} de los $X$-esquemas $X'_{ij}=\operatorname{Spec}(\mathcal{O}^{(ij)}_X)$ ($1\leq j\leq s_i$).
Además, bajo estas hipótesis y con estas notaciones, se tiene:
\end{env}
% Fuente francesa: PMIHES_1961__8__5_0, printed p. 119.

\oldpage[II]{119}
\begin{proposition}[6.3.7]
\label{II.6.3.7-fr}
Cada uno de los $X'_{ij}$ es un $X$-preesquema íntegro y normal, y su cuerpo de funciones racionales $R(X'_{ij})$ se identifica canónicamente con la clausura algebraica $K'_{ij}$ de $\boldsymbol{k}(\xi_i)$ en $K_{ij}$.
\end{proposition}

\begin{proof}
En virtud de lo anterior, se puede suponer que $X$ es íntegro, de modo que $r=1$ y $s_1=1$, y que la única álgebra $A_1$ es un cuerpo $K$;
sea $\xi$ el punto genérico de $X$, y sea $f:X'\to X$ el morfismo estructural.
Para todo abierto afín no vacío $U$ de $X$, $f^{-1}(U)$ se identifica con el espectro de la clausura integral $B'_U$ en el cuerpo $K$ del anillo íntegro $B_U=\Gamma(U,\mathcal{O}_X)$ \sref{II.6.3.4-fr};
como el anillo $B'_U$ es íntegro e íntegramente cerrado, lo mismo ocurre con los anillos locales de los puntos de su espectro, de modo que $f^{-1}(U)$ es, por definición, un esquema \emph{íntegro} y \emph{normal} (\sref[0]{0.4.1.4-fr} y \sref[I]{I.5.1.4-fr}).
Además, como $(0)$ es el único ideal primo de $B'_U$ situado sobre el ideal primo $(0)$ de $B_U$ \cite[t.~I, p.~259]{I-13}, $f^{-1}(\xi)$ se reduce a un único punto $\xi'$, y $\boldsymbol{k}(\xi')$ es el cuerpo de fracciones $K'$ de $B'_U$, que no es sino la clausura algebraica de $\boldsymbol{k}(\xi)$ en $K$.
Finalmente, $X'$ es irreducible, pues, cuando $U$ recorre el conjunto de los abiertos afines no vacíos de $X$, los $f^{-1}(U)$ forman una cubierta abierta de $X'$ por abiertos irreducibles;
además, la intersección $f^{-1}(U\cap V)$ de dos cualesquiera de estos abiertos contiene $\xi'$ y, por tanto, no es vacía; se concluye mediante \sref[0]{0.2.1.4-fr}.
\end{proof}
% Fuente francesa: PMIHES_1961__8__5_0, printed p. 119.

\begin{corollary}[6.3.8]
\label{II.6.3.8-fr}
Sea $X$ un preesquema reducido que solo tiene un número finito de componentes irreducibles $X_i$ ($1\leq i\leq r$), y sea $\xi_i$ el punto genérico de $X_i$.
Entonces la clausura integral $X'$ de $X$ respecto de $\mathcal{R}(X)$ es la suma de los $r$ $X$-esquemas $X'_i$, que son íntegros y normales.
Si $f:X'\to X$ es el morfismo estructural, $f^{-1}(\xi_i)$ se reduce al punto genérico $\xi'_i$ de $X'_i$, y se tiene $\boldsymbol{k}(\xi'_i)=\boldsymbol{k}(\xi_i)$; dicho de otro modo, $f$ es birracional.
\end{corollary}

En este caso se dice que $X'$ es la \emph{normalización} del preesquema reducido $X$;
se señala que $f$, por ser birracional y entero, es \emph{sobreyectivo} \sref{II.6.1.10-fr}.
Para que se tenga entonces $X'=X$, es necesario y suficiente que $X$ sea \emph{normal}.
Si $X$ es un preesquema íntegro, se sigue de \sref{II.6.3.8-fr} que su normalización $X'$ es íntegra.
% Fuente francesa: PMIHES_1961__8__5_0, printed pp. 119-120.

\begin{env}[6.3.9]
\label{II.6.3.9-fr}
Sean $X$ e $Y$ preesquemas íntegros, $f:X\to Y$ un morfismo dominante, y sean $L=R(X)$ (resp. $K=R(Y)$) el cuerpo de funciones racionales de $X$ (resp. $Y$);
a $f$ corresponde una inyección canónica $K\to L$, y, si se identifica $K$ (resp. $L$) con el haz simple $\mathcal{R}(Y)$ (resp. $\mathcal{R}(X)$), esta inyección es un $f$-morfismo.
Sea $K_1$ (resp. $L_1$) una extensión de $K$ (resp. $L$), y supongamos dado un morfismo $K_1\to L_1$ tal que el diagrama
\[
\xymatrix{
K_1\ar[r]&L_1\\
K\ar[u]\ar[r]&L\ar[u]
}
\]
sea conmutativo;
si $K_1$ (resp. $L_1$) se considera como haz simple sobre $Y$ (resp. $X$), y, por tanto, como $\mathcal{R}(Y)$-álgebra (resp. $\mathcal{R}(X)$-álgebra), esto implica que $K_1\to L_1$ es un $f$-morfismo.
Dicho esto, si $X'$ (resp. $Y'$) es la clausura integral de $X$ (resp. $Y$) respecto de $L_1$ (resp. $K_1$), entonces $X'$ (resp. $Y'$) es un preesquema íntegro y normal \sref{II.6.3.6-fr} cuyo cuerpo de funciones racionales se identifica canónicamente con la clausura algebraica
\oldpage[II]{120}
 $L'$ (resp. $K'$) de $L$ (resp. $K$) en $L_1$ (resp. $K_1$), y existe un morfismo canónico (necesariamente dominante) $f':X'\to Y'$ que hace conmutativo el diagrama \sref{II.6.3.5.1-fr}.


El caso más importante es aquel en que se toma $L_1=L$, siendo entonces $K_1$ una extensión de $K$ contenida en $L$, y en que se supone $X$ íntegro y \emph{normal}, de modo que $X'=X$.
Lo anterior muestra entonces que, como $X$ es normal, si $Y'$ es la clausura integral de $Y$ respecto de un cuerpo $K_1\subset L=R(X)$, todo morfismo dominante $f:X\to Y$ se factoriza como
\[
f:X\xrightarrow{f'}Y'\to Y
\]
donde $f'$ es dominante;
además, como el monomorfismo $K_1\to L$ está dado, $f'$ es necesariamente único, como se ve reduciéndose al caso en que $X$ e $Y$ son afines.
Se ve así que, dados $Y$, $L$ y un $K$-monomorfismo $K_1\to L$, la clausura integral $Y'$ de $Y$ respecto de $K_1$ es la solución de un \emph{problema universal}.
\end{env}
% Fuente francesa: PMIHES_1961__8__5_0, printed p. 120.

\begin{remark}[6.3.10]
\label{II.6.3.10-fr}
Retomemos las hipótesis de \sref{II.6.3.6-fr}, suponiendo además que cada una de las álgebras $A_i$ es \emph{de rango finito} sobre $\boldsymbol{k}(\xi_i)$ (lo que implica que $A_i$ es suma directa de un número finito de cuerpos);
puede demostrarse entonces, en ciertos casos, que el morfismo estructural $X'\to X$ no solo es entero, sino incluso \emph{finito}.
Limitémonos al caso en que $X$ es \emph{reducido};
como la cuestión es local sobre $X$, se puede suponer además que $X$ es afín, de anillo $C$, y que $C$ solo tiene un número finito de ideales primos minimales $\mathfrak{p}_i$ ($1\leq i\leq r$), tales que los $C_i=C/\mathfrak{p}_i$ sean íntegros;
entonces $X'$ es finito sobre $X$ si la clausura integral de cada $C_i$ en toda extensión de grado finito de su cuerpo de fracciones es un \emph{$C$-módulo de tipo finito} \sref{II.6.3.4-fr}.
Se sabe que esta condición se satisface siempre si $C$ es un \emph{álgebra de tipo finito sobre un cuerpo} \cite[t.~I, p.~267, th.~9]{I-13}, o \emph{sobre $\mathbf{Z}$} \cite[I,p.~93, th.~3]{I-9}, o \emph{sobre un anillo local noetheriano completo} \cite[p.~298]{II-25}.
Se concluye así que $X'\to X$ será un morfismo finito siempre que $X$ sea un esquema \emph{de tipo finito} sobre un cuerpo, o sobre $\mathbf{Z}$, o sobre un anillo local noetheriano completo.
\end{remark}
% Fuente francesa: PMIHES_1961__8__5_0, printed pp. 120-121.

\subsection{Determinante de un endomorfismo de $\mathcal{O}_X$-módulos}
\label{subsection:II.6.4-fr}

\begin{env}[6.4.1]
\label{II.6.4.1-fr}
Sean $A$ un anillo (conmutativo), $E$ un $A$-módulo libre de rango~$n$, y $u$ un endomorfismo de $E$;
recordemos que, para definir el \emph{polinomio característico} de $u$, se considera el endomorfismo $u\otimes1$ del $A[T]$-módulo libre de rango~$n$, $E\otimes_AA[T]$ (donde $T$ es una indeterminada), y se pone
\[
\label{II.6.4.1.1-fr}
P(u,T)=\det(T.I-(u\otimes1))
\tag{6.4.1.1}
\]
(donde $I$ es el automorfismo identidad de $E\otimes_AA[T]$).
Se tiene
\[
\label{II.6.4.1.2-fr}
P(u,T)=T^n-\sigma_1(u)T^{n-1}+\ldots+(-1)^n\sigma_n(u)
\tag{6.4.1.2}
\]
donde $\sigma_i(u)$ es un elemento de $A$, igual a un polinomio homogéneo de grado~$i$ (con coeficientes enteros) respecto de los elementos de la matriz de $u$ relativa a una base cualquiera de $E$;
se dice que los $\sigma_i(u)$ son las \emph{funciones simétricas elementales} de $u$;
en particular, se tiene $\sigma_1(u)=\operatorname{Tr}u$ y $\sigma_n(u)=\det u$.
Recordemos que, en virtud del teorema de Hamilton--Cayley, se tiene
\[
\label{II.6.4.1.3-fr}
P(u,u)=u^n-\sigma_1(u)u^{n-1}+\ldots+(-1)^n\sigma_n(u)=0
\tag{6.4.1.3}
\]
\oldpage[II]{121}
lo que también se escribe
\[
\label{II.6.4.1.4-fr}
(\det u). 1_E=uQ(u)
\tag{6.4.1.4}
\]
(donde $1_E$ es el automorfismo identidad de $E$), con
\[
\label{II.6.4.1.5-fr}
Q(u)=(-1)^{n+1}(u^{n-1}-\sigma_1(u)u^{n-2}+\ldots+(-1)^{n-1}\sigma_{n-1}(u)).
\tag{6.4.1.5}
\]

Sea $\varphi:A\to B$ un homomorfismo de anillos que hace de $B$ un $A$-álgebra;
consideremos el $B$-módulo $E_{(B)}=E\otimes_AB$, que es libre de rango~$n$, y la extensión $u\otimes1$ de $u$ a un endomorfismo de $E_{(B)}$;
es inmediato que $\sigma_i(u\otimes1)=\varphi(\sigma_i(u))$ para todo índice $i$.
\end{env}
% Fuente francesa: PMIHES_1961__8__5_0, printed p. 121.

\begin{env}[6.4.2]
\label{II.6.4.2-fr}
Supongamos ahora que $A$ es un anillo \emph{íntegro}, de cuerpo de fracciones $K$, y sea $E$ un $A$-módulo \emph{de tipo finito} (pero no necesariamente libre).
Sea $n$ el \emph{rango} de $E$, es decir, el rango del $K$-módulo libre $E\otimes_AK$;
a todo endomorfismo $u$ de $E$ corresponde canónicamente el endomorfismo $u\otimes1$ de $E\otimes_AK$;
por abuso de lenguaje, se denomina también \emph{polinomio característico} de $u$, y se denota por $P(u,T)$, el polinomio $P(u\otimes1,T)$, cuyos coeficientes $\sigma_i(u\otimes1)$ (que pertenecen a $K$) se denotan también por $\sigma_i(u)$ y se llaman \emph{funciones simétricas elementales} de $u$;
en particular, $\det u=\det(u\otimes1)$ por definición.
Con estas notaciones, las ecuaciones~\sref{II.6.4.1.3-fr} a \sref{II.6.4.1.5-fr} tienen sentido y siguen siendo válidas, interpretando $u^j$ como el homomorfismo
$E\to E\otimes_AK$ compuesto del endomorfismo $u^j\otimes1=(u\otimes1)^j$ de $E\otimes_AK$ y del homomorfismo canónico $x\to x\otimes1$.

Si $F$ es el submódulo de torsión de $E$, y si $E_0=E/F$, entonces $u(F)\subset F$ y, pasando al cociente, $u$ induce un endomorfismo $u_0$ de $E_0$;
además, $E\otimes_AK$ se identifica con $E_0\otimes_AK$, y $u\otimes1$ con $u_0\otimes1$, de modo que $\sigma_i(u)=\sigma_i(u_0)$ para $1\leq i\leq n$.

Cuando $E$ es \emph{sin torsión}, $E$ se identifica canónicamente con un sub-$A$-módulo de $E\otimes_AK$, y la relación $u\otimes1=0$ equivale a $u=0$.
Si $E$ es un $A$-módulo libre, las dos definiciones de $\sigma_i(u)$ dadas en \sref{II.6.4.1-fr} y \sref{II.6.4.2-fr} coinciden por lo anterior, lo que justifica la notación adoptada.

Se señala también que, si $E$ es un \emph{módulo de torsión}, o, dicho de otro modo, si $E_0=\{0\}$, entonces el álgebra exterior de $E_0$ se reduce a $K$, y el determinante del único endomorfismo $u_0$ de $E_0$ es \emph{igual a $1$}.
\end{env}
% Fuente francesa: PMIHES_1961__8__5_0, printed pp. 121-122.

\begin{proposition}[6.4.3]
\label{II.6.4.3-fr}
Sean $A$ un anillo íntegro, $E$ un $A$-módulo de tipo finito, y $u$ un endomorfismo de $E$;
entonces las funciones simétricas elementales $\sigma_i(u)$ de $u$ (y, en particular, $\det u$) son elementos de $K$ enteros sobre $A$.
\end{proposition}

\begin{proof}
Sea $A'$ la clausura integral de $A$;
como $A'[T]$ es un anillo íntegramente cerrado \cite[p.~99]{II-24}, es la clausura integral de $A[T]$ en su cuerpo de fracciones $K(T)$.
Reemplazando $u$ por $T.I-u\otimes1$, y $A$ por $A[T]$, se ve que se puede reducir a probar que $\det u$ es entero sobre $A$.
Si $n$ es el rango de $E$, entonces $\det u=\det(\bigwedge^n u)$ y $(\bigwedge^n u)\otimes1=\bigwedge^n(u\otimes1)$, de modo que se puede suponer que $n=1$.
Pero entonces la aplicación $u\to\det u$ es un homomorfismo del $A$-módulo $\operatorname{Hom}_A(E,E)$ en $K$;
como $E$ es de tipo finito, $\operatorname{Hom}_A(E,E)$ es isomorfo a un submódulo del $A$-módulo de tipo finito $E^n$ si $E$ admite una familia de $n$ generadores; por tanto, los elementos $\det u$ pertenecen a un sub-$A$-módulo de tipo finito de $K$ y son, en consecuencia, enteros sobre $A$.
\end{proof}

\oldpage[II]{122}
\begin{corollary}[6.4.4]
\label{II.6.4.4-fr}
Bajo las hipótesis de \sref{II.6.4.3-fr}, si se supone además que $A$ es normal, entonces los $\sigma_i(u)$ (y, en particular, $\operatorname{Tr}u$ y $\det u$) pertenecen a $A$.
\end{corollary}
% Fuente francesa: PMIHES_1961__8__5_0, printed p. 122.

\begin{proposition}[6.4.5]
\label{II.6.4.5-fr}
Sean $A$ un anillo íntegro, $E$ un $A$-módulo de tipo finito y de rango~$n$, y $u$ un endomorfismo de $E$ tal que los $\sigma_i(u)$ pertenezcan a $A$.
Para que $u$ sea un automorfismo de $E$, es necesario que $\det u$ sea invertible en $A$;
esta condición es suficiente si $E$ es sin torsión.
\end{proposition}

\begin{proof}
La condición es \emph{suficiente}, pues la hipótesis y las ecuaciones~\sref{II.6.4.1.4-fr} y \sref{II.6.4.1.5-fr} (que son válidas en $E$, y no solo en $E\otimes_AK$, ya que $E$ es sin torsión) prueban que $(\det u)^{-1}Q(u)$ es el inverso de $u$.

La condición es \emph{necesaria}, pues, si $u$ es invertible, se sigue de \sref{II.6.4.3-fr} que $\det(u^{-1})$ pertenece a la clausura integral $A'$ de $A$ en su cuerpo de fracciones $K$, y es evidentemente el inverso de $\det u$ \emph{en $A'$}.
Nuestra afirmación se sigue entonces de:
\end{proof}

\begin{lemma}[6.4.5.1]
\label{II.6.4.5.1-fr}
Sea $A$ un subanillo de un anillo $A'$ tal que $A'$ sea \emph{entero} sobre $A$.
Si un elemento $x\in A$ es invertible en $A'$, también es invertible en $A$.
\end{lemma}

\begin{proof}
En caso contrario, $x$ pertenecería a un ideal maximal $\mathfrak{m}$ de $A$, y se seguiría del primer teorema de Cohen--Seidenberg \cite[t.~I, p.~257, th.~3]{I-13} que existiría un ideal maximal $\mathfrak{m}'$ de $A'$ tal que $\mathfrak{m}=\mathfrak{m}'\cap A$;
se tendría entonces $x\in\mathfrak{m}'$, lo cual es absurdo.
\end{proof}

\begin{corollary}[6.4.6]
\label{II.6.4.6-fr}
Sean $A$ un anillo íntegro e íntegramente cerrado, $E$ un $A$-módulo sin torsión de tipo finito, y $u$ un endomorfismo de $E$.
Para que $u$ sea un automorfismo de $E$, es necesario y suficiente que $\det u$ sea invertible en $A$.
\end{corollary}

\begin{proof}
Esto se sigue de \sref{II.6.4.4-fr} y \sref{II.6.4.5-fr}.
\end{proof}
% Fuente francesa: PMIHES_1961__8__5_0, printed p. 122.

\begin{remark}[6.4.7]
\label{II.6.4.7-fr}
Más adelante necesitaremos una generalización de los resultados anteriores.
Consideremos un anillo noetheriano \emph{reducido} $A$;
sean $\mathfrak{p}_\alpha$ ($1\leq\alpha\leq r$) sus ideales primos minimales, $K_\alpha$ el cuerpo de fracciones del anillo íntegro $A/\mathfrak{p}_\alpha$, y $K$ el anillo total de fracciones de $A$, que es la \emph{suma directa} de los cuerpos $K_\alpha$.
Sea $E$ un $A$-módulo de tipo finito, y \emph{supongamos} que $E\otimes_AK$ es un $K$-módulo \emph{libre} de rango~$n$ (lo cual ya no es aquí consecuencia de las demás hipótesis);
equivalentemente, puede exigirse que todos los $K_\alpha$-espacios vectoriales $E\otimes_AK_\alpha=E_\alpha$ tengan \emph{la misma dimensión~$n$};
si $u$ es entonces un endomorfismo de $E$, se pone de nuevo $P(u,T)=P(u\otimes1,T)$ y $\sigma_j(u)=\sigma_j(u\otimes1)$ y, en particular, $\det u=\det(u\otimes1)$;
los $\sigma_j(u)$ son, por tanto, elementos de $K$.
Es inmediato que $E\otimes_AK$ es la suma directa de los $E_\alpha$, y que estos son estables por $u\otimes1$;
la restricción de $u\otimes1$ a $E_\alpha$ es precisamente la extensión $u_\alpha$ de $u$ a este $K_\alpha$-espacio vectorial;
se concluye que $\sigma_j(u)$ es el elemento de $K$ cuyas componentes en los $K_\alpha$ son los $\sigma_j(u_\alpha)$.
Como la clausura integral de $A$ en $K$ es la \emph{suma directa} de las clausuras integrales de $A$ en los $K_\alpha$, los $\sigma_j(u)$ son \emph{enteros} sobre $A$.

\begin{lemma}[6.4.7.1]
\label{II.6.4.7.1-fr}
La sub-$A$-álgebra de $K$ generada por todos los elementos $\sigma_j(u)$ ($1\leq j\leq n$), cuando $u$ recorre $\operatorname{Hom}_A(E,E)$, es un $A$-módulo de tipo finito.
\end{lemma}
% Fuente francesa: PMIHES_1961__8__5_0, printed pp. 122-123.

\begin{proof}
Basta demostrar que la sub-$A[T]$-álgebra de $K[T]$ generada por los $P(u,T)$ es un $A[T]$-módulo de tipo finito, pues, si los $F_i(T)$ ($1\leq i\leq m$) forman una familia de generadores de este $A[T]$-módulo, los coeficientes de los $F_i(T)$ son enteros sobre $A$ y generan, por tanto, un $A$-álgebra que es un $A$-módulo de tipo finito \cite[t.~I, p.~255, th.~1]{I-13}.
\oldpage[II]{123}
Se puede, pues, reemplazar $A$ por $A[T]$ (que es noetheriano), y $E$ por $E\otimes_AA[T]=E'$, que es tal que $E'\otimes_{A[T]}K[T]=E\otimes_AK[T]$ sea un $K[T]$-módulo libre de rango~$n$.
Retomando las notaciones iniciales, se ve así que basta probar que el $A$-módulo generado por los elementos $\det u$, cuando $u$ recorre $\operatorname{Hom}_A(E,E)$, es de tipo finito;
\emph{a fortiori} (puesto que todo submódulo de un $A$-módulo de tipo finito es de tipo finito), basta probar que, cuando $v$ recorre el conjunto de los endomorfismos de $\bigwedge^nE$, el $A$-módulo generado por los $\det v$ es de tipo finito;
dicho de otro modo, se puede reducir de nuevo al caso en que $n=1$.
Pero entonces la proposición se sigue de que $\operatorname{Hom}_A(E,E)$ es un $A$-módulo de tipo finito y de que $v\to\det v$ es un homomorfismo de este $A$-módulo en $K$.
\end{proof}

Sea $F$ el núcleo del homomorfismo canónico $E\to E\otimes_AK$, y sea $E_0=E/F$;
se ve, como antes, que $E\otimes_AK$ se identifica con $E_0\otimes_AK$, que $u(F)\subset F$ y que, si $u_0$ es el endomorfismo de $E_0$ deducido de $u$ por paso al cociente, entonces $u\otimes1$ se identifica con $u_0\otimes1$, y $\sigma_j(u)=\sigma_j(u_0)$ para todo $j$.
\emph{Si se tiene $F=0$}, entonces las ecuaciones~\sref{II.6.4.1.3-fr} y \sref{II.6.4.1.5-fr} siguen teniendo sentido y siendo válidas cuando $E$ se identifica con un submódulo de $E\otimes_AK$, y los $u^j$ con homomorfismos $E\to E\otimes_AK$;
la Proposición~\sref{II.6.4.5-fr} se extiende, por tanto, también a este caso, con la misma demostración.
\end{remark}
% Fuente francesa: PMIHES_1961__8__5_0, printed p. 123.

\begin{env}[6.4.8]
\label{II.6.4.8-fr}
Sean $(X,\mathcal{A})$ un espacio anillado, $\mathcal{E}$ un $\mathcal{A}$-módulo \emph{localmente libre} (de rango finito), y $u$ un endomorfismo de $\mathcal{E}$.
Por hipótesis, existe una base $\mathfrak{B}$ de la topología de $X$ tal que, para todo $V\in\mathfrak{B}$, $\mathcal{E}|V$ sea isomorfo a algún $\mathcal{A}^n|V$ (donde $n$ puede variar con $V$).
Sea $u$ un endomorfismo de $\mathcal{E}$; para todo $V\in\mathfrak{B}$, $u_V$ es, por tanto, un endomorfismo del $\Gamma(V,\mathcal{A})$-módulo $\Gamma(V,\mathcal{E})$, que es libre por hipótesis;
el determinante $\det u_V$ está, pues, definido y pertenece a $\Gamma(V,\mathcal{A})$.
Además, si $e_1,\ldots,e_n$ forman una base de $\Gamma(V,\mathcal{E})$, sus restricciones a todo abierto $W\subset V$ forman una base de $\Gamma(W,\mathcal{E})$ sobre $\Gamma(W,\mathcal{A})$, de modo que $\det u_W$ es la restricción de $\det u_V$ a $W$.
Existe, por tanto, una única sección de $\mathcal{A}$ sobre $X$, que denotaremos por $\det u$ y se llama el \emph{determinante} de $u$, tal que la restricción de $\det u$ a todo $V\in\mathfrak{B}$ sea $\det u_V$.
Es claro que, para todo $x\in X$, se tiene $(\det u)_x=\det u_x$;
para dos endomorfismos $u$ y $v$ de $\mathcal{E}$, se tiene
\[
\label{II.6.4.8.1-fr}
\det(u\circ v)=(\det u)(\det v)
\tag{6.4.8.1}
\]
así como
\[
\label{II.6.4.8.2-fr}
\det(1_{\mathcal{E}})=1_{\mathcal{A}}
\tag{6.4.8.2}
\]
y, si el rango de $\mathcal{E}$ es constante (lo cual ocurre \sref[0]{0.5.4.1-fr} si $X$ es conexo) e igual a $n$,
\[
\label{II.6.4.8.3-fr}
\det(s.u)=s^n\det u
\tag{6.4.8.3}
\]
para todo $s\in\Gamma(X,\mathcal{A})$ (se señala que $\det(0)=0_{\mathcal{A}}$ si $n\geq1$, pero $\det(0)=1_{\mathcal{A}}$ si $n=0$).
Además, para que $u$ sea un \emph{automorfismo} de $\mathcal{E}$, es necesario y suficiente que $\det u$ sea \emph{invertible} en $\Gamma(X,\mathcal{A})$.

Si el rango de $\mathcal{E}$ es constante, pueden definirse de igual modo las funciones simétricas elementales $\sigma_i(u)$, que son elementos de $\Gamma(X,\mathcal{A})$;
se tienen de nuevo las relaciones de \sref{II.6.4.1.3-fr} y \sref{II.6.4.1.5-fr}.
% Fuente francesa: PMIHES_1961__8__5_0, printed p. 124.

\oldpage[II]{124}
Se ha definido así un homomorfismo $\det:\operatorname{Hom}_{\mathcal{A}}(\mathcal{E},\mathcal{E})
\to\Gamma(X,\mathcal{A})$ de monoides multiplicativos;
si se señala que $\operatorname{Hom}_{\mathcal{A}}(\mathcal{E},\mathcal{E})=
\Gamma(X,\mathcal{H}om_{\mathcal{A}}(\mathcal{E},\mathcal{E}))$ por definición, se ve que puede reemplazarse $X$ en esta definición por un abierto cualquiera $U$ de $X$, lo que define inmediatamente un homomorfismo $\det:\mathcal{H}om_{\mathcal{A}}(\mathcal{E},\mathcal{E})\to\mathcal{A}$ de haces de monoides multiplicativos.
Si $\mathcal{E}$ tiene rango constante, se definen de igual modo homomorfismos $\sigma_i:\mathcal{H}om_{\mathcal{A}}(\mathcal{E},\mathcal{E})
\to\mathcal{A}$ de haces de conjuntos;
para $i=1$, el homomorfismo $\sigma_1=\operatorname{Tr}$ es un homomorfismo de $\mathcal{A}$-módulos.

Sea $(Y,\mathcal{B})$ otro espacio anillado, y sea $f:(X,\mathcal{A})\to(Y,\mathcal{B})$ un morfismo de espacios anillados;
si $\mathcal{F}$ es un $\mathcal{B}$-módulo localmente libre, entonces $f^*(\mathcal{F})$ es un $\mathcal{A}$-módulo localmente libre (del mismo rango que $\mathcal{F}$ si este último es de rango constante) \sref[0]{0.5.4.5-fr}.
Para todo endomorfismo $v$ de $\mathcal{F}$, $f^*(v)$ es un endomorfismo de $f^*(\mathcal{F})$, y se sigue inmediatamente de las definiciones que $\det f^*(v)$ es la sección de $\mathcal{A}=f^*(\mathcal{B})$ sobre $X$ que corresponde canónicamente a $\det v\in\Gamma(Y,\mathcal{B})$.
Puede decirse además que el homomorfismo $f^*(\det):f^*(\mathcal{H}om_{\mathcal{B}}(\mathcal{F},\mathcal{F}))
\to f^*(\mathcal{B})=\mathcal{A}$ es el compuesto
\[
f^*(\mathcal{H}om_{\mathcal{B}}(\mathcal{F},\mathcal{F}))
\xrightarrow{\gamma^\sharp}
\mathcal{H}om_{\mathcal{A}}(f^*(\mathcal{F}),f^*(\mathcal{F}))
\xrightarrow{\det}\mathcal{A}
\]
\sref[0]{0.4.4.6-fr}.
Se tienen resultados análogos para los $\sigma_i$.
\end{env}
% Fuente francesa: PMIHES_1961__8__5_0, printed p. 124.

\begin{env}[6.4.9]
\label{II.6.4.9-fr}
Supongamos ahora que $X$ es un \emph{preesquema localmente íntegro}, de modo que el haz $\mathcal{R}(X)$ de funciones racionales sobre $X$ es un haz de cuerpos localmente simple \sref[I]{I.7.3.4-fr}, y cuasicoherente en cuanto $\mathcal{O}_X$-módulo.
Si $\mathcal{E}$ es un $\mathcal{O}_X$-módulo cuasicoherente \emph{de tipo finito}, entonces $\mathcal{E}'=\mathcal{E}\otimes_{\mathcal{O}_X}\mathcal{R}(X)$ es un $\mathcal{R}(X)$-módulo localmente libre \sref[I]{I.7.3.6-fr};
para todo endomorfismo $u$ de $\mathcal{E}$, $u\otimes1_{\mathcal{R}(X)}$ es entonces un endomorfismo de $\mathcal{E}'$, y $\det(u\otimes1)$ es una sección de $\mathcal{R}(X)$ sobre $X$, que se denomina también el \emph{determinante} de $u$ y se denota igualmente por $\det u$.
Se sigue de \sref{II.6.4.3-fr} que $\det u$ es una sección de la \emph{clausura integral} de $\mathcal{O}_X$ en $\mathcal{R}(X)$ \sref{II.6.3.2-fr};
si, además, $X$ es \emph{normal}, entonces $\det u$ es una \emph{sección de $\mathcal{O}_X$} sobre $X$ y, si se supone también que $\mathcal{E}$ es \emph{sin torsión}, para que $u$ sea un automorfismo de $\mathcal{E}$ es necesario y suficiente que $\det u$ sea \emph{invertible}, en virtud de \sref{II.6.4.6-fr}.
Las ecuaciones~\sref{II.6.4.8.1-fr} a \sref{II.6.4.8.3-fr} siguen siendo válidas;
del homomorfismo $u\to\det u$, aplicado a los módulos de secciones de $\mathcal{H}om_{\mathcal{O}_X}(\mathcal{E},\mathcal{E})$, se deduce un homomorfismo de haces $\det:\mathcal{H}om_{\mathcal{O}_X}(\mathcal{E},\mathcal{E})
\to\mathcal{R}(X)$, que toma sus valores en $\mathcal{O}_X$ si $X$ es normal.
Se tienen definiciones y resultados análogos para las demás funciones simétricas elementales $\sigma_j(u)$, si $\mathcal{E}'$ es de \emph{rango constante};
si, además, $X$ es \emph{normal}, entonces los $\sigma_j(u)$ son \emph{secciones de $\mathcal{O}_X$} sobre $X$.
% Fuente francesa: PMIHES_1961__8__5_0, printed pp. 124-125.

Finalmente, sean $X$ e $Y$ preesquemas íntegros, y sea $f:X\to Y$ un morfismo \emph{dominante}.
Se sabe que existe un homomorfismo canónico $f^*(\mathcal{R}(Y))\to\mathcal{R}(X)$ \sref[I]{I.7.3.8-fr}\footnote{Véase \emph{Errata and Addenda}, lista 1.}, de donde se obtiene, para todo $\mathcal{O}_Y$-módulo cuasicoherente de tipo finito $\mathcal{F}$, un homomorfismo canónico $\theta:f^*(\mathcal{F}\otimes_{\mathcal{O}_Y}\mathcal{R}(Y))
\to f^*(\mathcal{F})\otimes_{\mathcal{O}_X}\mathcal{R}(X)$.

Si $v$ es un endo%
\oldpage[II]{125}
morfismo de $\mathcal{F}$, entonces $f^*(v\otimes1_{\mathcal{R}(Y)})$ es un endomorfismo de $f^*(\mathcal{F}\otimes_{\mathcal{O}_Y}\mathcal{R}(Y))$, y se tiene un diagrama conmutativo
\[
\xymatrix{
f^*(\mathcal{F}\otimes_{\mathcal{O}_Y}\mathcal{R}(Y))
  \ar[r]^{f^*(v\otimes1)}\ar[d]_{\theta}
& f^*(\mathcal{F}\otimes_{\mathcal{O}_Y}\mathcal{R}(Y))\ar[d]^{\theta}\\
f^*(\mathcal{F})\otimes_{\mathcal{O}_X}\mathcal{R}(X)
  \ar[r]_{f^*(v)\otimes1}
& f^*(\mathcal{F})\otimes_{\mathcal{O}_X}\mathcal{R}(X)
}
\]

Se concluye fácilmente que $\det f^*(v)$ es la imagen canónica, por el homomorfismo $f^*(\mathcal{R}(Y))\to\mathcal{R}(X)$, de la sección $\det v$ de $\mathcal{R}(Y)$;
en efecto, se puede reducir inmediatamente al caso en que $X=\operatorname{Spec}(A)$ e $Y=\operatorname{Spec}(B)$ son afines, siendo $A$ y $B$ anillos íntegros de cuerpos de fracciones respectivos $K$ y $L$, y siendo el homomorfismo $B\to A$ inyectivo y extendiéndose, por tanto, a un monomorfismo $L\to K$;
si $\mathcal{F}=\widetilde{M}$, donde $M$ es un $B$-módulo de tipo finito, entonces el rango sobre $L$ de $M\otimes_B L$ es \emph{igual} al de $(M\otimes_B A)\otimes_A K$ sobre $K$, y $\det((u\otimes1)\otimes1)$ es la imagen en $K$ de $\det(u\otimes1)$ para todo endomorfismo $u$ de $M$, de donde se sigue la conclusión.
\end{env}
% Fuente francesa: PMIHES_1961__8__5_0, printed p. 125.

\begin{env}[6.4.10]
\label{II.6.4.10-fr}
Supongamos, por último, que $X$ es un \emph{preesquema localmente noetheriano reducido}, de modo que el haz $\mathcal{R}(X)$ de funciones racionales sobre $X$ es también un $\mathcal{O}_X$-módulo cuasicoherente \sref[I]{I.7.3.4-fr};
sea $\mathcal{E}$ un $\mathcal{O}_X$-módulo \emph{coherente} tal que $\mathcal{E}'=\mathcal{E}\otimes_{\mathcal{O}_X}\mathcal{R}(X)$ sea \emph{localmente libre} (de rango finito).
En virtud de \sref{II.6.4.7-fr}, si $\mathcal{E}'$ es de rango constante, pueden definirse, para todo endomorfismo $u$ de $\mathcal{E}$, los $\sigma_j(u)$, que son secciones de $\mathcal{R}(X)$ sobre $X$.
Si no se supone que $\mathcal{E}'$ sea de rango constante, todavía puede definirse el homomorfismo $\det:\mathcal{H}om_{\mathcal{O}_X}(\mathcal{E},\mathcal{E})
\to\mathcal{R}(X)$.
\end{env}
% Fuente francesa: PMIHES_1961__8__5_0, printed pp. 125-126.

\subsection{Norma de un haz invertible}
\label{subsection:II.6.5-fr}

\begin{env}[6.5.1]
\label{II.6.5.1-fr}
Sean $(X,\mathcal{A})$ un espacio anillado y $\mathcal{B}$ una $\mathcal{A}$-álgebra (conmutativa).
El $\mathcal{A}$-módulo $\mathcal{B}$ se identifica canónicamente con un sub-$\mathcal{A}$-módulo de $\mathcal{H}om_{\mathcal{A}}(\mathcal{B},\mathcal{B})$, identificándose una sección $f$ de $\mathcal{B}$ sobre un abierto $U$ de $X$ con la multiplicación por esta sección.
Si $(X,\mathcal{A})$ y $\mathcal{B}$ satisfacen una de las condiciones enumeradas en \sref{II.6.4.8-fr}, \sref{II.6.4.9-fr} o \sref{II.6.4.10-fr}, pueden definirse $\det(f)$ (y, en ciertos casos, los $\sigma_j(f)$) como secciones de $\mathcal{A}$ o de $\mathcal{R}(X)$ sobre $U$, que llamaremos la \emph{norma} de $f$ (resp. las \emph{funciones simétricas elementales} de $f$), y denotaremos por $N_{\mathcal{B}/\mathcal{A}}(f)$.
Supondremos satisfecha \emph{una} de las condiciones siguientes:
\begin{enumerate}
  \item[{\rm(I)}] \emph{$\mathcal{B}$ es un $\mathcal{A}$-módulo localmente libre (de rango finito).}
  \item[{\rm(II)}] \emph{$(X,\mathcal{A})$ es un preesquema localmente noetheriano reducido, $\mathcal{B}$ es un $\mathcal{A}$-módulo coherente tal que $\mathcal{B}\otimes_{\mathcal{A}}\mathcal{R}(X)$ es un $\mathcal{R}(X)$-módulo localmente libre y, para toda sección $f\in\Gamma(U,\mathcal{B})$ sobre un abierto $U\subset X$, la norma $N_{\mathcal{B}/\mathcal{A}}(f)$ es una sección de $\mathcal{A}$ sobre $U$.}
\end{enumerate}

\oldpage[II]{126}
Se señala que la última condición se satisface automáticamente cuando el preesquema localmente noetheriano $X$ es \emph{normal} \sref{II.6.4.9-fr}.

La hipótesis de que $\mathcal{B}\otimes_{\mathcal{A}}\mathcal{R}(X)$ sea localmente libre puede expresarse también de la manera siguiente: designemos por $X_\alpha$ los subpreesquemas cerrados reducidos de $X$ cuyos espacios subyacentes son las componentes irreducibles de $X$ \sref[I]{I.5.2.1-fr}, que son, por tanto, preesquemas íntegros localmente noetherianos.
Todo $x\in X$ pertenece a un número finito de los subespacios $X_\alpha$;
por otra parte, $\mathcal{B}\otimes_{\mathcal{A}}\mathcal{R}(X_\alpha)$ es un $\mathcal{R}(X_\alpha)$-módulo localmente libre de rango constante~$k_\alpha$ \sref[I]{I.7.3.6-fr};
decir que $\mathcal{B}\otimes_{\mathcal{A}}\mathcal{R}(X)$ es un $\mathcal{R}(X)$-módulo localmente libre significa que, para todo $x\in X$, \emph{son iguales los rangos $k_\alpha$ correspondientes a los índices tales que $x\in X_\alpha$}.
La cuestión es local, y se puede reducir al caso $X=\operatorname{Spec}(C)$, donde $C$ es un anillo noetheriano reducido y $\mathcal{B}=\widetilde{D}$, donde $D$ es un $C$-álgebra que es un $C$-módulo de tipo finito;
si $\mathfrak{p}_i$ ($1\leq i\leq m$) son los ideales primos minimales de $C$, el anillo total de fracciones $L$ de $C$ es la suma directa de los cuerpos de fracciones $K_i$ de los anillos íntegros $C/\mathfrak{p}_i$, y $D\otimes_C L$ es la suma directa de los $D\otimes_C K_i$, de donde se sigue la conclusión.
% Fuente francesa: PMIHES_1961__8__5_0, printed p. 126.

Es claro que, bajo la hipótesis~(I) o (II), se ha definido así un \emph{homomorfismo de haces de monoides multiplicativos} $N_{\mathcal{B}/\mathcal{A}}:\mathcal{B}\to\mathcal{A}$, denotado también por $N$ si no puede haber confusión, y llamado homomorfismo \emph{norma}.
Para dos secciones $f$ y $g$ de $\mathcal{B}$ sobre un mismo abierto $U$, se tiene, pues,
\[
\label{II.6.5.1.1-fr}
N_{\mathcal{B}/\mathcal{A}}(fg)
=N_{\mathcal{B}/\mathcal{A}}(f)N_{\mathcal{B}/\mathcal{A}}(g)
\tag{6.5.1.1}
\]
para las secciones correspondientes de $\mathcal{A}$ sobre $U$;
\[
\label{II.6.5.1.2-fr}
N_{\mathcal{B}/\mathcal{A}}(1_{\mathcal{B}})=1_{\mathcal{A}}~;
\tag{6.5.1.2}
\]
por último, para toda sección $s$ de $\mathcal{A}$ sobre $U$,
\[
\label{II.6.5.1.3-fr}
N_{\mathcal{B}/\mathcal{A}}(s.1_{\mathcal{B}})=s^n
\tag{6.5.1.3}
\]
si el rango de $\mathcal{B}$ es constante e igual a $n$ (para $s=0_{\mathcal{A}}$, esta fórmula da $N(0_{\mathcal{B}})=0_{\mathcal{A}}$ si $n\geq1$, y $N(0_{\mathcal{B}})=N(1_{\mathcal{B}})=1_{\mathcal{A}}$ si $n=0$).

Bajo la hipótesis~(I), para que $f\in\Gamma(U,\mathcal{B})$ sea invertible es necesario y suficiente que $N(f)\in\Gamma(U,\mathcal{A})$ sea invertible.
Bajo la hipótesis~(II), esta condición es necesaria;
es suficiente (por \sref{II.6.4.7-fr}) si se supone que $\mathcal{B}\to\mathcal{B}\otimes_{\mathcal{A}}\mathcal{R}(X)$ es \emph{inyectivo} y que se satisface la hipótesis más restrictiva siguiente:
\begin{enumerate}
  \item[{\rm(II~\emph{bis})}] \emph{$(X,\mathcal{A})$ es un preesquema localmente noetheriano reducido, $\mathcal{B}$ es un $\mathcal{A}$-módulo coherente tal que $\mathcal{B}\otimes_{\mathcal{A}}\mathcal{R}(X)$ es un $\mathcal{R}(X)$-módulo localmente libre y, para toda sección $f\in\Gamma(U,\mathcal{B})$ sobre un abierto tal que $\mathcal{B}\otimes_{\mathcal{A}}\mathcal{R}(X)|U$ sea de rango constante~$n$ sobre $\mathcal{R}(X)|U$, los $\sigma_j(f)$ ($1\leq j\leq n$) son secciones de $\mathcal{A}$ sobre $U$.}
\end{enumerate}

(Se señala de nuevo que esta condición se satisface si $X$ es \emph{normal}.)
\end{env}
% Fuente francesa: PMIHES_1961__8__5_0, printed pp. 126-127.

\begin{env}[6.5.2]
\label{II.6.5.2-fr}
Supongamos satisfecha una de las hipótesis (I) o (II) de \sref{II.6.5.1-fr}, y sea $\mathcal{L}'$ un $\mathcal{B}$-\emph{módulo invertible}.
Le asociaremos canónicamente (salvo isomorfismo único) un $\mathcal{A}$-\emph{módulo invertible} $\mathcal{L}$ del modo siguiente.
Denotemos por $\mathcal{A}^*$ (resp. $\mathcal{B}^*$) el subhaz de $\mathcal{A}$ (resp. de $\mathcal{B}$) tal que $\Gamma(U,\mathcal{A}^*)$ (resp. $\Gamma(U,\mathcal{B}^*)$) sea el conjunto de los 
\oldpage[II]{127}
elementos invertibles de $\Gamma(U,\mathcal{A})$ (resp. de $\Gamma(U,\mathcal{B})$) para todo abierto $U\subset X$;

son haces de grupos multiplicativos, y la restricción de $N_{\mathcal{B}/\mathcal{A}}$ a $\mathcal{B}^*$ es un \emph{homomorfismo} $\mathcal{B}^*\to\mathcal{A}^*$ de haces de grupos \sref{II.6.5.1-fr}.
Sea $\mathfrak{L}$ el conjunto de los pares $(U_\lambda,\eta_\lambda)$ que tienen la propiedad siguiente: $U_\lambda$ es un abierto de $X$, y $\eta_\lambda$ es un isomorfismo $\mathcal{L}'|U_\lambda\xrightarrow{\sim}\mathcal{B}|U_\lambda$ de $(\mathcal{B}|U_\lambda)$-módulos.
Por hipótesis, los $U_\lambda$ forman una cubierta de $X$;
para dos índices cualesquiera $\lambda$ y $\mu$, se pone $\omega_{\lambda\mu}=(\eta_\lambda|U_\lambda\cap U_\mu)
\circ(\eta_\mu|U_\lambda\cap U_\mu)^{-1}$, que es un automorfismo de $\mathcal{B}|U_\lambda\cap U_\mu$ identificado canónicamente con una sección de $\mathcal{B}^*$ sobre $U_\lambda\cap U_\mu$, y $(\omega_{\lambda\mu})$ es un $1$-cociclo para la cubierta $\mathfrak{U}=(U_\lambda)$ con valores en $\mathcal{B}^*$ \sref[0]{0.5.4.7-fr}.
El hecho de que $N_{\mathcal{B}/\mathcal{A}}:\mathcal{B}^*\to\mathcal{A}^*$ sea un homomorfismo implica que $(N_{\mathcal{B}/\mathcal{A}}\omega_{\lambda\mu})$ es un $1$-cociclo de $\mathfrak{U}$ con valores en $\mathcal{A}^*$, al que corresponde (salvo isomorfismo único) un $\mathcal{A}$-módulo invertible $\mathcal{L}$ \sref[0]{0.5.4.7-fr}, que denotaremos por $N_{\mathcal{B}/\mathcal{A}}(\mathcal{L}')$ y se llama la \emph{norma} del $\mathcal{B}$-módulo invertible $\mathcal{L}'$.

Sea $\mathfrak{M}$ un subconjunto de $\mathfrak{L}$ tal que los $U_\lambda$ correspondientes formen todavía una cubierta de $X$, y sea $\mathfrak{B}$ la cubierta obtenida;
la restricción del cociclo $(\omega_{\lambda\mu})$ a $\mathfrak{B}$ define un $1$-cociclo $(N_{\mathcal{B}/\mathcal{A}}\omega_{\lambda\mu})$ de $\mathfrak{B}$ con valores en $\mathcal{A}^*$, restricción del $1$-cociclo $(N_{\mathcal{B}/\mathcal{A}}\omega_{\lambda\mu})$ de $\mathfrak{U}$;
es claro que existe un isomorfismo canónico del $\mathcal{A}$-módulo invertible definido por este $1$-cociclo de $\mathfrak{B}$ en $N_{\mathcal{B}/\mathcal{A}}(\mathcal{L}')$, lo que permite definir este $\mathcal{A}$-módulo invertible mediante una subcubierta cualquiera de $\mathfrak{U}$.
% Fuente francesa: PMIHES_1961__8__5_0, printed p. 127.

Esta posibilidad muestra inmediatamente que, si $\mathcal{L}'_1$ y $\mathcal{L}'_2$ son dos $\mathcal{B}$-módulos invertibles, entonces, en virtud de \sref{II.6.5.1.1-fr} y \sref{II.6.5.1.2-fr},
\[
\label{II.6.5.2.1-fr}
N(\mathcal{L}'_1\otimes_{\mathcal{B}}\mathcal{L}'_2)
=N(\mathcal{L}'_1)\otimes_{\mathcal{A}}N(\mathcal{L}'_2)
\tag{6.5.2.1}
\]
y
\[
\label{II.6.5.2.2-fr}
N_{\mathcal{B}/\mathcal{A}}(\mathcal{B})=\mathcal{A}
\tag{6.5.2.2}
\]
así como
\[
\label{II.6.5.2.3-fr}
N(\mathcal{L}'^{-1})=(N(\mathcal{L}'))^{-1}
\tag{6.5.2.3}
\]
salvo isomorfismos canónicos.
Además, se sigue de \sref{II.6.5.1.3-fr} que, si $\mathcal{L}$ es un $\mathcal{A}$-módulo invertible y si $\mathcal{B}$ es de rango constante~$n$ sobre $\mathcal{A}$ en el caso~(I) (resp. $\mathcal{B}\otimes_{\mathcal{A}}\mathcal{R}(X)$ es de rango constante~$n$ sobre $\mathcal{R}(X)$ en el caso~(II)), entonces, salvo isomorfismo canónico,
\[
\label{II.6.5.2.4-fr}
N_{\mathcal{B}/\mathcal{A}}(\mathcal{L}\otimes_{\mathcal{A}}\mathcal{B})
=\mathcal{L}^{\otimes n}.
\tag{6.5.2.4}
\]
\end{env}
% Fuente francesa: PMIHES_1961__8__5_0, printed pp. 127-128.

\begin{env}[6.5.3]
\label{II.6.5.3-fr}
Mostremos ahora que $N_{\mathcal{B}/\mathcal{A}}(\mathcal{L}')$ es un \emph{funtor} covariante en la categoría de los $\mathcal{B}$-módulos invertibles.
Sea $h':\mathcal{L}'_1\to\mathcal{L}'_2$ un homomorfismo de $\mathcal{B}$-módulos invertibles, y sea $\mathfrak{B}=(U_\lambda)$ una cubierta abierta de $X$ tal que, para todo $\lambda$, existan $(\mathcal{B}|U_\lambda)$-isomorfismos $\eta^{(1)}_\lambda:\mathcal{L}'_1|U_\lambda\xrightarrow{\sim}\mathcal{B}|U_\lambda$ y $\eta^{(2)}_\lambda:\mathcal{L}'_2|U_\lambda\xrightarrow{\sim}\mathcal{B}|U_\lambda$;
existe entonces, para cada $\lambda$, un endomorfismo $h'_\lambda$ de $\mathcal{B}|U_\lambda$ tal que $h'_\lambda\circ\eta^{(1)}_\lambda
=\eta^{(2)}_\lambda\circ(h'|U_\lambda)$, y puede identificarse evidentemente $h'_\lambda$ con una sección de $\mathcal{B}$ sobre $U_\lambda$ \sref[0]{0.5.1.1-fr}.
Para todo par $(\lambda,\mu)$ de índices, las restricciones a $U_\lambda\cap U_\mu$ de $(\eta^{(2)}_\lambda)^{-1}\circ h'_\lambda\circ\eta^{(1)}_\lambda$ y $(\eta^{(2)}_\mu)^{-1}\circ h'_\mu\circ\eta^{(1)}_\mu$ coinciden.
Se obtiene así, para los $1$-cociclos $(\omega^{(1)}_{\lambda\mu})$ y $(\omega^{(2)}_{\lambda\mu})$ con valores en $\mathcal{B}^*$ correspondientes a $\mathcal{L}'_1$ y $\mathcal{L}'_2$, la relación
\[
\omega^{(2)}_{\lambda\mu}h'_\mu
=h'_\lambda\omega^{(1)}_{\lambda\mu}.
\]

\oldpage[II]{128}
Si se pone $h_\lambda=N(h'_\lambda)$, se tienen las relaciones análogas
\[
N(\omega^{(2)}_{\lambda\mu})h_\mu
=h_\lambda N(\omega^{(1)}_{\lambda\mu})
\]
y, por tanto, los $h_\lambda$ definen un homomorfismo $N(\mathcal{L}'_1)\to N(\mathcal{L}'_2)$, que denotaremos por $N_{\mathcal{B}/\mathcal{A}}(h')$, o $N(h')$.
Bajo la hipótesis~(I), para que $h'$ sea un \emph{isomorfismo}, es necesario y suficiente (pues se trata de una cuestión local) que $N_{\mathcal{B}/\mathcal{A}}(h')$ sea un isomorfismo.
Bajo la hipótesis~(II), esta condición sigue siendo necesaria;
es suficiente si se satisface la hipótesis~(II~\emph{bis}) y $\mathcal{B}\to\mathcal{B}\otimes_{\mathcal{A}}\mathcal{R}(X)$ es inyectivo.

Tomemos, en particular, $\mathcal{L}'_1=\mathcal{B}$;
los homomorfismos $\mathcal{B}\to\mathcal{L}'$ se identifican entonces \sref[0]{0.5.1.1-fr} con las secciones de $\mathcal{L}'$ sobre $X$, de donde resulta una aplicación canónica
\[
N_{\mathcal{B}/\mathcal{A}}:\Gamma(X,\mathcal{L}')
\to\Gamma(X,N_{\mathcal{B}/\mathcal{A}}(\mathcal{L}')).
\]

Se sigue de nuevo de \sref{II.6.5.1.1-fr} que, si $f'_1\in\Gamma(X,\mathcal{L}'_1)$ y $f'_2\in\Gamma(X,\mathcal{L}'_2)$, entonces
\[
\label{II.6.5.3.1-fr}
N(f'_1\otimes f'_2)=N(f'_1)\otimes N(f'_2).
\tag{6.5.3.1}
\]

Para todo $\mathcal{A}$-módulo invertible $\mathcal{L}$ y toda sección $f\in\Gamma(X,\mathcal{L})$, se tiene, teniendo en cuenta \sref{II.6.5.2.4-fr},
\[
\label{II.6.5.3.2-fr}
N_{\mathcal{B}/\mathcal{A}}(f\otimes1_{\mathcal{B}})=f^{\otimes n}
\tag{6.5.3.2}
\]
siempre que $\mathcal{B}$ sea de rango constante~$n$ bajo la hipótesis~(I) (resp. siempre que $\mathcal{B}\otimes_{\mathcal{A}}\mathcal{R}(X)$ sea de rango constante~$n$ bajo la hipótesis~(II)).
Finalmente, para que el homomorfismo $\mathcal{B}\to\mathcal{L}'$ correspondiente a una sección $f'$ de $\mathcal{L}'$ sobre $X$ sea un isomorfismo, es necesario y suficiente que $f'_x$ sea una base de $\mathcal{L}'_x$ para todo $x\in X$;
bajo la hipótesis~(I), esta condición equivale, pues, a decir que $(N(f'))_x$ es una base de $(N(\mathcal{L}'))_x$ para todo $x$;
bajo la hipótesis~(II), esta condición sigue siendo necesaria, y es suficiente cuando $\mathcal{B}$ satisface la hipótesis~(II~\emph{bis}) y $\mathcal{B}\to\mathcal{B}\otimes_{\mathcal{A}}\mathcal{R}(X)$ es inyectivo.
\end{env}
% Fuente francesa: PMIHES_1961__8__5_0, printed pp. 128-129.

\begin{env}[6.5.4]
\label{II.6.5.4-fr}
Sean $(X,\mathcal{A})$ y $(X',\mathcal{A}')$ dos espacios anillados, \[
f:X'\to X
\] un morfismo, $\mathcal{B}$ una $\mathcal{A}$-álgebra y $\mathcal{B}'=f^*(\mathcal{B})$ la $\mathcal{A}'$-álgebra imagen inversa.
Se supone satisfecha una de las hipótesis siguientes:
\begin{enumerate}
  \item[$1^\circ$] $\mathcal{B}$ satisface la hipótesis~(I) de \sref{II.6.5.1-fr}.
  \item[$2^\circ$] $(X,\mathcal{A})$ y $\mathcal{B}$ satisfacen la hipótesis~(II) de \sref{II.6.5.1-fr}, $(X',\mathcal{A}')$ es un preesquema localmente noetheriano reducido y, si se denotan por $X_\alpha$ y $X'_\beta$ los subpreesquemas cerrados reducidos respectivos de $X$ y $X'$ cuyos espacios subyacentes son las componentes irreducibles de estos espacios, la restricción de $f$ a cada $X'_\beta$ es un morfismo \emph{dominante} de $X'_\beta$ sobre un $X_\alpha$.
\end{enumerate}

En estas condiciones, $\mathcal{B}'$ satisface, respectivamente, la hipótesis~(I) o la hipótesis~(II) de \sref{II.6.5.1-fr}.
La primera afirmación es inmediata;
para establecer la segunda, basta ver que, para todo $x'\in X'$, los rangos de los $\mathcal{B}'\otimes_{\mathcal{O}_{X'}}\mathcal{R}(X'_\beta)$ son \emph{iguales} para los $\beta$ tales que $x'\in X'_\beta$.
Ahora bien, si la restricción de $f$ a $X'_\beta$ es un morfismo dominante sobre $X_\alpha$, el rango de $\mathcal{B}'\otimes_{\mathcal{O}_{X'}}\mathcal{R}(X'_\beta)$ es igual al de $\mathcal{B}\otimes_{\mathcal{O}_X}\mathcal{R}(X_\alpha)$ (como se ve de inmediato reduciéndose al caso afín, como en \sref{II.6.4.9-fr}), de donde resulta nuestra afirmación en virtud de la hipótesis~(II) y de \sref{II.6.5.1-fr}.

\oldpage[II]{129}
Sentado esto, se sigue de \sref{II.6.4.8-fr}, \sref{II.6.4.9-fr} y \sref{II.6.4.10-fr} que, si $s$ es una sección de $\mathcal{B}$ sobre un abierto $U\subset X$ y $s'$ la sección correspondiente de $\mathcal{B}'$ sobre $f^{-1}(U)$, entonces $N_{\mathcal{B}'/\mathcal{A}'}(s')$ es la sección de $\mathcal{A}'$ sobre $f^{-1}(U)$ que corresponde a la sección $N_{\mathcal{B}/\mathcal{A}}(s)$ de $\mathcal{A}$ sobre $U$.

Si $\mathcal{M}$ es un $\mathcal{B}$-módulo invertible, se deduce de lo anterior que, si $\mathcal{M}'=f^*(\mathcal{M})$ (que es un $\mathcal{B}'$-módulo invertible), se tiene, salvo isomorfismo canónico,
$N_{\mathcal{B}'/\mathcal{A}'}(\mathcal{M}')
=f^*(N_{\mathcal{B}/\mathcal{A}}(\mathcal{M}))$
\end{env}

% Fuente francesa: PMIHES_1961__8__5_0, printed p. 129.

\begin{env}[6.5.5]
\label{II.6.5.5-fr}
Supongamos en adelante que $(X,\mathcal{A})$ es un \emph{preesquema}.
El dato de una $\mathcal{A}$-álgebra cuasicoherente $\mathcal{B}$, que es un \emph{$\mathcal{A}$-módulo de tipo finito}, equivale entonces, como se sabe, al de un morfismo \emph{finito} $g:X'\to X$ tal que $g_*(\mathcal{O}_{X'})=\mathcal{B}$, definido salvo $X$-isomorfismo (\sref{II.6.1.2-fr} y \sref{II.1.3.1-fr}).
Además, darse un $\mathcal{O}_{X'}$-módulo cuasicoherente $\mathcal{F}'$ equivale a darse un $\mathcal{B}$-módulo cuasicoherente $\mathcal{F}$ tal que $g_*(\mathcal{F}')=\mathcal{F}$ \sref{II.1.4.3-fr}, y para que $\mathcal{F}'$ sea invertible es necesario y suficiente que $\mathcal{F}$ lo sea \sref{II.6.1.12-fr}.
Para traducir los resultados anteriores en términos del morfismo finito $g$, habrá que suponer, pues, o bien que $g_*(\mathcal{O}_{X'})$ es un $\mathcal{O}_X$-módulo \emph{localmente libre} (de tipo finito), o bien que $(X,\mathcal{O}_X)$ y $g_*(\mathcal{O}_{X'})$ satisfacen la hipótesis~(II).
Para todo $\mathcal{O}_{X'}$-módulo invertible $\mathcal{L}'$, se pondrá entonces
\[
\label{II.6.5.5.1-fr}
N_{X'/X}(\mathcal{L}')
=N_{g_*(\mathcal{O}_{X'})/\mathcal{O}_X}(g_*(\mathcal{L}'))
\tag{6.5.5.1}
\]
y se dirá que es la \emph{norma} (relativa a $g$) de $\mathcal{L}'$.
Del mismo modo, si $h':\mathcal{L}'_1\to\mathcal{L}'_2$ es un homomorfismo de $\mathcal{O}_{X'}$-módulos invertibles, se pondrá
\[
\label{II.6.5.5.2-fr}
N_{X'/X}(h')
=N_{g_*(\mathcal{O}_{X'})/\mathcal{O}_X}(g_*(h')):
N_{X'/X}(\mathcal{L}'_1)\to N_{X'/X}(\mathcal{L}'_2).
\tag{6.5.5.2}
\]

En particular, para $\mathcal{L}'_1=\mathcal{O}_{X'}$, se obtiene así una aplicación canónica
\[
\label{II.6.5.5.3-fr}
N_{X'/X}:\Gamma(X',\mathcal{L}')
\to\Gamma(X,N_{X'/X}(\mathcal{L}')).
\tag{6.5.5.3}
\]

Dejaremos al lector la mayor parte de estas traducciones y nos limitaremos a explicitar las siguientes:
\end{env}

% Fuente francesa: PMIHES_1961__8__5_0, printed pp. 129-130.

\begin{proposition}[6.5.6]
\label{II.6.5.6-fr}
Sea $g:X'\to X$ un morfismo finito, y supongamos, o bien que $g_*(\mathcal{O}_{X'})$ es un $\mathcal{O}_X$-módulo localmente libre, o bien que $(X,\mathcal{O}_X)$ y $g_*(\mathcal{O}_{X'})$ satisfacen \textup{(II~bis)} (lo que ocurrirá, en particular, si $X$ es localmente noetheriano y normal).
Para que un homomorfismo $h':\mathcal{L}'_1\to\mathcal{L}'_2$ de $\mathcal{O}_{X'}$-módulos invertibles sea un isomorfismo, es necesario y suficiente, en la primera hipótesis, que $N_{X'/X}(h')$ sea un isomorfismo;
en la segunda hipótesis, esta condición sigue siendo necesaria, y es suficiente si el homomorfismo
$g_*(\mathcal{O}_{X'})\to
g_*(\mathcal{O}_{X'})\otimes_{\mathcal{O}_X}\mathcal{R}(X)$
es inyectivo.
\end{proposition}

\begin{proof}
Se señala que aquí se utiliza el hecho de que, para que $g_*(h')$ sea un isomorfismo, es necesario y suficiente que $h'$ lo sea \sref{II.1.4.2-fr}.
\end{proof}

\begin{corollary}[6.5.7]
\label{II.6.5.7-fr}
Sea $g:X'\to X$ un morfismo finito, y supongamos, o bien que $g_*(\mathcal{O}_{X'})$ es un $\mathcal{O}_X$-módulo localmente libre, o bien que $(X,\mathcal{O}_X)$ y $g_*(\mathcal{O}_{X'})$ satisfacen \textup{(II~bis)} y que
$g_*(\mathcal{O}_{X'})\to
g_*(\mathcal{O}_{X'})\otimes_{\mathcal{O}_X}\mathcal{R}(X)$
es inyectivo.
Sea $\mathcal{L}'$ un $\mathcal{O}_{X'}$-módulo invertible, $f'$ una sección de $\mathcal{L}'$ sobre $X'$ y $f=N_{X'/X}(f')$ la sección correspondiente de $\mathcal{L}=N_{X'/X}(\mathcal{L}')$ sobre $X$ \sref{II.6.5.5.3-fr}.
\oldpage[II]{130}
Entonces
$g(X'\setmin X'_{f'})=X\setmin X_f$
y $X_f$ es el mayor abierto $U$ de $X$ tal que $g^{-1}(U)\subset X'_{f'}$.
\end{corollary}

\begin{proof}
En efecto, $g(X'\setmin X'_{f'})$ es cerrado en $X$ \sref{II.6.1.10-fr};
basta, pues, probar la última afirmación.
Ahora bien, la relación $U\subset X_f$ equivale a que el homomorfismo $\mathcal{O}_X|U\to\mathcal{L}|U$ definido por $f|U$ sea un isomorfismo.
En virtud de \sref{II.6.5.6-fr}, esto equivale a decir que el homomorfismo
$\mathcal{O}_{X'}|g^{-1}(U)\to\mathcal{L}'|g^{-1}(U)$
definido por $f'|g^{-1}(U)$ sea un isomorfismo, es decir, a la relación $g^{-1}(U)\subset X'_{f'}$.
\end{proof}

% Fuente francesa: PMIHES_1961__8__5_0, printed p. 130.

\begin{proposition}[6.5.8]
\label{II.6.5.8-fr}
Sea $g:X'\to X$ un morfismo finito y $f:Y\to X$ un morfismo;
sean $Y'=X'_{(Y)}$, $g'=g_{(Y)}$ y $f'=f_{(X')}$, de suerte que se tenga el diagrama conmutativo
\[
\xymatrix{
X'\ar[d]_g & Y'\ar[l]_{f'}\ar[d]^{g'}\\
X & Y\ar[l]^f
}
\]

Se supone, o bien que $g_*(\mathcal{O}_{X'})$ es localmente libre, o bien que $(X,\mathcal{O}_X)$ y $g_*(\mathcal{O}_{X'})$ satisfacen \textup{(II)}, que $Y$ es un preesquema localmente noetheriano reducido y que la restricción de $f$ a toda componente irreducible de $Y$ es un morfismo dominante sobre una componente irreducible de $X$.
Entonces, para todo $\mathcal{O}_{X'}$-módulo invertible $\mathcal{L}'$, se tiene, salvo isomorfismo canónico,
\[
N_{Y'/Y}({f'}^*(\mathcal{L}'))=f^*(N_{X'/X}(\mathcal{L}'))
\]
\end{proposition}

\begin{proof}
Señalemos que se tiene
$f^*(g_*(\mathcal{L}'))=g'_*({f'}^*(\mathcal{L}'))$
en virtud de \sref{II.1.5.2-fr}, y, en particular,
$g'_*(\mathcal{O}_{Y'})=f^*(g_*(\mathcal{O}_{X'}))$
si $g_*(\mathcal{O}_{X'})$ es localmente libre, ocurre lo mismo con $g'_*(\mathcal{O}_{Y'})$.
La conclusión resulta entonces de las definiciones y de \sref{II.6.5.4-fr}.
\end{proof}

\begin{remark}[6.5.9]
\label{II.6.5.9-fr}
Generalizaremos ulteriormente la noción de norma desarrollada anteriormente, poniéndola en relación con la noción de imagen directa de un divisor.
\end{remark}
% Fuente francesa: PMIHES_1961__8__5_0, printed pp. 130-131.

\subsection{Aplicación: criterios de amplitud}
\label{subsection:II.6.6-fr}

\begin{proposition}[6.6.1]
\label{II.6.6.1-fr}
Sean $Y$ un preesquema, $f:X\to Y$ un morfismo cuasicompacto y $g:X'\to X$ un morfismo finito y sobreyectivo.
Se supone, o bien que $g_*(\mathcal{O}_{X'})$ es un $\mathcal{O}_X$-módulo localmente libre, o bien que $(X,\mathcal{O}_X)$ y $g_*(\mathcal{O}_{X'})$ satisfacen \textup{(II~bis)}.
Entonces, para todo $\mathcal{O}_{X'}$-módulo invertible $\mathcal{L}'$ amplio para $f\circ g$, $N_{X'/X}(\mathcal{L}')=\mathcal{L}$ es amplio para $f$.
\end{proposition}

\begin{proof}
Se puede suponer que $Y$ es afín \sref{II.4.6.4-fr}, y entonces, en virtud de \sref{II.4.6.6-fr}, el enunciado equivale al siguiente.
\end{proof}

\begin{corollary}[6.6.2]
\label{II.6.6.2-fr}
Sean $X$ un preesquema cuasicompacto y $g:X'\to X$ un morfismo finito sobreyectivo tal que, o bien $g_*(\mathcal{O}_{X'})$ sea un $\mathcal{O}_X$-módulo localmente libre, o bien $(X,\mathcal{O}_X)$ y $g_*(\mathcal{O}_{X'})$ satisfagan \textup{(II~bis)}.
Entonces, para todo $\mathcal{O}_{X'}$-módulo amplio $\mathcal{L}'$, $\mathcal{L}=N_{X'/X}(\mathcal{L}')$ es amplio.
\end{corollary}

\begin{proof}
En la segunda hipótesis, se puede suponer además que el homomorfismo canónico
$g_*(\mathcal{O}_{X'})\to
g_*(\mathcal{O}_{X'})\otimes_{\mathcal{O}_X}\mathcal{R}(X)$
es \emph{inyectivo}.

En efecto, si no es así, sea $\mathcal{T}$ 
\oldpage[II]{131}
el núcleo de este homomorfismo, que es un ideal coherente de $g_*(\mathcal{O}_{X'})=\mathcal{B}$ \sref[I]{I.6.1.1-fr}, y pongamos $X''=\operatorname{Spec}(\mathcal{B}/\mathcal{T})$;
se tiene así un diagrama conmutativo
\[
\xymatrix{
X''\ar[rr]^h\ar[dr]_{g'} && X'\ar[dl]^g\\
& X
}
\]
en el que $h$ es una inmersión cerrada \sref{II.1.4.10-fr}.
Además, se sabe que el soporte de $\mathcal{T}$ es un conjunto cerrado \sref[0]{0.5.2.2-fr} no denso en $X$ \sref[I]{I.7.4.6-fr}, de donde se concluye que, para el punto genérico $x$ de una componente irreducible de $X$, existe una vecindad abierta afín $U$ de $x$ tal que
$\mathcal{B}|U=(\mathcal{B}/\mathcal{T})|U$.
Como $g$ es sobreyectivo por hipótesis, se deduce que $x\in g'(X'')$;
por tanto, $g'$ es dominante y, como es un morfismo finito, es \emph{sobreyectivo} \sref{II.6.1.10-fr}.
Por definición,
$g'_*(\mathcal{O}_{X''})\otimes_{\mathcal{O}_X}\mathcal{R}(X)
=(\mathcal{B}/\mathcal{T})\otimes_{\mathcal{O}_X}\mathcal{R}(X)
=g_*(\mathcal{O}_{X'})\otimes_{\mathcal{O}_X}\mathcal{R}(X)$,
así, $(X,\mathcal{O}_X)$ y $g'_*(\mathcal{O}_{X''})$ satisfacen (II~\emph{bis}) y, además,
$g'_*(\mathcal{O}_{X''})\to
g'_*(\mathcal{O}_{X''})\otimes_{\mathcal{O}_X}\mathcal{R}(X)$
es inyectivo.
Finalmente, $h^*(\mathcal{L}')=\mathcal{L}''$ es un $\mathcal{O}_{X''}$-módulo amplio \sref{II.4.6.13-fr}[(i~\emph{bis})], y se tiene
$N_{X''/X}(\mathcal{L}'')=N_{X'/X}(\mathcal{L}')$.
% Fuente francesa: PMIHES_1961__8__5_0, printed pp. 131-132.

En efecto, para definir estos dos $\mathcal{O}_X$-módulos invertibles se puede utilizar una misma cubierta abierta afín $(U_\lambda)$ de $X$ tal que las restricciones de $g_*(\mathcal{L}')$ y $g'_*(\mathcal{L}'')$ a $U_\lambda$ sean isomorfas, respectivamente, a $\mathcal{B}|U_\lambda$ y $(\mathcal{B}/\mathcal{T})|U_\lambda$.
Se ve de inmediato que a todo isomorfismo
$\eta_\lambda:g_*(\mathcal{L}')|U_\lambda\to\mathcal{B}|U_\lambda$
corresponde canónicamente un isomorfismo
\[
\eta'_\lambda:g'_*(\mathcal{L}'')|U_\lambda\longrightarrow
(\mathcal{B}/\mathcal{T})|U_\lambda
\]
de suerte que, si $(\omega_{\lambda\mu})$ y $(\omega'_{\lambda\mu})$ son los $1$-cociclos correspondientes a los sistemas de isomorfismos $(\eta_\lambda)$ y $(\eta'_\lambda)$ \sref{II.6.5.2-fr}, $\omega'_{\lambda\mu}$ sea la imagen canónica en
$\Gamma(U_\lambda\cap U_\mu,\mathcal{B}/\mathcal{T})$
de $\omega_{\lambda\mu}\in\Gamma(U_\lambda\cap U_\mu,\mathcal{B})$.
En virtud de la definición de $\mathcal{T}$, se concluye que
\[
N_{\mathcal{B}/\mathcal{A}}\omega_{\lambda\mu}
=N_{(\mathcal{B}/\mathcal{T})/\mathcal{A}}\omega'_{\lambda\mu}
\]
(donde $\mathcal{A}=\mathcal{O}_X$), de donde resulta la igualdad anunciada.

Supongamos, pues, que el homomorfismo
$g_*(\mathcal{O}_{X'})\to
g_*(\mathcal{O}_{X'})\otimes_{\mathcal{O}_X}\mathcal{R}(X)$
es inyectivo cuando se está en la hipótesis~(II~\emph{bis}).
Basta probar que, cuando $f$ recorre las secciones de $\mathcal{L}^{\otimes n}$ ($n>0$) sobre $X$, los $X_f$ forman una base de la topología de $X$ \sref{II.4.5.2-fr}.
Sea, en efecto, $x\in X$, y sea $U$ una vecindad cualquiera de $x$;
como $g^{-1}(x)$ es finito \sref{II.6.1.7-fr} y $\mathcal{L}'$ es amplio, existe un entero $n>0$ y una sección $f'$ de ${\mathcal{L}'}^{\otimes n}$ sobre $X'$ tales que $X'_{f'}$ sea una vecindad de $g^{-1}(x)$ contenida en $g^{-1}(U)$ \sref{II.4.5.4-fr}.
Como
\[
\mathcal{L}^{\otimes n}=N_{X'/X}({\mathcal{L}'}^{\otimes n})
\]
basta tomar $f=N_{X'/X}(f')$;
en efecto, entonces
$X\setmin X_f=g(X'\setmin X'_{f'})$
\sref{II.6.5.7-fr}, y, por tanto, $x\in X_f\subset U$.
\end{proof}

\begin{corollary}[6.6.3]
\label{II.6.6.3-fr}
Bajo las hipótesis de \sref{II.6.6.1-fr}, para que un $\mathcal{O}_X$-módulo invertible $\mathcal{L}$ sea amplio para $f$, es necesario y suficiente que $\mathcal{L}'=g^*(\mathcal{L})$ sea amplio para $f\circ g$.
\end{corollary}

\begin{proof}
La condición es necesaria, pues $g$ es afín \sref{II.5.1.12-fr}.
Para probar que es suficiente, se puede suponer que $Y$ es afín \sref{II.4.6.4-fr};
entonces $X$ y $X'$ son cuasicompactos y $\mathcal{L}'$ es amplio \sref{II.4.6.6-fr}, y hay que probar que $\mathcal{L}$ es amplio.
\oldpage[II]{132}
Ahora bien, el conjunto de los puntos $x\in X$ que admiten una vecindad en la que $g_*(\mathcal{O}_{X'})$ (resp. $g_*(\mathcal{O}_{X'})\otimes_{\mathcal{O}_X}\mathcal{R}(X)$) tiene un rango dado~$n$, en la primera (resp. segunda) hipótesis, es a la vez abierto y cerrado en $X$;
por tanto, $X$ es el preesquema suma de un número finito de estos abiertos, y se puede suponer que es igual a uno de ellos \sref{II.4.6.17-fr}.
Pero entonces
$N_{X'/X}(\mathcal{L}')=\mathcal{L}^{\otimes n}$,
de modo que $\mathcal{L}^{\otimes n}$ es amplio por \sref{II.6.6.2-fr}, y, por consiguiente, también lo es $\mathcal{L}$ \sref{II.4.5.6-fr}.
\end{proof}
% Fuente francesa: PMIHES_1961__8__5_0, printed p. 132.

\begin{corollary}[6.6.4]
\label{II.6.6.4-fr}
Supongamos satisfechas las hipótesis de \sref{II.6.6.1-fr}, y supongamos además que $f:X\to Y$ es de tipo finito.
Entonces, para que $f$ sea cuasiproyectivo, es necesario y suficiente que $f\circ g$ sea cuasiproyectivo.
Si se supone además que $Y$ es un esquema cuasicompacto, o un preesquema cuyo espacio subyacente es noetheriano, para que $f$ sea proyectivo es necesario y suficiente que $f\circ g$ lo sea.
\end{corollary}

\begin{proof}
La hipótesis implica que $f\circ g$ es de tipo finito.
Teniendo en cuenta la definición de los morfismos cuasiproyectivos \sref{II.5.3.1-fr}, la primera afirmación resulta de \sref{II.6.6.1-fr} y \sref{II.6.6.3-fr}.
Teniendo en cuenta este resultado y \sref{II.5.5.3-fr}[(ii)], queda por probar que, si $f$ es cuasiproyectivo, para que $f$ sea propio es necesario y suficiente que $f\circ g$ lo sea.
Pero $f$ es entonces separado \sref{II.5.3.1-fr} y de tipo finito;
como $g$ es sobreyectivo, nuestra afirmación resulta de \sref{II.5.4.2-fr}[(ii)] y \sref{II.5.4.3-fr}[(ii)].
\end{proof}

En particular:

\begin{corollary}[6.6.5]
\label{II.6.6.5-fr}
Sean $X$ un preesquema de tipo finito sobre un cuerpo $K$ y $K'$ una extensión de grado finito de $K$.
Para que $X$ sea proyectivo \textup{(resp.} cuasiproyectivo\textup{)} sobre $K$, es necesario y suficiente que $X'=X\otimes_K K'$ sea proyectivo \textup{(resp.} cuasiproyectivo\textup{)} sobre $K'$.
\end{corollary}

\begin{proof}
La condición es necesaria (\sref{II.5.3.4-fr}[(iii)] y \sref{II.5.5.5-fr}[(iii)]).
Recíprocamente, supongámosla satisfecha, y sea $g:X'\to X$ la proyección canónica.
Es claro que $g$ es un morfismo finito \sref{II.6.1.5-fr}[(iii)] y sobreyectivo \sref[I]{I.3.5.2-fr}[(ii)].
Además, $g_*(\mathcal{O}_{X'})$ es un $\mathcal{O}_X$-módulo localmente libre, pues es isomorfo a $\mathcal{O}_X\otimes_K K'$ \sref{II.1.5.2-fr}.
Se sigue entonces de la hipótesis y de \sref{II.6.1.11-fr} y \sref{II.5.5.5-fr}[(ii)] que $X'$ es proyectivo (resp. cuasiproyectivo) \emph{sobre $K$};
se deduce entonces de \sref{II.6.6.4-fr} que $X$ es proyectivo (resp. cuasiproyectivo) sobre $K$.
\end{proof}

En el capítulo~IV mostraremos que el enunciado de \sref{II.6.6.5-fr} sigue siendo válido cuando $K'$ es una extensión \emph{cualquiera} de $K$.

El resto de este número está dedicado a la demostración del criterio de \sref{II.6.6.11-fr}, que es un refinamiento bastante técnico de \sref{II.6.6.1-fr};
puede omitirse en una primera lectura.
% Fuente francesa: PMIHES_1961__8__5_0, printed pp. 132-133.

\begin{lemma}[6.6.6]
\label{II.6.6.6-fr}
Sean $X$ un preesquema noetheriano reducido y $\mathcal{E}$ un $\mathcal{O}_X$-módulo coherente tal que $\mathcal{E}\otimes_{\mathcal{O}_X}\mathcal{R}(X)$ sea un $\mathcal{R}(X)$-módulo localmente libre de rango constante~$n$.
Existe entonces un preesquema noetheriano reducido $Z$ y un morfismo finito birracional $h:Z\to X$ que tiene la propiedad siguiente: los morfismos de haces de conjuntos
$\sigma_i:\mathcal{H}om_{\mathcal{O}_X}(\mathcal{E},\mathcal{E})
\to\mathcal{R}(X)$ \textup{($1\leq i\leq n$)}
\textup{(cf. \sref{II.6.4.10-fr})} envían $\mathcal{H}om_{\mathcal{O}_X}(\mathcal{E},\mathcal{E})$ en la $\mathcal{O}_X$-álgebra coherente $h_*(\mathcal{O}_Z)$.
\end{lemma}

\begin{proof}
Consideremos un abierto afín $U$ de $X$, de anillo $A(U)=A$;
sea $E=\Gamma(U,\mathcal{E})$, y sea $C_U$ la subálgebra de $R(U)$ generada por los $\sigma_i(u)$ cuando $u$ recorre $\operatorname{Hom}_A(E,E)$;
se ha visto en \sref{II.6.4.7.1-fr} que esta $A$-álgebra es de \emph{rango finito}.
Además, es claro que la formación de las álgebras $C_U$ conmuta con las operaciones de restricción de un abierto afín $U$ a un abierto afín $U'\subset U$.
Se ha definido así una sub-$\mathcal{O}_X$-álgebra \emph{finita} $\mathcal{C}$ de $\mathcal{R}(X)$ tal que $\Gamma(U,\mathcal{C})=C_U$ para todo abierto afín $U$ de $X$.
\oldpage[II]{133}
Se tomará $Z=\operatorname{Spec}(\mathcal{C})$, y para $h$ el morfismo estructural, que es, por tanto, finito \sref{II.6.1.2-fr};
como $\mathcal{C}$ es reducida, $Z$ es un preesquema noetheriano reducido \sref{II.1.3.8-fr}.
Finalmente, el anillo total de fracciones de $C_U$ es $R(U)$ por definición, y como $C_U$ está contenido en la clausura integral de $A(U)$ en $R(U)$, existe una correspondencia biyectiva entre los ideales primos minimales de $A(U)$ y los ideales primos minimales de $C_U$ \cite[t.~I, p.~259]{I-13}, lo que prueba que $h$ es birracional y termina la demostración.
\end{proof}
% Fuente francesa: PMIHES_1961__8__5_0, printed p. 133.

\begin{corollary}[6.6.7]
\label{II.6.6.7-fr}
Bajo las hipótesis de \sref{II.6.6.6-fr}, sea $W$ un abierto de $X$ tal que, para todo $x\in W$, o bien $X$ sea normal en el punto $x$, o bien $\mathcal{E}_x$ sea un $\mathcal{O}_x$-módulo libre.
Puede entonces suponerse que $h$ está definido de tal modo que la restricción de $h$ a $h^{-1}(W)$ sea un isomorfismo de $h^{-1}(W)$ sobre $W$.
\end{corollary}

\begin{proof}
Una u otra hipótesis implica que, si $U\subset W$ es un abierto afín, se tiene, con las notaciones de \sref{II.6.6.6-fr},
$(\sigma_i(u))_x\in A_x$
para todo $x\in U$ \sref{II.6.4.3-fr}, luego $\sigma_i(u)\in A$, y la conclusión resulta de la definición de $h$ dada en \sref{II.6.6.6-fr}.
\end{proof}

\begin{env}[6.6.8]
\label{II.6.6.8-fr}
Sean $X$ un preesquema noetheriano reducido y $g:X'\to X$ un morfismo finito sobreyectivo, de suerte que $\mathcal{B}=g_*(\mathcal{O}_{X'})$ sea una $\mathcal{O}_X$-álgebra \emph{coherente};
supongamos además que $\mathcal{B}\otimes_{\mathcal{O}_X}\mathcal{R}(X)$ es un $\mathcal{R}(X)$-módulo \emph{localmente libre de rango constante~$n$}.
Puede aplicarse entonces el lema \sref{II.6.6.6-fr}, tomando $\mathcal{E}=\mathcal{B}$;
de donde, con las notaciones de \sref{II.6.6.6-fr}, se obtiene un homomorfismo de haces de monoides multiplicativos
$\sigma_n:\mathcal{H}om_{\mathcal{O}_X}(\mathcal{B},\mathcal{B})
\to h_*(\mathcal{O}_Z)$,
y, componiendo este homomorfismo con el homomorfismo canónico
$\mathcal{B}\to
\mathcal{H}om_{\mathcal{O}_X}(\mathcal{B},\mathcal{B})$
\sref{II.6.5.1-fr}, se obtiene un homomorfismo de haces de monoides multiplicativos
\[
\label{II.6.6.8.1-fr}
N':\mathcal{B}=g_*(\mathcal{O}_{X'})\longrightarrow
h_*(\mathcal{O}_Z)=\mathcal{C}
\tag{6.6.8.1}
\]

Sentado esto, para todo $\mathcal{O}_{X'}$-módulo invertible $\mathcal{L}'$, $g_*(\mathcal{L}')$ es un $\mathcal{B}$-módulo invertible \sref{II.6.1.12-fr}, y el método de \sref{II.6.5.2-fr} permite asociar funtorialmente a $\mathcal{L}'$ un $\mathcal{C}$-módulo invertible, que denotaremos por
$N'(g_*(\mathcal{L}'))$.
\end{env}
% Fuente francesa: PMIHES_1961__8__5_0, printed pp. 133-134.

\begin{lemma}[6.6.9]
\label{II.6.6.9-fr}
Sean $X$ un preesquema noetheriano reducido y $g:X'\to X$ un morfismo finito sobreyectivo tal que
$g_*(\mathcal{O}_{X'})\otimes_{\mathcal{O}_X}\mathcal{R}(X)$
sea un $\mathcal{R}(X)$-módulo localmente libre de rango constante.
Existe entonces un preesquema noetheriano reducido $Z$ y un morfismo finito birracional $h:Z\to X$ que tiene la propiedad siguiente: para todo $\mathcal{O}_{X'}$-módulo amplio $\mathcal{L}'$, el $\mathcal{O}_Z$-módulo invertible $\mathcal{M}$ tal que
$h_*(\mathcal{M})=N'(g_*(\mathcal{L}'))$
\textup{(notaciones de \sref{II.6.6.8-fr})} es amplio.
\end{lemma}

\begin{proof}
Supongamos primero que el homomorfismo
$\mathcal{B}\to
\mathcal{B}\otimes_{\mathcal{O}_X}\mathcal{R}(X)$
es \emph{inyectivo}.
Definamos $Z$ y $h$ como en \sref{II.6.6.6-fr} (con $\mathcal{E}=g_*(\mathcal{O}_{X'})$).
Sea $z\in Z$;
hay que probar que existe un entero $m>0$ y una sección $t$ de $\mathcal{M}^{\otimes m}$ sobre $Z$ tales que $Z_t$ sea un abierto afín que contenga a $z$ \sref{II.4.5.2-fr}.
Sea $x=h(z)$, y sea $U$ un abierto afín de $X$ que contenga a $x$;
entonces $h^{-1}(U)$ es una vecindad abierta afín de $z$, y basta encontrar $t$ tal que
$z\in Z_t\subset h^{-1}(U)$,
pues $Z_t$ será entonces necesariamente afín \sref{II.5.5.8-fr}.
Existe, por hipótesis, un entero $n>0$ y una sección $s'$ de ${\mathcal{L}'}^{\otimes n}$ sobre $X'$ tales que
\[
\label{II.6.6.9.1-fr}
g^{-1}(x)\subset X'_{s'}\subset g^{-1}(U)
\tag{6.6.9.1}
\]
en virtud de \sref{II.4.5.4-fr}.
Por definición, $s'$ es también una sección de $g_*({\mathcal{L}'}^{\otimes n})$ sobre $X$, y le corresponde, como en \sref{II.6.5.2-fr}, una sección
$s=N'(s')$
de
$N'(g_*({\mathcal{L}'}^{\otimes n}))$
sobre
\oldpage[II]{134}
 $X$.

Vamos a ver que, si $t$ es la sección $s$ considerada como sección de $\mathcal{M}^{\otimes n}$ sobre $Z$, entonces $t$ responde a la cuestión.
Pongamos
\[
\label{II.6.6.9.2-fr}
V=X\setmin g(X'\setmin X'_{s'})
\tag{6.6.9.2}
\]
que es un abierto de $X$ que contiene a $x$ y está contenido en $U$, en virtud de \sref{II.6.6.9.1-fr} y \sref{II.6.1.10-fr}.
Vamos a demostrar que
\[
\label{II.6.6.9.3-fr}
h^{-1}(V)\subset Z_t\subset h^{-1}(U),
\tag{6.6.9.3}
\]
lo que terminará la demostración.
% Fuente francesa: PMIHES_1961__8__5_0, printed p. 134.

Esto equivale a decir que el conjunto $T$ de los $y\in X$ tales que $s_y$ sea invertible contiene a $V$ y está contenido en $U$.
Para ello, consideremos primero un abierto afín $W$ contenido en $V$;
entonces $g^{-1}(W)$ es un abierto afín de $X'$ y, en virtud de \sref{II.6.6.9.2-fr}, $s'_{y'}$ es invertible para todo $y'\in g^{-1}(W)$;
en virtud de las hipótesis sobre $X$ y $\mathcal{B}$, pueden aplicarse los resultados de \sref{II.6.4.7-fr}, y se ve que, si $y=g(y')$, entonces $s_y$ es invertible;
en otras palabras, $V\subset T$.
Por otra parte, se sigue también de \sref{II.6.4.7-fr} que, recíprocamente, si $s_y$ es invertible, ocurre lo mismo con $s'_{y'}$, lo que implica que $y'\in g^{-1}(U)$ por \sref{II.6.6.9.1-fr}, luego $y\in U$, de donde $T\subset U$ en este caso.

Se pasa de aquí al caso general mediante el mismo razonamiento que en \sref{II.6.6.2-fr}, sustituyendo $X'$ por $X''$ de modo que
$g'_*(\mathcal{O}_{X''})\to
g'_*(\mathcal{O}_{X''})\otimes_{\mathcal{O}_X}\mathcal{R}(X)$
sea inyectivo, y $\mathcal{L}'$ por un $\mathcal{O}_{X''}$-módulo amplio $\mathcal{L}''$ tal que
$N'(g_*(\mathcal{L}'))=N'(g'_*(\mathcal{L}''))$.
El lema \sref{II.6.6.9-fr} queda así demostrado en todos los casos (con una elección conveniente de $h$).
\end{proof}

\begin{corollary}[6.6.10]
\label{II.6.6.10-fr}
Supongamos satisfechas las hipótesis de \sref{II.6.6.9-fr};
para todo $\mathcal{O}_X$-módulo invertible $\mathcal{L}$ tal que $g^*(\mathcal{L})$ sea amplio, $h^*(\mathcal{L})$ es amplio.
\end{corollary}

\begin{proof}
En efecto, si se pone $\mathcal{L}'=g^*(\mathcal{L})$, se tiene
$g_*(\mathcal{L}')=\mathcal{L}\otimes_{\mathcal{O}_X}\mathcal{B}$
\sref[0]{0.5.4.10-fr}, luego
\[
N'(g_*(\mathcal{L}'))=
(\mathcal{L}\otimes_{\mathcal{O}_X}\mathcal{C})^{\otimes n}
\]
(mediante el mismo razonamiento que para \sref{II.6.5.2.4-fr}).
Se concluye que
$\mathcal{M}=(h^*(\mathcal{L}))^{\otimes n}$,
y, como $\mathcal{M}$ es amplio, también lo es $h^*(\mathcal{L})$ \sref{II.4.5.6-fr}.
\end{proof}
% Fuente francesa: PMIHES_1961__8__5_0, printed pp. 134-135.

\begin{proposition}[6.6.11]
\label{II.6.6.11-fr}
Sean $Y$ un esquema afín, $X$ un preesquema noetheriano reducido, $f:X\to Y$ un morfismo cuasicompacto y $g:X'\to X$ un morfismo finito sobreyectivo.
Sea $W$ una parte abierta de $X$ tal que, para todo $x\in W$, o bien $X$ sea normal en el punto $x$, o bien exista una vecindad abierta $T\subset W$ de $x$ tal que
$(g_*(\mathcal{O}_{X'}))|T$
sea un $(\mathcal{O}_X|T)$-módulo localmente libre.
Existe entonces un $Y$-preesquema reducido $Z$ y un $Y$-morfismo \emph{finito birracional} $h:Z\to X$ tales que la restricción de $h$ a $h^{-1}(W)$ sea un isomorfismo $h^{-1}(W)\xrightarrow{\sim}W$ y que se tenga la propiedad siguiente: para todo $\mathcal{O}_X$-módulo invertible $\mathcal{L}$ tal que $g^*(\mathcal{L})$ sea amplio relativamente a $f\circ g$, $h^*(\mathcal{L})$ es amplio relativamente a $f\circ h$.
\end{proposition}

\begin{proof}
Como $Y$ es afín, $g^*(\mathcal{L})$ es amplio, y se trata entonces de probar que, para una elección conveniente de $h$, $h^*(\mathcal{L})$ es amplio \sref{II.4.6.6-fr}.
Vamos a ver que se puede sustituir $g$ por un morfismo finito sobreyectivo $g':X''\to X$ tal que ${g'}^*(\mathcal{L})$ sea amplio y
$g'_*(\mathcal{O}_{X''})\otimes_{\mathcal{O}_X}\mathcal{R}(X)$
sea un $\mathcal{R}(X)$-módulo \emph{localmente libre de rango constante};
se estará entonces reducido a las condiciones de \sref{II.6.6.10-fr}, y la proposición quedará demostrada.

Para ello, sea $\mathcal{B}=g_*(\mathcal{O}_{X'})$;
designemos por $X_i$ ($1\leq i\leq n$) los subpreesquemas cerrados reducidos de $X$ cuyos espacios subyacentes son las componentes irreducibles de $X$ \sref[I]{I.5.2.1-fr};
son íntegros por hipótesis.

Sea $X'_i$ el subpreesquema cerrado
\oldpage[II]{135}
 $g^{-1}(X_i)$ de $X'$, y sea $g_i:X'_i\to X_i$ el morfismo restricción de $g$ a $X'_i$, que es finito \sref{II.6.1.5-fr}[(iii)] y sobreyectivo;
sea $k_i$ el \emph{rango} de la $\mathcal{O}_{X_i}$-álgebra
$\mathcal{B}_i=(g_i)_*(\mathcal{O}_{X'_i})$.
Como $\mathcal{B}_i\otimes_{\mathcal{O}_{X_i}}\mathcal{R}(X_i)$ es un prehaz constante \sref[I]{I.7.3.5-fr}, el rango $k_i$ es también el rango de la $(\mathcal{O}_X|U)$-álgebra $\mathcal{B}|U$ para todo abierto $U$ de $X$ que no encuentre más componente irreducible que $X_i$.
Si $T$ es un conjunto abierto de $X$ tal que $\mathcal{B}|T$ sea isomorfo a $\mathcal{O}_X^m|T$, resulta de la observación anterior que los números $k_i$ son iguales a $m$ para todos los índices $i$ tales que $T\cap X_i\neq\varnothing$.
% Fuente francesa: PMIHES_1961__8__5_0, printed p. 135.

Sea entonces $U$ el abierto de $X$ formado por los puntos que admiten una vecindad en la que $\mathcal{B}$ sea un $\mathcal{O}_X$-módulo localmente libre, y sean $U_j$ ($1\leq j\leq s$) sus componentes conexas, que son abiertas en $X$ y en número finito (pues $U$ es noetheriano);
designemos por $V_j$ el subpreesquema cerrado de $X'$, clausura del subpreesquema inducido sobre el abierto $g^{-1}(U_j)$ \sref[I]{I.9.5.11-fr}.
Por lo anterior, para todos los índices $i$ tales que $X_i\cap U_j\neq\varnothing$, los rangos $k_i$ son iguales a un mismo entero $m_j$;
se señala, además, que un mismo $X_i$ no puede encontrar dos $U_j$ de índices distintos.
Sean $i_\lambda$ los índices $i$ tales que $X_i\cap U=\varnothing$.
Consideremos el producto $k$ de todos los $k_i$, pongamos $n_i=k/k_i$, y sea $X''$ el preesquema definido del modo siguiente.
Para cada $j$ ($1\leq j\leq s$), se consideran $k/m_j$ preesquemas isomorfos a $V_j$, y para cada $\lambda$, $k/k_{i_\lambda}=n_{i_\lambda}$ preesquemas isomorfos a $X'_{i_\lambda}$;
$X''$ es la \emph{suma} de todos estos preesquemas.
Se define un morfismo $g'':X''\to X'$ que se reduce a la inyección canónica en cada uno de los sumandos de $X''$;
es claro que $g''$ es un morfismo finito dominante, luego sobreyectivo (pues un morfismo finito es cerrado \sref{II.6.1.10-fr});
pongamos $g'=g\circ g''$, que es un morfismo finito sobreyectivo $X''\to X$.
Se tiene
${g'}^*(\mathcal{L})={g''}^*(g^*(\mathcal{L}))$,
de modo que ${g'}^*(\mathcal{L})$ es un $\mathcal{O}_{X''}$-módulo amplio \sref{II.5.1.12-fr}.
Es claro entonces que, para este nuevo preesquema $X''$, los rangos definidos como los $k_i$ para $X'$ son todos \emph{iguales a $k$};
teniendo en cuenta \sref[I]{I.7.3.3-fr}, se concluye inmediatamente que, para todo abierto afín $T$ de $X$,
$\bigl(g'_*(\mathcal{O}_{X''})\otimes_{\mathcal{O}_X}\mathcal{R}(X)\bigr)|T$
es un $(\mathcal{R}(X)|T)$-módulo isomorfo a $(\mathcal{R}(X)|T)^k$.
\end{proof}

\begin{corollary}[6.6.12]
\label{II.6.6.12-fr}
Si, en el enunciado de \sref{II.6.6.11-fr}, se tiene $W=X$, entonces, para que un $\mathcal{O}_X$-módulo invertible $\mathcal{L}$ sea amplio relativamente a $f$, es necesario y suficiente que $g^*(\mathcal{L})$ sea amplio relativamente a $f\circ g$.
\end{corollary}

\begin{remark}[6.6.13]
\label{II.6.6.13-fr}
En el capítulo~III veremos que, si $Y$ es noetheriano, $f$ es de tipo finito y la restricción de $f$ al subpreesquema cerrado reducido de $X$ cuyo espacio subyacente es $X\setmin W$ es \emph{propia}, la conclusión de \sref{II.6.6.12-fr} sigue siendo válida.
Pero daremos en el capítulo~V ejemplos de esquemas algebraicos $X$ sobre un cuerpo $K$ (sin que el morfismo estructural $X\to\operatorname{Spec}(K)$ sea propio) cuya normalización $X'$ es cuasiafín, pero que no son cuasiafines (de suerte que $\mathcal{O}_X$ no es amplio, aunque $\mathcal{O}_{X'}$ lo sea \textup{(\hyperref[II.5.1.2-fr]{5.1.2})} y el morfismo $X'\to X$ sea finito sobreyectivo \sref{II.6.3.10-fr}).
Veremos en el número siguiente que esta circunstancia no puede producirse cuando se sustituye «cuasiafín» por «afín».
\end{remark}
% Fuente francesa: PMIHES_1961__8__5_0, printed pp. 135-136.

\subsection{Teorema de Chevalley}
\label{subsection:II.6.7-fr}

Vamos a establecer (con ayuda del criterio de Serre \sref{II.5.2.1-fr}) el teorema siguiente, demostrado por C.~Chevalley mediante otros métodos en el caso de los esquemas \emph{algebraicos}.

\oldpage[II]{136}
\begin{theorem}[6.7.1]
\label{II.6.7.1-fr}
Sean $X$ un esquema afín, $Y$ un preesquema noetheriano y $f:X\to Y$ un morfismo finito sobreyectivo.
Entonces $Y$ es un esquema afín.
\end{theorem}

\begin{proof}
Es claro que
$f_{\mathrm{red}}:X_{\mathrm{red}}\to Y_{\mathrm{red}}$
es finito \sref{II.6.1.5-fr}[(vi)];
como $X_{\mathrm{red}}$ es un esquema afín y decir que $Y$ es afín equivale a decir que $Y_{\mathrm{red}}$ es afín (pues $Y$ es noetheriano \sref[I]{I.6.1.7-fr}), se ve que puede suponerse que $Y$ es \emph{reducido}.
Para toda parte cerrada $Y'$ de $Y$, existe un único subpreesquema reducido de $Y$ que tiene $Y'$ por espacio subyacente \sref[I]{I.5.1.2-fr};
su imagen inversa $f^{-1}(Y')$, canónicamente isomorfa a $X\times_Y Y'$ \sref[I]{I.4.4.1-fr}, es afín como subpreesquema cerrado de $X$, y la restricción de $f$ a $f^{-1}(Y')$, que se identifica con $f\times_Y 1_{Y'}$, es un morfismo finito sobreyectivo \sref{II.6.1.5-fr}[(iii)].
En virtud del principio de recurrencia noetheriana \sref[0]{0.2.2.2-fr}, puede reducirse, pues (teniendo en cuenta \sref[I]{I.6.1.7-fr}), a probar el teorema bajo la hipótesis de que, para \emph{toda} parte cerrada $Y'\neq Y$, todo subpreesquema cerrado de $Y$ que tenga $Y'$ por espacio subyacente sea afín.
Se concluye entonces que, \emph{para todo $\mathcal{O}_Y$-módulo coherente $\mathcal{F}$ cuyo soporte (cerrado) $Z$ sea distinto de $Y$, se tiene $\mathrm{H}^1(Y,\mathcal{F})=0$}.
En efecto, existe un subpreesquema cerrado $Y'$ de $Y$ que tiene $Z$ por espacio subyacente y tal que, si $j:Y'\to Y$ es la inyección canónica, se tenga
$\mathcal{F}=j_*(j^*(\mathcal{F}))$
\sref[I]{I.9.3.5-fr};
entonces
$\mathrm{H}^1(Y,\mathcal{F})
=\mathrm{H}^1(Y',j^*(\mathcal{F}))=0$
por \sref{II.5.2.3-fr} y \sref[I]{I.5.1.9.2-fr}.

Supongamos primero que $Y$ no es irreducible, y sean $Y'$ una componente irreducible de $Y$ y $Y''=Y\setmin Y'$;
se denota también por $Y'$ el subpreesquema cerrado reducido de $Y$ que tiene $Y'$ por espacio subyacente, y por $j$ la inyección canónica $Y'\to Y$.
Sea $\mathcal{F}$ un $\mathcal{O}_Y$-módulo coherente, y consideremos el homomorfismo canónico
\[
\rho:\mathcal{F}\longrightarrow
\mathcal{F}'=j_*(j^*(\mathcal{F}))
\]
\sref[0]{0.4.4.3-fr};
$\mathcal{F}'$ es un $\mathcal{O}_Y$-módulo coherente en virtud de \sref[0]{0.5.3.10-fr} y \sref[0]{0.5.3.12-fr}, pues
$j_*(\mathcal{O}_{Y'})=\mathcal{O}_Y/\mathcal{J}$,
donde $\mathcal{J}$ denota el haz de ideales de $\mathcal{O}_Y$ que define el subpreesquema $Y'$.
Entonces $\mathcal{G}=\operatorname{Ker}\rho$ y $\mathcal{K}=\operatorname{Im}\rho$ son también $\mathcal{O}_Y$-módulos coherentes \sref[0]{0.5.3.4-fr};
pero, por definición, el tallo $\mathcal{F}'_y$ de $\mathcal{F}'$ en el punto genérico $y$ de $Y'$ es igual a $\mathcal{F}_y$, pues $y$ es interior a $Y'$ y se tiene, por tanto, $\mathcal{J}_y=0$, ya que $Y$ es reducido.
Se concluye que $y$ no pertenece al soporte (cerrado) de $\mathcal{G}$;
por otra parte, el soporte de $\mathcal{F}'$ (y \emph{a fortiori} el de $\mathcal{K}$) está contenido en $Y'$;
en otras palabras, los soportes de $\mathcal{G}$ y $\mathcal{K}$ son \emph{distintos de $Y$}.
Se deduce que
$\mathrm{H}^1(Y,\mathcal{G})=\mathrm{H}^1(Y,\mathcal{K})=0$,
y la sucesión exacta de cohomología aplicada a
$0\to\mathcal{G}\to\mathcal{F}\to\mathcal{K}\to0$
implica que $\mathrm{H}^1(Y,\mathcal{F})=0$.
Se concluye entonces mediante el criterio de Serre \sref{II.5.2.1-fr}.
% Fuente francesa: PMIHES_1961__8__5_0, printed pp. 136-137.

Supongamos, pues, que $Y$ es irreducible y, por tanto, \emph{íntegro}.
También puede suponerse que $X$ es \emph{íntegro}:
si se denotan por $X_i$ los subpreesquemas cerrados reducidos de $X$ cuyos espacios subyacentes son las componentes irreducibles de $X$ \sref[I]{I.5.2.1-fr}, y por $f_i$ la restricción de $f$ a $X_i$, al menos uno de los $f_i$ es dominante y, como es un morfismo finito \sref{II.6.1.5-fr}, es sobreyectivo \sref{II.6.1.10-fr};
como $X_i$ es además un esquema afín, se ve que puede sustituirse $X$ por $X_i$ en el enunciado.

\begin{lemma}[6.7.1.1]
\label{II.6.7.1.1-fr}
Sean $X$ e $Y$ dos preesquemas íntegros noetherianos, $x$ \textup{(resp. $y$)} el punto genérico de $X$ \textup{(resp. de $Y$)}, y $f:X\to Y$ un morfismo finito sobreyectivo.
Sea $\mathcal{L}$ un $\mathcal{O}_X$-módulo invertible tal que exista una vecindad abierta afín $U$ de $y$ y una sección $g\in\Gamma(X,\mathcal{L})$ para las cuales
\oldpage[II]{137}

$x\in X_g\subset f^{-1}(U)$.

Existen entonces enteros $m>0$, $n>0$, un homomorfismo
$u:\mathcal{O}_Y^m\to f_*(\mathcal{L}^{\otimes n})$
y una vecindad abierta $V$ de $y$ tales que la restricción $u|V$ sea un isomorfismo de $\mathcal{O}_Y^m|V$ sobre $f_*(\mathcal{L}^{\otimes n})|V$.
\end{lemma}
% Fuente francesa: PMIHES_1961__8__5_0, printed p. 137.

\begin{proof}
Sea $C$ el anillo (íntegro) de $U$, y sea $k=\mathcal{O}_y$ su cuerpo de fracciones:
como $f$ es finito, $U'=f^{-1}(U)$ es afín \sref{II.1.3.2-fr};
sea $D$ su anillo (íntegro), de cuerpo de fracciones $K=\mathcal{O}_x$;
por hipótesis, $D$ es un $C$-módulo de tipo finito \sref{II.6.1.4-fr}, luego $K$ es una extensión de rango finito de $k$.
La fibra
$f^{-1}(y)=X\times_Y\operatorname{Spec}(\boldsymbol{k}(y))=X\times_Y\operatorname{Spec}(\mathcal{O}_y)$
se identifica con $\operatorname{Spec}(K)$ (\hyperref[I.3.6.5-fr]{\textbf{I}, 3.6.6});
sean $s_i$ ($1\leq i\leq m$) elementos de $D$ que formen una base de $K$ sobre $k$.
Existe un entero $n>0$ tal que las secciones
$(s_i|X_g)g^{\otimes n}$
de $\mathcal{L}^{\otimes n}$ sobre $X_g$ se prolonguen en secciones $b_i$ ($1\leq i\leq m$) de $\mathcal{L}^{\otimes n}$ sobre $X$ \sref[I]{I.9.3.1-fr}.
Los $b_i$ son también, por definición, secciones de $f_*(\mathcal{L}^{\otimes n})$ sobre $Y$, y definen, pues, un homomorfismo
$u:\mathcal{O}_Y^m\to f_*(\mathcal{L}^{\otimes n})$
\sref[0]{0.5.1.1-fr};
vamos a ver que $u$ responde a la cuestión.
Se tiene
$\mathcal{L}^{\otimes n}|U'=\widetilde{M}$,
donde $M$ es un $D$-módulo de tipo finito;
si $\varphi$ es la inyección $C\to D$ correspondiente al morfismo $f^{-1}(U)\to U$ dado por la restricción de $f$, entonces $M_{[\varphi]}$ es un $C$-módulo de tipo finito.
Por tanto,
\[
f_*(\mathcal{L}^{\otimes n})|U=(M_{[\varphi]})^\sim
\]
\sref[I]{I.1.6.3-fr} es coherente y, como $U$ es un abierto afín cualquiera de $Y$, $f_*(\mathcal{L}^{\otimes n})$ es coherente.
Además, $u|U=\widetilde{\theta}$, donde $\theta$ es un $C$-homomorfismo
$C^m\to M_{[\varphi]}$,
y $u_y=\theta_y$ es el homomorfismo
$\theta\otimes1:K^m=C^m\otimes K\longrightarrow M_{[\varphi]}\otimes K$
este último es, por definición, un \emph{isomorfismo}, pues los $(b_i)_x$ forman una base de $M_{[\varphi]}\otimes K$ sobre $k$, siendo $g_x$, por hipótesis, un generador de $(\mathcal{L}^{\otimes n})_x$.
Se concluye que los soportes de los $\mathcal{O}_Y$-módulos $\operatorname{Ker}u$ y $\operatorname{Coker}u$ no contienen a $y$;
como estos $\mathcal{O}_Y$-módulos son coherentes \sref[0]{0.5.3.3-fr}, sus soportes son cerrados \sref[0]{0.5.2.2-fr}, de donde resulta el lema.
\end{proof}
% Fuente francesa: PMIHES_1961__8__5_0, printed pp. 137-138.

Sentado esto, las hipótesis del lema \sref{II.6.7.1.1-fr} se satisfacen en el caso que consideramos tomando $\mathcal{L}=\mathcal{O}_X$, pues $X$ es afín \sref[I]{I.1.1.10-fr};
pongamos $\mathcal{A}=\mathcal{O}_Y$ y $\mathcal{B}=f_*(\mathcal{O}_X)$.
En virtud del criterio de Serre \sref{II.5.2.1.1-fr}, basta probar que, para todo $\mathcal{O}_Y$-módulo coherente $\mathcal{F}$, se tiene $\mathrm{H}^1(Y,\mathcal{F})=0$;
basta incluso probarlo cuando $\mathcal{F}\subset\mathcal{O}_Y$, lo que implica que $\mathcal{F}$ es sin torsión, pues $Y$ es íntegro;
de hecho, vamos a demostrar que $\mathrm{H}^1(Y,\mathcal{F})=0$ para \emph{todo} $\mathcal{O}_Y$-módulo coherente \emph{sin torsión} $\mathcal{F}$.
Ahora bien, el homomorfismo $u:\mathcal{A}^m\to\mathcal{B}$ define un homomorfismo
\[
v:\mathcal{G}
=\mathcal{H}om_{\mathcal{A}}(\mathcal{B},\mathcal{F})
\longrightarrow
\mathcal{H}om_{\mathcal{A}}(\mathcal{A}^m,\mathcal{F})
=\mathcal{F}^m.
\]

Mostremos primero que $v$ es \emph{inyectivo}:
por hipótesis, $\mathcal{T}=\operatorname{Coker}u$ tiene un soporte que no encuentra a $V$, y es, por tanto, un $\mathcal{O}_Y$-módulo \emph{de torsión} \sref[I]{I.7.4.6-fr};
la sucesión exacta
\[
\mathcal{A}^m\longrightarrow\mathcal{B}\longrightarrow\mathcal{T}
\longrightarrow0
\]
da, por exactitud por la izquierda del funtor $\mathcal{H}om_{\mathcal{A}}$, la sucesión exacta
\[
0\longrightarrow
\mathcal{H}om_{\mathcal{A}}(\mathcal{T},\mathcal{F})
\longrightarrow\mathcal{G}\xrightarrow{v}\mathcal{F}^m.
\]
Pero, como $\mathcal{F}$ es sin torsión, se tiene
$\mathcal{H}om_{\mathcal{A}}(\mathcal{T},\mathcal{F})=0$
\sref[0]{0.5.2.6-fr}, de donde resulta nuestra afirmación.
Se tiene, pues, la sucesión exacta
\[
0\longrightarrow\mathcal{G}\longrightarrow\mathcal{F}^m
\longrightarrow\operatorname{Coker}v\longrightarrow0
\]
\oldpage[II]{138}
donde $\mathcal{G}$ y $\operatorname{Coker}v$ son $\mathcal{O}_Y$-módulos coherentes (\sref[0]{0.5.3.4-fr} y \sref[0]{0.5.3.5-fr});
en virtud de la sucesión exacta de cohomología, bastaría probar que
$\mathrm{H}^1(Y,\mathcal{G})
=\mathrm{H}^1(Y,\operatorname{Coker}v)=0$
para deducir que
$\mathrm{H}^1(Y,\mathcal{F}^m)
=(\mathrm{H}^1(Y,\mathcal{F}))^m=0$,
y, por tanto, que $\mathrm{H}^1(Y,\mathcal{F})=0$.
Pero la restricción $v|V$ es un isomorfismo, por lo que el soporte de $\operatorname{Coker}v$ es distinto de $Y$, de donde
$\mathrm{H}^1(Y,\operatorname{Coker}v)=0$
por la hipótesis inicial.
Por otra parte, $\mathcal{G}$ es un \emph{$\mathcal{B}$-módulo} coherente \sref[I]{I.9.6.4-fr};
como $X$ es afín sobre $Y$, existe un $\mathcal{O}_X$-módulo cuasicoherente $\mathcal{K}$ tal que $\mathcal{G}$ es isomorfo a $f_*(\mathcal{K})$ \sref{II.1.4.3-fr};
como $X$ es afín, se tiene $\mathrm{H}^1(X,\mathcal{K})=0$ \sref[I]{I.5.1.9.2-fr}, luego también $\mathrm{H}^1(Y,\mathcal{G})=0$ por \sref{II.5.2.3-fr}, lo que termina la demostración del teorema \sref{II.6.7.1-fr}.
\end{proof}
% Fuente francesa: PMIHES_1961__8__5_0, printed p. 138.

\begin{corollary}[6.7.2]
\label{II.6.7.2-fr}
Sean $X$ un preesquema noetheriano y $(X_i)_{1\leq i\leq n}$ una cubierta finita del espacio $X$ formada por conjuntos cerrados.
Para que $X$ sea afín, es necesario y suficiente que, para cada $i$, exista un subpreesquema cerrado de $X$, cuyo espacio subyacente sea $X_i$, que sea afín.
\end{corollary}

\begin{proof}
En efecto, si así ocurre, sea $X'$ el esquema \emph{suma} de los $X_i$;
es claro que $X'$ es afín, y se define un morfismo sobreyectivo $f:X'\to X$ tomando como restricción de $f$ a $X_i$ la inyección canónica.
En virtud de \sref{II.6.7.1-fr}, todo se reduce a verificar que $f$ es \emph{finito}, y esto ya se ha visto en \sref{II.6.1.5-fr}.
\end{proof}

\begin{corollary}[6.7.3]
\label{II.6.7.3-fr}
Para que un preesquema noetheriano $X$ sea afín, es necesario y suficiente que sean afines los subpreesquemas cerrados reducidos cuyos espacios subyacentes son las componentes irreducibles de $X$.
\end{corollary}
% Fuente francesa: PMIHES_1961__8__5_0, printed pp. 138-139.

\section{Criterios valuativos}
\label{section:II.7-fr}

Damos en este parágrafo criterios valuativos de separación y de la propiedad de ser propio para un morfismo, es decir, criterios que hacen intervenir un esquema auxiliar variable de la forma $\operatorname{Spec}(A)$, donde $A$ es un anillo de valuación.
Mediante hipótesis «noetherianas» convenientes, estos criterios pueden restringirse al caso en que $A$ es un anillo de valuación \emph{discreta}.
Este será, sin duda, el único caso que tendremos que aplicar en lo sucesivo, y no hemos introducido los anillos de valuación arbitrarios, en el caso general, sino para establecer el vínculo con los desarrollos clásicos.

\subsection{Recordatorios sobre los anillos de valuación}
\label{subsection:II.7.1-fr}

\begin{env}[7.1.1]
\label{II.7.1.1-fr}
Entre las diversas propiedades equivalentes que caracterizan los anillos de valuación, utilizaremos la siguiente:
se dice que un anillo $A$ es un \emph{anillo de valuación} si es un anillo íntegro que no es un cuerpo y si, en el conjunto de los anillos locales contenidos en el cuerpo de fracciones $K$ de $A$ y distintos de $K$, $A$ es \emph{maximal} para la relación de dominación \sref[I]{I.8.1.1-fr}.
Recordemos que un anillo de valuación es \emph{íntegramente cerrado}.
Si $A$ es un anillo de valuación, $A_{\mathfrak{p}}$ es también un anillo de valuación para todo ideal primo $\mathfrak{p}\neq0$ de $A$.
\end{env}
% Fuente francesa: PMIHES_1961__8__5_0, printed p. 139.

\begin{env}[7.1.2]
\label{II.7.1.2-fr}
Sean $K$ un cuerpo y $A$ un subanillo local de $K$ que no sea un cuerpo;
\oldpage[II]{139}
existe entonces un anillo de valuación que tiene $K$ por cuerpo de fracciones y domina a $A$ (\cite[p.~1-07, lema~2]{I-1}).

Por otra parte, sean $B$ un anillo de valuación, $k$ su cuerpo residual, $K$ su cuerpo de fracciones y $L$ una extensión de $k$.
Existe entonces un anillo de valuación \emph{completo} $C$ que domina a $B$ y cuyo cuerpo residual es $L$.
En efecto, $L$ es extensión algebraica de una extensión trascendente pura
$L'=k(T_\mu)_{\mu\in M}$
se sabe que puede prolongarse la valuación de $K$ correspondiente a $B$ a una valuación de
$K'=K(T_\mu)_{\mu\in M}$
de suerte que $L'$ sea el cuerpo residual de esta valuación (\cite[p.~98]{II-24});
sustituyendo $B$ por el completado del anillo de esta valuación prolongada, se ve que puede limitarse al caso en que $B$ es completo y $L$ es una clausura algebraica de $k$.
Si $\overline{K}$ es una clausura algebraica de $K$, puede prolongarse a $\overline{K}$ la valuación que define $B$, y el cuerpo residual correspondiente es una clausura algebraica de $k$, como se ve levantando a $\overline{K}$ los coeficientes de un polinomio mónico de $k[T]$.
Se queda finalmente reducido al caso $L=k$, y basta entonces tomar para $C$ el completado de $B$ para responder a la cuestión.
\end{env}

\begin{env}[7.1.3]
\label{II.7.1.3-fr}
Sean $K$ un cuerpo y $A$ un subanillo de $K$;
la clausura integral $A'$ de $A$ en $K$ es la intersección de los anillos de valuación que contienen a $A$ y tienen a $K$ por cuerpo de fracciones (\cite[p.~51, th.~2]{I-11}).
La proposición \sref{II.7.1.2-fr} se expresa entonces bajo la forma geométrica equivalente siguiente.
\end{env}
% Fuente francesa: PMIHES_1961__8__5_0, printed p. 139.

\begin{proposition}[7.1.4]
\label{II.7.1.4-fr}
Sean $Y$ un preesquema, $p:X\to Y$ un morfismo, $x$ un punto de $X$, $y=p(x)$ e $y'\neq y$ una especialización \sref[0]{0.2.1.2-fr} de $y$.
Existe entonces un esquema local $Y'$, espectro de un anillo de valuación, y un morfismo separado $f:Y'\to Y$ tales que, si se denota por $a$ el único punto cerrado de $Y'$ y por $b$ el punto genérico de $Y'$, se tenga
$f(a)=y'$ y $f(b)=y$.
Puede suponerse además que se satisface una de las dos propiedades adicionales siguientes:
\begin{enumerate}
  \item[\rm{(i)}] $Y'$ es el espectro de un anillo de valuación completo cuyo cuerpo residual es algebraicamente cerrado, y existe un $\boldsymbol{k}(y)$-homomorfismo $\boldsymbol{k}(x)\to \boldsymbol{k}(b)$.
  \item[\rm{(ii)}] Existe un $\boldsymbol{k}(y)$-isomorfismo $\boldsymbol{k}(x)\xrightarrow{\sim}\boldsymbol{k}(b)$.
\end{enumerate}
\end{proposition}

\begin{proof}
Sea $Y_1$ el subpreesquema cerrado reducido de $Y$ que tiene $\overline{\{y\}}$ por espacio subyacente \sref[I]{I.5.2.1-fr}, y sea $X_1$ el subpreesquema cerrado imagen inversa $p^{-1}(Y_1)$;
como $y'\in\overline{\{y\}}$ por hipótesis y $\boldsymbol{k}(x)$ es el mismo en $X$ y en $X_1$, puede suponerse que $Y$ es \emph{íntegro}, de punto genérico $y$.
$\mathcal{O}_{y'}$ es entonces un anillo local íntegro que no es un cuerpo, cuyo cuerpo de fracciones es $\mathcal{O}_y=\boldsymbol{k}(y)$, y $\boldsymbol{k}(x)$ es una extensión de $\boldsymbol{k}(y)$.

Para realizar las condiciones $f(a)=y'$ y $f(b)=y$, así como la condición adicional (i) (resp. (ii)), se toma $Y'=\operatorname{Spec}(A')$, donde $A'$ es un anillo de valuación que domina a $\mathcal{O}_{y'}$ y que es completo y tiene un cuerpo residual algebraicamente cerrado, extensión de $\boldsymbol{k}(x)$ (resp. un anillo de valuación $A'$ que domina a $\mathcal{O}_{y'}$ y cuyo cuerpo de fracciones es $\boldsymbol{k}(x)$); la existencia de $A'$ está garantizada por (\hyperref[II.7.1.2-fr]{7.1.2}).

\end{proof}
% Fuente francesa: PMIHES_1961__8__5_0, printed pp. 139-140.

\begin{env}[7.1.5]
\label{II.7.1.5-fr}
Recordemos que un anillo local $A$ se dice \emph{de dimensión $1$} si existe un ideal primo distinto del ideal maximal $\mathfrak{m}$ y si todo ideal primo de $A$ distinto de $\mathfrak{m}$ es un ideal primo \emph{minimal};
cuando $A$ es \emph{íntegro}, esto equivale a exigir que $\mathfrak{m}$ y $(0)$ sean los únicos ideales primos y que $\mathfrak{m}\neq(0)$;
en otros términos, $Y=\operatorname{Spec}(A)$ consta de dos
\oldpage[II]{140}
puntos $a$ y $b$: $a$ es el único punto \emph{cerrado}, se tiene $\mathfrak{j}_a=\mathfrak{m}$, y $\boldsymbol{k}(a)=k$ es el \emph{cuerpo residual} $k=A/\mathfrak{m}$;
$b$ es el \emph{punto genérico} de $Y$, $\mathfrak{j}_b=(0)$, el conjunto $\{b\}$ es el único abierto de $Y$ distinto de $\varnothing$ y de $Y$ (abierto que, por tanto, es \emph{denso en todas partes}), y $\boldsymbol{k}(b)=K$ es el \emph{cuerpo de fracciones} de $A$.
\end{env}

\begin{env}[7.1.6]
\label{II.7.1.6-fr}
Para un anillo local $A$, noetheriano y de dimensión $1$, se sabe (\cite[pp.~2-08 y 17-01]{I-1}) que las condiciones siguientes son equivalentes:
\begin{enumerate}
  \item[\rm{(a)}] $A$ es normal;
  \item[\rm{(b)}] $A$ es regular;
  \item[\rm{(c)}] $A$ es un anillo de valuación;
\end{enumerate}
además, $A$ es entonces un \emph{anillo de valuación discreta}.
Las Proposiciones~\sref{II.7.1.2-fr} y \sref{II.7.1.3-fr} tienen entonces los análogos siguientes para los anillos de valuación discreta:
\end{env}

% Fuente francesa: PMIHES_1961__8__5_0, printed p. 140.

\begin{proposition}[7.1.7]
\label{II.7.1.7-fr}
Sean $A$ un anillo local íntegro noetheriano que no sea un cuerpo, $K$ su cuerpo de fracciones y $L$ una extensión de tipo finito de $K$;
existe entonces un anillo de valuación discreta que tiene a $L$ por cuerpo de fracciones y domina a $A$.
\end{proposition}

\begin{proof}
Supongamos primero que $L=K$.
Sean $\mathfrak{m}$ el ideal maximal de $A$, $(x_1,\ldots,x_n)$ un sistema de generadores no nulos de $\mathfrak{m}$ y $B$ el subanillo
$A[x_2/x_1,\ldots,x_n/x_1]$
de $K$, que es noetheriano.
Es inmediato que el ideal $\mathfrak{m}B$ de $B$ es idéntico al ideal principal $x_1B$;
si $\mathfrak{p}$ es un ideal primo minimal por encima de $x_1B$, entonces $\mathfrak{p}$ tiene rango $1$ (\cite[t.~I, p.~277]{I-13});
en otros términos, $B_{\mathfrak{p}}$ es un anillo local noetheriano \emph{de dimensión $1$};
es claro que $\mathfrak{p}B_{\mathfrak{p}}\cap A$ es un ideal de $A$ que contiene a $\mathfrak{m}$ y no contiene a $1$, luego es igual a $\mathfrak{m}$ y, por consiguiente, $B_{\mathfrak{p}}$ \emph{domina} a $A$ \sref[I]{I.8.1.1-fr}.
Del teorema de Krull--Akizuki (\cite[p.~293]{II-25}) resulta que la clausura integral $C$ de $B_{\mathfrak{p}}$ es un anillo noetheriano (aunque $C$ no sea necesariamente un $B_{\mathfrak{p}}$-módulo de tipo finito);
si $\mathfrak{n}$ es un ideal maximal de $C$, entonces $C_{\mathfrak{n}}$ es un anillo local noetheriano normal de dimensión $1$ (\cite[p.~295]{II-25}), y por tanto un anillo de valuación discreta que domina a $B_{\mathfrak{p}}$ y, \emph{a fortiori}, a $A$.

Si ahora $L$ es una extensión de tipo finito de $K$, por lo precedente puede reducirse al caso en que $A$ es ya un anillo de valuación discreta.
Sea $w$ una valuación de $K$ asociada a $A$;
existe una valuación discreta $w'$ de $L$ que \emph{prolonga} a $w$: por inducción sobre el número de generadores de $L$, puede reducirse al caso en que $L=K(\alpha)$, y entonces la proposición es clásica (\cite[p.~106]{II-24}).
\end{proof}

% Fuente francesa: PMIHES_1961__8__5_0, printed p. 140.

\begin{corollary}[7.1.8]
\label{II.7.1.8-fr}
Sean $A$ un anillo íntegro noetheriano, $K$ su cuerpo de fracciones y $L$ una extensión de tipo finito de $K$.
Entonces la clausura integral de $A$ en $L$ es la intersección de los anillos de valuación discreta que tienen a $L$ por cuerpo de fracciones y contienen a $A$.
\end{corollary}

\begin{proof}
En efecto, tal anillo de valuación discreta, al ser normal, contiene \emph{a fortiori} todo elemento de $L$ entero sobre $A$.
Basta, pues, demostrar que, si $x\in L$ no es entero sobre $A$, existe un anillo de valuación discreta $C$ que tiene a $L$ por cuerpo de fracciones, contiene a $A$ y no contiene a $x$.
La hipótesis sobre $x$ significa que $x\notin B=A[1/x]$ o, en otros términos, que $1/x$ no es invertible en el anillo noetheriano $B$.
Existe por tanto un ideal primo $\mathfrak{p}$ de $B$ que contiene a $1/x$.
El anillo local íntegro $B_{\mathfrak{p}}$ es noetheriano y está contenido en $L$, que es una extensión de tipo finito del cuerpo de fracciones de $B_{\mathfrak{p}}$ (el cual contiene a $K$).
En virtud de \sref{II.7.1.7-fr}, existe entonces un anillo de valuación discreta $C$ que domina a $B_{\mathfrak{p}}$ y tiene a $L$ por cuerpo de fracciones;
como $1/x\in\mathfrak{p}B_{\mathfrak{p}}$ pertenece al ideal maximal de $C$, se tiene $x\notin C$, lo que concluye la demostración.
\end{proof}

La forma geométrica de \sref{II.7.1.7-fr} es la siguiente:

% Fuente francesa: PMIHES_1961__8__5_0, printed p. 141.

\oldpage[II]{141}
\begin{proposition}[7.1.9]
\label{II.7.1.9-fr}
Sean $Y$ un preesquema localmente noetheriano, $p:X\to Y$ un morfismo localmente de tipo finito, $x$ un punto de $X$, $y=p(x)$ e $y'\neq y$ una especialización de $y$.
Existe entonces un esquema local $Y'$, espectro de un anillo de valuación discreta, un morfismo separado $f:Y'\to Y$ y una $Y$-aplicación racional $g$ de $Y'$ en $X$, tales que, si $a$ designa el punto cerrado de $Y'$ y $b$ su punto genérico, se tiene $f(a)=y'$, $f(b)=y$, $g(b)=x$, y tales que, en el diagrama conmutativo
\[
\tag{7.1.9.1}
\label{II.7.1.9.1-fr}
\begin{tikzcd}[column sep=large,row sep=large]
& \boldsymbol{k}(x) \arrow[dl,"\gamma"'] \\
\boldsymbol{k}(b) & \boldsymbol{k}(y) \arrow[u,"\pi"'] \arrow[l,"\varphi"]
\end{tikzcd}
\]
(donde $\pi$, $\varphi$ y $\gamma$ son los homomorfismos correspondientes, respectivamente, a $p$, $f$ y $g$), el homomorfismo $\gamma$ es biyectivo.
\end{proposition}

\begin{proof}
Como en \sref{II.7.1.4-fr}, puede reducirse al caso en que $Y$ es íntegro de punto genérico $y$ (teniendo en cuenta \sref[I]{I.6.3.4-fr}[iv]) y, como la cuestión es local sobre $X$ y $Y$, puede suponerse que $p$ es de tipo finito;
se está entonces en la situación de \sref{II.7.1.4-fr}, con la propiedad adicional de que $\boldsymbol{k}(x)$ es una extensión \emph{de tipo finito} de $\boldsymbol{k}(y)$ \sref[I]{I.6.4.11-fr} y de que $\mathcal{O}_{y'}$ es noetheriano;
esto permite aplicar \sref{II.7.1.7-fr} y tomar $Y'=\operatorname{Spec}(A')$, donde $A'$ es un anillo de valuación discreta que domina a $\mathcal{O}_{y'}$ y cuyo cuerpo de fracciones es $\boldsymbol{k}(x)$.
Así queda definido un diagrama conmutativo \sref{II.7.1.9.1-fr} en el que $\gamma$ es biyectivo, y $\pi$ y $\varphi$ corresponden a los morfismos $p$ y $f$.
Además, como $X$ e $Y$ son localmente noetherianos \sref[I]{I.6.6.2-fr} e $Y'$ es íntegro, existe una única $Y$-aplicación racional $g$ de $Y'$ en $X$ a la que corresponde el isomorfismo $\gamma$ \sref[I]{I.7.1.15-fr}, lo que termina la demostración.
\end{proof}

% Fuente francesa: PMIHES_1961__8__5_0, printed pp. 141-142.

\subsection{Criterio valuativo de separación}
\label{subsection:II.7.2-fr}

\begin{proposition}[7.2.1]
\label{II.7.2.1-fr}
Sean $X$ e $Y$ dos preesquemas y $f:X\to Y$ un morfismo cuasicompacto.
Las dos condiciones siguientes son equivalentes:
\begin{enumerate}
  \item[\rm{(a)}] El morfismo $f$ es cerrado.
  \item[\rm{(b)}] Para todo $x\in X$ y toda especialización $y'$ de $y=f(x)$, distinta de $y$, existe una especialización $x'$ de $x$ tal que $f(x')=y'$.
\end{enumerate}
\end{proposition}

\begin{proof}
La condición~\emph{(b)} expresa que
$f(\overline{\{x\}})=\overline{\{y\}}$
y es, por tanto, consecuencia de~\emph{(a)}.
Para demostrar que \emph{(b)} implica \emph{(a)}, consideremos una parte cerrada $X'$ del espacio subyacente $X$;
sea $Y'=\overline{f(X')}$, y demostremos que $Y'=f(X')$.
Consideremos los subpreesquemas cerrados reducidos de $X$ e $Y$ que tienen, respectivamente, a $X'$ e $Y'$ por espacios subyacentes \sref[I]{I.5.2.1-fr};
existe entonces un morfismo $f':X'\to Y'$ tal que el diagrama
\[
\begin{tikzcd}[column sep=large,row sep=large]
X' \arrow[r,"f'"] \arrow[d] & Y' \arrow[d] \\
X \arrow[r,"f"'] & Y
\end{tikzcd}
\]
es conmutativo \sref[I]{I.5.2.2-fr}, y, como $f$ es cuasicompacto, también lo es $f'$.
Se queda así reducido a demostrar que, si $f$ es un morfismo cuasicompacto y \emph{dominante}, entonces
\oldpage[II]{142}
la condición~\emph{(b)} implica que $f(X)=Y$.
Ahora bien, sea $y'$ un punto de $Y$ y sea $y$ el punto genérico de una componente irreducible de $Y$ que contiene a $y'$;
en virtud de~\emph{(b)}, basta demostrar que $f^{-1}(y)$ no es vacío.
Pero se sabe que esta propiedad es consecuencia de que $f$ es cuasicompacto y dominante \sref[I]{I.6.6.5-fr}.
\end{proof}

\begin{corollary}[7.2.2]
\label{II.7.2.2-fr}
Sea $f:X\to Y$ una inmersión cuasicompacta.
Para que el espacio subyacente $X$ sea cerrado en $Y$, es necesario y suficiente que contenga toda especialización, en $Y$, de cada uno de sus puntos.
\end{corollary}

% Fuente francesa: PMIHES_1961__8__5_0, printed p. 142.

\begin{proposition}[7.2.3]
\label{II.7.2.3-fr}
Sean $Y$ un preesquema \textup{(resp.} un preesquema localmente noetheriano\textup{)} y $f:X\to Y$ un morfismo \textup{(resp.} un morfismo localmente de tipo finito\textup{)}.
Las condiciones siguientes son equivalentes:
\begin{enumerate}
  \item[\rm{(a)}] $f$ es separado.
  \item[\rm{(b)}] El morfismo diagonal $X\to X\times_YX$ es cuasicompacto y, para todo $Y$-preesquema de la forma $Y'=\operatorname{Spec}(A)$, donde $A$ es un anillo de valuación \textup{(resp.} un anillo de valuación discreta\textup{)}, dos $Y$-morfismos de $Y'$ en $X$ que coinciden en el punto genérico de $Y'$ son iguales.
  \item[\rm{(c)}] El morfismo diagonal $X\to X\times_YX$ es cuasicompacto y, para todo $Y$-preesquema de la forma $Y'=\operatorname{Spec}(A)$, donde $A$ es un anillo de valuación \textup{(resp.} un anillo de valuación discreta\textup{)}, dos $Y'$-secciones de $X'=X_{(Y')}$ que coinciden en el punto genérico de $Y'$ son iguales.
\end{enumerate}
\end{proposition}

\begin{proof}
La equivalencia entre \emph{(b)} y \emph{(c)} resulta de la correspondencia biyectiva entre los $Y$-morfismos de $Y'$ en $X$ y las $Y'$-secciones de $X'$ \sref[I]{I.3.3.14-fr}.
Si $X$ es separado sobre $Y$, la condición~\emph{(b)} se satisface en virtud de \sref[I]{I.7.2.2.1-fr}, puesto que $Y'$ es íntegro.
Queda por demostrar que \emph{(b)} implica que el morfismo diagonal
$\Delta:X\to X\times_YX$
es cerrado, y equivale a demostrar que satisface el criterio de \sref{II.7.2.2-fr}.
Sea, pues, $z$ un punto de la diagonal $\Delta(X)$ y $z'\neq z$ una especialización de $z$ en $X\times_YX$.
Existe entonces, por \sref{II.7.1.4-fr}, un anillo de valuación $A$ y un morfismo $f$ de $Y'=\operatorname{Spec}(A)$ en $X\times_YX$ que aplica el punto cerrado $a$ de $Y'$ en $z'$ y el punto genérico $b$ de $Y'$ en $z$;
este morfismo hace de $Y'$ un $(X\times_YX)$-preesquema y, \emph{a fortiori}, un $Y$-preesquema.
Componiendo con $f$ las dos proyecciones de $X\times_YX$, se obtienen dos $Y$-morfismos $g_1$, $g_2$ de $Y'$ en $X$ que, por hipótesis, coinciden en el punto $b$;
son, por tanto, iguales a un mismo morfismo $g$, lo que significa \sref[I]{I.5.3.1-fr} que $f$ se factoriza en $f=\Delta\circ g$ y, por consiguiente, $z'\in\Delta(X)$.
Cuando se supone que $Y$ es localmente noetheriano y $f$ localmente de tipo finito, $X\times_YX$ es localmente noetheriano \sref[I]{I.6.6.7-fr};
puede entonces repetirse el razonamiento precedente suponiendo que $A$ es un anillo de valuación discreta, en virtud de \sref{II.7.1.9-fr}.
\end{proof}

\begin{remark}[7.2.4]
\label{II.7.2.4-fr}
\begin{enumerate}
  \item[\rm{(i)}] La hipótesis de que el morfismo $\Delta$ sea cuasicompacto se satisface siempre que $Y$ sea localmente noetheriano y $f$ localmente de tipo finito, pues $X\times_YX$ es entonces localmente noetheriano \sref[I]{I.6.6.4-fr}[i].
  En el caso general, significa también que, para toda cubierta $(U_\alpha)$ de $X$ por abiertos afines, los conjuntos $U_\alpha\cap U_\beta$ son \emph{cuasicompactos}.
  \item[\rm{(ii)}] Para que $f$ sea separado, \emph{basta} exigir la condición~\emph{(b)} o la condición~\emph{(c)} solo para anillos de valuación $A$ \emph{completos} cuyos cuerpos residuales sean \emph{algebraicamente cerrados};
  esto resulta, en efecto, de la demostración de \sref{II.7.2.3-fr} y de \sref{II.7.1.4-fr}.
\end{enumerate}
\end{remark}

% Fuente francesa: PMIHES_1961__8__5_0, printed p. 143.

\oldpage[II]{143}

\subsection{Criterio valuativo de la propiedad de ser propio}
\label{subsection:II.7.3-fr}

\begin{proposition}[7.3.1]
\label{II.7.3.1-fr}
Sean $A$ un anillo de valuación, $Y=\operatorname{Spec}(A)$, $b$ el punto genérico de $Y$, $X$ un \emph{esquema} íntegro y $f:X\to Y$ un morfismo \emph{cerrado} tal que $f^{-1}(b)$ se reduzca a un único punto $x$ y que el homomorfismo correspondiente $\boldsymbol{k}(b)\to \boldsymbol{k}(x)$ sea biyectivo.
Entonces $f$ es un isomorfismo.
\end{proposition}

\begin{proof}
Como $f$ es cerrado y dominante, se tiene $f(X)=Y$;
basta \sref[I]{I.4.2.2-fr} demostrar que, para todo $y'\neq b$ en $Y$, existe \emph{un único} punto $x'$ tal que $f(x')=y'$ y que el homomorfismo correspondiente $\mathcal{O}_{y'}\to\mathcal{O}_{x'}$ es biyectivo, pues entonces $f$ será un homeomorfismo.
Ahora bien, si $f(x')=y'$, entonces $\mathcal{O}_{x'}$ es un anillo local contenido en
$K=\boldsymbol{k}(x)=\boldsymbol{k}(b)$
y que domina a $\mathcal{O}_{y'}$;
este último es el anillo local $A_{y'}$ y, por tanto, es un anillo de valuación \sref{II.7.1.1-fr} que tiene a $K$ por cuerpo de fracciones.
Por otra parte, $\mathcal{O}_{x'}\neq K$, puesto que $x'$ no es el punto genérico de $X$ \sref[0]{0.2.1.3-fr};
se concluye entonces que $\mathcal{O}_{x'}=\mathcal{O}_{y'}$.
Como $X$ es un esquema íntegro, la igualdad $\mathcal{O}_{x'}=\mathcal{O}_{x''}$ implica $x'=x''$ \sref[I]{I.8.2.2-fr}, lo que termina la demostración.
\end{proof}

\begin{env}[7.3.2]
\label{II.7.3.2-fr}
Sean $A$ un anillo de valuación, $Y=\operatorname{Spec}(A)$ y $b$ el punto genérico de $Y$, de modo que $\mathcal{O}_b=\boldsymbol{k}(b)$ es igual a $K$, cuerpo de fracciones de $A$;
sea $f:X\to Y$ un morfismo.
Se sabe \sref[I]{I.7.1.4-fr} que las \emph{$Y$-secciones racionales} de $X$ están en correspondencia biyectiva con los \emph{gérmenes} de $Y$-secciones (definidas en vecindades de $b$) en el punto $b$, de donde se obtiene una aplicación canónica
\[
\tag{7.3.2.1}
\label{II.7.3.2.1-fr}
\Gamma_{\mathrm{rat}}(X/Y)\longrightarrow
\Gamma(f^{-1}(b)/\operatorname{Spec}(K))
\]
identificándose por definición \sref[I]{I.3.4.5-fr} los elementos de $\Gamma(f^{-1}(b)/\operatorname{Spec}(K))$ con los \emph{puntos de $f^{-1}(b)=X\otimes_AK$ racionales sobre $K$}.
Cuando $f$ es \emph{separado}, de \sref[I]{I.5.4.7-fr} resulta que la aplicación \sref{II.7.3.2.1-fr} es \emph{inyectiva}, puesto que $Y$ es un esquema íntegro.

Componiendo \sref{II.7.3.2.1-fr} con la aplicación canónica
$\Gamma(X/Y)\to\Gamma_{\mathrm{rat}}(X/Y)$
de \sref[I]{I.7.1.2-fr}, se obtiene una aplicación canónica
\[
\tag{7.3.2.2}
\label{II.7.3.2.2-fr}
\Gamma(X/Y)\longrightarrow \Gamma(f^{-1}(b)/\operatorname{Spec}(K)).
\]
Cuando $f$ es \emph{separado}, esta aplicación es también \emph{inyectiva} \sref[I]{I.5.4.7-fr}.
\end{env}

% Fuente francesa: PMIHES_1961__8__5_0, printed pp. 143-144.

\begin{proposition}[7.3.3]
\label{II.7.3.3-fr}
Sean $A$ un anillo de valuación de cuerpo de fracciones $K$, $Y=\operatorname{Spec}(A)$, $b$ el punto genérico de $Y$ y $f:X\to Y$ un morfismo \emph{separado} y \emph{cerrado}.
Entonces la aplicación canónica \textup{\sref{II.7.3.2.2-fr}} es biyectiva \textup{(lo que equivale a decir que es \emph{sobreyectiva} e implica que las $Y$-secciones racionales de $X$ están \emph{definidas en todas partes})}.
\end{proposition}

\begin{proof}
Sea, pues, $x$ un punto de $f^{-1}(b)$ \emph{racional} sobre $K$.
Como $f$ es separado, también lo es el morfismo $f^{-1}(b)\to\operatorname{Spec}(K)$ correspondiente a $f$ \sref[I]{I.5.5.1-fr}[iv], y, puesto que toda sección de $f^{-1}(b)$ es una inmersión cerrada \sref[I]{I.5.4.6-fr}, $\{x\}$ es cerrado \emph{en $f^{-1}(b)$}.
Consideremos el subpreesquema cerrado reducido $X'$ de $X$ que tiene por espacio subyacente la clausura $\overline{\{x\}}$ de $\{x\}$ \emph{en $X$}.
Es claro que la restricción de $f$ a $X'$ satisface las hipótesis de \sref{II.7.3.1-fr} y es, por tanto, un \emph{isomorfismo} de $X'$ sobre $Y$, cuyo isomorfismo inverso es la $Y$-sección buscada de $X$.
\end{proof}

\begin{env}[7.3.4]
\label{II.7.3.4-fr}
Para enunciar los dos resultados siguientes, utilizaremos una terminología que se justificará y discutirá en el capítulo~IV:
si $F$ es una parte
\oldpage[II]{144}
de un preesquema $Y$, se llama \emph{codimensión} de $F$ en $Y$, y se denota por $\operatorname{codim}_YF$, al ínfimo de los enteros $\dim(\mathcal{O}_z)$ cuando $z$ recorre $F$.
\end{env}

\begin{corollary}[7.3.5]
\label{II.7.3.5-fr}
Sean $Y$ un preesquema localmente noetheriano reducido y $N$ el conjunto de los puntos $y\in Y$ en los que $Y$ no es regular \textup{\sref[0]{0.4.1.4-fr}};
se supone que $\operatorname{codim}_YN\geq2$.
Sea $f:X\to Y$ un morfismo de tipo finito, \emph{separado} y \emph{cerrado}, y sea $g$ una $Y$-sección racional de $X$;
si $Y'$ es el conjunto de los puntos de $Y$ en los que $g$ no está definida \textup{(conjunto que es \emph{cerrado} \sref[I]{I.7.2.1-fr})}, entonces $\operatorname{codim}_YY'\geq2$.
\end{corollary}

\begin{proof}
Basta demostrar que $g$ está definida en todo punto $z\in Y$ tal que $\dim\mathcal{O}_z\leq1$.
Si $\dim\mathcal{O}_z=0$, entonces $z$ es el punto genérico de una componente irreducible de $Y$ \sref[I]{I.1.1.14-fr}, por lo que pertenece a todo abierto denso en todas partes de $Y$ y, en particular, al dominio de definición de $g$.
Supongamos, pues, que $\dim\mathcal{O}_z=1$;
por hipótesis, $\mathcal{O}_z$ es entonces un anillo local noetheriano regular y, por \sref{II.7.1.6-fr}, un \emph{anillo de valuación discreta}.
Sea $Z=\operatorname{Spec}(\mathcal{O}_z)$;
como $U=Y\setmin Y'$ es denso en todas partes, $U\cap Z$ no es vacío \sref[I]{I.2.4.2-fr};
sea $g'$ la aplicación racional de $Z$ en $X$ inducida por $g$ \sref[I]{I.7.2.8-fr};
basta demostrar que $g'$ es un \emph{morfismo} \sref[I]{I.7.2.9-fr}.
Ahora bien, $g'$ puede considerarse como una $Z$-sección racional del $Z$-preesquema
$f^{-1}(Z)=X\times_YZ$
es claro que el morfismo $f^{-1}(Z)\to Z$ correspondiente a $f$ es cerrado, y de \sref[I]{I.5.5.1-fr}[i] resulta que es separado;
se concluye entonces de \sref{II.7.3.3-fr} que $g'$ está definida en todas partes;
como $Z$ es reducido y $X$ es separado sobre $Y$, $g'$ es un morfismo \sref[I]{I.7.2.2-fr}.
\end{proof}

% Fuente francesa: PMIHES_1961__8__5_0, printed p. 144.

\begin{corollary}[7.3.6]
\label{II.7.3.6-fr}
Sean $S$ un preesquema localmente noetheriano y $X$ e $Y$ dos $S$-preesquemas;
se supone que $Y$ es reducido y, además, que el conjunto $N$ de los puntos $y\in Y$ en los que $Y$ no es regular satisface $\operatorname{codim}_YN\geq2$;
se supone finalmente que el morfismo estructural $X\to S$ es propio.
Sea $f$ una $S$-aplicación racional de $Y$ en $X$, y sea $Y'$ el conjunto de los puntos de $Y$ en los que $f$ no está definida;
entonces $\operatorname{codim}_YY'\geq2$.
\end{corollary}

\begin{proof}
Se sabe \sref[I]{I.7.1.2-fr} que pueden identificarse las $S$-aplicaciones racionales de $Y$ en $X$ con las $Y$-secciones racionales de $X\times_SY$;
como el morfismo estructural $X\times_SY\to Y$ es cerrado \sref{II.5.4.1-fr}, puede aplicarse \sref{II.7.3.5-fr}, de donde resulta el corolario.
\end{proof}

\begin{remark}[7.3.7]
\label{II.7.3.7-fr}
Las hipótesis impuestas a $Y$ en \sref{II.7.3.5-fr} y \sref{II.7.3.6-fr} se satisfacen, en particular, cuando $Y$ es \emph{normal} \sref[0]{0.4.1.4-fr}, en virtud de \sref{II.7.1.6-fr}.
\end{remark}

% Fuente francesa: PMIHES_1961__8__5_0, printed pp. 144-145.

Podemos caracterizar los morfismos universalmente cerrados (resp. propios) mediante un recíproco de \sref{II.7.3.3-fr}:

\begin{theorem}[7.3.8]
\label{II.7.3.8-fr}
Sean $Y$ un preesquema \textup{(resp.} un preesquema localmente noetheriano\textup{)} y $f:X\to Y$ un morfismo cuasicompacto separado \textup{(resp.} un morfismo de tipo finito\textup{)}.
Las condiciones siguientes son equivalentes:
\begin{enumerate}
  \item[\rm{(a)}] $f$ es universalmente cerrado \textup{(resp.} propio\textup{)}.
  \item[\rm{(b)}] Para todo $Y$-esquema de la forma $Y'=\operatorname{Spec}(A)$, donde $A$ es un anillo de valuación \textup{(resp.} un anillo de valuación discreta) de cuerpo de fracciones $K$, la aplicación canónica
  \[
\operatorname{Hom}_Y(Y',X)\longrightarrow
\operatorname{Hom}_Y(\operatorname{Spec}(K),X)
\]
  correspondiente a la inyección canónica $A\to K$ es sobreyectiva \textup{(resp.} biyectiva\textup{)}.
\oldpage[II]{145}
  \item[\rm{(c)}] Para todo $Y$-esquema de la forma $Y'=\operatorname{Spec}(A)$, donde $A$ es un anillo de valuación \textup{(resp.} un anillo de valuación discreta), la aplicación canónica \textup{\sref{II.7.3.2.2-fr}} relativa al $Y'$-preesquema $X_{(Y')}$ es sobreyectiva \textup{(resp.} biyectiva).
\end{enumerate}
\end{theorem}

\begin{proof}
La equivalencia entre \emph{(b)} y \emph{(c)} resulta inmediatamente de \sref[I]{I.3.3.14-fr};
\emph{(a)} implica \emph{(c)}, pues, en uno u otro caso, \emph{(a)} implica que $f_{(Y')}$ es separado \sref[I]{I.5.5.1-fr}[iv] y cerrado, y basta aplicar \sref{II.7.3.3-fr}.
Queda por demostrar que \emph{(b)} implica \emph{(a)}.

Consideremos primero el caso en que $Y$ es arbitrario y $f$ es separado y cuasicompacto.
Si $f$ satisface la condición~\emph{(b)}, también la satisface
$f_{(Y'')}:X_{(Y'')}\to Y''$,
donde $Y''$ es un $Y$-preesquema arbitrario, en virtud de la equivalencia entre \emph{(b)} y \emph{(c)} y del hecho de que
$X_{(Y'')}\times_{Y''}Y'=X\times_YY'$
para todo morfismo $Y'\to Y''$ \sref[I]{I.3.3.9.1-fr};
como, además, $f_{(Y'')}$ es separado y cuasicompacto cuando $f$ lo es (\sref[I]{I.5.5.1-fr}[iv] y \sref[I]{I.6.6.4-fr}[iii]), se queda reducido a demostrar que \emph{(b)} implica que $f$ es \emph{cerrado}.
Para ello, basta verificar la condición~\emph{(b)} de \sref{II.7.2.1-fr}.
Sean, pues, $x\in X$ e $y'$ una especialización de $y=f(x)$, distinta de $y$;
en virtud de \sref{II.7.1.4-fr}, existen un esquema $Y'$, espectro de un anillo de valuación, y un morfismo separado $g:Y'\to Y$ tales que, si $a$ designa el punto cerrado y $b$ el punto genérico de $Y'$, se tiene $g(a)=y'$, $g(b)=y$, y existe un $\boldsymbol{k}(y)$-homomorfismo $\boldsymbol{k}(x)\to \boldsymbol{k}(b)$.
Este último corresponde canónicamente a un $Y$-morfismo
$\operatorname{Spec}(\boldsymbol{k}(b))\to X$
\sref[I]{I.2.4.6-fr}, y de \emph{(b)} resulta entonces que existe un $Y$-morfismo $h:Y'\to X$ al que corresponde el morfismo precedente.
Se tiene entonces $h(b)=x$;
si se pone $h(a)=x'$, entonces $x'$ es una especialización de $x$ y
$f(x')=f(h(a))=g(a)=y'$.

Si ahora $Y$ es localmente noetheriano y $f$ de tipo finito, la hipótesis~\emph{(b)} implica primero que $f$ es \emph{separado}, en virtud de \sref{II.7.2.3-fr}, pues el morfismo diagonal $X\to X\times_YX$ es cuasicompacto \sref{II.7.2.4-fr}.
Además, para demostrar que $f$ es propio, basta demostrar que
$f_{(Y'')}:X_{(Y'')}\to Y''$
es \emph{cerrado} para todo $Y$-preesquema $Y''$ \emph{de tipo finito}, teniendo en cuenta \sref{II.5.6.3-fr}.
Como $Y''$ es entonces localmente noetheriano, puede repetirse el razonamiento del primer caso tomando para $Y'$ el espectro de un anillo de valuación discreta y aplicando \sref{II.7.1.9-fr} en lugar de \sref{II.7.1.4-fr}.
\end{proof}

% Fuente francesa: PMIHES_1961__8__5_0, printed pp. 145-146.

\begin{remarks}[7.3.9]
\label{II.7.3.9-fr}
\begin{enumerate}
  \item[\rm{(i)}] Cuando $Y$ es un preesquema arbitrario y $f$ un morfismo separado, para que $f$ sea universalmente cerrado \emph{basta} exigir la condición~\emph{(b)} o la condición~\emph{(c)} solo para anillos de valuación $A$ \emph{completos} cuyos cuerpos residuales sean \emph{algebraicamente cerrados};
  esto resulta, en efecto, de la demostración precedente y de \sref{II.7.1.4-fr}.

  \item[\rm{(ii)}] Del criterio~\emph{(c)} de \sref{II.7.3.8-fr} se deduce una nueva demostración, más próxima a los métodos clásicos, del hecho de que un morfismo proyectivo $X\to Y$ es cerrado \sref{II.5.5.3-fr}.
  Puede suponerse, en efecto, que $Y$ es afín y, por consiguiente, identificar $X$ con un subpreesquema cerrado de un fibrado proyectivo $\mathbf{P}_Y^n$ \sref{II.5.3.3-fr};
  para demostrar que $X\to Y$ es cerrado, basta verificar que lo es el morfismo estructural $\mathbf{P}_Y^n\to Y$, y el criterio~\emph{(c)} de \sref{II.7.3.8-fr}, junto con \sref{II.4.1.3.1-fr}, muestra que se queda reducido a demostrar el hecho siguiente:
  \emph{si $Y$ es el espectro de un anillo de valuación $A$, de cuerpo de fracciones $K$, todo punto de $\mathbf{P}_Y^n$ con valores en $K$ proviene}, por restricción al punto genérico de $Y$, \emph{de un punto de $\mathbf{P}_Y^n$ con valores en $A$.}
  Ahora bien, todo $\mathcal{O}_Y$-módulo invertible es trivial \sref[I]{I.2.4.8-fr};
  de \sref{II.4.2.6-fr} resulta, pues, que un punto de $\mathbf{P}_Y^n$ con valores en $K$ se identifica con una clase de sistemas
  $(\zeta c_0,\zeta c_1,\ldots,\zeta
c_n)$
  de $K$, donde $\zeta\neq0$ y los $c_i$ son elementos no todos nulos de $K$.
  Multiplicando los $c_i$ por un elemento de $A$ de
\oldpage[II]{146}
  valuación conveniente, puede suponerse que todos los $c_i$ pertenecen a $A$ y que al menos uno de ellos es invertible.
  Pero entonces, por \sref{II.4.2.6-fr}, el sistema $(c_0,\ldots,c_n)$ define también un punto de $\mathbf{P}_Y^n$ con valores en $A$, lo que demuestra la afirmación.

  \item[\rm{(iii)}] Los criterios \sref{II.7.2.3-fr} y \sref{II.7.3.8-fr} son especialmente cómodos cuando se considera el dato de un $Y$-preesquema $X$ como equivalente al dato del funtor
  \[
X(Y')=\operatorname{Hom}_Y(Y',X)
\]
  sobre los $Y$-preesquemas $Y'$;
  estos criterios nos permitirán, por ejemplo, demostrar que, bajo ciertas condiciones, los «esquemas de Picard» son propios.
\end{enumerate}
\end{remarks}

% Fuente francesa: PMIHES_1961__8__5_0, printed pp. 146-147.

\begin{corollary}[7.3.10]
\label{II.7.3.10-fr}
Sean $Y$ un esquema íntegro \textup{(resp.} un esquema íntegro localmente noetheriano), $X$ un esquema íntegro y $f:X\to Y$ un morfismo dominante.
\begin{enumerate}
  \item[\rm{(i)}] Si $f$ es cuasicompacto y universalmente cerrado, todo anillo de valuación cuyo cuerpo de fracciones sea el cuerpo $R(X)$ de las funciones racionales sobre $X$ y que domine a un anillo local de $Y$ domina también a un anillo local de $X$.
  \item[\rm{(ii)}] Recíprocamente, supongamos que $f$ es de tipo finito y que la propiedad enunciada en~\textup{(i)} se verifica para todo anillo de valuación \textup{(resp.} todo anillo de valuación discreta) que tenga a $R(X)$ por cuerpo de fracciones.
  Entonces $f$ es propio.
\end{enumerate}
\end{corollary}

\begin{proof}
Observemos ante todo que las hipótesis implican en todos los casos que $f$ es separado \sref[I]{I.5.5.9-fr}.
\begin{enumerate}
  \item[\rm{(i)}] Sean $K=R(Y)$, $L=R(X)$, $y$ un punto de $Y$ y $A$ un anillo de valuación que tenga a $L$ por cuerpo de fracciones y domine a $\mathcal{O}_y$;
  la inyección $\mathcal{O}_y\to A$ define, por tanto, un morfismo $h$ de $Y'=\operatorname{Spec}(A)$ en $Y$ \sref[I]{I.2.4.4-fr} tal que $h(a)=y$, donde se denota por $a$ el punto cerrado de $Y'$;
  además, si $\eta$ es el punto genérico de $Y$, que es también el de $\operatorname{Spec}(\mathcal{O}_y)$, se tiene $h(b)=\eta$, donde se denota por $b$ el punto genérico de $Y'$ (pues $K\subset L$ por hipótesis).
  Si $\xi$ es el punto genérico de $X$, se tiene por hipótesis $\boldsymbol{k}(\xi)=\boldsymbol{k}(b)=L$, de donde resulta un $Y$-morfismo $g:\operatorname{Spec}(L)\to X$ tal que $g(b)=\xi$;
  en virtud de \sref{II.7.3.8-fr}, $g$ proviene de un $Y$-morfismo $g':Y'\to X$.
  Si $x=g'(a)$, es claro que $A$ domina a $\mathcal{O}_x$.

  \item[\rm{(ii)}] Como la cuestión es local sobre $Y$, puede suponerse que $Y$ es afín (resp. afín y noetheriano).
  Puesto que $f$ es de tipo finito, puede aplicarse en ambos casos el lema de Chow \sref{II.5.6.1-fr}.
  Existen, pues, un morfismo proyectivo $p:P\to Y$, una inmersión $j:X'\to P$ y un morfismo proyectivo, sobreyectivo y birracional $g:X'\to X$, siendo $X'$ íntegro, tales que el diagrama
  \[
\begin{tikzcd}[column sep=large,row sep=large]
P \arrow[d,"p"'] & X' \arrow[l,"j"'] \arrow[d,"g"] \\
Y & X \arrow[l,"f"]
\end{tikzcd}
\]
  es conmutativo.
  Basta probar que $j$ es una inmersión \emph{cerrada}, pues entonces $f\circ g=p\circ j$ será un morfismo proyectivo, y por tanto propio, y, como $g$ es sobreyectivo, $f$ será también propio \sref{II.5.4.3-fr}.
  Sea $Z$ el subpreesquema cerrado reducido de $P$ cuyo espacio subyacente es $\overline{j(X')}$ \sref[I]{I.5.2.1-fr};
  como $X'$ es íntegro, $j$ se factoriza en $i\circ h$, donde $i:Z\to P$ es la inyección canónica, $h:X'\to Z$ una inmersión abierta dominante \sref[I]{I.5.2.3-fr}, y $Z$ es íntegro;
\oldpage[II]{147}
  además, $Z$ es proyectivo sobre $Y$, y se ve que puede uno limitarse al caso en que $P$ es \emph{íntegro} y $j$ \emph{dominante} y \emph{birracional}, reduciéndose todo a demostrar que $j$ es \emph{sobreyectivo}.
  Sea, en efecto, $z\in P$;
  $\mathcal{O}_z$ es un anillo local íntegro (resp. íntegro y noetheriano) cuyo cuerpo de fracciones es
  \[
L=R(P)=R(X')=R(X).
\]
  Puede uno limitarse al caso en que $z$ no es el punto genérico de $P$.
  Existe entonces, por \sref{II.7.1.2-fr} y \sref{II.7.1.7-fr}, un anillo de valuación (resp. un anillo de valuación discreta) $A$ que tiene a $L$ por cuerpo de fracciones y domina a $\mathcal{O}_z$.
  \emph{A fortiori}, $A$ domina a $\mathcal{O}_y$, donde $y=p(z)$, y, por hipótesis, existe entonces un $x\in X$ tal que $A$ domina a $\mathcal{O}_x$.
  Como $g$ es propio, la primera parte de la demostración prueba que $A$ domina también a $\mathcal{O}_{x'}$ para cierto $x'\in X'$;
  se sigue que $\mathcal{O}_z$ y $\mathcal{O}_{j(x')}=\mathcal{O}_{x'}$ están emparentados \sref[I]{I.8.1.4-fr} y, como $P$ es un esquema, esto implica $z=j(x')$ \sref[I]{I.8.2.2-fr} y termina la demostración.
\end{enumerate}
\end{proof}
% Fuente francesa: PMIHES_1961__8__5_0, printed p. 147 (PDF physical page 144).

\begin{corollary}[7.3.11]
\label{II.7.3.11-fr}
Sean $X$ e $Y$ dos esquemas íntegros, y sea $f:X\to Y$ un morfismo dominante, cuasicompacto y universalmente cerrado.
Supongamos además que $Y$ es afín, de anillo (íntegro) $B$.
Entonces $\Gamma(X,\mathcal{O}_X)$ es canónicamente isomorfo a un subanillo de la clausura integral de $B$ en $R(X)$.
\end{corollary}

\begin{proof}
En efecto, por \sref[I]{I.8.2.1.1-fr}, $\Gamma(X,\mathcal{O}_X)$ se identifica con la intersección de los $\mathcal{O}_x$ para $x\in X$;
en virtud de \sref{II.7.3.10-fr}, \sref{II.7.1.2-fr} y \sref{II.7.1.3-fr}, $\Gamma(X,\mathcal{O}_X)$ queda, por tanto, contenido en la intersección de los anillos de valuación que contienen a $B$ y tienen a $R(X)$ por cuerpo de fracciones;
la conclusión se sigue entonces de \sref{II.7.1.3-fr}.
\end{proof}
% Fuente francesa: PMIHES_1961__8__5_0, printed p. 147 (PDF physical page 144).

\begin{remarks}[7.3.12]
\label{II.7.3.12-fr}
Bajo las hipótesis de \sref{II.7.3.11-fr}, y suponiendo que $R(X)$ es una extensión de tipo finito de $R(Y)$, se podrá, en muchos casos, concluir que $\Gamma(X,\mathcal{O}_X)$ es un \emph{módulo de tipo finito} sobre el anillo $B=\Gamma(Y,\mathcal{O}_Y)$.
Será este el caso, por ejemplo, cuando $B$ sea un \emph{álgebra de tipo finito sobre un cuerpo}, pues se sabe entonces que la clausura integral de $B$ en una extensión de tipo finito de su cuerpo de fracciones es un $B$-módulo de tipo finito (\cite[t.~I, p.~267, th.~9]{I-13});
la conclusión resulta entonces de \sref{II.7.3.11-fr} y de que $B$ es noetheriano.

En particular, \emph{un esquema $X$ propio y afín sobre un cuerpo $K$ es finito}.
En efecto, en virtud de \sref{II.1.6.4-fr}, \sref{II.5.4.6-fr} y \sref[I]{I.6.4.4-fr}[\emph{(c)}], puede uno limitarse al caso en que $X$ es \emph{reducido}.
Además, basta probar que cada uno de los subpreesquemas cerrados de $X$ cuyo espacio subyacente es una componente irreducible de $X$ (que son en número finito) es finito sobre $K$, con lo cual, habida cuenta de \sref{II.5.4.5-fr}, se reduce uno finalmente al caso en que $X$ es \emph{íntegro}.
Pero entonces el resultado se sigue de las observaciones precedentes.

En el capítulo~III volveremos a demostrar la proposición precedente por otros métodos y como consecuencia de resultados más generales, mostrando que, si $f:X\to Y$ es propio e $Y$ es localmente noetheriano, $f_*(\mathcal{F})$ es coherente para todo $\mathcal{O}_X$-módulo coherente $\mathcal{F}$ \sref[III]{III.4.4.2-fr}.

Señalemos finalmente que el criterio~\sref{II.7.3.10-fr} se toma como \emph{definición} de los morfismos propios en la geometría algebraica clásica.
Solo lo hemos mencionado por esta razón, pues el criterio~\sref{II.7.3.8-fr} parece más manejable en todas las aplicaciones que conocemos.
\end{remarks}
% Fuente francesa: PMIHES_1961__8__5_0, printed p. 148 (PDF physical page 145).

\oldpage[II]{148}

\subsection{Curvas algebraicas y cuerpos de funciones de dimensión 1}
\label{subsection:II.7.4-fr}

El objeto de este número es mostrar cómo se formula, en el lenguaje de los esquemas, la noción clásica de curva algebraica (tal como se expone, por ejemplo, en el libro de C.~Chevalley \cite{II-23}).
En todo este número, \emph{se designa por $k$ un cuerpo; todos los esquemas considerados son $k$-esquemas de tipo finito, y todos los morfismos son $k$-morfismos}.

\begin{proposition}[7.4.1]
\label{II.7.4.1-fr}
Sea $X$ un preesquema de tipo finito sobre $k$ (por tanto, noetheriano);
sean $x_i$ ($(1\leq i\leq n)$) los puntos genéricos de las componentes irreducibles $X_i$ de $X$, y sea $K_i=\boldsymbol{k}(x_i)$ ($(1\leq i\leq n)$).
Las condiciones siguientes son equivalentes:
\begin{enumerate}
    \item[\rm{(a)}] Cada uno de los $K_i$ es una extensión de $k$ cuyo grado de trascendencia es igual a 1.
    \item[\rm{(b)}] Para todo punto cerrado $x$ de $X$, el anillo local $\mathcal{O}_x$ es de dimensión 1 \textup{\sref{II.7.1.5-fr}}.
    \item[\rm{(c)}] Las partes cerradas irreducibles de $X$ distintas de los $X_i$ son exactamente los puntos cerrados de $X$.
\end{enumerate}
\end{proposition}

\begin{proof}
Como $X$ es cuasicompacto, toda parte cerrada irreducible $F$ de $X$ contiene un punto cerrado \sref[0]{0.2.1.3-fr}.
En virtud de \sref[I]{I.2.4.2-fr}, hay una correspondencia biyectiva entre los ideales primos de $\mathcal{O}_x$ y las partes cerradas irreducibles de $X$ que contienen a $x$ \sref[I]{I.1.1.14-fr};
de ello se deduce inmediatamente la equivalencia entre \emph{(b)} y \emph{(c)}.
Por otra parte, si $\mathfrak{p}_\alpha$ ($(1\leq\alpha\leq r)$) son los ideales primos minimales del anillo local noetheriano $\mathcal{O}_x$, los anillos locales $\mathcal{O}_x/\mathfrak{p}_\alpha$ son íntegros y tienen por cuerpos de fracciones aquellos $K_i$ tales que $x\in X_i$.
Además, se sabe (\cite[p.~4-06, th.~2]{I-1}) que la dimensión de una $k$-álgebra local íntegra de tipo finito es igual al grado de trascendencia sobre $k$ de su cuerpo de fracciones.
Finalmente, la dimensión de $\mathcal{O}_x$ es el supremo de las dimensiones de los $\mathcal{O}_x/\mathfrak{p}_\alpha$;
ahora bien, la condición~\emph{(a)} implica que estas dimensiones son iguales a 1, y por tanto \emph{(a)} implica \emph{(b)};
recíprocamente, si $\mathcal{O}_x$ es de dimensión 1, ninguno de los $\mathfrak{p}_\alpha$ puede ser igual al ideal maximal de $\mathcal{O}_x$, pues de lo contrario $\mathcal{O}_x$ sería de dimensión 0;
así, cada uno de los $\mathcal{O}_x/\mathfrak{p}_\alpha$ es de dimensión 1, lo que demuestra que \emph{(b)} implica \emph{(a)}.
\end{proof}
% Fuente francesa: PMIHES_1961__8__5_0, printed p. 148 (PDF physical page 145).

Se señala que, bajo las condiciones de \sref{II.7.4.1-fr}, el conjunto $X$ es \emph{vacío} o \emph{infinito}, como resulta inmediatamente de \sref[I]{I.6.4.4-fr}.

\begin{definition}[7.4.2]
\label{II.7.4.2-fr}
Se llama \emph{curva algebraica sobre $k$} a un \emph{esquema} algebraico no vacío sobre $k$ que verifica las condiciones de \textup{\sref{II.7.4.1-fr}}.
\end{definition}

En el lenguaje de la dimensión, que será introducido en el capítulo~IV, esto se expresará diciendo que una curva algebraica sobre $k$ es un $k$-esquema algebraico no vacío \emph{cuyas componentes irreducibles son todas de dimensión 1}.

Se señala que, si $X$ es una curva algebraica sobre $k$, los subpreesquemas cerrados reducidos $X_i$ ($(1\leq i\leq n)$) de $X$ cuyos espacios subyacentes son las componentes irreducibles de $X$ son también curvas algebraicas sobre $k$.
% Fuente francesa: PMIHES_1961__8__5_0, printed pp. 148-149 (PDF physical pages 145-146).

\begin{corollary}[7.4.3]
\label{II.7.4.3-fr}
Sea $X$ una curva algebraica irreducible.
El único punto no cerrado de $X$ es su punto genérico.
Las partes cerradas de $X$ distintas de $X$ son los conjuntos finitos de puntos cerrados;
son también las únicas partes de $X$ que no son densas en todas partes.
\end{corollary}

\begin{proof}
Si un punto $x\in X$ no es cerrado, su adherencia en $X$ es una parte cerrada irreducible de $X$, y por tanto es necesariamente todo $X$, en virtud de \sref{II.7.4.1-fr}; así, $x$ es el punto genérico de $X$.
Una parte cerrada $F$ de $X$ distinta de $X$ no puede contener
\oldpage[II]{149}
el punto genérico de $X$, por lo que todos sus puntos son cerrados (en $X$ y, \emph{a fortiori}, en $F$);
considerando el subpreesquema cerrado reducido de $X$ cuyo espacio subyacente es $F$ \sref[I]{I.5.2.1-fr}, se deduce entonces de \sref[I]{I.6.2.2-fr} que $F$ es finito y discreto.
La adherencia en $X$ de toda parte infinita de $X$ es, por tanto, necesariamente igual a $X$.
\end{proof}
% Fuente francesa: PMIHES_1961__8__5_0, printed p. 149 (PDF physical page 146).

Cuando $X$ es una curva algebraica cualquiera, aplicando \sref{II.7.4.3-fr} a las componentes irreducibles de $X$, se ve que los únicos puntos no cerrados de $X$ son los puntos genéricos de esas componentes.

\begin{corollary}[7.4.4]
\label{II.7.4.4-fr}
Sean $X$ e $Y$ dos curvas algebraicas irreducibles sobre $k$, y sea $f:X\to Y$ un $k$-morfismo.
Para que $f$ sea dominante, es necesario y suficiente que $f^{-1}(y)$ sea finito para todo $y\in Y$.
\end{corollary}

\begin{proof}
En efecto, si $f$ no es dominante, $f(X)$ es necesariamente una parte \emph{finita} de $Y$ en virtud de \sref{II.7.4.3-fr}, de modo que no es posible que $f^{-1}(y)$ sea finito para todo punto de $Y$, pues de lo contrario $X$ sería finito, lo cual es absurdo \sref{II.7.4.1-fr}.
Recíprocamente, si $f$ es dominante, para todo $y\in Y$ distinto del punto genérico $\eta$ de $Y$, el conjunto $f^{-1}(y)$ es cerrado en $X$, puesto que $\{y\}$ es cerrado en $Y$ \sref{II.7.4.3-fr};
por otra parte, por hipótesis, $f^{-1}(y)$ no contiene el punto genérico $\xi$ de $X$ y, por tanto, es finito en virtud de \sref{II.7.4.3-fr}.
Finalmente, para ver que, cuando $f$ es dominante, $f^{-1}(\eta)$ es finito, se observa que la fibra $f^{-1}(\eta)$ es un esquema irreducible de tipo finito sobre $\boldsymbol{k}(\eta)$ y tiene por punto genérico a $\xi$ (\sref[I]{I.6.3.9-fr} y \sref[I]{I.6.4.11-fr}).
Como $\boldsymbol{k}(\xi)$ y $\boldsymbol{k}(\eta)$ son extensiones de tipo finito de $k$, ambas de grado de trascendencia 1, $\boldsymbol{k}(\xi)$ es necesariamente una extensión de grado finito de $\boldsymbol{k}(\eta)$; por consiguiente, $\xi$ es cerrado en $f^{-1}(\eta)$ \sref[I]{I.6.4.2-fr}, y $f^{-1}(\eta)$ consta, por tanto, del único punto $\xi$.
\end{proof}
% Fuente francesa: PMIHES_1961__8__5_0, printed p. 149 (PDF physical page 146).

Veremos en el Capítulo~III que un morfismo \emph{propio} $f:X\to Y$ de preesquemas noetherianos, tal que $f^{-1}(y)$ sea finito para todo $y\in Y$, es necesariamente \emph{finito};
se deducirá, por tanto, de \sref{II.7.4.4-fr} que un morfismo dominante propio de una curva algebraica irreducible en una curva algebraica es \emph{finito}.

\begin{corollary}[7.4.5]
\label{II.7.4.5-fr}
Sea $X$ una curva algebraica sobre $k$.
Para que $X$ sea regular, es necesario y suficiente que $X$ sea normal o, equivalentemente, que los anillos locales de sus puntos cerrados sean anillos de valuación discreta.
\end{corollary}

\begin{proof}
Esto resulta inmediatamente de la condición~\emph{(b)} de \sref{II.7.4.1-fr} y de \sref{II.7.1.6-fr}.
\end{proof}
% Fuente francesa: PMIHES_1961__8__5_0, printed p. 149 (PDF physical page 146).

\begin{corollary}[7.4.6]
\label{II.7.4.6-fr}
Sean $X$ una curva algebraica reducida y $\mathcal{A}$ una $\mathcal{R}(X)$-álgebra coherente reducida;
entonces la clausura integral $X'$ de $X$ respecto de $\mathcal{A}$ \textup{\sref{II.6.3.4-fr}} es una curva algebraica normal, y el morfismo canónico $X'\to X$ es finito.
\end{corollary}

\begin{proof}
El hecho de que $X'\to X$ sea finito resulta de \sref{II.6.3.10-fr};
$X'$ es, pues, un $k$-esquema algebraico;
además, si $x_i$ ($(1\leq i\leq n)$) son los puntos genéricos de las componentes irreducibles de $X$, y $x'_j$ ($(1\leq j\leq m)$) los puntos genéricos de las componentes irreducibles de $X'$, cada $\boldsymbol{k}(x'_j)$ es una extensión algebraica finita de uno de los $\boldsymbol{k}(x_i)$ \sref{II.6.3.6-fr} y tiene, por tanto, grado de trascendencia 1 sobre $k$.
Así, $X'$ es efectivamente una curva algebraica sobre $k$ y, además, se sabe que $X'$ es suma de un número finito de esquemas íntegros y normales (\sref{II.6.3.6-fr} y \sref{II.6.3.7-fr}).
\end{proof}
% Fuente francesa: PMIHES_1961__8__5_0, printed pp. 149-150 (PDF physical pages 146-147).

\begin{env}[7.4.7]
\label{II.7.4.7-fr}
Se dice que una curva algebraica $X$ sobre $k$ es \emph{completa} si es \emph{propia} sobre $k$.
\end{env}

\begin{corollary}[7.4.8]
\label{II.7.4.8-fr}
Para que una curva algebraica reducida $X$ sobre $k$ sea completa, es necesario y suficiente que su normalización $X'$ sea completa.
\end{corollary}

\oldpage[II]{150}
\begin{proof}
En efecto, el morfismo canónico $f:X'\to X$ es finito \sref{II.7.4.6-fr} y, por tanto, propio \sref{II.6.1.11-fr} y sobreyectivo \sref{II.6.3.8-fr};
si $g:X\to\operatorname{Spec}(k)$ es el morfismo estructural, entonces $g$ y $g\circ f$ son propios simultáneamente, como resulta de \sref{II.5.4.2-fr}[(ii)] y \sref{II.5.4.3-fr}[(ii)], pues $g$ es separado por hipótesis.
\end{proof}
% Fuente francesa: PMIHES_1961__8__5_0, printed p. 150 (PDF physical page 147).

\begin{proposition}[7.4.9]
\label{II.7.4.9-fr}
Sean $X$ una curva algebraica \emph{normal} sobre $k$ e $Y$ un $k$-esquema algebraico \textup{propio} sobre $k$.
Entonces toda $k$-aplicación racional de $X$ en $Y$ está definida en todas partes o, dicho de otro modo, es un morfismo.
\end{proposition}

\begin{proof}
En efecto, se sigue de \sref{II.7.3.7-fr} que, en los puntos $x\in X$ donde tal aplicación no está definida, la dimensión de $\mathcal{O}_x$ debería ser $\geq2$, por lo que el conjunto de esos puntos es vacío;
la última afirmación se deduce de \sref[I]{I.7.2.3-fr}.
\end{proof}
% Fuente francesa: PMIHES_1961__8__5_0, printed p. 150 (PDF physical page 147).

\begin{corollary}[7.4.10]
\label{II.7.4.10-fr}
Una curva algebraica normal sobre $k$ es cuasiproyectiva sobre $k$.
\end{corollary}

\begin{proof}
Como $X$ es suma de un número finito de curvas algebraicas íntegras y normales \sref{II.6.3.8-fr}, se puede reducir al caso en que $X$ es íntegra \sref{II.5.3.6-fr}.
Como $X$ es cuasicompacta, está recubierta por un número finito de abiertos afines $U_i$ ($(1\leq i\leq n)$) y, como cada uno de ellos es de tipo finito sobre $k$, existen un entero $n_i$ y una $k$-inmersión $f_i:U_i\to\mathbf{P}_k^{n_i}$ \sref{II.5.3.3-fr} y \sref{II.5.3.4-fr}[(i)].
Como $U_i$ es denso en $X$, se sigue de \sref{II.7.4.9-fr} que $f_i$ se prolonga a un $k$-morfismo $g_i:X\to\mathbf{P}_k^{n_i}$, de donde un $k$-morfismo $g=(g_1,\ldots,g_n)_k$ de $X$ en el producto $P$ de los $\mathbf{P}_k^{n_i}$ sobre $k$.
Además, para cada $i$, puesto que la restricción de $g_i$ a $U_i$ es una inmersión, también lo es la restricción de $g$ a $U_i$ \sref[I]{I.5.3.14-fr}.
Como los $U_i$ recubren $X$ y $g$ es separado \sref[I]{I.5.5.1-fr}[(v)], $g$ es una inmersión de $X$ en $P$ \sref[I]{I.8.2.8-fr}.
Como el morfismo de Segre \sref{II.4.3.3-fr} proporciona una inmersión de $P$ en $\mathbf{P}_k^N$, esto termina de probar que $X$ es cuasiproyectiva.
\end{proof}
% Fuente francesa: PMIHES_1961__8__5_0, printed p. 150 (PDF physical page 147).

\begin{corollary}[7.4.11]
\label{II.7.4.11-fr}
Toda curva algebraica normal $X$ es isomorfa al esquema inducido sobre un abierto denso en todas partes de una curva algebraica normal y completa $\widehat{X}$, determinada salvo isomorfismo único.
\end{corollary}

\begin{proof}
Si $X_1$ y $X_2$ son dos curvas normales y completas, se sigue inmediatamente de \sref{II.7.4.9-fr} que todo isomorfismo de un abierto $U_1$ denso en $X_1$ sobre un abierto $U_2$ denso en $X_2$ se prolonga de manera única a un isomorfismo de $X_1$ sobre $X_2$;
de aquí resulta la afirmación de unicidad.
Para probar la existencia de $\widehat{X}$, basta observar que se puede considerar $X$ como un subesquema de un fibrado proyectivo $\mathbf{P}_k^n$ \sref{II.7.4.10-fr}.
Sea $\overline{X}$ la clausura de $X$ en $\mathbf{P}_k^n$ \sref[I]{I.9.5.11-fr};
como $X$ está inducido por $\overline{X}$ sobre un abierto denso en $\overline{X}$ \sref[I]{I.9.5.10-fr}, los puntos genéricos $x_i$ de las componentes irreducibles de $X$ son también los de las componentes irreducibles de $\overline{X}$, y los $\boldsymbol{k}(x_i)$ son los mismos para ambos esquemas; por tanto, según \sref{II.7.4.1-fr}, $\overline{X}$ es una curva algebraica sobre $k$, reducida \sref[I]{I.9.5.9-fr} y proyectiva sobre $k$ \sref{II.5.5.1-fr}, y por ello completa \sref{II.5.5.3-fr}.
Tomemos entonces como $\widehat{X}$ la \emph{normalización} de $\overline{X}$, que sigue siendo completa \sref{II.7.4.8-fr};
además, si $h:\widehat{X}\to\overline{X}$ es el morfismo canónico, la restricción de $h$ a $h^{-1}(X)$ es un isomorfismo sobre $X$, pues $X$ es normal \sref{II.6.3.4-fr}, y, como $h^{-1}(X)$ contiene los puntos genéricos de las componentes irreducibles de $\widehat{X}$ \sref{II.6.3.8-fr}, es denso en $\widehat{X}$, lo que termina la demostración.
\end{proof}
% Fuente francesa: PMIHES_1961__8__5_0, printed p. 151 (PDF physical page 148).

\oldpage[II]{151}
\begin{remark}[7.4.12]
\label{II.7.4.12-fr}
Probaremos en el capítulo~III que la conclusión de \sref{II.7.4.10-fr} sigue siendo válida sin suponer que la curva sea normal (ni siquiera reducida);
mostraremos también que, para que una curva algebraica (reducida o no) sea afín, es necesario y suficiente que sus componentes irreducibles (reducidas) no sean completas.
\end{remark}
% Fuente francesa: PMIHES_1961__8__5_0, printed p. 151 (PDF physical page 148).

\begin{corollary}[7.4.13]
\label{II.7.4.13-fr}
Sean $X$ una curva \emph{normal} irreducible de cuerpo $R(X)=K$, e $Y$ una curva íntegra \textup{completa} de cuerpo $L=R(Y)$.
Existe una correspondencia biyectiva canónica entre los $k$-morfismos \emph{dominantes} $X\to Y$ y los $k$-monomorfismos $L\to K$.
\end{corollary}

\begin{proof}
En virtud de \sref{II.7.4.9-fr}, las $k$-aplicaciones racionales de $X$ en $Y$ se identifican con los $k$-morfismos $u:X\to Y$.
Como los morfismos dominantes $u$ se caracterizan por la condición $u(x)=y$ (designando por $x$ e $y$ los puntos genéricos respectivos de $X$ e $Y$), el corolario resulta de estas observaciones y de \sref[I]{I.7.1.13-fr}.
\end{proof}
% Fuente francesa: PMIHES_1961__8__5_0, printed p. 151 (PDF physical page 148).

\begin{env}[7.4.14]
\label{II.7.4.14-fr}
Puede precisarse el resultado de \sref{II.7.4.13-fr} cuando se toma por $Y$ la \emph{recta proyectiva} $\mathbf{P}_k^1=\operatorname{Proj}(k[T_0,T])$, siendo $T_0$ y $T$ dos indeterminadas.
Este es, en efecto, un esquema íntegro \sref{II.2.4.4-fr}, y el esquema inducido sobre el abierto $D_+(T_0)$ de $Y$ es isomorfo a $\operatorname{Spec}(k[T])$ \sref{II.2.3.6-fr}; por tanto, el punto genérico de $Y$ es el ideal $(0)$ de $k[T]$, y su cuerpo de funciones es $k(T)$, lo cual muestra que $Y$ es una curva algebraica completa sobre $k$.
Además, el único ideal primo graduado de $S=k[T_0,T]$ que contiene $T_0$ y es distinto de $S_+$ es el ideal principal $(T_0)$; por consiguiente, el complementario de $D_+(T_0)$ en $Y=\mathbf{P}_k^1$ se reduce a \emph{un punto cerrado}, llamado «punto en el infinito», que denotaremos por $\infty$ (para un estudio general de las relaciones entre fibrados vectoriales y fibrados proyectivos, véase \sref{subsection:II.8.4-fr}).
Con estas notaciones:
\end{env}
% Fuente francesa: PMIHES_1961__8__5_0, printed pp. 151-152 (PDF physical pages 148-149).

\begin{corollary}[7.4.15]
\label{II.7.4.15-fr}
Sea $X$ una curva \emph{normal} irreducible de cuerpo $R(X)=K$.
Existe una correspondencia biyectiva canónica entre el conjunto $K$ y el conjunto de los morfismos $u$ de $X$ en $\mathbf{P}_k^1$ distintos del morfismo constante de valor $\infty$.
Para que $u$ sea dominante, es necesario y suficiente que el elemento correspondiente de $K$ sea trascendente sobre $k$.
\end{corollary}

\begin{proof}
El enunciado resulta inmediatamente de \sref{II.7.4.9-fr} y del lema siguiente:
\begin{lemma}[7.4.15.1]
\label{II.7.4.15.1-fr}
Sea $X$ un preesquema íntegro sobre $k$, y sea $K=R(X)$ su cuerpo de funciones.
Existe una correspondencia biyectiva canónica entre el conjunto $K$ y el conjunto de las aplicaciones racionales $u$ de $X$ en $\mathbf{P}_k^1$ distintas del morfismo constante de valor $\infty$.
Para que tal aplicación racional sea dominante, es necesario y suficiente que el elemento correspondiente de $K$ sea trascendente sobre $k$.
\end{lemma}
Ante todo, las aplicaciones racionales de $X$ en $\mathbf{P}_k^1$ corresponden biyectivamente a los puntos de $\mathbf{P}_k^1$ con valores en la extensión $K$ de $k$ \sref[I]{I.7.1.12-fr}.
Si tal punto está situado en el punto genérico de $\mathbf{P}_k^1$ \sref[I]{I.3.4.5-fr}, la aplicación racional correspondiente es evidentemente dominante.
En el caso contrario, como todo punto de $\mathbf{P}_k^1$ distinto del punto genérico es cerrado \sref{II.7.4.3-fr}, la imagen del dominio de definición $U$ de $u$ por el único morfismo $U\to\mathbf{P}_k^1$ de la clase $u$ \sref[I]{I.7.2.2-fr} se reduce a un punto cerrado $y$ de $\mathbf{P}_k^1$, y este morfismo (que no está necesariamente definido en todo $X$) no es, por tanto, dominante;
por abuso de lenguaje, se dice entonces que la aplicación racional $u$ es «constante, de valor $y$».
Queda por poner en correspondencia biyectiva los puntos de $\mathbf{P}_k^1$ con valores en $K$ que no están situados en $\infty$ \sref[I]{I.3.4.5-fr} y los elementos de $K$, y comprobar que un punto así está situado en el punto genérico de $\mathbf{P}_k^1$ si y solo si corresponde a un
\oldpage[II]{152}
elemento trascendente sobre $k$.
Pero esto es inmediato \sref{II.4.2.6-fr}[ejemplo~1].
\end{proof}
% Fuente francesa: PMIHES_1961__8__5_0, printed p. 152 (PDF physical page 149).

\begin{corollary}[7.4.16]
\label{II.7.4.16-fr}
Sean $X$ e $Y$ dos curvas algebraicas sobre $k$, \textup{normales, completas} e irreducibles;
sean $K=R(X)$ y $L=R(Y)$ sus cuerpos de funciones.
Existe una correspondencia biyectiva canónica entre el conjunto de los $k$-isomorfismos $X\xrightarrow{\sim}Y$ y el conjunto de los $k$-isomorfismos $L\xrightarrow{\sim}K$.
\end{corollary}

\begin{proof}
Es una consecuencia evidente de \sref{II.7.4.13-fr}.
\end{proof}
% Fuente francesa: PMIHES_1961__8__5_0, printed p. 152 (PDF physical page 149).

\begin{env}[7.4.17]
\label{II.7.4.17-fr}
Este corolario \sref{II.7.4.16-fr} muestra que una curva algebraica sobre $k$, normal, completa e irreducible, está \emph{determinada por su cuerpo $K$ de funciones racionales salvo isomorfismo único};
por definición, $K$ es una extensión de tipo finito de $k$, de grado de trascendencia $1$ (lo que clásicamente se denomina un \emph{cuerpo de funciones algebraicas de una variable}).
Además:
\end{env}
% Fuente francesa: PMIHES_1961__8__5_0, printed p. 152 (PDF physical page 149).

\begin{proposition}[7.4.18]
\label{II.7.4.18-fr}
Para toda extensión $K$ de $k$, de tipo finito y de grado de trascendencia $1$, existe una curva algebraica $X$ (determinada salvo isomorfismo único), normal, completa e irreducible, tal que $R(X)=K$.
El conjunto de los anillos locales de $X$ se identifica \sref[I]{I.8.2.1-fr} con el conjunto formado por el cuerpo $K$ y los anillos de valuación que contienen a $k$ y tienen a $K$ por cuerpo de fracciones.
\end{proposition}

\begin{proof}
En efecto, $K$ es una extensión de grado finito de una extensión trascendente pura $k(T)$ de $k$, que se identifica, como se ha visto, con el cuerpo de funciones racionales de la recta proyectiva $Y=\mathbf{P}_k^1$.
Sea $X$ la \emph{clausura integral} de $Y$ respecto de $K$ \sref{II.6.3.4-fr};
entonces $X$ es una curva algebraica normal de cuerpo $K$ \sref{II.6.3.7-fr}, y es completa, puesto que el morfismo $X\to Y$ es finito \sref{II.7.4.6-fr}.
Los anillos locales $\mathcal{O}_x$ de $X$ son el cuerpo $K$, cuando $x$ es el punto genérico; si $x$ es distinto del punto genérico, $\mathcal{O}_x$ es un anillo de valuación discreta que contiene a $k$ y tiene a $K$ por cuerpo de fracciones \sref{II.7.4.5-fr}.
Recíprocamente, sea $A$ un tal anillo;
como el morfismo $X\to\operatorname{Spec}(k)$ es propio y $A$ domina a $k$, también domina un anillo local $\mathcal{O}_x$ de $X$ \sref{II.7.3.10-fr};
siendo este último un anillo de valuación que tiene a $K$ por cuerpo de fracciones, es necesariamente igual a $A$.
\end{proof}
% Fuente francesa: PMIHES_1961__8__5_0, printed p. 152 (PDF physical page 149).

\begin{namedtheorem}[7.4.19]{Observaciones}
\label{II.7.4.19-fr}
De \sref{II.7.4.16-fr} y \sref{II.7.4.18-fr} resulta que \emph{el dato de una curva algebraica sobre $k$, normal, completa e irreducible, es esencialmente equivalente al dato de una extensión $K$ de $k$, de tipo finito y de grado de trascendencia $1$}.
Se observa que, si $k'$ es una extensión del cuerpo base $k$, $X\otimes_k k'$ seguirá siendo una curva algebraica completa sobre $k'$ \sref{II.5.4.2-fr}[(iii)], pero, en general, no será ni reducida ni irreducible.
Será, sin embargo, reducida e irreducible si $K$ es una extensión separable de $k$ y $k$ es algebraicamente cerrado en $K$ (lo que, en una terminología clásica que no seguiremos, se expresa diciendo que $K$ es una «extensión regular» de $k$).
Pero incluso en este caso puede ocurrir que $X\otimes_k k'$ no sea normal.
El lector encontrará detalles sobre estas cuestiones en el capítulo~IV.
\end{namedtheorem}
% Fuente francesa: PMIHES_1961__8__5_0, printed pp. 152-153 (PDF physical pages 149-150).

\section{Esquemas obtenidos por explosión; conos proyectantes; clausura proyectiva}
\label{section:II.8-fr}

\subsection{Preesquemas obtenidos por explosión}
\label{subsection:II.8.1-fr}


\begin{env}[8.1.1]
\label{II.8.1.1-fr}
Sea $Y$ un preesquema y, para todo entero $n\geq0$, sea $\mathcal{I}_n$ un haz cuasicoherente de ideales de $\mathcal{O}_Y$; se supone verificadas las condiciones siguientes:
\begin{equation}
\label{II.8.1.1.1-fr}
  \mathcal{I}_0=\mathcal{O}_Y,\qquad
  \mathcal{I}_n\subset\mathcal{I}_m\quad\text{para }m\leq n
\tag{8.1.1.1}
\end{equation}
\begin{equation}
\label{II.8.1.1.2-fr}
  \mathcal{I}_m\mathcal{I}_n\subset\mathcal{I}_{m+n}
  \quad\text{cualesquiera que sean }m,n.
\tag{8.1.1.2}
\end{equation}

\oldpage[II]{153}
Se señala que estas hipótesis implican
\begin{equation}
\label{II.8.1.1.3-fr}
  \mathcal{I}_1^n\subset\mathcal{I}_n.
\tag{8.1.1.3}
\end{equation}

Pongamos
\begin{equation}
\label{II.8.1.1.4-fr}
  \mathcal{S}=\bigoplus_{n\geq0}\mathcal{I}_n.
\tag{8.1.1.4}
\end{equation}

De \sref{II.8.1.1.1-fr} y \sref{II.8.1.1.2-fr} resulta que $\mathcal{S}$ es una $\mathcal{O}_Y$-álgebra graduada cuasicoherente y define, por tanto, un $Y$-esquema $X=\operatorname{Proj}(\mathcal{S})$.
Si $\mathcal{J}$ es un haz invertible de ideales de $\mathcal{O}_Y$, entonces $\mathcal{I}_n\otimes_{\mathcal{O}_Y}\mathcal{J}^{\otimes n}$ se identifica canónicamente con $\mathcal{I}_n\mathcal{J}^n$.
Si se sustituyen entonces los $\mathcal{I}_n$ por los $\mathcal{I}_n\mathcal{J}^n$ y, con ello, $\mathcal{S}$ por la $\mathcal{O}_Y$-álgebra cuasicoherente $\mathcal{S}_{(\mathcal{J})}$, entonces $X_{(\mathcal{J})}=\operatorname{Proj}(\mathcal{S}_{(\mathcal{J})})$ es canónicamente isomorfo a $X$ \sref{II.3.1.8-fr}.
\end{env}
% Fuente francesa: PMIHES_1961__8__5_0, printed p. 153 (PDF physical page 150).

\begin{env}[8.1.2]
\label{II.8.1.2-fr}
Supongamos que $Y$ es \emph{localmente íntegro}, de modo que el haz $\mathcal{R}(Y)$ de funciones racionales es una $\mathcal{O}_Y$-álgebra cuasicoherente \sref[I]{I.7.3.7-fr}.
Diremos que un $\mathcal{O}_Y$-submódulo $\mathcal{I}$ de $\mathcal{R}(Y)$ es un \emph{ideal fraccional} de $\mathcal{R}(Y)$ si es \emph{de tipo finito} \sref[0]{0.5.2.1-fr}.
Supongamos dado, para todo $n\geq0$, un ideal fraccional cuasicoherente $\mathcal{I}_n$ de $\mathcal{R}(Y)$, tal que $\mathcal{I}_0=\mathcal{O}_Y$ y que se satisface la condición \sref{II.8.1.1.2-fr} (pero no necesariamente la segunda condición de \sref{II.8.1.1.1-fr});
se puede entonces definir de nuevo, mediante la fórmula \sref{II.8.1.1.4-fr}, una $\mathcal{O}_Y$-álgebra graduada cuasicoherente y el $Y$-esquema correspondiente $X=\operatorname{Proj}(\mathcal{S})$;
se tendrá de nuevo un isomorfismo canónico de $X$ sobre $X_{(\mathcal{J})}$ para todo ideal fraccional \emph{invertible} $\mathcal{J}$ de $\mathcal{R}(Y)$.
\end{env}
% Fuente francesa: PMIHES_1961__8__5_0, printed p. 153 (PDF physical page 150).

\begin{definition}[8.1.3]
\label{II.8.1.3-fr}
Sea $Y$ un preesquema \textup{(resp.} un preesquema localmente íntegro), y sea $\mathcal{I}$ un ideal cuasicoherente de $\mathcal{O}_Y$ \textup{(resp.} un ideal fraccional cuasicoherente de $\mathcal{R}(Y)$).
Se dice que el $Y$-esquema
$X=\operatorname{Proj}(\bigoplus_{n\geq0}\mathcal{I}^n)$
se obtiene por explosión del ideal $\mathcal{I}$, o que es el preesquema obtenido por explosión de $Y$ relativamente a $\mathcal{I}$.
Cuando $\mathcal{I}$ es un ideal cuasicoherente de $\mathcal{O}_Y$ y $Y'$ es el subpreesquema cerrado de $Y$ definido por $\mathcal{I}$, se dice también que $X$ es el $Y$-esquema obtenido por explosión de $Y'$.
\end{definition}

Por definición, $\mathcal{S}=\bigoplus_{n\geq0}\mathcal{I}^n$ está entonces generada por $\mathcal{S}_1=\mathcal{I}$;
si $\mathcal{I}$ es un $\mathcal{O}_Y$-módulo \emph{de tipo finito}, $X$ es, por tanto, \emph{proyectivo} sobre $Y$ \sref{II.5.5.2-fr}.
Sin hipótesis alguna sobre $\mathcal{I}$, el $\mathcal{O}_X$-módulo $\mathcal{O}_X(1)$ es \emph{invertible} \sref{II.3.2.5-fr} y \emph{muy amplio}, en virtud de \sref{II.4.4.3-fr} aplicado al morfismo estructural $X\to Y$.

Se señala que, si $f:X\to Y$ es el morfismo estructural, la restricción de $f$ a $f^{-1}(Y\setmin Y')$ es un \emph{isomorfismo} sobre $Y\setmin Y'$ cuando $\mathcal{I}$ es un \emph{ideal de $\mathcal{O}_Y$} y $Y'$ es el subpreesquema cerrado que define:
en efecto, como la cuestión es local sobre $Y$, basta suponer que $\mathcal{I}=\mathcal{O}_Y$, y nuestra afirmación resulta de \sref{II.3.1.7-fr}.

Cuando se sustituye $\mathcal{I}$ por $\mathcal{I}^d$ ($d>0$), el $Y$-esquema obtenido por explosión $X$ es sustituido por un $Y$-esquema $X'$ canónicamente \emph{isomorfo} \sref{II.8.1.1-fr};
igualmente, para todo ideal (resp. ideal fraccional) \emph{invertible} $\mathcal{J}$, el preesquema obtenido por explosión $X_{(\mathcal{J})}$ relativamente al ideal $\mathcal{I}\mathcal{J}$ es canónicamente isomorfo a $X$ \sref{II.8.1.1-fr}.

En particular, cuando $\mathcal{I}$ es un ideal (resp. un ideal fraccional) \emph{invertible}, el $Y$-esquema obtenido por explosión de $\mathcal{I}$ es \emph{isomorfo a $Y$} \sref{II.3.1.7-fr}.
% Fuente francesa: PMIHES_1961__8__5_0, printed p. 154 (PDF physical page 151).

\begin{proposition}[8.1.4]
\label{II.8.1.4-fr}
Sea $Y$ un preesquema íntegro.
\begin{enumerate}
    \item[\rm{(i)}] Para toda sucesión $(\mathcal{I}_n)$ de ideales fraccionales cuasicoherentes de $\mathcal{R}(Y)$ que satisfaga \sref{II.8.1.1.2-fr}
\oldpage[II]{154}
        y tal que $\mathcal{I}_0=\mathcal{O}_Y$, el $Y$-esquema $X=\operatorname{Proj}(\bigoplus_{n\geq0}\mathcal{I}_n)$ es íntegro, y el morfismo estructural $f:X\to Y$ es dominante.
    \item[\rm{(ii)}] Sea $\mathcal{I}$ un ideal fraccional cuasicoherente de $\mathcal{R}(Y)$, y sea $X$ el $Y$-esquema obtenido por explosión de $Y$ relativamente a $\mathcal{I}$.
        Si $\mathcal{I}\neq0$, entonces el morfismo estructural $f:X\to Y$ es birracional y sobreyectivo.
\end{enumerate}
\end{proposition}

\begin{proof}
\medskip\noindent
\begin{enumerate}
    \item[\rm{(i)}] Esto resulta del hecho de que $\mathcal{S}=\bigoplus_{n\geq0}\mathcal{I}_n$ es una $\mathcal{O}_Y$-álgebra \emph{íntegra} (\sref{II.3.1.12-fr} y \sref{II.3.1.14-fr}), pues, para todo $y\in Y$, $\mathcal{O}_y$ es un anillo íntegro \sref[I]{I.5.1.4-fr}.
    \item[\rm{(ii)}] Por (i), $X$ es íntegro; si, además, $x$ e $y$ son los puntos genéricos de $X$ e $Y$, respectivamente, se tiene $f(x)=y$, y queda por demostrar que $\boldsymbol{k}(x)$ es de rango $1$ sobre $\boldsymbol{k}(y)$.
        Pero $x$ es también el punto genérico de la fibra $f^{-1}(y)$; si $\psi$ es el morfismo canónico $Z\to Y$, donde $Z=\operatorname{Spec}(\boldsymbol{k}(y))$, entonces el preesquema $f^{-1}(y)$ puede identificarse con $\operatorname{Proj}(\mathcal{S}')$, donde $\mathcal{S}'=\psi^*(\mathcal{S})$ \sref{II.3.5.3-fr}.
        Pero es claro que $\mathcal{S}'=\bigoplus_{n\geq0}(\mathcal{I}_y)^n$, y, como $\mathcal{I}$ es un ideal fraccional cuasicoherente de $\mathcal{R}(Y)$ que no es nulo, $\mathcal{I}_y\neq0$ \sref[I]{I.7.3.6-fr}, de donde $\mathcal{I}_y=\boldsymbol{k}(y)$;
        entonces $\operatorname{Proj}(\mathcal{S}')$ puede identificarse con $\operatorname{Spec}(\boldsymbol{k}(y))$ \sref{II.3.1.7-fr}, de donde se sigue la conclusión.
\end{enumerate}
\end{proof}
% Fuente francesa: PMIHES_1961__8__5_0, printed pp. 154-155 (PDF physical pages 151-152).

Demostraremos una recíproca de \sref{II.8.1.4-fr} en \hyperref[III.2.3.7-fr]{\textup{(\textbf{III}, 2.3.7)}}.

\begin{env}[8.1.5]
\label{II.8.1.5-fr}
Volvamos a la situación y a las notaciones de \sref{II.8.1.1-fr}.
Por definición, los homomorfismos de inclusión $\mathcal{I}_{n+1}\to\mathcal{I}_n$ \sref{II.8.1.1.1-fr} definen, para cada $k\in\mathbf{Z}$, un homomorfismo inyectivo de grado $0$ de $\mathcal{S}$-módulos graduados
\begin{equation}
\label{II.8.1.5.1-fr}
  u_k:\mathcal{S}_+(k+1)\to\mathcal{S}(k)~;
\tag{8.1.5.1}
\end{equation}
como $\mathcal{S}_+(k+1)$ y $\mathcal{S}(k+1)$ son canónicamente \textbf{(TN)}-isomorfos, a $u_k$ le corresponde canónicamente un homomorfismo inyectivo de $\mathcal{O}_X$-módulos \sref{II.3.4.2-fr}:
\begin{equation}
\label{II.8.1.5.2-fr}
  \widetilde{u}_k:\mathcal{O}_X(k+1)\to\mathcal{O}_X(k).
\tag{8.1.5.2}
\end{equation}

Recordemos, por otra parte, \sref{II.3.2.6-fr} que se han definido homomorfismos canónicos
\begin{equation}
\label{II.8.1.5.3-fr}
  \lambda:\mathcal{O}_X(h)\otimes_{\mathcal{O}_X}\mathcal{O}_X(k)
  \to\mathcal{O}_X(h+k)
\tag{8.1.5.3}
\end{equation}
y, como el diagrama
\[
\xymatrix{
  \mathcal{S}(h)\otimes_{\mathcal{S}}\mathcal{S}(k)
    \otimes_{\mathcal{S}}\mathcal{S}(l)\ar[r]\ar[d]
  & \mathcal{S}(h+k)\otimes_{\mathcal{S}}\mathcal{S}(l)\ar[d]\\
  \mathcal{S}(h)\otimes_{\mathcal{S}}\mathcal{S}(k+l)\ar[r]
  & \mathcal{S}(h+k+l)
}
\]
es conmutativo, de la funtorialidad de las $\lambda$ \sref{II.3.2.6-fr} resulta que los homomorfismos \sref{II.8.1.5.3-fr} definen sobre
\begin{equation}
\label{II.8.1.5.4-fr}
  \mathcal{S}_X=\bigoplus_{n\in\mathbf{Z}}\mathcal{O}_X(n)
\tag{8.1.5.4}
\end{equation}
\oldpage[II]{155}
una estructura de \emph{$\mathcal{O}_X$-álgebra graduada cuasicoherente}.

Además, el diagrama
\[
\xymatrix{
  \mathcal{S}(h)\otimes_{\mathcal{S}}\mathcal{S}_+(k+1)\ar[r]\ar[d]_{1\otimes u_k}
  & \mathcal{S}_+(h+k+1)\ar[d]^{u_{k+h}}\\
  \mathcal{S}(h)\otimes_{\mathcal{S}}\mathcal{S}(k)\ar[r]
  & \mathcal{S}(h+k)
}
\]
es conmutativo; la funtorialidad de las $\lambda$ muestra entonces que se tiene un diagrama conmutativo
\begin{equation}
\label{II.8.1.5.5-fr}
\xymatrix{
  \mathcal{O}_X(h)\otimes_{\mathcal{O}_X}\mathcal{O}_X(k+1)
    \ar[r]^{\lambda}\ar[d]_{1\otimes\widetilde{u}_k}
  & \mathcal{O}_X(h+k+1)\ar[d]^{\widetilde{u}_{k+h}}\\
  \mathcal{O}_X(h)\otimes_{\mathcal{O}_X}\mathcal{O}_X(k)
    \ar[r]^{\lambda}
  & \mathcal{O}_X(h+k)
}
\tag{8.1.5.5}
\end{equation}
siendo las flechas horizontales los homomorfismos canónicos.
Se puede decir, por tanto, que los $\widetilde{u}_k$ definen un \emph{homomorfismo inyectivo} (de grado $0$) \emph{de $\mathcal{S}_X$-módulos graduados}
\begin{equation}
\label{II.8.1.5.6-fr}
  \widetilde{u}:\mathcal{S}_X(1)\to\mathcal{S}_X.
\tag{8.1.5.6}
\end{equation}
\end{env}
% Fuente francesa: PMIHES_1961__8__5_0, printed p. 155 (PDF physical page 152).

\begin{env}[8.1.6]
\label{II.8.1.6-fr}
Conservando las notaciones de \sref{II.8.1.5-fr}, observemos ahora que, para $n\geq0$, el homomorfismo compuesto $\widetilde{v}_n=\widetilde{u}_{n-1}\circ\widetilde{u}_{n-2}\circ\ldots\circ\widetilde{u}_0$ es un homomorfismo \emph{inyectivo} $\mathcal{O}_X(n)\to\mathcal{O}_X$;
denotaremos por $\mathcal{I}_{n,X}$ su imagen, que es, por tanto, un ideal cuasicoherente de $\mathcal{O}_X$, \emph{isomorfo} a $\mathcal{O}_X(n)$.
Además, el diagrama
\[
\xymatrix{
  \mathcal{O}_X(m)\otimes_{\mathcal{O}_X}\mathcal{O}_X(n)
    \ar[r]^{\lambda}\ar[d]_{\widetilde{v}_m\otimes\widetilde{v}_n}
  & \mathcal{O}_X(m+n)\ar[d]^{\widetilde{v}_{m+n}}\\
  \mathcal{O}_X\ar[r]_{\mathrm{id}}
  & \mathcal{O}_X
}
\]
es conmutativo para $m\geq0$, $n\geq0$.
Se deducen, por tanto, las inclusiones siguientes:
\begin{equation}
\label{II.8.1.6.1-fr}
  \mathcal{I}_{0,X}=\mathcal{O}_X,\qquad
  \mathcal{I}_{n,X}\subset\mathcal{I}_{m,X}
  \quad\text{para }0\leq m\leq n
\tag{8.1.6.1}
\end{equation}
\begin{equation}
\label{II.8.1.6.2-fr}
  \mathcal{I}_{m,X}\mathcal{I}_{n,X}\subset\mathcal{I}_{m+n,X}
  \quad\text{para }m\geq0,\ n\geq0.
\tag{8.1.6.2}
\end{equation}
\end{env}
% Fuente francesa: PMIHES_1961__8__5_0, printed p. 156 (PDF physical page 153).

\oldpage[II]{156}

\begin{proposition}[8.1.7]
\label{II.8.1.7-fr}
Sea $Y$ un preesquema, $\mathcal{I}$ un ideal cuasicoherente de $\mathcal{O}_Y$, y $X=\operatorname{Proj}(\bigoplus_{n\geq0}\mathcal{I}^n)$ el $Y$-esquema obtenido por explosión de $\mathcal{I}$.
Se tiene entonces, para todo $n>0$, un isomorfismo canónico
\begin{equation}
\label{II.8.1.7.1-fr}
  \mathcal{O}_X(n)\xrightarrow{\sim}\mathcal{I}^n\mathcal{O}_X
  =\mathcal{I}_{n,X}
\tag{8.1.7.1}
\end{equation}
\textup{(cf. \sref[0]{0.4.3.5-fr})}, y, por tanto, $\mathcal{I}^n\mathcal{O}_X$ es un $\mathcal{O}_X$-módulo invertible muy amplio si $n>0$.
\end{proposition}

\begin{proof}
La última afirmación es inmediata, puesto que $\mathcal{O}_X(1)$ es invertible \sref{II.3.2.5-fr} y muy amplio para $Y$ por definición (\sref{II.4.4.3-fr} y \sref{II.4.4.9-fr}).
Por otra parte, por definición, la imagen de $v_n$ es precisamente $\mathcal{I}^n\mathcal{S}$, y \sref{II.8.1.7.1-fr} resulta entonces de la exactitud del funtor $\widetilde{\mathcal{M}}$ \sref{II.3.2.4-fr} y de la fórmula~\sref{II.3.2.4.1-fr}.
\end{proof}
% Fuente francesa: PMIHES_1961__8__5_0, printed p. 156 (PDF physical page 153).

\begin{corollary}[8.1.8]
\label{II.8.1.8-fr}
Bajo las hipótesis de \sref{II.8.1.7-fr}, si $f:X\to Y$ es el morfismo estructural e $Y'$ el subpreesquema cerrado de $Y$ definido por $\mathcal{I}$, entonces el subpreesquema cerrado $X'=f^{-1}(Y')$ de $X$ está definido por $\mathcal{I}\mathcal{O}_X$ (canónicamente isomorfo a $\mathcal{O}_X(1)$), de donde se obtiene una sucesión exacta corta canónica
\begin{equation}
\label{II.8.1.8.1-fr}
  0\to\mathcal{O}_X(1)\to\mathcal{O}_X\to\mathcal{O}_{X'}\to0
\tag{8.1.8.1}
\end{equation}
\end{corollary}

\begin{proof}
Esto resulta de \sref{II.8.1.7.1-fr} y de \sref[I]{I.4.4.5-fr}.
\end{proof}
% Fuente francesa: PMIHES_1961__8__5_0, printed p. 156 (PDF physical page 153).

\begin{env}[8.1.9]
\label{II.8.1.9-fr}
Bajo las hipótesis de \sref{II.8.1.7-fr}, se puede precisar la estructura de los $\mathcal{I}_{n,X}$.
Observemos, en efecto, que el homomorfismo
\[
  \widetilde{u}_{-1}:\mathcal{O}_X\to\mathcal{O}_X(-1)
\]
corresponde canónicamente a una sección $s$ de $\mathcal{O}_X(-1)$ sobre $X$, que llamaremos la \emph{sección canónica} (relativa a $\mathcal{I}$) \sref[0]{0.5.1.1-fr}.
Por otra parte, en el diagrama~\sref{II.8.1.5.5-fr}, las flechas horizontales son isomorfismos \sref{II.3.2.7-fr}; sustituyendo en este diagrama $h$ por $k$ y $k$ por $-1$, se obtiene $\widetilde{u}_{k-1}=1_k\otimes\widetilde{u}_{-1}$ (donde $1_h$ designa la identidad de $\mathcal{O}_X(h)$), es decir, el homomorfismo $\widetilde{u}_k$ no es otro que la \emph{multiplicación tensorial por la sección canónica $s$} (para todo $k\in\mathbf{Z}$).
El homomorfismo $\widetilde{u}$ \sref{II.8.1.5.6-fr} se interpreta también de la misma manera.

Por tanto, para todo $n\geq0$, el homomorfismo $\widetilde{v}_n:\mathcal{O}_X(n)\to\mathcal{O}_X$ no es otro que la multiplicación tensorial por $s^{\otimes n}$; se deduce:
\end{env}
% Fuente francesa: PMIHES_1961__8__5_0, printed p. 156 (PDF physical page 153).

\begin{corollary}[8.1.10]
\label{II.8.1.10-fr}
Con las notaciones de \sref{II.8.1.8-fr}, el espacio subyacente a $X'$ es el conjunto de los $x\in X$ tales que $s(x)=0$, donde $s$ designa la sección canónica de $\mathcal{O}_X(-1)$.
\end{corollary}

\begin{proof}
En efecto, si $c_x$ es un generador del tallo $(\mathcal{O}_X(1))_x$ en un punto $x$, entonces $s_x\otimes c_x$ se identifica canónicamente con un generador del tallo de $\mathcal{I}_{1,X}$ en el punto $x$, y es por tanto invertible si y solo si $s_x\notin\mathfrak{m}_x(\mathcal{O}_X(-1))_x$, es decir, si y solo si $s(x)\neq0$.
\end{proof}
% Fuente francesa: PMIHES_1961__8__5_0, printed pp. 156-157 (PDF physical pages 153-154).

\begin{proposition}[8.1.11]
\label{II.8.1.11-fr}
Sea $Y$ un preesquema íntegro, $\mathcal{I}$ un ideal fraccional cuasicoherente de $\mathcal{R}(Y)$, y $X$ el $Y$-esquema obtenido por explosión de $\mathcal{I}$.
Entonces $\mathcal{I}\mathcal{O}_X$ es un $\mathcal{O}_X$-módulo invertible muy amplio para $Y$.
\end{proposition}

\begin{proof}
Como la cuestión es local sobre $Y$ \sref{II.4.4.5-fr}, se puede reducir al caso en que $Y=\operatorname{Spec}(A)$, donde $A$ es un anillo íntegro cuyo cuerpo de fracciones es $K$, y $\mathcal{I}=\widetilde{\mathfrak{J}}$, siendo $\mathfrak{J}$ un ideal fraccional de $K$;
existe entonces un elemento $a\neq0$ de $A$ tal que $a\mathfrak{J}\subset A$.
Pongamos $S=\bigoplus_{n\geq0}\mathfrak{J}^n$;
la aplicación $x\to ax$ es un $A$-isomorfismo de $\mathfrak{J}^{n+1}=(S(1))_n$ sobre $a\mathfrak{J}^{n+1}=a\mathfrak{J}S_n\subset\mathfrak{J}^n=S_n$,
\oldpage[II]{157}
y define por tanto un \textbf{(TN)}-isomorfismo de grado $0$ de $S$-módulos graduados $S_+(1)\to a\mathfrak{J}S$.
Por otra parte, $x\to a^{-1}x$ es un isomorfismo de grado $0$ de $S$-módulos graduados $a\mathfrak{J}S\xrightarrow{\sim}\mathfrak{J}S$.
Se obtiene así, por composición \sref{II.3.2.4-fr}, un isomorfismo de $\mathcal{O}_X$-módulos $\mathcal{O}_X(1)\xrightarrow{\sim}\mathcal{I}\mathcal{O}_X$ y, como $S$ está generada por $S_1=\mathfrak{J}$, $\mathcal{O}_X(1)$ es invertible \sref{II.3.2.5-fr} y muy amplio (\sref{II.4.4.3-fr} y \sref{II.4.4.9-fr}), de donde se sigue la afirmación.
\end{proof}
% Fuente francesa: PMIHES_1961__8__5_0, printed p. 157 (PDF physical page 154).

\subsection{Resultados preliminares sobre la localización de anillos graduados}
\label{subsection:II.8.2-fr}

\begin{env}[8.2.1]
\label{II.8.2.1-fr}
Sea $S$ un anillo graduado, que por el momento no suponemos concentrado en grados no negativos.
Se pondrá
\begin{equation}
\label{II.8.2.1.1-fr}
  S^{\geq}=\bigoplus_{n\geq0}S_n,
  \qquad
  S^{\leq}=\bigoplus_{n\leq0}S_n
\tag{8.2.1.1}
\end{equation}
que son subanillos graduados de $S$, concentrados respectivamente en grados no negativos y no positivos.
Si $f$ es un elemento homogéneo de grado $d$ (positivo o negativo) de $S$, el anillo de fracciones $S_f=S'$ queda dotado de nuevo de una estructura de anillo graduado, tomando como $S'_n$ ($n\in\mathbf{Z}$) el conjunto de los $x/f^k$, donde $x\in S_{n+kd}$ ($k\geq0$);
se pondrá además $S_{(f)}=S'_0$, y se escribirá $S_f^{\geq}$ y $S_f^{\leq}$ por $S'^{\geq}$ y $S'^{\leq}$, respectivamente.
Si $d>0$, se tiene
\begin{equation}
\label{II.8.2.1.2-fr}
  (S^{\geq})_f=S_f
\tag{8.2.1.2}
\end{equation}
pues, si $x\in S_{n+kd}$ y $n+kd<0$, puede escribirse $x/f^k=xf^h/f^{h+k}$, y se tiene $n+(h+k)d>0$ para $h$ suficientemente grande y $>0$.
Se concluye, por definición, que
\begin{equation}
\label{II.8.2.1.3-fr}
  (S^{\geq})_{(f)}=(S_f^{\geq})_0=S_{(f)}.
\tag{8.2.1.3}
\end{equation}

Si $M$ es un $S$-módulo graduado, se pondrá análogamente
\begin{equation}
\label{II.8.2.1.4-fr}
  M^{\geq}=\bigoplus_{n\geq0}M_n,
  \qquad
  M^{\leq}=\bigoplus_{n\leq0}M_n
\tag{8.2.1.4}
\end{equation}
que son respectivamente un $S^{\geq}$-módulo graduado y un $S^{\leq}$-módulo graduado, y tienen por intersección el $S_0$-módulo $M_0$.
Si $f\in S_d$, se define asimismo $M_f$ como el $S_f$-módulo graduado cuyos elementos de grado $n$ son los $z/f^k$, donde $z\in M_{n+kd}$ ($k\geq0$);
se designará por $M_{(f)}$ el conjunto de los elementos de grado $0$ de $M_f$, que es un $S_{(f)}$-módulo, y se escribirá $M_f^{\geq}$ y $M_f^{\leq}$ en lugar de $(M_f)^{\geq}$ y $(M_f)^{\leq}$, respectivamente.
Si $d>0$, se ve, como antes, que
\begin{equation}
\label{II.8.2.1.5-fr}
  (M^{\geq})_f=M_f
\tag{8.2.1.5}
\end{equation}
y
\begin{equation}
\label{II.8.2.1.6-fr}
  (M^{\geq})_{(f)}=(M_f^{\geq})_0=M_{(f)}.
\tag{8.2.1.6}
\end{equation}
\end{env}
% Fuente francesa: PMIHES_1961__8__5_0, printed pp. 157-158 (PDF physical pages 154-155).

\begin{env}[8.2.2]
\label{II.8.2.2-fr}
Sea $\mathbf{z}$ una indeterminada, que llamaremos \emph{variable de homogeneización}.
Si $S$ es un anillo graduado (con grados enteros arbitrarios), entonces el álgebra de polinomios\footnote{No debe confundirse con el uso de la notación $\widehat{S}$ para designar la completación separada de un anillo.}
\begin{equation}
\label{II.8.2.2.1-fr}
  \widehat{S}=S[\mathbf{z}]
\tag{8.2.2.1}
\end{equation}
\oldpage[II]{158}
es una $S$-álgebra graduada, donde se define el grado de $f\mathbf{z}^n$ ($n\geq0$), con $f$ homogéneo, mediante
\begin{equation}
\label{II.8.2.2.2-fr}
  \deg(f\mathbf{z}^n)=n+\deg f.
\tag{8.2.2.2}
\end{equation}
\end{env}
% Fuente francesa: PMIHES_1961__8__5_0, printed p. 158 (PDF physical page 155).

\begin{lemma}[8.2.3]
\label{II.8.2.3-fr}
\begin{enumerate}
  \item[\rm{(i)}] Existen isomorfismos canónicos de anillos (no graduados)
    \begin{equation}
\label{II.8.2.3.1-fr}
  \widehat{S}_{(\mathbf{z})}\xrightarrow{\sim}
  \widehat{S}/(\mathbf{z}-1)\widehat{S}\xrightarrow{\sim}S
\tag{8.2.3.1}
\end{equation}
  \item[\rm{(ii)}] Existe un isomorfismo canónico de anillos (no graduados)
    \begin{equation}
\label{II.8.2.3.2-fr}
  \widehat{S}_{(f)}\xrightarrow{\sim}S_f^{\leq}
\tag{8.2.3.2}
\end{equation}
    para todo $f\in S_d$ con $d>0$.
\end{enumerate}
\end{lemma}

\begin{proof}
El primero de los isomorfismos de \sref{II.8.2.3.1-fr} se definió en \sref{II.2.2.5-fr}, y el segundo es trivial;
el isomorfismo $\widehat{S}_{(\mathbf{z})}\xrightarrow{\sim}S$ así definido hace corresponder a $x\mathbf{z}^n/\mathbf{z}^{n+k}$ (donde $\deg(x)=k$, con $k\geq-n$) el elemento $x$.
El homomorfismo \sref{II.8.2.3.2-fr} hace corresponder a $x\mathbf{z}^n/f^k$ (donde $\deg(x)=kd-n$) el elemento $x/f^k$, de grado $-n$ en $S_f^{\leq}$, y es de nuevo evidente que se obtiene efectivamente un isomorfismo.
\end{proof}
% Fuente francesa: PMIHES_1961__8__5_0, printed p. 158 (PDF physical page 155).

\begin{env}[8.2.4]
\label{II.8.2.4-fr}
Sea $M$ un $S$-módulo graduado.
Es claro que el $S$-módulo
\begin{equation}
\label{II.8.2.4.1-fr}
  \widehat{M}=M\otimes_S\widehat{S}=M\otimes_S S[\mathbf{z}]
\tag{8.2.4.1}
\end{equation}
es suma directa de los $S$-módulos $M\otimes S\mathbf{z}^n$, y por tanto de los grupos abelianos $M_k\otimes S\mathbf{z}^n$ ($k\in\mathbf{Z}$, $n\geq0$);
se define sobre $\widehat{M}$ una estructura de $\widehat{S}$-módulo graduado tomando
\begin{equation}
\label{II.8.2.4.2-fr}
  \deg(x\otimes\mathbf{z}^n)=n+\deg x
\tag{8.2.4.2}
\end{equation}
para todo $x$ homogéneo en $M$.
Dejamos al lector el cuidado de demostrar el análogo de \sref{II.8.2.3-fr}:
\end{env}
% Fuente francesa: PMIHES_1961__8__5_0, printed p. 158 (PDF physical page 155).

\begin{lemma}[8.2.5]
\label{II.8.2.5-fr}
\begin{enumerate}
  \item[\rm{(i)}] Existe un diisomorfismo canónico de módulos (no graduados)
    \begin{equation}
\label{II.8.2.5.1-fr}
  \widehat{M}_{(\mathbf{z})}\xrightarrow{\sim}M.
\tag{8.2.5.1}
\end{equation}
  \item[\rm{(ii)}] Para todo $f\in S_d$ ($d>0$), existe un diisomorfismo de módulos (no graduados)
    \begin{equation}
\label{II.8.2.5.2-fr}
  \widehat{M}_{(f)}\xrightarrow{\sim}M_f^{\leq}.
\tag{8.2.5.2}
\end{equation}
\end{enumerate}
\end{lemma}
% Fuente francesa: PMIHES_1961__8__5_0, printed p. 158 (PDF physical page 155).

\begin{env}[8.2.6]
\label{II.8.2.6-fr}
Sea $S$ un anillo graduado en grados \emph{positivos}, y consideremos en $S$ la sucesión decreciente de ideales graduados
\begin{equation}
\label{II.8.2.6.1-fr}
  S_{[n]}=\bigoplus_{m\geq n}S_m\qquad(n\geq0)
\tag{8.2.6.1}
\end{equation}
(en particular, se tiene $S_{[0]}=S$ y $S_{[1]}=S_+$).
Como es evidente que $S_{[m]}S_{[n]}\subset S_{[m+n]}$, se puede definir un \emph{anillo graduado} $S^\natural$ tomando
\begin{equation}
\label{II.8.2.6.2-fr}
  S^\natural=\bigoplus_{n\geq0}S_n^\natural
  \quad\text{con}\quad S_n^\natural=S_{[n]}.
\tag{8.2.6.2}
\end{equation}
$S_0^\natural$ es entonces el anillo $S$ considerado como anillo \emph{no graduado}, y por tanto $S^\natural$ es una $S_0^\natural$-álgebra.
Para todo elemento homogéneo $f\in S_d$ ($d>0$), denotaremos por $f^\natural$ el elemento $f$ considerado como perteneciente a $S_{[d]}=S_d^\natural$.
Con estas notaciones:
\end{env}
% Fuente francesa: PMIHES_1961__8__5_0, printed p. 159 (PDF physical page 156).

\oldpage[II]{159}
\begin{lemma}[8.2.7]
\label{II.8.2.7-fr}
Sea $S$ un anillo graduado en grados positivos, y sea $f$ un elemento homogéneo de $S_d$ ($d>0$).
Existen isomorfismos canónicos de anillos
\begin{equation}
\label{II.8.2.7.1-fr}
  S_f\xrightarrow{\sim}\bigoplus_{n\in\mathbf{Z}}S(n)_{(f)}
\tag{8.2.7.1}
\end{equation}
\begin{equation}
\label{II.8.2.7.2-fr}
  (S_f^{\geq})_{f/1}\xrightarrow{\sim}S_f
\tag{8.2.7.2}
\end{equation}
\begin{equation}
\label{II.8.2.7.3-fr}
  S^\natural_{(f^\natural)}\xrightarrow{\sim}S_f^{\geq}
\tag{8.2.7.3}
\end{equation}
donde los dos primeros son isomorfismos de anillos graduados.
\end{lemma}

\begin{proof}
Por definición, se tiene inmediatamente $(S_f)_n=(S(n)_f)_0$, de donde resulta el isomorfismo de \sref{II.8.2.7.1-fr}, que no es sino la identidad.
Por otra parte, como $f/1$ es invertible en $S_f$, existe un isomorfismo canónico $S_f\xrightarrow{\sim}(S_f^{\geq})_{f/1}=(S_f)_{f/1}$, en virtud de \sref{II.8.2.1.2-fr} aplicado a $S_f$;
el isomorfismo inverso es, por definición, el isomorfismo de \sref{II.8.2.7.2-fr}.
Finalmente, si $x=\sum_{m\geq n}y_m$ es un elemento de $S_{[n]}$, con $n=kd$, al elemento $x/(f^\natural)^k$ le corresponde el elemento $\sum_m y_m/f^k$ de $S_f^{\geq}$, y se comprueba de inmediato que esto define un isomorfismo \sref{II.8.2.7.3-fr}.
\end{proof}
% Fuente francesa: PMIHES_1961__8__5_0, printed p. 159 (PDF physical page 156).

\begin{env}[8.2.8]
\label{II.8.2.8-fr}
Si $M$ es un $S$-módulo graduado, se define análogamente, para todo $n\in\mathbf{Z}$,
\begin{equation}
\label{II.8.2.8.1-fr}
  M_{[n]}=\bigoplus_{m\geq n}M_m
\tag{8.2.8.1}
\end{equation}
y, como $S_{[m]}M_{[n]}\subset M_{[m+n]}$ ($m\geq0$), se puede definir un $S^\natural$-módulo graduado $M^\natural$ tomando
\begin{equation}
\label{II.8.2.8.2-fr}
  M^\natural=\bigoplus_{n\in\mathbf{Z}}M_n^\natural,
  \quad\text{con}\quad M_n^\natural=M_{[n]}.
\tag{8.2.8.2}
\end{equation}
\end{env}
% Fuente francesa: PMIHES_1961__8__5_0, printed p. 159 (PDF physical page 156).

Dejamos de nuevo al lector el cuidado de demostrar el siguiente lema:
\begin{lemma}[8.2.9]
\label{II.8.2.9-fr}
Con las notaciones de \sref{II.8.2.7-fr} y \sref{II.8.2.8-fr}, existen diisomorfismos canónicos de módulos
\begin{equation}
\label{II.8.2.9.1-fr}
  M_f\xrightarrow{\sim}\bigoplus_{n\in\mathbf{Z}}M(n)_{(f)}
\tag{8.2.9.1}
\end{equation}
\begin{equation}
\label{II.8.2.9.2-fr}
  (M_f^{\geq})_{f/1}\xrightarrow{\sim}M_f
\tag{8.2.9.2}
\end{equation}
\begin{equation}
\label{II.8.2.9.3-fr}
  M^\natural_{(f^\natural)}\xrightarrow{\sim}M_f^{\geq}
\tag{8.2.9.3}
\end{equation}
donde los dos primeros son diisomorfismos de módulos graduados.
\end{lemma}
% Fuente francesa: PMIHES_1961__8__5_0, printed pp. 159--160 (PDF physical pages 156--157).

\begin{lemma}[8.2.10]
\label{II.8.2.10-fr}
Sea $S$ un anillo graduado en grados positivos.
\begin{enumerate}
  \item[\rm{(i)}] Para que $S^\natural$ sea una $S_0^\natural$-álgebra de tipo finito \textup{(resp.} un anillo noetheriano\textup{)}, es necesario y suficiente que $S$ sea una $S_0$-álgebra de tipo finito \textup{(resp.} un anillo noetheriano\textup{)}.
  \item[\rm{(ii)}] Para que $S_{n+1}^\natural=S_1^\natural S_n^\natural$ ($n\geq n_0$), es necesario y suficiente que $S_{n+1}=S_1S_n$ ($n\geq n_0$).
  \item[\rm{(iii)}] Para que $S_n^\natural=S_1^{\natural n}$ ($n\geq n_0$), es necesario y suficiente que $S_n=S_1^n$ ($n\geq n_0$).
  \item[\rm{(iv)}] Si $(f_\alpha)$ es un conjunto de elementos homogéneos de $S_+$ tal que $S_+$ está contenido en el radical en $S$ del ideal de $S$ generado por los $f_\alpha$, entonces $S_+^\natural$ está contenido en el radical en $S^\natural$ del ideal de $S^\natural$ generado por los $f_\alpha^\natural$.
\end{enumerate}
\end{lemma}

\begin{proof}
\medskip\noindent
\begin{enumerate}
  \item[\rm{(i)}] Si $S^\natural$ es una $S_0^\natural$-álgebra de tipo finito, entonces $S_+=S_1^\natural$ es un módulo de tipo finito sobre $S=S_0^\natural$, por \sref{II.2.1.6-fr}[i], y por tanto $S$ es una $S_0$-álgebra de tipo finito \sref{II.2.1.4-fr};
    si $S^\natural$ es un anillo noetheriano, también lo es $S_0^\natural=S$ \sref{II.2.1.5-fr}.
    Recíprocamente, si $S$ es una $S_0$-álgebra
\oldpage[II]{160}
    de tipo finito, se sabe \sref{II.2.1.6-fr}[ii] que existen $h>0$ y $m_0>0$ tales que $S_{n+h}=S_hS_n$ para $n\geq m_0$;
    evidentemente, puede suponerse que $m_0\geq h$.
    Además, los $S_m$ son $S_0$-módulos de tipo finito \sref{II.2.1.6-fr}[i].
    Así, si $n\geq m_0+h$, se tiene $S_n^\natural=S_hS_{n-h}^\natural
=S_h^\natural S_{n-h}^\natural$;
    y si $m<m_0+h$, poniendo $E=S_{m_0}+\ldots+S_{m_0+h-1}$, se tiene
    $S_m^\natural=S_m+\ldots+S_{m_0+h-1}+S_hE+S_h^2E+\ldots$
    Para $1\leq m\leq m_0$, sea $G_m$ la reunión de los sistemas finitos de generadores de los $S_0$-módulos $S_i$ para $m\leq i\leq m_0+h-1$, considerada como subconjunto de $S_{[m]}$.
    Para $m_0+1\leq m\leq m_0+h-1$, sea $G_m$ la reunión de los sistemas finitos de generadores de los $S_0$-módulos $S_i$ para $m\leq i\leq m_0+h-1$ y de $S_hE$, considerada como subconjunto de $S_{[m]}$.
    Es claro que $S_m^\natural=S_0^\natural G_m$ para $1\leq m\leq m_0+h-1$, y por consiguiente la reunión $G$ de los $G_m$ para $1\leq m\leq m_0+h-1$ es un sistema de generadores de la $S_0^\natural$-álgebra $S^\natural$.
    Se concluye así que, si $S=S_0^\natural$ es un anillo noetheriano, también lo es $S^\natural$.
  \item[\rm{(ii)}] Es claro que, si $S_{n+1}=S_1S_n$ para $n\geq n_0$, entonces $S_{n+1}^\natural=S_1S_n^\natural$, y, \emph{a fortiori}, $S_{n+1}^\natural=S_1^\natural S_n^\natural$ para $n\geq n_0$.
    Recíprocamente, esta última igualdad puede escribirse
    \[
  S_{n+1}+S_{n+2}+\ldots=(S_1+S_2+\ldots)(S_n+S_{n+1}+\ldots)
\]
    y comparando los términos de grado $n+1$ (en $S$) en ambos miembros se obtiene $S_{n+1}=S_1S_n$.
  \item[\rm{(iii)}] Si $S_n=S_1^n$ para $n\geq n_0$, entonces $S_n^\natural=S_1^n+S_1^{n+1}+\ldots$;
    como $S_1^\natural$ contiene $S_1+S_1^2+\ldots$, se tiene $S_n^\natural\subset S_1^{\natural n}$, y por tanto $S_n^\natural=S_1^{\natural n}$ para $n\geq n_0$.
    Recíprocamente, los únicos términos de $S_1^{\natural n}=(S_1+S_2+\ldots)^n$ que son de grado $n$ en $S$ son los de $S_1^n$;
    la igualdad $S_n^\natural=S_1^{\natural n}$ implica, pues, que $S_n=S_1^n$.
  \item[\rm{(iv)}] Basta probar que, si un elemento $g\in S_{k+h}$ se considera como elemento de $S_k^\natural$ ($k>0$, $h\geq0$), existe un entero $n>0$ tal que en $S_{kn}^\natural$, $g^n$ es una combinación lineal de los $f_\alpha^\natural$ con coeficientes en $S^\natural$.
    Por hipótesis, existe un entero $m_0$ tal que, para $m\geq m_0$, se tiene, \emph{en $S$}, $g^m=\sum_\alpha c_{\alpha m}f_\alpha$, donde los índices $\alpha$ que figuran aquí son \emph{independientes de $m$};
    además, puede suponerse evidentemente que los $c_{\alpha m}$ son homogéneos, con
    \[
  \deg(c_{\alpha m})=m(k+h)-\deg f_\alpha
\]
    en $S$.
    Tomemos entonces $m_0$ suficientemente grande para que $km_0>\deg f_\alpha$ para todos los $f_\alpha$ que aparecen en $g^{m_0}$;
    para todo $\alpha$, sea $c'_{\alpha m}$ el elemento $c_{\alpha m}$ considerado como de grado $km-\deg(f_\alpha)$ \emph{en $S^\natural$};
    se tiene entonces, en $S^\natural$, $g^m=\sum_\alpha c'_{\alpha m}f_\alpha^\natural$, lo que termina la demostración.
\end{enumerate}
\end{proof}
% Fuente francesa: PMIHES_1961__8__5_0, printed p. 160 (PDF physical page 157).

\begin{env}[8.2.11]
\label{II.8.2.11-fr}
Consideremos la $S_0$-álgebra graduada
\begin{equation}
\label{II.8.2.11.1-fr}
  S^\natural\otimes_S S_0=S^\natural/S_+S^\natural
  =\bigoplus_{n\geq0}S_{[n]}/S_+S_{[n]}.
\tag{8.2.11.1}
\end{equation}

Como $S_n$ es un $S_0$-módulo cociente de $S_{[n]}/S_+S_{[n]}$, existe un homomorfismo canónico de $S_0$-álgebras graduadas
\begin{equation}
\label{II.8.2.11.2-fr}
  S^\natural\otimes_S S_0\to S
\tag{8.2.11.2}
\end{equation}
que es claramente \emph{sobreyectivo}, y corresponde por tanto \sref{II.2.9.2-fr} a una \emph{inmersión cerrada} canónica
\begin{equation}
\label{II.8.2.11.3-fr}
  \operatorname{Proj}(S)\to
  \operatorname{Proj}(S^\natural\otimes_S S_0)
\tag{8.2.11.3}
\end{equation}
\end{env}
% Fuente francesa: PMIHES_1961__8__5_0, printed p. 161 (PDF physical page 158).

\oldpage[II]{161}
\begin{proposition}[8.2.12]
\label{II.8.2.12-fr}
El morfismo canónico \sref{II.8.2.11.3-fr} es biyectivo.
Para que el homomorfismo \sref{II.8.2.11.2-fr} sea \textup{(TN)}-biyectivo, es necesario y suficiente que exista un $n_0$ tal que $S_{n+1}=S_1S_n$ para $n\geq n_0$.
Si se verifica esta última condición, \sref{II.8.2.11.3-fr} es un isomorfismo;
la recíproca es cierta cuando $S$ es noetheriano.
\end{proposition}

\begin{proof}
Para demostrar la primera afirmación, basta \sref{II.2.8.3-fr} probar que el núcleo $\mathfrak{J}$ del homomorfismo \sref{II.8.2.11.2-fr} está formado por elementos \emph{nilpotentes}.
Ahora bien, si $f\in S_{[n]}$ es un elemento cuya clase módulo $S_+S_{[n]}$ pertenece a este núcleo, esto significa que $f\in S_{[n+1]}$;
entonces $f^{n+1}$, considerado como elemento de $S_{[n(n+1)]}$, pertenece también a $S_+S_{[n(n+1)]}$, pues puede escribirse como $f.f^n$;
por tanto, la clase de $f^{n+1}$ módulo $S_+S_{[n(n+1)]}$ es nula, lo que prueba nuestra afirmación.
Como la hipótesis $S_{n+1}=S_1S_n$ para $n\geq n_0$ equivale a $S_{n+1}^{\natural}=S_1^{\natural}S_n^{\natural}$ para $n\geq n_0$ \sref{II.8.2.10-fr}[ii], esta hipótesis equivale, por definición, a que \sref{II.8.2.11.2-fr} sea (TN)-inyectivo, y por tanto (TN)-biyectivo, y entonces \sref{II.8.2.11.3-fr} es un isomorfismo, en virtud de \sref{II.2.9.1-fr}.
Recíprocamente, si \sref{II.8.2.11.3-fr} es un isomorfismo, el haz $\widetilde{\mathfrak{J}}$ sobre $\operatorname{Proj}(S^{\natural}\otimes_S S_0)$ es nulo \sref{II.2.9.2-fr}[i];
como $S^{\natural}\otimes_S S_0$ es noetheriano, por ser un cociente de $S^{\natural}$ \sref{II.8.2.10-fr}[i], se deduce de \sref{II.2.7.3-fr} que $\mathfrak{J}$ satisface la condición (TN), y por tanto $S_{n+1}^{\natural}=S_1^{\natural}S_n^{\natural}$ para $n\geq n_0$, lo que termina la demostración en virtud de \sref{II.8.2.10-fr}[ii].
\end{proof}
% Fuente francesa: PMIHES_1961__8__5_0, printed p. 161 (PDF physical page 158).

\begin{env}[8.2.13]
\label{II.8.2.13-fr}
Consideremos ahora las inyecciones canónicas $(S_+)^n\to S_{[n]}$, que definen un homomorfismo inyectivo de grado $0$ de anillos graduados
\begin{equation}
\label{II.8.2.13.1-fr}
  \bigoplus_{n\geq0}(S_+)^n\to S^{\natural}.
\tag{8.2.13.1}
\end{equation}
\end{env}
% Fuente francesa: PMIHES_1961__8__5_0, printed pp. 161--162 (PDF physical pages 158--159).

\begin{proposition}[8.2.14]
\label{II.8.2.14-fr}
Para que el homomorfismo \sref{II.8.2.13.1-fr} sea un \textup{(TN)}-isomorfismo, es necesario y suficiente que exista un $n_0$ tal que $S_n=S_1^n$ para todo $n\geq n_0$.
Cuando así ocurre, el morfismo correspondiente a \sref{II.8.2.13.1-fr} está definido en todas partes y es un isomorfismo
\[
  \operatorname{Proj}(S^{\natural})\to
  \operatorname{Proj}\left(\bigoplus_{n\geq0}(S_+)^n\right)~;
\]
la recíproca es cierta cuando $S$ es noetheriano.
\end{proposition}

\begin{proof}
Las dos primeras afirmaciones son evidentes, en virtud de \sref{II.8.2.10-fr}[iii] y \sref{II.2.9.1-fr}.
La tercera se deducirá de \sref{II.8.2.10-fr}[i y iii] y del lema siguiente:
\begin{lemma}[8.2.14.1]
\label{II.8.2.14.1-fr}
Sea $T$ un anillo graduado en grados positivos que sea además una $T_0$-álgebra de tipo finito.
Si el morfismo correspondiente al homomorfismo inyectivo $\bigoplus_{n\geq0}T_1^n\to T$ está definido en todas partes y es también un isomorfismo $\operatorname{Proj}(T)\to\operatorname{Proj}\left(\bigoplus_{n\geq0}T_1^n\right)$, existe un $n_0$ tal que $T_n=T_1^n$ para $n\geq n_0$.
\end{lemma}

Sean $g_i$ ($1\leq i\leq r$) generadores del $T_0$-módulo $T_1$.
La hipótesis implica ante todo que los $D_+(g_i)$ recubren $\operatorname{Proj}(T)$ \sref{II.2.8.1-fr}.
Sea $(h_j)_{1\leq j\leq s}$ un sistema de elementos homogéneos de $T_+$, con $\deg(h_j)=n_j$, que forman, junto con los $g_i$, un sistema de generadores del ideal $T_+$, o, lo que es equivalente \sref{II.2.1.3-fr}, un sistema de generadores de $T$ como $T_0$-álgebra;
si se pone $T'=\bigoplus_{n\geq0}T_1^n$, el elemento $h_j/g_i^{n_j}$ del anillo $T_{(g_i)}$ debe, por hipótesis, pertenecer al subanillo $T'_{(g_i)}$, y por tanto existe un entero $k$ tal que $T_1^k h_j\subset T_1^{k+n_j}$ para todo $j$.
Se concluye así, por inducción sobre el exponente $r$, que $T_1^k h_j^r\subset T'$ para todo $r\geq1$, y, por definición de los $h_j$, se tiene por tanto $T_1^kT\subset T'$.
Por otra parte, para todo $j$ existe un entero $m_j$ tal que $h_j^{m_j}$ pertenece al ideal de $T$ generado por los $g_i$ \sref{II.2.3.14-fr}, de modo que $h_j^{m_j}\in T_1T$, y
\oldpage[II]{162}
$h_j^{m_jk}\in T_1^kT\subset T'$.
Existe, pues, un entero $m_0\geq k$ tal que $h_j^m\in T_1^{mn_j}$ para $m\geq m_0$.
Así, si $q$ es el mayor de los enteros $n_j$, el número $n_0=qsm_0+k$ es el buscado.
En efecto, un elemento de $T_n$, para $n\geq n_0$, es suma de monomios pertenecientes a $T_1^\alpha u$, donde $u$ es un producto de potencias de los $h_j$;
si $\alpha\geq k$, se deduce de lo anterior que $T_1^\alpha u\subset T_1^n$;
en el caso contrario, uno de los exponentes de los $h_j$ es $\geq m_0$, por lo que $u\in T_1^\beta v$, donde $\beta\geq k$ y $v$ es de nuevo un producto de potencias de los $h_j$;
se reduce entonces al caso anterior, y se concluye que $T_1^\alpha u\subset T_1^n$ en todos los casos.
\end{proof}
% Fuente francesa: PMIHES_1961__8__5_0, printed p. 162 (PDF physical page 159).

\begin{remark}[8.2.15]
\label{II.8.2.15-fr}
La condición $S_n=S_1^n$ para $n\geq n_0$ implica evidentemente que $S_{n+1}=S_1S_n$ para $n\geq n_0$, pero la recíproca no es cierta, aun suponiendo que $S$ sea noetheriano.
Por ejemplo, sean $K$ un cuerpo, $A=K[\mathbf{x}]$ y $B=K[\mathbf{y}]/\mathbf{y}^2K[\mathbf{y}]$, donde $\mathbf{x}$ y $\mathbf{y}$ son indeterminadas, tomándose $\mathbf{x}$ de grado $1$ y $\mathbf{y}$ de grado $2$, y sea $S=A\otimes_KB$, de modo que $S$ es un álgebra graduada sobre $K$ que tiene una base formada por los elementos $1$, $\mathbf{x}^n$ ($n\geq1$) y $\mathbf{x}^n\mathbf{y}$ ($n\geq0$).
Es inmediato que $S_{n+1}=S_1S_n$ para $n\geq2$, pero $S_1^n=K\mathbf{x}^n$, mientras que $S_n=K\mathbf{x}^n+K\mathbf{x}^{n-2}\mathbf{y}$ para $n\geq2$.
\end{remark}
% Fuente francesa: PMIHES_1961__8__5_0, printed p. 162 (PDF physical page 159).

\subsection{Conos proyectantes}
\label{subsection:II.8.3-fr}

\begin{env}[8.3.1]
\label{II.8.3.1-fr}
Sea $Y$ un preesquema;
en todo este número, solo se considerarán \emph{$Y$-preesquemas} y \emph{$Y$-morfismos}.
Sea $\mathcal{S}$ una $\mathcal{O}_Y$-álgebra cuasicoherente graduada en grados \emph{positivos};
\emph{supondremos además que $\mathcal{S}_0=\mathcal{O}_Y$}.
Conforme a las notaciones introducidas en \sref{II.8.2.2-fr}, pondremos
\begin{equation}
\label{II.8.3.1.1-fr}
  \widehat{\mathcal{S}}=\mathcal{S}[\mathbf{z}]
  =\mathcal{S}\otimes_{\mathcal{O}_Y}\mathcal{O}_Y[\mathbf{z}]
\tag{8.3.1.1}
\end{equation}
que se considera como una $\mathcal{O}_Y$-álgebra graduada en grados positivos definiendo los grados mediante la fórmula \sref{II.8.2.2.2-fr}, de modo que, para todo abierto afín $U$ de $Y$, se tiene
\[
  \Gamma(U,\widehat{\mathcal{S}})=(\Gamma(U,\mathcal{S}))[\mathbf{z}].
\]
En lo que sigue, pondremos
\begin{equation}
\label{II.8.3.1.2-fr}
  X=\operatorname{Proj}(\mathcal{S}),\qquad
  C=\operatorname{Spec}(\mathcal{S}),\qquad
  \widehat{C}=\operatorname{Proj}(\widehat{\mathcal{S}})
\tag{8.3.1.2}
\end{equation}
(donde, en la definición de $C$, se considera $\mathcal{S}$ como una $\mathcal{O}_Y$-álgebra no graduada), y diremos que $C$ (resp. $\widehat{C}$) es el \emph{cono afín} (resp. el \emph{cono proyectivo}) definido por $\mathcal{S}$;
también se dirá «cono» en lugar de «cono afín».
Por abuso de lenguaje, diremos asimismo que $C$ (resp. $\widehat{C}$) es el \emph{cono proyectante afín de $X$} (resp. el \emph{cono proyectante proyectivo de $X$}), quedando sobrentendido que el preesquema $X$ está dado en la forma $\operatorname{Proj}(\mathcal{S})$;
por último, diremos que $\widehat{C}$ es la \emph{clausura proyectiva} de $C$ (quedando sobrentendido que el dato de $\mathcal{S}$ figura en la estructura de $C$).
\end{env}
\begin{proposition}[8.3.2]
\label{II.8.3.2-fr}
Existen $Y$-morfismos canónicos
\begin{equation}
\label{II.8.3.2.1-fr}
  Y\xrightarrow{\varepsilon}C\xrightarrow{i}\widehat{C}
\tag{8.3.2.1}
\end{equation}
\begin{equation}
\label{II.8.3.2.2-fr}
  X\xrightarrow[j]{}\widehat{C}
\tag{8.3.2.2}
\end{equation}

\oldpage[II]{163}
tales que $\varepsilon$ y $j$ son inmersiones cerradas, e $i$ es un morfismo afín que es una inmersión abierta dominante, para la cual
\begin{equation}
\label{II.8.3.2.3-fr}
  i(C)=\widehat{C}\setmin j(X)~;
\tag{8.3.2.3}
\end{equation}
Además, $\widehat{C}$ es el menor subpreesquema cerrado de $\widehat{C}$ que contiene a $i(C)$.
\end{proposition}

\begin{proof}
Para definir $i$, se considera el abierto de $\widehat{C}$ dado por
\begin{equation}
\label{II.8.3.2.4-fr}
  \widehat{C}_{\mathbf{z}}=
  \operatorname{Spec}(\widehat{\mathcal{S}}/(\mathbf{z}-1)\widehat{\mathcal{S}})
\tag{8.3.2.4}
\end{equation}
\sref{II.3.1.4-fr}, donde $\mathbf{z}$ se identifica canónicamente con una sección de $\widehat{\mathcal{S}}$ sobre $Y$.
El isomorfismo $i:C\xrightarrow{\sim}\widehat{C}_{\mathbf{z}}$ corresponde entonces al isomorfismo canónico \sref{II.8.2.3.1-fr}
\[
  \widehat{\mathcal{S}}/(\mathbf{z}-1)\widehat{\mathcal{S}}
  \xrightarrow{\sim}\mathcal{S}.
\]

El morfismo $\varepsilon$ corresponde al homomorfismo de aumentación $\mathcal{S}\to\mathcal{S}_0=\mathcal{O}_Y$, cuyo núcleo es $\mathcal{S}_+$ \sref{II.1.2.7-fr}, y, como este homomorfismo es sobreyectivo, $\varepsilon$ es una inmersión cerrada \sref{II.1.4.10-fr}.
Finalmente, $j$ corresponde \sref{II.3.5.1-fr} al homomorfismo sobreyectivo de grado $0$ $\widehat{\mathcal{S}}\to\mathcal{S}$ cuya restricción a $\mathcal{S}$ es la identidad y que es cero sobre $\mathbf{z}\widehat{\mathcal{S}}$, que es su núcleo;
$j$ está definido en todas partes y es una inmersión cerrada, por \sref{II.3.6.2-fr}.

Para demostrar las demás afirmaciones de \sref{II.8.3.2-fr}, se puede limitar evidentemente al caso en que $Y=\operatorname{Spec}(A)$ es afín y $\mathcal{S}=\widetilde{S}$, siendo $S$ una $A$-álgebra graduada, de donde $\widehat{\mathcal{S}}=(\widehat{S})^{\sim}$;
los elementos homogéneos $f$ de $S_+$ pueden identificarse entonces con secciones de $\widehat{\mathcal{S}}$ sobre $Y$, y el abierto de $\widehat{C}$ denotado por $D_+(f)$ en \sref{II.2.3.3-fr} puede escribirse $\widehat{C}_f$ \sref{II.3.1.4-fr};
de modo análogo, el abierto de $C$ denotado por $D(f)$ en \sref[I]{I.1.1.1-fr} puede escribirse $C_f$ \sref[0]{0.5.5.2-fr}.
Teniendo esto en cuenta, resulta de \sref{II.2.3.14-fr} y de la definición de $\widehat{S}$ que, en este caso, los abiertos $\widehat{C}_{\mathbf{z}}=i(C)$ y $\widehat{C}_f$ (con $f$ homogéneo en $S_+$) forman una \emph{cubierta} de $\widehat{C}$.
Además, con estas notaciones,
\begin{equation}
\label{II.8.3.2.5-fr}
  i^{-1}(\widehat{C}_f)=C_f~;
\tag{8.3.2.5}
\end{equation}
en efecto, $\widehat{C}_f\cap i(C)=\widehat{C}_f\cap\widehat{C}_{\mathbf{z}}
=\widehat{C}_{f\mathbf{z}}=\operatorname{Spec}(\widehat{S}_{(f\mathbf{z})})$.
Pero, si $d=\deg(f)$, entonces $\widehat{S}_{(f\mathbf{z})}$ es canónicamente isomorfo a $(\widehat{S}_{(\mathbf{z})})_{f/\mathbf{z}^d}$ \sref{II.2.2.2-fr}, y de la definición del isomorfismo de \sref{II.8.2.3.1-fr} resulta que la imagen de $(\widehat{S}_{(\mathbf{z})})_{f/\mathbf{z}^d}$ por el isomorfismo correspondiente de anillos de fracciones es exactamente $S_f$.
Como $C_f=\operatorname{Spec}(S_f)$, esto prueba \sref{II.8.3.2.5-fr} y muestra al mismo tiempo que el morfismo $i$ es afín;
además, la restricción de $i$ a $C_f$, considerada como morfismo hacia $\widehat{C}_f$, corresponde \sref[I]{I.1.7.3-fr} al homomorfismo canónico $\widehat{S}_{(f)}\to\widehat{S}_{(f\mathbf{z})}$, y, por lo anterior y por \sref{II.8.2.3.2-fr}, se puede enunciar el resultado siguiente:
\begin{env}[8.3.2.6]
\label{II.8.3.2.6-fr}
\emph{Si $Y=\operatorname{Spec}(A)$ es afín y $\mathcal{S}=\widetilde{S}$, entonces, para todo $f$ homogéneo en $S_+$, $\widehat{C}_f$ se identifica canónicamente con $\operatorname{Spec}(S_f^{\leq})$, y el morfismo $C_f\to\widehat{C}_f$ obtenido restringiendo $i$ corresponde entonces a la inyección canónica $S_f^{\leq}\to S_f$}.
\end{env}

Observemos ahora que (para $Y$ afín) el complemento de $\widehat{C}_{\mathbf{z}}$ en $\widehat{C}=\operatorname{Proj}(\widehat{S})$
\oldpage[II]{164}
es, por definición, el conjunto de ideales primos graduados de $\widehat{S}$ que contienen a $\mathbf{z}$, que es precisamente $j(X)$, por la definición de $j$, lo que prueba \sref{II.8.3.2.3-fr}.

Finalmente, para demostrar la última afirmación de \sref{II.8.3.2-fr}, se puede suponer que $Y$ es afín.
Con las notaciones anteriores, observemos que, en el anillo $\widehat{S}$, $\mathbf{z}$ no es divisor de $0$;
como $i(C)=\widehat{C}_{\mathbf{z}}$, basta demostrar el lema siguiente:
\begin{lemma}[8.3.2.7]
\label{II.8.3.2.7-fr}
Sea $T$ un anillo graduado en grados positivos, $Z=\operatorname{Proj}(T)$ y $g$ un elemento homogéneo de $T$ de grado $d>0$.
Si $g$ no es divisor de $0$ en $T$, entonces $Z$ es el menor subpreesquema cerrado de $Z$ que contiene a $Z_g=D_+(g)$.
\end{lemma}

Por \sref[I]{I.4.1.9-fr}, la cuestión es local sobre $Z$;
para todo elemento homogéneo $h\in T_e$ ($e>0$), basta, pues, demostrar que $Z_h$ es el menor subpreesquema cerrado de $Z_h$ que contiene a $Z_{gh}$;
de las definiciones y de \sref[I]{I.4.3.2-fr} resulta que esta condición equivale a exigir que el homomorfismo canónico $T_{(h)}\to T_{(gh)}$ sea \emph{inyectivo}.
Pero este homomorfismo puede identificarse con el homomorfismo canónico $T_{(h)}\to(T_{(h)})_{g^e/h^d}$ \sref{II.2.2.3-fr}.
Ahora bien, como $g^e$ no es divisor de $0$ en $T$, $g^e/h^d$ no lo es en $T_h$ (ni, \emph{a fortiori}, en $T_{(h)}$), pues el hecho de que $(g^e/h^d)(t/h^m)=0$ (para $t\in T$ y $m>0$) implica la existencia de un $n>0$ tal que $h^ng^et=0$, de donde $h^nt=0$, y por tanto $t/h^m=0$ en $T_h$.
Esto termina la demostración \sref[0]{0.1.2.2-fr}.
\end{proof}
\begin{env}[8.3.3]
\label{II.8.3.3-fr}
Se identificará a menudo el cono afín $C$ con el subpreesquema inducido por el cono proyectivo $\widehat{C}$ sobre el abierto $i(C)$ mediante la inmersión abierta $i$.
El subpreesquema cerrado de $C$ asociado a la inmersión cerrada $\varepsilon$ se llama el \emph{preesquema de los vértices} de $C$;
también se dice que $\varepsilon$, que es una $Y$-sección de $C$, es la \emph{sección vértice}, o la \emph{sección nula}, de $C$;
se puede identificar $Y$ con el preesquema de los vértices de $C$ mediante $\varepsilon$.
Además, $i\circ\varepsilon$ es una $Y$-sección de $\widehat{C}$, y por tanto también una inmersión cerrada \sref[I]{I.5.4.6-fr}, correspondiente al homomorfismo canónico sobreyectivo de grado $0$ $\widehat{\mathcal{S}}=\mathcal{S}[\mathbf{z}]\to\mathcal{O}_Y[\mathbf{z}]$ \sref{II.3.1.7-fr}, cuyo núcleo es $\mathcal{S}_+[\mathbf{z}]=\mathcal{S}_+\widehat{\mathcal{S}}$;
el subpreesquema de $\widehat{C}$ asociado a esta inmersión cerrada se llama también el \emph{preesquema de los vértices} de $\widehat{C}$, e $i\circ\varepsilon$ la \emph{sección vértice} de $\widehat{C}$;
puede identificarse mediante $i\circ\varepsilon$ con $Y$.
Por último, el subpreesquema cerrado de $\widehat{C}$ asociado a $j$ se llama el \emph{lugar en el infinito} de $\widehat{C}$, y puede identificarse con $X$ mediante $j$.
\end{env}

\begin{env}[8.3.4]
\label{II.8.3.4-fr}
Los subpreesquemas de $C$ (resp. $\widehat{C}$) inducidos sobre los \emph{abiertos}
\begin{equation}
\label{II.8.3.4.1-fr}
  E=C\setmin\varepsilon(Y),\qquad
  \widehat{E}=\widehat{C}\setmin i(\varepsilon(Y))
\tag{8.3.4.1}
\end{equation}
se llaman, respectivamente (por abuso de lenguaje), el \emph{cono afín sin vértice} y el \emph{cono proyectivo sin vértice} definidos por $\mathcal{S}$;
se observa que, pese a esta terminología, \emph{$E$ no es necesariamente afín} sobre $Y$, ni $\widehat{E}$ proyectivo sobre $Y$ \sref{II.8.4.3-fr}.
Cuando se identifica $C$ con $i(C)$, se tiene, pues, para los espacios subyacentes
\begin{equation}
\label{II.8.3.4.2-fr}
  C\cup\widehat{E}=\widehat{C},\qquad C\cap\widehat{E}=E
\tag{8.3.4.2}
\end{equation}
de modo que $\widehat{C}$ puede considerarse obtenido por \emph{pegado} de los subpreesquemas abiertos $C$ y $\widehat{E}$;
además, por \sref{II.8.3.2.3-fr},
\begin{equation}
\label{II.8.3.4.3-fr}
  E=\widehat{E}\setmin j(X).
\tag{8.3.4.3}
\end{equation}

\oldpage[II]{165}
Si $Y=\operatorname{Spec}(A)$ es afín, entonces, con las notaciones de \sref{II.8.3.2-fr},
\begin{equation}
\label{II.8.3.4.4-fr}
  E=\bigcup C_f,\qquad
  \widehat{E}=\bigcup\widehat{C}_f,\qquad
  C_f=C\cap\widehat{C}_f
\tag{8.3.4.4}
\end{equation}
donde $f$ recorre el conjunto de los elementos homogéneos de $S_+$ (o solamente un subconjunto $M$ de este conjunto que genere un ideal de $S_+$ cuyo radical en $S_+$ sea el propio $S_+$, o, equivalentemente, tal que los $X_f$, para $f\in M$, recubran $X$ \sref{II.2.3.14-fr}).
El pegado de $C$ y $\widehat{C}_f$ a lo largo de $C_f$ queda entonces determinado por los morfismos de inyección $C_f\to C$ y $C_f\to\widehat{C}_f$, que, como se ha visto \sref{II.8.3.2.6-fr}, corresponden respectivamente a los homomorfismos canónicos $S\to S_f$ y $S_f^{\leq}\to S_f$.
\end{env}
\begin{proposition}[8.3.5]
\label{II.8.3.5-fr}
Con las notaciones de \sref{II.8.3.1-fr} y \sref{II.8.3.4-fr}, el morfismo asociado \sref{II.3.5.1-fr} a la inyección canónica $\varphi:\mathcal{S}\to\widehat{\mathcal{S}}=\mathcal{S}[\mathbf{z}]$ es un morfismo afín sobreyectivo \textup{(llamado} retracción canónica)
\begin{equation}
\label{II.8.3.5.1-fr}
  p:\widehat{E}\to X
\tag{8.3.5.1}
\end{equation}
tal que
\begin{equation}
\label{II.8.3.5.2-fr}
  p\circ j=1_X.
\tag{8.3.5.2}
\end{equation}
\end{proposition}

\begin{proof}
Para demostrar la proposición, se puede limitar al caso en que $Y$ es afín.
Teniendo en cuenta la expresión \sref{II.8.3.4.4-fr} de $\widehat{E}$, el hecho de que el dominio de definición $G(\varphi)$ de $p$ sea igual a $\widehat{E}$ resultará de la primera de las afirmaciones siguientes:
\begin{env}[8.3.5.3]
\label{II.8.3.5.3-fr}
\itshape Si $Y=\operatorname{Spec}(A)$ es afín y $\mathcal{S}=\widetilde{S}$, entonces, para todo $f\in S_+$ homogéneo,
\begin{equation}
\label{II.8.3.5.4-fr}
  p^{-1}(X_f)=\widehat{C}_f
\tag{8.3.5.4}
\end{equation}
y la restricción de $p$ a $\widehat{C}_f=\operatorname{Spec}(S_f^{\leq})$, considerada como morfismo de $\widehat{C}_f$ a $X_f$, corresponde a la inyección canónica $S_{(f)}\to S_f^{\leq}$.
Si, además, $f\in S_1$, entonces $\widehat{C}_f$ es isomorfo a $X_f\otimes_{\mathbf{Z}}\mathbf{Z}[T]$ (donde $T$ es una indeterminada).
\end{env}

En efecto, la fórmula~\sref{II.8.3.5.4-fr} no es más que un caso particular de \sref{II.2.8.1.1-fr}, y la segunda afirmación no es otra que la definición de $\operatorname{Proj}(\varphi)$ cuando $Y$ es afín \sref{II.2.8.1-fr}.
La fórmula~\sref{II.8.3.5.2-fr} y la sobreyectividad de $p$ resultan entonces de que el compuesto $\mathcal{S}\to\widehat{\mathcal{S}}\to\mathcal{S}$ de los homomorfismos canónicos es la identidad en $\mathcal{S}$.
Finalmente, la última afirmación de \sref{II.8.3.5.3-fr} resulta de que $S_f^{\leq}$ es isomorfo a $S_{(f)}[T]$ cuando $f\in S_1$ \sref{II.2.2.1-fr}.
\end{proof}
\begin{corollary}[8.3.6]
\label{II.8.3.6-fr}
La restricción
\begin{equation}
\label{II.8.3.6.1-fr}
  \pi:E\to X
\tag{8.3.6.1}
\end{equation}
de $p$ a $E$ es un morfismo afín sobreyectivo.
Si $Y$ es afín y $f$ homogéneo en $S_+$, entonces
\begin{equation}
\label{II.8.3.6.2-fr}
  \pi^{-1}(X_f)=C_f
\tag{8.3.6.2}
\end{equation}
y la restricción de $\pi$ a $C_f$ corresponde a la inyección canónica $S_{(f)}\to S_f$.
Si, además, $f\in S_1$, entonces $C_f$ es isomorfo a $X_f\otimes_{\mathbf{Z}}\mathbf{Z}[T,T^{-1}]$ (donde $T$ es una indeterminada).
\end{corollary}

\begin{proof}
La fórmula~\sref{II.8.3.6.2-fr} resulta inmediatamente de \sref{II.8.3.5.3-fr} y \sref{II.8.3.2.5-fr}, y demuestra la sobreyectividad de $\pi$;
ya se ha visto que la inmersión $i$, restringida a $C_f$, corresponde
\oldpage[II]{166}
a la inyección $S_f^{\leq}\to S_f$ \sref{II.8.3.2-fr}.
Finalmente, la última afirmación es consecuencia de que, para $f\in S_1$, $S_f$ es isomorfo a $S_{(f)}[T,T^{-1}]$ \sref{II.2.2.1-fr}.
\end{proof}
\begin{remark}[8.3.7]
\label{II.8.3.7-fr}
Cuando $Y$ es afín, los elementos del espacio subyacente a $E$ son los ideales primos $\mathfrak{p}$ de $S$ (no necesariamente graduados) que no contienen a $S_+$, por la definición de la inmersión $\varepsilon$ \sref{II.8.3.2-fr}.
Para un tal ideal $\mathfrak{p}$, los $\mathfrak{p}\cap S_n$ satisfacen evidentemente las condiciones de \sref{II.2.1.9-fr}, por lo que existe un único ideal primo \emph{graduado} $\mathfrak{q}$ de $S$ tal que $\mathfrak{q}\cap S_n=\mathfrak{p}\cap S_n$ para todo $n$;
la aplicación $\pi:E\to X$ de los espacios subyacentes se interpreta entonces mediante la relación
\begin{equation}
\label{II.8.3.7.1-fr}
  \pi(\mathfrak{p})=\mathfrak{q}.
\tag{8.3.7.1}
\end{equation}

En efecto, para demostrar esta relación, basta considerar un elemento $f$ homogéneo en $S_+$ tal que $\mathfrak{p}\in D(f)$, y observar que $\mathfrak{q}_{(f)}$ es la imagen inversa de $\mathfrak{p}_f$ por la inyección $S_{(f)}\to S_f$.
\end{remark}

\begin{corollary}[8.3.8]
\label{II.8.3.8-fr}
Si $\mathcal{S}$ está generada por $\mathcal{S}_1$, los morfismos $p$ y $\pi$ son de tipo finito;
para todo $x\in X$, la fibra $p^{-1}(x)$ es isomorfa a $\operatorname{Spec}(\boldsymbol{k}(x)[T])$, y la fibra $\pi^{-1}(x)$ es isomorfa a $\operatorname{Spec}(\boldsymbol{k}(x)[T,T^{-1}])$.
\end{corollary}

\begin{proof}
Esto resulta inmediatamente de \sref{II.8.3.5-fr} y \sref{II.8.3.6-fr}, observando que, cuando $Y$ es afín y $S$ está generada por $S_1$, los $X_f$, para $f\in S_1$, forman una cubierta de $X$ \sref{II.2.3.14-fr}.
\end{proof}
\begin{remark}[8.3.9]
\label{II.8.3.9-fr}
El cono afín sin vértice correspondiente a la $\mathcal{O}_Y$-álgebra graduada $\mathcal{O}_Y[T]$ (donde $T$ es una indeterminada) se identifica con $G_m=\operatorname{Spec}(\mathcal{O}_Y[T,T^{-1}])$, pues no es otro que $C_T$, como se ha visto en \sref{II.8.3.2-fr} (véase \sref{II.8.4.4-fr} para un resultado más general).
Este preesquema está provisto canónicamente de una estructura de \textquotedblleft\emph{$Y$-esquema en grupos conmutativos}\textquotedblright.
Esta noción se desarrollará en detalle más adelante, pero puede resumirse aquí brevemente como sigue.
Un $Y$-esquema en grupos es un $Y$-esquema $G$ provisto de dos $Y$-morfismos, $p:G\times_YG\to G$ y $s:G\to G$, que satisfacen condiciones formalmente análogas a los axiomas de la ley de composición y de la aplicación de inversión de un grupo: el diagrama
\[
  \xymatrix{
    G\times G\times G \ar[r]^{p\times1}\ar[d]_{1\times p}
      & G\times G \ar[d]^{p} \\
    G\times G \ar[r]_{p} & G
  }
\]
debe ser conmutativo (\textquotedblleft asociatividad\textquotedblright), y debe cumplirse una condición que corresponda al hecho de que, en los grupos, las aplicaciones
\[
  (x,y)\to(x,x^{-1},y)\to(x,x^{-1}y)\to x(x^{-1}y)
\]
y
\[
  (x,y)\to(x,x^{-1},y)\to(x,yx^{-1})\to(yx^{-1})x
\]
deben reducirse ambas a $(x,y)\to y$;
la sucesión de morfismos que corresponde, por ejemplo, a la primera aplicación compuesta es
\[
  G\times G\xrightarrow{(1,s)\times1}G\times G\times G
  \xrightarrow{1\times p}G\times G\xrightarrow{p}G
\]
y el lector escribirá de manera análoga la segunda sucesión.

\oldpage[II]{167}
Es inmediato \sref[I]{I.3.4.3-fr} que dar una estructura de $Y$-esquema en grupos sobre un $Y$-esquema $G$ equivale a dar, para \emph{todo} $Y$-preesquema $Z$, una estructura de \emph{grupo} sobre el conjunto $\operatorname{Hom}_Y(Z,G)$, de modo que, para todo $Y$-morfismo $Z\to Z'$, la aplicación correspondiente $\operatorname{Hom}_Y(Z',G)\to\operatorname{Hom}_Y(Z,G)$ sea un homomorfismo de grupos.
En el caso particular de $G_m$ que consideramos aquí, $\operatorname{Hom}_Y(Z,G)$ se identifica con el conjunto de las $Z$-secciones de $Z\times_YG_m$ \sref[I]{I.3.3.14-fr}, y por tanto con el conjunto de las $Z$-secciones de $\operatorname{Spec}(\mathcal{O}_Z[T,T^{-1}])$;
finalmente, el mismo razonamiento que en \sref[I]{I.3.3.15-fr} muestra que este conjunto se identifica canónicamente con el conjunto de los elementos \emph{invertibles} del anillo $\Gamma(Z,\mathcal{O}_Z)$, y la estructura de grupo sobre este conjunto proviene de la multiplicación del anillo $\Gamma(Z,\mathcal{O}_Z)$.
El lector comprobará que los morfismos $p$ y $s$ anteriores se obtienen como sigue: corresponden, por \sref{II.1.2.7-fr} y \sref{II.1.4.6-fr}, a los homomorfismos de $\mathcal{O}_Y$-álgebras
\begin{align*}
  \pi &: \mathcal{O}_Y[T,T^{-1}]
  \to\mathcal{O}_Y[T,T^{-1},T',T'^{-1}]\\
  \sigma &: \mathcal{O}_Y[T,T^{-1}]\to\mathcal{O}_Y[T,T^{-1}]
\end{align*}
y están completamente definidos por los datos $\pi(T)=TT'$ y $\sigma(T)=T^{-1}$.

Dicho esto, $G_m$ puede considerarse como un \textquotedblleft\emph{dominio universal de operadores}\textquotedblright\ para todo \emph{cono afín} $C=\operatorname{Spec}(\mathcal{S})$, donde $\mathcal{S}$ es una $\mathcal{O}_Y$-álgebra cuasicoherente graduada en grados positivos.
Esto significa que se puede definir canónicamente un $Y$-morfismo $G_m\times_YC\to C$ que tenga las propiedades formales de una ley externa de un conjunto provisto de un grupo de operadores;
o, como antes para los esquemas en grupos, que se puede dar, para todo $Y$-preesquema $Z$, una ley externa sobre $\operatorname{Hom}_Y(Z,C)$ cuyo conjunto de operadores sea el grupo $\operatorname{Hom}_Y(Z,G_m)$, con los axiomas usuales de los conjuntos provistos de un grupo de operadores y una condición de compatibilidad respecto de los $Y$-morfismos $Z\to Z'$.
En el caso presente, el morfismo $G_m\times_YC\to C$ se define dando un homomorfismo de $\mathcal{O}_Y$-álgebras $\mathcal{S}\to\mathcal{S}\otimes_{\mathcal{O}_Y}\mathcal{O}_Y[T,T^{-1}]=\mathcal{S}[T,T^{-1}]$, que a toda sección $s_n\in\Gamma(U,\mathcal{S}_n)$ (donde $U$ es un abierto de $Y$) asocia la sección $s_nT^n\in\Gamma(U,\mathcal{S}\otimes_{\mathcal{O}_Y}
\mathcal{O}_Y[T,T^{-1}])$.

Recíprocamente, supongamos dada una $\mathcal{O}_Y$-álgebra cuasicoherente $\mathcal{S}$, \emph{a priori no graduada}, y, sobre $C=\operatorname{Spec}(\mathcal{S})$, una estructura de \textquotedblleft\emph{$Y$-esquema en conjuntos provistos de un grupo de operadores}\textquotedblright\ cuyo dominio de operadores sea el $Y$-esquema en grupos $G_m$;
se obtiene entonces canónicamente una \emph{graduación} de $\mathcal{O}_Y$-álgebra sobre $\mathcal{S}$.
En efecto, dar un $Y$-morfismo $G_m\times_YC\to C$ equivale a dar un homomorfismo de $\mathcal{O}_Y$-álgebras $\psi:\mathcal{S}\to\mathcal{S}[T,T^{-1}]$, que se escribe $\psi=\sum_{n\in\mathbf{Z}}\psi_nT^n$, donde los $\psi_n:\mathcal{S}\to\mathcal{S}$ son homomorfismos de $\mathcal{O}_Y$-módulos (para toda sección $s\in\Gamma(U,\mathcal{S})$, $U$ abierto de $Y$, y todo punto $x$ de $U$, existe una vecindad abierta $V$ de $x$ en $U$ tal que $\psi_n(s)|_V=0$ salvo para un número finito de índices $n$).
Se comprueba entonces que los axiomas de los conjuntos provistos de un grupo de operadores implican que los $\psi_n(\mathcal{S})=\mathcal{S}_n$ definen una graduación (en grados enteros arbitrarios) de $\mathcal{O}_Y$-álgebra sobre $\mathcal{S}$, siendo los $\psi_n$ los proyectores correspondientes.
Se obtiene así una noción de estructura de \textquotedblleft\emph{cono afín}\textquotedblright\ sobre todo $Y$-esquema afín, definida de manera \textquotedblleft geométrica\textquotedblright\ sin recurrir a una graduación previa.
\oldpage[II]{168}
No desarrollaremos más este punto de vista aquí, y dejamos al lector el cuidado de formular con precisión las definiciones y los resultados correspondientes a las indicaciones anteriores.
\end{remark}
\subsection{Clausura proyectiva de un fibrado vectorial}
\label{subsection:II.8.4-fr}

\begin{env}[8.4.1]
\label{II.8.4.1-fr}
Sean $Y$ un preesquema y $\mathcal{E}$ un $\mathcal{O}_Y$-módulo cuasicoherente.
Si se toma como $\mathcal{S}$ la $\mathcal{O}_Y$-álgebra graduada $\mathbf{S}_{\mathcal{O}_Y}(\mathcal{E})$, la definición~\sref{II.8.3.1.1-fr} muestra que $\widehat{\mathcal{S}}$ se identifica con $\mathbf{S}_{\mathcal{O}_Y}(\mathcal{E}\oplus\mathcal{O}_Y)$.
Como el cono afín $\operatorname{Spec}(\mathcal{S})$ definido por $\mathcal{S}$ es, por definición, $\mathbf{V}(\mathcal{E})$, y $\operatorname{Proj}(\mathcal{S})$ es, por definición, $\mathbf{P}(\mathcal{E})$, se obtiene:
\end{env}
\begin{proposition}[8.4.2]
\label{II.8.4.2-fr}
La clausura proyectiva de un fibrado vectorial $\mathbf{V}(\mathcal{E})$ sobre $Y$ es canónicamente isomorfa a $\mathbf{P}(\mathcal{E}\oplus\mathcal{O}_Y)$, y el lugar en el infinito de este último es canónicamente isomorfo a $\mathbf{P}(\mathcal{E})$.
\end{proposition}
\begin{remark}[8.4.3]
\label{II.8.4.3-fr}
Tomemos, por ejemplo, $\mathcal{E}=\mathcal{O}_Y^r$ con $r\geq2$;
los conos sin vértice $E$ y $\widehat{E}$ definidos por $\mathcal{S}$ no son entonces ni afines ni proyectivos sobre $Y$ si $Y\neq\varnothing$.
La segunda afirmación es inmediata, pues $\widehat{C}=\mathbf{P}(\mathcal{O}_Y^{r+1})$ es proyectivo sobre $Y$, y los espacios subyacentes a $E$ y $\widehat{E}$ son abiertos no cerrados de $\widehat{C}$; por tanto, las inmersiones canónicas $E\to\widehat{C}$ y $\widehat{E}\to\widehat{C}$ no son proyectivas \sref{II.5.5.3-fr}, y se concluye por \sref{II.5.5.5-fr}[v].
Por otra parte, suponiendo $Y=\operatorname{Spec}(A)$ afín y, por ejemplo, $r=2$, se tiene $C=\operatorname{Spec}(A[T_1,T_2])$, y $E$ es el preesquema inducido por $C$ sobre el abierto $D(T_1)\cup D(T_2)$;
pero ya hemos visto que este último no es afín \sref[I]{I.5.5.11-fr};
\emph{a fortiori}, $\widehat{E}$ no puede ser afín, pues $E$ es el abierto donde no se anula la sección $\mathbf{z}$ sobre $\widehat{E}$ \sref{II.8.3.2-fr}.

Sin embargo:
\end{remark}
\begin{proposition}[8.4.4]
\label{II.8.4.4-fr}
Si $\mathcal{L}$ es un $\mathcal{O}_Y$-módulo invertible, existen isomorfismos canónicos para los conos sin vértice $E$ y $\widehat{E}$ correspondientes a $C=\mathbf{V}(\mathcal{L})$:
\begin{equation}
\label{II.8.4.4.1-fr}
  \operatorname{Spec}\left(
    \bigoplus_{n\in\mathbf{Z}}\mathcal{L}^{\otimes n}
  \right)\xrightarrow{\sim}E
\tag{8.4.4.1}
\end{equation}
\begin{equation}
\label{II.8.4.4.2-fr}
  \mathbf{V}(\mathcal{L}^{-1})\xrightarrow{\sim}\widehat{E}.
\tag{8.4.4.2}
\end{equation}

Además, existe un isomorfismo canónico de la clausura proyectiva de $\mathbf{V}(\mathcal{L})$ sobre la de $\mathbf{V}(\mathcal{L}^{-1})$ que transforma la sección nula \textup{(resp.} el lugar en el infinito\textup{)} de la primera en el lugar en el infinito \textup{(resp.} la sección nula\textup{)} de la segunda.
\end{proposition}

\begin{proof}
Aquí se tiene $\mathcal{S}=\bigoplus_{n\geq0}\mathcal{L}^{\otimes n}$;
la inyección canónica
\[
  \mathcal{S}\to\bigoplus_{n\in\mathbf{Z}}\mathcal{L}^{\otimes n}
\]
define un morfismo dominante canónico
\begin{equation}
\label{II.8.4.4.3-fr}
  \operatorname{Spec}\left(
    \bigoplus_{n\in\mathbf{Z}}\mathcal{L}^{\otimes n}
  \right)
  \to\mathbf{V}(\mathcal{L})
  =\operatorname{Spec}\left(
    \bigoplus_{n\geq0}\mathcal{L}^{\otimes n}
  \right)
\tag{8.4.4.3}
\end{equation}
y basta demostrar que este morfismo es un isomorfismo del esquema $\operatorname{Spec}(\bigoplus_{n\in\mathbf{Z}}\mathcal{L}^{\otimes n})$ sobre $E$.
Como la cuestión es local sobre $Y$, se puede suponer que $Y=\operatorname{Spec}(A)$ es afín
\oldpage[II]{169}
y que $\mathcal{L}=\mathcal{O}_Y$, de modo que $\mathcal{S}=(A[T])^{\sim}$ y $\bigoplus_{n\in\mathbf{Z}}\mathcal{L}^{\otimes n}
=(A[T,T^{-1}])^{\sim}$.
Pero $A[T,T^{-1}]$ es el anillo de fracciones $A[T]_T$ de $A[T]$, y por tanto \eref{II.8.4.4.3-fr} identifica el primer miembro con el preesquema inducido por $C=\mathbf{V}(\mathcal{L})$ sobre el abierto $D(T)$;
el complementario $V(T)$ de este abierto en $C$ es el espacio subyacente al subpreesquema cerrado de $C$ definido por el ideal $TA[T]$, que es precisamente la sección nula de $C$, de donde $E=D(T)$.

Por otra parte, el isomorfismo \eref{II.8.4.4.2-fr} será consecuencia de la última afirmación, pues $\mathbf{V}(\mathcal{L}^{-1})$ es el complementario del lugar en el infinito de su clausura proyectiva, y $\widehat{E}$ es el complementario de la sección nula de la clausura proyectiva de $C=\mathbf{V}(\mathcal{L})$.
Estas clausuras proyectivas son, respectivamente, $\mathbf{P}(\mathcal{L}^{-1}\oplus\mathcal{O}_Y)$ y $\mathbf{P}(\mathcal{L}\oplus\mathcal{O}_Y)$;
pero se puede escribir $\mathcal{L}\oplus\mathcal{O}_Y
=\mathcal{L}\otimes(\mathcal{L}^{-1}\oplus\mathcal{O}_Y)$.
La existencia del isomorfismo canónico buscado resulta entonces de \sref{II.4.1.4-fr}, y todo se reduce a comprobar que este isomorfismo intercambia las secciones nulas y los lugares en el infinito.
Para ello, se puede reducir al caso en que $Y=\operatorname{Spec}(A)$ es afín, $L=Ac$ y $L^{-1}=Ac'$, con el isomorfismo canónico $L\otimes L^{-1}\to A$ que envía $c\otimes c'$ al elemento $1$ de $A$.
Entonces $\mathbf{S}(L\oplus A)$ es el producto tensorial de $A[\mathbf{z}]$ con $\bigoplus_{n\geq0}Ac^{\otimes n}$, y $\mathbf{S}(L^{-1}\oplus A)$ es el producto tensorial de $A[\mathbf{z}]$ con $\bigoplus_{n\geq0}Ac'^{\otimes n}$, y el isomorfismo definido en \sref{II.4.1.4-fr} envía $\mathbf{z}^h\otimes c'^{\otimes(n-h)}$ al elemento $\mathbf{z}^{n-h}\otimes c^{\otimes h}$.
Ahora bien, en $\mathbf{P}(\mathcal{L}^{-1}\oplus\mathcal{O}_Y)$, el lugar en el infinito es el conjunto de los puntos donde se anula la sección $\mathbf{z}$, y la sección nula es el conjunto de los puntos donde se anula la sección $c'$;
como se tienen definiciones análogas para $\mathbf{P}(\mathcal{L}\oplus\mathcal{O}_Y)$, nuestra conclusión resulta inmediatamente de la descripción anterior.
\end{proof}
\subsection{Comportamiento funtorial}
\label{subsection:II.8.5-fr}

\begin{env}[8.5.1]
\label{II.8.5.1-fr}
Sean $Y$ y $Y'$ preesquemas, $q:Y'\to Y$ un morfismo, y $\mathcal{S}$ (resp. $\mathcal{S}'$) una $\mathcal{O}_Y$-álgebra cuasicoherente graduada en grados \emph{positivos} (resp. una $\mathcal{O}_{Y'}$-álgebra cuasicoherente graduada en grados \emph{positivos}).
Consideremos un $q$-morfismo de álgebras graduadas
\begin{equation}
\label{II.8.5.1.1-fr}
  \varphi:\mathcal{S}\to\mathcal{S}'.
\tag{8.5.1.1}
\end{equation}

Se sabe \sref{II.1.5.6-fr} que le corresponde canónicamente un morfismo
\[
  \Phi=\operatorname{Spec}(\varphi):
  \operatorname{Spec}(\mathcal{S}')\to\operatorname{Spec}(\mathcal{S})
\]
tal que el diagrama
\begin{equation}
\label{II.8.5.1.2-fr}
  \xymatrix{
    C'\ar[r]^{\Phi}\ar[d] & C\ar[d]\\
    Y'\ar[r]_{q} & Y
  }
\tag{8.5.1.2}
\end{equation}
es conmutativo, donde se escribe $C=\operatorname{Spec}(\mathcal{S})$ y $C'=\operatorname{Spec}(\mathcal{S}')$.
\emph{Supongamos además que $\mathcal{S}_0=\mathcal{O}_Y$ y $\mathcal{S}'_0=\mathcal{O}_{Y'}$};
sean $\varepsilon:Y\to C$ y $\varepsilon':Y'\to C'$ las inmersiones canónicas \sref{II.8.3.2-fr};
se tiene entonces el diagrama conmutativo
\begin{equation}
\label{II.8.5.1.3-fr}
  \xymatrix{
    Y'\ar[r]^{q}\ar[d]_{\varepsilon'} & Y\ar[d]^{\varepsilon}\\
    C'\ar[r]_{\Phi} & C
  }
\tag{8.5.1.3}
\end{equation}
\oldpage[II]{170}
que corresponde al diagrama
\[
  \xymatrix{
    \mathcal{S}\ar[r]^{\varphi}\ar[d] & \mathcal{S}'\ar[d]\\
    \mathcal{O}_Y\ar[r] & \mathcal{O}_{Y'}
  }
\]
donde las flechas verticales son los homomorfismos de aumentación; la conmutatividad resulta, por tanto, de que $\varphi$ es un homomorfismo de álgebras \emph{graduadas}.
\end{env}
\begin{proposition}[8.5.2]
\label{II.8.5.2-fr}
Si $E$ \textup{(resp. $E'$)} es el cono afín sin vértice definido por $\mathcal{S}$ \textup{(resp. $\mathcal{S}'$)}, se tiene $\Phi^{-1}(E)\subset E'$;
si además $\operatorname{Proj}(\varphi):G(\varphi)\to\operatorname{Proj}(\mathcal{S})$ está definido en todas partes (o, equivalentemente, si $G(\varphi)=\operatorname{Proj}(\mathcal{S}')$), entonces $\Phi^{-1}(E)=E'$, y recíprocamente.
\end{proposition}

\begin{proof}
La primera afirmación resulta de la conmutatividad de \eref{II.8.5.1.3-fr}.
Para demostrar la segunda, se puede reducir al caso en que $Y=\operatorname{Spec}(A)$ y $Y'=\operatorname{Spec}(A')$ son afines, $\mathcal{S}=\widetilde{S}$ y $\mathcal{S}'=\widetilde{S}'$.
Para todo $f$ homogéneo de $S_+$, poniendo $f'=\varphi(f)$, se tiene $\Phi^{-1}(C_f)=C'_{f'}$ \sref[I]{I.2.2.4.1-fr};
decir que $G(\varphi)=\operatorname{Proj}(S')$ significa que el radical (en $S'_+$) del ideal generado por los $f'=\varphi(f)$ es el propio $S'_+$ (\sref{II.2.8.1-fr} y \sref{II.2.3.14-fr}), lo que equivale a decir que los $C'_{f'}$ cubren $E'$ \sref{II.8.3.4.4-fr}.
\end{proof}
\begin{env}[8.5.3]
\label{II.8.5.3-fr}
El $q$-morfismo $\varphi$ se prolonga canónicamente a un $q$-morfismo de álgebras graduadas
\begin{equation}
\label{II.8.5.3.1-fr}
  \widehat{\varphi}:\widehat{\mathcal{S}}
  \to\widehat{\mathcal{S}}'
\tag{8.5.3.1}
\end{equation}
poniendo $\widehat{\varphi}(\mathbf{z})=\mathbf{z}$.
Se deduce de ello un morfismo
\[
  \widehat{\Phi}=\operatorname{Proj}(\widehat{\varphi}):
  G(\widehat{\varphi})\to\widehat{C}
  =\operatorname{Proj}(\widehat{\mathcal{S}})
\]
tal que el diagrama
\[
  \xymatrix{
    G(\widehat{\varphi})\ar[r]^{\widehat{\Phi}}\ar[d] &
    \widehat{C}\ar[d]\\
    Y'\ar[r]_{q} & Y
  }
\]
es conmutativo \sref{II.3.5.6-fr}.
Resulta inmediatamente de las definiciones que, si se denota por $i:C\to\widehat{C}$ e $i':C'\to\widehat{C}'$ las inmersiones abiertas canónicas \sref{II.8.3.2-fr}, entonces $i'(C')\subset G(\widehat{\varphi})$, y el diagrama
\begin{equation}
\label{II.8.5.3.2-fr}
  \xymatrix{
    C'\ar[r]^{\Phi}\ar[d]_{i'} & C\ar[d]^{i}\\
    G(\widehat{\varphi})\ar[r]_{\widehat{\Phi}} & \widehat{C}
  }
\tag{8.5.3.2}
\end{equation}
es conmutativo.
Finalmente, si se pone $X=\operatorname{Proj}(\mathcal{S})$ y $X'=\operatorname{Proj}(\mathcal{S}')$, y si $j:X\to\widehat{C}$ y $j':X'\to\widehat{C}'$ son las inmersiones cerradas canónicas \sref{II.8.3.2-fr}, resulta asimismo de la definición de estas inmersiones que $j'(G(\varphi))\subset G(\widehat{\varphi})$, y que el diagrama
\oldpage[II]{171}
\begin{equation}
\label{II.8.5.3.3-fr}
  \xymatrix{
    G(\varphi)\ar[r]^{\operatorname{Proj}(\varphi)}\ar[d]_{j'} &
    X\ar[d]^{j}\\
    G(\widehat{\varphi})\ar[r]_{\widehat{\Phi}} & \widehat{C}
  }
\tag{8.5.3.3}
\end{equation}
es conmutativo.
\end{env}
\begin{proposition}[8.5.4]
\label{II.8.5.4-fr}
Si $\widehat{E}$ \textup{(resp. $\widehat{E}'$)} es el cono proyectivo sin vértice definido por $\mathcal{S}$ (resp. por $\mathcal{S}'$), entonces $\widehat{\Phi}^{-1}(\widehat{E})\subset\widehat{E}'$;
además, si $p:\widehat{E}\to X$ y $p':\widehat{E}'\to X'$ son las retracciones canónicas, entonces $p'(\widehat{\Phi}^{-1}(\widehat{E}))\subset G(\varphi)$, y el diagrama
\begin{equation}
\label{II.8.5.4.1-fr}
  \xymatrix{
    \widehat{\Phi}^{-1}(\widehat{E})
      \ar[r]^{\widehat{\Phi}}\ar[d]_{p'} &
    \widehat{E}\ar[d]^{p}\\
    G(\varphi)\ar[r]_{\operatorname{Proj}(\varphi)} & X
  }
\tag{8.5.4.1}
\end{equation}
es conmutativo.
Si $\operatorname{Proj}(\varphi)$ está definido en todas partes, también lo está $\widehat{\Phi}$, y se tiene $\widehat{\Phi}^{-1}(\widehat{E})=\widehat{E}'$.
\end{proposition}

\begin{proof}
La primera afirmación resulta de la conmutatividad de los diagramas \eref{II.8.5.1.3-fr} y \eref{II.8.5.3.2-fr}, y las dos afirmaciones siguientes, de la definición de las retracciones canónicas \sref{II.8.3.5-fr} y de la definición de $\widehat{\varphi}$.
Para ver que $\widehat{\Phi}$ está definido en todas partes cuando lo está $\operatorname{Proj}(\varphi)$, se puede limitar al caso en que $Y=\operatorname{Spec}(A)$ y $Y'=\operatorname{Spec}(A')$ son afines, y en que $\mathcal{S}=\widetilde{S}$ y $\mathcal{S}'=\widetilde{S}'$;
la hipótesis es que, cuando $f$ recorre el conjunto de los elementos homogéneos de $S_+$, el ideal generado en $S'_+$ por los $\varphi(f)$ tiene por radical en $S'_+$ al propio ideal $S'_+$;
se concluye inmediatamente que el radical en $(S'[\mathbf{z}])_+$ del ideal generado por $\mathbf{z}$ y los $\varphi(f)$ es el propio $(S'[\mathbf{z}])_+$, de donde resulta nuestra afirmación;
esto muestra al mismo tiempo que $\widehat{E}'$ es la unión de los $\widehat{C}'_{\varphi(f)}$, y por tanto es igual a $\widehat{\Phi}^{-1}(\widehat{E})$.
\end{proof}
\begin{corollary}[8.5.5]
\label{II.8.5.5-fr}
Siempre que $\operatorname{Proj}(\varphi)$ está definido en todas partes, la imagen inversa por $\widehat{\Phi}$ del espacio subyacente al lugar en el infinito \textup{(resp.} al preesquema de los vértices\textup{)} de $\widehat{C}$ es el espacio subyacente al lugar en el infinito \textup{(resp.} al preesquema de los vértices\textup{)} de $\widehat{C'}$.
\end{corollary}

\begin{proof}
Esto resulta inmediatamente de \sref{II.8.5.4-fr} y \sref{II.8.5.2-fr}, teniendo en cuenta las relaciones \sref{II.8.3.4.1-fr} y \sref{II.8.3.4.2-fr}.
\end{proof}
\subsection{Un isomorfismo canónico para los conos sin vértice}
\label{subsection:II.8.6-fr}

\begin{env}[8.6.1]
\label{II.8.6.1-fr}
Sean $Y$ un preesquema, $\mathcal{S}$ una $\mathcal{O}_Y$-álgebra cuasicoherente graduada en grados positivos \emph{tal que $\mathcal{S}_0=\mathcal{O}_Y$}, y $X$ el $Y$-esquema $\operatorname{Proj}(\mathcal{S})$.
Vamos a aplicar los resultados de \sref{subsection:II.8.5-fr} al caso en que $Y'=X$ y $q:X\to Y$ es el morfismo estructural;
sea
\begin{equation}
\label{II.8.6.1.1-fr}
  \mathcal{S}_X=\bigoplus_{n\in\mathbf{Z}}\mathcal{O}_X(n)
\tag{8.6.1.1}
\end{equation}
\oldpage[II]{172}
que es una $\mathcal{O}_X$-álgebra graduada cuasicoherente, cuya multiplicación se define mediante los homomorfismos canónicos \sref{II.3.2.6.1-fr}
\[
  \mathcal{O}_X(m)\otimes_{\mathcal{O}_X}\mathcal{O}_X(n)
  \to\mathcal{O}_X(m+n)
\]
cuya asociatividad está asegurada por el diagrama conmutativo de \sref{II.2.5.11.4-fr}.
Tomemos por $\mathcal{S}'$ la $\mathcal{O}_X$-subálgebra cuasicoherente graduada en grados positivos $\mathcal{S}_X^{\geq}=\bigoplus_{n\geq0}\mathcal{O}_X(n)$ de $\mathcal{S}_X$.

Finalmente, consideremos el $q$-morfismo canónico
\begin{equation}
\label{II.8.6.1.2-fr}
  \alpha:\mathcal{S}\to\mathcal{S}_X^{\geq}
\tag{8.6.1.2}
\end{equation}
definido en \sref{II.3.3.2.3-fr} como un homomorfismo $\mathcal{S}\to q_*(\mathcal{S}_X)$, pero que evidentemente envía $\mathcal{S}$ a $q_*(\mathcal{S}_X^{\geq})$.
Pongamos
\begin{equation}
\label{II.8.6.1.3-fr}
  C_X=\operatorname{Spec}(\mathcal{S}_X^{\geq}),\qquad
  \widehat{C}_X=\operatorname{Proj}(\mathcal{S}_X^{\geq}[\mathbf{z}]),\qquad
  X'=\operatorname{Proj}(\mathcal{S}_X^{\geq})
\tag{8.6.1.3}
\end{equation}
y denotemos por $E_X$ y $\widehat{E}_X$ el cono afín sin vértice y el cono proyectivo sin vértice correspondientes (respectivamente);
se denotan por $\varepsilon_X:X\to C_X$, $i_X:C_X\to\widehat{C}_X$, $j_X:X'\to\widehat{C}_X$, $p_X:\widehat{E}_X\to X'$ y $\pi_X:E_X\to X'$ los morfismos canónicos definidos en \sref{subsection:II.8.3-fr}.
\end{env}
\begin{proposition}[8.6.2]
\label{II.8.6.2-fr}
El morfismo estructural $u:X'\to X$ es un \emph{isomorfismo}, y el morfismo $\operatorname{Proj}(\alpha)$ está definido en todas partes y es idéntico a $u$.
El morfismo $\operatorname{Proj}(\widehat{\alpha}):\widehat{C}_X\to\widehat{C}$ está definido en todas partes, y sus restricciones a $\widehat{E}_X$ y $E_X$ son \emph{isomorfismos} sobre $\widehat{E}$ y $E$, respectivamente.
Finalmente, si se identifican $X'$ y $X$ mediante $u$, los morfismos $p_X$ y $\pi_X$ se identifican con los morfismos estructurales de los $X$-preesquemas $\widehat{E}_X$ y $E_X$.
\end{proposition}

\begin{proof}
Se puede limitar evidentemente al caso en que $Y=\operatorname{Spec}(A)$ es afín y $\mathcal{S}=\widetilde{S}$;
entonces $X$ es la unión de los abiertos afines $X_f$, donde $f$ recorre el conjunto de los elementos homogéneos de $S_+$, y el anillo de cada $X_f$ es $S_{(f)}$.
Resulta de \sref{II.8.2.7.1-fr} que
\begin{equation}
\label{II.8.6.2.1-fr}
  \Gamma(X_f,\mathcal{S}_X^{\geq})=S_f^{\geq}.
\tag{8.6.2.1}
\end{equation}

Por tanto, $u^{-1}(X_f)=\operatorname{Proj}(S_f^{\geq})$.
Pero, si $f\in S_d$ ($d>0$), entonces $\operatorname{Proj}(S_f^{\geq})$ es canónicamente isomorfo a $\operatorname{Proj}((S_f^{\geq})^{(d)})$ \sref{II.2.4.7-fr}, y, por otra parte, $(S_f^{\geq})^{(d)}=(S^{(d)})_f^{\geq}$ puede identificarse con $S_{(f)}[T]$ \sref{II.2.2.1-fr} mediante la aplicación $T\to f/1$;
se concluye entonces \sref{II.3.1.7-fr} que el morfismo estructural $u^{-1}(X_f)\to X_f$ es un isomorfismo, de donde resulta la primera afirmación.
Para demostrar la segunda, observemos que la restricción $u^{-1}(X_f)\cap G(\alpha)\to X=\operatorname{Proj}(S)$ de $\operatorname{Proj}(\alpha)$ corresponde a la aplicación canónica $x\to x/1$ de $S$ en $S_f^{\geq}$ \sref{II.2.6.2-fr};
se deduce ante todo que $G(\alpha)=X'$, y luego, teniendo en cuenta que $u^{-1}(X_f)=(u^{-1}(X_f))_{f/1}$, resulta de \sref{II.2.8.1.1-fr} que la imagen de $u^{-1}(X_f)$ por $\operatorname{Proj}(\alpha)$ está contenida en $X_f$, y que la restricción de $\operatorname{Proj}(\alpha)$ a $u^{-1}(X_f)$, considerada como morfismo hacia $X_f=\operatorname{Spec}(S_{(f)})$, es efectivamente idéntica a la de $u$.
Finalmente, aplicando \sref{II.8.3.5.4-fr} a $p_X$ en lugar de a $p$, se ve que $p_X^{-1}(u^{-1}(X_f))=
\operatorname{Spec}((S_f^{\geq})_{f/1}^{\leq})$, y este abierto es, por \sref{II.8.5.4.1-fr}, la imagen inversa por $\operatorname{Proj}(\widehat{\alpha})$ de $p^{-1}(X_f)=\operatorname{Spec}(S_f^{\leq})$ \sref{II.8.3.5.3-fr}.
Teniendo en cuenta \sref{II.8.2.3.2-fr}, la restricción de $\operatorname{Proj}(\widehat{\alpha})$ a $p_X^{-1}(u^{-1}(X_f))$ corresponde al isomorfismo inverso de \sref{II.8.2.7.2-fr}, restringido a $S_f^{\leq}$, de donde resulta la tercera afirmación;
la última afirmación es evidente por definición.

\oldpage[II]{173}
Se observa asimismo que del diagrama conmutativo \eref{II.8.5.3.2-fr} resulta que \emph{la restricción a $C_X$ de $\operatorname{Proj}(\widehat{\alpha})$ no es otra que el morfismo $\operatorname{Spec}(\alpha)$}.
\end{proof}
\begin{corollary}[8.6.3]
\label{II.8.6.3-fr}
Considerados como $X$-esquemas, $\widehat{E}_X$ es canónicamente isomorfo a $\operatorname{Spec}(\mathcal{S}_X^{\leq})$, y $E_X$ a $\operatorname{Spec}(\mathcal{S}_X)$.
\end{corollary}

\begin{proof}
Como se sabe que los morfismos $p_X$ y $\pi_X$ son afines (\sref{II.8.3.5-fr} y \sref{II.8.3.6-fr}), basta (en virtud de \sref{II.1.3.1-fr}) demostrar el corolario en el caso en que $Y=\operatorname{Spec}(A)$ es afín y $\mathcal{S}=\widetilde{S}$.
La primera afirmación resulta de la existencia de los isomorfismos canónicos \sref{II.8.2.7.2-fr} $(S_f^{\geq})_{f/1}^{\leq}\xrightarrow{\sim}S_f^{\leq}$ y del hecho de que estos isomorfismos son compatibles con la aplicación que envía $f$ a $fg$ (donde $f$ y $g$ son homogéneos en $S_+$).
Análogamente, la fórmula \sref{II.8.3.6.2-fr}, aplicada a $\pi_X$ en lugar de a $\pi$, muestra que $\pi_X^{-1}(u^{-1}(X_f))=
\operatorname{Spec}((S_f^{\geq})_{f/1})$ para $f$ homogéneo en $S_+$, y la segunda afirmación resulta entonces de la existencia de los isomorfismos canónicos \sref{II.8.2.7.2-fr} $(S_f^{\geq})_{f/1}\xrightarrow{\sim}S_f$.

Se puede decir entonces que $\widehat{C}_X$, considerado como un $X$-esquema, se obtiene por \emph{pegado} de los dos $X$-esquemas afines sobre $X$, $C_X=\operatorname{Spec}(\mathcal{S}_X^{\geq})$ y $\widehat{E}_X=\operatorname{Spec}(\mathcal{S}_X^{\leq})$ cuya intersección es el abierto $E_X=\operatorname{Spec}(\mathcal{S}_X)$.
\end{proof}
\begin{corollary}[8.6.4]
\label{II.8.6.4-fr}
Supongamos que $\mathcal{O}_X(1)$ es un $\mathcal{O}_X$-módulo invertible y que $\mathcal{S}_X$ es isomorfa a $\bigoplus_{n\in\mathbf{Z}}(\mathcal{O}_X(1))^{\otimes n}$ (lo que ocurrirá, en particular, siempre que $\mathcal{S}$ esté generada por $\mathcal{S}_1$ \textup{(}\sref{II.3.2.5-fr} y \sref{II.3.2.7-fr}\textup{))}.
Entonces el cono proyectivo sin vértice $\widehat{E}$ se identifica con el fibrado vectorial de rango 1 $\mathbf{V}(\mathcal{O}_X(-1))$ sobre $X$, y el cono afín sin vértice $E$ con el subpreesquema de este fibrado vectorial inducido sobre el complementario de la sección nula.
Con esta identificación, la retracción canónica $\widehat{E}\to X$ se identifica con el morfismo estructural del $X$-esquema $\mathbf{V}(\mathcal{O}_X(-1))$.
Por último, existe un $Y$-morfismo canónico $\mathbf{V}(\mathcal{O}_X(1))\to C$, cuya restricción al complementario de la sección nula de $\mathbf{V}(\mathcal{O}_X(1))$ es un isomorfismo de este complementario sobre el cono afín sin vértice $E$.
\end{corollary}

\begin{proof}
Si se pone $\mathcal{L}=\mathcal{O}_X(1)$, entonces $\mathcal{S}_X^{\geq}$ es idéntica a $\mathbf{S}_{\mathcal{O}_X}(\mathcal{L})$, por lo que $\widehat{E}_X$ se identifica canónicamente con $\mathbf{V}(\mathcal{L}^{-1})$, en virtud de \sref{II.8.6.3-fr}, y $C_X$ con $\mathbf{V}(\mathcal{L})$.
El morfismo $\mathbf{V}(\mathcal{L})\to C$ es la restricción de $\operatorname{Proj}(\widehat{\alpha})$, y las afirmaciones del corolario son entonces casos particulares de \sref{II.8.6.2-fr}.
\end{proof}

Se observa que la imagen inversa por el morfismo $\mathbf{V}(\mathcal{O}_X(1))\to C$ del espacio subyacente al preesquema de los vértices de $C$ es el espacio subyacente a la sección nula de $\mathbf{V}(\mathcal{O}_X(1))$ \sref{II.8.5.5-fr};
pero, en general, los subpreesquemas correspondientes de $C$ y de $\mathbf{V}(\mathcal{O}_X(1))$ no son isomorfos.
Esta cuestión se estudiará más adelante.
\subsection{Explosión de conos proyectantes}
\label{subsection:II.8.7-fr}

\begin{env}[8.7.1]
\label{II.8.7.1-fr}
Bajo las condiciones de \sref{II.8.6.1-fr}, se tiene, poniendo $r=\operatorname{Proj}(\widehat{\alpha})$, un diagrama conmutativo
\begin{equation}
\label{II.8.7.1.1-fr}
  \xymatrix{
    X\ar[r]^{i_X\circ\varepsilon_X}\ar[d]_{q} &
    \widehat{C}_X\ar[d]^{r}\\
    Y\ar[r]_{i\circ\varepsilon} & \widehat{C}
  }
\tag{8.7.1.1}
\end{equation}
\oldpage[II]{174}
en virtud de \eref{II.8.5.1.3-fr} y \eref{II.8.5.3.2-fr};
además, la restricción de $r$ al complementario $\widehat{C}_X\setmin i_X(\varepsilon_X(X))$ de la sección nula es un \emph{isomorfismo} sobre el complementario $\widehat{C}\setmin i(\varepsilon(Y))$ de la sección nula, en virtud de \sref{II.8.6.2-fr}.
Si se supone, para simplificar, que $Y$ es afín y que $\mathcal{S}$ es de tipo finito y está generada por $\mathcal{S}_1$, entonces $X$ es proyectivo sobre $Y$ y $\widehat{C}_X$ es proyectivo sobre $X$ \sref{II.5.5.1-fr}; por tanto, $\widehat{C}_X$ es proyectivo sobre $Y$ \sref{II.5.5.5-fr}[ii], y \emph{a fortiori} sobre $\widehat{C}$ \sref{II.5.5.5-fr}[v].
Se tiene así un $Y$-morfismo proyectivo $r:\widehat{C}_X\to\widehat{C}$ (cuya restricción a $C_X$ es un $Y$-morfismo proyectivo $C_X\to C$) que \emph{contrae $X$ sobre $Y$} e induce un \emph{isomorfismo} al restringirlo a los \emph{complementarios de $X$ y de $Y$}.
Se tiene, pues, una relación entre $C_X$ y $C$, análoga a la que existe entre un preesquema obtenido por explosión y su preesquema de partida \sref{II.8.1.3-fr}.
Mostraremos, en efecto, que $C_X$ puede identificarse con el espectro homogéneo de una $\mathcal{O}_C$-álgebra graduada.
\end{env}
\begin{env}[8.7.2]
\label{II.8.7.2-fr}
Conservando las notaciones de \sref{II.8.6.1-fr}, consideremos, para todo $n\geqslant 0$, el ideal cuasicoherente
\begin{equation}
\label{II.8.7.2.1-fr}
  \mathcal{S}_{[n]}=\bigoplus_{m\geqslant n}\mathcal{S}_m
\tag{8.7.2.1}
\end{equation}
de la $\mathcal{O}_Y$-álgebra graduada $\mathcal{S}$.
Está claro que
\begin{equation}
\label{II.8.7.2.2-fr}
  \mathcal{S}_{[0]}=\mathcal{S},\qquad
  \mathcal{S}_{[n]}\subset\mathcal{S}_{[m]}
  \qquad\text{para }m\leqslant n
\tag{8.7.2.2}
\end{equation}
\begin{equation}
\label{II.8.7.2.3-fr}
  \mathcal{S}_{[n]}\mathcal{S}_{[m]}\subset\mathcal{S}_{[m+n]}.
\tag{8.7.2.3}
\end{equation}

Consideremos el $\mathcal{O}_C$-módulo asociado a $\mathcal{S}_{[n]}$, que es un ideal cuasicoherente de $\mathcal{O}_C=\widetilde{\mathcal{S}}$ \sref{II.1.4.4-fr}
\begin{equation}
\label{II.8.7.2.4-fr}
  \mathcal{I}_n=(\mathcal{S}_{[n]})^{\sim}
\tag{8.7.2.4}
\end{equation}

Se deducen entonces de \sref{II.8.7.2.2-fr} y \sref{II.8.7.2.3-fr}, teniendo en cuenta \sref{II.1.4.4-fr} y \sref{II.1.4.8.1-fr}, las fórmulas análogas
\begin{equation}
\label{II.8.7.2.5-fr}
  \mathcal{I}_0=\mathcal{O}_C,\qquad
  \mathcal{I}_n\subset\mathcal{I}_m
  \qquad\text{para }m\leqslant n
\tag{8.7.2.5}
\end{equation}
\begin{equation}
\label{II.8.7.2.6-fr}
  \mathcal{I}_n\mathcal{I}_m\subset\mathcal{I}_{n+m}.
\tag{8.7.2.6}
\end{equation}

Se está así en la situación de \sref{II.8.1.1-fr}, lo que lleva a introducir la $\mathcal{O}_C$-álgebra graduada cuasicoherente
\begin{equation}
\label{II.8.7.2.7-fr}
  \mathcal{S}^{\natural}
  =\bigoplus_{n\geqslant 0}\mathcal{I}_n
  =\left(\bigoplus_{n\geqslant 0}\mathcal{S}_{[n]}\right)^{\sim}.
\tag{8.7.2.7}
\end{equation}
\end{env}
\begin{proposition}[8.7.3]
\label{II.8.7.3-fr}
Existe un $C$-isomorfismo canónico
\begin{equation}
\label{II.8.7.3.1-fr}
  h:C_X\xrightarrow{\sim}\operatorname{Proj}(\mathcal{S}^{\natural}).
\tag{8.7.3.1}
\end{equation}
\end{proposition}

\begin{proof}
Supongamos primero que $Y=\operatorname{Spec}(A)$ es afín, de modo que $\mathcal{S}=\widetilde{S}$, siendo $S$ una $A$-álgebra graduada en grados positivos, y $C=\operatorname{Spec}(S)$.
La definición~\sref{II.8.2.6.2-fr} muestra entonces, con las notaciones de \sref{II.8.2.6-fr}, que $\mathcal{S}^{\natural}=(S^{\natural})^{\sim}$.
Para definir \sref{II.8.7.3.1-fr}, consideremos un elemento homogéneo $f\in S_d$ ($d>0$) y el elemento correspondiente $f^{\natural}\in S^{\natural}$ \sref{II.8.2.6-fr};
el $S$-isomorfismo de \sref{II.8.2.7.3-fr} define entonces un $C$-isomorfismo
\begin{equation}
\label{II.8.7.3.2-fr}
  \operatorname{Spec}(S_f^{\geqslant})
  \xrightarrow{\sim}
  \operatorname{Spec}(S_{(f^{\natural})}^{\natural}).
\tag{8.7.3.2}
\end{equation}

\oldpage[II]{175}
Pero, con las notaciones de \sref{II.8.6.2-fr}, si $v:C_X\to X$ es el morfismo estructural, de \sref{II.8.6.2.1-fr} resulta que $v^{-1}(X_f)=\operatorname{Spec}(S_f^{\geqslant})$.
Se tiene asimismo $\operatorname{Spec}(S_{(f^{\natural})}^{\natural})=D_+(f^{\natural})$, lo que significa que \sref{II.8.7.3.2-fr} define un isomorfismo $v^{-1}(X_f)\to D_+(f^{\natural})$.
Además, si $g\in S_e$ ($e>0$), entonces el diagrama
\[
  \xymatrix{
    v^{-1}(X_{fg})\ar[r]^{\sim}\ar[d] &
    D_+(f^{\natural}g^{\natural})\ar[d]\\
    v^{-1}(X_f)\ar[r]^{\sim} & D_+(f^{\natural})
  }
\]
es conmutativo, por la definición del isomorfismo de \sref{II.8.2.7.3-fr}.
Finalmente, por definición, $S_+$ está engendrado por los elementos homogéneos $f$, y se deduce entonces de \sref{II.8.2.10-fr}[iv] y de \sref{II.2.3.14-fr} que los $D_+(f^{\natural})$ forman un recubrimiento de $\operatorname{Proj}(S^{\natural})$, y que los $v^{-1}(X_f)$ forman un recubrimiento de $C_X$, puesto que los $X_f$ forman un recubrimiento de $X$;
esto termina, pues, de definir en este caso el isomorfismo \sref{II.8.7.3.1-fr}.

Para demostrar \sref{II.8.7.3-fr} en el caso general, basta mostrar que, si $U$ y $U'$ son abiertos afines de $Y$ tales que $U'\subset U$, de anillos $A$ y $A'$, entonces, poniendo $\mathcal{S}|U=\widetilde{S}$ y $\mathcal{S}|U'=\widetilde{S}'$, el diagrama
\begin{equation}
\label{II.8.7.3.3-fr}
  \xymatrix{
    C_{U'}\ar[r]\ar[d] & \operatorname{Proj}(S'^{\natural})\ar[d]\\
    C_U\ar[r] & \operatorname{Proj}(S^{\natural})
  }
\tag{8.7.3.3}
\end{equation}
es conmutativo.
Pero $S'$ se identifica canónicamente con $S\otimes_A A'$, y, por tanto, $S'^{\natural}$ se identifica canónicamente con
\[
  S^{\natural}\otimes_S S'=S^{\natural}\otimes_A A'~;
\]
así, $\operatorname{Proj}(S'^{\natural})=\operatorname{Proj}(S^{\natural})
\times_U U'$ \sref{II.2.8.10-fr};
análogamente, si $X=\operatorname{Proj}(S)$ y $X'=\operatorname{Proj}(S')$, entonces $X'=X\times_U U'$ y $\mathcal{S}_{X'}=\mathcal{S}_X\otimes_{\mathcal{O}_U}\mathcal{O}_{U'}$ \sref{II.3.5.4-fr}, o, equivalentemente, $\mathcal{S}_{X'}=j^*(\mathcal{S}_X)$, donde $j$ es la proyección $X'\to X$.
Se tiene entonces, por \sref{II.1.5.2-fr}, que $C_{U'}=C_U\times_X X'=C_U\times_U U'$, y la conmutatividad de \sref{II.8.7.3.3-fr} es inmediata.
\end{proof}
\begin{remark}[8.7.4]
\label{II.8.7.4-fr}
\begin{enumerate}
  \item[\rm{(i)}] El final de la demostración de \sref{II.8.7.3-fr} puede generalizarse inmediatamente de la manera siguiente.
    Sea $g:Y'\to Y$ un morfismo, $\mathcal{S}'=g^*(\mathcal{S})$, y $X'=\operatorname{Proj}(\mathcal{S}')$;
    se tiene entonces un diagrama conmutativo
    \begin{equation}
  \label{II.8.7.4.1-fr}
    \xymatrix{
      C_{X'}\ar[r]\ar[d] &
      \operatorname{Proj}(\mathcal{S}'^{\natural})\ar[d]\\
      C_X\ar[r] & \operatorname{Proj}(\mathcal{S}^{\natural})
    }
  \tag{8.7.4.1}
  \end{equation}

    Sea ahora $\varphi:\mathcal{S}''\to\mathcal{S}$ un homomorfismo de $\mathcal{O}_Y$-álgebras graduadas tal que, si se pone $X''=\operatorname{Proj}(\mathcal{S}'')$, entonces $u=\operatorname{Proj}(\varphi):X\to X''$ esté definido en todas partes;
    se tiene también
\oldpage[II]{176}
    un $Y$-morfismo $v:C\to C''$ (con $C''=\operatorname{Spec}(\mathcal{S}'')$) tal que $\mathcal{A}(v)=\varphi$, y, puesto que $\varphi$ es un homomorfismo de álgebras graduadas, $\varphi$ induce un $v$-morfismo de álgebras graduadas $\psi:\mathcal{S}''^{\natural}\to\mathcal{S}^{\natural}$ \sref{II.1.4.1-fr}.
    Además, se deduce de \sref{II.8.2.10-fr}[iv] y de la hipótesis sobre $\varphi$ que $\operatorname{Proj}(\psi)$ está definido en todas partes.
    Finalmente, teniendo en cuenta \sref{II.3.5.6.1-fr}, existe un $u$-morfismo canónico $\mathcal{S}_{X''}\to\mathcal{S}_X$, de donde se obtiene, por \sref{II.1.5.6-fr}, un morfismo $w:C_X\to C_{X''}$.
    En estas condiciones, el diagrama
    \begin{equation}
  \label{II.8.7.4.2-fr}
    \xymatrix{
      C_X\ar[r]^{\sim}\ar[d]_{w} &
      \operatorname{Proj}(\mathcal{S}^{\natural})
        \ar[d]^{\operatorname{Proj}(\psi)}\\
      C_{X''}\ar[r]^{\sim} & \operatorname{Proj}(\mathcal{S}''^{\natural})
    }
  \tag{8.7.4.2}
  \end{equation}
    es conmutativo, como se comprueba inmediatamente reduciéndose al caso en que $Y$ es afín.
  \item[\rm{(ii)}] Observemos que, por \sref{II.8.7.2.5-fr} y \sref{II.8.7.2.6-fr}, se tiene $\mathcal{I}_1^m\subset\mathcal{I}_m\subset\mathcal{I}_1$ para todo $m>0$.
    Ahora bien, por definición, $\mathcal{I}_1=(\mathcal{S}_+)^{\sim}$, y, por tanto, $\mathcal{I}_1$ define en $C$ el subpreesquema cerrado $\varepsilon(Y)$ (\sref{II.1.4.10-fr} y \sref{II.8.3.2-fr});
    se concluye así que, para todo $m>0$, \emph{el soporte de $\mathcal{O}_C/\mathcal{I}_m$ está contenido en el espacio subyacente al preesquema de los vértices $\varepsilon(Y)$};
    sobre la imagen inversa del cono afín sin vértice $E$, el morfismo estructural $\operatorname{Proj}(\mathcal{S}^{\natural})\to C$ se restringe, por tanto, a un \emph{isomorfismo} (por \sref{II.8.7.3-fr} y \sref{II.8.7.1-fr}).
    Además, identificando canónicamente $C$ con un abierto de $\widehat{C}$ \sref{II.8.3.3-fr}, se puede evidentemente prolongar los ideales $\mathcal{I}_m$ de $\mathcal{O}_C$ a ideales $\mathcal{J}_m$ de $\mathcal{O}_{\widehat{C}}$, imponiendo que coincidan con $\mathcal{O}_{\widehat{C}}$ sobre el abierto $\widehat{E}$ de $\widehat{C}$.
    Si se define $\mathcal{T}=\bigoplus_{n\geqslant 0}\mathcal{J}_n$, que es una $\mathcal{O}_{\widehat{C}}$-álgebra graduada cuasicoherente, se puede prolongar el isomorfismo \sref{II.8.7.3.1-fr} a un $\widehat{C}$-isomorfismo
    \begin{equation}
  \label{II.8.7.4.3-fr}
    \widehat{C}_X\xrightarrow{\sim}\operatorname{Proj}(\mathcal{T}).
  \tag{8.7.4.3}
  \end{equation}

    En efecto, sobre $\widehat{E}$, se deduce de lo anterior que $\operatorname{Proj}(\mathcal{T})$ se identifica canónicamente con $\widehat{E}$, y se define, por tanto, el isomorfismo \sref{II.8.7.4.3-fr} sobre $\widehat{E}$ exigiendo que coincida con el isomorfismo canónico $\widehat{E}_X\to\widehat{E}$ \sref{II.8.6.2-fr};
    está claro que este isomorfismo y \sref{II.8.7.3.1-fr} coinciden entonces sobre $E$.
\end{enumerate}
\end{remark}
\begin{corollary}[8.7.5]
\label{II.8.7.5-fr}
Supongamos que existe un $n_0>0$ tal que
\begin{equation}
\label{II.8.7.5.1-fr}
  \mathcal{S}_{n+1}=\mathcal{S}_1\mathcal{S}_n
  \qquad\text{para }n\geqslant n_0.
\tag{8.7.5.1}
\end{equation}

Entonces el subpreesquema de los vértices de $C_X$ (isomorfo a $X$) es la imagen inversa, por el morfismo canónico $r:C_X\to C$, del subpreesquema de los vértices de $C$ (isomorfo a $Y$).
Recíprocamente, si se verifica esta propiedad, y si se supone además que $Y$ es noetheriano y que $\mathcal{S}$ es de tipo finito, existe un $n_0>0$ tal que se cumple \sref{II.8.7.5.1-fr}.
\end{corollary}

\begin{proof}
Como la primera afirmación es local sobre $Y$, se puede suponer que $Y=\operatorname{Spec}(A)$ es afín, de modo que $\mathcal{S}=\widetilde{S}$, con $S$ una $A$-álgebra graduada en grados positivos.
La afirmación se deduce entonces de \sref{II.8.2.12-fr}, pues $\operatorname{Proj}(S^{\natural}\otimes_S S_0)=
C_X\times_C\varepsilon(Y)$ (mediante la identificación de \sref{II.8.7.3.1-fr}), o, dicho de otro modo, porque este preesquema es la imagen inversa de $\varepsilon(Y)$ en $C_X$ \sref[I]{I.4.4.1-fr}.
La recíproca se deduce asimismo de \sref{II.8.2.12-fr} cuando $Y$ es noetheriano afín y $S$ es de tipo finito.
\oldpage[II]{177}
Si $Y$ es noetheriano (aunque no necesariamente afín) y $\mathcal{S}$ es de tipo finito, existe un recubrimiento finito de $Y$ por abiertos afines noetherianos $U_i$, y se deduce entonces de lo anterior que, para todo $i$, existe un entero $n_i$ tal que $\mathcal{S}_{n+1}|U_i=(\mathcal{S}_1|U_i)(\mathcal{S}_n|U_i)$ para $n\geqslant n_i$;
el mayor de los $n_i$ garantiza entonces que se cumple \sref{II.8.7.5.1-fr}.
\end{proof}
\begin{env}[8.7.6]
\label{II.8.7.6-fr}
Consideremos ahora el $C$-preesquema $Z$ obtenido por \emph{explosión} en el cono afín $C$ del \emph{subpreesquema de los vértices} $\varepsilon(Y)$;
por la Definición~\sref{II.8.1.3-fr}, es precisamente el preesquema $\operatorname{Proj}(\bigoplus_{n\geqslant 0}\mathcal{S}_+^n)$;
la inyección canónica
\begin{equation}
\label{II.8.7.6.1-fr}
  \iota:\bigoplus_{n\geqslant 0}\mathcal{S}_+^n
  \longrightarrow\mathcal{S}^{\natural}
\tag{8.7.6.1}
\end{equation}
define (gracias a la identificación de \sref{II.8.7.3-fr}) un $C$-morfismo dominante canónico
\begin{equation}
\label{II.8.7.6.2-fr}
  G(\iota)\longrightarrow Z
\tag{8.7.6.2}
\end{equation}
donde $G(\iota)$ es un abierto de $C_X$ \sref{II.3.5.1-fr};
se observa que puede ocurrir que $G(\iota)\ne C_X$, como muestra el ejemplo en que $Y=\operatorname{Spec}(K)$, siendo $K$ un cuerpo, y $\mathcal{S}=\widetilde{S}$, con $S=K[\mathbf{y}]$, donde $\mathbf{y}$ es una indeterminada \emph{de grado} 2;
si $R_n$ denota el conjunto $(S_+)^n$, considerado como subconjunto de $S_{[n]}=S_n^{\natural}$, entonces $S_+^{\natural}$ no es el radical en $S_+^{\natural}$ del ideal generado por la unión de los $R_n$ (cf. \sref{II.2.3.14-fr}).
\end{env}
\begin{corollary}[8.7.7]
\label{II.8.7.7-fr}
Supongamos que existe un $n_0>0$ tal que
\begin{equation}
\label{II.8.7.7.1-fr}
  \mathcal{S}_n=\mathcal{S}_1^n
  \qquad\text{para }n\geqslant n_0.
\tag{8.7.7.1}
\end{equation}

Entonces el morfismo canónico \sref{II.8.7.6.2-fr} está definido en todas partes y es un isomorfismo $C_X\xrightarrow{\sim}Z$.
Recíprocamente, si se verifica esta propiedad, y si se supone además que $Y$ es noetheriano y que $\mathcal{S}$ es de tipo finito, existe un $n_0$ tal que se cumple \sref{II.8.7.7.1-fr}.
\end{corollary}

\begin{proof}
La primera afirmación es local sobre $Y$ y, por tanto, se deduce de \sref{II.8.2.14-fr};
lo mismo ocurre con la recíproca, razonando como en \sref{II.8.7.5-fr}.
\end{proof}
\begin{remark}[8.7.8]
\label{II.8.7.8-fr}
Como la condición \sref{II.8.7.7.1-fr} implica \sref{II.8.7.5.1-fr}, se ve que, cuando se cumple, no solo $C_X$ se identifica con el preesquema obtenido por explosión del vértice (identificado con $Y$) del cono afín $C$, sino que además el vértice (identificado con $X$) de $C_X$ se identifica con el subpreesquema cerrado que es la imagen inversa del vértice $Y$ de $C$.
Además, la hipótesis~\sref{II.8.7.7.1-fr} implica que, sobre $X=\operatorname{Proj}(\mathcal{S})$, los $\mathcal{O}_X$-módulos $\mathcal{O}_X(n)$ son invertibles (\sref{II.3.2.5-fr} y \sref{II.3.2.9-fr}), y que $\mathcal{O}_X(n)=\mathcal{L}^{\otimes n}$, con $\mathcal{L}=\mathcal{O}_X(1)$ (\sref{II.3.2.7-fr} y \sref{II.3.2.9-fr});
por la Definición~\sref{II.8.6.1.1-fr}, $C_X$ es, por tanto, el \emph{fibrado vectorial} $\mathbf{V}(\mathcal{L})$ sobre $X$, y su vértice es la \emph{sección nula} de este fibrado vectorial.
\end{remark}
\subsection{Haces amplios y contracciones}
\label{subsection:II.8.8-fr}

\begin{env}[8.8.1]
\label{II.8.8.1-fr}
Sean $Y$ un preesquema, $f:X\to Y$ un morfismo \emph{separado} y \emph{cuasicompacto}, y $\mathcal{L}$ un $\mathcal{O}_X$-módulo invertible \emph{amplio relativamente a $f$}.
Consideremos la $\mathcal{O}_Y$-álgebra graduada en grados positivos
\begin{equation}
\label{II.8.8.1.1-fr}
  \mathcal{S}=\mathcal{O}_Y\oplus
  \bigoplus_{n\geqslant 1}f_*(\mathcal{L}^{\otimes n})
\tag{8.8.1.1}
\end{equation}
\oldpage[II]{178}
que es cuasicoherente \sref[I]{I.9.2.2-fr}[\emph{a}].
Hay un homomorfismo canónico de $\mathcal{O}_X$-álgebras graduadas
\begin{equation}
\label{II.8.8.1.2-fr}
  \tau:f^*(\mathcal{S})\longrightarrow
  \bigoplus_{n\geqslant 0}\mathcal{L}^{\otimes n}
\tag{8.8.1.2}
\end{equation}
que, en grados $\geqslant 1$, coincide con el homomorfismo canónico $\sigma:f^*(f_*(\mathcal{L}^{\otimes n}))\to\mathcal{L}^{\otimes n}$ \sref[0]{0.4.4.3-fr}, y es la identidad en grado $0$.
La hipótesis de que $\mathcal{L}$ es $f$-amplio implica entonces (\sref{II.4.6.3-fr} y \sref{II.3.6.1-fr}) que el $Y$-morfismo correspondiente
\begin{equation}
\label{II.8.8.1.3-fr}
  r=r_{\mathcal{L},\tau}:X\longrightarrow
  P=\operatorname{Proj}(\mathcal{S})
\tag{8.8.1.3}
\end{equation}
está definido en todas partes y es una \emph{inmersión abierta dominante}, y que
\begin{equation}
\label{II.8.8.1.4-fr}
  r^*(\mathcal{O}_P(n))=\mathcal{L}^{\otimes n}
  \qquad\text{para todo }n\in\mathbf{Z}.
\tag{8.8.1.4}
\end{equation}
\end{env}
\begin{proposition}[8.8.2]
\label{II.8.8.2-fr}
Sea $C=\operatorname{Spec}(\mathcal{S})$ el cono afín definido por $\mathcal{S}$;
si $\mathcal{L}$ es $f$-amplio, existe un $Y$-morfismo canónico
\begin{equation}
\label{II.8.8.2.1-fr}
  g:V=\mathbf{V}(\mathcal{L})\longrightarrow C
\tag{8.8.2.1}
\end{equation}
tal que el diagrama
\begin{equation}
\label{II.8.8.2.2-fr}
  \xymatrix{
    X\ar[r]^{j}\ar[d]_{f} &
    \mathbf{V}(\mathcal{L})\ar[r]^{\pi}\ar[d]^{g} & X\ar[d]^{f}\\
    Y\ar[r]_{\varepsilon} & C\ar[r]_{\psi} & Y
  }
\tag{8.8.2.2}
\end{equation}
es conmutativo, donde $\psi$ y $\pi$ son los morfismos estructurales, y $j$ y $\varepsilon$ las inmersiones canónicas que llevan, respectivamente, $X$ e $Y$ a la sección nula de $\mathbf{V}(\mathcal{L})$ y al preesquema de los vértices de $C$.
Además, la restricción de $g$ a $\mathbf{V}(\mathcal{L})\setmin j(X)$ es una inmersión abierta
\begin{equation}
\label{II.8.8.2.3-fr}
  \mathbf{V}(\mathcal{L})\setmin j(X)\longrightarrow
  E=C\setmin\varepsilon(Y)
\tag{8.8.2.3}
\end{equation}
en el cono afín sin vértice $E$ correspondiente a $\mathcal{S}$.
\end{proposition}

\begin{proof}
Con las notaciones de \sref{II.8.8.1-fr}, sean $\mathcal{S}_P^{\geqslant}=\bigoplus_{n\geqslant 0}\mathcal{O}_P(n)$ y $C_P=\operatorname{Spec}(\mathcal{S}_P^{\geqslant})$.
Se sabe por \sref{II.8.6.2-fr} que existe un morfismo canónico $h=\operatorname{Spec}(\alpha):C_P\to C$ tal que el diagrama
\begin{equation}
\label{II.8.8.2.4-fr}
  \xymatrix{
    C_P\ar[r]\ar[d]_{h} & P\ar[d]^{p}\\
    C\ar[r]_{\psi} & Y
  }
\tag{8.8.2.4}
\end{equation}
es conmutativo; además, si $\varepsilon_P:P\to C_P$ es la inmersión canónica, entonces el diagrama
\begin{equation}
\label{II.8.8.2.5-fr}
  \xymatrix{
    P\ar[r]^{\varepsilon_P}\ar[d]_{p} & C_P\ar[d]^{h}\\
    Y\ar[r]_{\varepsilon} & C
  }
\tag{8.8.2.5}
\end{equation}
es conmutativo \sref{II.8.7.1.1-fr}, y, finalmente, la restricción de $h$ al cono afín sin vértice $E_P$ es un \emph{isomorfismo} $E_P\xrightarrow{\sim}E$ \sref{II.8.6.2-fr}.
Se deduce de \sref{II.8.8.1.4-fr} que
\[
  r^*(\mathcal{S}_P^{\geqslant})=
  \mathbf{S}_{\mathcal{O}_X}(\mathcal{L})
\]
\oldpage[II]{179}
y, por tanto, existe un $P$-morfismo canónico $q:\mathbf{V}(\mathcal{L})\to C_P$, con el diagrama conmutativo
\begin{equation}
\label{II.8.8.2.6-fr}
  \xymatrix{
    \mathbf{V}(\mathcal{L})\ar[r]^{\pi}\ar[d]_{q} & X\ar[d]^{r}\\
    C_P\ar[r] & P
  }
\tag{8.8.2.6}
\end{equation}
que identifica $\mathbf{V}(\mathcal{L})$ con el producto $C_P\times_P X$ \sref{II.1.5.2-fr};
como $r$ es una inmersión abierta, también lo es $q$ \sref[I]{I.4.3.2-fr}.
Además, la restricción de $q$ a $\mathbf{V}(\mathcal{L})\setmin j(X)$ envía este preesquema a $E_P$, por \sref{II.8.5.2-fr}, y el diagrama
\begin{equation}
\label{II.8.8.2.7-fr}
  \xymatrix{
    X\ar[r]^{j}\ar[d]_{r} & \mathbf{V}(\mathcal{L})\ar[d]^{q}\\
    P\ar[r]_{\varepsilon_P} & C_P
  }
\tag{8.8.2.7}
\end{equation}
es conmutativo (pues es un caso particular de \eref{II.8.5.1.3-fr}).
Las afirmaciones de \sref{II.8.8.2-fr} se deducen inmediatamente de estos hechos, tomando por $g$ el morfismo compuesto $h\circ q$.
\end{proof}
\begin{remark}[8.8.3]
\label{II.8.8.3-fr}
Supongamos además que $Y$ es un preesquema \emph{noetheriano} y que $f$ es un morfismo \emph{propio}.
Como $r$ es entonces \emph{propio} \sref{II.5.4.4-fr}, y por tanto cerrado, y es además una inmersión abierta dominante, $r$ es necesariamente un \emph{isomorfismo} $X\xrightarrow{\sim}P$.
Se verá además en el Capítulo~III \sref[III]{III.2.3.5.1-fr} que $\mathcal{S}$ es entonces necesariamente una $\mathcal{O}_Y$-álgebra \emph{de tipo finito}.
Se sigue que $\mathcal{S}^{\natural}$ es una $\mathcal{S}_0^{\natural}$-álgebra \emph{de tipo finito} (\sref{II.8.2.10-fr}[i] y \sref{II.8.7.2.7-fr});
como $C_P$ es $C$-isomorfo a $\operatorname{Proj}(\mathcal{S}^{\natural})$ \sref{II.8.7.3-fr}, se ve que el morfismo $h:C_P\to C$ es \emph{proyectivo};
como el morfismo $r$ es un isomorfismo, también lo es $q:\mathbf{V}(\mathcal{L})\to C_P$, y se concluye que el morfismo $g:\mathbf{V}(\mathcal{L})\to C$ es \emph{proyectivo}.
Además, como la restricción de $h$ a $E_P$ es un isomorfismo sobre $E$ y $q$ es un isomorfismo, la restricción \sref{II.8.8.2.3-fr} de $g$ es un \emph{isomorfismo} $\mathbf{V}(\mathcal{L})\setmin j(X)\xrightarrow{\sim}E$.

Si se supone además que $\mathcal{L}$ es \emph{muy amplio} para $f$, se verá también en el Capítulo~III \sref[III]{III.2.3.5.1-fr} que existe un entero $n_0>0$ tal que $\mathcal{S}_n=\mathcal{S}_1^n$ para $n\geqslant n_0$.
Se concluirá entonces por \sref{II.8.7.7-fr} que $\mathbf{V}(\mathcal{L})$ se identifica con el preesquema $Z$ obtenido \emph{por explosión en el cono afín $C$ del subpreesquema de los vértices} (identificado con $Y$),
y que la \emph{sección nula} de $\mathbf{V}(\mathcal{L})$ (identificada con $X$) es la \emph{imagen inversa} del subpreesquema de los vértices $Y$ de $C$.

Una parte de los resultados precedentes puede establecerse de hecho sin hipótesis noetheriana:
\end{remark}
\begin{corollary}[8.8.4]
\label{II.8.8.4-fr}
Sean $Y$ un preesquema \textup{(resp.} un esquema cuasicompacto\textup{)}, $f:X\to Y$ un morfismo propio y $\mathcal{L}$ un $\mathcal{O}_X$-módulo invertible amplio relativamente a $f$.
Entonces el morfismo de \sref{II.8.8.2.1-fr} es propio \textup{(resp.} proyectivo\textup{)}, y su restricción \sref{II.8.8.2.3-fr} es un isomorfismo.
\end{corollary}

\begin{proof}
Para probar que $g$ es propio, se puede reducir al caso en que $Y$ es afín, y basta entonces considerar el caso en que $Y$ es un esquema cuasicompacto.
El mismo razonamiento que en \sref{II.8.8.3-fr} muestra primero que $r$ es un \emph{isomorfismo} $X\xrightarrow{\sim}P$;
por tanto, $q$ es también un isomorfismo y, como la restricción de $h$ a $E_P$ es un isomorfismo $E_P\xrightarrow{\sim}E$, ya se ve que \sref{II.8.8.2.3-fr} es un isomorfismo.
Solo queda probar que $g$ es \emph{proyectivo}.

Como $f$ es de tipo finito por hipótesis, se puede aplicar \sref{II.3.8.5-fr} al homomorfismo
\oldpage[II]{180}
$\tau$ de \sref{II.8.8.1.2-fr}:
existen un entero $d>0$ y un $\mathcal{O}_Y$-submódulo cuasicoherente de tipo finito $\mathcal{E}$ de $\mathcal{S}_d$ tales que, si $\mathcal{S}'$ es la $\mathcal{O}_Y$-subálgebra de $\mathcal{S}$ generada por $\mathcal{E}$ y $\tau'=\tau\circ f^*(\varphi)$ (donde $\varphi$ es la inyección canónica $\mathcal{S}'\to\mathcal{S}$), entonces $r'=r_{\mathcal{L},\tau'}$ es una inmersión
\[
  X\longrightarrow P'=\operatorname{Proj}(\mathcal{S}').
\]
Además, como $\varphi$ es inyectiva, $r'$ es también una \emph{inmersión dominante} \sref{II.3.7.6-fr};
el mismo razonamiento que para $r$ muestra, pues, que $r'$ es una \emph{inmersión cerrada sobreyectiva};
como $r'$ se factoriza en $X\xrightarrow{r}\operatorname{Proj}(\mathcal{S})
\xrightarrow{\Phi}\operatorname{Proj}(\mathcal{S}')$, donde $\Phi=\operatorname{Proj}(\varphi)$, se concluye que $\Phi$ es también una \emph{inmersión cerrada sobreyectiva}.
Pero esto implica que $\Phi$ es un \emph{isomorfismo};
en efecto, se puede reducir al caso en que $Y=\operatorname{Spec}(A)$ es afín y $\mathcal{S}=\widetilde{S}$, $\mathcal{S}'=\widetilde{S}'$, donde $S$ es una $A$-álgebra graduada y $S'$ una subálgebra graduada de $S$.
Para todo elemento homogéneo $t\in S'$, $S'_{(t)}$ es entonces un subanillo de $S_{(t)}$;
volviendo a la definición de $\operatorname{Proj}(\varphi)$ \sref{II.2.8.1-fr}, se ve que basta probar que, si $B'$ es un subanillo de un anillo $B$ y el morfismo $\operatorname{Spec}(B)\to\operatorname{Spec}(B')$ correspondiente a la inyección canónica $B'\to B$ es una inmersión cerrada, este morfismo es necesariamente un \emph{isomorfismo};
pero esto se sigue de \sref[I]{I.4.2.3-fr}.
Además, $\Phi^*(\mathcal{O}_{P'}(n))=\mathcal{O}_P(n)$ (\sref{II.3.5.2-fr}[ii] y \sref{II.3.5.4-fr}), y por tanto $r'^*(\mathcal{O}_{P'}(n))$ es isomorfo a $\mathcal{L}^{\otimes n}$ \sref{II.4.6.3-fr}.
Pongamos $\mathcal{S}''=\mathcal{S}'^{(d)}$, de modo que \sref{II.3.1.8-fr}[i] $X$ se identifica canónicamente con $P''=\operatorname{Proj}(\mathcal{S}'')$, y $\mathcal{L}''=\mathcal{L}^{\otimes d}$ con $\mathcal{O}_{P''}(1)$ \sref{II.3.2.9-fr}[ii].

Ahora bien, si $C''=\operatorname{Spec}(\mathcal{S}'')$, entonces $\mathcal{S}_{P''}^{\geqslant}=
\bigoplus_{n\geqslant 0}\mathcal{O}_{P''}(n)$ se identifica con $\bigoplus_{n\geqslant 0}\mathcal{L}''^{\otimes n}$, y por tanto $C_{P''}=\operatorname{Spec}(\mathcal{S}_{P''}^{\geqslant})$ con $\mathbf{V}(\mathcal{L}'')$;
se sabe además por \sref{II.8.7.3-fr} que $C_{P''}$ es $C''$-isomorfo a $\operatorname{Proj}(\mathcal{S}''^{\natural})$;
por definición de $\mathcal{S}''$, $\mathcal{S}''^{\natural}$ está generada por $\mathcal{S}_1''^{\natural}$, y $\mathcal{S}_1''^{\natural}$ es de tipo finito sobre $\mathcal{S}_0''^{\natural}=\mathcal{S}''$ (\sref{II.8.2.10-fr}[i y iii]), por lo que $\operatorname{Proj}(\mathcal{S}''^{\natural})$ es \emph{proyectivo} sobre $C''$ \sref{II.5.5.1-fr}.
Consideremos el diagrama
\begin{equation}
\label{II.8.8.4.1-fr}
  \xymatrix{
    \mathbf{V}(\mathcal{L})\ar[r]^{g}\ar[d]_{u} &
    \operatorname{Spec}(\mathcal{S})=C\ar[d]^{v}\\
    \mathbf{V}(\mathcal{L}'')\ar[r]_{g''} &
    \operatorname{Spec}(\mathcal{S}'')=C''
  }
\tag{8.8.4.1}
\end{equation}
donde $g$ y $g''$ corresponden, por \sref{II.1.5.6-fr}, a los $f$-morfismos canónicos
\[
  \mathcal{S}\longrightarrow
  \bigoplus_{n\geqslant 0}\mathcal{L}^{\otimes n}
  \quad\text{y}\quad
  \mathcal{S}''\longrightarrow
  \bigoplus_{n\geqslant 0}\mathcal{L}''^{\otimes n}
\]
de \sref{II.3.3.2.3-fr} (véase \sref{II.8.8.5-fr} más abajo), y $v$ al morfismo de inclusión $\mathcal{S}''\to\mathcal{S}$ y $u$ al morfismo de inclusión $\bigoplus_{n\geqslant 0}\mathcal{L}^{\otimes nd}\to
\bigoplus_{n\geqslant 0}\mathcal{L}^{\otimes n}$;
es inmediato por \sref{II.3.3.2-fr} que este diagrama es conmutativo.
Acabamos de ver que $g''$ es un morfismo proyectivo;
por otra parte, $u$ es un morfismo \emph{finito}.
En efecto, como la cuestión es local sobre $X$, se puede suponer que $X$ es afín de anillo $A$ y que $\mathcal{L}=\mathcal{O}_X$;
todo se reduce entonces a observar que el anillo $A[T]$ es un módulo de tipo finito sobre su subanillo $A[T^d]$ (donde $T$ es una indeterminada).
Como $Y$ es un esquema cuasicompacto y $C''$ es afín sobre $Y$, $C''$ es también un esquema cuasicompacto,
\oldpage[II]{181}
y por tanto $g''\circ u$ es un morfismo proyectivo \sref{II.5.5.5-fr}[ii];
por la conmutatividad de \sref{II.8.8.4.1-fr}, $v\circ g$ es también proyectivo y, como $v$ es afín, y por tanto separado, se concluye finalmente que $g$ es proyectivo \sref{II.5.5.5-fr}[v].
\end{proof}
\begin{env}[8.8.5]
\label{II.8.8.5-fr}
Retomemos la situación de \sref{II.8.8.1-fr}.
Veremos que el morfismo $g:\mathbf{V}(\mathcal{L})\to C$ también puede definirse de una manera válida para todo $\mathcal{O}_X$-módulo invertible $\mathcal{L}$ (no necesariamente amplio).
Para ello, consideremos el $f$-morfismo
\begin{equation}
\label{II.8.8.5.1-fr}
  \tau^\flat:\mathcal{S}\longrightarrow
  \bigoplus_{n\geqslant 0}\mathcal{L}^{\otimes n}
\tag{8.8.5.1}
\end{equation}
correspondiente al morfismo $\tau$ de \sref{II.8.8.1.2-fr}.
Este induce, por \sref{II.1.5.6-fr}, un morfismo $g':V\to C$ tal que, si $\pi:V\to X$ y $\psi:C\to Y$ son los morfismos estructurales, los diagramas
\begin{equation}
\label{II.8.8.5.2-fr}
  \xymatrix{
    X\ar[d]_{f} & V\ar[l]_{\pi}\ar[d]^{g'}\\
    Y & C\ar[l]^{\psi}
  }
  \qquad
  \xymatrix{
    X\ar[r]^{j}\ar[d]_{f} & V\ar[d]^{g'}\\
    Y\ar[r]_{\varepsilon} & C
  }
\tag{8.8.5.2}
\end{equation}
son conmutativos (\eref{II.8.5.1.2-fr} y \eref{II.8.5.1.3-fr}).
Mostraremos que (si se supone que $\mathcal{L}$ es $f$-amplio) \emph{los morfismos $g$ y $g'$ son idénticos}.

Como la cuestión es local sobre $Y$, se puede suponer que $Y=\operatorname{Spec}(A)$ es afín y, por \sref{II.8.8.1.3-fr}, identificar $X$ con un abierto de $P=\operatorname{Proj}(S)$, donde $S=A\oplus\bigoplus_{n\geqslant 1}\Gamma(X,\mathcal{L}^{\otimes n})$;
se deduce entonces de \sref{II.8.8.1.4-fr} que $\Gamma(X,\mathcal{O}_P(n))=\Gamma(X,\mathcal{L}^{\otimes n})$ para todo $n\in\mathbf{Z}$.
Teniendo en cuenta la definición de $h=\operatorname{Spec}(\alpha)$, donde $\alpha$ es el $p$-morfismo canónico $\widetilde{S}\to\mathcal{S}_P^{\geqslant}$ \sref{II.8.6.1.2-fr}, hay que demostrar que la restricción a $X$ de $\alpha^\sharp:p^*(\widetilde{S})\to\mathcal{S}_P^{\geqslant}$ es idéntica a $\tau$.
Teniendo en cuenta \sref[0]{0.4.4.3-fr}, basta demostrar que, si se compone el homomorfismo canónico $\alpha_n:S_n\to\Gamma(P,\mathcal{O}_P(n))$ con el homomorfismo de restricción $\Gamma(P,\mathcal{O}_P(n))\longrightarrow\Gamma(X,\mathcal{O}_P(n))=\Gamma(X,\mathcal{L}^{\otimes n})$, se obtiene la identidad para todo $n>0$;
pero esto se desprende inmediatamente de la definición del álgebra $S$ y de la de $\alpha_n$ \sref{II.2.6.2-fr}.
\end{env}
\begin{proposition}[8.8.6]
\label{II.8.8.6-fr}
Supongamos (con las notaciones de \sref{II.8.8.5-fr}) que, si se escribe $f=(f_0,\lambda)$, el homomorfismo $\lambda:\mathcal{O}_Y\to f_*(\mathcal{O}_X)$ es biyectivo;
entonces:
\begin{enumerate}
  \item[\rm{(i)}] si se escribe $g=(g_0,\mu)$, entonces $\mu:\mathcal{O}_C\to g_*(\mathcal{O}_V)$ es un isomorfismo; y
  \item[\rm{(ii)}] si $X$ es íntegro \textup{(resp. \emph{localmente íntegro y normal})}, entonces $C$ es íntegro \textup{(resp. \emph{normal})}.
\end{enumerate}
\end{proposition}

\begin{proof}
En efecto, el $f$-morfismo $\tau^\flat$ es entonces un \emph{isomorfismo}
\[
  \tau^\flat:\mathcal{S}=\psi_*(\mathcal{O}_C)
  \longrightarrow f_*(\pi_*(\mathcal{O}_V))
  =\psi_*(g_*(\mathcal{O}_V))
\]
y el $Y$-morfismo $g$ puede considerarse como aquel para el cual el homomorfismo $\mathcal{A}(g)$ \sref{II.1.1.2-fr} es igual a $\tau^\flat$.
Para ver que $\mu$ es un isomorfismo de $\mathcal{O}_C$-módulos, basta, por \sref{II.1.4.2-fr}, ver que $\mathcal{A}(\mu):\psi_*(\mathcal{O}_C)\longrightarrow\psi_*(g_*(\mathcal{O}_V))$ es un isomorfismo.
Pero, por la Definición~\sref{II.1.1.2-fr}, se tiene $\mathcal{A}(\mu)=\mathcal{A}(g)$, de donde se sigue la conclusión de (i).

Para demostrar (ii), se puede reducir al caso en que $Y$ es afín, y entonces $\mathcal{S}=\widetilde{S}$, con
\oldpage[II]{182}
$S=\bigoplus_{n\geqslant 0}\Gamma(X,\mathcal{L}^{\otimes n})$;
la hipótesis de que $X$ es íntegro implica que el anillo $S$ es íntegro \sref[I]{I.7.4.4-fr}, y por tanto también lo es $C$ \sref[I]{I.5.1.4-fr}.
Para demostrar que $C$ es normal, utilizaremos el lema siguiente:

\begin{lemma}[8.8.6.1]
\label{II.8.8.6.1-fr}
Sea $Z$ un preesquema íntegro normal.
Entonces el anillo $\Gamma(Z,\mathcal{O}_Z)$ es íntegro e íntegramente cerrado.
\end{lemma}

\begin{proof}
Se sigue de \sref[I]{I.8.2.1.1-fr} que $\Gamma(Z,\mathcal{O}_Z)$ es la intersección, en el cuerpo de funciones racionales $R(Z)$, de los anillos íntegramente cerrados $\mathcal{O}_z$ para todos los $z\in Z$.
\end{proof}

Dicho esto, demostremos primero que $V$ es \emph{localmente íntegro} y \emph{normal};
para ello, se puede reducir al caso en que $X=\operatorname{Spec}(A)$ es afín, de anillo $A$ íntegro e íntegramente cerrado \sref{II.6.3.8-fr}, y en que $\mathcal{L}=\mathcal{O}_X$.
Como entonces $V=\operatorname{Spec}(A[T])$, y $A[T]$ es íntegro e íntegramente cerrado \cite[p.~99]{II-24}, esto demuestra nuestra afirmación.
Para todo abierto afín $U$ de $C$, $g^{-1}(U)$ es cuasicompacto, pues el morfismo $g$ es cuasicompacto;
como $V$ es localmente íntegro, las componentes conexas de $g^{-1}(U)$ son preesquemas íntegros abiertos en $g^{-1}(U)$ y, por tanto, son finitas en número; además, como $V$ es normal, estos preesquemas son también normales \sref{II.6.3.8-fr}.
Entonces $\Gamma(U,\mathcal{O}_C)$, que es igual a $\Gamma(g^{-1}(U),\mathcal{O}_V)$ por (i), es el producto directo finito de anillos íntegros e íntegramente cerrados \sref{II.8.8.6.1-fr}, lo que demuestra que $C$ es normal \sref{II.6.3.4-fr}.
\end{proof}
\subsection{El criterio de amplitud de Grauert: enunciado}
\label{subsection:II.8.9-fr}

Nos proponemos demostrar que las propiedades probadas en \sref{II.8.8.2-fr} \emph{caracterizan} los $\mathcal{O}_X$-módulos amplios para $f$ y, más precisamente, demostrar el criterio siguiente:
\begin{theorem}[8.9.1]
\label{II.8.9.1-fr}
\emph{(Criterio de Grauert).}
Sean $Y$ un preesquema, $p:X\to Y$ un morfismo separado y cuasicompacto, y $\mathcal{L}$ un $\mathcal{O}_X$-módulo invertible.
Para que $\mathcal{L}$ sea amplio relativamente a $p$, es necesario y suficiente que existan un $Y$-preesquema $C$, una $Y$-sección $\varepsilon:Y\to C$ de $C$ y un $Y$-morfismo $q:\mathbf{V}(\mathcal{L})\to C$ que satisfagan las propiedades siguientes:
\begin{enumerate}
  \item[\rm{(i)}] el diagrama
    \begin{equation}
  \label{II.8.9.1.1-fr}
    \xymatrix{
      X\ar[r]^{j}\ar[d]_{p} & \mathbf{V}(\mathcal{L})\ar[d]^{q}\\
      Y\ar[r]_{\varepsilon} & C
    }
  \tag{8.9.1.1}
  \end{equation}
    es conmutativo, donde $j$ es la sección nula del fibrado vectorial $\mathbf{V}(\mathcal{L})$; y
  \item[\rm{(ii)}] la restricción de $q$ a $\mathbf{V}(\mathcal{L})\setmin j(X)$ es una inmersión abierta cuasicompacta
    \[
    \mathbf{V}(\mathcal{L})\setmin j(X)\longrightarrow C
  \]
    cuya imagen no interseca $\varepsilon(Y)$.
\end{enumerate}
\end{theorem}

Se observa que, si $C$ es \emph{separado} sobre $Y$, en la condición~(ii) se puede suprimir la hipótesis de que la inmersión abierta sea cuasicompacta;
para ver que esta propiedad (la cuasicompacidad) es, en efecto, consecuencia de las demás condiciones, se puede reducir al caso en que $Y$ es afín, y la afirmación se sigue entonces de \sref[I]{I.5.5.1-fr}[i] y \sref[I]{I.5.5.10-fr}.
También se puede suprimir
\oldpage[II]{183}
la misma hipótesis si se supone que $X$ es noetheriano, pues entonces $V$ también es noetheriano, y la afirmación se sigue de \sref[I]{I.6.3.5-fr}.

\begin{corollary}[8.9.2]
\label{II.8.9.2-fr}
Si el morfismo $p:X\to Y$ es propio, en el enunciado del Teorema~\sref{II.8.9.1-fr} se puede suponer que $q$ es propio y sustituir «inmersión abierta» por «isomorfismo».
\end{corollary}

De manera más sugestiva, se puede decir (siempre que $p:X\to Y$ sea \emph{propio}) que $\mathcal{L}$ \emph{es amplio relativamente a $p$ si y solo si se puede} «contraer» \emph{la sección nula del fibrado vectorial $\mathbf{V}(\mathcal{L})$ al preesquema de base} $Y$.
Un caso particular importante es aquel en que $Y$ es el espectro de un cuerpo y en que la operación de «contracción» consiste en contraer la sección nula de $\mathbf{V}(\mathcal{L})$ \emph{a un solo punto}.

\begin{env}[8.9.3]
\label{II.8.9.3-fr}
Tanto la necesidad de las condiciones del Teorema \sref{II.8.9.1-fr} como el Corolario \sref{II.8.9.2-fr} se siguen inmediatamente de \sref{II.8.8.2-fr} y \sref{II.8.8.4-fr}.

Para demostrar que las condiciones de \sref{II.8.9.1-fr} son suficientes, consideremos una situación ligeramente más general.
Para ello, pongamos (con las notaciones de \sref{II.8.8.2-fr})
\[
  \mathcal{S}'=\bigoplus_{n\geqslant 0}\mathcal{L}^{\otimes n}
\]
y
\[
  V=\mathbf{V}(\mathcal{L})=\operatorname{Spec}(\mathcal{S}').
\]

El subpreesquema cerrado $j(X)$, sección nula de $\mathbf{V}(\mathcal{L})$, está definido por el haz cuasicoherente de ideales $\mathcal{J}=(\mathcal{S}'_+)^{\sim}$ de $\mathcal{O}_V$ \sref{II.1.4.10-fr}.
Este $\mathcal{O}_V$-módulo es \emph{invertible}, pues esta propiedad es local sobre $X$, y esto se reduce a observar que el ideal $TA[T]$ de un anillo de polinomios $A[T]$ es un $A[T]$-módulo libre cíclico.
Además, es inmediato (de nuevo porque la cuestión es local sobre $X$) que
\[
  \mathcal{L}=j^*(\mathcal{J})
\]
y
\[
  j_*(\mathcal{L})=\mathcal{J}/\mathcal{J}^2.
\]

Por otra parte, si
\[
  \pi:\mathbf{V}(\mathcal{L})\longrightarrow X
\]
es el morfismo estructural, entonces $\pi_*(\mathcal{J})=\mathcal{S}'_+$ y $\pi_*(\mathcal{J}/\mathcal{J}^2)=\mathcal{L}$;
hay, por tanto, homomorfismos canónicos $\mathcal{L}\to\pi_*(\mathcal{J})\to\mathcal{L}$, siendo el primero la inyección canónica $\mathcal{L}\to\mathcal{S}'_+$ y el segundo la proyección canónica de $\mathcal{S}'_+$ sobre $\mathcal{S}'_1=\mathcal{L}$, y su composición es la identidad.
También se puede incluir canónicamente $\pi_*(\mathcal{J})=\mathcal{S}'_+=\bigoplus_{n\geqslant 1}\mathcal{L}^{\otimes n}$ en el \emph{producto} $\prod_{n\geqslant 1}\mathcal{L}^{\otimes n}=\varprojlim_n\pi_*(\mathcal{J}/\mathcal{J}^{n+1})$ (pues $\pi_*(\mathcal{J}/\mathcal{J}^{n+1})=
\mathcal{L}\oplus\mathcal{L}^{\otimes2}\oplus\ldots\oplus
\mathcal{L}^{\otimes n}$), y se tienen así dos homomorfismos canónicos
\begin{equation}
\label{II.8.9.3.1-fr}
  \mathcal{L}\longrightarrow
  \varprojlim\pi_*(\mathcal{J}/\mathcal{J}^{n+1})
  \longrightarrow\mathcal{L}
\tag{8.9.3.1}
\end{equation}
cuya composición es la identidad.

Con esto, la generalización de \sref{II.8.9.1-fr} que vamos a demostrar es la siguiente:
\end{env}
\begin{proposition}[8.9.4]
\label{II.8.9.4-fr}
Sean $Y$ un preesquema, $V$ un $Y$-preesquema y $X$ un subpreesquema cerrado de $V$ definido por un ideal $\mathcal{J}$ de $\mathcal{O}_V$, que es un $\mathcal{O}_V$-módulo \emph{invertible};
si $j:X\to V$ es
\oldpage[II]{184}
la inyección canónica, se pone entonces $\mathcal{L}=j^*(\mathcal{J})=
\mathcal{J}\otimes_{\mathcal{O}_V}\mathcal{O}_X$, de modo que $j_*(\mathcal{L})=\mathcal{J}/\mathcal{J}^2$.
Se supone que el morfismo estructural $p:X\to Y$ es separado y cuasicompacto, y que se cumplen las condiciones siguientes:
\begin{enumerate}
  \item[\rm{(i)}] existe un $Y$-morfismo $\pi:V\to X$ de tipo finito tal que $\pi\circ j=1_X$, y por tanto $\pi_*(\mathcal{J}/\mathcal{J}^2)=\mathcal{L}$;
  \item[\rm{(ii)}] existe un homomorfismo de $\mathcal{O}_X$-módulos $\varphi:\mathcal{L}\longrightarrow\varprojlim\pi_*(\mathcal{J}/\mathcal{J}^{n+1})$ tal que la composición
    \[
    \mathcal{L}\xrightarrow{\varphi}
    \varprojlim\pi_*(\mathcal{J}/\mathcal{J}^{n+1})
    \xrightarrow{\alpha}\pi_*(\mathcal{J}/\mathcal{J}^2)=\mathcal{L}
  \]
    (donde $\alpha$ es el homomorfismo canónico) es la identidad;
  \item[\rm{(iii)}] existen un $Y$-preesquema $C$, una $Y$-sección $\varepsilon$ de $C$ y un $Y$-morfismo $q:V\to C$ tales que el diagrama
    \begin{equation}
  \label{II.8.9.4.1-fr}
    \xymatrix{
      X\ar[r]^{j}\ar[d]_{p} & V\ar[d]^{q}\\
      Y\ar[r]_{\varepsilon} & C
    }
  \tag{8.9.4.1}
  \end{equation}
    es conmutativo; y
  \item[\rm{(iv)}] la restricción de $q$ a $W=V\setmin j(X)$ es una inmersión abierta cuasicompacta en $C$, cuya imagen no interseca $\varepsilon(Y)$.
\end{enumerate}
Entonces $\mathcal{L}$ es amplio relativo a $p$.
\end{proposition}


\subsection{Criterio de amplitud de Grauert: demostración}
\label{subsection:II.8.10-fr}

\begin{lemma}[8.10.1]
\label{II.8.10.1-fr}
Sean $\pi:V\to X$ un morfismo, $j:X\to V$ una $X$-sección de $V$ que es además una inmersión cerrada, y $\mathcal{J}$ un ideal cuasicoherente de $\mathcal{O}_V$ que define el subpreesquema cerrado de $V$ asociado a $j$.
Entonces se cumplen todas las afirmaciones siguientes.
\begin{enumerate}
  \item[\rm{(i)}] Para todo $n\geqslant 0$, $\pi_*(\mathcal{O}_V/\mathcal{J}^{n+1})$ y $\pi_*(\mathcal{J}/\mathcal{J}^{n+1})$ son $\mathcal{O}_X$-módulos cuasicoherentes, y $\pi_*(\mathcal{O}_V/\mathcal{J})=\mathcal{O}_X$ y $\pi_*(\mathcal{J}/\mathcal{J}^2)=j^*(\mathcal{J})$.
  \item[\rm{(ii)}] Si $X=\{\xi\}=\operatorname{Spec}(k)$, donde $k$ es un cuerpo, entonces $\varprojlim\pi_*(\mathcal{O}_V/\mathcal{J}^{n+1})$ es isomorfo a la completación separada del anillo local $\mathcal{O}_{j(\xi)}$ para la topología $\mathfrak{m}_{j(\xi)}$-ádica.
  \item[\rm{(iii)}] Supongamos que $\mathcal{J}$ es un $\mathcal{O}_V$-módulo invertible (lo cual implica que
    \[
    \mathcal{L}=j^*(\mathcal{J})=
    \pi_*(\mathcal{J}/\mathcal{J}^2)
  \]
    es un $\mathcal{O}_X$-módulo invertible), y que existe un homomorfismo $\varphi:\mathcal{L}\longrightarrow\varprojlim\pi_*(\mathcal{J}/\mathcal{J}^{n+1})$ tal que la composición $\mathcal{L}\xrightarrow{\varphi}\varprojlim\pi_*(\mathcal{J}/\mathcal{J}^{n+1})\xrightarrow{\alpha}\pi_*(\mathcal{J}/\mathcal{J}^2)$ (donde $\alpha$ es el homomorfismo canónico) es la identidad.
    Si se escribe $\mathcal{S}=\bigoplus_{n\geqslant0}\mathcal{L}^{\otimes n}$, entonces $\varphi$ induce canónicamente un isomorfismo de $\mathcal{O}_X$-álgebras del completado $\widehat{\mathcal{S}}$ de $\mathcal{S}$ respecto de su filtración canónica (dicho completado es isomorfo al producto $\prod_{n\geqslant0}\mathcal{L}^{\otimes n}$) a $\varprojlim\pi_*(\mathcal{O}_V/\mathcal{J}^{n+1})$.
\end{enumerate}
\end{lemma}

\begin{proof}
Observemos ante todo que el soporte del $\mathcal{O}_V$-módulo $\mathcal{O}_V/\mathcal{J}^{n+1}$ es $j(X)$, y que el soporte de $\mathcal{J}/\mathcal{J}^{n+1}$ está contenido en $j(X)$.
En el caso de (ii), $j(X)$ es un punto cerrado $j(\xi)$ de $V$,
\oldpage[II]{185}
y, por definición, $\pi_*(\mathcal{O}_V/\mathcal{J}^{n+1})$ es el tallo de $\mathcal{O}_V/\mathcal{J}^{n+1}$ en el punto $j(\xi)$, es decir, poniendo $C=\mathcal{O}_{j(\xi)}$ y denotando por $\mathfrak{m}$ el ideal maximal de $C$, el $C$-módulo $C/\mathfrak{m}^{n+1}$;
la afirmación (ii) es entonces evidente.

Para demostrar (i), observemos que la cuestión es local sobre $X$;
se puede, por tanto, restringir al caso en que $X$ es afín.
Sea $U$ un abierto afín de $V$;
entonces $j(X)\cap U$ es un abierto afín de $j(X)$, por lo que $U_0=\pi(j(X)\cap U)$, que le es isomorfo, es un abierto afín de $X$;
para todo abierto afín $W_0\subset U_0$ de $X$, $W=\pi^{-1}(W_0)\cap U$ es un abierto afín de $V$, puesto que $X$ es un esquema \sref[I]{I.5.5.10-fr};
en particular, $U'=U\cap\pi^{-1}(U_0)$ es un abierto afín de $V$, y es claro que $\pi(U')=U_0$ y $j(U_0)=j(X)\cap U$.
Ahora bien, por definición,
\[
  \Gamma(W_0,\pi_*(\mathcal{O}_V/\mathcal{J}^{n+1}))
  =\Gamma(\pi^{-1}(W_0),\mathcal{O}_V/\mathcal{J}^{n+1})~;
\]
pero, como todo punto de $\pi^{-1}(W_0)$ que no pertenece a $j(W_0)$ posee una vecindad abierta en $\pi^{-1}(W_0)$ que no interseca $j(X)$ y en la cual $\mathcal{O}_V/\mathcal{J}^{n+1}$ es, por tanto, nulo, es claro que las secciones de $\mathcal{O}_V/\mathcal{J}^{n+1}$ sobre $\pi^{-1}(W_0)$ y sobre $W$ se corresponden biunívocamente.
En otros términos, si $\pi'$ es la restricción de $\pi$ a $U'$, los $(\mathcal{O}_X|U_0)$-módulos
\[
  \pi_*(\mathcal{O}_V/\mathcal{J}^{n+1})|U_0
  \quad\hbox{y}\quad
  \pi'_*\bigl((\mathcal{O}_V/\mathcal{J}^{n+1})|U'\bigr)
\]
son idénticos.
Como $U'$ y $U_0$ son afines y los $U_0$ recubren $X$, se concluye entonces \sref[I]{I.1.6.3-fr} que $\pi_*(\mathcal{O}_V/\mathcal{J}^{n+1})$ es cuasicoherente, y la demostración es idéntica para $\pi_*(\mathcal{J}/\mathcal{J}^{n+1})$.

Para demostrar finalmente (iii), observemos que $\mathcal{S}$ es precisamente $\mathbf{S}_{\mathcal{O}_X}(\mathcal{L})$;
por consiguiente, $\varphi$ induce canónicamente un homomorfismo de $\mathcal{O}_X$-álgebras
\[
  \psi:\mathcal{S}\longrightarrow
  \varprojlim\pi_*(\mathcal{O}_V/\mathcal{J}^{n+1})
\]
\sref{II.1.7.4-fr};
además, este homomorfismo envía $\mathcal{L}^{\otimes n}$ a $\varprojlim_m\pi_*(\mathcal{J}^n/\mathcal{J}^{m+1})$, por lo que es continuo para las topologías consideradas y, en consecuencia, se prolonga en un homomorfismo
\[
  \widehat{\psi}:\widehat{\mathcal{S}}\longrightarrow
  \varprojlim\pi_*(\mathcal{O}_V/\mathcal{J}^{n+1}).
\]
Para ver que se trata efectivamente de un isomorfismo, se puede, como en la demostración de (i), reducir al caso en que $X=\operatorname{Spec}(A)$ y $V=\operatorname{Spec}(B)$ son afines, con $\mathcal{J}=\widetilde{\mathfrak{J}}$, donde $\mathfrak{J}$ es un ideal de $B$;
a $\pi$ corresponde una inyección $A\to B$ que identifica $A$ con un subanillo de $B$ \emph{complementario} a $\mathfrak{J}$, y $\mathcal{L}$ (resp. $\pi_*(\mathcal{O}_V/\mathcal{J}^{n+1})$) es el $\mathcal{O}_X$-módulo cuasicoherente asociado al $A$-módulo $L=\mathfrak{J}/\mathfrak{J}^2$ (resp. $B/\mathfrak{J}^{n+1}$).
Como $\mathcal{J}$ es un $\mathcal{O}_V$-módulo \emph{invertible}, se puede suponer además que $\mathfrak{J}=Bt$, donde $t$ no es divisor de cero en $B$.
De la relación $B=A\oplus Bt$ se deduce entonces que, para todo $n>0$,
\[
  B=A\oplus At\oplus At^2\oplus\ldots\oplus At^n\oplus Bt^{n+1}
\]
y, por tanto, existe un $A$-isomorfismo canónico del anillo de series formales $A[[T]]$ sobre $C=\varprojlim B/\mathfrak{J}^{n+1}$ que envía $T$ a $t$.
También se tiene $L=A\bar t$, donde $\bar t$ es la clase de $t$ módulo $Bt^2$, y, por hipótesis, el homomorfismo $\varphi$ envía $\bar t$ a un elemento $t'\in C$ congruente con $t$ módulo $Ct^2$.
Se deduce entonces por inducción sobre $n$ que
\[
  A\oplus At'\oplus\ldots\oplus At'{}^n\oplus Ct^{n+1}
  =A\oplus At\oplus\ldots\oplus At^n\oplus Ct^{n+1}
\]
lo que prueba que el homomorfismo $\widehat{\psi}$ corresponde efectivamente a un isomorfismo de $\prod_{n\geqslant0}L^{\otimes n}$ sobre $C$.
\end{proof}

\begin{lemma}[8.10.2]
\label{II.8.10.2-fr}
Bajo las hipótesis del Lema~\sref{II.8.10.1-fr}, sea $g:X'\to X$ un morfismo,
\oldpage[II]{186}
sea $V'=V\times_X X'$, y sean $\pi':V'\to X'$ y $g':V'\to V$ las proyecciones canónicas, de modo que se tiene el diagrama conmutativo
\[
  \xymatrix{
    V\ar[d]_{\pi} & V'\ar[l]_{g'}\ar[d]^{\pi'}\\
    X & X'\ar[l]^{g}
  }
\]

Entonces $j'=j\times1_{X'}$ es una $X'$-sección de $V'$ que es además una inmersión cerrada, y $\mathcal{J}'=g'^*(\mathcal{J})\mathcal{O}_{V'}$ es el ideal cuasicoherente de $\mathcal{O}_{V'}$ que define el subpreesquema cerrado de $V'$ asociado a $j'$.
Además,
\[
  \pi'_*(\mathcal{O}_{V'}/\mathcal{J}'^{n+1})
  =g^*(\pi_*(\mathcal{O}_V/\mathcal{J}^{n+1})).
\]
Por último, $\mathcal{J}'$ es un $\mathcal{O}_{V'}$-módulo canónicamente isomorfo a $g'^*(\mathcal{J})$ y, en particular, es invertible si $\mathcal{J}$ es un $\mathcal{O}_V$-módulo invertible.
\end{lemma}

\begin{proof}
El hecho de que $j'$ sea una inmersión cerrada se deduce de \sref[I]{I.4.3.1-fr}, y es una $X'$-sección de $V'$ por la funtorialidad del cambio de base de preesquemas.
Además, si $Z$ (resp. $Z'$) es el subpreesquema cerrado de $V$ (resp. $V'$) asociado a $j$ (resp. $j'$), entonces $Z'=g'{}^{-1}(Z)$ \sref[I]{I.4.3.1-fr}, y la segunda afirmación se deduce de \sref[I]{I.4.4.5-fr}.
Para demostrar las demás afirmaciones, se ve, como en \sref{II.8.10.1-fr}, que se puede reducir al caso en que $X$, $V$ y $X'$ (y, por tanto, también $V'$) sean afines;
conservemos las notaciones de la demostración de \sref{II.8.10.1-fr} y sea $X'=\operatorname{Spec}(A')$.
Entonces $V'=\operatorname{Spec}(B')$, donde $B'=B\otimes_A A'$, y $\mathcal{J}'=\widetilde{\mathfrak{J}}'$, donde $\mathfrak{J}'=\operatorname{Im}(\mathfrak{J}\otimes_A A')$.
Por tanto,
\[
  B'/\mathfrak{J}'^{n+1}=(B/\mathfrak{J}^{n+1})\otimes_A A'~;
\]
además, como $\mathfrak{J}$ es un factor directo (como $A$-módulo) de $B$, $\mathfrak{J}\otimes_A A'$ es un factor directo (como $A'$-módulo) de $B'$ y, por consiguiente, se identifica canónicamente con $\mathfrak{J}'$.
\end{proof}

\begin{corollary}[8.10.3]
\label{II.8.10.3-fr}
Supongamos verificadas las hipótesis del Lema~\sref{II.8.10.1-fr}, y supongamos además que $\pi$ es de tipo finito y que $\mathcal{J}$ es un $\mathcal{O}_V$-módulo invertible.
Entonces, para todo $x\in X$, el anillo local en el punto $j(x)$ de la fibra $\pi^{-1}(x)$ es un anillo regular (por tanto, íntegro) de dimensión 1, cuyo completado es isomorfo al anillo de series formales $\boldsymbol{k}(x)[[T]]$ (donde $T$ es una indeterminada);
además, existe exactamente una componente irreducible de $\pi^{-1}(x)$ que contiene a $j(x)$.
\end{corollary}

\begin{proof}
Como $\pi^{-1}(x)=V\times_X\operatorname{Spec}(\boldsymbol{k}(x))$, \sref{II.8.10.2-fr} reduce al caso en que $X$ es el espectro de un cuerpo $K$.
Como $\pi$ es de tipo finito \sref[I]{I.6.3.4-fr}[iv], $\mathcal{O}_{j(x)}$ es entonces un anillo local noetheriano y, por tanto, separado para la topología $\mathfrak{m}_{j(x)}$-ádica \sref[0]{0.7.3.5-fr};
de \sref{II.8.10.1-fr}[ii y iii] se sigue que el completado de este anillo es isomorfo a $K[[T]]$ y, por consiguiente, $\mathcal{O}_{j(x)}$ es regular y de dimensión 1 (\cite[p.~17-01, th.~1]{I-1});
por último, como $\mathcal{O}_{j(x)}$ es íntegro, $j(x)$ pertenece exactamente a una de las componentes irreducibles (en número finito) de $V$ \sref[I]{I.5.1.4-fr}.
\end{proof}

\begin{corollary}[8.10.4]
\label{II.8.10.4-fr}
Supongamos verificadas las hipótesis del Lema~\sref{II.8.10.1-fr}, y supongamos además que $\mathcal{J}$ es un $\mathcal{O}_V$-módulo invertible.
Sea $W=V\setmin j(X)$;
para todo ideal cuasicoherente $\mathcal{K}$ de $\mathcal{O}_X$, pongamos $\mathcal{K}_V=\pi^*(\mathcal{K})\mathcal{O}_V$ y $\mathcal{K}_W=\mathcal{K}_V|W$.
Entonces $\mathcal{K}_V$ es el mayor ideal cuasicoherente de $\mathcal{O}_V$ cuya restricción a $W$ es $\mathcal{K}_W$.
\end{corollary}

\begin{proof}
En efecto, se ve, como en \sref{II.8.10.1-fr}, que la cuestión es local sobre $X$ y sobre $V$;
se puede, pues, conservar las notaciones de la demostración de \sref{II.8.10.1-fr}, con $\mathfrak{J}=Bt$, donde $t$ no es divisor de cero en $B$.
Además, $W=\operatorname{Spec}(B_t)$ y $\mathcal{K}=\widetilde{\mathfrak{K}}$, donde $\mathfrak{K}$ es un ideal de $A$;
de donde $\pi^*(\mathcal{K})\mathcal{O}_V=(\mathfrak{K}.B)^{\sim}$ \sref[I]{I.1.6.9-fr}, $\mathcal{K}_W=(\mathfrak{K}.B_t)^{\sim}$, y el mayor ideal
\oldpage[II]{187}
de $B$ cuya imagen canónica en $B_t$ sea $\mathfrak{K}.B_t$ es la imagen inversa de $\mathfrak{K}.B_t$, es decir, el conjunto de los $s\in B$ tales que, para algún entero $n>0$, se tiene $t^ns\in\mathfrak{K}.B$.
Hay que demostrar que esta última relación implica que $s\in\mathfrak{K}.B$, o también que la imagen canónica de $t$ no es divisor de cero en $B/\mathfrak{K}B=(A/\mathfrak{K})\otimes_A B$, lo que se deduce de \sref{II.8.10.2-fr} aplicado a $X'=\operatorname{Spec}(A/\mathfrak{K})$.
\end{proof}

\begin{corollary}[8.10.5]
\label{II.8.10.5-fr}
Supongamos verificadas las hipótesis de \sref{II.8.10.3-fr};
sean $W=V\setmin j(X)$, $x$ un punto de $X$, $\mathcal{K}$ un ideal cuasicoherente de $\mathcal{O}_X$, y $z$ el punto genérico de la componente irreducible de $\pi^{-1}(x)$ que contiene a $j(x)$ \sref{II.8.10.3-fr}.
\begin{enumerate}
    \item[\rm{(i)}] Sea $g$ una sección de $\mathcal{O}_V$ sobre $V$ tal que $g|W$ sea una sección de $\mathcal{K}_W$ sobre $W$ (con las notaciones de \sref{II.8.10.4-fr}).
        Entonces $g$ es una sección de $\mathcal{K}_V$;
        si además $g(z)\neq0$, y si, para todo entero $m>0$, se denota por $g_m^x$ el germen en el punto $x$ de la imagen canónica $g_m$ de $g$ en $\Gamma(X,\pi_*(\mathcal{O}_V/\mathcal{J}^{m+1}))$, entonces existe un entero $m>0$ tal que la imagen de $g_m^x$ en
        \[
    (\pi_*(\mathcal{O}_V/\mathcal{J}^{m+1}))_x
    \otimes_{\mathcal{O}_x}\boldsymbol{k}(x)
  \]
        sea $\neq0$.
    \item[\rm{(ii)}] Supongamos además que se cumplen las condiciones de \sref{II.8.10.1-fr}[iii].
        Entonces, si existe una sección $g$ de $\mathcal{K}_V$ sobre $V$ tal que $g(z)\neq0$, existe un entero $n\geqslant0$ y una sección $f$ de $\mathcal{K}.\mathcal{L}^{\otimes n}=
  \mathcal{K}\otimes\mathcal{L}^{\otimes n}\subset\mathcal{L}^{\otimes n}$ tal que $f(x)\neq0$.
        Si $g$ es una sección de $\mathcal{J}$, se puede tomar $n>0$.
\end{enumerate}
\end{corollary}

\begin{proof}
\begin{enumerate}
    \item[\rm{(i)}] Como el ideal de $\mathcal{O}_W$ engendrado por $g|W$ está contenido en $\mathcal{K}_W$ por hipótesis, el ideal de $\mathcal{O}_V$ engendrado por $g$ está contenido en $\mathcal{K}_V$ por \sref{II.8.10.4-fr}; en otras palabras, $g$ es una sección de $\mathcal{K}_V$.
        Para demostrar la segunda afirmación de (i), se puede suponer de nuevo que $X$ y $V$ son afines y conservar las notaciones de \sref{II.8.10.1-fr};
        la fibra $\pi^{-1}(x)$ es entonces afín, con anillo $B'=B\otimes_A \boldsymbol{k}(x)$, y existe en $B'$ un elemento $t'$ que no es divisor de cero y tal que $B'=\boldsymbol{k}(x)\oplus B't'$.
        Como $j(x)$ es una especialización de $z$ y $g(z)\neq0$, necesariamente se tiene $g_{j(x)}\neq0$.
        Pero $\mathcal{O}_{j(x)}$ es un anillo local separado \sref{II.8.10.3-fr}, y por tanto se inyecta en su completado; la imagen de $g$ en este completado no es, pues, nula.
        Ahora bien, este completado es isomorfo a $\varprojlim_n(B'/B't'^{n+1})$ \sref{II.8.10.3-fr};
        si $g'=g\otimes1\in B'$, existe entonces un entero $m$ tal que $g'\notin B't'^{m+1}$, o, lo que es lo mismo, la imagen $g'_m$ de $g'$ en $B'/B't'^{m+1}$ es distinta de cero.
        Pero, como $g'_m$ es precisamente la imagen de $g_m^x$, nuestra afirmación queda demostrada.
    \item[\rm{(ii)}] En virtud de \sref{II.8.10.1-fr}[iii], $\pi_*(\mathcal{O}_V/\mathcal{J}^{m+1})$ es isomorfo a la suma directa de los $\mathcal{L}^{\otimes k}$ para $0\leqslant k\leqslant m$;
        denotemos por $f_k$ la sección de $\mathcal{L}^{\otimes k}$ sobre $X$ que es la componente del elemento de $\bigoplus_{k=0}^m\Gamma(X,\mathcal{L}^{\otimes k})$ correspondiente a $g_m$ mediante este isomorfismo.
        Eligiendo $m$ como en (i), existe, por tanto, un índice $k$ tal que $f_k(x)\neq0$, por (i).
        Para ver que $f_k$ es una sección de $\mathcal{K}\mathcal{L}^{\otimes k}$, basta considerar, como antes, el caso en que $X$ y $V$ son afines, y esto se deduce inmediatamente del hecho de que $g\in\mathfrak{K}.B$ (con las notaciones de \sref{II.8.10.4-fr}).
        La última afirmación se deduce del hecho de que la hipótesis $g\in\Gamma(V,\mathcal{J})$ implica que $f_0=0$.
\end{enumerate}
\end{proof}

\begin{env}[8.10.6]
\label{II.8.10.6-fr}
\emph{Demostración de} \sref{II.8.9.4-fr}.
La cuestión es local sobre $Y$ \sref{II.4.6.4-fr};
como $\varepsilon$ es una $Y$-sección, se puede reemplazar $C$ por una vecindad abierta afín $U$ de un punto de $\varepsilon(Y)$ tal que $\varepsilon(Y)\cap U$ sea cerrado en $U$.
En otras palabras, se puede suponer que $C$ es afín y que $Y$ es un subpreesquema cerrado de $C$ (y, por tanto, también afín) definido por un haz cuasicoherente
\oldpage[II]{188}
$\mathcal{I}$ de ideales de $\mathcal{O}_C$.
Como $p$ es separado y cuasicompacto, $X$ es entonces un esquema cuasicompacto, y se reduce a demostrar que $\mathcal{L}$ es \emph{amplio} \sref{II.4.6.4-fr}.
En virtud del criterio~\sref{II.4.5.2-fr}[\emph{a}], hay que demostrar lo siguiente:
para todo ideal cuasicoherente $\mathcal{K}$ de $\mathcal{O}_X$ y todo punto $x\in X$ que no pertenezca al soporte de $\mathcal{O}_X/\mathcal{K}$, existe un entero $n>0$ y una sección $f$ de $\mathcal{K}\otimes\mathcal{L}^{\otimes n}$ sobre $X$ tal que $f(x)\neq0$.

Para ello, pongamos
\[
  \mathcal{K}_V=\pi^*(\mathcal{K})\mathcal{O}_V
\]
donde \[
  \mathcal{K}_W=\mathcal{K}_V|W,\qquad\hbox{donde}\qquad W=V\setmin j(X)~;
\];
como la restricción de $q$ a $W$ es una inmersión abierta cuasicompacta en $C$, se sigue de \sref[I]{I.9.4.2-fr} que $\mathcal{K}_W$ es la restricción a $W$ de un ideal cuasicoherente $\mathcal{K}'_V$ de $\mathcal{O}_V$ de la forma
\[
  \mathcal{K}'_V=q^*(\mathcal{K}_C)\mathcal{O}_V
\]
donde $\mathcal{K}_C$ es un ideal cuasicoherente de $\mathcal{O}_C$.
Además, como, por hipótesis, $q^{-1}(Y)\subset j(X)$, y como $Y$ está definido por el ideal $\mathcal{I}$, la restricción a $W$ de $q^*(\mathcal{I})\mathcal{O}_V$ es igual a la de $\mathcal{O}_V$; por tanto, $\mathcal{K}_W$ es también la restricción a $W$ de $q^*(\mathcal{I}\mathcal{K}_C)\mathcal{O}_V$, y se puede suponer que $\mathcal{K}_C\subset\mathcal{I}$, de donde
\begin{equation}
\label{II.8.10.6.1-fr}
  \mathcal{K}'_V\subset q^*(\mathcal{I})\mathcal{O}_V\subset\mathcal{J}
\tag{8.10.6.1}
\end{equation}
teniendo en cuenta \sref[I]{I.4.4.6-fr} y la conmutatividad de \sref{II.8.9.4.1-fr}.
Además, se deduce de \sref{II.8.10.4-fr} que
\begin{equation}
\label{II.8.10.6.2-fr}
  \mathcal{K}'_V\subset\mathcal{K}_V.
\tag{8.10.6.2}
\end{equation}
Así las cosas, se sigue de \sref{II.8.10.3-fr} que $j(x)$ pertenece a una sola componente irreducible de $\pi^{-1}(x)$;
sea $z$ el punto genérico de esta componente, y pongamos $z'=q(z)$.
En virtud de \sref{II.8.10.5-fr}, la demostración quedará terminada (teniendo en cuenta \sref{II.8.10.6.1-fr} y \sref{II.8.10.6.2-fr}) si probamos la existencia de una sección $g$ de $\mathcal{K}'_V$ sobre $V$ tal que $g(z)\neq0$.
Pero, por hipótesis, $\mathcal{K}$ tiene una restricción igual a la de $\mathcal{O}_X$ en una vecindad abierta de $x$;
por otra parte, se sigue de \sref{II.8.10.3-fr} que $z\neq j(x)$, y por tanto $z\in W$, de modo que $(\mathcal{K}_W)_z=\mathcal{O}_{V,z}$, y, por definición, $(\mathcal{K}_C)_{z'}=\mathcal{O}_{C,z'}$.
Como $C$ es afín, existe, pues, una sección $g'$ de $\mathcal{K}_C$ sobre $C$ tal que $g'(z')\neq0$; tomando por $g$ la sección de $\mathcal{K}'_V$ que corresponde canónicamente a $g'$, se tiene efectivamente $g(z)\neq0$, lo que termina la demostración.
\end{env}

\begin{remark}[8.10.7]
\label{II.8.10.7-fr}
Ignoramos si, en el enunciado de \sref{II.8.9.4-fr}, la condición~(ii) es o no superflua.
En todo caso, la conclusión deja de ser válida si no se supone la existencia de un $Y$-morfismo $\pi:V\to X$ tal que $\pi\circ j=1_X$;
indiquemos brevemente cómo se puede construir, en efecto, un contraejemplo cuyos detalles no se desarrollarán hasta más adelante.
Se toma $Y=\operatorname{Spec}(k)$, donde $k$ es un cuerpo, y $C=\operatorname{Spec}(A)$, donde $A=k[T_1,T_2]$, con la $Y$-sección $\varepsilon$ correspondiente al homomorfismo de aumentación $A\to k$.
Se denota por $C'$ el esquema obtenido de $C$ por explosión del punto cerrado $a=\varepsilon(Y)$ de $C$;
si $D$ es la imagen inversa de $a$ en $C'$, se considera en $D$ un punto cerrado $b$,
\oldpage[II]{189}
y se denota por $V$ el esquema obtenido de $C'$ por explosión de $b$;
$X$ es el subpreesquema cerrado de $V$ dado por la imagen inversa de $a$ mediante el morfismo estructural $q:V\to C$.
Se muestra ahora que $X$ es la unión de dos componentes irreducibles, $X_1$ y $X_2$, donde $X_1$ es la imagen inversa de $b$ en $V$.
Es inmediato que el ideal $\mathcal{J}$ de $\mathcal{O}_V$ que define $X$ sigue siendo invertible, pero se puede probar que $j^*(\mathcal{J})=\mathcal{L}$ (donde $j$ es la inyección canónica $X\to V$) no es amplio, considerando el «grado» de la imagen inversa de $\mathcal{L}$ sobre $X_1$, que sería $>0$ si $\mathcal{L}$ fuera amplio, pero que se puede demostrar (mediante un cálculo elemental de intersecciones) que es igual a 0.
\end{remark}

\subsection{Unicidad de las contracciones}
\label{subsection:II.8.11-fr}

\begin{lemma}[8.11.1]
\label{II.8.11.1-fr}
Sean $U$ y $V$ dos preesquemas, y sea $h=(h_0,\lambda):U\to V$ un morfismo sobreyectivo.
Se supone que:
\begin{enumerate}
  \item[$1^\circ$] $\lambda:\mathcal{O}_V\to h_*(\mathcal{O}_U)
  =(h_0)_*(\mathcal{O}_U)$ es un isomorfismo.
  \item[$2^\circ$] El espacio subyacente a $V$ se identifica con el espacio cociente del espacio subyacente a $U$ por la relación
  $R:h_0(x)=h_0(y)$
  \textup{(condición que siempre se cumple cuando el morfismo $h$ es \emph{abierto} o \emph{cerrado}, o, \emph{a fortiori}, cuando $h$ es \emph{propio})}.
\end{enumerate}
Entonces, para todo preesquema $W$, la aplicación
\begin{equation}
\label{II.8.11.1.1-fr}
  \operatorname{Hom}(V,W)\longrightarrow\operatorname{Hom}(U,W)
\tag{8.11.1.1}
\end{equation}
que a cada morfismo $v=(v_0,\nu)$ de $V$ en $W$ asocia el morfismo $u=v\circ h=(u_0,\mu)$, es una biyección de $\operatorname{Hom}(V,W)$ sobre el conjunto de los $u$ tales que $u_0$ sea constante en toda fibra $h_0^{-1}(x)$.
\end{lemma}

\begin{proof}
Es claro que, si $u=v\circ h$, de modo que $u_0=v_0\circ h_0$, entonces $u_0$ es constante en todo conjunto $h_0^{-1}(x)$.
Recíprocamente, si $u$ tiene esta propiedad, mostraremos que existe un único $v\in\operatorname{Hom}(V,W)$ tal que $u=v\circ h$.
La existencia y unicidad de la aplicación continua $v_0:V\to W$ tal que $u_0=v_0\circ h_0$ resultan de las hipótesis, pues $h_0$ se identifica con la aplicación canónica de $U$ sobre $U/R$.
También se puede, sustituyendo $V$ por un preesquema isomorfo si es necesario, suponer que $\lambda$ es la identidad;
por hipótesis, $\mu$ es entonces un homomorfismo
\[
  \mu:\mathcal{O}_W\longrightarrow(u_0)_*(\mathcal{O}_U)
  =(v_0)_*((h_0)_*(\mathcal{O}_U))
\]
tal que el homomorfismo correspondiente $\mu^\sharp:u_0^*(\mathcal{O}_W)\to\mathcal{O}_U$ es \emph{local} en cada tallo.
Como $(v_0)_*((h_0)_*(\mathcal{O}_U))=(v_0)_*(\mathcal{O}_V)$, se tiene necesariamente $\nu=\mu$, y todo se reduce a demostrar que el homomorfismo correspondiente $\nu^\sharp:v_0^*(\mathcal{O}_W)\to\mathcal{O}_V$ es local en cada tallo.
Ahora bien, todo $y\in V$ es de la forma $h_0(x)$ para algún $x\in U$;
pongamos $z=v_0(y)=u_0(x)$.
Entonces, por \sref[0]{0.3.5.5-fr}, el homomorfismo $\mu_x^\sharp$ se factoriza como
\[
  \mu_x^\sharp:\mathcal{O}_z
  \xrightarrow{\nu_y^\sharp}\mathcal{O}_y
  \xrightarrow{\lambda_x^\sharp}\mathcal{O}_x.
\]
Por hipótesis, $\lambda_x^\sharp$ y $\mu_x^\sharp$ son homomorfismos locales;
por tanto, $\lambda_x^\sharp$ transforma todo elemento invertible de $\mathcal{O}_y$ en un elemento invertible de $\mathcal{O}_x$;
si $\nu_y^\sharp$ transformase un elemento no invertible de $\mathcal{O}_z$ en un elemento invertible de $\mathcal{O}_y$, entonces $\mu_x^\sharp$ transformaría ese elemento de $\mathcal{O}_z$ en un elemento invertible de $\mathcal{O}_x$, en contradicción con la hipótesis, de donde resulta el lema.
\end{proof}

\begin{corollary}[8.11.2]
\label{II.8.11.2-fr}
Sean $U$ un preesquema íntegro y $V$ un preesquema normal;
entonces todo morfismo $h:U\to V$ que sea universalmente cerrado, birracional y radicial es un isomorfismo.
\end{corollary}

\begin{proof}
Si $h=(h_0,\lambda)$, se sigue de las hipótesis que $h_0$ es inyectiva y cerrada, y que $h_0(U)$
\oldpage[II]{190}
es densa en $V$, por lo que $h_0$ es un \emph{homeomorfismo} de $U$ sobre $V$.
Para demostrar el corolario, bastará probar que $\lambda:\mathcal{O}_V\to(h_0)_*(\mathcal{O}_U)$ es un isomorfismo: se podrá aplicar entonces \sref{II.8.11.1-fr}, que probará que la aplicación \sref{II.8.11.1.1-fr} es biyectiva (pues cada una de las fibras $h_0^{-1}(x)$ se reduce a un solo punto);
por tanto, $h$ será un isomorfismo.
Como la cuestión es evidentemente local sobre $V$, se puede suponer que $V=\operatorname{Spec}(A)$ es afín de anillo íntegro e integralmente cerrado \sref{II.8.8.6.1-fr};
$h$ corresponde entonces \sref[I]{I.2.2.4-fr} a un homomorfismo $\varphi:A\to\Gamma(U,\mathcal{O}_U)$, y todo se reduce a probar que $\varphi$ es un isomorfismo.
Ahora bien, si $K$ es el cuerpo de fracciones de $A$, entonces $\Gamma(U,\mathcal{O}_U)$ tiene, por hipótesis, a $K$ como cuerpo de fracciones, y $A$ es un subanillo de $\Gamma(U,\mathcal{O}_U)$, siendo $\varphi$ la inyección canónica \sref[I]{I.8.2.7-fr}.
Como el morfismo $h$ satisface las hipótesis de \sref{II.7.3.11-fr}, $\Gamma(U,\mathcal{O}_U)$ es un subanillo de la clausura integral de $A$ en $K$ y, por tanto, es idéntico a $A$ por hipótesis.
\end{proof}

\begin{remark}[8.11.3]
\label{II.8.11.3-fr}
Se verá en el Capítulo~III \sref[III]{III.4.4.11-fr} que, cuando $V$ es un preesquema \emph{localmente noetheriano}, todo morfismo $h:U\to V$ que sea propio y cuasifinito (en particular, todo morfismo que satisfaga las hipótesis de \sref{II.8.11.2-fr}) es necesariamente \emph{finito}.
La conclusión de \sref{II.8.11.2-fr} resulta entonces, en este caso, de \sref{II.6.1.15-fr}.
\end{remark}

\begin{env}[8.11.4]
\label{II.8.11.4-fr}
Vamos a ver ahora que, en el criterio de Grauert \sref{II.8.9.1-fr}, se puede afirmar a menudo que el preesquema $C$ y la «contracción» $q$ están determinados de manera \emph{esencialmente única}.
\end{env}

\begin{lemma}[8.11.5]
\label{II.8.11.5-fr}
Sean $Y$ un preesquema, $p:X\to Y$ un morfismo propio, $\mathcal{L}$ un $\mathcal{O}_X$-módulo invertible $p$-amplio, $C$ un $Y$-preesquema, $\varepsilon:Y\to C$ una $Y$-sección y $q:V=\mathbf{V}(\mathcal{L})\to C$ un $Y$-morfismo, tales que el diagrama de \sref{II.8.9.1.1-fr} sea conmutativo.
Se supone además que, si $p=(p_0,\theta)$, entonces $\theta:\mathcal{O}_Y\to p_*(\mathcal{O}_X)$ es un isomorfismo.
Sean $\mathcal{S}'=\bigoplus_{n\geqslant0}p_*(\mathcal{L}^{\otimes n})$ y $C'=\operatorname{Spec}(\mathcal{S}')$, y sea $q':\mathbf{V}(\mathcal{L})\to C'$ el $Y$-morfismo canónico \sref{II.8.8.5-fr}.
Entonces existe un único $Y$-morfismo $u:C'\to C$ tal que $q=u\circ q'$.
\end{lemma}

\begin{proof}
La hipótesis sobre $\theta$ implica, en particular, que $p$ es sobreyectivo;
como, en virtud de \sref{II.8.8.4-fr}, la restricción de $q'$ a $\mathbf{V}(\mathcal{L})\setmin j(X)$ es un \emph{isomorfismo} sobre $C'\setmin\varepsilon'(Y)$ (donde $\varepsilon'$ es la sección vértice de $C'$), se sigue de \sref{II.8.8.4-fr} que $q'$ es \emph{propio} y \emph{sobreyectivo};
además, en virtud de \sref{II.8.8.6-fr}, si se pone $q'=(q'_0,\tau)$, entonces $\tau:\mathcal{O}_{C'}\to q'_*(\mathcal{O}_V)$ es un isomorfismo.
Se está, pues, en condiciones de aplicar \sref{II.8.11.1-fr}, y el lema quedará demostrado si se prueba que $q$ es constante en toda fibra $q'{}^{-1}(z')$, donde $z'\in C'$.
Ahora bien, esta condición se satisface trivialmente para $z'\notin\varepsilon'(Y)$.
Si $z'\in\varepsilon'(Y)$, existe un único $y\in Y$ tal que $z'=\varepsilon'(y)$ y, en virtud de la conmutatividad de \sref{II.8.8.5.2-fr} y del hecho de que $q'$ aplica $\mathbf{V}(\mathcal{L})\setmin j(X)$ en $C'\setmin\varepsilon'(Y)$, se tiene $q'{}^{-1}(z')=j(p^{-1}(y))$;
la conmutatividad del diagrama de \sref{II.8.9.1.1-fr} demuestra entonces nuestra afirmación.
\end{proof}

\begin{corollary}[8.11.6]
\label{II.8.11.6-fr}
Bajo las hipótesis de \sref{II.8.11.5-fr}, supongamos además que $q$ es propio y que la restricción de $q$ a $\mathbf{V}(\mathcal{L})\setmin j(X)$ es un isomorfismo sobre $C\setmin\varepsilon(Y)$.
Entonces el morfismo $u$ es universalmente cerrado, sobreyectivo y radicial, y su restricción a $C'\setmin\varepsilon'(Y)$ es un isomorfismo sobre $C\setmin\varepsilon(Y)$.
\end{corollary}

\begin{proof}
Como $q'$ es un isomorfismo de $\mathbf{V}(\mathcal{L})\setmin j(X)$ sobre $C'\setmin\varepsilon'(Y)$ \sref{II.8.8.4-fr}, la última afirmación resulta inmediatamente de que $q=u\circ q'$.
Además, la conmutatividad de los diagramas
\oldpage[II]{191}
de \sref{II.8.8.5.2-fr} y \sref{II.8.9.1.1-fr} demuestra que la restricción de $u$ al subpreesquema cerrado $\varepsilon'(Y)$ de $C'$ es un isomorfismo sobre el subpreesquema cerrado $\varepsilon(Y)$ de $C$, de donde se deduce inmediatamente que, para todo $z'\in\varepsilon'(Y)$, si $z=u(z')$, entonces $u$ define un isomorfismo de $\boldsymbol{k}(z)$ sobre $\boldsymbol{k}(z')$.
Estas observaciones prueban que $u$ es biyectivo y radicial;
además, si $\psi$ y $\psi'$ son los morfismos estructurales $C\to Y$, $C'\to Y$, entonces $\psi'=\psi\circ u$ y, como $\psi'$ es separado \sref{II.1.2.4-fr}, también lo es $u$ \sref[I]{I.5.5.1-fr}[v].
Ya se ha visto, en la demostración de \sref{II.8.11.5-fr}, que $q'$ es sobreyectivo;
como $q=u\circ q'$ es propio, se concluye finalmente, por \sref{II.5.4.3-fr} y \sref{II.5.4.9-fr}, que $u$ es universalmente cerrado.
\end{proof}

\begin{proposition}[8.11.7]
\label{II.8.11.7-fr}
Sean $Y$ un preesquema, $X$ un preesquema \emph{íntegro}, $p:X\to Y$ un morfismo propio, $\mathcal{L}$ un $\mathcal{O}_X$-módulo invertible $p$-amplio, $C$ un $Y$-preesquema \emph{normal}, $\varepsilon:Y\to C$ una $Y$-sección y $q:V=\mathbf{V}(\mathcal{L})\to C$ un $Y$-morfismo, tales que el diagrama de \sref{II.8.9.1.1-fr} sea conmutativo.
Se supone que $q$ es propio y que la restricción de $q$ a $V-j(X)$ es un isomorfismo sobre $C-\varepsilon(Y)$.
Se supone además que, si $p=(p_0,\theta)$, entonces $\theta:\mathcal{O}_Y\to p_*(\mathcal{O}_X)$ es un isomorfismo.
Sean $\mathcal{S}'=\bigoplus_{n\geqslant0}p_*(\mathcal{L}^{\otimes n})$ y $C'=\operatorname{Spec}(\mathcal{S}')$, \textup{y} sea $q':\mathbf{V}(\mathcal{L})\to C'$ el $Y$-morfismo canónico \sref{II.8.8.5-fr}.
Entonces el único $Y$-morfismo $u:C'\to C$ tal que $q=u\circ q'$ es un \emph{isomorfismo}.
\end{proposition}

\begin{proof}
De \sref{II.8.8.6-fr} se sigue que $C'$ es íntegro;
como $u$ es un homeomorfismo de los espacios subyacentes $C'\to C$ (pues $u$ es biyectivo y cerrado, por \sref{II.8.11.6-fr}), $C$ es irreducible, luego íntegro y, puesto que la restricción de $u$ a un abierto no vacío de $C'$ es un isomorfismo sobre un abierto de $C$, $u$ es birracional.
Como se supone que $C$ es normal, basta aplicar \sref{II.8.11.2-fr} para obtener la conclusión.
\end{proof}

\begin{remark}[8.11.8]
\label{II.8.11.8-fr}
\begin{enumerate}
    \item[\rm{(i)}] Observemos que la hipótesis de que $C$ es normal implica que $X$ también es normal.
        En efecto, $C'=\operatorname{Spec}(\mathcal{S}')$ es entonces normal, por ser isomorfo a $C$, e íntegro, por \sref{II.8.8.6-fr};
        se concluye, pues, que $\operatorname{Proj}(\mathcal{S}')$ es \emph{normal}.
        En efecto, la cuestión es local sobre $Y$;
        si $Y$ es afín, con $\mathcal{S}'=\widetilde{S}'$, entonces el anillo $S'=\Gamma(C',\mathcal{O}_{C'})$ es íntegro e integralmente cerrado \sref{II.8.8.6.1-fr} y, por tanto, para todo elemento homogéneo $f\in S'_+$, el anillo graduado $S'_f$ es íntegro e integralmente cerrado \cite[t.~I, p.~257 y 261]{I-13}, y también lo es el anillo $S'_{(f)}$ de sus términos de grado cero, porque la intersección de $S'_f$ con el cuerpo de fracciones de $S'_{(f)}$ es igual a $S'_{(f)}$;
        esto prueba nuestra afirmación \sref{II.6.3.4-fr}.
        Por último, como $X$ es isomorfo a un subpreesquema abierto de $\operatorname{Proj}(\mathcal{S}')$ \sref{II.8.8.1-fr}, $X$ es, en efecto, normal.
        Se puede, por tanto, expresar \sref{II.8.11.7-fr} en la forma siguiente:
        \emph{Si $X$ es íntegro y normal, y $p=(p_0,\theta):X\to Y$ es un morfismo propio tal que $\theta:\mathcal{O}_Y\to p_*(\mathcal{O}_X)$ es un isomorfismo, entonces, para todo $\mathcal{O}_X$-módulo $p$-amplio $\mathcal{L}$, existe una manera y solo una de contraer la sección nula de $V=\mathbf{V}(\mathcal{L})$ de modo que se obtenga un $Y$-esquema normal $C$ y un $Y$-morfismo propio $q:V\to C$.}
    \item[\rm{(ii)}] Cuando $p$ es propio, la hipótesis $p_*(\mathcal{O}_X)=\mathcal{O}_Y$ puede considerarse una hipótesis auxiliar, que no restringe verdaderamente la generalidad del resultado.
        En efecto, si no se satisface, basta sustituir $Y$ por el $Y$-esquema $Y'=\operatorname{Spec}(p_*(\mathcal{O}_X))$ y considerar $X$ como un $Y'$-esquema.
        Volveremos sobre este método general en el capítulo~III, \textsection~4.
\end{enumerate}
\end{remark}


\subsection{Haces cuasicoherentes sobre los conos proyectantes}
\label{subsection:II.8.12-fr}

\begin{env}[8.12.1]
\label{II.8.12.1-fr}
Retomemos las hipótesis y notaciones de \sref{II.8.3.1-fr}.
Sea $\mathcal{M}$ un \emph{$\mathcal{S}$-módulo graduado cuasicoherente}; para evitar toda confusión, denotaremos por $\widetilde{\mathcal{M}}$ el $\mathcal{O}_C$-módulo cuasicoherente
\oldpage[II]{192}
asociado a $\mathcal{M}$ \sref{II.1.4.3-fr} cuando $\mathcal{M}$ se considera como un $\mathcal{S}$-módulo \emph{no graduado}, y por $\mathcal{P}\mathit{roj}_0(\mathcal{M})$ el $\mathcal{O}_X$-módulo cuasicoherente asociado a $\mathcal{M}$, considerado $\mathcal{M}$ esta vez como un $\mathcal{S}$-módulo graduado (en otras palabras, el $\mathcal{O}_X$-módulo denotado por $\widetilde{\mathcal{M}}$ en \sref{II.3.2.2-fr}).
Además, ponemos
\begin{equation}
\label{II.8.12.1.1-fr}
  \mathcal{M}_X=\mathcal{P}\mathit{roj}(\mathcal{M})
  =\bigoplus_{n\in\mathbf{Z}}\mathcal{P}\mathit{roj}_0(\mathcal{M}(n))~;
\tag{8.12.1.1}
\end{equation}
estando definida la $\mathcal{O}_X$-álgebra graduada cuasicoherente $\mathcal{S}_X$ por \sref{II.8.6.1.1-fr}, $\mathcal{P}\mathit{roj}(\mathcal{M})$ está dotado de una estructura de \emph{$\mathcal{S}_X$-módulo graduado (cuasicoherente)}, mediante los homomorfismos canónicos \sref{II.3.2.6.1-fr}
\begin{equation}
\label{II.8.12.1.2-fr}
  \mathcal{O}_X(m)\otimes_{\mathcal{O}_X}
  \mathcal{P}\mathit{roj}_0(\mathcal{M}(n))
  \to \mathcal{P}\mathit{roj}_0(\mathcal{S}(m)\otimes_{\mathcal{S}}\mathcal{M}(n))
  \to \mathcal{P}\mathit{roj}_0(\mathcal{M}(m+n))
\tag{8.12.1.2}
\end{equation}
efectuándose la verificación de los axiomas de los módulos con ayuda del diagrama conmutativo de \sref{II.2.5.11.4-fr}.

Si $Y=\operatorname{Spec}(A)$ es afín, $\mathcal{S}=\widetilde{S}$ y $\mathcal{M}=\widetilde{M}$, donde $S$ es un $A$-álgebra graduada y $M$ un $S$-módulo graduado, entonces, para todo elemento homogéneo $f\in S_+$, se tiene
\begin{equation}
\label{II.8.12.1.3-fr}
  \Gamma(X_f,\mathcal{P}\mathit{roj}(\widetilde{M}))=M_f
\tag{8.12.1.3}
\end{equation}
por las definiciones y \sref{II.8.2.9.1-fr}.

Consideremos ahora el $\widehat{\mathcal{S}}$-módulo graduado cuasicoherente
\begin{equation}
\label{II.8.12.1.4-fr}
  \widehat{\mathcal{M}}
  =\mathcal{M}\otimes_{\mathcal{S}}\widehat{\mathcal{S}}
\tag{8.12.1.4}
\end{equation}
($\widehat{\mathcal{S}}$ está definido por \sref{II.8.3.1.1-fr}); de él se deduce un $\mathcal{O}_{\widehat{C}}$-módulo graduado cuasicoherente $\mathcal{P}\mathit{roj}_0(\widehat{\mathcal{M}})$, que denotaremos también por
\begin{equation}
\label{II.8.12.1.5-fr}
  \mathcal{M}^{\square}=\mathcal{P}\mathit{roj}_0(\widehat{\mathcal{M}}).
\tag{8.12.1.5}
\end{equation}

Es claro \sref{II.3.2.4-fr} que $\mathcal{M}^{\square}$ es un funtor aditivo \emph{exacto} en $\mathcal{M}$, que conmuta con las sumas directas y con los límites inductivos.
\end{env}

\begin{proposition}[8.12.2]
\label{II.8.12.2-fr}
Con las notaciones de \sref{II.8.3.2-fr}, se tienen isomorfismos canónicos funtoriales
\begin{equation}
\label{II.8.12.2.1-fr}
  i^*(\mathcal{M}^{\square})\xrightarrow{\sim}\widetilde{\mathcal{M}},
  \qquad
  j^*(\mathcal{M}^{\square})\xrightarrow{\sim}
  \mathcal{P}\mathit{roj}_0(\mathcal{M}).
\tag{8.12.2.1}
\end{equation}
\end{proposition}

\begin{proof}
En efecto, $i^*(\mathcal{M}^{\square})$ se identifica canónicamente con $\left(\widehat{\mathcal{M}}/(\mathbf{z}-1)\widehat{\mathcal{M}}\right)^{\sim}$ sobre $\operatorname{Spec}(\widehat{\mathcal{S}}/(\mathbf{z}-1)\widehat{\mathcal{S}})$ en virtud de \sref{II.3.2.3-fr};
el primero de los isomorfismos canónicos \sref{II.8.12.2.1-fr} se deduce entonces inmediatamente \sref{II.1.4.1-fr} del isomorfismo canónico $\widehat{\mathcal{M}}/(\mathbf{z}-1)\widehat{\mathcal{M}}
\xrightarrow{\sim}\mathcal{M}$.
Por otra parte, la inmersión canónica $j:X\to\widehat{C}$ corresponde al homomorfismo canónico $\widehat{\mathcal{S}}\to\mathcal{S}$ de núcleo $\mathbf{z}\widehat{\mathcal{S}}$ \sref{II.8.3.2-fr};
el segundo homomorfismo de \sref{II.8.12.2.1-fr} es el caso particular del homomorfismo canónico de \sref{II.3.5.2-fr}[ii], pues aquí se tiene $\widehat{\mathcal{M}}\otimes_{\widehat{\mathcal{S}}}\mathcal{S}
=\mathcal{M}$;
para comprobar que es un isomorfismo, se puede limitar al caso en que $Y=\operatorname{Spec}(A)$ es afín, $\mathcal{S}=\widetilde{S}$ y $\mathcal{M}=\widetilde{M}$;
remitiéndonos a \sref{II.2.8.8-fr}, la comprobación de que, para todo $f$ homogéneo en $S_+$, el homomorfismo precedente, restringido a $X_f$, se reduce a un isomorfismo es entonces inmediata.
\end{proof}
\oldpage[II]{193}

Por abuso de lenguaje, se dirá también, gracias a la existencia del primero de los isomorfismos \sref{II.8.12.2.1-fr}, que $\mathcal{M}^{\square}$ es la \emph{clausura proyectiva} del $\mathcal{O}_C$-módulo $\widetilde{\mathcal{M}}$ (quedando sobrentendido que entre los datos del $\mathcal{O}_C$-módulo $\widetilde{\mathcal{M}}$ figura la graduación del $\mathcal{S}$-módulo $\mathcal{M}$).

\begin{env}[8.12.3]
\label{II.8.12.3-fr}
Con las notaciones de \sref{II.8.3.5-fr}, se tiene un homomorfismo canónico funtorial
\begin{equation}
\label{II.8.12.3.1-fr}
  p^*(\mathcal{P}\mathit{roj}_0(\mathcal{M}))
  \to\mathcal{M}^{\square}|\widehat{E}.
\tag{8.12.3.1}
\end{equation}

En efecto, se trata de un caso particular del homomorfismo $\upsilon^{\sharp}$ definido de manera general en \sref{II.3.5.6-fr}.
Si $Y=\operatorname{Spec}(A)$ es afín, $\mathcal{S}=\widetilde{S}$ y $\mathcal{M}=\widetilde{M}$, entonces, remitiéndose a \sref{II.2.8.8-fr}, se ve que la restricción de \sref{II.8.12.3.1-fr} a $p^{-1}(X_f)=\widehat{C}_f$ (para un $f$ homogéneo en $S_+$) corresponde al homomorfismo canónico
\begin{equation}
\label{II.8.12.3.2-fr}
  M_{(f)}\otimes_{S_{(f)}}S_f^{\leq}\to M_f^{\leq}
\tag{8.12.3.2}
\end{equation}
teniendo en cuenta \sref{II.8.2.3.2-fr} y \sref{II.8.2.5.2-fr}.
\end{env}

\begin{env}[8.12.4]
\label{II.8.12.4-fr}
Situémonos en las condiciones de \sref{II.8.5.1-fr}, cuyas hipótesis y notaciones conservamos.
De \sref{II.1.5.6-fr} se sigue que, para todo $\mathcal{S}$-módulo graduado cuasicoherente $\mathcal{M}$, se tiene, por una parte, un isomorfismo canónico
\begin{equation}
\label{II.8.12.4.1-fr}
  \Phi^*(\widetilde{\mathcal{M}})
  \xrightarrow{\sim}
  \left(q^*(\mathcal{M})\otimes_{q^*(\mathcal{S})}\mathcal{S}'\right)^{\sim}
\tag{8.12.4.1}
\end{equation}
de $\mathcal{O}_{C'}$-módulos;
por otra parte, \sref{II.3.5.6-fr} implica la existencia de un $\operatorname{Proj}(\varphi)$-morfismo canónico
\begin{equation}
\label{II.8.12.4.2-fr}
  \mathcal{P}\mathit{roj}_0\mathcal{M}\to
  \left(\mathcal{P}\mathit{roj}_0
  (q^*(\mathcal{M})\otimes_{q^*(\mathcal{S})}\mathcal{S}')\right)
  |G(\varphi)
\tag{8.12.4.2}
\end{equation}
y también de un $\widehat{\Phi}$-morfismo canónico
\begin{equation}
\label{II.8.12.4.3-fr}
  \mathcal{P}\mathit{roj}_0\widehat{\mathcal{M}}\to
  \left(\mathcal{P}\mathit{roj}_0
  (q^*(\widehat{\mathcal{M}})
  \otimes_{q^*(\widehat{\mathcal{S}})}\widehat{\mathcal{S}}')\right)
  |G(\widehat{\varphi}).
\tag{8.12.4.3}
\end{equation}
\end{env}

\begin{env}[8.12.5]
\label{II.8.12.5-fr}
Consideremos ahora la situación de \sref{II.8.6.1-fr}, con las mismas notaciones;
se toma, pues, $Y'=X$, siendo el morfismo $q:X\to Y$ el morfismo estructural y $\varphi$ el $q$-morfismo canónico de \sref{II.8.6.1.2-fr}.
Se tiene entonces un isomorfismo canónico
\begin{equation}
\label{II.8.12.5.1-fr}
  q^*(\mathcal{M})\otimes_{q^*(\mathcal{S})}\mathcal{S}_X^{\geq}
  \xrightarrow{\sim}\mathcal{M}_X^{\geq}
\tag{8.12.5.1}
\end{equation}
poniendo $\mathcal{M}_X^{\geq}=\bigoplus_{n\geq0}\mathcal{P}\mathit{roj}_0(\mathcal{M}(n))$.
En efecto, se puede limitar al caso en que $Y=\operatorname{Spec}(A)$ es afín, $\mathcal{S}=\widetilde{S}$ y $\mathcal{M}=\widetilde{M}$, y definir el isomorfismo \sref{II.8.12.5.1-fr} en cada uno de los abiertos afines $X_f$ (donde $f$ es homogéneo en $S_+$), comprobando la compatibilidad al pasar a un múltiplo homogéneo de $f$.
Ahora bien, la restricción a $X_f$ del primer miembro de \sref{II.8.12.5.1-fr} es $\widetilde{M}'=
\left((M\otimes_A S_{(f)})
\otimes_{S\otimes_A S_{(f)}}S_f^{\geq}\right)^{\sim}$ por \sref{II.8.6.2.1-fr};
como se tiene un isomorfismo canónico de $M\otimes_A S_{(f)}$ sobre $M\otimes_S(S\otimes_A S_{(f)})$, se deduce un isomorfismo de $\widetilde{M}'$ sobre $\left(M\otimes_S S_f^{\geq}\right)^{\sim}$, y este último es canónicamente isomorfo, por \sref{II.8.2.9.1-fr}, a la restricción a $X_f$ del segundo miembro de \sref{II.8.12.5.1-fr}, y satisface las condiciones de compatibilidad requeridas.
\oldpage[II]{194}

Sustituyendo $\mathcal{M}$ por $\widehat{\mathcal{M}}$, $\mathcal{S}$ por $\widehat{\mathcal{S}}$ y $\mathcal{S}_X$ por $(\mathcal{S}_X^{\geq})^{\widehat{\,}}$ en el razonamiento precedente, se tiene asimismo un isomorfismo canónico
\begin{equation}
\label{II.8.12.5.2-fr}
  q^*(\widehat{\mathcal{M}})
  \otimes_{q^*(\widehat{\mathcal{S}})}(\mathcal{S}_X^{\geq})^{\widehat{\,}}
  \xrightarrow{\sim}(\mathcal{M}_X^{\geq})^{\widehat{\,}}.
\tag{8.12.5.2}
\end{equation}

Si se recuerda \sref{II.8.6.2-fr} que el morfismo estructural $u:\operatorname{Proj}(\mathcal{S}_X^{\geq})\to X$ es un isomorfismo, se deduce en primer lugar de lo anterior que se tiene un $u$-isomorfismo canónico
\begin{equation}
\label{II.8.12.5.3-fr}
  \mathcal{P}\mathit{roj}_0\mathcal{M}
  \xrightarrow{\sim}\mathcal{P}\mathit{roj}_0(\mathcal{M}_X^{\geq})
\tag{8.12.5.3}
\end{equation}
como caso particular de \sref{II.8.12.4.2-fr}.
Se observa que, con las notaciones de la demostración de \sref{II.8.6.2-fr}, esto se reduce a ver que el homomorfismo canónico $M_{(f)}\otimes_{S_{(f)}}(S_f^{\geq})^{(d)}\to(M_f^{\geq})^{(d)}$ es un isomorfismo cuando $f\in S_d$, lo cual es inmediato.

En segundo lugar, el isomorfismo \sref{II.8.12.5.2-fr} proporciona, aplicando esta vez \sref{II.8.12.4.3-fr} al morfismo canónico $r=\operatorname{Proj}(\widehat{\alpha}):\widehat{C}_X\to
\widehat{C}$, un $r$-morfismo canónico
\begin{equation}
\label{II.8.12.5.4-fr}
  \mathcal{M}^{\square}\to(\mathcal{M}_X^{\geq})^{\square}.
\tag{8.12.5.4}
\end{equation}

Recordemos ahora \sref{II.8.6.2-fr} que las restricciones de $r$ a los conos sin vértice $\widehat{E}_X$ y $E_X$ son \emph{isomorfismos} sobre $\widehat{E}$ y $E$, respectivamente.
Además:
\end{env}

\begin{proposition}[8.12.6]
\label{II.8.12.6-fr}
Las restricciones a $\widehat{E}_X$ y $E_X$ del $r$-morfismo canónico de \sref{II.8.12.5.4-fr} son isomorfismos
\begin{equation}
\label{II.8.12.6.1-fr}
  \mathcal{M}^{\square}|\widehat{E}
  \xrightarrow{\sim}(\mathcal{M}_X^{\geq})^{\square}|\widehat{E}_X
\tag{8.12.6.1}
\end{equation}
\begin{equation}
\label{II.8.12.6.2-fr}
  \mathcal{M}^{\sim}|E
  \xrightarrow{\sim}(\mathcal{M}_X^{\geq})^{\sim}|E_X.
\tag{8.12.6.2}
\end{equation}
\end{proposition}

\begin{proof}
Se reduce al caso en que $Y$ es afín, como en la demostración de \sref{II.8.6.2-fr} ; con las notaciones de esta última, y reduciéndose a las definiciones \sref{II.2.8.8-fr}, se debe demostrar que el homomorfismo canónico
\[
  \widehat{M}_{(f)}\otimes_{\widehat{S}_{(f)}}
  (S_f^{\geq})_{(f/1)}^{\widehat{\,}}
  \to(M\otimes_S S_f^{\geq})_{(f/1)}^{\widehat{\,}}
\]
es un isomorfismo;
pero, en virtud de \sref{II.8.2.3.2-fr} y \sref{II.8.2.5.2-fr}, el primer miembro se identifica canónicamente con $M_f^{\leq}\otimes_{S_f^{\leq}}(S_f^{\geq})_{f/1}^{\leq}$ y, por tanto, con $M_f^{\leq}$, por \sref{II.8.2.7.2-fr}, mientras que el segundo se identifica con $(M_f^{\geq})_{f/1}^{\leq}$ y, por tanto, también con $M_f^{\leq}$, por \sref{II.8.2.9.2-fr}, de donde se obtiene la conclusión relativa a \sref{II.8.12.6.1-fr};
\sref{II.8.12.6.2-fr} se sigue entonces de \sref{II.8.12.6.1-fr} y \sref{II.8.12.2.1-fr}.
\end{proof}

\begin{corollary}[8.12.7]
\label{II.8.12.7-fr}
Con las identificaciones de \sref{II.8.6.3-fr}, la restricción de $(\mathcal{M}_X^{\geq})^{\square}$ a $\widehat{E}_X$ puede identificarse con $(\mathcal{M}_X^{\leq})^{\sim}$, y la restricción de $(\mathcal{M}_X^{\geq})^{\square}$ a $E_X$ con $\widetilde{\mathcal{M}}_X$.
\end{corollary}

\begin{proof}
Se puede limitar al caso afín, y esto se sigue de la identificación de $(M_f^{\geq})_{f/1}^{\leq}$ con $M_f^{\leq}$, y de $(M_f^{\geq})_{f/1}$ con $M_f$ \sref{II.8.2.9.2-fr}.
\end{proof}

\begin{proposition}[8.12.8]
\label{II.8.12.8-fr}
Bajo las hipótesis de \sref{II.8.6.4-fr}, el homomorfismo canónico de \sref{II.8.12.3.1-fr} es un isomorfismo.
\end{proposition}

\begin{proof}
Teniendo en cuenta que $\operatorname{Proj}(\mathcal{S}_X^{\geq})\to X$ es un isomorfismo \sref{II.8.6.2-fr}, y los
\oldpage[II]{195}
isomorfismos \sref{II.8.12.5.4-fr} y \sref{II.8.12.6.1-fr}, se reduce a demostrar la proposición correspondiente para el homomorfismo canónico $p_X^*(\mathcal{P}\mathit{roj}_0(\mathcal{M}_X^{\geq}))
\to(\mathcal{M}_X^{\geq})^{\square}|\widehat{E}_X$; en otras palabras, se puede limitar al caso en que $\mathcal{S}_1$ es un $\mathcal{O}_Y$-módulo invertible y $\mathcal{S}$ está generada por $\mathcal{S}_1$.
Con las notaciones de \sref{II.8.12.3-fr}, se tiene entonces, para algún $f\in S_1$, que $S_f^{\leq}=S_{(f)}[1/f]$, y el homomorfismo canónico $M_{(f)}\otimes_{S_{(f)}}S_f^{\leq}\to M_f^{\leq}$ es un isomorfismo, por la definición de $M_f^{\leq}$.
\end{proof}

\begin{env}[8.12.9]
\label{II.8.12.9-fr}
Consideremos ahora los $\mathcal{S}$-módulos cuasicoherentes
\[
  \mathcal{M}_{[n]}=\bigoplus_{m\geqslant n}\mathcal{M}_m
\]
y (con las notaciones de \sref{II.8.7.2-fr}) el $\mathcal{S}^{\natural}$-módulo graduado cuasicoherente
\begin{equation}
\label{II.8.12.9.1-fr}
  \mathcal{M}^{\natural}
  =\left(\bigoplus_{n\geqslant0}\mathcal{M}_{[n]}\right)^{\sim}.
\tag{8.12.9.1}
\end{equation}

Se ha visto en \sref{II.8.7.3-fr} que existe un $C$-isomorfismo canónico $h:C_X\xrightarrow{\sim}\operatorname{Proj}(\mathcal{S}^{\natural})$.
Además:
\end{env}

\begin{proposition}[8.12.10]
\label{II.8.12.10-fr}
Existe un $h$-isomorfismo canónico
\begin{equation}
\label{II.8.12.10.1-fr}
  \mathcal{P}\mathit{roj}_0(\mathcal{M}^{\natural})
  \xrightarrow{\sim}\widetilde{\mathcal{M}}_X.
\tag{8.12.10.1}
\end{equation}
\end{proposition}

\begin{proof}
Se razona como en \sref{II.8.7.3-fr}, utilizando esta vez la existencia del diisomorfismo de \sref{II.8.2.9.3-fr} en lugar del de \sref{II.8.2.7.3-fr}.
Dejamos los detalles al lector.
\end{proof}


\subsection{Clausuras proyectivas de subhaces y subesquemas cerrados}
\label{subsection:II.8.13-fr}

\begin{env}[8.13.1]
\label{II.8.13.1-fr}
Con las hipótesis y notaciones de \sref{II.8.12.1-fr}, consideremos un sub-$\mathcal{S}$-módulo cuasicoherente $\mathcal{N}$ de $\mathcal{M}$, \emph{no necesariamente graduado}.
Se puede entonces considerar el $\mathcal{O}_C$-módulo cuasicoherente $\widetilde{\mathcal{N}}$ asociado a $\mathcal{N}$, que es un sub-$\mathcal{O}_C$-módulo de $\widetilde{\mathcal{M}}$.
Se ha visto, por otra parte, en \sref{II.8.12.2.1-fr}, que $\widetilde{\mathcal{M}}$ puede identificarse con la restricción de $\mathcal{M}^{\square}$ a $C$.
Como la inyección canónica $i:C\to\widehat{C}$ es un morfismo afín \sref{II.8.3.2-fr} y, \emph{a fortiori}, cuasicompacto, la \emph{prolongación canónica} $(\widetilde{\mathcal{N}})^{-}$, el mayor sub-$\mathcal{O}_{\widehat{C}}$-módulo contenido en $\mathcal{M}^{\square}$ que induce $\widetilde{\mathcal{N}}$ sobre $C$, es un $\mathcal{O}_{\widehat{C}}$-módulo \emph{cuasicoherente} \sref[I]{I.9.4.2-fr}.
Daremos una descripción más explícita por medio de un $\widehat{\mathcal{S}}$-módulo graduado.
\end{env}

\begin{env}[8.13.2]
\label{II.8.13.2-fr}
Para ello, consideremos, para todo entero $n\geqslant0$, el homomorfismo $\bigoplus_{i\leqslant n}\mathcal{M}_i\to\mathcal{M}$ que, para todo abierto $U$ de $Y$, envía la familia
\[
  (s_i)\in\bigoplus_{i\leqslant n}\Gamma(U,\mathcal{M}_i)
\]
a la sección $\sum_i s_i\in\Gamma(U,\mathcal{M})$.
Denotemos por $\mathcal{N}'_n$ la imagen inversa de $\mathcal{N}$ por este homomorfismo, que es un sub-$\mathcal{O}_Y$-módulo cuasicoherente de $\bigoplus_{i\leqslant n}\mathcal{M}_i$.
Consideremos ahora el homomorfismo $\bigoplus_{i\leqslant n}\mathcal{M}_i\to
\widehat{\mathcal{M}}=\mathcal{M}[\mathbf{z}]$ que envía $(s_i)$ a la sección $\sum_{i\leqslant n}s_i \mathbf{z}^{n-i}\in\Gamma(U,\widehat{\mathcal{M}}_n)$, y sea $\mathcal{N}_n$ la imagen de $\mathcal{N}'_n$ por este homomorfismo;
se tiene inmediatamente que $\overline{\mathcal{N}}=\bigoplus_{n\geqslant0}\mathcal{N}_n$ es un sub-$\widehat{\mathcal{S}}$-módulo graduado (cuasicoherente) de $\widehat{\mathcal{M}}$;
se dice que $\overline{\mathcal{N}}$ se obtiene de $\mathcal{N}$ por \emph{homogeneización}, mediante la «variable de homogeneización» $\mathbf{z}$.
Se observará
\oldpage[II]{196}
que, si $\mathcal{N}$ ya es un sub-$\mathcal{S}$-módulo \emph{graduado} de $\mathcal{M}$, entonces $\overline{\mathcal{N}}$ se identifica con la suma directa de las componentes $\widehat{\mathcal{N}}_n$ de grado $n\geqslant0$ en $\widehat{\mathcal{N}}=\mathcal{N}[\mathbf{z}]$.
\end{env}

\begin{proposition}[8.13.3]
\label{II.8.13.3-fr}
El $\mathcal{O}_{\widehat{C}}$-módulo $\mathcal{P}\mathit{roj}_0(\overline{\mathcal{N}})$ es la prolongación canónica $(\widetilde{\mathcal{N}})^{-}$ de $\widetilde{\mathcal{N}}$ a $\widehat{C}$.
\end{proposition}

\begin{proof}
La cuestión es local sobre $Y$ y $\widehat{C}$, en virtud de la definición de la prolongación canónica \sref[I]{I.9.4.1-fr}.
Se puede, pues, suponer ya que $Y=\operatorname{Spec}(A)$ es afín, con $\mathcal{S}=\widetilde{S}$, $\mathcal{M}=\widetilde{M}$ y $\mathcal{N}=\widetilde{N}$, donde $N$ es un sub-$S$-módulo de $M$ no necesariamente graduado.
Además, por \sref{II.8.3.2.6-fr}, $\widehat{C}$ es unión de los abiertos afines $\widehat{C}_{\mathbf{z}}=C$ y $\widehat{C}_f=\operatorname{Spec}(S_f^{\leq})$ (con $f$ homogéneo en $S_+$).
Basta, por tanto, demostrar que: $1^\circ$ la restricción de $\mathcal{P}\mathit{roj}_0(\overline{\mathcal{N}})$ a $C$ es $\widetilde{\mathcal{N}}$; $2^\circ$ la restricción de $\mathcal{P}\mathit{roj}_0(\overline{\mathcal{N}})$ a cada $\widehat{C}_f$ es la prolongación canónica de la restricción de $\widetilde{\mathcal{N}}$ a $C\cap\widehat{C}_f=\operatorname{Spec}(S_f)$ \sref{II.8.3.2.6-fr}.
Para el primer punto, observemos que $\mathcal{P}\mathit{roj}_0(\overline{\mathcal{N}})|C$ se identifica con $(\overline{N}_{(\mathbf{z})})^{\sim}$ \sref{II.8.3.2.4-fr};
pero $\overline{N}_{(\mathbf{z})}$ se identifica canónicamente \sref{II.2.2.5-fr} con la imagen de $\overline{N}$ en $\widehat{M}/(\mathbf{z}-1)\widehat{M}$ y, mediante el isomorfismo canónico de este último con $M$ \sref{II.8.2.5-fr}, dicha imagen se identifica con $N$, por la definición de $\overline{N}$ dada en \sref{II.8.13.2-fr}.

Para demostrar el segundo punto, observemos que la inyección $i:C\cap\widehat{C}_f\to\widehat{C}_f$ corresponde a la inyección canónica $S_f^{\leq}\to S_f$ \sref{II.8.3.2.6-fr};
se tiene asimismo $\Gamma(\widehat{C}_f,\mathcal{M}^{\square})=M_f^{\leq}$, $\Gamma(\widehat{C}_f,i_*(\widetilde{\mathcal{N}}))=N_f$ y, por \sref{II.8.12.2.1-fr}, $\Gamma(\widehat{C}_f,i_*(i^*(\mathcal{M}^{\square})))=M_f$.
Teniendo en cuenta \sref[I]{I.9.4.2-fr}, se reduce a demostrar que $\overline{N}_{(f)}\subset\widehat{M}_{(f)}=M_f^{\leq}$ se identifica canónicamente con la imagen inversa de $N_f$ por la inyección canónica $M_f^{\leq}\to M_f$.
En efecto, sea $d=\deg(f)>0$, y supongamos que un elemento $(\sum_{k\leqslant md}x_k)/f^m$ de $M_f$ (con $x_k\in M_k$) es de la forma $y/f^m$, con $y\in N$.
Multiplicando $y$ y los $x_k$ por una misma potencia conveniente $f^h$, se puede suponer ya que $\sum_{k\leqslant md}x_k=y$.
Pero, en la identificación de \sref{II.8.2.5.2-fr}, $(\sum_{k\leqslant md}x_k)/f^m$ corresponde a $\sum_{k\leqslant md}x_k\mathbf{z}^{md-k}/f^m$, y este es efectivamente un elemento de $\overline{N}_{(f)}$, puesto que $\sum_{k\leqslant md}x_k\in N$;
el recíproco es evidente.
\end{proof}

\begin{remark}[8.13.4]
\label{II.8.13.4-fr}
\begin{enumerate}
    \item[\rm{(i)}] El caso de aplicación más importante de \sref{II.8.13.3-fr} es aquel en que $\mathcal{M}=\mathcal{S}$, siendo entonces $\widetilde{\mathcal{N}}$ un haz \emph{arbitrario} cuasicoherente de ideales $\mathcal{J}$ de $\mathcal{O}_C$ \sref{II.1.4.3-fr}, que corresponde biyectivamente a un \emph{subpreesquema cerrado} $Z$ de $C$.
        Entonces la prolongación canónica $\overline{\mathcal{J}}$ de $\mathcal{J}$ es el haz cuasicoherente de ideales de $\mathcal{O}_{\widehat{C}}$ que define la \emph{adherencia} $\overline{Z}$ de $Z$ en $\widehat{C}$ \sref[I]{I.9.5.10-fr};
        la Proposición~\sref{II.8.13.3-fr} proporciona un modo canónico de definir $\overline{Z}$ mediante un ideal graduado de $\widehat{\mathcal{S}}=\mathcal{S}[\mathbf{z}]$.
    \item[\rm{(ii)}] Supongamos, para simplificar, que $Y$ es afín y adoptemos la notación de la demostración de \sref{II.8.13.3-fr}.
        Para cada $x\in N$ no nulo, sea $d(x)$ el mayor grado de las componentes homogéneas $x_i$ de $x$ en $M$;
        por definición, $\overline{N}$ es el submódulo de $\widehat{M}$ formado por $0$ y por los elementos de la forma $h(x,k)=\mathbf{z}^k\sum_{i\leqslant d(x)}x_i \mathbf{z}^{d(x)-i}$ ($k$ entero $\geqslant0$);
        por tanto, está generado, como módulo sobre $\widehat{S}=S[\mathbf{z}]$, por los elementos de la forma
        \[
    h(x,0)=\sum_{i\leqslant d(x)}x_i \mathbf{z}^{d(x)-i}.
  \]
\oldpage[II]{197}
        Se dice que $h(x,0)$ se deduce de $x$ por \emph{homogeneización} mediante la «variable homogeneizante» $\mathbf{z}$.
        Pero, como $h(x,0)$ no depende aditivamente de $x$ (ni, \emph{a fortiori}, $S$-linealmente), \emph{se evitará creer} (incluso cuando $M=S$) que los $h(x,0)$ forman un \emph{sistema de generadores} del $\widehat{S}$-módulo graduado $\overline{N}$ cuando $x$ recorre un \emph{sistema de generadores} del $S$-módulo $N$.
        Esto sí sucede, sin embargo, cuando $N$ es un $S$-módulo \emph{libre cíclico} (el único caso considerado en la geometría algebraica elemental), pues, si $t$ es una base de $N$, entonces $h(t,0)$ genera el $\widehat{S}$-módulo $\overline{N}$.
\end{enumerate}
\end{remark}


\subsection{Complementos sobre los haces asociados a $\mathcal{S}$-módulos graduados}
\label{subsection:II.8.14-fr}

\begin{env}[8.14.1]
\label{II.8.14.1-fr}
Sean $Y$ un preesquema, $\mathcal{S}$ una $\mathcal{O}_Y$-álgebra cuasicoherente \emph{graduada en grados positivos}, $X=\operatorname{Proj}(\mathcal{S})$ y $q:X\to Y$ el morfismo estructural (que es separado, por \sref{II.3.1.3-fr}).
Con la notación de \sref{II.8.12.1-fr}, hemos definido un funtor $\mathcal{M}_X=\mathcal{P}\mathit{roj}(\mathcal{M})$ en $\mathcal{M}$, de la categoría de $\mathcal{S}$-módulos graduados cuasicoherentes a la categoría de $\mathcal{S}_X$-módulos graduados cuasicoherentes;
además, está claro \sref{II.3.2.4-fr} que este funtor es \emph{aditivo} y \emph{exacto}, y que conmuta con los límites inductivos.

Observemos, además, que de la definición \sref{II.8.12.1.1-fr} se sigue inmediatamente que
\begin{equation}
\label{II.8.14.1.1-fr}
  \mathcal{P}\mathit{roj}(\mathcal{M}(n))
  =(\mathcal{P}\mathit{roj}(\mathcal{M}))(n)
  \qquad\text{para todo }n\in\mathbf{Z}.
\tag{8.14.1.1}
\end{equation}
\end{env}

\begin{env}[8.14.2]
\label{II.8.14.2-fr}
Vamos a extender primero a los $\mathcal{S}_X$-módulos de la forma $\mathcal{P}\mathit{roj}(\mathcal{M})$ los homomorfismos canónicos $\lambda$ y $\mu$ definidos en (\hyperref[II.3.2.6-fr]{3.2.6}).
Para ello, observemos que, para cualesquiera $m\in\mathbf{Z}$ y $n\in\mathbf{Z}$, se tiene, por \sref{II.2.1.2.1-fr}, un homomorfismo canónico de $\mathcal{O}_X$-módulos
\begin{equation}
\label{II.8.14.2.1-fr}
  \lambda_{mn}:\mathcal{P}\mathit{roj}_0(\mathcal{M}(m))
  \otimes_{\mathcal{O}_X}\mathcal{P}\mathit{roj}_0(\mathcal{N}(n))
  \to\mathcal{P}\mathit{roj}_0
  ((\mathcal{M}\otimes_{\mathcal{S}}\mathcal{N})(m+n))
\tag{8.14.2.1}
\end{equation}
y, asimismo, por \sref{II.2.1.2.2-fr}, un homomorfismo canónico de $\mathcal{O}_X$-módulos
\begin{equation}
\label{II.8.14.2.2-fr}
  \mu_{mn}:\mathcal{P}\mathit{roj}_0
  ((\mathcal{H}om_{\mathcal{S}}(\mathcal{M},\mathcal{N}))(n-m))
  \to\mathcal{H}om_{\mathcal{O}_X}
  (\mathcal{P}\mathit{roj}_0(\mathcal{M}(m)),
   \mathcal{P}\mathit{roj}_0(\mathcal{N}(n)))
\tag{8.14.2.2}
\end{equation}
para cualesquiera $\mathcal{S}$-módulos graduados cuasicoherentes $\mathcal{M}$ y $\mathcal{N}$.
Esto induce un homomorfismo
\[
  \mu_k:\mathcal{P}\mathit{roj}_0
  ((\mathcal{H}om_{\mathcal{S}}(\mathcal{M},\mathcal{N}))(k))
  \to
  (\mathcal{H}om_{\mathcal{S}_X}
  (\mathcal{P}\mathit{roj}(\mathcal{M}),\mathcal{P}\mathit{roj}(\mathcal{N})))_k
\]
que envía cada $u\in\Gamma(U,\mathcal{P}\mathit{roj}_0
((\mathcal{H}om_{\mathcal{S}}(\mathcal{M},\mathcal{N}))(k)))$ al homomorfismo $\mu_k(u)$, de grado $k$, de $\mathbf{Z}$-módulos graduados $\Gamma(U,\mathcal{P}\mathit{roj}(\mathcal{M}))\to
\Gamma(U,\mathcal{P}\mathit{roj}(\mathcal{N}))$ (donde $U$ es abierto en $X$), cuya restricción a cada $\Gamma(U,\mathcal{P}\mathit{roj}_0(\mathcal{M}(m)))$ coincide con $\mu_{m,m+k}(u)$;
además, volviendo a la definición de los $\mu_{mn}$ \sref{II.2.5.12.1-fr}, se ve inmediatamente que $\mu_k(u)$ es de hecho un homomorfismo de grado $k$ de $\Gamma(U,\mathcal{S}_X)$-módulos graduados y que, además, los $\mu_k$ definen un homomorfismo de \emph{$\mathcal{S}_X$-módulos graduados}
\begin{equation}
\label{II.8.14.2.3-fr}
  \mathcal{P}\mathit{roj}
  (\mathcal{H}om_{\mathcal{S}}(\mathcal{M},\mathcal{N}))
  \to\mathcal{H}om_{\mathcal{S}_X}
  (\mathcal{P}\mathit{roj}(\mathcal{M}),\mathcal{P}\mathit{roj}(\mathcal{N})).
\tag{8.14.2.3}
\end{equation}

De modo análogo, teniendo en cuenta el diagrama de asociatividad \sref{II.2.5.11.4-fr}, los homomorfismos \sref{II.8.14.2.1-fr} dan un homomorfismo de \emph{$\mathcal{S}_X$-módulos graduados}
\begin{equation}
\label{II.8.14.2.4-fr}
  \lambda:\mathcal{P}\mathit{roj}(\mathcal{M})
  \otimes_{\mathcal{S}_X}\mathcal{P}\mathit{roj}(\mathcal{N})
  \to\mathcal{P}\mathit{roj}(\mathcal{M}\otimes_{\mathcal{S}}\mathcal{N}).
\tag{8.14.2.4}
\end{equation}
\end{env}
\oldpage[II]{198}

\begin{proposition}[8.14.3]
\label{II.8.14.3-fr}
El homomorfismo \sref{II.8.14.2.4-fr} es biyectivo;
también lo es \sref{II.8.14.2.3-fr} siempre que el $\mathcal{S}$-módulo graduado $\mathcal{M}$ admita una presentación finita \sref{II.3.1.1-fr}.
\end{proposition}

\begin{proof}
La cuestión es claramente local sobre $X$ y $Y$;
se puede, por tanto, suponer que $Y=\operatorname{Spec}(A)$ es afín, con $\mathcal{S}=\widetilde{S}$, $\mathcal{M}=\widetilde{M}$ y $\mathcal{N}=\widetilde{N}$, donde $S$ es una $A$-álgebra graduada en grados positivos y $M$ y $N$ son $S$-módulos graduados.
Si $f$ es un elemento homogéneo de $S_+$, los homomorfismos \sref{II.8.14.2.1-fr} y \sref{II.8.14.2.2-fr}, restringidos al abierto afín $D_+(f)$, corresponden a los homomorfismos canónicos \sref{II.2.5.11.1-fr} y \sref{II.2.5.12.1-fr}:
\begin{align*}
  M(m)_{(f)}\otimes_{S_{(f)}}N(n)_{(f)}
  &\to(M\otimes_S N)(m+n)_{(f)}\\
  (\operatorname{Hom}_S(M,N))(n-m)_{(f)}
  &\to\operatorname{Hom}_{S_{(f)}}(M(m)_{(f)},N(n)_{(f)}).
\end{align*}

Si se remite a las definiciones de estos homomorfismos, se ve entonces (teniendo en cuenta \sref{II.8.2.9.1-fr}) que la restricción de \sref{II.8.14.2.4-fr} a $D_+(f)$ corresponde al homomorfismo canónico
\[
  M_f\otimes_{S_f}N_f\to(M\otimes_S N)_f
\]
definido en \sref[0]{0.1.3.4-fr}, y se sabe que este último homomorfismo es un isomorfismo.
De modo análogo, la restricción de \sref{II.8.14.2.3-fr} a $D_+(f)$ corresponde al homomorfismo canónico \sref[0]{0.1.3.5-fr}
\[
  (\operatorname{Hom}_S(M,N))_f\to
  \operatorname{Hom}_{S_f}(M_f,N_f)
\]
teniendo en cuenta que, como $M$ es de tipo finito, el módulo $\operatorname{Hom}_S(M,N)$, suma directa de los subgrupos formados por los homomorfismos \emph{homogéneos} de $S$-módulos \sref{II.2.1.2-fr}, coincide con el conjunto de \emph{todos} los homomorfismos $M\to N$ de $S$-módulos.
La hipótesis de que $M$ admite una presentación finita implica entonces \sref[0]{0.1.3.5-fr} que el homomorfismo canónico considerado es, en efecto, un isomorfismo.
\end{proof}

\begin{proposition}[8.14.4]
\label{II.8.14.4-fr}
Si $U$ es un abierto cuasicompacto de $X$, existe un entero $d$ tal que, para todo entero $n$ múltiplo de $d$, $\mathcal{O}_X(n)|U$ es invertible y su inverso es $\mathcal{O}_X(-n)|U$.
\end{proposition}

\begin{proof}
Como $q(U)$ es cuasicompacto, está recubierto por un número finito de abiertos afines $V_i$, y así todo $x\in U$ está contenido en algún abierto afín de la forma $D_+(f)$, donde $f$ es un elemento homogéneo de grado $>0$ de uno de los anillos $\Gamma(V_i,\mathcal{S})$.
Como $U$ es cuasicompacto, se puede recubrirlo por un número finito de tales abiertos $D_+(f_j)$;
sea $d$ un múltiplo común de los grados de los $f_j$.
Este $d$ posee la propiedad requerida, por \sref{II.2.5.17-fr}.
\end{proof}

\begin{env}[8.14.5]
\label{II.8.14.5-fr}
Bajo las hipótesis y con la notación de \sref{II.8.14.1-fr}, se han definido, en \sref{II.3.3.2-fr}, homomorfismos canónicos de $\mathcal{O}_Y$-módulos
\begin{equation}
\label{II.8.14.5.1-fr}
  \alpha_n:\mathcal{M}_n\to
  q_*(\mathcal{P}\mathit{roj}_0(\mathcal{M}(n)))
  \qquad(n\in\mathbf{Z})
\tag{8.14.5.1}
\end{equation}

Generalizando la notación de \sref{II.3.3.1-fr}, ponemos, para todo \emph{$\mathcal{S}_X$-módulo graduado $\mathcal{F}$},
\begin{equation}
\label{II.8.14.5.2-fr}
  \bbGamma_*(\mathcal{F})=
  \bigoplus_{n\in\mathbf{Z}}q_*(\mathcal{F}_n).
\tag{8.14.5.2}
\end{equation}

En particular, $\bbGamma_*(\mathcal{S}_X)=\bigoplus_{n\in\mathbf{Z}}q_*(\mathcal{O}_X(n))$ es la $\mathcal{O}_Y$-álgebra graduada denotada por $\bbGamma_*(\mathcal{O}_X)$ en \sref{II.3.3.1.2-fr};
está claro que $\bbGamma_*(\mathcal{F})$ es un \emph{$\bbGamma_*(\mathcal{S}_X)$-módulo graduado} \sref[0]{0.4.2.2-fr}.
Cuando
\oldpage[II]{199}
se toma $\mathcal{M}=\mathcal{S}$ en los homomorfismos \sref{II.8.14.5.1-fr}, se obtiene el homomorfismo de $\mathcal{O}_Y$-álgebras graduadas
\begin{equation}
\label{II.8.14.5.3-fr}
  \alpha:\mathcal{S}\to\bbGamma_*(\mathcal{S}_X)
\tag{8.14.5.3}
\end{equation}
definido anteriormente en \sref{II.3.3.2-fr}, y que hace de $\bbGamma_*(\mathcal{F})$ un \emph{$\mathcal{S}$-módulo graduado};
los homomorfismos \sref{II.8.14.5.1-fr} definen entonces un homomorfismo (de grado $0$) de \emph{$\mathcal{S}$-módulos graduados}
\begin{equation}
\label{II.8.14.5.4-fr}
  \alpha:\mathcal{M}\to\bbGamma_*(\mathcal{P}\mathit{roj}(\mathcal{M})).
\tag{8.14.5.4}
\end{equation}
\end{env}

\begin{env}[8.14.6]
\label{II.8.14.6-fr}
En general, para un $\mathcal{S}_X$-módulo graduado cuasicoherente $\mathcal{F}$, no es seguro que el $\mathcal{S}$-módulo graduado $\bbGamma_*(\mathcal{F})$ sea necesariamente cuasicoherente.
Consideremos un abierto $X'$ de $X$ tal que la restricción $q':X'\to Y$ de $q$ a $X'$ sea un morfismo \emph{cuasicompacto}.
Como $q'$ es además separado, $q'_*(\mathcal{F}')$ es entonces un $\mathcal{O}_Y$-módulo cuasicoherente para todo $\mathcal{O}_{X'}$-módulo cuasicoherente $\mathcal{F}'$ \sref[I]{I.9.2.2-fr}[\emph{b}].
Ponemos
\begin{equation}
\label{II.8.14.6.1-fr}
  \mathcal{S}_{X'}=\mathcal{S}_X|X'
  =\bigoplus_{n\in\mathbf{Z}}\mathcal{O}_X(n)|X'
\tag{8.14.6.1}
\end{equation}
y, para todo $\mathcal{S}_{X'}$-módulo graduado $\mathcal{F}'$,
\begin{equation}
\label{II.8.14.6.2-fr}
  \bbGamma'_*(\mathcal{F}')=
  \bigoplus_{n\in\mathbf{Z}}q'_*(\mathcal{F}'_n).
\tag{8.14.6.2}
\end{equation}

La observación anterior muestra entonces que, si $\mathcal{F}'$ es un $\mathcal{S}_{X'}$-módulo cuasicoherente, $\bbGamma'_*(\mathcal{F}')$ es un $\mathcal{S}$-módulo graduado \emph{cuasicoherente} \sref[I]{I.9.6.1-fr}.

Se observa también que la inyección canónica $j:X'\to X$ es \emph{cuasicompacta}, porque $q'=q\circ j$ es cuasicompacto y $q$ es separado \sref[I]{I.6.6.4-fr}[v].
Entonces $\mathcal{F}=j_*(\mathcal{F}')$ es un $\mathcal{S}_X$-módulo graduado cuasicoherente para todo $\mathcal{S}_{X'}$-módulo graduado cuasicoherente $\mathcal{F}'$, y de las definiciones anteriores se sigue que
\begin{equation}
\label{II.8.14.6.3-fr}
  \bbGamma'_*(\mathcal{F}')=\bbGamma_*(\mathcal{F}).
\tag{8.14.6.3}
\end{equation}

Bajo las mismas hipótesis sobre $X'$, para todo $\mathcal{S}$-módulo graduado cuasicoherente $\mathcal{M}$, ponemos
\begin{equation}
\label{II.8.14.6.4-fr}
  \mathcal{P}\mathit{roj}'(\mathcal{M})
  =\mathcal{P}\mathit{roj}(\mathcal{M})|X'
\tag{8.14.6.4}
\end{equation}
que es un $\mathcal{S}_{X'}$-módulo graduado cuasicoherente.
El homomorfismo canónico
\[
  \mathcal{P}\mathit{roj}(\mathcal{M})\to
  j_*(\mathcal{P}\mathit{roj}'(\mathcal{M}))
\]
\sref[0]{0.4.4.3-fr} proporciona, pues, un homomorfismo canónico $\bbGamma_*(\mathcal{P}\mathit{roj}(\mathcal{M}))\to
\bbGamma'_*(\mathcal{P}\mathit{roj}'(\mathcal{M}))$ de $\mathcal{S}$-módulos graduados y, componiéndolo con \sref{II.8.14.5.4-fr}, se obtiene un homomorfismo canónico funtorial (de grado $0$) de $\mathcal{S}$-módulos graduados cuasicoherentes
\begin{equation}
\label{II.8.14.6.5-fr}
  \alpha':\mathcal{M}\to
  \bbGamma'_*(\mathcal{P}\mathit{roj}'(\mathcal{M})).
\tag{8.14.6.5}
\end{equation}
\end{env}

\begin{env}[8.14.7]
\label{II.8.14.7-fr}
Conservemos sobre $X'$ las hipótesis de \sref{II.8.14.6-fr}, y sea $\mathcal{F}'$ un \emph{$\mathcal{S}_{X'}$-módulo graduado cuasicoherente}, de modo que $\mathcal{P}\mathit{roj}'(\bbGamma'_*(\mathcal{F}'))$ es también un $\mathcal{S}_{X'}$-módulo graduado \emph{cuasicoherente}.
\oldpage[II]{200}
Vamos a definir un homomorfismo canónico funtorial (de grado $0$) de $\mathcal{S}_{X'}$-módulos graduados
\begin{equation}
\label{II.8.14.7.1-fr}
  \beta':\mathcal{P}\mathit{roj}'(\bbGamma'_*(\mathcal{F}'))\to\mathcal{F}'.
\tag{8.14.7.1}
\end{equation}

Supongamos primero que $Y=\operatorname{Spec}(A)$ es afín y que $\mathcal{S}=\widetilde{S}$, donde $S$ es una $A$-álgebra graduada en grados positivos;
entonces $\bbGamma'_*(\mathcal{F}')=\widetilde{M}$, donde $M=\bigoplus_{n\in\mathbf{Z}}\Gamma(X',\mathcal{F}'_n)$ es un $S$-módulo graduado.
Sea $f\in S_d$ tal que $D_+(f)\subset X'$;
por definición \sref{II.2.6.2-fr}, $\alpha_d(f)$ restringida a $D_+(f)$ es la sección de $\mathcal{O}_X(d)$ sobre $D_+(f)$ correspondiente al elemento $f/1$ de $(S(d))_{(f)}$ y, por tanto, es invertible;
por consiguiente, también lo es $\alpha_{nd}(f^n)$ para todo $n>0$.
De aquí se concluye inmediatamente que se ha definido un $S_f$-homomorfismo (de grado $0$) de módulos graduados $\beta_f:M_f\to\Gamma(D_+(f),\mathcal{F}')$ enviando cada elemento $z/f^n\in M_f$ (donde $z\in M$) a la sección $(z|D_+(f))(\alpha_{nd}(f^n)|D_+(f))^{-1}$ de $\mathcal{F}'$ sobre $D_+(f)$.
Además, se tiene un diagrama conmutativo correspondiente a \sref{II.2.6.4.1-fr}, de donde resulta la definición de $\beta'$ en este caso.
Para pasar al caso general, hay que considerar una $A$-álgebra $A'$, la $A'$-álgebra graduada $S'=S\otimes_A A'$, y utilizar el diagrama conmutativo análogo a \sref{II.2.8.13.2-fr};
dejamos los detalles al lector.
\end{env}

\begin{proposition}[8.14.8]
\label{II.8.14.8-fr}
Si $X'$ es un abierto de $X=\operatorname{Proj}(\mathcal{S})$ tal que $q':X'\to Y$ es cuasicompacto, el homomorfismo $\beta'$ definido en \sref{II.8.14.7-fr} es biyectivo.
\end{proposition}

\begin{proof}
Se puede restringir claramente al caso en que $Y$ es afín, y todo se reduce entonces a demostrar (con la notación de \sref{II.8.14.7-fr}) que el homomorfismo $\beta_f:M_f\to\Gamma(D_+(f),\mathcal{F}')$ es un isomorfismo.
Pero sustituir $f$ por una de sus potencias no cambia ni $D_+(f)$ ni $\beta_f$;
como $X'$ es \emph{cuasicompacto} por hipótesis, se puede suponer siempre, por \sref{II.8.14.4-fr}, que el haz $\mathcal{O}_X(d)$ es \emph{invertible}.
Como $X'$ es un esquema (porque $q'$ es separado), la proposición es entonces exactamente \sref[I]{I.9.3.1-fr}.
\end{proof}

\begin{corollary}[8.14.9]
\label{II.8.14.9-fr}
Bajo las hipótesis de \sref{II.8.14.8-fr}, todo $\mathcal{S}_{X'}$-módulo graduado cuasicoherente $\mathcal{F}'$ es isomorfo a un $\mathcal{S}_{X'}$-módulo graduado de la forma $\mathcal{P}\mathit{roj}'(\mathcal{M})$, donde $\mathcal{M}$ es un $\mathcal{S}$-módulo graduado cuasicoherente.
Además, si $\mathcal{F}'$ es de tipo finito y se supone que $Y$ es un esquema cuasicompacto o un preesquema cuyo espacio subyacente es noetheriano, se puede suponer que $\mathcal{M}$ es de tipo finito.
\end{corollary}

\begin{proof}
La demostración, a partir de \sref{II.8.14.8-fr}, sigue exactamente el mismo camino que la demostración de \sref{II.3.4.5-fr} a partir de \sref{II.3.4.4-fr}, y dejamos los detalles al lector.
\end{proof}

\begin{proposition}[8.14.10]
\label{II.8.14.10-fr}
Bajo las hipótesis de \sref{II.8.14.7-fr}, sean $\mathcal{M}$ un $\mathcal{S}$-módulo graduado cuasicoherente y $\mathcal{F}'$ un $\mathcal{S}_{X'}$-módulo graduado cuasicoherente;
los homomorfismos compuestos
\begin{equation}
\label{II.8.14.10.1-fr}
  \mathcal{P}\mathit{roj}'(\mathcal{M})
  \xrightarrow{\mathcal{P}\mathit{roj}'(\alpha')}
  \mathcal{P}\mathit{roj}'(\bbGamma'_*(\mathcal{P}\mathit{roj}'(\mathcal{M})))
  \xrightarrow{\beta'}\mathcal{P}\mathit{roj}'(\mathcal{M})
\tag{8.14.10.1}
\end{equation}
\begin{equation}
\label{II.8.14.10.2-fr}
  \bbGamma'_*(\mathcal{F}')
  \xrightarrow{\alpha'}
  \bbGamma'_*(\mathcal{P}\mathit{roj}'(\bbGamma'_*(\mathcal{F}')))
  \xrightarrow{\bbGamma'_*(\beta')}\bbGamma'_*(\mathcal{F}')
\tag{8.14.10.2}
\end{equation}
son los isomorfismos identidad.
\end{proposition}

\begin{proof}
La cuestión es local sobre $Y$, y la demostración se hace como en \sref{II.2.6.5-fr};
dejamos los detalles al lector.
\end{proof}

\begin{remark}[8.14.11]
\label{II.8.14.11-fr}
En el capítulo~III \sref[III]{III.2.3.1-fr}, veremos que, cuando $Y$ es \emph{localmente noetheriano} y $\mathcal{S}$ es una $\mathcal{O}_Y$-álgebra graduada cuasicoherente \emph{de tipo finito} (en cuyo caso
\oldpage[II]{201}
se puede tomar $X'=X$), entonces el homomorfismo $\alpha$ \sref{II.8.14.5.4-fr} es \textbf{(TN)}-\emph{biyectivo} para todo $\mathcal{S}$-módulo graduado cuasicoherente $\mathcal{M}$ que satisfaga la condición~\textbf{(TF)}.
\end{remark}

\begin{remark}[8.14.12]
\label{II.8.14.12-fr}
La situación descrita en \sref{II.8.14.4-fr} es un caso particular de la siguiente.
Sea $X$ un espacio anillado y sea $\mathcal{S}$ una $\mathcal{O}_X$-álgebra graduada (en grados positivos y negativos);
supongamos que existe un entero $d>0$ tal que $\mathcal{S}_d$ y $\mathcal{S}_{-d}$ son \emph{invertibles}, y que el homomorfismo canónico
\begin{equation}
\label{II.8.14.12.1-fr}
  \mathcal{S}_d\otimes_{\mathcal{O}_X}\mathcal{S}_{-d}
  \to\mathcal{O}_X
\tag{8.14.12.1}
\end{equation}
es un \emph{isomorfismo} (de modo que $\mathcal{S}_{-d}$ se identifica con $\mathcal{S}_d^{-1}$).
Se dice entonces que la $\mathcal{O}_X$-álgebra graduada $\mathcal{S}$ es \emph{periódica}, \emph{de período $d$}.
Esta terminología proviene de la propiedad siguiente:
\emph{bajo las hipótesis precedentes, para todo $\mathcal{S}$-módulo graduado $\mathcal{F}$, el homomorfismo canónico}
\begin{equation}
\label{II.8.14.12.2-fr}
  \mathcal{S}_d\otimes\mathcal{F}_n\to\mathcal{F}_{n+d}
\tag{8.14.12.2}
\end{equation}
\emph{es un isomorfismo para todo $n\in\mathbf{Z}$.}
En efecto, la cuestión es local sobre $X$, y se puede suponer que $\mathcal{S}_d$ posee una sección \emph{invertible} $s$ sobre $X$, cuya inversa $s'$ es una sección de $\mathcal{S}_{-d}$.
El homomorfismo $\mathcal{F}_{n+d}\to\mathcal{S}_d\otimes\mathcal{F}_n$ que envía cada sección $z\in\Gamma(U,\mathcal{F}_{n+d})$ a la sección $(s|U)\otimes(s'|U)z$ de $\mathcal{S}_d\otimes\mathcal{F}_n$ sobre $U$ es entonces el inverso de \sref{II.8.14.12.2-fr}, de donde se sigue la afirmación.
Esto induce, para todo $k\in\mathbf{Z}$, un isomorfismo canónico
\[
  (\mathcal{S}_d)^{\otimes k}\otimes\mathcal{F}_n
  \xrightarrow{\sim}\mathcal{F}_{n+kd}.
\]
Por tanto, \emph{dar un $\mathcal{S}$-módulo graduado $\mathcal{F}$ equivale a dar los $\mathcal{S}_0$-módulos $\mathcal{F}_i$ ($(0\leqslant i\leqslant d-1)$) y los homomorfismos canónicos}
\[
  \mathcal{S}_i\otimes\mathcal{F}_j\to\mathcal{F}_{i+j}
  \qquad\text{para }0\leqslant i,j\leqslant d-1
\]
(poniendo $\mathcal{F}_{i+j}=\mathcal{S}_d\otimes_{\mathcal{S}_0}
\mathcal{F}_{i+j-d}$ cuando $i+j\geqslant d$).
Naturalmente, para que estos homomorfismos definan una estructura bien determinada de $\mathcal{S}$-módulo sobre la suma directa de los $(\mathcal{S}_d)^{\otimes k}\otimes\mathcal{F}_i$ $(k\in\mathbf{Z},\ 0\leqslant i\leqslant d-1)$, deben satisfacer ciertas condiciones de asociatividad que no explicitaremos.

En el caso $d=1$ (que es el considerado en \sref{subsection:II.3.3-fr}), se puede decir, por tanto, que la categoría de los $\mathcal{S}$-módulos graduados (resp. cuasicoherentes si $X$ es un preesquema y $\mathcal{S}$ es cuasicoherente) es \emph{equivalente} a la categoría de los $\mathcal{S}_0$-módulos arbitrarios (resp. cuasicoherentes);
en este sentido pueden considerarse los resultados de este número como una generalización de los del §\hyperref[section:II.3-fr]{3}.
Además, se ve que, bajo condiciones de finitud adecuadas, los resultados de este número (junto con \sref{II.8.14.11-fr}) reducen, en cierto sentido, el estudio de las álgebras graduadas cuasicoherentes sobre un preesquema, y de los módulos graduados «módulo \textbf{(TN)}» sobre tales álgebras, al estudio del caso particular en que las álgebras consideradas son \emph{periódicas} (y en el que la condición~(\textbf{TN}) para $\mathcal{M}$ \sref{II.3.4.2-fr} implica, por tanto, que $\mathcal{M}=0$).
\end{remark}

\begin{remark}[8.14.13]
\label{II.8.14.13-fr}
Bajo las hipótesis de \sref{II.8.14.1-fr}, sea $d$ un entero $>0$;
se ha definido un $Y$-isomorfismo canónico $h$ de $X$ sobre $X^{(d)}=\operatorname{Proj}(\mathcal{S}^{(d)})$ \sref{II.3.1.8-fr}.
Para todo
\oldpage[II]{202}
$\mathcal{S}$-módulo graduado cuasicoherente $\mathcal{M}$ y todo entero $k$ tal que $0\leqslant k\leqslant d-1$, se tiene también (con la notación de \sref{II.3.1.1-fr}) un $h$-isomorfismo canónico
\begin{equation}
\label{II.8.14.13.1-fr}
  (\mathcal{P}\mathit{roj}(\mathcal{M}))^{(d,k)}
  \xleftarrow{\sim}\mathcal{P}\mathit{roj}(\mathcal{M}^{(d,k)}).
\tag{8.14.13.1}
\end{equation}

Supongamos primero que $Y=\operatorname{Spec}(A)$ es afín, $\mathcal{S}=\widetilde{S}$ y $\mathcal{M}=\widetilde{M}$, donde $S$ es una $A$-álgebra graduada en grados positivos y $M$ un $S$-módulo graduado.
Se sabe que, para todo $f\in S_e$ ($e>0$), $h$ envía $D_+(f)$ sobre $D_+(f^d)$ y corresponde al isomorfismo canónico $S_{(f^d)}\xrightarrow{\sim}S_{(f)}$ \sref{II.2.2.2-fr}.
La restricción de \sref{II.8.14.13.1-fr} a $D_+(f^d)$ corresponde entonces al di-isomorfismo canónico $M_{f^d}\xrightarrow{\sim}M_f$, restringido a los elementos de $M_{f^d}$ cuyo grado es congruente con $k$ módulo $d$.
Dejamos al lector la tarea de comprobar que estos isomorfismos son compatibles con el paso de $f$ a uno de sus múltiplos homogéneos $fg$, y después que existe una compatibilidad análoga con el paso de $S$ a una $A'$-álgebra graduada $S'=S\otimes_A A'$, donde $A'$ es una $A$-álgebra.
En particular, se obtiene así un $h$-isomorfismo
\begin{equation}
\label{II.8.14.13.2-fr}
  (\mathcal{S}^{(d)})_{X^{(d)}}
  \xrightarrow{\sim}(\mathcal{S}_X)^{(d)}
\tag{8.14.13.2}
\end{equation}
que respeta las estructuras multiplicativas de la fuente y del destino y que, gracias a \sref{II.8.14.13.1-fr}, se convierte en un $h$-di-isomorfismo de un $(\mathcal{S}^{(d)})_{X^{(d)}}$-módulo graduado sobre un $(\mathcal{S}_X)^{(d)}$-módulo graduado.
Análogamente, se tiene un $h$-isomorfismo
\begin{equation}
\label{II.8.14.13.3-fr}
  \mathcal{P}\mathit{roj}_0(\mathcal{S}^{(d,k)}(n))
  \xrightarrow{\sim}\mathcal{O}_X(nd+k)
\tag{8.14.13.3}
\end{equation}
que completa el resultado de \sref{II.3.2.9-fr}[ii].

El isomorfismo de \sref{II.8.14.13.1-fr} induce inmediatamente un isomorfismo de $\mathcal{S}^{(d)}$-módulos graduados
\begin{equation}
\label{II.8.14.13.4-fr}
  \bbGamma_*^{(d)}(\mathcal{P}\mathit{roj}(\mathcal{M}^{(d,k)}))
  \xrightarrow{\sim}
  \bbGamma_*((\mathcal{P}\mathit{roj}(\mathcal{M}))^{(d,k)})
\tag{8.14.13.4}
\end{equation}
donde $\bbGamma_*^{(d)}$ corresponde al morfismo estructural $q^{(d)}:X^{(d)}\to Y$;
se comprueba de inmediato que el homomorfismo canónico $\alpha$ \sref{II.8.14.5.4-fr} y el homomorfismo análogo $\alpha^{(d)}$ para $X^{(d)}$ hacen conmutativo el diagrama siguiente:
\begin{equation}
\label{II.8.14.13.5-fr}
  \xymatrix{
    & \mathcal{M}^{(d,k)}\ar[dl]_{\alpha^{(d)}}\ar[dr]^{\alpha} & \\
    \bbGamma_*^{(d)}(\mathcal{P}\mathit{roj}(\mathcal{M}^{(d,k)}))
      \ar[rr]^{\sim}
      && \bbGamma_*((\mathcal{P}\mathit{roj}(\mathcal{M}))^{(d,k)})
  }
\tag{8.14.13.5}
\end{equation}
La comprobación se efectúa suponiendo que $Y$ es afín y calculando las restricciones de las imágenes por $\alpha^{(d)}$ y $\alpha$ de un mismo elemento de $M^{(d,k)}$ a los abiertos $D_+(f^d)$ y $D_+(f)$ (con la misma notación que antes).
\end{remark}

\begin{proposition}[8.14.14]
\label{II.8.14.14-fr}
Sea $Y$ un preesquema cuasicompacto, sea $\mathcal{S}$ una $\mathcal{O}_Y$-álgebra graduada cuasicoherente de tipo finito y sea $\mathcal{M}$ un $\mathcal{S}$-módulo graduado cuasicoherente que satisfaga la condición~\textup{\textbf{(TF)}};
sea $X=\operatorname{Proj}(\mathcal{S})$.
Entonces $\mathcal{S}_X$ es una $\mathcal{O}_X$-álgebra graduada periódica \sref{II.8.14.12-fr}, y existe
\oldpage[II]{203}
algún período $d$ de $\mathcal{S}_X$ tal que los $(\mathcal{P}\mathit{roj}(\mathcal{M}))^{(d,k)}$ ($(0\leqslant k\leqslant d-1)$) son $(\mathcal{S}_X)^{(d)}$-módulos de tipo finito.
\end{proposition}

\begin{proof}
En efecto, \sref{II.3.1.10-fr} prueba que existe un $d$ tal que $\mathcal{S}^{(d)}$ está engendrada por $\mathcal{S}_d=(\mathcal{S}^{(d)})_1$, siendo esta última un $\mathcal{S}_0$-módulo de tipo finito.
Para probar la primera afirmación se puede, pues, en virtud de \sref{II.8.14.13.2-fr}, reducir al caso $d=1$, y entonces la proposición se sigue de \sref{II.3.2.7-fr}.
Además, teniendo en cuenta \sref{II.8.14.13.1-fr}, la segunda afirmación es consecuencia de \sref{II.2.1.6-fr}[iii] y \sref{II.3.4.3-fr}.
\end{proof}
\clearpage
\oldpage[II]{205}
\phantomsection
\addcontentsline{toc}{section}{Bibliografía}
\renewcommand{\refname}{Bibliografía}
\noindent\textit{Nota editorial. Para su uso autónomo, se reproducen aquí las entradas [1], [9],
[11] y [13], citadas en este volumen, a partir de la bibliografía de la serie precedente.
La edición original de EGA~II solo imprime la continuación [23]--[26] en la página impresa~205.}
\par\medskip

\begin{thebibliography}{26}

\bibitem[1]{I-1}
H.~Cartan y C.~Chevalley,
\emph{Séminaire de l'École Normale Supérieure, 8th year (1955--56):
Géométrie algébrique}.

\bibitem[9]{I-9}
M.~Nagata,
\emph{A general theory of algebraic geometry over Dedekind domains},
Amer. Math. J. \textbf{78} (1956), 78--116, y \textbf{80} (1958),
382--420.

\bibitem[11]{I-11}
P.~Samuel,
\emph{Commutative algebra} (notas de D.~Herzig), Cornell University,
1953.

\bibitem[13]{I-13}
P.~Samuel y O.~Zariski,
\emph{Commutative algebra}, 2 vols., Van Nostrand, Nueva York, 1958--60.


\bibitem[23]{II-23}
C.~\textsc{Chevalley},
Introduction to the theory of algebraic functions of one variable,
\emph{Math. Surveys}, n.º~6, Amer. Math. Soc., Nueva York, 1951.

\bibitem[24]{II-24}
P.~\textsc{Jaffard},
\emph{Les systèmes d'idéaux}, París (Dunod), 1960.

\bibitem[25]{II-25}
M.~\textsc{Nagata},
On the derived normal rings of Noetherian integral domains,
\emph{Mem. Coll. Sci. Kyoto}, sér.~A, t.~XXIX (1955), pp.~293-303.

\bibitem[26]{II-26}
M.~\textsc{Nagata},
Existence theorems for non projective complete algebraic varieties,
\emph{Ill. J. Math.}, t.~II (1958), pp.~490-498.

\end{thebibliography}
\clearpage

\clearpage
\oldpage[II]{207}
\phantomsection
\addcontentsline{toc}{section}{Índice de notaciones}
\section*{Índice de notaciones}

\begingroup
\small
\setlength{\parindent}{0pt}
\setlength{\parskip}{2pt}

$\mathcal{A}(X),\mathcal{A}(\mathcal{F}),\mathcal{A}(\mathcal{B})$
($X$ un $S$-preesquema, $\mathcal{F}$ un $\mathcal{O}_X$-módulo,
$\mathcal{B}$ una $\mathcal{O}_X$-álgebra): \hyperref[II.1.1.1-fr]{1.1.1}.

$\mathcal{A}(h)$ ($h$ un $S$-morfismo): \hyperref[II.1.1.2-fr]{1.1.2}.

$\Spec(\mathcal{B})$ ($\mathcal{B}$ una $\mathcal{O}_S$-álgebra): \hyperref[II.1.3.1-fr]{1.3.1}.

$\widetilde{\mathcal{M}}$ ($\mathcal{M}$ un $\mathcal{B}$-módulo,
$\mathcal{B}$ una $\mathcal{O}_S$-álgebra): \hyperref[II.1.4.3-fr]{1.4.3}.

$\mathbf{T}(E),\mathbf{S}(E),\mathbf{S}_A(E)$ ($E$ un $A$-módulo): \hyperref[II.1.7.1-fr]{1.7.1}.

$\mathbf{S}(\mathcal{E}),\mathbf{S}_{\mathcal{A}}(\mathcal{E})$
($\mathcal{E}$ un $\mathcal{A}$-módulo): \hyperref[II.1.7.4-fr]{1.7.4}.

$\mathbf{V}(\mathcal{E})$ ($\mathcal{E}$ un $\mathcal{O}_S$-módulo cuasicoherente): \hyperref[II.1.7.8-fr]{1.7.8}.

$S_n,S_+,M_n,S^{(d)},M^{(d,k)},M^{(d)}$
($S$ un anillo graduado, $M$ un $S$-módulo graduado): \hyperref[II.2.1.1-fr]{2.1.1}.

$M(n),S(n)$ ($S$ un anillo graduado, $M$ un $S$-módulo graduado): \hyperref[II.2.1.1-fr]{2.1.1}.

$\Hom_S(M,N)$ ($S$ un anillo graduado, $M,N$ $S$-módulos graduados): \hyperref[II.2.1.2-fr]{2.1.2}.

$\mathfrak N_+,\mathfrak r_+(\mathfrak J)$
($\mathfrak J$ un ideal de $S_+$): \hyperref[II.2.1.10-fr]{2.1.10}.

$S_{(f)},M_{(f)}$ ($f$ un elemento homogéneo de $S_+$): \hyperref[II.2.2.1-fr]{2.2.1}.

$S_{(T)},M_{(T)},S_{(\mathfrak p)},M_{(\mathfrak p)}$
($T$ una parte multiplicativa de $S_+$ formada por elementos homogéneos,
$\mathfrak p$ un ideal primo graduado de $S_+$): \hyperref[II.2.2.7-fr]{2.2.7}.

$\Proj(S)$ ($S$ un anillo graduado): \hyperref[II.2.3.1-fr]{2.3.1}.

$V_+(E)$ ($E$ una parte de $S_+$): \hyperref[II.2.3.2-fr]{2.3.2}.

$D_+(f)$ ($f$ un elemento de $S$): \hyperref[II.2.3.3-fr]{2.3.3}.

$\mathfrak j_+(Y)$ ($Y$ una parte de $\Proj(S)$): \hyperref[II.2.3.10-fr]{2.3.10}.

$\widetilde M$ ($M$ un $S$-módulo graduado): \hyperref[II.2.5.2-fr]{2.5.2}.

$\widetilde u$ ($u$ un homomorfismo de grado $0$): \hyperref[II.2.5.4-fr]{2.5.4}.

$\mathcal{O}_X(n),\mathcal{F}(n)$
($X=\Proj(S)$, $\mathcal{F}$ un $\mathcal{O}_X$-módulo): \hyperref[II.2.5.10-fr]{2.5.10}.

$\lambda$ (homomorfismo canónico): \hyperref[II.2.5.11-fr]{2.5.11}.

$\mu$ (homomorfismo canónico): \hyperref[II.2.5.12-fr]{2.5.12}.

$X^{(d)}$ ($X=\Proj(S)$, $d$ un entero $>0$): \hyperref[II.2.5.16-fr]{2.5.16}.

$\Gamma_*(\mathcal{F})$ ($\mathcal{F}$ un $\mathcal{O}_X$-módulo,
$X=\Proj(S)$): \hyperref[II.2.6.1-fr]{2.6.1}.

$\alpha_n,\alpha,\alpha_M$ (homomorfismos canónicos): \hyperref[II.2.6.2-fr]{2.6.2}.

$\beta,\beta_{\mathcal{F}}$ (homomorfismos canónicos): \hyperref[II.2.6.4-fr]{2.6.4}.

$G(\vphi),{}^a\vphi$ (por abuso de lenguaje;
$\vphi$ un homomorfismo de anillos graduados): \hyperref[II.2.8.1-fr]{2.8.1}.

$\Proj(\vphi)$ ($\vphi$ un homomorfismo de anillos graduados): \hyperref[II.2.8.2-fr]{2.8.2}.

$\mathcal{S}^{(d)},\mathcal{M}^{(d,k)},\mathcal{M}^{(d)},\mathcal{M}(n)$
($\mathcal{S}$ una $\mathcal{O}_Y$-álgebra graduada,
$\mathcal{M}$ un $\mathcal{S}$-módulo graduado): \hyperref[II.3.1.1-fr]{3.1.1}.

$\Proj(\mathcal{S})$ ($\mathcal{S}$ una $\mathcal{O}_Y$-álgebra graduada cuasicoherente): \hyperref[II.3.1.3-fr]{3.1.3}.

$X_f$ ($X=\Proj(\mathcal{S})$, $f$ una sección de $\mathcal{S}_d$ sobre $Y$): \hyperref[II.3.1.4-fr]{3.1.4}.

$\widetilde{\mathcal{M}}$ ($\mathcal{M}$ un $\mathcal{S}$-módulo graduado): \hyperref[II.3.2.2-fr]{3.2.2}.

$\mathcal{O}_X(n),\mathcal{F}(n)$
($X=\Proj(\mathcal{S})$, $\mathcal{F}$ un $\mathcal{O}_X$-módulo): \hyperref[II.3.2.5-fr]{3.2.5}.

$\lambda,\mu$ (homomorfismos canónicos): \hyperref[II.3.2.6-fr]{3.2.6}.

$\bbGamma_*(\mathcal{F})$ ($\mathcal{F}$ un $\mathcal{O}_X$-módulo,
$X=\Proj(\mathcal{S})$): \hyperref[II.3.3.1-fr]{3.3.1}.

$\alpha_n,\alpha,\alpha_{\mathcal{M}}$ (homomorfismos canónicos): \hyperref[II.3.3.2-fr]{3.3.2}.
\oldpage[II]{208}
$\beta,\beta_{\mathcal{F}}$ (homomorfismos canónicos): \hyperref[II.3.3.4-fr]{3.3.4}.

$G(\vphi),\Proj(\vphi)$
($\vphi$ un homomorfismo de $\mathcal{O}_Y$-álgebras graduadas): \hyperref[II.3.5.1-fr]{3.5.1}.

$\Proj(u)$ ($u$ un $g$-morfismo de una $\mathcal{O}_Y$-álgebra graduada
en una $\mathcal{O}_{Y'}$-álgebra graduada): \hyperref[II.3.5.6-fr]{3.5.6}.

$r_{\mathcal{L},\psi}$ ($\mathcal{L}$ un $\mathcal{O}_X$-módulo invertible): \hyperref[II.3.7.1-fr]{3.7.1}.

$\mathbf{P}(\mathcal{E}),\mathbf{P}(E),\mathbf{P}_Y^n,\mathbf{P}_A^n$: \hyperref[II.4.1.1-fr]{4.1.1}.

$\mathbf{P}(u)$ ($u$ un homomorfismo sobreyectivo de $\mathcal{O}_Y$-módulos): \hyperref[II.4.1.2-fr]{4.1.2}.

$\Hyp_Y(X,\mathcal{E})$ ($\mathcal{E}$ un $\mathcal{O}_Y$-módulo cuasicoherente): \hyperref[II.4.2.5-fr]{4.2.5}.

$\varsigma$ (morfismo de Segre): \hyperref[II.4.3.1-fr]{4.3.1}.

$\mathcal{F}(n)$ ($\mathcal{F}$ un $\mathcal{O}_X$-módulo, $X$ provisto de un
$\mathcal{O}_X$-módulo invertible): \hyperref[II.4.5.1-fr]{4.5.1}.

$\varepsilon$ (automorfismo identidad): \hyperref[II.4.5.1-fr]{4.5.1}.

$P(u,T),\sigma_i(u)$ ($u$ un endomorfismo de un módulo libre de base finita): \hyperref[II.6.4.1-fr]{6.4.1}.

$P(u,T),\sigma_i(u),\det u$
($u$ un endomorfismo de un módulo de tipo finito sobre un anillo íntegro o un
anillo noetheriano reducido): \hyperref[II.6.4.2-fr]{6.4.2} y \hyperref[II.6.4.7-fr]{6.4.7}.

$\sigma_i(u),\det u$ ($u$ un endomorfismo de un $\mathcal{A}$-módulo localmente libre): \hyperref[II.6.4.8-fr]{6.4.8}.

$\sigma_i,\det$ (homomorfismos de haces): \hyperref[II.6.4.8-fr]{6.4.8}, \hyperref[II.6.4.9-fr]{6.4.9} y \hyperref[II.6.4.10-fr]{6.4.10}.

$\sigma_i(u),\det u$
($u$ un endomorfismo de un $\mathcal{O}_X$-módulo de tipo finito sobre un preesquema
localmente íntegro o un preesquema localmente noetheriano reducido):
\hyperref[II.6.4.9-fr]{6.4.9} y \hyperref[II.6.4.10-fr]{6.4.10}.

$N_{\mathcal{B}/\mathcal{A}}(f)$
($f$ una sección de la $\mathcal{A}$-álgebra $\mathcal{B}$): \hyperref[II.6.5.1-fr]{6.5.1}.

$N_{\mathcal{B}/\mathcal{A}}(\mathcal{L}')$
($\mathcal{L}'$ un $\mathcal{B}$-módulo invertible): \hyperref[II.6.5.2-fr]{6.5.2}.

$N_{\mathcal{B}/\mathcal{A}}(h')$
($h'$ un homomorfismo de $\mathcal{B}$-módulos invertibles): \hyperref[II.6.5.3-fr]{6.5.3}.

$N_{X'/X}(\mathcal{L}')$
($X'$ finito sobre $X$, $\mathcal{L}'$ un $\mathcal{O}_{X'}$-módulo invertible): \hyperref[II.6.5.5-fr]{6.5.5}.

$N_{X'/X}(h')$
($h'$ un homomorfismo de $\mathcal{O}_{X'}$-módulos invertibles): \hyperref[II.6.5.5-fr]{6.5.5}.

$\mathcal{S}_X$ ($X=\Proj(\mathcal{S})$, $\mathcal{S}$ una suma de ideales de
$\mathcal{O}_Y$): \hyperref[II.8.1.5-fr]{8.1.5}.

$S^\geq,S^\leq,S_f^\geq,S_f^\leq,
M^\geq,M^\leq,M_f^\geq,M_f^\leq$
($S$ un anillo graduado, $M$ un $S$-módulo graduado,
$f$ un elemento homogéneo de $S_+$): \hyperref[II.8.2.1-fr]{8.2.1}.

$\widehat S$ ($S$ un anillo graduado): \hyperref[II.8.2.2-fr]{8.2.2}.

$\widehat M$ ($M$ un $S$-módulo graduado): \hyperref[II.8.2.4-fr]{8.2.4}.

$S_{[n]},S^\natural,f^\natural$
($S$ un anillo graduado, $f$ homogéneo en $S$): \hyperref[II.8.2.6-fr]{8.2.6}.

$M_{[n]},M^\natural$ ($M$ un $S$-módulo graduado): \hyperref[II.8.2.8-fr]{8.2.8}.

$\widehat{\mathcal{S}}$ ($\mathcal{S}$ una $\mathcal{O}_Y$-álgebra graduada): \hyperref[II.8.3.1-fr]{8.3.1}.

$G_m$ (esquema de grupos): \hyperref[II.8.3.9-fr]{8.3.9}.

$\mathcal{S}_X$ ($X=\Proj(\mathcal{S})$, $\mathcal{S}$ una $\mathcal{O}_Y$-álgebra graduada): \hyperref[II.8.6.1-fr]{8.6.1}.

$\mathcal{S}_X,C_X,\widehat C_X,E_X,\widehat E_X$
($X=\Proj(\mathcal{S})$, $\mathcal{S}$ una $\mathcal{O}_Y$-álgebra graduada
con $\mathcal{S}_0=\mathcal{O}_Y$): \hyperref[II.8.6.1-fr]{8.6.1}.

$\mathcal{S}_{[n]},\mathcal{S}^{\natural}$
($\mathcal{S}$ una $\mathcal{O}_Y$-álgebra graduada): \hyperref[II.8.7.2-fr]{8.7.2}.

$\egaShProj_0(\mathcal{M}),\egaShProj(\mathcal{M}),\mathcal{M}_X,
\widehat{\mathcal{M}},\mathcal{M}^{\square}$
($\mathcal{M}$ un $\mathcal{S}$-módulo graduado): \hyperref[II.8.12.1-fr]{8.12.1}.

$\mathcal{M}_X^\geq$
($X=\Proj(\mathcal{S})$, $\mathcal{M}$ un $\mathcal{S}$-módulo graduado): \hyperref[II.8.12.5-fr]{8.12.5}.

$\mathcal{M}_{[n]},\mathcal{M}^{\natural}$
($\mathcal{M}$ un $\mathcal{S}$-módulo graduado): \hyperref[II.8.12.9-fr]{8.12.9}.

$(\widetilde{\mathcal{N}})^-$
($\mathcal{N}$ un sub-$\mathcal{S}$-módulo de un $\mathcal{S}$-módulo graduado): \hyperref[II.8.13.1-fr]{8.13.1}.

$\overline{\mathcal{N}}$
($\mathcal{N}$ un sub-$\mathcal{S}$-módulo de un $\mathcal{S}$-módulo graduado): \hyperref[II.8.13.2-fr]{8.13.2}.

$\bbGamma_*(\mathcal{F})$
($X=\Proj(\mathcal{S})$, $\mathcal{F}$ un $\mathcal{S}_X$-módulo graduado): \hyperref[II.8.14.5-fr]{8.14.5}.

$\mathcal{S}_{X'},\bbGamma'_*(\mathcal{F}'),\egaShProj'(\mathcal{M}),
\alpha',\beta'$
($X'$ abierto en $X=\Proj(\mathcal{S})$, $\mathcal{F}'$ un $\mathcal{S}_{X'}$-módulo graduado,
$\mathcal{M}$ un $\mathcal{S}$-módulo graduado): \hyperref[II.8.14.6-fr]{8.14.6}.

\endgroup

\clearpage
\oldpage[II]{209}
\phantomsection
\addcontentsline{toc}{section}{Índice terminológico}
\section*{Índice terminológico}

\begingroup
\small
\setlength{\parindent}{0pt}
\setlength{\parskip}{2pt}
\newcommand{\egaTerm}[2]{#1: #2.\par}

\egaTerm{Afín (morfismo)}{\hyperref[II.1.6.1-fr]{1.6.1}}
\egaTerm{Afín (preesquema) sobre un preesquema}{\hyperref[II.1.2.1-fr]{1.2.1}}
\egaTerm{Álgebra graduada sobre un anillo graduado}{\hyperref[II.2.1.2-fr]{2.1.2}}
\egaTerm{Álgebra simétrica de un $A$-módulo}{\hyperref[II.1.7.1-fr]{1.7.1}}
\egaTerm{$\mathcal{A}$-álgebra simétrica de un $\mathcal{A}$-módulo}{\hyperref[II.1.7.4-fr]{1.7.4}}
\egaTerm{$\mathcal{O}_Y$-álgebra graduada esencialmente reducida}{\hyperref[II.3.1.12-fr]{3.1.12}}
\egaTerm{$\mathcal{O}_Y$-álgebra graduada íntegra}{\hyperref[II.3.1.12-fr]{3.1.12}}
\egaTerm{Amplio ($\mathcal{O}_X$-módulo invertible)}{\hyperref[II.4.5.3-fr]{4.5.3}}
\egaTerm{Amplio relativamente a $f$; $f$-amplio; $Y$-amplio ($\mathcal{O}_X$-módulo invertible)}{\hyperref[II.4.6.1-fr]{4.6.1}}
\egaTerm{Anillo graduado esencialmente íntegro}{\hyperref[II.2.1.11-fr]{2.1.11}}
\egaTerm{Anillo graduado esencialmente reducido}{\hyperref[II.2.1.10-fr]{2.1.10}}
\egaTerm{Asociado a un $\mathcal{B}$-módulo (haz)}{\hyperref[II.1.4.3-fr]{1.4.3}}
\egaTerm{Asociado a un $S$-módulo graduado (haz)}{\hyperref[II.2.5.3-fr]{2.5.3}}
\egaTerm{Asociado a un $\mathcal{S}$-módulo graduado (haz)}{\hyperref[II.3.2.2-fr]{3.2.2}}
\egaTerm{Asociado a una $\mathcal{O}_S$-álgebra ($S$-esquema)}{\hyperref[II.1.3.1-fr]{1.3.1}}

\egaTerm{Condición (TF); condición (TN)}{\hyperref[II.2.7.2-fr]{2.7.2}}
\egaTerm{Condición (\textbf{TF}); condición (\textbf{TN})}{\hyperref[II.3.4.2-fr]{3.4.2}}
\egaTerm{Cono afín, cono proyectivo definido por una $\mathcal{O}_Y$-álgebra graduada}{\hyperref[II.8.3.1-fr]{8.3.1}}
\egaTerm{Cono afín sin vértice, cono proyectivo sin vértice definido por una $\mathcal{O}_Y$-álgebra graduada}{\hyperref[II.8.3.4-fr]{8.3.4}}
\egaTerm{Cono proyectante afín, cono proyectante proyectivo de un preesquema $\Proj(\mathcal{S})$}{\hyperref[II.8.3.1-fr]{8.3.1}}
\egaTerm{Curva algebraica sobre un cuerpo}{\hyperref[II.7.4.2-fr]{7.4.2}}
\egaTerm{Curva algebraica completa}{\hyperref[II.7.4.7-fr]{7.4.7}}
\egaTerm{Criterio de Serre}{\hyperref[II.5.2.1-fr]{5.2.1}}

\egaTerm{Determinante de un endomorfismo de un módulo de tipo finito sobre un anillo íntegro}{\hyperref[II.6.4.2-fr]{6.4.2}}
\egaTerm{Determinante de un endomorfismo de un módulo de tipo finito sobre un anillo noetheriano reducido}{\hyperref[II.6.4.7-fr]{6.4.7}}
\egaTerm{Determinante de un endomorfismo de un $\mathcal{A}$-módulo localmente libre}{\hyperref[II.6.4.8-fr]{6.4.8}}
\egaTerm{Determinante de un endomorfismo de un $\mathcal{O}_X$-módulo de tipo finito sobre un preesquema $X$ localmente íntegro}{\hyperref[II.6.4.9-fr]{6.4.9}}
\egaTerm{Determinante de un endomorfismo de un $\mathcal{O}_X$-módulo de tipo finito sobre un preesquema $X$ localmente noetheriano y reducido}{\hyperref[II.6.4.10-fr]{6.4.10}}

\egaTerm{Explosión (preesquema)}{\hyperref[II.8.1.3-fr]{8.1.3}}
\egaTerm{Entero (morfismo)}{\hyperref[II.6.1.1-fr]{6.1.1}}
\egaTerm{Entero (preesquema) sobre un preesquema}{\hyperref[II.6.1.1-fr]{6.1.1}}
\egaTerm{Entera ($\mathcal{O}_S$-álgebra cuasicoherente)}{\hyperref[II.6.1.2-fr]{6.1.2}}
\egaTerm{Entera (sección de una $\mathcal{A}$-álgebra) sobre $\mathcal{A}$}{\hyperref[II.6.3.1-fr]{6.3.1}}

\egaTerm{Haz fundamental de un fibrado proyectivo}{\hyperref[II.4.1.1-fr]{4.1.1}}
\egaTerm{Clausura integral de un haz de anillos $\mathcal{A}$ en una $\mathcal{A}$-álgebra}{\hyperref[II.6.3.2-fr]{6.3.2}}
\egaTerm{Clausura integral de un preesquema $X$ relativamente a una $\mathcal{O}_X$-álgebra}{\hyperref[II.6.3.4-fr]{6.3.4}}
\egaTerm{Clausura proyectiva de un cono afín}{\hyperref[II.8.3.1-fr]{8.3.1}}
\oldpage[II]{210}
\egaTerm{Fibra geométrica racional de un fibrado proyectivo sobre una extensión $K$ de $\boldsymbol{k}(y)$}{\hyperref[II.4.2.6-fr]{4.2.6}}
\egaTerm{Fibrado proyectivo definido por un $\mathcal{O}_Y$-módulo}{\hyperref[II.4.1.1-fr]{4.1.1}}
\egaTerm{Fibrado vectorial definido por un $\mathcal{O}_S$-módulo}{\hyperref[II.1.7.8-fr]{1.7.8}}
\egaTerm{Finito (morfismo)}{\hyperref[II.6.1.1-fr]{6.1.1}}
\egaTerm{Finito (preesquema) sobre un preesquema}{\hyperref[II.6.1.1-fr]{6.1.1}}
\egaTerm{Finita ($\mathcal{O}_S$-álgebra cuasicoherente)}{\hyperref[II.6.1.2-fr]{6.1.2}}
\egaTerm{Funciones simétricas elementales de un endomorfismo de un módulo localmente libre}{\hyperref[II.6.4.1-fr]{6.4.1}}
\egaTerm{Funciones simétricas elementales de un endomorfismo de un módulo de tipo finito sobre un anillo íntegro}{\hyperref[II.6.4.2-fr]{6.4.2}}
\egaTerm{Funciones simétricas elementales de un endomorfismo de un módulo de tipo finito sobre un anillo noetheriano reducido}{\hyperref[II.6.4.7-fr]{6.4.7}}
\egaTerm{Funciones simétricas elementales de un endomorfismo de un $\mathcal{A}$-módulo localmente libre de rango constante}{\hyperref[II.6.4.8-fr]{6.4.8}}
\egaTerm{Funciones simétricas elementales de un endomorfismo de un $\mathcal{O}_X$-módulo de tipo finito sobre un preesquema $X$ localmente íntegro}{\hyperref[II.6.4.9-fr]{6.4.9}}
\egaTerm{Funciones simétricas elementales de un endomorfismo de un $\mathcal{O}_X$-módulo de tipo finito sobre un preesquema $X$ localmente noetheriano reducido}{\hyperref[II.6.4.10-fr]{6.4.10}}
\egaTerm{Funciones simétricas elementales de una sección de una $\mathcal{A}$-álgebra}{\hyperref[II.6.5.1-fr]{6.5.1}}

\egaTerm{Homogeneización de un sub-$\mathcal{S}$-módulo de un $\mathcal{S}$-módulo graduado}{\hyperref[II.8.13.2-fr]{8.13.2}}
\egaTerm{Homomorfismo de anillos graduados}{\hyperref[II.2.1.2-fr]{2.1.2}}
\egaTerm{Homomorfismo de grado $k$ de módulos graduados}{\hyperref[II.2.1.2-fr]{2.1.2}}

\egaTerm{Ideal de $S_+$; ideal primo graduado de $S_+$ ($S$ un anillo graduado)}{\hyperref[II.2.1.10-fr]{2.1.10}}

\egaTerm{Lema de Chow}{\hyperref[II.5.6.1-fr]{5.6.1}}
\egaTerm{Lugar en el infinito de un cono proyectivo}{\hyperref[II.8.3.3-fr]{8.3.3}}

\egaTerm{Morfismo canónico $G(\varepsilon)\to\Proj(S)$ $\bigl(S=\bigoplus_{n\geq0}\Gamma(X,\mathcal{L}^{\otimes n})\bigr)$}{\hyperref[II.4.5.1-fr]{4.5.1}}
\egaTerm{Morfismo canónico $X\to\Spec(\Gamma(X,\mathcal{O}_X))$}{\hyperref[II.5.1.1-fr]{5.1.1}}
\egaTerm{Morfismo de Segre}{\hyperref[II.4.3.1-fr]{4.3.1}}

\egaTerm{Nilradical de $S_+$ ($S$ un anillo graduado)}{\hyperref[II.2.1.10-fr]{2.1.10}}
\egaTerm{Nilradical de $\mathcal{S}_+$ ($\mathcal{S}$ una $\mathcal{O}_Y$-álgebra graduada cuasicoherente)}{\hyperref[II.3.1.12-fr]{3.1.12}}
\egaTerm{Norma de una sección de una $\mathcal{A}$-álgebra}{\hyperref[II.6.5.1-fr]{6.5.1}}
\egaTerm{Norma de un $\mathcal{B}$-módulo invertible}{\hyperref[II.6.5.2-fr]{6.5.2}}
\egaTerm{Norma de una sección de un $\mathcal{B}$-módulo invertible}{\hyperref[II.6.5.3-fr]{6.5.3}}
\egaTerm{Norma de un $\mathcal{O}_{X'}$-módulo invertible}{\hyperref[II.6.5.5-fr]{6.5.5}}

\egaTerm{Periódica (álgebra graduada)}{\hyperref[II.8.14.12-fr]{8.14.12}}
\egaTerm{Polinomio característico de un endomorfismo de un módulo de tipo finito sobre un anillo íntegro}{\hyperref[II.6.4.2-fr]{6.4.2}}
\egaTerm{Polinomio característico de un endomorfismo de un módulo de tipo finito sobre un anillo noetheriano reducido}{\hyperref[II.6.4.7-fr]{6.4.7}}
\egaTerm{Preesquema de vértices de un cono afín (proyectivo)}{\hyperref[II.8.3.3-fr]{8.3.3}}
\egaTerm{Presentación finita (módulo graduado que admite una)}{\hyperref[II.2.1.1-fr]{2.1.1}}
\egaTerm{Presentación finita ($\mathcal{S}$-módulo graduado que admite una)}{\hyperref[II.3.1.1-fr]{3.1.1}}
\egaTerm{Producto tensorial graduado de dos módulos graduados}{\hyperref[II.2.1.2-fr]{2.1.2}}
\egaTerm{Proyectivo (morfismo)}{\hyperref[II.5.5.2-fr]{5.5.2}}
\egaTerm{Proyectivo (preesquema) sobre un preesquema}{\hyperref[II.5.5.2-fr]{5.5.2}}
\egaTerm{Propio (morfismo)}{\hyperref[II.5.4.1-fr]{5.4.1}}
\egaTerm{Propia (parte)}{\hyperref[II.5.4.10-fr]{5.4.10}}
\egaTerm{Propio (preesquema) sobre un preesquema}{\hyperref[II.5.4.1-fr]{5.4.1}}

\egaTerm{Cuasiafín (morfismo)}{\hyperref[II.5.1.1-fr]{5.1.1}}
\egaTerm{Cuasiafín (preesquema) sobre un preesquema}{\hyperref[II.5.1.1-fr]{5.1.1}}
\egaTerm{Cuasiafín (esquema)}{\hyperref[II.5.1.1-fr]{5.1.1}}
\oldpage[II]{211}
\markboth{CAPÍTULO II}{CAPÍTULO II}
\egaTerm{Cuasifinito (morfismo)}{\hyperref[II.6.2.3-fr]{6.2.3}}
\egaTerm{Cuasifinito (preesquema) sobre un preesquema}{\hyperref[II.6.2.3-fr]{6.2.3}}
\egaTerm{Cuasiproyectivo (morfismo)}{\hyperref[II.5.3.1-fr]{5.3.1}}
\egaTerm{Cuasiproyectivo (preesquema) sobre un preesquema}{\hyperref[II.5.3.1-fr]{5.3.1}}

\egaTerm{Radical de un ideal de $S_+$ ($S$ anillo graduado)}{\hyperref[II.2.1.10-fr]{2.1.10}}
\egaTerm{Retracción canónica de un cono proyectivo privado de su vértice sobre su lugar en el infinito}{\hyperref[II.8.3.5-fr]{8.3.5}}

\egaTerm{Sección canónica de $\mathcal{O}_X(-1)$ ($X$ preesquema obtenido por explosión)}{\hyperref[II.8.1.9-fr]{8.1.9}}
\egaTerm{Sección nula de un cono afín}{\hyperref[II.8.3.3-fr]{8.3.3}}
\egaTerm{Sección nula de un fibrado vectorial}{\hyperref[II.1.7.9-fr]{1.7.9}}
\egaTerm{Sección vértice de un cono afín (proyectivo)}{\hyperref[II.8.3.3-fr]{8.3.3}}
\egaTerm{Espectro de una $\mathcal{O}_S$-álgebra}{\hyperref[II.1.3.1-fr]{1.3.1}}
\egaTerm{Espectro primo homogéneo de un anillo graduado}{\hyperref[II.2.3.1-fr]{2.3.1}}
\egaTerm{Espectro homogéneo de una $\mathcal{O}_S$-álgebra graduada cuasicoherente}{\hyperref[II.3.1.3-fr]{3.1.3}}

\egaTerm{(TN)-inyectivo, (TN)-sobreyectivo, (TN)-biyectivo (homomorfismo de módulos graduados)}{\hyperref[II.2.7.2-fr]{2.7.2}}
\egaTerm{(TN)-isomorfismo de módulos graduados}{\hyperref[II.2.7.2-fr]{2.7.2}}
\egaTerm{(TN)-inyectivo, (TN)-sobreyectivo, (TN)-biyectivo (homomorfismo de anillos graduados)}{\hyperref[II.2.9.1-fr]{2.9.1}}
\egaTerm{(TN)-isomorfismo de anillos graduados}{\hyperref[II.2.9.1-fr]{2.9.1}}
\egaTerm{(\textbf{TN})-inyectivo, (\textbf{TN})-sobreyectivo, (\textbf{TN})-biyectivo (homomorfismo de $\mathcal{S}$-módulos graduados)}{\hyperref[II.3.4.2-fr]{3.4.2}}
\egaTerm{(\textbf{TN})-isomorfismo de $\mathcal{S}$-módulos graduados}{\hyperref[II.3.4.2-fr]{3.4.2}}
\egaTerm{(\textbf{TN})-inyectivo, (\textbf{TN})-sobreyectivo, (\textbf{TN})-biyectivo (homomorfismo de $\mathcal{O}_Y$-álgebras graduadas)}{\hyperref[II.3.6.1-fr]{3.6.1}}
\egaTerm{(\textbf{TN})-isomorfismo de $\mathcal{O}_Y$-álgebras graduadas}{\hyperref[II.3.6.1-fr]{3.6.1}}
\egaTerm{Topología espectral sobre $\Proj(S)$}{\hyperref[II.2.3.3-fr]{2.3.3}}
\egaTerm{Muy amplio para $q$, muy amplio para $Y$ ($\mathcal{O}_X$-módulo invertible)}{\hyperref[II.4.4.2-fr]{4.4.2}}
\egaTerm{De tipo finito ($\mathcal{S}$-módulo graduado)}{\hyperref[II.3.1.1-fr]{3.1.1}}

\egaTerm{Universalmente cerrado (morfismo)}{\hyperref[II.5.4.9-fr]{5.4.9}}

\endgroup

\clearpage
\oldpage[II]{212}
\thispagestyle{empty}
\mbox{}

\clearpage
\oldpage[II]{213}
\markboth{ÍNDICE}{ÍNDICE}
\phantomsection
\addcontentsline{toc}{section}{Índice original}
\section*{Índice}
\thispagestyle{plain}

\begingroup
\small
\renewcommand{\arraystretch}{1.04}
\begin{longtable}{@{}p{0.09\textwidth}p{0.69\textwidth}r@{}}
\textbf{Unidad} & \textbf{Título} & \textbf{Página original}\\
\hline
\textbf{II} & \textbf{Estudio global elemental de algunas clases de morfismos} & \textbf{5}\\
\hyperref[section:II.1-fr]{\textsection1.} & Morfismos afines & 5\\
\hyperref[subsection:II.1.1-fr]{1.1} & $S$-preesquemas y $\mathcal{O}_S$-álgebras & 5\\
\hyperref[subsection:II.1.2-fr]{1.2} & Preesquemas afines sobre un preesquema & 6\\
\hyperref[subsection:II.1.3-fr]{1.3} & Preesquema afín sobre $S$ asociado a una $\mathcal{O}_S$-álgebra & 8\\
\hyperref[subsection:II.1.4-fr]{1.4} & Haces cuasicoherentes sobre un preesquema afín sobre $S$ & 9\\
\hyperref[II.1.5-fr]{1.5} & Cambio del preesquema base & 12\\
\hyperref[II.1.6-fr]{1.6} & Morfismos afines & 14\\
\hyperref[II.1.7-fr]{1.7} & Fibrado vectorial asociado a un haz de módulos & 14\\
\hyperref[section:II.2-fr]{\textsection2.} & Espectros primos homogéneos & 19\\
\hyperref[subsection:II.2.1-fr]{2.1} & Generalidades sobre los anillos y módulos graduados & 19\\
\hyperref[subsection:II.2.2-fr]{2.2} & Anillos de fracciones de un anillo graduado & 23\\
\hyperref[subsection:II.2.3-fr]{2.3} & Espectro primo homogéneo de un anillo graduado & 25\\
\hyperref[subsection:II.2.4-fr]{2.4} & La estructura de esquema sobre $\Proj(S)$ & 28\\
\hyperref[subsection:II.2.5-fr]{2.5} & El haz asociado a un módulo graduado & 30\\
\hyperref[subsection:II.2.6-fr]{2.6} & El $S$-módulo graduado asociado a un haz sobre $\Proj(S)$ & 36\\
\hyperref[subsection:II.2.7-fr]{2.7} & Condiciones de finitud & 38\\
\hyperref[subsection:II.2.8-fr]{2.8} & Comportamientos funtoriales & 41\\
\hyperref[subsection:II.2.9-fr]{2.9} & Subpreesquemas cerrados de un esquema $\Proj(S)$ & 48\\
\hyperref[section:II.3-fr]{\textsection3.} & Espectro homogéneo de un haz de álgebras graduadas & 49\\
\hyperref[subsection:II.3.1-fr]{3.1} & Espectro homogéneo de una $\mathcal{O}_Y$-álgebra graduada cuasicoherente & 49\\
\hyperref[subsection:II.3.2-fr]{3.2} & Haz sobre $\Proj(\mathcal{S})$ asociado a un $\mathcal{S}$-módulo graduado & 54\\
\hyperref[subsection:II.3.3-fr]{3.3} & $\mathcal{S}$-módulo graduado asociado a un haz sobre $\Proj(\mathcal{S})$ & 56\\
\hyperref[subsection:II.3.4-fr]{3.4} & Condiciones de finitud & 59\\
\hyperref[subsection:II.3.5-fr]{3.5} & Comportamiento funtorial & 61\\
\hyperref[subsection:II.3.6-fr]{3.6} & Subpreesquemas cerrados de un preesquema $\Proj(\mathcal{S})$ & 64\\
\hyperref[subsection:II.3.7-fr]{3.7} & Morfismos de un preesquema a un espectro homogéneo & 65\\
\hyperref[subsection:II.3.8-fr]{3.8} & Criterios para la inmersión en un espectro homogéneo & 69\\
\end{longtable}
\oldpage[II]{214}
\markboth{A. GROTHENDIECK --- Cap. II}{A. GROTHENDIECK --- Cap. II}
\begin{longtable}{@{}p{0.09\textwidth}p{0.69\textwidth}r@{}}
\hyperref[section:II.4-fr]{\textsection4.} & Fibrados proyectivos; haces amplios & 71\\
\hyperref[subsection:II.4.1-fr]{4.1} & Definición de los fibrados proyectivos & 71\\
\hyperref[subsection:II.4.2-fr]{4.2} & Morfismos de un preesquema a un fibrado proyectivo & 72\\
\hyperref[subsection:II.4.3-fr]{4.3} & El morfismo de Segre & 76\\
\hyperref[subsection:II.4.4-fr]{4.4} & Inmersiones en fibrados proyectivos; haces muy amplios & 78\\
\hyperref[subsection:II.4.5-fr]{4.5} & Haces amplios & 83\\
\hyperref[II.4.6-fr]{4.6} & Haces relativamente amplios & 89\\
\hyperref[section:II.5-fr]{\textsection5.} & Morfismos cuasiafines; morfismos cuasiproyectivos; morfismos propios; morfismos proyectivos & 94\\
\hyperref[subsection:II.5.1-fr]{5.1} & Morfismos cuasiafines & 94\\
\hyperref[subsection:II.5.2-fr]{5.2} & El criterio de Serre & 97\\
\hyperref[subsection:II.5.3-fr]{5.3} & Morfismos cuasiproyectivos & 99\\
\hyperref[subsection:II.5.4-fr]{5.4} & Morfismos propios y morfismos universalmente cerrados & 100\\
\hyperref[subsection:II.5.5-fr]{5.5} & Morfismos proyectivos & 103\\
\hyperref[subsection:II.5.6-fr]{5.6} & Lema de Chow & 106\\
\hyperref[section:II.6-fr]{\textsection6.} & Morfismos enteros y morfismos finitos & 110\\
\hyperref[subsection:II.6.1-fr]{6.1} & Preesquemas enteros sobre otro preesquema & 110\\
\hyperref[subsection:II.6.2-fr]{6.2} & Morfismos cuasifinitos & 114\\
\hyperref[subsection:II.6.3-fr]{6.3} & Clausura integral de un preesquema & 116\\
\hyperref[subsection:II.6.4-fr]{6.4} & Determinante de un endomorfismo de un $\mathcal{O}_X$-módulo & 120\\
\hyperref[subsection:II.6.5-fr]{6.5} & Norma de un haz invertible & 125\\
\hyperref[subsection:II.6.6-fr]{6.6} & Aplicación: criterios de amplitud & 130\\
\hyperref[subsection:II.6.7-fr]{6.7} & Teorema de Chevalley & 135\\
\hyperref[section:II.7-fr]{\textsection7.} & Criterios valuativos & 138\\
\hyperref[subsection:II.7.1-fr]{7.1} & Recordatorios sobre los anillos de valuación & 138\\
\hyperref[subsection:II.7.2-fr]{7.2} & Criterio valuativo de separación & 141\\
\hyperref[subsection:II.7.3-fr]{7.3} & Criterio valuativo de la propiedad de ser propio & 143\\
\hyperref[subsection:II.7.4-fr]{7.4} & Curvas algebraicas y cuerpos de funciones de dimensión uno & 148\\
\hyperref[section:II.8-fr]{\textsection8.} & Esquemas obtenidos por explosión; conos proyectantes; clausura proyectiva & 152\\
\hyperref[subsection:II.8.1-fr]{8.1} & Preesquemas obtenidos por explosión & 152\\
\hyperref[subsection:II.8.2-fr]{8.2} & Resultados preliminares sobre la localización de anillos graduados & 157\\
\hyperref[subsection:II.8.3-fr]{8.3} & Conos proyectantes & 162\\
\hyperref[subsection:II.8.4-fr]{8.4} & Clausura proyectiva de un fibrado vectorial & 168\\
\hyperref[subsection:II.8.5-fr]{8.5} & Comportamiento funtorial & 169\\
\hyperref[subsection:II.8.6-fr]{8.6} & Un isomorfismo canónico para los conos sin vértice & 171\\
\end{longtable}
\oldpage[II]{215}
\markboth{ÍNDICE}{ÍNDICE}
\begin{longtable}{@{}p{0.09\textwidth}p{0.69\textwidth}r@{}}
\hyperref[subsection:II.8.7-fr]{8.7} & Explosión de conos proyectantes & 173\\
\hyperref[subsection:II.8.8-fr]{8.8} & Haces amplios y contracciones & 177\\
\hyperref[subsection:II.8.9-fr]{8.9} & El criterio de amplitud de Grauert: enunciado & 182\\
\hyperref[subsection:II.8.10-fr]{8.10.} & El criterio de amplitud de Grauert: demostración & 184\\
\hyperref[subsection:II.8.11-fr]{8.11.} & Unicidad de las contracciones & 189\\
\hyperref[subsection:II.8.12-fr]{8.12.} & Haces cuasicoherentes sobre los conos proyectantes & 191\\
\hyperref[subsection:II.8.13-fr]{8.13.} & Clausura proyectiva de subhaces y subesquemas cerrados & 195\\
\hyperref[subsection:II.8.14-fr]{8.14.} & Complementos sobre los haces asociados a $\mathcal{S}$-módulos graduados & 197\\
\multicolumn{2}{@{}l}{Bibliografía (continuación)} & 205\\
\multicolumn{2}{@{}l}{Índice de notaciones} & 207\\
\multicolumn{2}{@{}l}{Índice terminológico} & 209\\
\multicolumn{2}{@{}l}{Erratas y adiciones (Lista 1)} & 217\\
\end{longtable}

\begin{flushright}
\emph{Recibido el 22 de julio de 1960.}
\end{flushright}
\endgroup
\oldpage[II]{216}
\oldpage[II]{217}
\phantomsection
\addcontentsline{toc}{section}{Erratas y adiciones (Lista 1)}
\section*{Erratas y adiciones}
\begin{center}
\large (Lista 1)
\end{center}

\subsection*{Capítulo 0}

\paragraph{(0, 2.1.3)}
En la línea 15 contando desde abajo de la p.~21, reemplazar
«inter-sección» por «intersección».

\paragraph{(0, 2.1.6)}
En la línea 13 de la p.~22, reemplazar
$Z_i\cap\complement(\bigcup\limits_{j\ne i}Z_j)$ par
$\bigcup\limits_{j\ne i}Z_j$.

\paragraph{(0, 3.1.6)}
En la línea 17 contando desde abajo de la p.~25, reemplazar
$\Gamma(U,\sh{F})$ por $\Gamma(V,\sh{F})$.

\paragraph{(0, 4.1.1)}
En la línea 12 de la p.~36, reemplazar $\sh{G}\to\sh{F}$ por
$\sh{B}\to\sh{A}$.

\paragraph{(0, 4.2.1)}
En la línea 11 contando desde abajo de la p.~39, reemplazar «haces» por
«haz».

\paragraph{(0, 5.1.2)}
Después de la línea 14 de la p.~45, añadir:

Si $\sh{G}$ es un $\sh{O}_Y$-módulo generado por sus secciones sobre
$Y$, entonces $f^*(\sh{G})$ está generado por sus secciones sobre $X$,
pues $f^*$ es un funtor exacto por la derecha.

\paragraph{(0, 5.2.4)}
En la línea 18 contando desde abajo de la p.~46, reemplazar $f_U$ por
$f_U^*$.

\paragraph{(0, 5.3.4)}
Antes de la línea 16 contando desde abajo de la p.~47, añadir:

Si
$\sh{A}\to\sh{B}\to\sh{C}\to\sh{D}\to\sh{E}$
es una sucesión exacta de $\sh{O}_X$-módulos, y si
$\sh{A},\sh{B},\sh{D},\sh{E}$ son coherentes, entonces $\sh{C}$ es
coherente.

\paragraph{(0, 5.4.1)}
Después de la línea 2 de la p.~49, añadir:

Supongamos que $\sh{O}_X$ es \emph{coherente}, y sea $\sh{F}$ un
$\sh{O}_X$-módulo \emph{coherente}. Si, en un punto $x\in X$, $\sh{F}_x$ es
un $\sh{O}_x$-módulo \emph{libre de rango $n$}, existe una vecindad $U$ de $x$
tal que $\sh{F}|U$ es \emph{localmente libre de rango $n$}; en efecto,
$\sh{F}_x$ es entonces isomorfo a $\sh{O}_x^n$, y la afirmación se
deduce de (5.2.7).

\paragraph{(0, 6.6.2)}
Reemplazar las líneas 10 y 11 de la p.~59 por:

En efecto, esto se deduce de (6.4.1, \emph{d)}).

\paragraph{(0, 6.7.1)}
Antes de la línea 10 contando desde abajo de la p.~59, añadir:

Si $f$ es un morfismo plano, también se dice que $X$ es
\emph{plano sobre $Y$}, o \emph{$Y$-plano}.

\paragraph{(0, 7.2.9)}
En la línea 11 contando desde abajo de la p.~65, reemplazar $M_n$ por
$M_{n-1}$.

\subsection*{Capítulo I}

\paragraph{(I, 1.2.1)}
En la línea 18 de la p.~83, reemplazar $\boldsymbol{K}$ por $\boldsymbol{k}$ (dos veces).

\paragraph{(I, 1.2.7)}
En la línea 16 contando desde abajo de la p.~84, reemplazar $A$ por
$A'$ (dos veces) y $\mathfrak N$ por $\mathfrak N'$; en la línea 17 contando desde
abajo, reemplazar $X$ por $X'$.

\paragraph{(I, 1.7.3) y (I, 2.2.4)}
Los enunciados dados bajo estos números pueden generalizarse, de
acuerdo con una observación de J.~Tate, como sigue.

\subsection*{1.8. Morfismos de espacios localmente anillados a esquemas afines}
\label{ega2:addendum:subsection:I.1.8-fr}

\begin{proposition}[1.8.1]
\label{ega2:addendum:I.1.8.1-fr}
Sean $(S,\sh{O}_S)$ un esquema afín y $(X,\sh{O}_X)$ un espacio
localmente anillado. Existe una biyección canónica del conjunto de
homomorfismos de anillos
\oldpage[II]{218}
$\Gamma(S,\sh{O}_S)\to\Gamma(X,\sh{O}_X)$
al conjunto de morfismos de espacios anillados
$(\psi,\theta):(X,\sh{O}_X)\to(S,\sh{O}_S)$ tales que, para todo
$x\in X$, $\theta_x^\sharp$ sea un homomorfismo local: $\sh{O}_{\psi(x)}\to\sh{O}_x$.
\end{proposition}

\begin{proof}
Observemos primero que, si $(X,\sh{O}_X)$ y $(S,\sh{O}_S)$ son dos
espacios anillados cualesquiera, un morfismo $(\psi,\theta)$ de
$(X,\sh{O}_X)$ a $(S,\sh{O}_S)$ define canónicamente un homomorfismo
de anillos
$\Gamma(\theta):\Gamma(S,\sh{O}_S)\to\Gamma(X,\sh{O}_X)$;
y, por tanto, una primera aplicación
\begin{equation}
\label{ega2:addendum:I.1.8.1.1-fr}
  \rho:\Hom((X,\sh{O}_X),(S,\sh{O}_S))
  \to\Hom(\Gamma(S,\sh{O}_S),\Gamma(X,\sh{O}_X)).
\tag{1.8.1.1}
\end{equation}
Recíprocamente, bajo las hipótesis de la proposición, pongamos
$A=\Gamma(S,\sh{O}_S)$ y consideremos un homomorfismo de anillos
$\vphi:A\to\Gamma(X,\sh{O}_X)$. Para todo $x\in X$, el conjunto de
los $f\in A$ tales que $\vphi(f)(x)=0$ es un \emph{ideal primo} de $A$, pues
$\sh{O}_x/\mathfrak m_x=\boldsymbol{k}(x)$ es un cuerpo. Es, por tanto, un
elemento de $S=\Spec(A)$, que denotaremos asimismo por
${}^a\vphi(x)$. Además, para todo $f\in A$, por definición (\textbf 0, 5.5.2),
${}^a\vphi^{-1}(D(f))=X_f$,
lo que prueba que ${}^a\vphi$ es una aplicación \emph{continua} $X\to S$.

Definamos después un homomorfismo
\[
  \widetilde{\vphi}:\sh{O}_S\to{}^a\vphi_*(\sh{O}_X)
\] de $\sh{O}_S$-módulos; para todo $f\in A$, se tiene
$\Gamma(D(f),\sh{O}_S)=A_f$ (1.3.6); para todo $s\in A$, se hace corresponder a $s/f\in A_f$ el elemento
$(\vphi(s)|X_f)(\vphi(f)|X_f)^{-1}$ de $\Gamma(X_f,\sh{O}_X)=\Gamma(D(f),{}^a\vphi_*(\sh{O}_X))$,
y se comprueba inmediatamente (pasando de $D(f)$ a $D(fg)$) que esto define un homomorfismo de $\sh{O}_S$-módulos, de donde, un morfismo de espacios
anillados $({}^a\vphi,\widetilde{\vphi})$. Con la misma notación, y
escribiendo $y={}^a\vphi(x)$, se tiene (\textbf 0, 3.7.1)
$\widetilde{\vphi}_x^\sharp(s_y/f_y)=(\vphi(s)_x)(\vphi(f)_x)^{-1}$;
Como $s_y\in\mathfrak m_y$ equivale, por definición, a
$\vphi(s)_x\in\mathfrak m_x$, se ve que $\widetilde{\vphi}_x^\sharp$ es un homomorfismo \emph{local} $\sh{O}_y\to\sh{O}_x$. Se ha
definido así una segunda aplicación
$\sigma:\Hom(\Gamma(S,\sh{O}_S),\Gamma(X,\sh{O}_X))\to\mathfrak L$,
donde $\mathfrak L$ denota el conjunto de morfismos
$(\psi,\theta):(X,\sh{O}_X)\to(S,\sh{O}_S)$ para los cuales
$\theta_x^\sharp$ es local para todo $x\in X$.

Queda probar que $\sigma$ y la restricción de $\rho$ a $\mathfrak L$
son inversas. La definición de $\widetilde{\vphi}$ da inmediatamente
$\Gamma(\widetilde{\vphi})=\vphi$, y por tanto $\rho\circ\sigma$ es la
identidad. Para ver que $\sigma\circ\rho$ es la identidad, partamos de
$(\psi,\theta)\in\mathfrak L$ y sea $\vphi=\Gamma(\theta)$. La
hipótesis sobre $\theta_x^\sharp$ proporciona, al pasar a los cuerpos
residuales, un monomorfismo
$\theta^x:\boldsymbol{k}(\psi(x))\to\boldsymbol{k}(x)$
tal que, para todo $f\in A=\Gamma(S,\sh{O}_S)$,
$\theta^x(f(\psi(x)))=\vphi(f)(x)$;
Así, $f(\psi(x))=0$ equivale a $\vphi(f)(x)=0$, lo que prueba que
${}^a\vphi=\psi$. Por otra parte, las definiciones hacen conmutativo
el diagrama
\[
  \xymatrix{
    A\ar[r]^\vphi\ar[d] &
    \Gamma(X,\sh{O}_X)\ar[d]\\
    A_{\psi(x)}\ar[r]_{\theta_x^\sharp} &
    \sh{O}_x
  }
\]
y también es conmutativo el diagrama análogo en el que
$\theta_x^\sharp$ se reemplaza por $\widetilde{\vphi}_x^\sharp$. De
(\textbf 0, 1.2.4) se sigue que
$\widetilde{\vphi}_x^\sharp=\theta_x^\sharp$, y, por consiguiente,
$\widetilde{\vphi}=\theta$.
\end{proof}

\begin{env}[1.8.2]
\label{ega2:addendum:I.1.8.2-fr}
Cuando $(X,\sh{O}_X)$ e $(Y,\sh{O}_Y)$ son espacios \emph{localmente}
anillados, consideraremos los morfismos
$(\psi,\theta):(X,\sh{O}_X)\to(Y,\sh{O}_Y)$ tales que, para todo
$x\in X$, $\theta_x^\sharp$ sea un homomorfismo \emph{local}: $\sh{O}_{\psi(x)}\to\sh{O}_x$.

Siempre que en lo sucesivo hablemos de un
\oldpage[II]{219}
\emph{morfismo de espacios localmente anillados}, entenderemos siempre
un morfismo de este tipo. Con esta definición de los morfismos, los
espacios localmente anillados forman una \emph{categoría}. Para dos objetos
$X,Y$ de esta categoría, $\Hom(X,Y)$ denotará, por tanto, el conjunto
de morfismos de espacios localmente anillados de $X$ a $Y$ (el
conjunto denotado por $\mathfrak L$ en \sref{ega2:addendum:I.1.8.1-fr}). Cuando queramos
considerar el conjunto de \emph{morfismos de espacios anillados} de $X$ a
$Y$, escribiremos $\Hom_{\mathit{an}}(X,Y)$ para evitar toda
ambigüedad. La aplicación \sref{ega2:addendum:I.1.8.1.1-fr} se escribe entonces
\begin{equation}
\label{ega2:addendum:I.1.8.2.1-fr}
  \rho:\Hom_{\mathit{an}}(X,Y)
  \to\Hom(\Gamma(Y,\sh{O}_Y),\Gamma(X,\sh{O}_X))
\tag{1.8.2.1}
\end{equation}
y su restricción
\begin{equation}
\label{ega2:addendum:I.1.8.2.2-fr}
  \rho':\Hom(X,Y)
  \to\Hom(\Gamma(Y,\sh{O}_Y),\Gamma(X,\sh{O}_X))
\tag{1.8.2.2}
\end{equation}
es \emph{funtorial} en $X$ e $Y$ en la categoría de los espacios localmente
anillados.
\end{env}

\begin{corollary}[1.8.3]
\label{ega2:addendum:I.1.8.3-fr}
Sea $(Y,\sh{O}_Y)$ un espacio localmente anillado. Para que $Y$ sea un
esquema afín, es necesario y suficiente que, para todo espacio
localmente anillado $(X,\sh{O}_X)$, la aplicación \sref{ega2:addendum:I.1.8.2.2-fr} sea
biyectiva.
\end{corollary}

\begin{proof}
La Proposición~\sref{ega2:addendum:I.1.8.1-fr} muestra que la condición es necesaria.
Recíprocamente, se supone que se cumple y se pone
$A=\Gamma(Y,\sh{O}_Y)$. La hipótesis y \sref{ega2:addendum:I.1.8.1-fr} muestran que los
funtores
$X\to\Hom(X,Y)$ y $X\to\Hom(X,\Spec(A))$
de los espacios localmente anillados a los conjuntos son \emph{isomorfos}.
Esto implica la existencia de un isomorfismo canónico
$Y\to\Spec(A)$ (cf.~\textbf 0, 8).
\end{proof}

\begin{env}[1.8.4]
\label{ega2:addendum:I.1.8.4-fr}
Sea $S=\Spec(A)$ un esquema afín. Denotemos por $(S',A')$ el espacio
anillado cuyo espacio subyacente \emph{consta de un punto} y cuyo haz
estructural $A'$ es el haz (necesariamente simple) sobre $S'$ definido
por el anillo $A$. Sea $\pi:S\to S'$ la única aplicación de $S$ en $S'$; observemos, por otra parte, que para todo
abierto $U$ de $S$ se tiene una aplicación canónica
$\Gamma(S',A')=\Gamma(S,\sh{O}_S)\to\Gamma(U,\sh{O}_S)$
que define un \emph{$\pi$-morfismo} $\iota:A'\to\sh{O}_S$ de haces de anillos.
Se obtiene así un \emph{morfismo} canónico \emph{de espacios anillados}
$i=(\pi,\iota):(S,\sh{O}_S)\to(S',A')$.
Para todo $A$-módulo $M$, denotemos por $M'$ el haz simple sobre $S'$
definido por $M$, que es evidentemente un $A'$-módulo. Está claro que
$i_*(\widetilde M)=M'$ (1.3.7).
\end{env}

\begin{lemma}[1.8.5]
\label{ega2:addendum:I.1.8.5-fr}
Con la notación de \sref{ega2:addendum:I.1.8.4-fr}, para todo $A$-módulo $M$, el
$\sh{O}_S$-homomorfismo canónico y funtorial \textup{(\textbf{0}, 4.4.3.3)}
\begin{equation}
\label{ega2:addendum:I.1.8.5.1-fr}
  i^*(i_*(\widetilde M))\to\widetilde M
\tag{1.8.5.1}
\end{equation}
es un isomorfismo.
\end{lemma}

\begin{proof}
Ambos miembros de \sref{ega2:addendum:I.1.8.5.1-fr} son exactos por la derecha
(siendo evidentemente exacto el funtor
$M\to i_*(\widetilde M)$) y conmutan con las sumas directas.
Considerando $M$ como el conúcleo de un homomorfismo
$A^{(I)}\to A^{(J)}$, se reduce uno al caso $M=A$, en el cual la
afirmación es inmediata.
\end{proof}

\begin{corollary}[1.8.6]
\label{ega2:addendum:I.1.8.6-fr}
Sea $(X,\sh{O}_X)$ un espacio anillado y sea $u:X\to S$ un morfismo de
\oldpage[II]{220}
espacios anillados. Para todo $A$-módulo $M$, existe, con la notación
de \sref{ega2:addendum:I.1.8.4-fr}, un isomorfismo canónico y funtorial de
$\sh{O}_X$-módulos
\begin{equation}
\label{ega2:addendum:I.1.8.6.1-fr}
  u^*(\widetilde M)\xrightarrow{\sim} u^*(i^*(M'))
\tag{1.8.6.1}
\end{equation}
\end{corollary}

\begin{corollary}[1.8.7]
\label{ega2:addendum:I.1.8.7-fr}
Bajo las hipótesis de \sref{ega2:addendum:I.1.8.6-fr}, para todo $A$-módulo $M$ y todo
$\sh{O}_X$-módulo $\sh{F}$ existe un isomorfismo canónico, funtorial
en $M$ y $\sh{F}$,
\begin{equation}
\label{ega2:addendum:I.1.8.7.1-fr}
  \Hom_{\sh{O}_S}(\widetilde M,u_*(\sh{F}))
  \xrightarrow{\sim}\Hom_A(M,\Gamma(X,\sh{F})).
\tag{1.8.7.1}
\end{equation}
\end{corollary}

\begin{proof}
Por (\textbf 0, 4.4.3) y el Lema~\sref{ega2:addendum:I.1.8.5-fr}, existe un isomorfismo
canónico de bifuntores
\[
  \Hom_{\sh{O}_S}(\widetilde M,u_*(\sh{F}))
  \xrightarrow{\sim}\Hom_{A'}(M',i_*(u_*(\sh{F})))
\]
El segundo miembro es claramente $\Hom_A(M,\Gamma(X,\sh{F}))$. Bajo
el homomorfismo canónico \sref{ega2:addendum:I.1.8.7.1-fr}, a un
$\sh{O}_S$-homomorfismo
$h:\widetilde M\to u_*(\sh{F})$ (equivalentemente, a un $u$-morfismo
$\widetilde M\to\sh{F}$) le corresponde el $A$-homomorfismo $\Gamma(h):M\to\Gamma(S,u_*(\sh{F}))=\Gamma(X,\sh{F})$.
\end{proof}

\begin{env}[1.8.8]
\label{ega2:addendum:I.1.8.8-fr}
Con la notación de \sref{ega2:addendum:I.1.8.4-fr}, se deduce claramente de
(\textbf 0, 4.1.1) que todo morfismo de espacios anillados
$(\psi,\theta):X\to S'$ equivale al dato de un homomorfismo de
anillos $A\to\Gamma(X,\sh{O}_X)$. Por tanto, la
Proposición~\sref{ega2:addendum:I.1.8.1-fr} puede interpretarse como una biyección
canónica
$\Hom(X,S)\xrightarrow{\sim}\Hom(X,S')$
donde en el segundo miembro se entienden morfismos de espacios
anillados, pues $A$ no es en general un anillo local.

Más generalmente, sean $X,Y$ espacios localmente anillados, y sea
$(Y',A')$ el espacio anillado cuyo espacio subyacente consta de un
punto y cuyo haz de anillos $A'$ es el haz simple definido por
$\Gamma(Y,\sh{O}_Y)$. Entonces \sref{ega2:addendum:I.1.8.2.1-fr} puede interpretarse
como una aplicación
\begin{equation}
\label{ega2:addendum:I.1.8.8.1-fr}
  \rho:\Hom_{\mathit{an}}(X,Y)\to\Hom(X,Y')
\tag{1.8.8.1}
\end{equation}
El Corolario~\sref{ega2:addendum:I.1.8.3-fr} dice que los esquemas afines se
caracterizan, entre los espacios localmente anillados, como aquellos
para los cuales la restricción de $\rho$ a $\Hom(X,Y)$,
\begin{equation}
\label{ega2:addendum:I.1.8.8.2-fr}
  \rho':\Hom(X,Y)\to\Hom(X,Y')
\tag{1.8.8.2}
\end{equation}
es \emph{biyectiva} para \emph{todo} espacio localmente anillado $X$. En un capítulo
posterior generalizaremos esta definición, asociando así a un espacio
anillado \emph{arbitrario} $Z$ (y no solamente a uno cuyo espacio subyacente
consta de un punto) un espacio localmente anillado que volveremos a
denotar por $\Spec(Z)$. Este será el punto de partida de una teoría
«relativa» de los preesquemas sobre espacios anillados arbitrarios,
que extiende los resultados del Capítulo~I.
\end{env}

\begin{env}[1.8.9]
\label{ega2:addendum:I.1.8.9-fr}
Se pueden considerar los pares $(X,\sh{F})$ formados por un espacio localmente anillado $X$ y un $\sh{O}_X$-módulo $\sh{F}$ como formando una categoría, siendo un \emph{morfismo} $(X,\sh{F})\to(Y,\sh{G})$ de esta categoría un par $(u,h)$ formado por un morfismo de
\oldpage[II]{221}
espacios localmente anillados $u:X\to Y$ y de un $u$-morfismo
$h:\sh{G}\to\sh{F}$ de módulos; estos
morfismos (para $(X,\sh{F})$ e $(Y,\sh{G})$ fijados) forman un conjunto
que denotaremos por $\Hom((X,\sh{F}),(Y,\sh{G}))$; la aplicación
$(u,h)\to(\rho'(u),\Gamma(h))$ es una aplicación canónica
\begin{equation}
\label{ega2:addendum:I.1.8.9.1-fr}
\begin{split}
  \Hom((X,\sh{F}),(Y,\sh{G}))\to
  \Hom\bigl(
    &(\Gamma(Y,\sh{O}_Y),\Gamma(Y,\sh{G})),\\
    &(\Gamma(X,\sh{O}_X),\Gamma(X,\sh{F}))
  \bigr)
\end{split}
\tag{1.8.9.1}
\end{equation}
\emph{funtorial} en $(X,\sh{F})$ y $(Y,\sh{G})$, siendo el segundo miembro
el conjunto de los dihomomorfismos correspondientes a los anillos y a
los módulos considerados (\textbf 0, 1.0.2).
\end{env}

\begin{corollary}[1.8.10]
\label{ega2:addendum:I.1.8.10-fr}
Sean $Y$ un espacio localmente anillado y $\sh{G}$ un
$\sh{O}_Y$-módulo. Para que $Y$ sea un esquema afín y $\sh{G}$ un
$\sh{O}_Y$-módulo cuasicoherente, es necesario y suficiente que, para
toda pareja $(X,\sh{F})$ formada por un espacio localmente anillado $X$
y un $\sh{O}_X$-módulo $\sh{F}$, la aplicación canónica
\sref{ega2:addendum:I.1.8.9.1-fr} sea biyectiva.
\end{corollary}

Dejamos al lector la tarea de desarrollar el razonamiento, modelado
sobre el de \sref{ega2:addendum:I.1.8.3-fr}, en el cual se utilizan \sref{ega2:addendum:I.1.8.1-fr} y
\sref{ega2:addendum:I.1.8.7-fr}.

\begin{remark}[1.8.11]
\label{ega2:addendum:I.1.8.11-fr}
Los enunciados (1.7.3), (1.7.4) y (2.2.4), así como la
definición (1.6.1), son casos particulares de \sref{ega2:addendum:I.1.8.1-fr}; del
mismo modo, (2.2.5) se deduce de \sref{ega2:addendum:I.1.8.7-fr}. El
Corolario~\sref{ega2:addendum:I.1.8.7-fr} proporciona también (1.6.3) (y, por tanto,
(1.6.4)) como caso particular. En efecto, si $X$ es afín y
$\Gamma(X,\sh{F})=N$, los funtores
$M\to
  \Hom_{\sh{O}_S}(\widetilde M,u_*(\widetilde N))$ y $M\to
  \Hom_{\sh{O}_S}(\widetilde M,(N_{[\vphi]})\supertilde)$
(donde $\vphi:A\to\Gamma(X,\sh{O}_X)$ corresponde a $u$) son
isomorfos por \sref{ega2:addendum:I.1.8.7-fr} y (1.3.8). Finalmente, (1.6.5) (y,
por tanto, (1.6.6)) se deduce de \sref{ega2:addendum:I.1.8.6-fr} y del hecho de que,
para todo $f\in A$, los $A_f$-módulos
$N'\otimes_{A'}A_f$ y $(N'\otimes_{A'}A)_f$
(con las notaciones de (1.6.5)) son canónicamente isomorfos.
\end{remark}

\paragraph{Insértese después de (I, 3.2.8):}

\begin{env}[3.2.9]
\label{ega2:addendum:I.3.2.9-fr}
De \sref{ega2:addendum:I.1.8.1-fr} se deduce que (3.2.2) puede precisarse del modo
siguiente. El preesquema
$Z=\Spec(B\otimes_A C)$
no solo es un producto de $X=\Spec(B)$ e $Y=\Spec(C)$ en la categoría
de los $S$-\emph{preesquemas}, sino también en la categoría de los \emph{espacios
localmente anillados} \emph{sobre} $S$ (con la definición de los
$S$-morfismos modelada sobre (2.5.2)). La demostración de (3.2.6)
muestra asimismo que, para cualesquiera $S$-preesquemas $X,Y$, el
preesquema $X\times_S Y$ no solo es el producto de $X$ y $Y$ en la categoría de los
$S$-preesquemas, sino también en la categoría de los espacios
localmente anillados sobre el preesquema $S$.
\end{env}

\paragraph{(I, 4.4.5)}
En la línea 14 contando desde abajo de la p.~126, sustitúyase $B$ por
$A$; en la línea 13 contando desde abajo, sustitúyase «álgebra sobre
$A$» por «álgebra sobre $B$».

\paragraph{(I, 5.3.5)}
En la línea 2 contando desde abajo de la p.~132, sustitúyase $Y\to S$
por $g:Y\to S$.

\paragraph{(I, 6.3.2.1)}
En la línea 6 contando desde abajo de la p.~144, léase:
$D(g_i)\subset W$.

\paragraph{(I, 7.3.8)}
Sustitúyase el texto actual, que es erróneo, por el siguiente:

Sean $X,Y$ dos preesquemas \emph{íntegros}; entonces $\sh{R}(X)$ y
$\sh{R}(Y)$ son, respectivamente, un $\sh{O}_X$-módulo
cuasicoherente y un $\sh{O}_Y$-módulo cuasicoherente (7.3.3). Si
$f:X\to Y$ es un morfismo \emph{dominante}, existe un homomorfismo canónico
de $\sh{O}_X$-módulos
\oldpage[II]{222}
\begin{equation}
\label{ega2:addendum:I.7.3.8.1-fr}
  \tau:f^*(\sh{R}(Y))\to\sh{R}(X).
\tag{7.3.8.1}
\end{equation}

Supongamos primero que $X=\Spec(A)$ e $Y=\Spec(B)$ son afines, con
$A$ y $B$ anillos íntegros. Entonces $f$ corresponde a un homomorfismo
inyectivo $B\to A$, que se prolonga en un monomorfismo $L\to K$ del
cuerpo de fracciones $L$ de $B$ en el cuerpo de fracciones $K$ de $A$.
El homomorfismo \sref{ega2:addendum:I.7.3.8.1-fr} corresponde entonces al homomorfismo
canónico
$L\otimes_B A\to K$
de (1.6.5).

En el caso general, para toda pareja de abiertos afines no vacíos
$U\subset X$, $V\subset Y$ tales que $f(U)\subset V$, la construcción
anterior define un homomorfismo $\tau_{U,V}$. Si $U'\subset U$,
$V'\subset V$ y $f(U')\subset V'$, entonces $\tau_{U,V}$ prolonga a
$\tau_{U',V'}$; la afirmación se obtiene por pegado. Si $x,y$ son los
puntos genéricos de $X$ y $Y$, respectivamente, entonces $f(x)=y$ y
\[
  \bigl(f^*(\sh{R}(Y))\bigr)_x
  =\sh{O}_y\otimes_{\sh{O}_y}\sh{O}_x
  =\sh{O}_x
\]
por (\textbf 0, 4.3.1) de modo que $\tau_x$ es un \emph{isomorfismo}.

\paragraph{(I, 9.5.2)}
En las tres últimas líneas de la p.~176 y las cinco primeras de la
p.~177, sustitúyase en todas partes $X'$ por $Y'$ y $X''$ por $Y''$.

\paragraph{(I, 10.5.4)}
En la línea 11 contando desde abajo de la p.~187, sustitúyase (10.3.4)
por (10.3.5).

\paragraph{(I, 10.6.3)}
En la línea 3 contando desde abajo de la p.~189, sustitúyase $X$ por
$\mathfrak X$.

\paragraph{(I, 10.14.2)}
En la línea 5 de la p.~210, sustitúyase «$\sh{O}_{\mathfrak X}$-\emph{módulo
coherente}» por «\emph{ideal coherente de} $\sh{O}_{\mathfrak X}$»; en la
línea 7, sustitúyase «$\sh{O}_{\mathfrak X}$-módulos coherentes» por
«ideales coherentes de $\sh{O}_{\mathfrak X}$».
\end{document}
